% ============ \emph{类域论}中文翻译 LaTeX 重排版 ============
% 原书: J.S. Milne, Class Field Theory (v4.03, August 6, 2020)
% 模板: AJbook.cls (李文威\emph{代数学方法}书系模板)
% 编号体系 (按原书):
%   - 引言语为第 0 章, 定理编号 0.1, 0.4, 0.5 ...
%   - 正文 I--VIII 章(Roman 数字), 定理/推论/引理/命题/定义/例/注记
%     共享一个计数器, 按节编号(如 命题 1.1, 定理 3.10)
% 结构: 封面 + 目录 + 记号/预备/参考文献/致谢/DRAMATIS PERSONAE
%       + 引言(第 0 章) + I--VIII 章 + 索引
%
\documentclass[geometry=a4, fontsetup=font-setup-open.tex, titlesetup=titles-setup, coverpage=cover, CJKthechapter=false]{AJbook}

% ============ 宏包 ============
\usepackage{mycommand}      % 常用数学命令(\R \Z \Q \CC \e \ii \Spec \Gm \Gal \Nm \Cl ...)
\usepackage{myarrows}       % 可伸缩箭头
\usepackage{amssymb}
\usepackage{enumerate}      % 支持 \begin{enumerate}[(a)] 形式的编号列表
\usepackage{chngcntr}       % \counterwithin
\usepackage{quiver}         % 交换图: 包装 tikz-cd, 并载入 amssymb/calc/decorations.pathmorphing/spath3,
                            % 额外提供 curve / between / 2tail 等样式, 使 q.uiver.app 导出的代码可直接编译
\usepackage{imakeidx}       % 索引(xindy)
\usepackage[bookmarksnumbered=true]{hyperref}
\usepackage{xurl}          % 允许 URL 在任意字符处断行, 消除书目条目的 overfull hbox
\usetikzlibrary{calc}       % 封面坐标计算

% ============ 参考文献库 ============
\addbibresource{references.bib}

% ============ 跨章共用的算子与范畴简写 ============
% 用 \providecommand, 与各章内部的重复定义不冲突
\providecommand{\Res}{\operatorname{Res}}    % 限制映射
\providecommand{\Cor}{\operatorname{Cor}}    % 余限制 (转移) 映射
\providecommand{\Inf}{\operatorname{Inf}}    % 膨胀映射
\providecommand{\Ver}{\operatorname{Ver}}    % Verlagerung
\providecommand{\inv}{\operatorname{inv}}    % 不变映射
\providecommand{\Coker}{\operatorname{Coker}}
\providecommand{\Ab}{\mathsf{Ab}}            % 交换群范畴
\providecommand{\Mod}{\mathsf{Mod}}          % 模范畴

% ============ 编号体系 ============
% 章用 Roman 数字显示: I, II, ..., VIII
\renewcommand{\thechapter}{\Roman{chapter}}
% 节只显示节号(随章重置), 如 1, 2, 3
\renewcommand{\thesection}{\arabic{section}}
% 定理类环境按节编号: 命题 1.1, 定理 3.10 ...
% 注意: theorem 计数器由 AJbook.cls 在 \AtBeginDocument 中定义,
% 故 \counterwithin 必须也放在 \AtBeginDocument 里, 待计数器存在后再执行.
\AtBeginDocument{\counterwithin{theorem}{section}}

% 章标题框与目录中只显示 Roman 数字(不要"附录 I"或"第 I 章")
\renewcommand{\AJchapterttl}[1]{#1}

% 天眉: "I 局部类域论: Lubin-Tate 理论". AJbook.cls 在 \AtBeginDocument 里定义 \chaptermark, 故此处覆盖须放在其之后
\AtBeginDocument{%
  \renewcommand{\chaptermark}[1]{\markboth{\thechapter\quad #1}{}}%
}

% ============ 超链接 ============
\hypersetup{
    pdfauthor={J.S. Milne},
    pdftitle={类域论 (Class Field Theory)},
    pdfsubject={代数数论},
    colorlinks=true,
    linkcolor=MidnightBlue,
    citecolor=BrickRed,
    urlcolor=MidnightBlue
}

% ============ 本重排版专用命令 ============
% 注: \N, \Z, \Q, \CC, \Gal, \Nm, \Cl, \Frob 等已在 mycommand.sty 中定义, 此处不再重复.
% 但原书的记号与模板默认不同, 需覆盖(不动 mycommand.sty):
%   范数在原书排作 Nm(mycommand 默认 \Nm= N), Frobenius 元排作 Frob(默认 \Frob= Fr).
\renewcommand{\Nm}{\operatorname{Nm}}        % 范数 (原书 Nm)
\renewcommand{\Frob}{\operatorname{Frob}}    % Frobenius 元 (原书 Frob)
\newcommand{\Spl}{\operatorname{Spl}}                 % 分裂素理想集
\newcommand{\ray}{\mathfrak{m}}                       % 模(modulus)
\newcommand{\rayC}{\mathrm{C}_{\mathfrak{m}}}          % ray class group
\newcommand{\ordp}[1]{\operatorname{ord}_{#1}}         % 赋值
% 原书用 sx<编号> / mo<编号> 标注 StackExchange / MathOverflow 网上讨论, 此处还原为可点击链接
\newcommand{\sx}[1]{\href{https://math.stackexchange.com/questions/#1/}{Math\,StackExchange\,#1}}
\newcommand{\mo}[1]{\href{https://mathoverflow.net/questions/#1/}{MathOverflow\,#1}}

% ============ 索引 (xindy) ============
\makeindex[%
    columns=2,
    program=texindy,
    intoc=true,
    options=-M texindy -I xelatex -C utf8,
    title={名词索引暨英译}]

% 全局断行弹性: 消除边界性 Overfull
\emergencystretch=3em

% ============ 文档开始 ============
\begin{document}
\frontmatter        % 封面(cover.tex) + 目录(模板自动处理, 罗马页码)

% 前置部分: 记号/预备/参考文献/致谢/DRAMATIS PERSONAE(原书罗马页码)
% ============ 前置部分: 记号/预备/参考文献/致谢/DRAMATIS PERSONAE ============
% 记号、预备知识与参考文献按原书连续排印, 不各自另起章节或页面.

\begingroup
\setstretch{1.0}
\clearpage
\noindent\textbf{记号.}\index{jihao@记号(notation)}

我们采用标准的(Bourbaki)记号:

$\N = \{0, 1, 2, \ldots\}$,
$\Z$ 为整数环,
$\Q$ 为有理数域,
$\R$ 为实数域,
$\CC$ 为复数域,
$\mathbb{F}_{p} = \Z/p\Z$ 为有 $p$ 个元素的域, 其中 $p$ 为素数.

对于整数 $m$ 与 $n$, $m\mid n$ 表示 $m$ 整除 $n$, 即 $n \in m\Z$. 通篇约定 $p$ 为素数, 即 $p = 2, 3, 5, \ldots$.

给定一个等价关系, $[*]$ 表示包含 $*$ 的等价类. 空集记为 $\emptyset$. 集合 $S$ 的基数记为 $|S|$ (于是当 $S$ 有限时 $|S|$ 为 $S$ 中元素的个数). 设 $I$ 与 $A$ 为集合; 以 $I$ 为指标集的 $A$ 中元素族, 记为 $(a_{i})_{i \in I}$, 是映射 $i \mapsto a_{i} : I \to A$.

$X \subset Y$ 表示 $X$ 是 $Y$ 的子集(不必是真子集);

$X \overset{\mathrm{def}}{=} Y$ 表示 $X$ 定义为 $Y$, 或按定义等于 $Y$;

$X \approx Y$ 表示 $X$ 同构于 $Y$;

$X \simeq Y$ 表示 $X$ 与 $Y$ 典范同构(或存在给定的或唯一的同构);

$\hookrightarrow$ 表示单射;

$\twoheadrightarrow$ 表示满射.

我们用哥特体(fraktur)字母表示理想, 例如
$\mathfrak{a}, \mathfrak{b}, \mathfrak{c}, \mathfrak{m}, \mathfrak{n}, \mathfrak{p}, \mathfrak{q},
 \mathfrak{A}, \mathfrak{B}, \mathfrak{C}, \mathfrak{M}, \mathfrak{N}, \mathfrak{P}, \mathfrak{Q}$
分别对应
$a, b, c, m, n, p, q, A, B, C, M, N, P, Q$.

\par\medskip\noindent\textbf{预备知识.}\index{yubeizhishi@预备知识(prerequisites)}\enspace
需要通常研究生一年级课程所覆盖的代数, 以及代数数论的一门课程, 例如我下面列出的课程讲义.

\par\medskip\noindent\textbf{参考文献.}\index{cankaowenxian@参考文献(references)}\enspace
除书末所列(以及脚注中所示的)参考文献外, 我还将引用我自己如下的课程讲义:

GT Group Theory (v3.16, 2020).

FT Fields and Galois Theory (v4.61, 2020).

ANT Algebraic Number Theory (v3.08, 2020).

\par\medskip\noindent\textbf{致谢.}\index{zhixie@致谢(acknowledgements)}\enspace

我感谢以下人士为这些讲义的早期版本提供勘误与意见: Vincenzo Acciaro; Ra\'ul Alonso Rodr\'iguez; Tom Bachmann; Oliver Braeunling; Chen, Bingxu; Kwangho Choiy; Brian Conrad; Keith Conrad; Giuseppe Canuto; Philip Dittmann; Ross Griffiths; Darij Grinberg; Florian Herzig; Chong Hui; Herv\'e Jacquet; Kiran Kedlaya; Timo Keller; Keenan Kidwell; Andrew Kirk; Michiel Kosters; Tyler Lawson; Jungin Lee; Franz Lemmermeyer; Yogesh More; Kim Nguyen; Jianing Li; Catherine O'Neil; Jack Petok; Kartik Prasanna; Nandini Ranganathan; Peter Roquette; Joshua Seaton; Corinne Sheridan; Jonah Sinick; Daniel Sparks; Jamie Tappenden; Wojtek Wawr\'ow; Yu Zhao; 以及其他人. 同时我感谢 math.stackexchange.com 与 mathoverflow.net 的贡献者, 尤其是那些在我的讲义上提问或作答的人.

我一直在读 Chevalley 关于类域论的新书; 我并非真正在做研究, 只不过是在自我修养罢了.\par\noindent Grothendieck, 1956.
\endgroup

\clearpage
\noindent\textbf{DRAMATIS PERSONAE.}\index{dramatispersonae@DRAMATIS PERSONAE(dramatis personae)}\index{renwubiao@人物表(dramatis personae)}

FERMAT (1601--1665). 提出他最后的``定理'', 并对 $m = 4$ 给出了证明. 他还提出了求方程
\begin{equation*}
X^{2} - AY^{2} = 1, \quad A \in \Z
\tag{1}
\end{equation*}
整数解的问题, 这本质上就是求 $\Z[\sqrt{A}]$ 中单位的问题\footnote{印度数学家婆什迦罗 (Bhaskara, 12 世纪) 已经知道求该方程解的一般法则。}. Brouncker 给出了求解该问题的算法, 但没有证明算法总有效.

EULER (1707--1783). 他将分析引入素数研究, 并发现了二次互反律的早期形式\footnote{Euler 发现了 [二次互反律], 而且看来有可能用散见于 Euler 的《全集》(Opera omnia) 或其遗稿中的不同片段拼出该定理的一个证明. 给出第一个完整证明的是 Gauss.$\ldots$ (Michael Berg, MR2131680).}.

LAGRANGE (1736--1813). 他证明了求解 (1) 的算法总能给出解, 并证明了每个正整数都是四个平方数之和.

LEGENDRE (1752--1833). 他引入了``Legendre 符号'' $\left(\frac{m}{p}\right)$, 并找到了二次互反律的完整形式
\[
\left(\frac{p}{q}\right)\left(\frac{q}{p}\right) = (-1)^{(p-1)(q-1)/4}, \quad p, q\ \text{为奇素数}.
\]
他证明了一个蕴含如下三元二次型局部--整体原理的结果: 对 $\Q$ 上的三元二次型 $Q(X, Y, Z)$, 它在 $\Q$ 中有非平凡零点, 当且仅当它在 $\R$ 中有, 且同余式 $Q \equiv 0 \pmod{p^{n}}$ 对所有 $p$ 与 $n$ 都有非平凡解.

GAUSS (1777--1855). 他给出了二次互反律的首批完整证明. 为寻求四次互反律, 他研究了 Gauss 整数 $\Z[\mathrm{i}]$. 他研究了 $\Z$ 上二元二次型的分类, 这与求二次域的类数问题密切相关.

DIRICHLET (1805--1859). 他引入了 $L$-级数, 并用它们证明了类数的解析公式以及算术级数中素数的密度定理. 他证明了如下的``单位定理'': 设 $\alpha$ 是首一不可约整系数多项式 $f(X)$ 的根; 若 $f(X)$ 有 $r$ 个实根与 $2s$ 个复根, 则 $\Z[\alpha]^{\times}$ 是秩为 $r + s - 1$ 的有限生成群.

KUMMER (1810--1893). 受寻求高次互反律的推动, 他对分圆域的算术作了深入研究, 并指出借助``理想数''可以恢复唯一分解. 他证明了费马大定理对正则素数成立.

HERMITE (1822--1901). 他对二次型有重要贡献, 并证明了五次多项式的根可以用椭圆函数表示. EISENSTEIN (1823--1852). 他发表了三次与四次互反律的首批完整证明.

KRONECKER (1823--1891). 他发展了 Dedekind 理想理论的一种替代方案. 他还提出了数学中最优美的想法之一, 用以生成数域的 Abel 扩张(Kronecker 青春之梦).

RIEMANN (1826--1866). 研究了 Riemann $\zeta$ 函数, 并提出 Riemann 假设.

DEDEKIND (1831--1916). 他通过给出数域中整数环的正确定义, 证明了此类环中理想可唯一分解为素理想之积, 并证明了模去主理想后它们落入有限多个类, 从而奠定了代数数论的现代基础. 他定义了数域的 $\zeta$ 函数.

WEBER (1842--1913). 他找到了``类群''允许分歧的正确推广. 在类域论与 Kronecker 青春之梦方面取得重要进展. HENSEL (1861--1941). 他在 1890 年代首次给出了 $p$-进数域的定义(作为无穷和 $\sum_{n=-k}^{\infty} a_{n}p^{n}$, $a_{n} \in \{0,1,\ldots,p-1\}$ 的集合)\footnote{赋值论由匈牙利数学家 K\"ursch\'ak 于 1912 年创立。}.

HILBERT (1862--1943). 他在 1897 年写了一本极具影响力的代数数论著作, 给出了该理论的首次系统叙述. 他的一些著名问题涉及数论, 也影响深远.

TAKAGI (1875--1960). 他证明了 Abel 类域论的基本定理, 即 Weber 与 Hilbert 所猜想者.

NOETHER (1882--1935). 她与 Artin 一道奠定了现代代数基础, 强调公理与概念性论证, 并对数域上中心单代数的分类作出贡献.

HECKE (1887--1947). 引入了 Hecke $L$-级数, 推广了 Dirichlet 的 $L$-级数与 Dedekind 的 $\zeta$ 函数.

ARTIN (1898--1962). 他发现了``Artin 互反律'', 即类域论的主定理(对 Takagi 结果的改进). 引入了 Artin $L$-级数.

HASSE (1898--1979). 他给出了局部类域论的首次证明, 证明了数域上所有二次型的 Hasse(局部--整体)原理, 并对数域上中心单代数的分类作出贡献.

BRAUER (1901--1977). 定义了 Brauer 群, 并对数域上中心单代数的分类作出贡献.

WEIL (1906--1998). 他说明了如何用 idèle 解释 Hecke 特征标. 定义了 Weil 群, 使他能够给出 Artin $L$-级数与 Hecke $L$-级数的统一推广.

CHEVALLEY (1909--1984). 类域论的主要陈述都是纯代数的, 但此前所有证明都用到了分析; Chevalley 给出了纯代数证明. 借助他引入的 idèle, 他得以对无穷 Abel 扩张给出类域论的自然表述.

IWASAWA (1917--1998). 他给代数数论引入了一种重要的新方法, 其灵感来自有限域上曲线的理论.

TATE (1925--2019). 他在群上同调中证明了新结果, 使他能对类域论给出优雅的重新表述. 他与 Lubin 一起找到了生成局部域 Abel 扩张的一种显式方法.

LANGLANDS (1936--). Langlands 纲领\footnote{不要与它的几何类比相混淆, 后者有时称为几何 Langlands 纲领, 它似乎缺乏算术与解析意义。} 是一系列宏大的猜想,其中除其他内容外, 包含了一个非 Abel 的类域论.


\mainmatter         % 正文(阿拉伯页码从 1 开始)

% 引言(第 0 章)
% ============ 引言(第 0 章, 原书 pp.1--18 之前半) ============
% 定理沿用原书编号 0.1, 0.4, 0.5; 例 0.6, 0.7.
\chapter*{引言}\index{yinyan@引言(introduction)}\index{leiyulun@类域论(class field theory)}

% 引言内定理按 0.x 编号
\renewcommand{\thetheorem}{0.\arabic{theorem}}
\setcounter{theorem}{0}

L'objet de la th\'eorie du corps de classes est de montrer comment les extensions ab\'eliennes d'un corps de nombres alg\'ebriques $K$ peuvent \^etre d\'etermin\'ees par des \'el\'ements tir\'es de la connaissance de $K$ lui-m\^eme; ou, si l'on veut pr\'esenter les choses en termes dialectiques, comment un corps poss\`ede en soi les \'el\'ements de son propre d\'epassement.\footnote{类域论的目标是展示一个代数数域 $K$ 的 Abel 扩张如何能够由取自 $K$ 自身知识的元素来确定; 或者, 如果愿意用辩证的措辞来说, 一个域如何在自身中就包含着超越自身的那些元素. \cite{Chevalley1940}.}

类域论的目标, 是用域本身的算术来描述局部域或整体域的 Galois 扩张. 对 Abel 扩张, 该理论在大约 1850 至 1930 年间由 Kronecker, Weber, Hilbert, Takagi, Artin, Hasse 等人发展起来. 对非 Abel 扩张, 理论应有之形态, 最早见于 1967 年 Langlands 致 Weil 的一封信\cite{LanglandsWeil1967}. 近些年来, 非 Abel 的局部域与函数域情形已取得许多进展, 但数域情形进展较少. 约自 1980 年起, Abel 类域论已成功推广到高维域.

除少数关于更一般情形的说明外, 这些讲义将集中于局部域与数域的 Abel 扩张这一基本情形.

在引言的其余部分, 我给出 Abel 类域论的一个简短历史概览. 通篇所说数域 $K$ 的一个扩张 $L$, 是指包含在某个固定的、含 $K$ 的代数闭域中的扩张域. 例如, 若 $K$ 是包含在 $\CC$ 中的数域, 则可考虑 $K$ 在 $\CC$ 中的扩张.

\subsection*{按分裂的素理想对扩张分类}\index{fenlie@分裂(素理想)(split (prime ideal))}

回忆(ANT, 3.34) $\mathcal{O}_{K}$ 的一个素理想 $\mathfrak{p}$ 在 $K$ 的扩张 $L$ 中作(分)分解
\begin{equation*}
\mathfrak{p}\mathcal{O}_{L} = \mathfrak{P}_{1}^{e_{1}} \cdots \mathfrak{P}_{g}^{e_{g}},
\tag{2}
\end{equation*}
其中 $\mathfrak{P}_{1}, \ldots, \mathfrak{P}_{g}$ 是 $\mathcal{O}_{L}$ 中满足 $\mathfrak{P}_{i} \cap \mathcal{O}_{K} = \mathfrak{p}$ 的素理想, 且 $e_{i} \geq 1$. 此外,
\begin{equation*}
n = e_{1}f_{1} + \cdots + e_{g}f_{g},
\tag{3}
\end{equation*}
这里 $n$ 是 $L$ 在 $K$ 上的次数, $f_{i}$ 是剩余类域扩张 $\mathcal{O}_{L}/\mathfrak{P}_{i} \supset \mathcal{O}_{K}/\mathfrak{p}$ 的次数. 当某个 $e_{i} > 1$ 时, 称素理想 $\mathfrak{p}$ 在 $L$ 中分歧; 当所有 $e_{i} > 1$ 时, 称它在 $L$ 中全分歧. 当对所有 $i$ 都有 $e_{i} = 1 = f_{i}$, 从而 $\mathfrak{p} = \mathfrak{P}_{1} \cdots \mathfrak{P}_{n}$ 时, 称素理想 $\mathfrak{p}$ 在 $L$ 中分裂(或强调为完全分裂). 取 $L/K$ 的一个本原元素 $\alpha \in \mathcal{O}_{L}$. 则
\[
L = K[\alpha] \simeq K[X]/(f(X)), \quad (f\ \text{为首一不可约多项式}),
\]
且分歧的素理想整除 $f$ 的判别式(故只有有限多个)\footnote{我们有
\[
\mathrm{disc}(f)=\mathrm{disc}(\mathcal{O}_{K}[\alpha]/\mathcal{O}_{K})
=(\mathcal{O}_{L}:\mathcal{O}_{K}[\alpha])^{2}\cdot\mathrm{disc}(\mathcal{O}_{L}/\mathcal{O}_{K})
\]
(见 ANT, 2.25). 因此, 若一个素理想在 $L$ 中分歧, 则它整除 $\mathrm{disc}(f)$; 但可能存在整除 $\mathrm{disc}(f)$ 却并不在 $L$ 中分歧的素理想. 例如, 设 $L$ 是由 $\sqrt{m}$ 生成的二次域. 若 $m\equiv1\mod 4$, 则 $2$ 整除 $f$ 的判别式, 但 $2$ 在 $L$ 中并不分歧. }. 此外, 不整除 $f$ 的判别式的一个素理想在 $L$ 中分裂, 当且仅当 $f$ 模 $p$ 完全分裂. 例如, 若 $(p)$ 在 $\Q[\sqrt{m}] \simeq \Q[X]/(X^{2} - m)$ 中分歧, 则 $p$ 整除 $4m$; 而不整除 $4m$ 的素理想 $(p)$ 在 $\Q[\sqrt{m}]$ 中分裂, 当且仅当 $m$ 模 $p$ 成为平方.

当 $L/K$ 为 Galois 扩张时, $e_{1} = e_{2} = \cdots = e$ 且 $f_{1} = f_{2} = \cdots = f$, 于是这些等式简化为
\begin{equation*}
\mathfrak{p}\mathcal{O}_{L} = (\mathfrak{P}_{1} \cdots \mathfrak{P}_{g})^{e}, \quad n = efg.
\tag{4}
\end{equation*}
设 $\Spl(L/K)$ 表示 $\mathcal{O}_{K}$ 中在 $L$ 中分裂的素理想集. 十九世纪末, Frobenius 用分析方法证明了如下陈述.

\begin{theorem}[FROBENIUS]\label{thm:0.1}
当 $L/K$ 为 Galois 扩张时, 集合 $\Spl(L/K)$ 在所有 $\mathcal{O}_{K}$ 的素理想集中具有密度 $1/[L:K]$.
\end{theorem}

由此可知 $\Spl(L/K)$ 决定了 $L$: 若在 $L$ 与 $L'$ 中分裂的是同一批素理想, 则推出
\[
\Spl(L/K) = \Spl(LL'/K) = \Spl(L'/K)
\]
(ANT, 8.38 的证明), 从而
\[
[L:K] = [LL':K] = [L':K].
\]
于是 $K$ 的 Galois 扩张由集合 $\Spl(L/K)$ 分类. 仍有以下两个基本问题.

\begin{npara}\label{para:0.2} 当 $L$ 跑遍 $K$ 的有限 Galois 扩张时, 确定这些集合 $\Spl(L/K)$.
\end{npara}

\begin{npara}\label{para:0.3} 用集合 $\Spl(L/K)$ 描述扩张 $L/K$.
\end{npara}

我们将看到, 当 $L$ 跑遍 $K$ 的 Abel 扩张时, \ref{para:0.2} 有精确的答案: 此时集合 $\Spl(L/K)$ 由某些同余条件刻画. 然而对非 Abel 的 Galois 扩张, 这些集合不被同余条件刻画, 所能期望的最好结果也不过是某种解析描述.

关于 \ref{para:0.3}, 注意 $\Spl(L/K)$ 的密度决定了 $L/K$ 的次数. 我们还希望确定 $L/K$ 的整个 Galois 群, 并且希望确定 $K$ 的每个素点如何在 $L$ 中分解. 例如, 我们希望确定在 $L$ 中分歧的素理想集, 并对那些不分歧的素理想, 确定 $\mathfrak{p}$ 上方的那些素点的剩余类次数 $f(\mathfrak{p})$.

在继续之前我作一点说明: 在上述讨论中, 我们可以忽略 $K$ 的有限个素理想. 例如, 若 $S$ 是 $\mathcal{O}_{K}$ 的素理想的一个固定有限集, 则我们可以将上述讨论中的 $\Spl(L/K)$ 替换为不在 $S$ 中且在 $L$ 中分裂的素理想集 $\Spl_{S}(L/K)$.

\begin{NOTES}
\phantomsection\label{pg:orig3}对一个数域的 Galois 扩张 $L/K$, Frobenius 给 $K$ 中每个在 $L$ 中不分歧的素理想 $\mathfrak{p}$ 对应了 $\Gal(L/K)$ 中的一个共轭类 $(\mathfrak{p}, L/K)$; 他猜想, 给出固定共轭类 $C$ 的素理想的密度是 $|C|/|G|$. 按构造, $(\mathfrak{p}, L/K)$ 中的元素在 $G$ 中的阶为 $f(\mathfrak{p})$, 因此把这一陈述用于平凡共轭类即给出定理\ref{thm:0.1}. Frobenius 只能证明比其猜想更弱的一个陈述(对 \ref{thm:0.1} 已足够), 其中某些共轭类被归并到一起\footnote{Frobenius, G., 1896, \emph{Die Beziehungen zwischen den Primidealen eines algebraischen K\"orpers und den Substitutionen einer Gruppe}, Berlin Akad.-Ber. 关于 Frobenius 定理的陈述, 见\hyperlink{item:viii.A-10}{练习 A-10} (原书第 268 页).}, 而完整的猜想直到 1926 年才由 Chebotarev 证明.
\end{NOTES}

\subsection*{$\Q$ 的二次扩张}\index{ercikuozhang@二次扩张(quadratic extension)}

对 $\Q$ 的二次扩张, \ref{para:0.2} 的答案由二次互反律给出. 例如, 设 $p$ 为奇素数, 令 $p^{*} = (-1)^{\frac{p-1}{2}}p$, 于是 $p^{*} \equiv 1 \pmod{4}$. 则 $\Q[\sqrt{p^{*}}]/\Q$ 仅在 $p$ 处分歧. 一个素数 $q \neq p$ 在 $\Q[\sqrt{p^{*}}]$ 中分裂, 当且仅当 $p^{*}$ 是模 $q$ 的平方, 即 $\left(\frac{p^{*}}{q}\right) = 1$. 但当 $q$ 为奇数时, 二次互反律说 $\left(\frac{p^{*}}{q}\right) = \left(\frac{q}{p}\right)$, 故 $q$ 在 $\Q[\sqrt{p^{*}}]$ 中分裂当且仅当 $q$ 是模 $p$ 的平方. 设 $H \subset (\Z/p\Z)^{\times}$ 为平方元集. 则 $\Spl(\Q[\sqrt{p^{*}}]/\Q)$ 由满足 $q \bmod p \in H$ 的素数 $q$ 组成. 回忆 $(\Z/p\Z)^{\times}$ 是阶为 $p-1$ 的循环群, 故 $H$ 是 $(\Z/p\Z)^{\times}$ 中唯一的指数为 $2$ 的子群. 于是
\[
(\Z/p\Z)^{\times}/H \simeq \Gal(\Q[\sqrt{p^{*}}]/\Q).
\]
更一般地, 设 $S$ 为素数集的一个有限集. 令 $m$ 为只被 $S$ 中素数整除的正整数, 并设 $m$ 为奇数或能被 $4$ 整除. 对 $(\Z/m\Z)^{\times}$ 中一个指数为 $2$ 的子群 $H$, 令
\[
p(H) = \{(p) \mid p\ \text{为素数},\ p \notin S,\ p \bmod m \in H\}.
\]
则当 $L$ 跑遍在 $S$ 外不分歧的 $\Q$ 的二次扩张时, 集合 $\Spl_{S}(L/\Q)$ 恰为这些集合 $p(H)$ (对变动的 $m$). 特别地, 我们看到每个这样的集合 $\Spl_{S}(L/\Q)$ 都由某个 $m$ 模 $m$ 的同余条件决定.

\begin{NOTES}
在初等数论的课程中, 二次互反律往往显得像是一件奇特的珍玩, 其深刻程度并不超过许多其他奇特的结论. 事实上, 正如上述讨论所表明的, 它应当被视为类域论中的第一个结果, 这有助于解释 Euler, Lagrange, Legendre 与 Gauss 为何对它倾注了如此多的注意力.
\end{NOTES}

\subsection*{非分歧 Abel 扩张}\index{feifenqikuozhang@非分歧扩张(unramified extension)}\index{aboerkuozhang@Abel 扩张(abelian extension)}

设 $I$ 为 $K$ 的分式理想群, 即由 $\mathcal{O}_{K}$ 的素理想生成的自由 Abel 群; 设 $i : K^{\times} \to I$ 为把 $a \in K^{\times}$ 映到主理想 $(a)$ 的映射. $K$ 的类群 $C$ 是 $I/i(K^{\times})$. 给出一个 $C$ 的子群 $H$, 等价于给出一个包含 $i(K^{\times})$ 的 $I$ 的子群 $\tilde{H}$.

我所说 $K$ 的一个``素点'', 是指 $K$ 上非平凡赋值的等价类\footnote{英文中 “place” 比 “prime” 更常用, 我以后会改用前者; 也有人用 “spot”. }. 于是每个非零素理想对应恰好一个素点, 每个嵌入 $K \hookrightarrow \R$ 对应一个, 每对共轭的非实嵌入 $K \hookrightarrow \CC$ 对应一个. 相应的素点分别称为有限素点、实素点与复素点. $K$ 的一个元素若在对应嵌入 $K \hookrightarrow \R$ 的实素点处映到 $\R$ 中的正元素, 则称它在该实素点处为正. 若 $K$ 的一个实素点 $v$ 在 $L$ 中的每个上方素点仍为实素点, 则称 $v$ 在扩张 $L/K$ 中分裂; 否则称 $v$ 在 $L/K$ 中分歧. 例如, $\Q[\sqrt{-5}]$ 在 $\Q$ 上恰好在素点 $(2)$, $(5)$ 与 $\infty$ 处分歧.

我们说有限扩张 $L/K$ 是非分歧的, 如果它在 $K$ 的所有素点处都不分歧. 这意味着 $\mathcal{O}_{K}$ 的每个素理想在 $L$ 中都不分歧, 且 $K$ 的每个实素点仍保持为实(即 $K$ 的每个实嵌入都延拓为 $L$ 的一个实嵌入).

设 $H$ 为 $K$ 的类群 $C$ 的一个子群. 若 $\mathcal{O}_{K}$ 中在 $L$ 中分裂的素理想恰为那些落在 $\tilde{H}$ 中的素理想, 则称 $K$ 的一个有限非分歧 Abel 扩张 $L$ 是 $H$ 的类域. 下一个定理对 $K$ 的非分歧 Abel 扩张同时回答了 \ref{para:0.2} 与 \ref{para:0.3}.

\begin{theorem}[FURTW\"ANGLER 1907]\label{thm:0.4}
对 $C$ 的每个子群都存在类域; 它是唯一的, 并且 $K$ 的每个有限非分歧 Abel 扩张都作为某个 $C$ 的子群的类域出现. 若 $L$ 是 $H$ 的类域, 则 $\Gal(L/K) \approx C/H$, 并且对 $K$ 的每个素理想 $\mathfrak{p}$, $f(\mathfrak{p})$ 等于 $\mathfrak{p}$ 在群 $C/H$ 中像的阶.
\end{theorem}

对应 $K$ 的有限非分歧 Abel 扩张 $L$ 的 $C$ 的子群 $H$, 是由在 $L$ 中分裂的素理想生成的子群. 类域的唯一性由 \ref{thm:0.1} 推出.

$C$ 的平凡子群的类域称为 $K$ 的 Hilbert 类域. 它是 $K$ 在 $K$ 的所有素点(包括无穷素点)处都非分歧的最大 Abel 扩张. 在其中分裂的素理想恰为主理想, 且 $\Gal(L/K) \approx C$. 例如, $K = \Q[\sqrt{-5}]$ 的类数为 $2$, 其 Hilbert 类域为 $K[\sqrt{-1}]$ (ANT, 4.11).

\begin{NOTES}
定理\ref{thm:0.4} 由 Hilbert 于 1897 年猜想, 并由他的学生 Furtw\"angler 于 1907 年证明.
\end{NOTES}

\subsection*{分歧 Abel 扩张}\index{fenqikuozhang@分歧扩张(ramified extension)}

\phantomsection\label{pg:orig4}为把定理\ref{thm:0.4} 推广到分歧扩张, 我们需要推广理想类群的概念. 为此应先考察我们已经理解的 Abel 扩张, 即 $\Q$ 的分圆扩张. 设 $m$ 为奇数或能被 $4$ 整除的正整数, 并设 $\zeta_{m}$ 为本原 $m$ 次单位根. 回忆(ANT, 第 6 章) $\Q[\zeta_{m}]$ 是 $\Q$ 的一个扩张, 其 Galois 群为 $(\Z/m\Z)^{\times}$, 其中 $[n]$ 的作用是 $\zeta_{m} \mapsto \zeta_{m}^{n}$. 此外, 在我们关于 $m$ 的条件下, 在 $\Q[\zeta_{m}]$ 中分歧的素理想恰为满足 $p \mid m$ 的理想 $(p)$ 以及 $\infty$. 我们如何能把 $(\Z/m\Z)^{\times}$ 实现为一个理想类群? 我们可以尝试把一个理想 $(p)$ 映到 $(\Z/m\Z)^{\times}$ 中 $[p]$ 的类, 但这只有当 $(p,m) = 1$ 时才可能. 于是令 $S$ 为满足条件 $p \mid m$ 的素理想 $(p)$ 的集合, 并令 $I^{S}$ 为 $\Q$ 的、由不在 $S$ 中的素理想生成的分式理想群. $I^{S}$ 的每个元素可表示为 $(r/s)$, 其中 $r$, $s$ 为与 $m$ 互素的正整数; 我们把它映到 $(\Z/m\Z)^{\times}$ 中的 $[r][s]^{-1}$. 这给出满同态 $I^{S} \to (\Z/m\Z)^{\times}$, 其核待描述. 设 $m = \prod p^{\operatorname{ord}_{p}(m)}$, 于是
\[
\Z/m\Z \simeq \prod_{p \mid m} \Z/p^{\operatorname{ord}_{p}(m)} \quad (\text{中国剩余定理}),
\]
\[
(\Z/m\Z)^{\times} \simeq \prod_{p \mid m} (\Z/p^{\operatorname{ord}_{p}(m)})^{\times}.
\]
一个理想 $(r)$(其中 $(r,m) = 1 = (s,m)$) 映到 $1$, 当且仅当 $r$ 与 $s$ 同号, 并且对一切整除 $m$ 的 $p$, $r$ 与 $s$ 在 $(\Z/p^{\operatorname{ord}_{p}(m)})^{\times}$ 中映到同一元素. 注意在 $\infty$ 处的条件是 $\frac{r}{s} > 0$, 而在 $p$ 处的条件是 $\operatorname{ord}_{p}(\frac{r}{s} - 1) \geq \operatorname{ord}_{p}(m)$.

为推广这一点, 定义数域 $K$ 的一个\emph{模}\index{mo@模(modulus)(modulus)}\index{shexianleiqun@射线类群(ray class group)}为一个(形式)乘积 $\mathfrak{m} = \mathfrak{m}_{0}\mathfrak{m}_{\infty}$, 其中 $\mathfrak{m}_{0}$ 是 $\mathcal{O}_{K}$ 的一个整理想, $\mathfrak{m}_{\infty}$ 是 $K$ 的若干实素点的乘积. 设 $S(\mathfrak{m})$ 为整除 $\mathfrak{m}$ 的素点集, 并定义 $I^{S(\mathfrak{m})}$ 为 $\mathcal{O}_{K}$ 的、由不在 $S(\mathfrak{m})$ 中的素理想生成的分式理想群. 以\emph{模}\index{mo@模(modulus)(modulus)} $\mathfrak{m}$ 为模的\emph{射线类群}\index{shexianleiqun@射线类群(ray class group)}\index{shexianleiqun@射线类群(ray class group)} $C_{\mathfrak{m}}$ 定义为: 把 $I^{S(\mathfrak{m})}$ 模去由如下元素 $a$ 生成的主理想子群所得到的商, 其中 $a$ 满足: 在所有整除 $\mathfrak{m}_{\infty}$ 的实素点处 $a > 0$, 并且对所有整除 $\mathfrak{m}_{0}$ 的素理想 $\mathfrak{p}$ 有 $\operatorname{ord}_{\mathfrak{p}}(a - 1) \geq \operatorname{ord}_{\mathfrak{p}}(\mathfrak{m}_{0})$(这里 $\operatorname{ord}_{\mathfrak{p}}(\mathfrak{m}_{0})$ 是 $\mathfrak{p}$ 在 $\mathfrak{m}_{0}$ 中的指数).

例如, 若 $\mathfrak{m} = 1$, 则 $C_{\mathfrak{m}} = C$, 即通常的类群. 若 $\mathfrak{m}$ 是 $K$ 的实素点的乘积(于是 $\mathfrak{m}_{0} = 1$), 则 $C_{\mathfrak{m}}$ 是 $I$ 模去由 $K$ 的全正元\footnote{称 $K$ 的元素 $a$ 是\emph{全正}的, 若对 $K$ 的每个实嵌入 $\sigma\colon K\hookrightarrow\R$ 都有 $\sigma(a)>0$. 例如 $-5$ 在 $\Q[i]$ 中全正, 但在 $\Q$ 中不全正. 这看起来有些奇怪, 但就我书架上的书而言, 这个定义是标准的. 除该域为全实域的情形外, 更好的说法是: 该元素“在所有实素点处为正”. }生成的理想所得的子群的商, 即窄类群\footnote{英文原词 narrow-class group 中的连字符是要表明“窄”修饰的是那些类, 而不是这个群. }. 若 $K = \Q$ 且 $\mathfrak{m} = (m)\infty$, 则 $C_{\mathfrak{m}} \simeq (\Z/m\Z)^{\times}$, 如第一段所述.

设 $H$ 为某个模 $\mathfrak{m}$ 的射线类群 $C_{\mathfrak{m}}$ 的一个子群, 并设 $\tilde{H}$ 为它在 $I^{S(\mathfrak{m})}$ 中的原像. 称在整除 $\mathfrak{m}$ 的素点处非分歧的 $K$ 的 Abel 扩张 $L$ 是 $H$ 的类域, 若 $\mathcal{O}_{K}$ 中不整除 $\mathfrak{m}_{0}$ 且在 $L$ 中分裂的那些素理想恰为落在 $\tilde{H}$ 中的素理想. 下一个定理把定理\ref{thm:0.4} 推广到 $K$ 的所有 Abel 扩张.

\begin{theorem}[TAKAGI]\label{thm:0.5}
对射线类群 $C_{\mathfrak{m}}$ 的每个子群都存在类域; 它是唯一的, 并且 $K$ 的每个有限 Abel 扩张都作为某个射线类群的某个子群的类域出现. 若 $L$ 是 $H \subset C_{\mathfrak{m}}$ 的类域, 则 $\Gal(L/K) \approx C_{\mathfrak{m}}/H$, 并且 $K$ 的不整除 $\mathfrak{m}$ 的素理想 $\mathfrak{p}$ 在 $L$ 中不分歧, 其剩余类次数 $f(\mathfrak{p})$ 等于 $\mathfrak{p}$ 在群 $C_{\mathfrak{m}}/H$ 中像的阶.
\end{theorem}

对 $K$ 的任意有限素集 $S$, 定理\ref{thm:0.5} 完全解决了确定仅在 $S$ 中的素处分歧的 Abel Galois 扩张 $L/K$ 的集合 $\Spl_S(L/K)$ 这一问题.

该定理可以更精确. 设 $\mathfrak{m}$ 与 $\mathfrak{m}'$ 为 $K$ 的模. 若 $\mathfrak{m}' \mid \mathfrak{m}$, 则包含 $I^{S(\mathfrak{m})} \hookrightarrow I^{S(\mathfrak{m}')}$ 定义满同态 $C_{\mathfrak{m}} \twoheadrightarrow C_{\mathfrak{m}'}$; 若 $H \subset C_{\mathfrak{m}}$ 是某个子群 $H' \subset C_{\mathfrak{m}'}$ 的原像, 则 $H'$ 的每个类域也将是 $H$ 的类域. 若 $L$ 是 $H \subset C_{\mathfrak{m}}$ 的类域, 且 $H$ 并非以这种方式来自某个真整除 $\mathfrak{m}$ 的模, 则 $S(\mathfrak{m})$ 恰由在 $L$ 中分歧的素理想组成. 这个尽可能小的模称为 $L/K$ 的\emph{导子}\index{daozi@导子(conductor)}.

设 $L$ 为 $K$ 的一个 Abel 扩张. 则对每个只被在 $L$ 中分歧的素点整除的、充分大的模 $\mathfrak{m}$, $L$ 都是由在 $L$ 中分裂的素理想的类生成的子群 $H \subset C_{\mathfrak{m}}$ 的类域.

对 $K = \Q$, 该定理由 ANT 第 6 章证明的结果推出, 除了``$\Q$ 的每个 Abel 扩张都作为类域出现''这一断言, 它本质上就是 Kronecker--Weber 定理: $\Q$ 的每个 Abel 扩张都包含于某个分圆扩张.

\begin{example}\label{ex:0.6}
域 $\Q[\sqrt{6}]$ 的类数为 $1$, 故等于它自己的 Hilbert 类域. 然而它的窄类群阶为 $2$, 并且确实有 $\Q[\sqrt{6}]$ 的一个二次扩张, 即 $\Q[\sqrt{-2}, \sqrt{-3}]$, 在它所有有限素点处都非分歧于 $\Q[\sqrt{6}]$(但在两个无穷素点处都分歧).
\end{example}

当 Takagi 证明定理\ref{thm:0.5} 时, 他只证明了群 $C_{\mathfrak{m}}/H$ 与 $\Gal(L/K)$ 同构——他并未构造一个具体的同构. 有一个明显的映射 $I^{S(\mathfrak{m})} \to \Gal(L/K)$, 即把$S(\mathfrak{m})$ 中的一个素理想 $\mathfrak{p}$ 映到 $\Gal(L/K)$ 中相应的 Frobenius 元. Frobenius 定理表明这是满射, 问题是要证明它对某个 $\mathfrak{m}$ 经由 $C_{\mathfrak{m}}$ 分解, 即该映射容许一个模. Artin 于 1923 年猜想这一点, 并于 1927 年证明(见下).

\begin{example}\label{ex:0.7}
设 $d = p_{1}^{*} \cdots p_{t}^{*}$, 其中 $p_{i}$ 为互不相同的奇素数, 且 $p_{i}^{*} = (-1)^{\frac{p_{i}-1}{2}}p_{i}$ 如前. 域 $K = \Q[\sqrt{d}]$ 恰好在素理想 $(p_{1}), \ldots, (p_{t})$ 处分歧. 考虑 $K[\sqrt{p_{i}^{*}}]$. 它包含 $\Q[\sqrt{p_{1}^{*} \cdots p_{i-1}^{*} p_{i+1}^{*} \cdots p_{t}^{*}}]$, 后者在 $(p_{i})$ 处非分歧, 故 $(p_{i})$ 不能在 $K[\sqrt{p_{i}^{*}}]$ 中全分歧. 由于 $p_{i}$ 在 $K/\Q$ 中分歧, 推出 $K$ 中位于 $(p_{i})$ 之上的素点在 $K[\sqrt{p_{i}^{*}}]$ 中不分歧. 没有其他素点分歧, 故 $K[\sqrt{p_{i}^{*}}]$ 在 $K$ 上非分歧. 由 Kummer 理论(FT 5.29)知
\[
L \overset{\mathrm{def}}{=} K[\sqrt{p_{1}^{*}}, \sqrt{p_{2}^{*}}, \ldots, \sqrt{p_{t-1}^{*}}]
\]
在 $K$ 上的次数为 $2^{t-1}$, 且 $\Gal(L/K) \approx (\Z/2\Z)^{t-1}$. 设 $C_{\infty}$ 为 $K$ 的窄类群. 上面的构造表明 $(C_{\infty} : C_{\infty}^{2}) \geq 2^{t-1}$. 事实上, 只需稍多努力就能证明如下陈述:

设 $K$ 是在其中分歧的有限素点有 $t$ 个的 $\Q$ 的二次扩张. 则 $(C_{\infty} : C_{\infty}^{2}) = 2^{t-1}$ (\cite{Koch1992}, 定理 2.114).

这一结果 Gauss 已知(用不同的方法, 不同的语言). 特别地, 我们看到借助类域论, 容易构造出满足 $(C : C^{2})$ 很大的 $\Q$ 的二次扩张. 相比之下, 截至 1991 年, 尚不知道哪个二次域满足 $(C : C^{3}) > 3^{6}$. 所有在二次数域的类群中构造 $3$ 阶元素的方法似乎都涉及椭圆曲线.
\end{example}

\begin{NOTES}
射线类群 $C_{\mathfrak{m}}$ 由 Weber 于 1897 年引入. 定理\ref{thm:0.5} 由 Hilbert 与 Weber 猜想, 并由 Takagi 在 1915 至 1922 年间发表的一系列论文中证明——尤其参见他在 1920 年国际数学家大会上的报告\footnote{见他的全集. }, 其中他陈述的结果几乎与我们的完全一致.
\end{NOTES}

% ============ 以下为原书引言 pp.16--33 概述子节(此前缺失, 现补译) ============
\setcounter{theorem}{7}

\subsection*{密度定理}\index{midu@密度定理(density theorem)}\index{Dirichlet 定理@Dirichlet 定理(Dirichlet's theorem)}

设 $m$ 为 $\geq 2$ 的整数. 设 $a$ 为与 $m$ 互素的整数, 考虑序列
\begin{equation*}
\ldots, a-m, a, a+m, a+2m, a+3m, \ldots, a+km, \ldots.\tag{5}
\end{equation*}
此序列只依赖于 $a \bmod m$, 故恰有 $\varphi(m) \overset{\mathrm{def}}{=} |(\mathbb{Z}/m\mathbb{Z})^{\times}|$ 个不同的序列. Dirichlet 证明素数在这些序列中均匀分布, 即每个序列中的素数集在所有素数集中具有密度 $1/\varphi(m)$.

我们想在 Galois 群的意义下解释这一点. 设 $L/K$ 是数域的有限扩张, Galois 群为 $G$. 回忆(ANT, 第 8 章)对 $\mathcal{O}_{L}$ 中每个素理想 $\mathfrak{P}$, 存在 $G$ 中一个 Frobenius 元 $\sigma$ 使得
\begin{itemize}
\item $\sigma\mathfrak{P} = \mathfrak{P}$, 且
\item 对所有 $\alpha \in \mathcal{O}_{L}$, $\sigma\alpha \equiv \alpha^{q} \bmod \mathfrak{P}$, 其中 $q$ 是剩余类域 $\mathcal{O}_{K}/\mathfrak{p}$ 的元素个数, $\mathfrak{p} = \mathfrak{P} \cap \mathcal{O}_{K}$.
\end{itemize}
当 $\mathfrak{P}$ 在 $\mathfrak{p}$ 上不分歧时, 此元素由这些条件唯一确定——此时记作 $(\mathfrak{P}, L/K)$. 若 $\mathfrak{P}'$ 是 $\mathcal{O}_{L}$ 中另一个整除 $\mathfrak{p}$ 的素理想, 则 $\mathfrak{P}' = \tau\mathfrak{P}$ 对某个 $\tau \in G$, 且
\[
(\mathfrak{P}', L/K) = (\tau\mathfrak{P}, L/K) = \tau(\mathfrak{P}, L/K)\tau^{-1}.
\]
因此, 对 $K$ 中每个在 $L$ 中不分歧的素理想 $\mathfrak{p}$, 集合
\[
\{(\mathfrak{P}, L/K) \mid \mathfrak{P} \cap \mathcal{O}_{K} = \mathfrak{p}\}
\]
是 $G$ 中一个完整的共轭类——记作 $(\mathfrak{p}, L/K)$. 当 $G$ 交换时, 共轭类只含一个元素, 故 $(\mathfrak{p}, L/K)$ 可视为 $G$ 的元素.

现在考虑 $\mathbb{Q}[\zeta_{m}]/\mathbb{Q}$, 其中 $m$ 为奇数 $>1$ 或能被 $4$ 整除的正整数. 回忆 $\Gal(\mathbb{Q}[\zeta_{m}]/\mathbb{Q}) \simeq (\mathbb{Z}/m\mathbb{Z})^{\times}$, 其中 $[n]$ 的作用为 $\zeta \mapsto \zeta^{n}$. 不整除 $m$ 的素理想 $(p)$ 在 $\mathbb{Q}[\zeta_{m}]$ 中不分歧, 且 $(p, \mathbb{Q}[\zeta_{m}]/\mathbb{Q}) = [p]$(正如所料). 令 $\tau = [a] \in \Gal(\mathbb{Q}[\zeta_{m}]/\mathbb{Q})$. 则 $(p, \mathbb{Q}[\zeta_{m}]/\mathbb{Q}) = \tau$ 当且仅当 $p$ 落在序列 (5) 中, 于是 Dirichlet 定理说
\[
\{p \mid (p, \mathbb{Q}[\zeta_{m}]/\mathbb{Q}) = \tau\}
\]
在所有素数中具有密度 $1/\varphi(m)$. 对 $\mathbb{Q}$ 的分圆扩张, 这证明了 Frobenius 的猜想(本书 p.\pageref{pg:orig3}).

为证明他关于算术级数中的素数的定理, Dirichlet 引入了如今所谓的 Dirichlet $L$-级数. 设 $\chi \colon (\mathbb{Z}/m\mathbb{Z})^{\times} \to \mathbb{C}^{\times}$ 是 $(\mathbb{Z}/m\mathbb{Z})^{\times}$ 的一个特征(即同态). 则附于 $\chi$ 的 Dirichlet $L$-级数为
\[
L(s, \chi) = \prod_{(p,m)=1,\,p\ \text{素数}} \frac{1}{1 - \chi(p)p^{-s}} = \sum_{(n,m)=1,\,n>0} \chi(n)n^{-s}.
\]
当 $\chi$ 为平凡(主)特征时, 此式给出 Riemann $\zeta$ 函数, 只是缺了有限多个因子. Dedekind 把 $\zeta$ 函数的概念推广到每个数域:
\[
\zeta_{K}(s) = \prod_{\mathfrak{p}} \frac{1}{1 - \mathbb{N}\mathfrak{p}^{-s}} = \sum_{\mathfrak{a} \subset \mathcal{O}_{K}} \mathbb{N}\mathfrak{a}^{-s}, \quad \mathbb{N}\mathfrak{a} = (\mathcal{O}_{K} : \mathfrak{a}),
\]
而 Weber 把 $L$-级数的概念推广到任意射线类群 $C_{\mathfrak{m}}$ 的特征 $\chi$:
\[
L(s, \chi) = \prod_{(\mathfrak{p},\mathfrak{m})=\mathcal{O}_{K}} \frac{1}{1 - \chi(\mathfrak{p})\mathbb{N}\mathfrak{p}^{-s}} = \sum_{(\mathfrak{a},\mathfrak{m})=1} \chi(\mathfrak{a})\mathbb{N}\mathfrak{a}^{-s}.
\]
Weber 证明 $L(s,\chi)$ 可解析延拓到整个复平面上的亚纯函数, 并且在 $\chi$ 不是平凡特征时甚至是全纯的.

在整个复平面上这个函数仍记作 $L(s,\chi)$, 称为 Dirichlet 或 Weber $L$-函数.

\subsection*{Artin $L$-函数}\index{artinl@Artin $L$-函数(Artin $L$-function)}

设 $L/K$ 是有限 Galois 扩张, Galois 群为 $G$, 并设 $S$ 是 $\mathcal{O}_{K}$ 的有限素理想集, 包含所有在 $L$ 中分歧者. 设
\[
\rho \colon G \to \mathrm{GL}_{n}(\mathbb{C})
\]
是 $G$ 的一个表示(即群同态). 对每个不在 $S$ 中的素理想, 我们得到 $G$ 的 Frobenius 元的一个共轭类 $(\mathfrak{p}, L/\mathbb{Q})$, 它被 $\rho$ 映到 $\mathrm{GL}_{n}(\mathbb{C})$ 中的一个共轭类 $\Phi_{\mathfrak{p}}(\rho)$. $\Phi_{\mathfrak{p}}(\rho)$ 中的元素可对角化(Maschke 定理, GT 7.5), 且 $\Phi_{\mathfrak{p}}(\rho)$ 由其中元素的公共特征多项式 $\det(I - \Phi_{\mathfrak{p}}(\rho)T)$ 决定. Artin 定义
\[
L_{S}(s, \rho) = \prod_{\mathfrak{p} \notin S} \frac{1}{\det(I - \Phi_{\mathfrak{p}}(\rho)(\mathbb{N}\mathfrak{p})^{-s})}, \quad s \in \mathbb{C}.
\]
此乘积在某个右半平面中收敛到一个全纯函数. 这类函数称为 Artin $L$-函数.

注意 $G$ 的一维表示就是一个特征 $\chi \colon G \to \mathbb{C}^{\times}$, 且
\[
L_{S}(s, \chi) = \prod_{\mathfrak{p} \notin S} \frac{1}{1 - \chi(\mathfrak{p})(\mathbb{N}\mathfrak{p})^{-s}}.
\]
关于 Dirichlet 级数的一条初等等价引理表明, 因子 $\frac{1}{\det(I - \Phi_{\mathfrak{p}}(\rho)(\mathbb{N}\mathfrak{p})^{-s})}$ 由全纯函数 $L_{S}(s, \rho)$ 唯一确定. 若 $\rho$ 是单射, 则
\[
\operatorname{Spl}_{S}(L/K) = \{\mathfrak{p} \mid \Phi_{\mathfrak{p}}(\rho) = \{I\}\},
\]
故函数 $L_{S}(s, \rho)$ 决定了 $\operatorname{Spl}_{S}(L/K)$. 对非 Abel 扩张, 与其直接研究 $\operatorname{Spl}_{S}(L/K)$, 不如研究 Artin $L$-函数 $L_{S}(s, \rho)$.

\begin{NOTES}
关于 $L$-级数、尤其是 Artin $L$-级数的历史介绍, 见 Noah Snyder, \emph{Artin's $L$-functions: A Historical Approach}, 2002\cite{Snyder2002}.
\end{NOTES}

\subsection*{Artin 映射}\index{artinys@Artin 映射(Artin map)}\index{hufanying@互反映射(reciprocity map)}

回忆对 $\mathbb{Q}$ 的分圆扩张 $\mathbb{Q}[\zeta_{m}]$, 存在一个同构
\[
(\mathbb{Z}/m\mathbb{Z})^{\times} \to \Gal(\mathbb{Q}[\zeta_{m}]/\mathbb{Q}),
\]
把 $[p]$ 映到 $(p, \mathbb{Q}[\zeta_{m}]/\mathbb{Q})$. $\Gal(\mathbb{Q}[\zeta_{m}]/\mathbb{Q})$ 的一个特征 $\chi$ 通过复合给出一个 $(\mathbb{Z}/m\mathbb{Z})^{\times}$ 的特征 $\chi'$. 由于
\[
\chi'(p) = \chi((p, \mathbb{Q}[\zeta_{m}]/\mathbb{Q})),
\]
我们有
\[
L_{S}(s, \chi) = L_{S}(s, \chi').
\]
于是, 在此情形 Artin $L$-级数 $L_{S}(s, \chi)$ 是一个 Dirichlet $L$-级数.

更一般地, 类似的事是否也成立? 例如, 由一维表示定义的 Artin $L$-级数是否是一个 Weber $L$-级数? 设 $L/K$ 是 Abel 扩张, Galois 群为 $G$, 并设 $S$ 是 $K$ 的素点集, 包含那些在 $L$ 中分歧者. 设 $I^{S}$ 是由不在 $S$ 中的素理想生成的 $K$ 的分式理想群. 则 $I^{S}$ 是自由群, 故映射 $\mathfrak{p} \mapsto (\mathfrak{p}, L/K)$ 唯一地延拓为一个同态
\begin{equation*}
I^{S} \to \Gal(L/K),\tag{6}
\end{equation*}
由 Frobenius 密度定理\ref{thm:0.1}知其满射. 若存在模 $\mathfrak{m}$ 使 $S(\mathfrak{m}) = S$, 使得此同态经 $C_{\mathfrak{m}}$ 分解, 则 $\Gal(L/K)$ 的一个特征 $\chi$ 通过复合给出 $C_{\mathfrak{m}}$ 的一个特征 $\chi'$, 且
\[
L_{S}(s, \chi) = L_{S}(s, \chi'),
\]
即 Artin $L$-级数 $L_{S}(s, \chi)$ 将是一个 Weber $L$-级数.

Artin 当然猜想这样的模确实存在, 但他直到看到 Chebotarev 对其密度定理的证明后才得以证明, 并由此得到如下更完整的定理. 在陈述之前, 我们需要一些记号. 对模 $\mathfrak{m}$, 令 $P_{\mathfrak{m}}$ 为 $I^{S(\mathfrak{m})} \to C_{\mathfrak{m}}$ 的核, 于是 $P_{\mathfrak{m}}$ 由满足某些依赖于 $\mathfrak{m}$ 的正性与同余条件的元素所生成的主理想组成. 对 $K$ 的扩张 $L$, 令 $S'(\mathfrak{m})$ 表示 $L$ 中位于 $S(\mathfrak{m})$ 中素点之上的素点. 存在一个范映射 $\Nm \colon I^{S'(\mathfrak{m})} \to I^{S(\mathfrak{m})}$(ANT, p. 68), 使得 $\Nm(\mathfrak{P}) = \mathfrak{p}^{f}$, 其中 $\mathfrak{p} = \mathfrak{P} \cap \mathcal{O}_{K}$, $f$ 为剩余类次数. 注意
\[
(\Nm(\mathfrak{P}), L/K) = (\mathfrak{p}^{f}, L/K) \overset{\mathrm{def}}{=} (\mathfrak{p}, L/K)^{f} = 1,
\]
故 $\Nm(I^{S'(\mathfrak{m})})$ 包含在 $I^{S(\mathfrak{m})} \to C_{\mathfrak{m}}$ 的核中.

\begin{theorem}[ARTIN 1927]\label{thm:0.8}
设 $L$ 是 $K$ 的 Abel 扩张, $S$ 为在 $L$ 中分歧的素点集. 则对某个满足 $S(\mathfrak{m}) = S$ 的模 $\mathfrak{m}$, 同态
\[
\mathfrak{p} \mapsto (\mathfrak{p}, L/K) \colon I^{S} \to \Gal(L/K)
\]
经过 $C_{\mathfrak{m}} \overset{\mathrm{def}}{=} I^{S(\mathfrak{m})}/P_{\mathfrak{m}}$ 分解, 并定义一个同构
\[
I^{S(\mathfrak{m})}/P_{\mathfrak{m}} \cdot \Nm(I^{S'(\mathfrak{m})}) \to \Gal(L/K).
\]
特别地, 在 $L$ 中分裂的素理想恰为子群
\[
\tilde{H} = P_{\mathfrak{m}} \cdot \Nm(I^{S'(\mathfrak{m})})
\]
中的那些(属于 $I^{S}$).
\end{theorem}

定理中的同态 $\psi_{L/K} \colon C_{\mathfrak{m}} \to \Gal(L/K)$ 称为 Artin 映射(或互反映射). 正如 Artin 所指出的, 该定理包含了所有已知的互反律(见第 \ref{ch:VIII} 章), 并且应被视为这些定律对无 $1$ 次单位根之域 $K$ 的正确推广.

注意, 若 Artin 映射不是满射, 则存在 $K$ 的一个真扩张, 其中 $K$ 在 $S$ 外的每个素点都分裂. 分析方法——具体说 Frobenius 的密度定理——表明这样的扩张不存在, 但要代数地证明这一点却很困难(见第 \ref{ch:VII} 章).

由于一维表示的 Artin $L$-函数都是 Weber $L$-级数, 它们可延拓为整个复平面上的亚纯函数. Artin 猜想有限群的每个特征都是接近 Abel 的特征的 $\mathbb{Z}$-线性组合, 并证明此猜想蕴含所有 Artin $L$-级数都延拓为整个复平面上的亚纯函数. 这一点由 Brauer 于 1946 年证明.

Artin 猜想, 若 $\rho$ 是非平凡不可约表示, 则 $L(s, \rho)$ 在整个复平面上全纯. 他指出证明此猜想需要全新的方法. 这一现在称为 Artin 猜想的命题, 是数论中最重要的猜想之一. 在定理 1.8 之后, 一维表示的 Artin $L$-函数是 Weber $L$-函数, 故 Artin 猜想由 Weber 的结果推出. 一般地, 如 Langlands 所解释的, 不可约表示的 Artin $L$-函数应是尖形式自守 $L$-函数, 而后者已知是全纯的. 对二维表示, Artin 猜想大体上已被证明(Artin, Langlands, Tunnell, Taylor 等). 参见 Taylor, Richard, \emph{Pacific J. Math.} 181, 1997, 337–347, 及其后续论文.

对函数域, Weil 在证明有限域上曲线的 Riemann 假设的同时证明了 Artin 猜想, 并认为这是更重要的结果.

\subsection*{局部类域论与无限扩张}\index{juubu@局部类域论(local class field theory)}\index{wuxiankuozhang@无限扩张(infinite extension)}

$K$ 的 Abel 扩张 $L$ 定义 $K$ 在素点 $v$ 处的完备化 $K_{v}$ 的一个 Abel 扩张 $L^{v}/K_{v}$. 反之, $K_{v}$ 的每个有限扩张都来自 $K$ 的同一个次数的扩张(ANT, 7.62). 因此, 数域的 Abel 扩张的分类中包含局部域的 Abel 扩张的分类, 这并不令人意外. Hasse (1930)\footnote{H. Hasse, \emph{Die Normenresttheorie relative-Abelscher Zahlkörper als Klassenkörper im Kleinen}, J. für Mathematik (Crelle) 162 (1930), 145–154.} 明确表述了这一点.

\begin{theorem}\label{thm:0.9}
设 $K$ 是 $\mathbb{Q}_{p}$ 的有限扩张. 对 $K$ 的每个有限 Abel 扩张 $L$, $\Nm(L^{\times})$ 是 $K^{\times}$ 中有限指数的子群, 且映射 $L \mapsto \Nm(L^{\times})$ 是从 $K$ 的有限 Abel 扩张集到 $K^{\times}$ 的有限指数子群集上的一个反序双射.
\end{theorem}

于是到 1930 年, 数域与局部域的 Abel 扩张都已被分类. 然而, 该理论有三方面被认为不够令人满意.
\begin{itemize}
\item 局部域比整体域简单, 人们期望先局部地证明一个命题, 再整体地推出, 而非相反. 至少, 应当有一个纯局部的局部类域论证明.
\item 类域论的主要陈述是代数性的, 故它们应容许一个纯代数证明.
\item 定理\ref{thm:0.8} 分类了导子整除某个固定模的、数域 $K$ 的 Abel 扩张. 要包含其他 Abel 扩张, 就不得不不断放大导子. 这就提出了如下问题: 同时给出所有 Abel 扩张的一个统一描述, 或者说, 无限 Abel 扩张的类域论.
\end{itemize}

Chevalley 解决了所有这三个问题. 令 $K^{\mathrm{ab}}$ 表示 $K$ 的所有有限 Abel 扩张的复合. 局部类域论的完满陈述是: 对每个 $p$-进域 $K$, 存在定义好的同态 $\phi \colon K^{\times} \to \Gal(K^{\mathrm{ab}}/K)$(现称局部 Artin 映射), 它对每个有限 Abel 扩张 $L$ 诱导同构
\[
K^{\times}/\Nm(L^{\times}) \to \Gal(L/K),
\]
并且所有(开)有限指数子群都是范群. 这表明有可能定义一个整体 Artin 映射, 其各局部分量为局部 Artin 映射, 即如下图表对所有 $K$ 的素点 $v$ 交换:
\[
\begin{tikzcd}
K_{v}^{\times}\ar[r,"\phi_{v}"]\ar[d] & \Gal(K_{v}^{\mathrm{ab}}/K_{v})\ar[d]\\
\prod_{v} K_{v}^{\times}\ar[r] & \Gal(K^{\mathrm{ab}}/K).
\end{tikzcd}
\]
这个问题在于 $\prod_{v} K_{v}^{\times}$ 不是局部紧的. 紧群的乘积是紧的, 但局部紧群不满足类似结论, 而群 $K_{v}^{\times}$ 只是局部紧的. 事实上 $\prod_{v} K_{v}^{\times}$ 太大, 以致不存在这样的 $\phi$. Chevalley 通过定义 $K$ 的\emph{伊代}群 $\mathbb{I}_{K}$ 解决了这个问题: $\mathbb{I}_{K}$ 是 $\prod_{v} K_{v}^{\times}$ 中由满足对几乎所有非阿基米德 素点有 $a_{v} \in \mathcal{O}_{v}^{\times}$ 的族 $(a_{v})$ 组成的子群. 赋予使 $\prod_{v\mid\infty} K_{v}^{\times} \times \prod_{v\ \text{有限}} \mathcal{O}_{v}^{\times}$ 为开子群之拓扑时, $\mathbb{I}_{K}$ 成为局部紧群. 把 $K^{\times}$ 作为对角线子群嵌入 $\mathbb{I}_{K}$.

在 Chevalley 的类域论进路中, 先直接证明局部类域论. 然后定义一个整体 Artin 映射 $\phi_{K} \colon \mathbb{I}_{K} \to \Gal(K^{\mathrm{ab}}/K)$, 其各局部分量为局部 Artin 映射, 并证明 $K^{\times}$ 包含在 $\phi_{K}$ 的核中, 且 $\phi$ 诱导的同态
\[
\mathbb{I}_{K}/K^{\times} \to \Gal(K^{\mathrm{ab}}/K)
\]
满射, 其核等于 $\mathbb{I}_{K}/i(K^{\times})$ 的恒等连通分支. 换言之, 同态 $\phi_{K}$ 适合如下交换图表
\[
\begin{tikzcd}
K_{v}^{\times}\ar[r,"\phi_{v}"]\ar[d] & \Gal(K_{v}^{\mathrm{ab}}/K_{v})\ar[d]\\
\mathbb{I}_{K}\ar[r,"\phi_{K}"] & \Gal(K^{\mathrm{ab}}/K)
\end{tikzcd}
\]
对所有 $K$ 的素点 $v$ 成立, 且 $\phi_{K}$ 诱导同构
\[
\mathbf{C}_{K}/\mathbf{C}_{K}^{\circ} \to \Gal(K^{\mathrm{ab}}/K), \quad \mathbf{C}_{K} \overset{\mathrm{def}}{=} \mathbb{I}_{K}/K^{\times}.
\]
此陈述与定理\ref{thm:0.9} 的关系如下: 存在典范同构 $C/C^{\circ} \simeq \varprojlim_{\mathfrak{m}} C_{\mathfrak{m}}$, 并且对 $K$ 的每个有限 Abel 扩张 $L$ 与充分大的模 $\mathfrak{m}$, 有交换图表
\[
\begin{tikzcd}
\mathbf{C}\ar[r,"\phi_{K}"]\ar[d] & \Gal(K^{\mathrm{ab}}/K)\ar[d,"\tau \mapsto \tau|_{L}"]\\
C_{\mathfrak{m}}\ar[r,"\psi_{L/K}"] & \Gal(L/K).
\end{tikzcd}
\]
到 1950 年, 局部 Artin 映射已能在局部被刻画, 或被视为整体 Artin 映射的局部分支, 但直到 1965 年 Lubin 与 Tate 才找到它的一个显式局部描述.

\begin{NOTES}
在数域情形, Artin 映射 $C \to \Gal(K^{\mathrm{ab}}/K)$ 满射, 其核为 $C$ 的恒等分支. 在函数域情形, 它是单射但不满射. Weil 曾猜想, 把数域的整个伊代类群解释为 Galois 群, 可能是证明 Riemann 假设的关键(见本书 p.\pageref{pg:orig180}).
\end{NOTES}

\subsection*{中心单代数}\index{zhongxin@中心单代数(central simple algebra)}\index{Brauer 群@Brauer 群(Brauer group)}

域 $K$ 上的中心代数是指中心为 $K$ 且在 $K$ 上有限维的代数 $A$; 若 $A$ 没有非零真双侧理想, 则称其为单的. 例如, 矩阵代数 $M_{n}(K)$ 是 $K$ 上的中心单代数. 第一个有趣的中心单代数例子, 即 $\mathbb{R}$ 上的四元数代数, 由 Hamilton 于 1843 年发现. 它在 $\mathbb{R}$ 上由 $1, i, j, ij$ 张成, 其乘法由
\[
i^{2} = -1, \quad j^{2} = -1, \quad ij = -ji.
\]
决定.

设 $K$ 为域. 1906 年, Dickson 说明如何把 $K$ 的一个循环扩张\footnote{即, 具有循环 Galois 群的 Galois 扩张.} $L/K$、$\Gal(L/K)$ 的一个生成元 $\sigma$ 以及 $K$ 的一个非零元 $a$ 附到一个中心单代数 $(L/K, \sigma, a)$ 上. 作为 $L$-向量空间, 该代数由元素 $1, \alpha, \ldots, \alpha^{n-1}$ 张成, 其中 $n = [L:K]$, 其乘法由
\[
\alpha^{n} = a, \quad \alpha c = (\sigma c)\alpha \quad (\text{对 } c \in L)
\]
决定.

于是 $(L/K, \sigma, a)$ 是一个包含 $L$ 作为极大子域的中心单代数, 其中 $\sigma$ 在 $L$ 上的作用由某个内自同构诱导.

Brauer 与 Noether 说明了如何用因子系描述中心单代数, 一旦定义了群上同调, 这些因子系即被认出是 2-上循环.

1907 年, Wedderburn 证明 $k$ 上的每个中心单代数都同构于某个中心可除代数上的矩阵代数. Brauer 定义 $K$ 上的两个中心单代数相似, 如果它们同构于同一个可除代数上的矩阵代数, 并证明相似类在张量积下构成一个群(现称 $K$ 的 Brauer 群). 于是, 分类域 $K$ 上的中心单代数的问题变成了确定该域的 Brauer 群 $\mathrm{Br}(K)$.

Frobenius 的结果确定了 $\mathbb{R}$ 的 Brauer 群, Hasse 确定了非阿基米德局部域的 Brauer 群. 1932 年, Albert, Brauer, Hasse 与 Noether 证明数域上的每个可除代数都是循环的, 并确定了这样的域的 Brauer 群. 这是一个与类域论基本定理同样深刻的定理, 在这些讲义中将与它们同时证明.

Brauer 群有许多应用, 例如对有限群的表示理论(这正是 Brauer 所感兴趣的), 以及对数域上半单代数群的分类.

\begin{NOTES}
数域的 Brauer 群的确定过程在 Fenster 与 Schwermer 2005\cite{FensterSchwermer2005} 以及 Roquette 2005\cite{Roquette2005} 中有精彩的叙述. 我引用前者的 MR 评论:

In February, 1931, Hasse began a joint effort with Richard Brauer and Emmy Noether to classify all division algebras over algebraic number fields. They achieved this in a paper that was published in early 1932. The young A. A. Albert was also working on the problem at the same time. He had been in correspondence with Hasse during 1930–1931 and he had come close to resolving the problem. Some of his remarks to Hasse had been helpful in the Brauer-Hasse-Noether solution. Within weeks of that event, Albert found an alternate proof, which utilized a suggestion from Hasse.
\end{NOTES}

\subsection*{上同调}\index{shangtong@上同调(cohomology)}

对群 $G$ 与 $G$-模 $M$, 有同调群 $H_{r}(G, M)$($r \geq 0$)与上同调群 $H^{r}(G, M)$($r \geq 0$). $G$-模的短正合列给出同调与上同调群的长正合列. Tate 证明, 当 $G$ 有限时, 这些长正合列可以拼接成一条很长的序列. 更确切地说, 对所有 $r \in \mathbb{Z}$ 定义的群
\[
H_{T}^{r}(G, M) \overset{\mathrm{def}}{=}
\begin{cases}
H^{r}(G, M), & r > 0, \\
M^{G}/\Nm(M), & r = 0, \\
\operatorname{Ker}(\Nm)/I_{G}M, & r = -1, \\
H_{-r-1}(G, M), & r < -1,
\end{cases}
\]
具有此性质. 特别地,
\[
H_{T}^{-2}(G, \mathbb{Z}) \overset{\mathrm{def}}{=} H_{1}(G, \mathbb{Z}) \simeq G^{\mathrm{ab}} \quad (\text{为 } G \text{ 的最大 Abel 商}).
\]
设 $C$ 是一个 $G$-模, 使得对所有 $G$ 的子群 $H$ 都有 $H^{2}(H, C)$ 是阶为 $|H|$ 的循环群. 若 $H^{1}(G, C) = 0$, 则 Tate 证明选取 $H^{2}(G, C)$ 的一个生成元确定了同构族
\[
H_{T}^{r}(G, M) \to H_{T}^{r+2}(G, M \otimes C), \quad \text{所有 } r \in \mathbb{Z}.
\]
对局部域的 Galois 扩张 $L/K$(Galois 群为 $G$), 计算 $K$ 的 Brauer 群表明 $H^{2}(G, L^{\times})$ 是阶为 $|G|$ 的循环群, 并且有一个典范生成元. 由于 Hilbert 定理 90 给出 $H^{1}(G, L^{\times}) = 0$, Tate 定理给了我们典范同构
\[
H_{T}^{r}(G, \mathbb{Z}) \to H_{T}^{r+2}(G, L^{\times})
\]
对所有 $r \in \mathbb{Z}$. 当 $r = -2$ 时, 这成为同构
\[
G^{\mathrm{ab}} \xrightarrow{\simeq} K^{\times}/\Nm(L^{\times}),
\]
其逆就是局部 Artin 映射. 对整体域有类似的论证, 此时以伊代模去主伊代群代替域的乘法群.

\begin{NOTES}
我引用 Mac Lane 1978\cite{MacLane1978}, p. 17:

\begin{quote}
Mac Lane 回忆说, Artin(约 1948 年)在谈话中指出, 群的上同调本应在类域论中发挥作用. Hochschild (1950) 与 Hochschild 和 Nakayama (1952) 说明了类域论的 Brauer 群论证如何能被上同调论证取代. 1952 年, Tate 证明有限群 $G$ 的同调群与上同调群可以恰当地合并为一条长正合列. 他利用此序列, 连同转移与限制的性质, 给出类域论的一个优美重构.
\end{quote}
\end{NOTES}

\subsection*{类域的显式构造(Hilbert 第十二问题)}\index{hilbert12@Hilbert 第十二问题(Hilbert's 12th problem)}\index{xianshi@显式类域论(explicit class field theory)}

Kronecker--Weber 定理表明 $\mathbb{Q}$ 的每个 Abel 扩张都包含于由指数函数的一个特殊值 $\zeta = e^{2\pi i/m}$ 生成的域中. 稍后的一个定理(Kronecker 青春之梦)表明, 虚二次数域的每个 Abel 扩张都包含于由椭圆模函数的某些特殊值生成的域中. 在著名的 23 个问题中的第 12 个里, Hilbert 问: 其他数域的 Abel 扩张能否由显式全纯函数的特殊值生成.

推广椭圆模函数的函数已被充分了解, 其特殊值的算术性质亦然——这就是复乘理论与 Shimura 簇理论——但它们只对构造 CM 域的 Abel 扩张有用, 并且即使对这些域, 它们也只给出(粗略地说)不从全实子域的 Abel 扩张来的那些 Abel 扩张(除 $\mathbb{Q}$ 外). 另一方面, Stark 的猜想为全实域的 Abel 扩张的生成元提供了候选.

\subsection*{显式类域论}

类域论中定理的证明优美但抽象. 例如, 局部类域论最初的局部证明只表明存在唯一的 Artin 映射 $\phi \colon K^{\times} \to \Gal(K^{\mathrm{ab}}/K)$ 具有某些性质, 却未解决该映射与 $K^{\mathrm{ab}}$ 的显式描述问题. 所幸, 在这种情况下, 我们将在第 \ref{ch:I} 章解释的 Lubin 与 Tate 的理论对此问题给出了优美的回答.

已有大量工作致力于寻找类域论中其他映射与对象的显式描述. 在当前版本的讲义中, 这一主题大体被忽略.

\subsection*{计算类域论}\index{jisuan@计算类域论(computational class field theory)}

除了寻找类域论中映射与对象的显式描述外, 还可以要求计算它们的算法. 对此主题, 我只能指引读者参阅 Cohen 2000\cite{Cohen2000}.

\subsection*{非 Abel 类域论}\index{fiabel@非 Abel 类域论(nonabelian class field theory)}

在 1920 年 ICM 上他向世人宣布 Abel 类域论基本定理证明的同一报告中, Takagi 提出了非 Abel 类域论的问题:
\begin{quote}
En m'arrêtant ici, je me permets d'attirer votre attention sur un problème important de la théorie des nombres algébriques: à savoir, rechercher s'il est possible de définir la classe d'idéaux d'un corps algébrique de telle manière que les corps supérieur relativement normal mais non abélien puisse être caractérisé par le groupe correspondant de classes d'idéaux du corps de fond.\footnote{意即: 在此停笔, 请允许我提请诸位注意代数数论中的一个重要问题, 即能否以这样的方式定义代数数域的理想类, 使得相对正规但非 Abel 的扩张域可由基域的相应理想类群刻画.}
\end{quote}
范限制定理表明, (射线)类群的子群不能区分一个扩张域与其最大的 Abel 子扩张.

几十年来, 非 Abel 类域论应采取何种形式, 甚至它是否存在, 都不清楚. 1946 年, Artin 猜测找到正确的陈述是唯一的问题: 一旦我们知道它们是什么, 就能由 Abel 类域论推出它们(\emph{A Century of Mathematics in America}, Part II, (Peter Duren 编), 1989, p. 312). Weil 转述说, 一年后 Artin 表示他已对存在非 Abel 类域论失去信心(Weil, A., \emph{Œuvres}, Vol. III, p. 457.)

与其研究集合 $\operatorname{Spl}_{S}(L/K)$, 不如研究 $\Gal(L/K)$ 的表示 $\rho$ 的 Artin $L$-级数 $L(s, \rho)$. 描述集合 $\operatorname{Spl}_{S}(L/\mathbb{Q})$ 的问题于是变成描述以这种方式出现的解析函数集的问题. Langlands 构造了一类 $L$-级数, 称为自守 $L$-级数, 并猜想\footnote{注意它与谷山猜想的相似之处——事实上二者都是一个更广泛猜想的特例.} 每个 $L_{S}(s, \rho)$ 都是自守的, 并且指明了哪些自守 $L$-级数以这种方式出现. 于是, 此猜想对所有有限 Galois 扩张 $L/\mathbb{Q}$ 回答了原问题. 对 $n = 1$(即 $G$ Abel)与所有 $K$, Artin 证明所有 Artin $L$-级数都是自守的——这是他证明定理\ref{thm:0.8} 的动机. 如前所述, 对 $n = 2$, 该猜想在大多数情形已被证明.

1967 年, Langlands 陈述了他猜想的\footnote{见 1967 年致 Weil 的一封信, 以及后来(对我们其余人)的 \emph{Problems in the theory of automorphic forms} (1970). 对这些著作的一个引人入胜的介绍, 见 Casselman 2001\cite{Casselman2001}. Langlands 著作(包括上述两篇)的 TeX 版本可见于 \url{http://publications.ias.edu/rpl/}.} 函子性原理, 它把非 Abel 类域论作为一个特例包含在内. 对局部域 $K$, 可表述如下. $K$ 的 Weil 群 $W_{K}$ 定义为 $\Gal(K^{\mathrm{al}}/K)$ 中那些在剩余类域上作用为 Frobenius 元的整数次幂的元素组成的子群. Abel 局部类域论中的局部 Artin 映射可视为从 $K^{\times}$ 到 $W_{K}$ 的最大 Abel 商 $W_{K}^{\mathrm{ab}}$ 的同构 $\phi_{K}$. Langlands 猜想从 $W_{K}$ 到 $\mathrm{GL}_{n}(\mathbb{C})$ 的同态对应于 $\mathrm{GL}_{n}(K)$ 的某些表示. 对 $n = 1$, $\mathrm{GL}_{1}(K) = K^{\times}$ 的表示恰为特征, 并且对应由与 $\phi_{K}$ 的复合给出. 对 $n > 1$, $\mathrm{GL}_{n}(K)$ 的表示通常是无限维的.

在自守一侧, 设 $\mathcal{A}_{n}(K)$ 为 $\mathrm{GL}_{n}(K)$ 在复向量空间上的、每个向量的稳定子开的不可约表示的等价类集. 在 Galois 侧, 设 $\mathcal{G}_{n}(K)$ 为对 $(r, N)$ 的等价类集, 其中 $r$ 是 $W_{K}$ 在 $n$ 维复向量空间 $V$ 上的半单表示, 在某个开子群上平凡, $N$ 是 $V$ 的一个幂零自同态, 使得用 $r(\sigma)$($\sigma \in W_{K}$)共轭 $N$ 会乘上 $\phi_{K}^{-1}(\sigma)$ 的绝对值. $K$ 的局部 Langlands 猜想断言存在双射族 $(\sigma_{n})_{n \geq 1}$,
\[
\pi \mapsto \sigma_{n}(\pi) \colon \mathcal{A}_{n}(K) \to \mathcal{G}_{n}(K),
\]
使得
\begin{enumerate}[(a)]
\item $\sigma_{n}(\pi)$ 的行列式视为 $W_{K}$ 的特征时, 经 $\phi_{K}$ 对应于 $\pi$ 的中心特征;
\item 映射 $\sigma_{n}$ 保持各对表示的 $L$-因子与 $\varepsilon$-因子(自守侧由 Jacquet, Piatetskii-Shapiro 与 Shalika 定义, Galois 侧由 Langlands 与 Deligne 定义);
\item 对 $\chi \in \mathcal{A}_{1}(K)$, $\sigma_{n}(\pi \otimes (\chi \circ \det)) = \sigma_{n}(\pi) \otimes \sigma_{1}(\chi)$;
\item $\sigma_{n}$ 与取对偶 $\pi \mapsto \pi^{\vee}$ 交换.
\end{enumerate}
对每个 $K$, Henniart 证明这样的族至多存在一个. 该猜想本身由 Harris 与 Taylor (2001)\cite{HarrisTaylor2001} 证明. 几个月后, Henniart (2000)\cite{Henniart2000} 找到了一个更简单的证明.

\subsection*{最后的注记}\index{zhuihou@最后的注记(final remark)}

回忆 Dedekind 的如下定理(ANT, 3.41).

\begin{theorem}\label{thm:0.10}
设 $A$ 是具有分式域 $K$ 的 Dedekind 整环, $B$ 为 $K$ 在有限可分扩张 $L$ 中的整闭包. 设 $L = K[\alpha]$, $\alpha \in B$, 并设 $f(X)$ 为 $\alpha$ 在 $K$ 上的极小多项式. 则下列关于 $A$ 的素理想 $\mathfrak{p}$ 的条件等价:
\begin{enumerate}[(a)]
\item $\mathfrak{p}$ 不整除 $\operatorname{disc}(f(X))$;
\item $\mathfrak{p}$ 在 $B$ 中不分歧, 且 $A_{\mathfrak{p}}[\alpha]$ 是 $A_{\mathfrak{p}}$ 在 $L$ 中的整闭包(这里 $A_{\mathfrak{p}} = \{a/b \mid b \notin \mathfrak{p}\}$);
\item 存在分解
\[
f(X) \equiv f_{1}(X) \cdots f_{g}(X) \mod \mathfrak{p}
\]
其中诸 $f_{i}$ 互异、首一且在模 $\mathfrak{p}$ 下不可约.
\end{enumerate}
\end{theorem}

当这些条件成立时, $\mathfrak{p}$ 在 $L$ 中分解为素理想的乘积为
\[
\mathfrak{p}B = (\mathfrak{p}, f_{1}(\alpha)) \cdots (\mathfrak{p}, f_{g}(\alpha)).
\]
于是类域论实际上关于系数在数域中的一元多项式及其根: $K$ 的 Abel 扩张对应于首一不可约多项式 $f(X) \in K[X]$, 使得给出域自同构的 $f(X)$ 的根的置换可交换; Dedekind 定理表明, 在 Abel 扩张 $L$ 中, $K$ 的除了有限多个外的所有素理想的分解对应于有限域上多项式的分解.

\subsection*{本讲义采用的方法}\index{benjiangyi@本讲义采用的方法(approach taken in these notes)}

在这些讲义中, 我们遵循 Chevalley (1940, 1954)\cite{Chevalley1940}\cite{Chevalley1954}、Artin 与 Tate (1961)\cite{ArtinTate1961} 以及 Tate (1952, 1967)\cite{Tate1952}\cite{Tate1967} 的代数/上同调进路, 证明局部与整体类域论的基本定理. 上同调进路清楚地表明了类域论基本定理与 Brauer 群计算之间的关系. 事后看来, 自 Takagi 的论文以来, 上同调计算就隐含在类域论之中.

然而, 类域论还有其他进路, 当能提供额外洞见时, 我不惜给出多于一种证明. 例如, 我说明如何借助分析简化整体情形基本定理证明的某些步骤. 此外, 我说明(遵循 Lubin 与 Tate)形式群律理论如何能使局部类域论的结果更显式、更令人满意.

尽管可以既不用分析也不用上同调来证明类域论的基本定理, 但有些主要定理甚至离不开其中之一才能陈述, 例如关于素数密度的定理, 或关于数域的上同调群的定理. 近些年来, 上同调结果在类域论的许多应用中至关重要.\footnote{例如, 据 Google Scholar, 我的《算术对偶定理》一书已被 1014 篇研究论文引用(自其开始统计以来). MR 列出 430 篇(约自 2000 年起).}

课程的核心是奇数编号的章节. 第 \ref{ch:II} 章(论群的上同调)是其余各章的基础, 但第 \ref{ch:IV} 章、第 \ref{ch:VI} 章与第 \ref{ch:VIII} 章并非阅读第 \ref{ch:III}、\ref{ch:V}、\ref{ch:VII} 章所必需. 除第一节外, 第 \ref{ch:I} 章可被那些希望尽快到达整体类域论基本定理的人跳过.

\begin{exercise}\label{ex:0.11}
完成证明: 二次互反律允许描述 $L/\mathbb{Q}$ 为二次扩张时的集合 $\operatorname{Spl}_{S}(L/\mathbb{Q})$.
\end{exercise}

\begin{exercise}\label{ex:0.12}
证明 $\mathbb{Q}[\sqrt{-5}, \sqrt{-1}]$ 是 $\mathbb{Q}[\sqrt{-5}]$ 的 Hilbert 类域.
\end{exercise}

\begin{exercise}\label{ex:0.13}
证明 本书 p.\pageref{pg:orig4} 上定义的映射 $I^{S} \to (\mathbb{Z}/m\mathbb{Z})^{\times}$ 具有所描述的核, 从而诱导同构 $C_{\mathfrak{m}} \to (\mathbb{Z}/m\mathbb{Z})^{\times}$.
\end{exercise}

\begin{exercise}\label{ex:0.14}
证明例 \ref{ex:0.6} 中的陈述.
\end{exercise}

\begin{exercise}\label{ex:0.15}
设 $L = \mathbb{Q}[\sqrt{-1}, \sqrt{-5}]$. 则 $\Gal(L/\mathbb{Q}) = \{1, \sigma, \tau, \sigma\tau\} \simeq \mathbb{Z}/2\mathbb{Z} \times \mathbb{Z}/2\mathbb{Z}$, 其中 $\sigma$ 固定 $\mathbb{Q}[\sqrt{-1}]$, $\tau$ 固定 $\mathbb{Q}[\sqrt{-5}]$, $\sigma\tau$ 固定 $\mathbb{Q}[\sqrt{5}]$.
\begin{enumerate}[(a)]
\item 证明在 $L$ 中仅 $2, 5, \infty$ 分歧.
\item 对所有 $p \neq 2, 5$, 计算 $(p, L/\mathbb{Q})$.
\item 设 $\mathfrak{m} = (20)\infty$. 证明对某个 $H \subset C_{\mathfrak{m}}$, 映射 $p \mapsto (p, L/\mathbb{Q})$ 定义同构 $C_{\mathfrak{m}}/H \to \Gal(L/\mathbb{Q})$, 并求 $H$.
\end{enumerate}
提示: 证明 $L \subset \mathbb{Q}[\zeta]$, 其中 $\zeta$ 为本原 20 次单位根.
\end{exercise}

我将给你讲一个关于互反律的故事. 在我完成博士论文后, 我有了为非 Abel 扩张定义 $L$-级数的想法. 但要使它们与 Abel 扩张的 $L$-级数一致, 某个同构必须成立. 我能证明它蕴含所有标准的互反律. 于是我称之为一般互反律, 并试图证明它, 却做不到, 即便尝试多次. 然后我把它拿给其他数论学家看, 但他们都嘲笑它, 我记得尤其是 Hasse 对我说这绝不可能是真的. 尽管如此, 我仍坚持, 但毫无进展. 连续三年, 没有一周我不去尝试证明互反律. 这令人气馁, 其间我也转向其他事情. 后来有一天下午我没什么特别的事做, 于是说: “好吧, 我再试一次证明互反律.” 于是我走出去, 在花园里坐下. 你看, 从一开始我就有用分圆域的想法, 但它们从未奏效, 而此刻我突然明白, 一直以来我用错了它们——半小时内我就证出来了.

Emil Artin, 由 Mattuck 回忆(见 \emph{Recountings: Conversations with MIT Mathematicians} 2009).

为何其他数学家如此难以相信它? 当时的互反律与幂剩余符号紧密相连. 实际上 Euler 已以正确的方式表述了二次互反律(即符号 $(\Delta/p)$ 只依赖于 $p$ 模 $\Delta$ 的剩余类), 但 Legendre 的表述占了上风. 我很确信, 如果 Artin 立刻向 Hasse 表明一般互反律等价于已知的幂剩余互反律, Hasse 就不会笑. Hasse 不必等 Chebotarev 的密度定律来说服自己 Artin 是对的——到 Chebotarev 的想法出现时, Artin 必已无疑. 顺便说, Chebotarev 的文章做的远不止证明 Frobenius 元的重要性——它为 Artin 提供了证明其互反律的关键思想, 即 Hilbert 的 Abel 交叉技术.

让我再补充一点: Felix Bernstein 于 1904 年猜想了一个互反律, 在非分歧 Abel 扩张的特殊情形大体等价于 Artin 的定律. 它的技术性表明, 从已知的幂剩余互反律猜出像 Artin 那样简单的定律绝非易事.\footnote{见 Franz Lemmermeyer 在 \mo{243590} 上的讨论.}

% 恢复正文章节定理编号(供后续各章使用)
\renewcommand{\thetheorem}{\thesection.\arabic{theorem}}


% ===== 以下各章将在后续翻译阶段依次加入 =====
% ============ 第一章 I 局部类域论: Lubin-Tate 理论 ============
% 原书 Chapter I, pp.19--52 (局部类域论: Lubin-Tate 理论)
% 编号: 定理/推论/引理/命题/定义/例/注记/习题 共享计数器, 按节编号(1.1, 1.2, ...)
% 节内说明性编号 1.6--1.8, 1.11, 1.12, 1.14, 1.15 亦占该计数器;
% 其中 1.9, 1.10 为子节标题, 不占计数器, 故在 1.11 前将计数器置为 10.
\chapter{局部类域论: Lubin-Tate 理论}\label{ch:I}%
\index{jubuleiyulun@局部类域论(local class field theory)}\index{lubintatexingshiqunlü@Lubin-Tate 形式群律(Lubin-Tate formal group)}

局部类域论对局部域的 Abel 扩张进行分类. 换一个角度, 它描述了整体 Artin 映射的局部分量.

所谓局部域, 我是指对一个非平凡绝对值而言局部紧的域 $K$. 于是 (ANT, 7.49) 它或者是

(a) 对某个 $p$, $\Q_p$ 的有限扩张;

(b) 具有 $p$ 个元素的域上的 Laurent 级数域 $\F_p((T))$ 的有限扩张; 或者

(c) $\R$ 或 $\CC$ (阿基米德情形).

当 $K$ 非阿基米德 时, $K$ 中的整数环 (又称赋值环) 记为 $\mathcal{O}_K$ (或 $A$), 其极大理想记为 $\mathfrak{m}_K$ (或简记 $\mathfrak{m}$), 其单位群记为 $\mathcal{O}_K^{\times}$ 或 $U_K$. $\mathfrak{m}$ 的一个生成元称为 $K$ 的素元 (或称一致素元, 或局部一致化参数). 若 $\pi$ 是 $K$ 的素元, 则 $K^{\times}$ 中每个元素都可唯一写成 $a=u\pi^m$ 的形式, 其中 $u\in\mathcal{O}_K^{\times}$, $m\in\Z$. 我们定义
\[
\mathrm{ord}_K(a)=m.
\]
存在一个正合列
\[
0\to U_K\to K^{\times}\xrightarrow{\mathrm{ord}_K}\Z\to 0.
\]
$K$ 的剩余类域 $k$ 有 $q$ 个元素, 其特征为 $p$. $K$ 上的标准绝对值定义为 $|a|=q^{-\mathrm{ord}_K(a)}$.

我们以 $K^{\mathrm{al}}$ 记 $K$ 的一个固定的代数闭包 (当 $K$ 的特征 $p\neq 0$ 时, 指可分代数闭包), “$K$ 的扩张”意即 “含 $K$ 的 $K^{\mathrm{al}}$ 的子域”. $\mathrm{ord}_K$ 与 $|\cdot|$ 都可唯一扩张到 $K^{\mathrm{al}}$ ($\mathrm{ord}_K$ 的扩张取值于 $\Q$).

\section{主定理的陈述}\label{sec:I.1}

两个 $K$ 的有限 Abel 扩张的复合仍是 $K$ 的有限 Abel 扩张 (FT, 3.20). 因此所有有限 Abel 扩张在 $K^{\mathrm{al}}$ 中的并 $K^{\mathrm{ab}}$ 是一个无限 Abel 扩张, 其 Galois 群是 $\Gal(K^{\mathrm{al}}/K)$ 模去其换位子群之闭包所得的商群. (见 FT 第 7 章及本章附录中关于无限扩张的 Galois 理论.)

设 $L$ 是 $K$ 的有限非分歧扩张. 则 $L/K$ 是 Galois 的, 且 $\Gal(L/K)$ 在 $\mathcal{O}_L$ 上的作用定义了一个同构 $\Gal(L/K)\simeq\Gal(l/k)$, 其中 $l$ 是 $L$ 的剩余类域. 因此 $\Gal(L/K)$ 是循环群, 由唯一的元素 $\sigma$ 生成, 该元素对所有 $\alpha\in\mathcal{O}_L$ 满足 $\sigma\alpha\equiv\alpha^q\bmod\mathfrak{m}_L$. 这个 $\sigma$ 称为 $\Gal(L/K)$ 的 \emph{Frobenius 元素}, 记为 $\Frob_{L/K}$. 见 ANT, 7.54.

\begin{theorem}[局部互反律]\label{thm:1.1}
对每个非阿基米德局部域 $K$, 存在唯一的同态
\[
\phi_K\colon K^{\times}\to\Gal(K^{\mathrm{ab}}/K)
\]
满足以下性质:

(a) 对 $K$ 的每个素元 $\pi$ 以及 $K$ 的每个有限非分歧扩张 $L$, $\phi_K(\pi)$ 在 $L$ 上的作用等于 $\Frob_{L/K}$;

(b) 对 $K$ 的每个有限 Abel 扩张 $L$, $\Nm_{L/K}(L^{\times})$ 包含在 $a\mapsto\phi_K(a)|_L$ 的核中, 且 $\phi_K$ 诱导出一个同构
\[
\phi_{L/K}\colon K^{\times}/\Nm_{L/K}(L^{\times})\to\Gal(L/K).
\]
特别地,
\[
(K^{\times}:\Nm_{L/K}(L^{\times}))=[L:K].
\]
\end{theorem}

记 $\Nm_{L/K}(L^{\times})$ 为 $\Nm(L^{\times})$. 陈述 (b) 表明, 对每个有限 Abel 扩张 $L/K$, 映射 $\phi_K$ 可分解为
\[
\begin{tikzcd}
K^{\times}\ar[r,"\phi_K"]\ar[d] & \Gal(K^{\mathrm{ab}}/K)\ar[d,"\tau\mapsto\tau|_L"]\\
K^{\times}/\Nm(L^{\times})\ar[r,"\simeq","\phi_{L/K}"'] & \Gal(L/K).
\end{tikzcd}
\]
我们称 $\phi_K$ 与 $\phi_{L/K}$ 为 $K$ 与 $L/K$ 的\emph{局部 Artin 映射}. 它们也常被称为\emph{局部互反律映射}, 记为 $\mathrm{rec}_K$ 与 $\mathrm{rec}_{L/K}$; 而 $\phi_{L/K}$ 也常称为\emph{范数剩余映射或范数剩余符号}, 记作 $a\mapsto(a,L/K)$.

形如 $\Nm(L^{\times})$ ($L$ 跑遍 $K$ 的有限 Abel 扩张) 的 $K^{\times}$ 的子群称为 $K^{\times}$ 中的\emph{范数群}\index{fanshuqun@范数群(norm group)}.

\begin{corollary}\label{cor:1.2}
设 $K$ 为非阿基米德局部域, 并假设存在同态 $\phi\colon K^{\times}\to\Gal(K^{\mathrm{ab}}/K)$ 满足定理的条件 (a) 与 (b).

(a) 映射 $L\mapsto\Nm(L^{\times})$ 是从 $K$ 的有限 Abel 扩张的集合到 $K^{\times}$ 中范数群的集合上的双射.

(b) $L\subset L'\iff \Nm(L^{\times})\supset\Nm(L'^{\times})$.

(c) $\Nm((L\cdot L')^{\times})=\Nm(L^{\times})\cap\Nm(L'^{\times})$.

(d) $\Nm((L\cap L')^{\times})=\Nm(L^{\times})\cdot\Nm(L'^{\times})$.\erratum{原书此处记作 $\Nm((L\cap L)^{\times})$, 依 Milne errata (v4.03) 应理解为 $L\cap L'$.}

(e) 每个包含某个范数群的 $K^{\times}$ 的子群自身也是范数群.
\end{corollary}

\begin{proof}
注意范次的传递性
\[
\Nm_{L'/K}=\Nm_{L/K}\circ\Nm_{L'/L}
\]
表明 $L\subset L'\implies \Nm(L^{\times})\supset\Nm(L'^{\times})$. 因此
\[
\Nm((L\cdot L')^{\times})\subset\Nm(L^{\times})\cap\Nm(L'^{\times}).
\]
反之, 若 $a\in\Nm(L^{\times})\cap\Nm(L'^{\times})$, 则
\[
\phi_{L/K}(a)=1=\phi_{L'/K}(a).
\]
但 $\phi_{LL'/K}(a)|_L=\phi_{L/K}(a)$, 且 $\phi_{LL'/K}(a)|_{L'}=\phi_{L'/K}(a)$. 由于映射
\[
\sigma\mapsto(\sigma|_L,\sigma|_{L'})\colon\Gal(LL'/K)\to\Gal(L/K)\times\Gal(L'/K)
\]
是单的 (FT, 3.20), 这就推出 $\phi_{LL'/K}(a)=1$, 从而 $a\in\Nm((L\cdot L')^{\times})$. 这证明了 (c); 现在来完成 (b) 的证明. 当 $\Nm(L^{\times})\supset\Nm(L'^{\times})$ 时, (c) 成为
\[
\Nm((LL')^{\times})=\Nm(L'^{\times}).
\]
由于一个范数群的指数等于定义它的 Abel 扩张的次数 (\ref{thm:1.1}b), 且 $LL'\supset L'$, 这表明 $LL'=L'$, 从而 $L'\supset L$.

现在证明 (a). 按定义, 映射 $L\mapsto\Nm(L^{\times})$ 是满的, 而由 (b) 知它是单的.

接着证明 (e). 设 $N$ 是一个范数群, 比如 $N=\Nm(L^{\times})$, 并设 $I\supset N$. 令 $M$ 为 $\phi_{L/K}(I)$ 的固定域, 于是 $\phi_{L/K}$ 把 $I/N$ 同构地映到 $\Gal(L/M)$ 上. 考虑交换图表
\[
\begin{tikzcd}
K^{\times}\ar[r,"\phi_{L/K}"]\ar[d,equal] & \Gal(L/K)\ar[d]\\
K^{\times}\ar[r,"\phi_{M/K}"'] & \Gal(M/K)
\end{tikzcd}
\]
其中右侧竖向箭头为限制映射（左侧为恒等映射）. $\phi_{M/K}$ 的核是 $\Nm(M^{\times})$. 另一方面, 映射
\[
K^{\times}\to\Gal(L/K)\to\Gal(M/K)
\]
的核是 $\phi_{L/K}^{-1}(\Gal(L/M))$, 按 Galois 理论它等于 $I$.

最后证明 (d). 在 (a) 中的两个集合之间有一个反序的双射. 由于 $L\cap L'$ 是同时含于 $L$ 与 $L'$ 中的 $K$ 的最大扩张, 而 $\Nm(L^{\times})\cdot\Nm(L'^{\times})$ 是包含 $\Nm(L^{\times})$ 与 $\Nm(L'^{\times})$ 的最小子群 (并且由 (e) 知它是范数群), 两者必相对应.
\end{proof}

为了对 $K$ 的 Abel 扩张进行分类, 还需确定范数群. 下一个引理表明 (在假定定理\ref{thm:1.1} 的前提下) 它们都是开集.

\begin{lemma}\label{lem:1.3}
设 $L$ 是 $K$ 的扩张. 若 $\Nm(L^{\times})$ 在 $K^{\times}$ 中指数为有限, 则它是开集.
\end{lemma}

\begin{proof}
$L$ 中的单位群 $U_L$ 是紧的, 故 $\Nm(U_L)$ 在 $K^{\times}$ 中闭. 考察赋值可知
\[
\Nm(L^{\times})\cap U_K=\Nm(U_L),
\]
于是 $U_K/\Nm(U_L)\hookrightarrow K^{\times}/\Nm(L^{\times})$. 因此 $\Nm(U_L)$ 在 $U_K$ 中指数为有限且闭, 从而在 $U_K$ 中开; 而 $U_K$ 自身在 $K^{\times}$ 中开. 于是 $\Nm(L^{\times})$ 含有 $K^{\times}$ 的一个开子群, 故自身为开.\footnote{这里用到了: 拓扑群中有限指数的闭子群是开的 (因为其补集是有限多个陪集的并, 而陪集也是闭的); 以及: 一个子群若含有一个开子群, 则它自身为开 (因为它是该开子群若干陪集的并)。}
\end{proof}

\begin{theorem}[局部存在定理]\label{thm:1.4}
$K^{\times}$ 中的范数群恰为有限指数的开子群.
\end{theorem}

\begin{corollary}\label{cor:1.5}
推论\ref{cor:1.2} 成立, 其中 “范数群” 替换为 “有限指数的开子群”.
\end{corollary}

在概述定理\ref{thm:1.1} 与 \ref{thm:1.4} 的证明之前, 我先给出关于这两个定理的一些补充.

\begin{npara}\label{para:1.6}
推论\ref{cor:1.5} 对阿基米德局部域也成立. $\R$ 的 Abel 扩张是 $\R$ 与 $\CC$, 它们的范数子群分别是 $\R^{\times}$ 与 $\R_{>0}$. 设 $H$ 是 $\R^{\times}$ 中有限指数的子群. 则对某个 $m$ 有 $H\supset\R^{\times m}$, 而依据介值定理, $\R^{\times m}=\R^{\times}$ 或 $\R_{>0}$ 分别视 $m$ 为奇或偶而定. 因此 $\R^{\times}$ 与 $\R_{>0}$ 是 $\R^{\times}$ 中仅有的两个有限指数子群. 唯一的同构
\[
\R^{\times}/\R_{>0}\to\Gal(\CC/\R)
\]
称为 $\R$ 的局部 Artin 映射.
\end{npara}

\begin{npara}\label{para:1.7}
当 $K$ 的特征为零时, $K^{\times}$ 的每个有限指数子群都是开的. 要证明这一点, 注意一个有限指数的子群 $H$ 将包含某个 $m$ 的 $K^{\times m}$; 而对满足 $|1-a|<|m|^2$ 的每个 $a\in\mathcal{O}_K^{\times}$, Newton 引理 (ANT, 7.32) 用于 $X^m-a$ 表明 $a$ 可写成 $a=u^m$, 其中 $u\in1+\mathfrak{m}$. 于是 $H$ 包含 $K^{\times}$ 中 $1$ 的一个开邻域, 而由于它是群, 这就推出它是开的.

当 $K$ 的特征 $p\neq 0$ 时, 并非 $K^{\times}$ 的每个有限指数子群都是开的. \cite{Weil1967}, II 3, 命题 10 表明 $1+\mathfrak{m}\approx\prod_{\N}\Z_p$ (以 $\N$ 为指标的 $\Z_p$ 的拷贝之积), 由此推出 $K^{\times}$ 有一个同构于 $\prod_{\N}\F_p$ 的商群. 子群 $\bigoplus_{\N}\F_p$ 在 $\prod_{\N}\F_p$ 中稠密. 因此 $\prod_{\N}\F_p$ 中每个包含 $\bigoplus_{\N}\F_p$ 的真子群都不闭. 选取这样一个有限指数的子群, 令 $U$ 为它在 $1+\mathfrak{m}$ 中的原像. 则 $U\cdot\pi^{\Z}$ 在 $K^{\times}$ 中指数为有限但不闭 (从而不开). 参见 FT 7.26.
\end{npara}

\begin{npara}\label{para:1.8}
$K$ 的两个有限非分歧扩张的复合仍是非分歧的, 因此所有有限非分歧扩张在 $K^{\mathrm{al}}$ 中的并 $K^{\mathrm{un}}$ 是 $K$ 的一个非分歧扩张. $K^{\mathrm{un}}$ 的剩余类域 $\bar{k}$ 是 $K$ 的剩余类域 $k$ 的代数闭包.

$K^{\mathrm{un}}$ 的每个固定 $K$ 的自同构 $\sigma$ 保持 $K^{\mathrm{un}}$ 上的绝对值 $|\cdot|$, 从而诱导出 $\bar{k}/k$ 的一个自同构 $\bar{\sigma}$. 映射 $\sigma\mapsto\bar{\sigma}$ 是一个同构
\[
\Gal(K^{\mathrm{un}}/K)\to\Gal(\bar{k}/k).
\]
因此存在唯一的元素 $\Frob_K\in\Gal(K^{\mathrm{un}}/K)$ 在 $\bar{k}$ 上诱导映射 $x\mapsto x^q$, 且映射 $\alpha\mapsto\Frob_K^{\alpha}\colon\hat{\Z}\to\Gal(K^{\mathrm{un}}/K)$ 是拓扑群的同构. 定理\ref{thm:1.1} 的条件 (a) 可以重新表述为: 对 $K$ 的每个素元 $\pi$, $\phi_K(\pi)$ 在 $K^{\mathrm{un}}$ 上的作用等于 $\Frob_K$. 特别地, 对任何单位 $u\in\mathcal{O}_K^{\times}$, $\phi_K(\pi)$ 与 $\phi_K(\pi u)$ 在 $K^{\mathrm{un}}$ 上有相同的作用, 因此 $\phi_K(u)$ 在 $K^{\mathrm{un}}$ 上起恒等映射的作用.
\end{npara}

\subsection*{1.9 群}

\begin{equation*}
U_K\supset 1+\mathfrak{m}\supset\cdots\supset 1+\mathfrak{m}^n\supset\cdots
\tag{7}\label{eqn:i.1.9.filt}
\end{equation*}
构成 $U_K$ 中 $1$ 处的一组基本邻域系. 设 $L$ 是 $K$ 的有限 Abel 扩张. 若 $L/K$ 非分歧, 则由 (\ref{para:1.8}) 知 $\Ker(\phi_{L/K})\supset U_K$, 我们说 $L/K$ 的\emph{导子}\index{daozi@导子(conductor)}为 $0$. 否则, 使得 $\Ker(\phi_{L/K})$ 包含 $1+\mathfrak{m}^f$ 的最小 $f$ 称为 $L/K$ 的\emph{导子}. 当且仅当导子为 $0$ 时 $L/K$ 非分歧, 当且仅当导子 $\le 1$ 时 $L/K$ 为顺分歧 (tamely ramified). [注意对 $i\ge 1$, 商 $1+\mathfrak{m}^i/1+\mathfrak{m}^{i+1}$ 是 $p$-群; 因此若 $L/K$ 顺分歧, 则 $\phi_{L/K}|_{U_K}$ 经过 $U_K/(1+\mathfrak{m})$. 更精细的论述见 \ref{para:1.11}.]

\subsection*{1.10 同态}

\[
\phi_{L/K}\colon K^{\times}/\Nm(L^{\times})\to\Gal(L/K)
\]
当 $L$ 跑遍 $K$ 的有限 Abel 扩张 (按包含排序) 时构成一个逆向系. 取极限, 我们得到一个同构
\[
\hat{\phi}_K\colon\widehat{K^{\times}}\to\Gal(K^{\mathrm{ab}}/K),
\]
其中 $\widehat{K^{\times}}$ 是 $K^{\times}$ 关于如下拓扑的完备化: 以范数群作为 $1$ 处的基本邻域系. $K^{\times}$ 上的这个拓扑称为\emph{范数拓扑}\index{fanshutopu@范数拓扑(norm topology)}. 按定理\ref{thm:1.4}, 范数群就是 $K^{\times}$ 中有限指数的开子群. 选取一个素元 $\pi$, 并写
\begin{equation*}
K^{\times}=U_K\cdot\pi^{\Z}\simeq U_K\times\Z.
\tag{8}
\end{equation*}
子群
\[
(1+\mathfrak{m}^n)\cdot\langle\pi^m\rangle\simeq(1+\mathfrak{m}^n)\times m\Z
\]
构成 $1$ 处的一组基本邻域系. 特别地, 对范数拓扑而言 $U_K$ 在 $K^{\times}$ 中不开, 因此该拓扑比 $K^{\times}$ 上通常的拓扑更粗. 完备化后, 上式成为
\[
\widehat{K^{\times}}=U_K\cdot\pi^{\hat{\Z}}\simeq U_K\times\hat{\Z},
\]
其中 $\hat{\Z}$ 是 $\Z$ 关于有限指数子群所定义拓扑的完备化. 更规范地说, 恒等映射
\[
K^{\times}(\text{通常拓扑})\to K^{\times}(\text{范数拓扑})
\]
连续, 并在完备化上诱导一个同态 $K^{\times}\to\widehat{K^{\times}}$, 它落在一个交换图表中:
\[
\begin{tikzcd}[column sep=small]
0\ar[r] & U_K\ar[r]\ar[d,"\simeq"] & K^{\times}\ar[r]\ar[d] & \Z\ar[r]\ar[d,hook] & 0\\
0\ar[r] & U_K\ar[r] & \widehat{K^{\times}}\ar[r] & \hat{\Z}\ar[r] & 0
\end{tikzcd}
\]
粗略地说, $\widehat{K^{\times}}$ 是由把 $K^{\times}$ 中的 $\Z$ 换成 $\hat{\Z}$ 而得到的.

\setcounter{theorem}{10} % 1.9, 1.10 为子节标题, 不占定理计数器; 此处将计数器调至 10, 使下一条显示为 1.11
\begin{npara}\label{para:1.11}
选取一个素元 $\pi$ 决定了一个分解
\[
\widehat{K^{\times}}=U_K\cdot\pi^{\hat{\Z}}
\]
为两个闭子群的乘积, 从而 (由无限 Galois 理论) 决定了一个分解
\begin{equation*}
K^{\mathrm{ab}}=K_{\pi}\cdot K^{\mathrm{un}},
\tag{9}
\end{equation*}
其中 $K_{\pi}$ 是 $K^{\mathrm{ab}}$ 中被 $\phi_K(\pi)$ 固定的子域, $K^{\mathrm{un}}$ 是被 $\phi_K(U_K)$ 固定的子域. 显然, $K_{\pi}$ 是全部满足 $\pi\in\Nm(L^{\times})$ 的有限 Abel 扩张 $L/K$ (必全分歧) 的并. 例如, 当在 $K=\Q_p$ 中选取素元 $p$ 时, 分解式成为
\[
\Q_p^{\mathrm{ab}}=\Bigl(\bigcup_n\Q_p[\zeta_{p^n}]\Bigr)\cdot\Bigl(\bigcup_{(m,p)=1}\Q_p[\zeta_m]\Bigr).\erratum{原书右端两处 $\Q$ 均漏了下标 $p$, 依 Milne errata 补正.}
\]
映射 $\phi_K$ 限制为一个同构
\[
U_K\to\Gal(K_{\pi}/K).
\]
这个同构把 \ref{eqn:i.1.9.filt} 式 (本书 p.\pageref{eqn:i.1.9.filt}) 中的滤列映到 $\Gal(K_{\pi}/K)$ 上按分歧群 (上标号; 见 \ref{prop:4.3}) 的滤列.
\end{npara}

\begin{npara}\label{para:1.12}
两个全分歧扩张的复合未必全分歧. 例如考虑 $\Q$ 的二次扩张 $\Q[\sqrt{p}]$ 与 $\Q[\sqrt{pq}]$, 其中 $p$, $q$ 为不同的奇素数. 此时 $p$ 在这两个扩张中都全分歧, 但在它们的复合 $\Q[\sqrt{p},\sqrt{q}]$ 中并非全分歧, 因为它在子域 $\Q[\sqrt{q}]$ 中非分歧. 当 $q$ 取为模 $p$ 的非平方时, 把 $\Q$ 换成 $\Q_p$ 这些陈述仍然成立.

因此, 与 Abel 扩张及非分歧扩张的情形相反, 在 $K^{\mathrm{ab}}$ (或 $K^{\mathrm{al}}$) 中不存在 $K$ 的 “最大” 全分歧扩张: 只有依赖于 $\pi$ 选取的极大全分歧扩张 $K_{\pi}$.
\end{npara}

\subsection*{主定理证明的概述}

我们先说明局部 Artin 映射的唯一性由其余陈述推出.

\begin{theorem}\label{thm:1.13}
假设局部存在定理\ref{thm:1.4}. 至多存在一个同态 $\phi\colon K^{\times}\to\Gal(K^{\mathrm{ab}}/K)$ 满足定理\ref{thm:1.1} 的条件 (a) 与 (b).
\end{theorem}

\begin{proof}
若不存在这样的 $\phi$, 则无可证明; 因此我们假定存在某个 $\phi$, 从而推论\ref{cor:1.2} 以 “有限指数的开子群” 替换 “范数群” 后成立. 设 $\pi$ 是 $\mathcal{O}_K$ 的素元, 令 $K_{\pi,n}$ 是满足 $\Nm(K_{\pi,n})=(1+\mathfrak{m}^n)\cdot\langle\pi\rangle$ 的 $K$ 的扩张, 并令 $K_{\pi}=\bigcup_n K_{\pi,n}$. 由于 $\pi$ 从每个 $K_{\pi,n}$ 都是范数, $\phi(\pi)$ 在 $K_{\pi}$ 上作用为 $1$. 由于它在 $K^{\mathrm{un}}$ 上作用为 Frobenius 自同构, $\phi(\pi)$ 在 $K_{\pi}\cdot K^{\mathrm{un}}$ 上的作用便完全确定. 但 $K_{\pi}\cdot K^{\mathrm{un}}=K^{\mathrm{ab}}$ (参见 \ref{para:1.11}), 因此这表明对任意两个满足条件的同态 $\phi$ 与 $\phi'$ 都有 $\phi(\pi)=\phi'(\pi)$. 由于 $\pi$ 任意, 我们已证明若 $\phi$ 与 $\phi'$ 满足条件, 则对所有 $\pi$ 都有 $\phi(\pi)=\phi'(\pi)$. 群 $K^{\times}$ 由 $\mathcal{O}_K$ 的素元生成, 这就证明了 $\phi=\phi'$.
\end{proof}

剩下要证明存在性, 即局部 Artin 映射与局部存在定理的存在性. 我们将给出三个证明. 前两个的优点在于它们显式地构造了具有范数群 $(1+\mathfrak{m}^n)\cdot\langle\pi\rangle$ 的扩张 $K_{\pi,n}$, 以及局部 Artin 映射 $K^{\times}/\Nm(K_{\pi,n})\to\Gal(K_{\pi,n}/K)$. 后两个的优点在于它们证明了整体类域论中用到的上同调结果. 第三个的优点是简短.

只求最快到达整体类域论主定理的读者, 只需理解第三个证明, 因此可以跳过本章其余部分直接进入第 \ref{ch:II} 章.

\paragraph{存在性的第一个证明 (Lubin-Tate 与 Hasse-Arf; I \textsection\textsection 2--4)}
Lubin 与 Tate 的理论显式地构造了每个素元对应的域 $K_{\pi,n}$, 从而构造了域 $K_{\pi}$, 并且显式地构造了一个同态 $\phi_{\pi}\colon U_K\to\Gal(K_{\pi}/K)$; 此外它还表明 $K_{\pi}\cdot K^{\mathrm{un}}$ 以及 $\phi_{\pi}$ 到同态 $K^{\times}\to\Gal(K_{\pi}\cdot K^{\mathrm{un}}/K)$ (满足 $\phi(\pi)|_{K^{\mathrm{un}}}=\Frob_K$) 的唯一扩张都与 $\pi$ 无关. 由此得到定理\ref{thm:1.1} 与 \ref{cor:1.2}, 但 $K^{\mathrm{ab}}$ 换成了 $K_{\pi}\cdot K^{\mathrm{un}}$.

还需在附注中补入这一点的证明. 令 $K^{\mathrm{LT}}=K_{\pi}\cdot K^{\mathrm{un}}$. 一个要点是, 若 $L$ 是 $K$ 的有限扩张, 则图表
\[
\begin{tikzcd}
L^{\times}\ar[r,"\phi_L"]\ar[d,"\Nm"'] & \Gal(L^{\mathrm{LT}}/L)\ar[d,"\tau\mapsto\tau|_{K^{\mathrm{LT}}}"]\\
K^{\times}\ar[r,"\phi_K"] & \Gal(K^{\mathrm{LT}}/K)
\end{tikzcd}
\]
交换 (\cite{Iwasawa1986}, 定理 6.9, p.\,89). 这表明, 对任何包含在 $K^{\mathrm{LT}}$ 中的 $K$ 的 Abel 扩张 $L$, $\phi_K$ 定义了一个满同态
\[
\phi_{L/K}\colon K^{\times}/\Nm(L^{\times})\to\Gal(L/K).
\]
下一步证明它是单的 (ibid., 推论, p.\,90). 这就证明了局部互反律 (以 $K^{\mathrm{LT}}$ 代 $K^{\mathrm{ab}}$), 而局部存在定理随之易得.

剩下只需证明 $K_{\pi}\cdot K^{\mathrm{un}}=K^{\mathrm{ab}}$. 这一点最直接的证明用到代数数论中的 Hasse-Arf 定理.

\paragraph{主定理的第二个证明 (Lubin-Tate 与上同调; I \textsection\textsection 2--3; III \textsection\textsection 1--3)}
按定理\ref{thm:1.1} 与 \ref{thm:1.4}, 对每个素元 $\pi\in\mathcal{O}_K$ 与整数 $n\ge 1$, 存在 $K$ 的一个 Abel 扩张 $K_{\pi,n}$ 满足
\[
\Nm(K_{\pi,n})=(1+\mathfrak{m}^n)\cdot\pi^{\Z}.
\]
这些扩张具有以下性质:

${}^*\mathrm{a}$ $[K_{\pi,n}:K]=(q-1)q^{n-1}$;

${}^*$b 对所有 $n$, $\pi$ 是 $K_{\pi,n}$ 中的范数.

此外 $K^{\mathrm{ab}}=K_{\pi}\cdot K^{\mathrm{un}}$, 其中 $K_{\pi}=\bigcup_{n\ge 1}K_{\pi,n}$.

在本章 \textsection 3, 我们将对每个素元 $\pi\in\mathcal{O}_K$ 与整数 $n\ge 1$ 显式构造一个满足 ${}^*a$ 与 ${}^*b$ 的全分歧 Abel 扩张 $K_{\pi,n}$. 令 $K_{\pi}=\bigcup_{n\ge 1}K_{\pi,n}$, 并对 $m\ge 1$ 令 $K_m$ 为 $K$ 的 $m$ 次非分歧扩张. 我们将显式构造一个同态
\[
\phi_{\pi}\colon K^{\times}\to\Gal((K_{\pi}\cdot K^{\mathrm{un}})/K)
\]
使得

${}^*\mathrm{c}$ $\phi_{\pi}(\pi)|_{K^{\mathrm{un}}}=\Frob_K$;

${}^*\mathrm{d}$ 对所有 $m$, $n$, 当 $a\in(1+\mathfrak{m}^n)\cdot\langle\pi^m\rangle$ 时 $\phi_{\pi}(a)|_{(K_{\pi,n}\cdot K_m)}=\mathrm{id}$.

此外, 我们将证明 $K_{\pi}\cdot K^{\mathrm{un}}$ 与 $\phi_{\pi}$ 都与 $\pi$ 的选取无关.

在第 \ref{ch:III} 章定理\ref{thm:iii.3.4} 中, 我们将证明存在同态 $\phi\colon K^{\times}\to\Gal(K^{\mathrm{ab}}/K)$ 满足定理\ref{thm:1.1} 的条件 (a) 与 (b).

在本小节的余下部分, 我们说明这两个结果如何推出定理\ref{thm:1.4} (从而也推出定理\ref{thm:1.1}), 并附带更精确的陈述: 对所有 $\pi$ 有 $K^{\mathrm{ab}}=K_{\pi}\cdot K^{\mathrm{un}}$ 且 $\phi=\phi_{\pi}$.

设 $K'=K_{\pi}\cdot K^{\mathrm{un}}$ 是 $K^{\mathrm{ab}}$ 的子域, 并令 $\phi'=\phi_{\pi}$ — 注意 $K'$ 与 $\phi'$ 都与素元 $\pi$ 的选取无关.

\begin{npara}\label{para:1.14}
对所有 $a\in K^{\times}$, $\phi(a)|_{K'}=\phi'(a)$.
\end{npara}

\begin{proof}
对 $K$ 的每个素元 $\pi$, $\phi(\pi)$ 在 $K_{\pi,n}$ 上平凡作用, 因为 $\pi$ 是 $K_{\pi,n}$ 中的范数; 而 $\phi'(\pi)$ 在 $K_{\pi,n}$ 上平凡作用, 是由于条件 ${}^*c$ 中取 $m=1$ (我们可假定定义 $\phi'$ 所用的素元就是 $\pi$). 由于 $\phi(\pi)$ 与 $\phi'(\pi)$ 在 $K^{\mathrm{un}}$ 上都作用为 $\Frob_K$, 它们必须在 $K'=\bigcup_n K_{\pi,n}\cdot K^{\mathrm{un}}$ 上一致. 但 $K$ 的素元作为乘法群生成 $K^{\times}$ (任 $a\in K^{\times}$ 可写成 $a=u\pi^r$, 而 $u=(u\pi)\pi^{-1}$), 故这证明了断言.
\end{proof}

\begin{npara}\label{para:1.15}
\textbf{局部 Kronecker-Weber 定理.} 对任何素元 $\pi$,
\[
K^{\mathrm{ab}}=K_{\pi}\cdot K^{\mathrm{un}}.
\]
\end{npara}

\begin{proof}
令
\[
K_{n,m}=K_{\pi,n}\cdot K_m,\qquad U_{n,m}=(1+\mathfrak{m}^n)\cdot\langle\pi^m\rangle.
\]
已知对所有 $a\in U_{n,m}$ 有 $\phi_{\pi}(a)|_{K_{n,m}}=1$. 因此对所有 $a\in U_{n,m}$ 有 $\phi(a)|_{K_{n,m}}=1$, 从而 $U_{n,m}\subset\Nm(K_{n,m}^{\times})$. 但
\[
\begin{aligned}
(K^{\times}:U_{n,m})&=(U_K:1+\mathfrak{m}^n)(\langle\pi\rangle:\langle\pi^m\rangle)\\
&=(q-1)q^{n-1}\cdot m\\
&=[K_{\pi,n}:K][K_m:K]\\
&=[K_{n,m}:K],
\end{aligned}
\]
且已知 $\phi$ 诱导一个同构
\[
K^{\times}/\Nm(K_{n,m}^{\times})\to\Gal(K_{n,m}/K).
\]
因此
\[
U_{n,m}=\Nm(K_{n,m}^{\times}).
\]
现在设 $L$ 是 $K$ 的有限 Abel 扩张. 已知 $\phi$ 定义同构 $K^{\times}/\Nm(L^{\times})\to\Gal(L/K)$, 故 $\Nm(L^{\times})$ 在 $K^{\times}$ 中指数为有限. 按引理\ref{lem:1.3}, 它也是开的, 从而对某个 $n,m\ge 0$ 包含 $U_{n,m}$. 映射
\[
\phi\colon K^{\times}\to\Gal(L\cdot K_{n,m}/K)
\]
是满的; 且对 $a\in K^{\times}$,
\[
\phi(a)\text{ 固定 }L\text{ 的元素}\iff a\in\Nm(L^{\times}),
\]
\[
\phi(a)\text{ 固定 }K_{n,m}\text{ 的元素}\iff a\in\Nm(K_{n,m}^{\times})=U_{n,m}.
\]
由于 $\Nm(L^{\times})\supset U_{n,m}$, 这推出 $L\subset K_{n,m}$.

由此得到 $K^{\mathrm{ab}}=K_{\pi}\cdot K^{\mathrm{un}}$.

我们现在知道, 对 $K$ 的每个素元 $\pi$, 都有 $K^{\mathrm{ab}}=K_{\pi}\cdot K^{\mathrm{un}}$ 且 $\phi=\phi_{\pi}$. 为完成局部存在定理 (\ref{thm:1.4}) 的证明, 我们必须证明 $K^{\times}$ 的每个有限指数开子群 $H$ 都是范数群; 但如上所观察到的, 每个这样的群都对某个 $n$, $m$ 包含 $U_{n,m}$, 而 $U_{n,m}=\Nm(K_{n,m})$. 于是 (\ref{cor:1.2}e) 表明 $H$ 是范数群.
\end{proof}

\paragraph{主定理的第三个证明 (上同调与 Hilbert 符号; III \textsection\textsection 1--5)}
在第 \ref{ch:III} 章, 我们将用上同调 (具体地, Tate 的一个定理) 来构造局部 Artin 映射, 并用 Hilbert 符号来证明 $K^{\times}$ 的每个有限指数开子群都是范数群.

\begin{exercise}\label{ex:1.16}
只用代数数论的结果 (例如 ANT) 证明: 局部域的有限扩张 $L/K$ 全分歧, 当且仅当 $\Nm(L/K)$ 包含一个素元.
\end{exercise}

\section{Lubin-Tate 形式群律}\label{sec:I.2}

关于 $A$ 上形式群律的一个重要事实是: 它把每个有限扩张 $L/K$ 的整数中的极大理想 $\mathfrak{M}$ 变成一个 Galois 模. 足够小心的话, 你甚至可以让 $L$ 是 $K$ 的一个无限完备扩张. 无论如何, 这个群可能与加法群 $(\mathfrak{M},+)$ 以及乘法群 $(1+\mathfrak{M},\times)$ 都大不相同. 特别地, 它的挠子群可以为 Galois 群提供一个非常有用的表示模. (Lubin, \sx{2061570})

\subsection*{幂级数}

设 $A$ 是一个环 (总假定交换且带 $1$). 一个系数在 $A$ 中的幂级数是一个无穷序列
\[
f=(a_0,a_1,a_2,\ldots),\quad a_i\in A,\ i\in\N.
\]
加法与乘法定义为
\[
(a_0,a_1,\ldots)+(b_0,b_1,\ldots)=(a_0+b_0,a_1+b_1,\ldots),
\]
\[
(a_0,a_1,\ldots)(b_0,b_1,\ldots)=(a_0b_0,\ldots,\textstyle\sum_{i+j=k}a_i b_j,\ldots).
\]
若写成
\[
f=\sum_{i\ge 0}a_i T^i,
\]
这些公式便易于记忆.

系数在 $A$ 中的幂级数构成一个交换环, 记为 $A[[T]]$. 幂级数可以像多项式一样处理, 但有几点需要注意. 例如, 一般我们不能把元素 $c\in A$ 代入幂级数 $f(T)\in A[[T]]$, 因为计算 $f(c)=\sum_{i\ge 0}a_i c^i$ 要求我们求和 $A$ 中无穷多个元素, 而我们不是分析学家, 做不到这一点. 同样的理由, 我们只能在一个幂级数 $g(T)$ 的常数项为零时把它代入另一个幂级数 $f(T)$; 此时 $f(g(T))$ 有定义, 记为 $f\circ g$.

\begin{lemma}\label{lem:2.1}
(a) 对所有幂级数 $f\in A[[T]]$ 与 $g,h\in TA[[T]]$,
\[
f\circ(g\circ h)=(f\circ g)\circ h.
\]
(b) 设 $f=\sum_{i\ge 1}^{\infty}a_i T^i\in TA[[T]]$. 存在 $g\in TA[[T]]$ 使得 $f\circ g=T$, 当且仅当 $a_1$ 是 $A$ 中的单位; 此时 $g$ 唯一, 且满足 $g\circ f=T$.
\end{lemma}

\begin{proof}
(a) 一般有 $(f_1f_2)\circ g=(f_1\circ g)(f_2\circ g)$, 从而 $f^n\circ g=(f\circ g)^n$. 因此当 $f=T^n$ 时, $f\circ(g\circ h)$ 与 $(f\circ g)\circ h$ 都等于 $(g\circ h)^n$; 当 $f=\sum a_i T^i$ 时, 两者都等于 $\sum a_i(g\circ h)^i$.

(b) 我们寻找 $g=\sum_{i\ge 1}b_i T^i$ 使得
\[
\sum_{i\ge 1}a_i g^i=T,
\]
即
\[
a_1b_1=1,
\]
\[
a_1b_2+a_2b_1^2=0,
\]
\[
a_1b_n+\bigl(\text{关于 }a_2,\ldots,a_n,b_1,\ldots,b_{n-1}\text{ 的多项式}\bigr)=0.
\]
第一个方程表明, 为使 $g$ 存在, $a_1$ 必须可逆. 当 $a_1$ 可逆时, 这些方程唯一地确定了 $b_i$. 现在, 由于 $b_1$ 可逆, 同样的论证表明存在 $h\in TA[[T]]$ 使得 $g\circ h=T$. 但
\[
f=f\circ T=f\circ g\circ h=T\circ h=h,
\]
故 $g\circ f=T$.
\end{proof}

注意: 一般有 $f\circ(g+h)\neq f\circ g+f\circ h$.

多个变元的幂级数可类似定义. 此外, 若 $f(X_1,\ldots,X_n)\in A[[X_1,\ldots,X_n]]$ 且 $g_1,g_2,\ldots,g_n\in A[[Y_1,\ldots,Y_m]]$, 则只要所有 $g_i$ 的常数项都为零, $f(g_1,\ldots,g_n)$ 就是 $A[[Y_1,\ldots,Y_m]]$ 中一个良定义的元素.

\begin{remark}\label{rem:2.2}
设 $A$ 是一个完备离散赋值环, $\mathfrak{m}$ 是 $A$ 中的极大理想. 对每个 $f=\sum_{i\ge 0}a_i T^i\in A[[T]]$ 与每个 $c\in\mathfrak{m}$, 当 $i\to\infty$ 时 $a_i c^i\to 0$. 因此级数 $\sum_{i\ge 0}a_i c^i$ 收敛到 $A$ 中的一个元素 $f(c)$, 且当 $a_0\in\mathfrak{m}$ 时它落在 $\mathfrak{m}$ 中.
\end{remark}

\subsection*{形式群律}

一个群是一个带有满足群公理的合成法则的非空集合. 一个\emph{形式群律}\index{xingshiqunlü@形式群律(formal group law)}是一个满足群公理的合成法则 (但没有集合). 更精确地:

\begin{definition}\label{def:2.3}
设 $A$ 是一个交换环. $A$ 上的一个\emph{一元交换形式群律}是一个幂级数 $F\in A[[X,Y]]$ 满足

(a) $F(X,Y)=X+Y+\text{次数}\ge 2\text{ 的项}$;

(b) $F(X,F(Y,Z))=F(F(X,Y),Z)$;

(c) 存在唯一的 $i_F(X)\in XA[[X]]$ 使得 $F(X,i_F(X))=0$;

(d) $F(X,Y)=F(Y,X)$.
\end{definition}

\begin{remark}\label{rem:2.4}
(a) 条件 (a) 保证 $F(X,Y)$ 无常数项, 因此公理 (b) 有意义: 我们在每个次数上比较有限和.

(b) 在公理 (a) 与 (b) 中取 $Y=Z=0$, 我们发现
\[
F(X,0)=X+\text{次数}\ge 2\text{ 的项},\qquad F(F(X,0),0)=F(X,0).
\]
把幂级数 $F(X,0)$ 记为 $f(X)$. 第一个等式表明存在 $g$ 使得 $f\circ g=X$, 第二个等式说 $f\circ f=f$. 在第二个等式上复合 $g$ 得 $f=X$. 于是 $F(X,0)=X$, 类似地 $F(0,Y)=Y$. 因此
\[
F(X,Y)=X+Y+\sum_{\substack{1\le i<\infty\\ 1\le j<\infty}}a_{i,j}X^iY^j.
\]
要得到 $n$ 元群律, 把 $X$ 与 $Y$ 各自换成 $n$ 个变元的序列即可. 公理 (d) 是交换性条件. 由于我们不再考虑其他情形, 以后简称一元交换形式群律为形式群律.
\end{remark}

设 $A=\mathcal{O}_K$ 是非阿基米德局部域 $K$ 中的整数环, 并设 $F=\sum_{i,j}a_{ij}X^iY^j$ 是 $\mathcal{O}_K$ 上的形式群律. 对每个 $x,y\in\mathfrak{m}_K$, 当 $(i,j)\to\infty$ 时 $a_{ij}x^i y^j\to 0$, 因此级数
\[
F(x,y)=\sum a_{ij}x^i y^j
\]
收敛到 $\mathfrak{m}_K$ 中的一个元素 $x+_F y$. 这样 $\mathfrak{m}_K$ 成为一个交换群 $(\mathfrak{m}_K,+_F)$. 类似地, 对每个有限扩张 $L\supset K$, $\mathfrak{m}_L$ 获得一个群结构, 且包含映射 $(\mathfrak{m}_K,+_F)\hookrightarrow(\mathfrak{m}_L,+_F)$ 是同态.

\begin{example}\label{ex:2.5}
(a) 设 $F(X,Y)=X+Y$. 则 $+_F$ 是 $\mathfrak{m}_K$ 上的通常加法.

(b) 设 $F(X,Y)=X+Y+XY$. 映射
\[
a\mapsto 1+a\colon\mathfrak{m}\to 1+\mathfrak{m}
\]
是一个同构 $(\mathfrak{m},+_F)\to(1+\mathfrak{m},\times)$. 验证:
\[
\begin{tikzcd}
(a,b)\ar[r,mapsto]\ar[d,"+_F"'] & (1+a,1+b)\ar[d,"\times"]\\
a+b+ab\ar[r,mapsto] & 1+a+b+ab.
\end{tikzcd}
\]

(c) 设 $E$ 是非阿基米德局部域 $K$ 上的一条椭圆曲线, 并设
\[
Y^2Z+a_1XYZ+a_3YZ^2=X^3+a_2X^2Z+a_4XZ^2+a_6Z^3
\]
是 $E$ 的一个极小 Weierstrass 模型. 令 $T=-X/Y$. 把 $E$ 上的群律展开为 $T_1$, $T_2$ 的幂级数, 我们得到 $\mathcal{O}_K$ 上的一个形式群律 $F_E(T_1,T_2)$. 见 \cite{Milne2020}, 2.6, 2.7.
\end{example}

\begin{definition}\label{def:2.6}
设 $F(X,Y)$ 与 $G(X,Y)$ 是形式群律. 一个\emph{同态}\index{xingshiqunlü@形式群律(formal group law)!tonghui@同态(homomorphism)} $F\to G$ 是一个幂级数 $h\in TA[[T]]$ 使得
\[
h(F(X,Y))=G(h(X),h(Y)).
\]
当 $h$ 有逆, 即存在同态 $h'\colon G\to F$ 使得
\[
h\circ h'=T=h'\circ h
\]
时, $h$ 称为\emph{同构}\index{xingshiqunlü@形式群律(formal group law)!tonggou@同构(isomorphism)}. 一个同态 $h\colon F\to F$ 称为 $F$ 的\emph{自同态}\index{xingshiqunlü@形式群律(formal group law)!zitongtai@自同态(endomorphism)}.
\end{definition}

在 $A=\mathcal{O}_K$ 的情形, 一个同态 $f\colon F\to G$ 对每个 $L\supset K$ 定义一个同态
\[
a\mapsto f(a)\colon(\mathfrak{m}_L,+_F)\to(\mathfrak{m}_L,+_G).
\]

\begin{example}\label{ex:2.7}
设 $F=X+Y+XY=(1+X)(1+Y)-1$. 则 $f(T)=(1+T)^p-1$ 是 $F$ 的一个自同态, 因为
\[
F(f(X),f(Y))=(1+X)^p(1+Y)^p-1=f(F(X,Y)).
\]
注意下列图表交换:
\[
\begin{tikzcd}
\mathfrak{m}\ar[r,"f"]\ar[d,"a\mapsto 1+a"'] & \mathfrak{m}\ar[d,"a\mapsto 1+a"]\\
1+\mathfrak{m}\ar[r,"a\mapsto a^p"] & 1+\mathfrak{m}
\end{tikzcd}
\]
即当把 $(\mathfrak{m},+_F)$ 与 $(1+\mathfrak{m},\times)$ 等同起来时, $f$ 就变成了 $a\mapsto a^p$.
\end{example}

设 $G$ 是一个形式群律. 对任意 $f,g\in TA[[T]]$, 我们定义
\[
f+_G g=G(f(T),g(T)).
\]
由公理 \ref{def:2.3}a,b,c,d, 这个合成法则使 $TA[[T]]$ 成为一个交换群. 特别地,
\[
f+_G(i_G\circ f)=0.
\]

\begin{lemma}\label{lem:2.8}
(a) 对任何形式群律 $F$ 与 $G$, 从 $F$ 到 $G$ 的同态的集合 $\Hom(F,G)$ 关于加法 $f+_G g$ 成为一个 Abel 群.

(b) 对任何形式群律 $F$, $F$ 的自同态的 Abel 群 $\End(F)$ 关于乘法 $f\circ g$ 成为一个环 (未必交换).
\end{lemma}

\begin{proof}
设 $f$ 与 $g$ 是 $F\to G$ 的同态, 并令 $h=f+_G g$. 则
\[
\begin{aligned}
h(F(X,Y))&\stackrel{\mathrm{def}}{=}G(f(F(X,Y)),g(F(X,Y)))\\
&=G(G(f(X),f(Y)),G(g(X),g(Y))).
\end{aligned}
\]
(至少在符号上) 可把这个最后的幂级数写成
\begin{equation*}
(f(X)+_G f(Y))+_G(g(X)+_G g(Y)),\tag{$*$}
\end{equation*}
而由结合律与交换律我们可把它改写为
\begin{equation*}
(f(X)+_G g(X))+_G(f(Y)+_G g(Y)),\tag{$**$}
\end{equation*}
即为 $G(h(X),h(Y))$. 更形式地说, 把 (*) 变成 (**) 的那些运算, 也把 $G(G(f(X),f(Y)),G(g(X),g(Y)))$ 变成 $G(h(X),h(Y))$. 这就证明了 $h\in\Hom(F,G)$. 类似地可证 $i_G\circ f\in\Hom(F,G)$. 由于 $0\in\Hom(F,G)$, 这完成了 $\Hom(F,G)$ 是 $(TA[[T]],+_G)$ 的子群的证明.

我们在引理\ref{lem:2.1} 中已证 $f,g\mapsto f\circ g$ 是结合的. 为证 $\End(F)$ 是环, 还需注意对 $f,g,h\in\End(F)$,
\[
\begin{aligned}
f\circ(g+_F h)&\stackrel{\mathrm{def}}{=}f(F(g(T),h(T)))=F((f\circ g)(T),(f\circ h)(T))=f\circ g+_F f\circ h,\\
(f+_F g)\circ h&=f\circ h+_F g\circ h\quad(\text{同理}),
\end{aligned}
\]
且 $\End(F)$ 有单位元, 即 $T$.
\end{proof}

形式群律类似于代数群. 主要差别在于: 由于它们由幂级数而非多项式定义, 其点的坐标必须 “接近 $1$” 才能使乘积有定义. 关于形式群律有一个非常广的理论 — 例如见 M.\,Hazewinkel, \emph{Formal Groups and Applications}, Academic Press, 1978.

\subsection*{Lubin-Tate 群律}

现在令 $A=\mathcal{O}_K$ 为非阿基米德局部域 $K$ 的整数环, 并选取 $A$ 的一个素元 $\pi$.

\begin{definition}\label{def:2.9}
设 $\mathcal{F}_{\pi}$ 是满足下列条件的 $f(X)\in A[[X]]$ 的集合:

(a) $f(X)=\pi X+\text{次数}\ge 2\text{ 的项}$;

(b) $f(X)\equiv X^q\bmod\pi$.
\end{definition}

\begin{example}\label{ex:2.10}
(a) 多项式 $f(X)=\pi X+X^q$ 属于 $\mathcal{F}_{\pi}$.

(b) 设 $K=\Q_p$, $\pi=p$; 则多项式
\[
f(X)=(1+X)^p-1=pX+\binom{p}{2}X^2+\cdots+pX^{p-1}+X^p
\]
属于 $\mathcal{F}_p$.
\end{example}

\begin{lemma}\label{lem:2.11}
设 $f,g\in\mathcal{F}_{\pi}$, 并设 $\phi_1(X_1,\ldots,X_n)$ 是一个系数在 $A$ 中的线性型. 存在唯一的 $\phi\in A[[X_1,\ldots,X_n]]$ 使得
\[
\begin{cases}
\phi(X_1,\ldots,X_n)=\phi_1+\text{次数}\ge 2\text{ 的项},\\
f(\phi(X_1,\ldots,X_n))=\phi(g(X_1),\ldots,g(X_n)).
\end{cases}
\]
\end{lemma}

\begin{proof}
我们对 $r$ 用归纳法证明: 存在唯一的 $r$ 次多项式 $\phi_r(X_1,\ldots,X_n)$ 使得
\[
\begin{cases}
\phi_r(X_1,\ldots,X_n)=\phi_1+\text{次数}\ge 2\text{ 的项},\\
f(\phi_r(X_1,\ldots,X_n))=\phi_r(g(X_1),\ldots,g(X_n))+\text{次数}\ge r+1\text{ 的项}.
\end{cases}
\]
第一个多项式的唯一候选就是 $\phi_1$ 本身. 它显然满足第一个条件; 若写 $\phi_1=\sum a_i X_i$, 则第二个条件说
\[
\pi\Bigl(\sum a_i X_i\Bigr)=\sum a_i(\pi X_i)+\text{次数}\ge 2\text{ 的项},
\]
这也成立.

设 $r\ge 1$ 且我们已定义 $\phi_r$. 由于 $\phi_r$ 唯一, $\phi_{r+1}$ 必等于 $\phi_r+Q$, 其中 $Q$ 是 $A[X_1,\ldots,X_n]$ 中一个 $r+1$ 次齐次多项式. 我们需要
\[
f(\phi_{r+1}(X_1,\ldots,X_n))\stackrel{?}{=}\phi_{r+1}(g(X_1),\ldots,g(X_n))+\text{次数}\ge r+2\text{ 的项}.
\]
左边为
\[
f(\phi_r(X_1,\ldots,X_n))+\pi Q(X_1,\ldots,X_n)+\text{次数}\ge r+2\text{ 的项},
\]
而右边为
\[
\phi_r(g(X_1),\ldots,g(X_n))+Q(\pi X_1,\ldots,\pi X_n)+\text{次数}\ge r+2\text{ 的项}.
\]
由于 $Q$ 是 $r+1$ 次齐次的, $Q(\pi X_1,\ldots)=\pi^{r+1}Q(X_1,\ldots)$,\erratum{原书右端误作 $\pi^{r+1}Q(\pi X_1,\ldots)$, 多出一个 $\pi$, 依 Milne errata (D. Craven) 删去.} 故我们需要
\[
(\pi^{r+1}-\pi)Q(X_1,\ldots,X_n)\stackrel{?}{=}f(\phi_r(X_1,\ldots,X_n))-\phi_r(g(X_1),\ldots,g(X_n))+\text{次数}\ge r+2\text{ 的项}.
\]
因此 $Q$ 必须是满足下式的唯一多项式:
\[
\frac{f(\phi_r(X_1,\ldots,X_n))-\phi_r(g(X_1),\ldots,g(X_n))}{(\pi^r-1)\pi}=Q+\text{次数}\ge r+2\text{ 的项}.
\]
注意, 由于二项式定理在特征 $p$ 下的简单形式,
\[
f\circ\phi_r-\phi_r\circ g\equiv\phi_r(X_1,\ldots,X_n)^q-\phi_r(X_1^q,\ldots,X_n^q)\equiv 0\bmod\pi.
\]
由于 $\pi$ 整除 $f\circ\phi_r-\phi_r\circ g$, 而 $\pi^r-1$ 在 $A$ 中可逆, $Q$ 的系数确实在 $A$ 中; 又因 $\phi_r$ 满足归纳假设, 它的次数确为 $r+1$.

定义了 $\phi_r$ ($r=1,2,\ldots$) 并注意
\[
\phi_{r+1}=\phi_r+\text{次数}\ge r+1\text{ 的项}
\]
之后, 我们可以定义 $\phi$ 为满足
\[
\phi=\phi_r+\text{次数}\ge r+1\text{ 的项}\quad(\text{对所有 }r)
\]
的唯一幂级数. 显然它具有所需的第一个性质; 且对任意 $r$,
\[
\begin{aligned}
f(\phi(X_1,\ldots,X_n))&=f(\phi_r(X_1,\ldots,X_n))+\text{次数}\ge r+1\text{ 的项}\\
&=\phi_r(g(X_1,\ldots,X_n))+\text{次数}\ge r+1\text{ 的项}\\
&=\phi(g(X_1,\ldots,X_n))+\text{次数}\ge r+1\text{ 的项}.
\end{aligned}
\]
由于这对所有 $r$ 成立, $\phi$ 也具有所需的第二个性质.
\end{proof}

\begin{proposition}\label{prop:2.12}
对每个 $f\in\mathcal{F}_{\pi}$, 存在唯一的系数在 $A$ 中的形式群律 $F_f$ 以 $f$ 为自同态.
\end{proposition}

\begin{proof}
按引理\ref{lem:2.11}, 存在唯一的幂级数 $F_f(X,Y)$ 使得
\[
\begin{cases}
F_f(X,Y)=X+Y+\text{次数}\ge 2\text{ 的项},\\
f(F_f(X,Y))=F_f(f(X),f(Y)).
\end{cases}
\]
剩下只需验证这是一个形式群律.

交换性: 令 $G=F_f(Y,X)$. 则
\[
\begin{cases}
G(X,Y)=X+Y+\text{次数}\ge 2\text{ 的项},\\
f(G(X,Y))=f(F_f(Y,X))=F_f(f(Y),f(X))=G(f(X),f(Y)).
\end{cases}
\]
由于 $F_f(X,Y)$ 是具有这些性质的唯一幂级数, 推出 $G(X,Y)=F_f(X,Y)$.

结合性: 令 $G_1(X,Y,Z)=F_f(X,F_f(Y,Z))$, $G_2(X,Y,Z)=F_f(F_f(X,Y),Z)$. 则对 $i=1,2$,
\[
\begin{cases}
G_i(X,Y,Z)=X+Y+Z+\text{次数}\ge 2\text{ 的项},\\
G_i(f(X),f(Y),f(Z))=f(G_i(X,Y,Z)),
\end{cases}
\]
而引理\ref{lem:2.11} 再次表明满足这些条件的幂级数只有一个.
\end{proof}

\begin{example}\label{ex:2.13}
设 $K=\Q_p$, $\pi=p$. 则 $f(T)=(1+T)^p-1$ 属于 $\mathcal{F}_p$, 且 $F(X,Y)=X+Y+XY$ 以 $f$ 为自同态 (见 \ref{ex:2.7}). 因此 $F=F_f$.
\end{example}

命题所定义的那些形式群律 $F_f$ 就是 \emph{Lubin-Tate 形式群律}\index{lubintatexingshiqunlü@Lubin-Tate 形式群律(Lubin-Tate formal group)}. 它们恰是以如下方式自同态的形式群律:

\begin{itemize}
\item 其在原点的导数等于 $K$ 的一个素元, 且
\item 其模 $\mathfrak{m}$ 约化到 Frobenius 映射 $T\mapsto T^q$.
\end{itemize}

\begin{proposition}\label{prop:2.14}
对 $f,g\in\mathcal{F}_{\pi}$ 与 $a\in A$, 令 $[a]_{g,f}$ 为 $A[[T]]$ 中满足
\[
\begin{cases}
[a]_{g,f}(T)=aT+\text{次数}\ge 2\text{ 的项},\\
g\circ[a]_{g,f}=[a]_{g,f}\circ f
\end{cases}
\]
的唯一元素. 则 $[a]_{g,f}$ 是一个同态 $F_f\to F_g$.
\end{proposition}

\begin{proof}
令 $h=[a]_{g,f}$ — 它的存在性由引理\ref{lem:2.11} 保证. 我们要证明
\[
h(F_f(X,Y))=F_g(h(X),h(Y)).
\]
显然两边都等于 $aX+aY+\text{次数}\ge 2\text{ 的项}$. 此外,
\[
h(F_f(f(X),f(Y)))=(h\circ f)(F_f(X,Y))=g(h(F_f(X,Y))),
\]
\[
F_g(h(f(X)),h(f(Y)))=F_g(g(h(X)),g(h(Y)))=g(F_g(h(X),h(Y))),
\]
我们可以再次应用引理\ref{lem:2.11} 的唯一性.
\end{proof}

\begin{proposition}\label{prop:2.15}
对任意 $a,b\in A$,
\[
[a+b]_{g,f}=[a]_{g,f}+_{F_g}[b]_{g,f}
\]
且
\[
[ab]_{h,f}=[a]_{h,g}\circ[b]_{g,f}.
\]
\end{proposition}

\begin{proof}
在每种情形, 右边的幂级数都满足刻画左边幂级数的条件.
\end{proof}

\begin{corollary}\label{cor:2.16}
对 $f,g\in\mathcal{F}_{\pi}$, $F_f\simeq F_g$.
\end{corollary}

\begin{proof}
对每个 $u\in A^{\times}$, $[u]_{f,g}$ 与 $[u^{-1}]_{g,f}$ 是互逆的同构.

事实上, 存在唯一的同构 $h\colon F_f\to F_g$ 使得 $h(T)=T+\cdots$ 且 $g\circ h=h\circ f$, 即 $[1]_{g,f}$. 因此 $F_f\simeq F_g$.
\end{proof}

\begin{corollary}\label{cor:2.17}
对每个 $a\in A$, 存在唯一的自同态 $[a]_f\colon F_f\to F_f$ 使得 $[a]_f=aT+\text{次数}\ge 2\text{ 的项}$ 且 $[a]_f$ 与 $f$ 交换. 映射
\[
a\mapsto[a]_f\colon A\hookrightarrow\End(F_f)
\]
是一个环同态.
\end{corollary}

\begin{proof}
取 $[a]_f=[a]_{f,f}$ — 它是满足 $aT+\cdots$ 且与 $f$ 交换的唯一幂级数, 并且是 $F_f$ 的自同态. $a\mapsto[a]_f$ 是环同态这一事实由引理\ref{prop:2.15} 以及显然的事实 $[1]_f=T$ 推出.
\end{proof}

由此可知, 对 $K$ 的有限扩张 $L$, Abel 群 $(\mathfrak{m}_L,+_{F_f})$ 有一个自然的 $A$-模结构.

\begin{example}\label{ex:2.18}
设 $K=\Q_p$, $f=(1+T)^p-1\in\mathcal{F}_p$, 于是 $F_f=X+Y+XY$. 对每个 $a\in\Z_p$, 定义
\[
(1+T)^a=\sum_{m\ge 0}\binom{a}{m}T^m,\qquad \binom{a}{m}\stackrel{\mathrm{def}}{=}\frac{a(a-1)\cdots(a-m+1)}{m(m-1)\cdots 1}.
\]
当 $a\in\Z$ 时, 这些定义与通常的相符; 若 $(a_i)_{i\ge 1}$ 是收敛到 $a\in\Z_p$ 的整数序列, 则当 $i\to\infty$ 时 $\binom{a_i}{m}\to\binom{a}{m}$. 因此 $\binom{a}{m}\in\Z_p$. 我断言
\[
[a]_f=(1+T)^a-1.
\]
显然 $(1+T)^a-1=aT+\cdots$, 且
\[
f\circ((1+T)^a-1)=(1+T)^{ap}-1=((1+T)^a-1)\circ f
\]
当 $a$ 为整数时成立, 由连续性它对所有 $a\in\Z_p$ 成立.

在同构 $(\mathfrak{m},+_{F_f})\xrightarrow{t\mapsto 1+t}(1+\mathfrak{m},\times)$ 下, $[a]_f$ 的作用对应于把 $1+\mathfrak{m}$ 中元素映到其 $a$ 次幂的映射.
\end{example}

\begin{remark}\label{rem:2.19}
(a) 注意 $[\pi]_f=f$, 因为 $f$ 满足那两个定义条件.

(b) 同态 $a\mapsto[a]_f\colon A\to\End(F_f)$ 是单的, 因为 $a$ 可作为 $[a]_f$ 的首项系数恢复出来.

(c) 标准同构 $[1]_{g,f}\colon F_f\to F_g$ 与 $A$ 在 $F_f$ 和 $F_g$ 上的作用交换, 因为
\[
[a]_g\circ[1]_{g,f}=[a]_{g,f}=[1]_{g,f}\circ[a]_f.
\]
\end{remark}

\begin{summary}\label{note:2.20}
对每个 $f\in\mathcal{F}_{\pi}$, 存在唯一一个以 $f$ 为自同态的形式群律 $F_f$. 此外, 在 $F_f$ 上存在一个唯一的 $A$-模结构 $a\mapsto[a]_f\colon A\to\End(F_f)$ 使得

(a) 对所有 $a\in A$, $[a]_f=aT+\text{次数}\ge 2\text{ 的项}$;

(b) $[a]_f$ 与 $f$ 交换.

我们有 $[\pi]_f=f$. 若 $g\in\mathcal{F}_{\pi}$, 则 $F_f\simeq F_g$ (通过一个典范的 $A$-同构).
\end{summary}

\begin{exercise}\label{ex:2.21}
设 $F(X,Y)$ 是一个满足 $F(X,0)=X$ 且 $F(0,Y)=Y$ 的幂级数. 证明存在唯一的幂级数 $G(X)=-X+\sum_{i=2}^{\infty}a_i X^i$ 使得 $F(X,G(X))=0$. 由此推出定义\ref{def:2.3} 中的公理 (c) 是多余的.
\end{exercise}

\begin{NOTES}\label{note:2.22}
设 $f\in\mathcal{F}_{\pi}$. 存在唯一的幂级数 $\log(T)=T+\cdots\in K[[T]]$ 使得 $\log(F(x,y))=\log(x)+\log(y)$. 它具有性质
\[
\log([a]_f(T))=a\log(T)
\]
对 $a\in A$ 成立. 见 \sx{3644403}.
\end{NOTES}

\section{扩张 \texorpdfstring{$K_{\pi}$}{K pi} 的构造}\label{sec:I.3}

在本节, $A=\mathcal{O}_K$, 其中 $K$ 是非阿基米德局部域, 其剩余类域 $A/\mathfrak{m}=k$ 有 $q$ (一个 $p$ 的幂) 个元素. 我们固定 $A$ 的一个素元 $\pi$.

按 \textsection 1 中的讨论, 在 $K^{\mathrm{al}}$ 中应存在唯一的 $K$ 的扩张 $K_{\pi}$ 使得 $K^{\mathrm{ab}}=K_{\pi}\cdot K^{\mathrm{un}}$, 且 $\pi$ 是 $K_{\pi}$ 的每个有限子扩张中的范数, 即 $K^{\mathrm{ab}}$ 中被 $\phi(\pi)$ 固定的子域 $K_{\pi}$.

构造 $K^{\mathrm{un}}$ 很容易. 令 $\mu_m$ 为 $K^{\mathrm{al}}$ 中 $m$ 次单位根的集合, 即 $X^m-1$ 的根的集合. 当 $m$ 不被 $p$ 整除时, $X^m-1$ 的判别式是 $\mathcal{O}_K$ 中的单位, 因此由 $\mu_m$ 中元素生成的域 $K[\mu_m]$ 在 $K$ 上非分歧; 此外 $K[\mu_m]$ 的剩余类域是 $X^m-1$ 在 $k$ 上的分裂域, 故有 $q^f$ 个元素, 其中 $f$ 是使得 $m\mid p^f-1$ 的最小正整数. 因此 $\bigcup_{p\nmid m}K[\mu_m]$ 是 $K$ 的一个非分歧扩张, 其剩余类域为 $k$ 的代数闭包, 从而等于 $K^{\mathrm{un}}$. Galois 群 $\Gal(K^{\mathrm{un}}/K)\simeq\hat{\Z}$, 且 $a\in\hat{\Z}$ 按如下方式作用 $K^{\mathrm{un}}$: 对每个 $\zeta\in\mu_m$ 以及每个充分接近 $a$ 的整数 $a_0$ (依赖于 $m$), $a*\zeta=\zeta^{q^{a_0}}$ ($=\Frob_K^{a_0}(\zeta)$).\erratum{原书误作 $\zeta^{a_0}$ ($=\Frob_K^{a}(\zeta)$), 依 Milne errata (R. Sundararajan) 改正.}

在 $K=\Q_p$, $\pi=p$ 的情形, 对 $K_{\pi}$ 有一个类似的构造, 即 $(\Q_p)_p=\bigcup\Q_p[\mu_{p^n}]$ — 我们以后将证明它具有所指出的性质. 作用
\[
([m],\zeta)\mapsto\zeta^m\colon\Z/p^n\Z\times\mu_{p^n}\to\mu_{p^n}
\]
使 $\mu_{p^n}$ 成为一个秩 $1$ 的自由 $\Z/p^n\Z$-模. 由于 $\Z/p^n\Z\simeq\Z_p/p^n\Z_p$, 我们可以把 $\mu_{p^n}$ 看成一个 $\Z_p$-模, 同构于 $\Z_p/(p^n)$. $\Z_p$ 在 $\mu_{p^n}$ 上的作用诱导一个同构 $(\Z_p/p^n\Z_p)^{\times}\to\Gal(\Q_p[\mu_{p^n}]/\Q_p)$, 取所有 $n$ 的极限得到同构
\[
\Z_p^{\times}\to\Gal((\Q_p)_p/\Q_p).
\]
因此, 对扩张 $K^{\mathrm{un}}/K$ 与 $(\Q_p)_p/\Q_p$ 两者, 我们都有扩张的一组显式生成元, Galois 群的一个显式描述, 以及 Galois 群在生成元集合上的一个显式描述. 引人注目的是, Lubin-Tate 群对每个 $K$ 与 $\pi$ 的 $K_{\pi}/K$ 给出了类似的结果.

$K$ 上的绝对值 $|\cdot|$ 唯一地扩张到每个有限次子域 $L\subset K^{\mathrm{al}}$, 从而扩张到 $K^{\mathrm{al}}$; 将 $K^{\mathrm{al}}$ 对 $|\cdot|$ 完备化, 记所得的完备域为 $\widehat{K^{\mathrm{al}}}$. 设 $f\in\mathcal{F}_{\pi}$. 对任何 $\alpha,\beta\in\widehat{K^{\mathrm{al}}}$ 满足 $|\alpha|,|\beta|<1$ 以及 $a\in A$, 级数 $F_f(\alpha,\beta)$ 与 $[a]_f(\alpha)$ 收敛.\erratum{原书直接在 $K^{\mathrm{al}}$ 中断言这些级数收敛; 但 $K^{\mathrm{al}}$ 不完备, 收敛性并无保证. 依 Milne errata (R. Sundararajan), 改在其完备化 $\widehat{K^{\mathrm{al}}}$ 中陈述; 各 $\pi^n$-挠子模 $\Lambda_n$ 作为 $f^{(n)}$ 的根之集仍含于 $K^{\mathrm{al}}$ (见 \ref{rem:3.1}).} 因此我们可以定义 $A$-模 $\Lambda_f$ 为:
\[
\begin{aligned}
\Lambda_f&=\{\alpha\in\widehat{K^{\mathrm{al}}}\mid |\alpha|<1\}\quad(\text{作为集合}),\\
\alpha+_{\Lambda_f}\beta&=\alpha+_F\beta=F_f(\alpha,\beta),\\
a*\alpha&=[a]_f(\alpha).
\end{aligned}
\]
我们定义 $\Lambda_n$ 为 $\Lambda_f$ 中被 $[\pi]_f^n$ 零化的元素的子模. 若 $g$ 是 $\mathcal{F}_{\pi}$ 中另一个元素, 则典范 $A$-同构 $F_f\to F_g$ 诱导一个 $A$-同构 $\Lambda_f\to\Lambda_g$.

\begin{remark}\label{rem:3.1}
回忆 $[\pi]_f(T)=f(T)$, 因此 $\Lambda_n$ 是
\[
f^{(n)}\stackrel{\mathrm{def}}{=}f\circ f\circ\cdots\circ f\qquad(n\text{ 个因子})
\]
在 $K^{\mathrm{al}}$ 中绝对值 $<1$ 的根的集合. 为简单起见, 我们取 $f$ 为多项式 $\pi T+a_2T^2+\cdots+T^q$ 而非更复杂的幂级数 — 按 \ref{cor:2.16} 甚至取 $f=\pi T+T^q$ 就足够.\footnote{或者, $p$-进 Weierstrass 预备定理 (\cite{Washington1997}, 7.3) 表明每个 $f\in F_{\pi}$ 都可分解为 $f_{1}(T)u(T)$, 其中 $f_{1}(T)$ 是多项式 $\pi T+\cdots+aT^{q}$ ($a\equiv1\mod\mathfrak{m}$), 而 $u(T)$ 是 $A[[T]]$ 中的单位. 显然 $f^{(n)}$ 的根就是 $f_{1}^{(n)}$ 的根。} 于是
\[
\begin{aligned}
(f\circ f)(T)&\stackrel{\mathrm{def}}{=}f(f(T))=\pi(\pi T+\cdots+T^q)+\cdots+(\pi T+\cdots+T^q)^q\\
&=\pi^2 T+\cdots+T^{q^2},
\end{aligned}
\]
且
\[
f^{(n)}(T)=\pi^n T+\cdots+T^{q^n}.
\]
由 $f^{(n)}$ 的 Newton 多边形 (参见 ANT, 7.44) 我们看到 $f^{(n)}$ 的根都有正的 $\mathrm{ord}_K$, 从而绝对值 $<1$, 因此 $\Lambda_n$ 是 $f^{(n)}$ 在 $K^{\mathrm{al}}$ 中的全部根, 赋予交换群结构
\[
\alpha+_{F_f}\beta=F_f(\alpha,\beta)=\alpha+\beta+\cdots
\]
以及 $A$-模结构
\[
[a]_f\alpha=a\alpha+\cdots.
\]
\end{remark}

\begin{example}\label{ex:3.2}
取 $K=\Q_p$, $f=(T+1)^p-1\in\mathcal{F}_p$; 则
\[
\Lambda_n=\{\alpha\in\Q_p^{\mathrm{al}}\mid(\alpha+1)^{p^n}=1\}\simeq\{\zeta\mid\zeta^{p^n}=1\}=\mu_{p^n}.\erratum{原书误作 $(\alpha-1)^{p^n}=1$, 依 Milne errata (R. Sundararajan) 改正.}
\]
$\Lambda_n$ 上的 $+_F$ 对应于 $\mu_{p^n}$ 上的乘法, 而 $\Z_p$-模结构如前所定义. 因此 $\Lambda_n\simeq\Z/p^n\Z$ (作为 $\Z_p$-模).
\end{example}

回忆一个循环模是指可由单个元素生成的模. 由于 $A$ 是主理想整环, 且除共轭外只有一个素元, 每个有限生成挠 $A$-模 $M$ 分解为循环模的直和
\[
M\simeq A/(\pi^{n_1})\oplus\cdots\oplus A/(\pi^{n_r}),\quad n_1\le n_2\le\cdots,
\]
且序列 $n_1,\ldots,n_r$ 唯一确定.

\begin{lemma}\label{lem:3.3}
设 $M$ 是一个 $A$-模, 并令 $M_n=\Ker(\pi^n\colon M\to M)$. 假设

(a) $M_1$ 有 $q\stackrel{\mathrm{def}}{=}(A:(\pi))$ 个元素, 且

(b) $\pi\colon M\to M$ 是满的.

则 $M_n\simeq A/(\pi^n)$; 特别地, 它有 $q^n$ 个元素.
\end{lemma}

\begin{proof}
我们对 $n$ 用归纳法. 由于 $A/\pi$ 的阶为 $q$, 条件 (a) 与结构定理推出 $M_1\simeq A/(\pi)$. 考虑正合列
\[
0\to M_1\to M_n\xrightarrow{\pi}M_{n-1}\to 0.
\]
条件 (b) 表明它在 $M_{n-1}$ 处正合, 因此整体正合. 由此推出 $M_n$ 有 $q^n$ 个元素. 若 $M_n$ 不是循环的,
\[
M_n\simeq A/(\pi^{n_1})\oplus A/(\pi^{n_2})\oplus\cdots,\quad n_1,n_2,\ldots\ge 1,
\]
则 $M_1$ 也不是循环的,
\[
M_1\simeq(\pi^{n_1-1})/(\pi^{n_1})\oplus(\pi^{n_2-1})/(\pi^{n_2})\oplus\cdots.
\]
因此 $M_n$ 是一个阶为 $q^n$ 的循环 $A$-模, 而每个这样的模都同构于 $A/(\pi^n)$.
\end{proof}

\begin{proposition}\label{prop:3.4}
$A$-模 $\Lambda_n$ 同构于 $A/(\pi^n)$. 因此 $\End_A(\Lambda_n)\simeq A/(\pi^n)$, 且 $\Aut_A(\Lambda_n)\simeq(A/(\pi^n))^{\times}$.
\end{proposition}

\begin{proof}
形式群律的一个 $A$-同构 $h\colon F_f\to F_g$ 诱导 $A$-模的同构 $\Lambda_f\to\Lambda_g$, 因此我们选哪个 $f\in\mathcal{F}_{\pi}$ 无关紧要. 我们取 $f\in\mathcal{F}_{\pi}$ 为形如 $\pi T+\cdots+T^q$ 的多项式. 这是一个 Eisenstein 多项式 (ANT, p.\,129), 因此有 $q$ 个不同的根, 每个绝对值 $<1$. 设 $\alpha\in K^{\mathrm{al}}$ 绝对值 $<1$. 由
\[
f(T)-\alpha=T^q+\cdots+\pi T-\alpha
\]
的 Newton 多边形我们看到它的根绝对值 $<1$, 从而落在 $\Lambda_f$ 中. 我们已证明 $\Lambda_f$ 满足引理的假设, 故 $\Lambda_n\simeq A/(\pi^n)$. 由此推出 $A$ 在 $\Lambda_n$ 上的作用诱导一个同构 $A/(\pi^n)\to\End_A(\Lambda_n)$.
\end{proof}

\begin{lemma}\label{lem:3.5}
设 $L$ 是局部域 $K$ 的有限 Galois 扩张, Galois 群为 $G$. 对每个 $F\in\mathcal{O}_K[[X_1,\ldots,X_n]]$ 与 $\alpha_1,\ldots,\alpha_n\in\mathfrak{m}_L$,
\[
F(\tau\alpha_1,\ldots,\tau\alpha_n)=\tau F(\alpha_1,\ldots,\alpha_n),\quad\text{所有 }\tau\in G.
\]
\end{lemma}

\begin{proof}
若 $F$ 是多项式, 这由 $\tau$ 是固定 $\mathcal{O}_K$ 中元素的域同构这一事实推出. 由于 $L$ 上的绝对值是 $K$ 上绝对值的唯一扩张 (ANT, 7.38), 它被 $\tau\in G$ 保持, 即对所有 $\alpha\in L$ 有 $|\tau\alpha|=|\alpha|$, 因此 $\tau$ 连续. 于是它保持极限:
\[
\lim_{m\to\infty}\alpha_m=L\implies\lim_{m\to\infty}\tau\alpha_m=\tau L.
\]
写 $F=F_m+\text{次数}\ge m+1\text{ 的项}$, 其中 $F_m$ 是 $m$ 次多项式. 则
\[
\begin{aligned}
\tau(F(\alpha_1,\ldots,\alpha_n))&=\tau\Bigl(\lim_{m\to\infty}F_m(\alpha_1,\ldots,\alpha_n)\Bigr)\\
&=\lim_{m\to\infty}\tau\bigl(F_m(\alpha_1,\ldots,\alpha_n)\bigr)\\
&=\lim_{m\to\infty}F_m(\tau\alpha_1,\ldots,\tau\alpha_n).
\end{aligned}
\]
\end{proof}

\begin{theorem}\label{thm:3.6}
设 $K_{\pi,n}$ 表示由 $\Lambda_n$ 中元素在 $K$ 上生成的子域 $K[\Lambda_n]\subset K^{\mathrm{al}}$.

(a) 对每个 $n$, $K_{\pi,n}/K$ 是全分歧的, 次数为 $(q-1)q^{n-1}$.

(b) $A$ 在 $\Lambda_n$ 上的作用定义一个同构
\[
(A/\mathfrak{m}^n)^{\times}\to\Gal(K_{\pi,n}/K).
\]
特别地, $K_{\pi,n}/K$ 是 Abel 的.

(c) 对每个 $n$, $\pi$ 是 $K_{\pi,n}$ 中的范数.
\end{theorem}

\begin{proof}
我们再次可假定 $f\in\mathcal{F}_{\pi}$ 是形如 $\pi T+\cdots+T^q$ 的多项式. (a), (b). 选取 $f(T)$ 的一个非零根 $\pi_1$, 并 (归纳地) 选取 $f(T)-\pi_{n-1}$ 的一个根 $\pi_n$. 考虑域的序列
\[
K[\Lambda_n]\supset K[\pi_n]\stackrel{q}{\supset}K[\pi_{n-1}]\stackrel{q}{\supset}\cdots\stackrel{q}{\supset}K[\pi_1]\stackrel{q-1}{\supset}K.
\]
每个扩张都是 Eisenstein 的 (ANT, 7.55), 其次数如所标出. 因此 $K[\pi_n]$ 在 $K$ 上全分歧, 次数为 $q^{n-1}(q-1)$.

回忆 $\Lambda_n$ 是 $f^{(n)}$ 在 $K^{\mathrm{al}}$ 中的根的集合, 故 $K[\Lambda_n]$ 是 $f^{(n)}$ 的分裂域. 因此 $\Gal(K[\Lambda_n]/K)$ 可等同于 $\Lambda_n$ 上置换群的一个子群; 但引理\ref{lem:3.5} 表明 $\Gal(K[\Lambda_n]/K)$ 的每个元素在 $\Lambda_n$ 上作为一个 $A$-模同构作用, 因此 $\Gal(K[\Lambda_n]/K)$ 在 $\Lambda_n$ 的自同构群中的像包含于
\[
\End_A(\Lambda_n)=(A/(\pi^n))^{\times}.
\]
于是
\[
(q-1)q^{n-1}\ge|\Gal(K[\Lambda_n]/K)|=[K[\Lambda_n]:K]\ge[K[\pi_n]:K]=(q-1)q^{n-1}.
\]
我们必须处处取等号, 故 $\Gal(K[\Lambda_n]/K)\simeq(A/\mathfrak{m}^n)^{\times}$ 且 $K[\Lambda_n]=K[\pi_n]$.

(c) 令 $f^{[n]}(T)=(f/T)\circ f\circ\cdots\circ f$ ($n$ 项), 于是
\[
f^{[n]}(T)=\pi+\cdots+T^{(q-1)q^{n-1}}.
\]
则 $f^{[n]}(\pi_n)=f^{[n-1]}(\pi_{n-1})=\cdots=f(\pi_1)=0$. 由于 $f^{[n]}$ 是次数为 $(q-1)q^{n-1}=[K[\pi_n]:K]$ 的首一多项式, 它必是 $\pi_n$ 在 $K$ 上的极小多项式. 因此
\[
\Nm_{K[\Lambda_n]/K}\pi_n=(-1)^{(q-1)q^{n-1}}\pi=\pi\quad\text{除非 }q\text{ 为偶数且 }n=1.
\]
在例外情形 ($q$ 为偶数且 $n=1$), $\pi$ 是元素 $-\pi_1$ 的范数.\erratum{原书误作 "除非 $q=2$ 且 $n=1$", 并以 "$K[\Lambda_1]=K$" 处理例外情形; 但只要 $q$ 为偶数 (即 $2$ 的幂) 且 $n=1$, 指数 $(q-1)q^{n-1}$ 就是奇数, 此时 $\Nm_{K[\Lambda_1]/K}(\pi_1)=-\pi$, 应改取 $-\pi_1$ 的范数. 事实上对任何 $q$ 与 $n$ 都有 $\Nm_{K[\Lambda_n]/K}(-\pi_n)=\pi$, 从而无须单列例外情形. 依 Milne errata (R. Sundararajan) 改正.}
\end{proof}

\begin{summary}\label{note:3.7}
设 $f(T)=\pi T+\cdots+T^q$, 并令 $\Lambda_n$ 为 $f^{(n)}$ 在 $K^{\mathrm{al}}$ 中的根的集合. 定义 $K_{\pi,n}=K[\Lambda_n]$. 则
\[
\begin{aligned}
K_{\pi}&=\bigcup K_{\pi,n},\\
K_{\pi,n}&=K[\pi_n]\quad (f(\pi_n)=\pi_{n-1},\ \mathfrak{m}_{K_{\pi,n}}=(\pi_n)),\\
K_{\pi,2}&=K[\pi_2]\quad (f(\pi_2)=\pi_1,\ \mathfrak{m}_{K_{\pi,2}}=(\pi_2)),\\
K_{\pi,1}&=K[\pi_1]\quad (f(\pi_1)=0,\ \mathfrak{m}_{K_{\pi,1}}=(\pi_1),\ \pi_1\neq 0,\ [K_{\pi,1}:K]=q-1).
\end{aligned}
\]
此外, 作用
\[
a*\lambda=[a]_f(\lambda),\quad a\in A,\ \lambda\in\Lambda_n,
\]
诱导一个同构
\[
(A/\mathfrak{m}^n)^{\times}\to\Gal(K_{\pi,n}/K).
\]
取逆向极限, 我们得到一个同构
\[
A^{\times}\to\Gal(K_{\pi}/K).
\]
\end{summary}

\begin{example}\label{ex:3.8}
设 $K=\Q_p$, $f=(T+1)^p-1$. 对每个 $r$, 选取一个 $p^r$ 次单位根 $\zeta_{p^r}$ 使得 $\zeta_p$ 本原且 $\zeta_{p^r}^p=\zeta_{p^{r-1}}$. 则 $\pi_r=\zeta_{p^r}-1$, 且 $(\Q_p)_{p,n}=\Q_p[\pi_r]=\Q_p[\zeta_{p^r}]$. 此外同构
\[
(\Z_p/(p^n))^{\times}\to\Gal(\Q_p[\zeta_{p^r}]/\Q_p)
\]
是标准的那个.
\end{example}

\subsection*{局部 Artin 映射}

我们如下定义一个同态
\[
\phi_{\pi}\colon K^{\times}\to\Gal((K_{\pi}\cdot K^{\mathrm{un}})/K).
\]
设 $a\in K^{\times}$. 由于 $K_{\pi}\cap K^{\mathrm{un}}=K$, 只需分别描述 $\phi_{\pi}(a)$ 在 $K_{\pi}$ 与 $K^{\mathrm{un}}$ 上的作用. 写 $a=u\pi^m$, $u\in U$. 我们规定 $\phi_{\pi}(a)$ 在 $K^{\mathrm{un}}$ 上作用为 $\Frob^m$, 而它在 $K_{\pi}$ 上按如下规则作用:
\[
\phi_{\pi}(a)(\lambda)=[u^{-1}]_f(\lambda),\quad\text{所有 }\lambda\in\bigcup\Lambda_n.
\]
插入 ${}^{-1}$ 是为了使下一个定理成立.

\begin{theorem}\label{thm:3.9}
$K_{\pi}\cdot K^{\mathrm{un}}$ 与 $\phi_{\pi}$ 都与 $\pi$ 的选取无关.
\end{theorem}

本节的余下部分将用于证明此定理.

回忆 $K^{\mathrm{al}}$ 不完备 (见 ANT, 习题 7-7); 事实上即使 $K^{\mathrm{un}}$ 也不完备. 我们以 $\hat{K}^{\mathrm{un}}$ 记其完备化, 以 $B$ 记 $\hat{K}^{\mathrm{un}}$ 的赋值环 (绝对值 $|\cdot|$ 唯一地扩张到 $\hat{K}^{\mathrm{un}}$, 而 $B$ 是绝对值 $\le 1$ 的元素组成的集合). 我们以 $\sigma$ 记 $K^{\mathrm{un}}/K$ 的 Frobenius 自同构 $\Frob_K$, 也记它在 $\hat{K}^{\mathrm{un}}$ 上的扩张. 对幂级数 $\theta(T)=\sum b_i T^i\in B[[T]]$, 我们定义 $(\sigma\theta)(T)$ 为幂级数 $\sum(\sigma b_i)T^i$.

\begin{proposition}\label{prop:3.10}
设 $F_f$ 与 $F_g$ 是由 $f\in\mathcal{F}_{\pi}$ 与 $g\in\mathcal{F}_{\varpi}$ 定义的形式群律, 其中 $\pi$ 与 $\varpi=u\pi$ 是 $K$ 的两个素元. 则 $F_f$ 与 $F_g$ 在 $B$ 上成为 $A$-同构. 更精确地, 存在 $\varepsilon\in B^{\times}$ 使得 $\sigma\varepsilon=\varepsilon u$, 以及幂级数 $\theta(T)\in B[[T]]$ 使得:

(a) $\theta(T)=\varepsilon T+\text{次数}\ge 2\text{ 的项}$;

(b) $\sigma\theta=\theta\circ[u]_f$;

(c) $\theta(F_f(X,Y))=F_g(\theta(X),\theta(Y))$;

(d) $\theta\circ[a]_f=[a]_g\circ\theta$.

最后两个条件说 $\theta$ 是一个与 $A$ 的作用交换的同态 $F_f\to F_g$, 而第一个条件表明 $\theta$ 是同构 (因为 $\varepsilon$ 是单位).
\end{proposition}

\begin{lemma}\label{lem:3.11}
同态
\[
b\mapsto\sigma b-b\colon B\to B,
\]
\[
b\mapsto\sigma b/b\colon B^{\times}\to B^{\times},
\]
都是满的, 核分别为 $A$ 与 $A^{\times}$.
\end{lemma}

\begin{proof}
设 $R$ 是 $K^{\mathrm{un}}$ 中的赋值环, $\mathfrak{n}$ 是其极大理想. 则 $R$ 是离散赋值环, 且 $\varprojlim R/\mathfrak{n}^n=B$ (见下面的 (\ref{ex:A.7})). 我们对 $n$ 用归纳法证明序列
\begin{equation*}
0\longrightarrow A/\mathfrak{m}_K^n\longrightarrow R/\mathfrak{n}^n\xrightarrow{\sigma-1}R/\mathfrak{n}^n\longrightarrow 0
\tag{10}
\end{equation*}
正合. 对 $n=1$, 序列成为
\[
0\xrightarrow{\quad}k\xrightarrow{\quad}\bar{k}\xrightarrow{x\mapsto x^q-x}\bar{k}\xrightarrow{\quad}0.
\]
这里 $\bar{k}$ 是 $k$ 的代数闭包. 这显然正合. 假设序列对 $n-1$ 正合, 并考虑图表
\[
\begin{tikzcd}
0\ar[r] & R/\mathfrak{n}\ar[r]\ar[d,"\sigma-1"] & R/\mathfrak{n}^n\ar[r]\ar[d,"\sigma-1"] & R/\mathfrak{n}^{n-1}\ar[r]\ar[d,"\sigma-1"] & 0\\
0\ar[r] & R/\mathfrak{n}\ar[r] & R/\mathfrak{n}^n\ar[r] & R/\mathfrak{n}^{n-1}\ar[r] & 0.
\end{tikzcd}
\]
由蛇形引理我们发现 $\sigma-1\colon R/\mathfrak{n}^n\to R/\mathfrak{n}^n$ 是满的, 且其核有 $q^n$ 个元素. 由于 $A/\mathfrak{m}_K^n$ 包含在核中且有 $q^n$ 个元素, 它必等于该核. 这证明了 (10) 正合; 当我们对 $n$ 取逆向极限时, 序列成为
\[
0\longrightarrow A\longrightarrow B\xrightarrow{\sigma-1}B\longrightarrow 0,
\]
因而正合 (见下面的命题 \ref{prop:A.8}).

对 $B^{\times}$ 的证明类似.
\end{proof}

一个幂级数 $h$ 关于复合的逆记为 $h^{-1}$. 于是 $h\circ h^{-1}=T=h^{-1}\circ h$.

命题\ref{prop:3.10} 的证明分四步.

第 1 步. 证明存在 $\theta(T)\in B[[T]]$ 满足 (a) 与 (b).

第 2 步. 证明第 1 步中的 $\theta$ 可如此选取, 使得 $g=\sigma\theta\circ f\circ\theta^{-1}$.

第 3 步. 证明幂级数 $\theta(F_f(\theta^{-1}(X),\theta^{-1}(Y)))$ 具有刻画 $F_g(X,Y)$ 的性质, 因此等于它.\erratum{原书误作 "刻画 $F_f(X,Y)$", 依 Milne errata (A. Jangle) 改正.}

第 4 步. 证明幂级数 $\theta\circ[a]_f\circ\theta^{-1}$ 具有刻画 $[a]_g$ 的性质, 因此等于它.

\begin{proof}[第 1 步的证明]
选取 $\varepsilon\in B^{\times}$ 使得 $\sigma\varepsilon=\varepsilon u$ — 它的存在性由引理\ref{lem:3.11} 保证. 从 $\theta_1(T)=\varepsilon T$ 出发, 我们将构造一列多项式 $\theta_r$ 使得
\[
\begin{aligned}
\theta_r(T)&=\theta_{r-1}(T)+bT^r,\ \text{某个 }b\in B,\\
\sigma\theta_r&=\theta_r\circ[u]_f+\text{次数}\ge r+1\text{ 的项}.
\end{aligned}
\]
注意, 对 $\theta_1(T)=\varepsilon T$, 第二个方程成为
\[
\sigma\varepsilon T=\varepsilon(uT+\cdots)+\text{次数}\ge 2\text{ 的项},
\]
这由我们对 $\varepsilon$ 的选取为真. 设已找到 $\theta_r$, 我们要寻找 $\theta_{r+1}(T)=\theta_r(T)+bT^{r+1}$. 写 $b=a\varepsilon^{r+1}$, $a\in B$. 则第二个方程成为
\[
(\sigma\theta_r)(T)+(\sigma a)(\sigma\varepsilon)^{r+1}T^{r+1}\stackrel{?}{=}\theta_r([u]_f(T))+a\varepsilon^{r+1}(uT)^{r+1}+\text{次数}\ge r+1\text{ 的项}.
\]
因此我们需要
\[
(\sigma a-a)(\varepsilon u)^{r+1}\stackrel{?}{=}c,
\]
其中 $c$ 是 $\theta_r\circ[u]_f-\sigma\theta_r$ 中 $T^{r+1}$ 的系数. 我们可以取 $a$ 为 $B$ 中任何满足 $\sigma a-a=c/(\varepsilon u)^{r+1}$ 的元素.
\end{proof}

\begin{proof}[第 2 步的证明]
定义
\[
h=\sigma\theta\circ f\circ\theta^{-1}=\theta\circ[u]_f\circ f\circ\theta^{-1}=\theta\circ f\circ[u]_f\circ\theta^{-1}.
\]
则由于 $f$ 与 $[u]_f$ 的系数在 $A$ 中,
\[
\sigma h=\sigma\theta\circ f\circ[u]_f\circ\sigma\theta^{-1}=\sigma\theta\circ f\circ\theta^{-1}=h.
\]
中间等号用了 $[u]_f\circ\sigma\theta^{-1}=\theta^{-1}$, 这由 $\theta\circ[u]_f\circ\sigma\theta^{-1}=T$ 推出. 由于 $\sigma h=h$, 它落在 $A[[T]]$ 中. 此外,
\[
h(T)=\sigma\varepsilon\cdot\pi\cdot\varepsilon^{-1}T+\cdots=\varpi T+\text{次数}\ge 2\text{ 的项},
\]
且
\[
h(T)\equiv\sigma\theta\circ(\theta^{-1})^q\equiv\sigma\theta(\sigma\theta^{-1}(T^q))\equiv T^q\bmod\mathfrak{m}_K.
\]
因此 $h\in\mathcal{F}_{\varpi}$. 令 $\theta'=[1]_{g,h}\circ\theta$. 则 $\theta'$ 显然仍满足命题的条件 (a), 且仍满足 (b), 因为 $[1]_{g,h}\in A[[T]]$. 此外,
\[
\sigma\theta'\circ f\circ\theta'^{-1}=[1]_{g,h}\circ h\circ[1]_{g,h}^{-1}=g.
\]
\end{proof}

第 3 步与第 4 步的证明是引理\ref{lem:2.11} 的直接应用, 留给读者.

\begin{proof}[$K_{\pi}\cdot K^{\mathrm{un}}$ 与 $\pi$ 无关的证明]
设 $\pi$ 与 $\varpi=\pi u$ 是 $K$ 的两个素元. 由命题\ref{prop:3.10} 我们得到
\[
(\sigma\theta)\circ[\pi]_f=\theta\circ[u]_f\circ[\pi]_f=\theta\circ[\varpi]_f=[\varpi]_g\circ\theta,
\]
即
\[
(\sigma\theta)(f(T))=g(\theta(T)).
\]
因此, 对任何 $\alpha\in K^{\mathrm{al}}$ (回忆这是 $K$ 的可分代数闭包),
\[
f(\alpha)=0\implies g(\theta(\alpha))=0,
\]
且类似地
\[
g(\alpha)=0\implies f(\theta^{-1}(\alpha))=0.
\]
因此 $\theta$ 定义双射 $\Lambda_{f,1}\to\Lambda_{g,1}$, 从而
\[
\hat{K}^{\mathrm{un}}[\Lambda_{g,1}]=\hat{K}^{\mathrm{un}}[\theta(\Lambda_{f,1})]\subset\hat{K}^{\mathrm{un}}[\Lambda_{f,1}]=\hat{K}^{\mathrm{un}}[\theta^{-1}(\Lambda_{g,1})]\subset\hat{K}^{\mathrm{un}}[\Lambda_{g,1}].
\]
因此
\[
\hat{K}^{\mathrm{un}}[\Lambda_{g,1}]=\hat{K}^{\mathrm{un}}[\Lambda_{f,1}].
\]
现在下一个引理表明
\[
\hat{K}^{\mathrm{un}}[\Lambda_{g,1}]\cap K^{\mathrm{al}}=K^{\mathrm{un}}[\Lambda_{g,1}],\qquad \hat{K}^{\mathrm{un}}[\Lambda_{f,1}]\cap K^{\mathrm{al}}=K^{\mathrm{un}}[\Lambda_{f,1}],
\]
于是
\[
K^{\mathrm{un}}[\Lambda_{g,1}]=K^{\mathrm{un}}[\Lambda_{f,1}].
\]
论证可无困难地推广, 表明对所有 $n$ 有
\[
K^{\mathrm{un}}[\Lambda_{g,n}]=K^{\mathrm{un}}[\Lambda_{f,n}],
\]
从而 $K^{\mathrm{un}}\cdot K_{\varpi}=K^{\mathrm{un}}\cdot K_{\pi}$.
\end{proof}

\begin{lemma}\label{lem:3.12}
$K^{\mathrm{al}}$ 中每个含 $K$ 的子域 $E$ 都是闭的 (在拓扑意义下).
\end{lemma}

\begin{proof}
令 $H=\Gal(K^{\mathrm{al}}/E)$. 则 $H$ 固定 $E$ 的每个元素, 因此由连续性它也固定 $E$ 的闭包中的每个元素. 按 Galois 理论, 这表明 $E$ 等于它在 $K^{\mathrm{al}}$ 中的闭包.
\end{proof}

\begin{proof}[$\phi_{\pi}$ 与 $\pi$ 无关的证明]
我们将证明, 对任意两个素元 $\pi$ 与 $\varpi$,
\[
\phi_{\pi}(\varpi)=\phi_{\varpi}(\varpi).
\]
由于 $\pi$ 任意, 这推出对任何其他的素元 $\pi'$,
\[
\phi_{\pi'}(\varpi)=\phi_{\varpi}(\varpi)=\phi_{\pi}(\varpi).
\]
由于 $\varpi$ 也任意, 且素元生成群 $K^{\times}$, 这就推出 $\phi_{\pi}=\phi_{\pi'}$.

在 $K^{\mathrm{un}}$ 上, $\phi_{\pi}(\varpi)$ 与 $\phi_{\varpi}(\varpi)$ 都诱导 Frobenius 自同构. 剩下要证明它们在 $K_{\varpi}$ 上有相同的作用.

设 $\theta$ 是如命题\ref{prop:3.10} 中在 $\hat{K}^{\mathrm{un}}$ 上的同构 $F_f\to F_g$. 它对所有 $n$ 诱导同构 $\Lambda_{f,n}\to\Lambda_{g,n}$. 按定义, $\phi_{\varpi}(\varpi)$ 在 $K_{\varpi}$ 上是恒等; 由于 $K_{\varpi,n}$ 由元素 $\theta(\lambda)$ ($\lambda\in\Lambda_{f,n}$) 在 $K$ 上生成, 剩下只需证明
\[
\phi_{\pi}(\varpi)(\theta(\lambda))=\theta(\lambda),\qquad\text{所有 }\lambda\in\Lambda_{f,n}.
\]
写 $\varpi=u\pi$. 则 $\phi_{\pi}(\varpi)=\phi_{\pi}(u)\cdot\phi_{\pi}(\pi)=\tau\sigma$, 其中
\[
\sigma=\begin{cases}\Frob_K&\text{在 }K^{\mathrm{un}}\text{ 上},\\ \mathrm{id}&\text{在 }\lambda\text{ 上},\end{cases}\qquad
\tau=\begin{cases}\mathrm{id}&\text{在 }K^{\mathrm{un}}\text{ 上},\\ [u^{-1}]_f&\text{在 }\lambda\text{ 上}.\end{cases}
\]
利用级数 $\theta$ 的系数在 $\hat{K}^{\mathrm{un}}$ 中以及 (\ref{prop:3.10}), 我们发现
\[
\phi_{\pi}(\varpi)(\theta(\lambda))=\tau\sigma(\theta(\lambda))=(\sigma\theta)(\tau\lambda)=(\sigma\theta)([u^{-1}]_f(\lambda))=\theta(\lambda).
\]
\end{proof}

\begin{example}\label{ex:3.13}
我们描述局部 Artin 映射 $\phi_p\colon\Q_p^{\times}\to\Gal(L/\Q_p)$ 在 $L=\Q_p[\zeta]$ 的情形, 其中 $\zeta$ 是一个本原 $n$ 次单位根.

(a) 设 $n$ 与 $p$ 互素. 则 $L$ 在 $\Q_p$ 上非分歧, 其次数等于剩余类域扩张的次数. 剩余类域是 $\F_{p^f}$, 其中 $p^f$ 是使得 $n\mid(p^f-1)$ 的 $p$ 的最小幂. 映射 $\phi_p\colon\Q_p^{\times}\to\Gal(L/\Q_p)$ 把 $u\cdot p^t$ 送到 Frobenius 元素的 $t$ 次幂, 其核为 $\Z_p^{\times}\cdot\langle p^f\rangle$.

(b) 设 $n$ 是 $p$ 的幂 $p^r$. 此时 $L$ 在 $K$ 上全分歧, 次数为 $(p-1)p^{r-1}$, 且 $L=(\Q_p)_{p,n}$ (见 \ref{ex:3.8}). 映射 $\phi_p\colon\Q_p^{\times}\to\Gal(L/\Q_p)$ 可描述如下: 令 $a=up^t$, 并令 $u_0\in\Z$ 表示 $u$ 在 $(\Z_p/p^r\Z_p)^{\times}$ 中的类; 则 $\phi_p(a)$ 把 $\zeta$ 送到 $\zeta^{u_0^{-1}}$, 其中 $u_0$ 是一个模 $p^r$ 同余于 $u$ 的整数. 其核为 $\{up^m\mid u\equiv 1\bmod p^r,\ m\in\Z\}$.

(c) 一般情形, 写 $n=m\cdot p^r$, 其中 $m$ 与 $p$ 互素. 则 $L=\Q_p[\zeta_m]\cdot\Q_p[\zeta_{p^r}]$, 且映射
\[
\Q_p^{\times}/\Nm(\Q_p[\zeta_n]^{\times})\to\Gal(\Q_p[\zeta_n]/\Q_p)
\]
可描述如下: 写 $a=up^t$, $u\in\Z_p^{\times}$; 则 $a$ 在 $\Q_p[\zeta_m]$ 上作用为 $\zeta_m\mapsto\zeta_m^{p^t}$,\erratum{原书误作 $\zeta_m\mapsto\zeta_m^{t}$, 依 Milne errata (R. Sundararajan) 改正.} 而在 $\Q_p[\zeta_{p^r}]$ 上作用为 $\zeta_{p^r}\mapsto\zeta_{p^r}^{u_0^{-1}}$, 其中 $u_0$ 是一个模 $p^r$ 同余于 $u$ 的整数.
\end{example}

\section{局部 Kronecker-Weber 定理}\label{sec:I.4}

本节证明的主要结果是 $K^{\mathrm{ab}}=K_{\pi}\cdot K^{\mathrm{un}}$. 由于局部类域论主定理的证明并不需要它, 而且它被这些定理所蕴含, 本节可以跳过.

\subsection*{$K_{\pi,n}/K$ 的分歧群}

设 $L/K$ 是有限 Galois 扩张, Galois 群为 $G$. 回忆 (ANT, 7.57, 7.58, \ldots) 第 $i$ 个分歧群定义为
\[
G_i=\{\tau\in G\mid\mathrm{ord}_L(\tau a-a)\ge i+1,\ \text{所有 }a\in\mathcal{O}_L\};
\]
此外, 对 $i\ge 1$,
\[
G_i=\{\tau\in G_0\mid\mathrm{ord}_L(\tau\varPi-\varPi)\ge i+1\},
\]
其中 $\varPi$ 是 $L$ 的一个素元. 这里 $\mathrm{ord}_L$ 是正规化赋值 $L^{\times}\to\Z$. 于是 $G/G_0=\Gal(l/k)$, 并有包含关系:
\[
\begin{aligned}
(\varPi\mapsto\tau\varPi/\varPi\bmod\varPi)&\colon G_0/G_1\hookrightarrow l^{\times},\\
(\varPi\mapsto(\tau\varPi-\varPi)/\varPi^{i+1}\bmod\varPi)&\colon G_i/G_{i+1}\hookrightarrow l,
\end{aligned}
\]
其中 $l\supset k$ 是 $L$ 与 $K$ 的剩余类域. 因此 $(G_0:G_1)\mid \#l-1$, 且对 $i\ge 1$ 有 $(G_i:G_{i+1})\mid \#l$ (当 $L/K$ 全分歧时 $l=k$, 从而 $\#l=q$).\erratum{原书此处直接写 $q$; 依 Milne errata (Zhang, Zhuo), 应改为 $l$ 的阶, 它只在引理\ref{lem:4.9} 的情形 (全分歧) 才等于 $q=|k|$.} 此外对充分大的 $i$ 有 $G_i=\{1\}$.

令
\[
U^{(0)}=U=A^{\times},\qquad U^{(i)}=1+\mathfrak{m}^i,\ i\ge 1.
\]
则在 $A^{\times}/(1+\mathfrak{m}^n)=U/U^{(n)}$ 上我们有一个滤列
\[
U/U^{(n)}\supset U^{(1)}/U^{(n)}\supset\cdots\supset U^{(n)}/U^{(n)}=0.
\]

\begin{proposition}\label{prop:4.1}
设 $L=K_{\pi,n}$. 在定理\ref{thm:3.6} 的同构 $A^{\times}/U^{(n)}\xrightarrow{\simeq}G$ 下, $U^{(i)}/U^{(n)}$ 映满到 $G_{q^i-1}$.
\end{proposition}

\begin{proof}
我们取 $f=\pi T+T^q$. 显然 $G=G_0$, 且 $U^{(0)}/U^{(n)}$ 映到 $G_0$. 现在取 $i\ge 1$, 并设 $u\in U^{(i)}\setminus U^{(i+1)}$. 则 $u=1+v\pi^i$, 其中 $v\in A^{\times}$, 且
\[
[u]_f(\pi_n)=[1]_f(\pi_n)+[v]_f[\pi^i]_f(\pi_n)=\pi_n+[v]_f(\pi_{n-i})=\pi_n+(\text{单位})\pi_{n-i}.
\]
对每个 $i\ge 1$, $\pi_i=\pi\pi_{i+1}+\pi_i^q=\pi_{i+1}^q(\frac{\pi\pi_{i+1}}{\pi_{i+1}^q}+1)=\pi_{i+1}^q\times\text{单位}$, 因为 $\mathrm{ord}(\frac{\pi}{\pi_{i+1}^{q-1}})>0$. 因此 $\pi_{n-i}=\pi_n^{q^i}\times\text{单位}$, 且
\[
[u]_f(\pi_n)-\pi_n=\pi_n^{q^i}\times\text{单位}.
\]
按定义, 这意味着 $[u]_f\in G_{q^i-1}$, $[u]_f\notin G_{q^i}$. 由于这对所有 $i$ 成立, 它推出 $U^{(i)}$ 映到 $G_{q^i-1}$.

于是
\[
\begin{aligned}
G_0&=G,\\
G_{q-1}&=G_{q-2}=\cdots=G_1,\\
G_{q^2-1}&=G_{q^2-2}=\cdots=G_q,\\
&\ \vdots\\
G_{q^n-1}&=1
\end{aligned}
\]
是 $K_{\pi,n}/K$ 的一组完整的互异分歧群.
\end{proof}

\subsection*{分歧群的上标号}

设 $L$ 是 $K$ 的有限 Galois 扩张, Galois 群为 $G$. 我们把 $G_u$ 的定义扩充到所有实数 $u\ge -1$, 令
\[
G_u=G_i,\quad\text{其中 }i\text{ 是 }\ge u\text{ 的最小整数}.
\]
对 $u>0$,
\[
G_u=\{\tau\in G_0\mid\mathrm{ord}_L(\tau\varPi-\varPi)\ge i+1\}.
\]
定义 $\varphi\colon\R_{\ge 0}\to\R$ 为唯一的连续分段线性函数使得
\[
\begin{cases}
\varphi(0)=0,\\
\varphi'(u)=(G_0:G_u)^{-1}\quad\text{若 }u\text{ 不是整数}.
\end{cases}
\]
定义 $G^v=G_u$, 若 $v=\varphi(u)$, 即 $G^v=G_{\varphi^{-1}(v)}$.

\begin{example}\label{ex:4.2}
设 $L=K_{\pi,n}$. 则
\[
(G_0:G_1)=q-1,\qquad G_1=G_2=\cdots=G_{q-1}.
\]
因此 $\varphi'(u)=\frac{1}{q-1}$ 对 $0<u<q-1$, 且 $\varphi$ 的图像第一段从 $(0,0)$ 跑到 $(q-1,1)$; 故 $G^1=G_{q-1}$. 接着
\[
(G_{q-1}:G_q)=q,\qquad G_q=G_{q+1}=\cdots=G_{q^2-1}.
\]
因此 $\varphi'(u)=\frac{1}{q(q-1)}$ 对 $q-1<u<q^2-1$, 且 $\varphi$ 的图像第二段从 $(q-1,1)$ 跑到 $(q^2-1,2)$. 于是 $G^2=G_{q^2-1}$. 继续这样做, 我们得到如下图像:
\[
\begin{tikzcd}[column sep=small, row sep=small]
G_0 \ar[d, no head, double] & \overset{q-1}{\supset} & G_{q-1} \ar[d, no head, double] & \overset{q}{\supset} & G_{q^2-1} \ar[d, no head, double] & \supset & \cdots & \supset & G_{q^n-1}=1 \ar[d, no head, double] \\
G^0 & & G^1 & & G^2 & & \cdots & & G^n
\end{tikzcd}
\]
\end{example}

\begin{proposition}\label{prop:4.3}
在同构 $A^{\times}/U^{(n)}\to G$ 下,
\[
U^{(i)}/U^{(n)}\xrightarrow{\approx}G^i.
\]
\end{proposition}

\begin{proof}
由 (\ref{prop:4.1}) 与 (\ref{ex:4.2}) 直接推出.
\end{proof}

上标号的定义使得它与过渡到商相容 (而下标号与过渡到子群相容).

\begin{proposition}\label{prop:4.4}
考虑 Galois 扩张 $M\supset L\supset K$, 并令 $G=\Gal(M/K)$, $H=\Gal(M/L)$ (假定为正规的), 于是 $G/H=\Gal(L/K)$. 则
\[
(G/H)^v=\mathrm{Im}(G^v\to G/H),
\]
即 $(G/H)^v=G^vH/H$.
\end{proposition}

\begin{proof}
见 \cite{Serre1962}, IV.3, 命题 14.
\end{proof}

现在考虑一个无限 Galois 扩张 $\Omega/K$. 利用 (\ref{prop:4.4}) 我们可以在 $G=\Gal(\Omega/K)$ 上定义一个滤列:
\[
\tau\in G^v\iff\tau\in\Gal(L/K)^v,\quad\text{所有有限 Galois 子扩张 }L\subset\Omega.
\]

\begin{definition}\label{def:4.5}
对有限 Galois 扩张 $L/K$, 若对所有 $\varepsilon>0$ 都有 $G^v\neq G^{v+\varepsilon}$, 则称 $v$ 是滤列 $\{G^v\}$ 中的一个\emph{跳跃}\index{fenqiqun@分歧群(ramification group)!tiaoyue@跳跃(jump)}.
\end{definition}

\begin{theorem}[Hasse-Arf]\label{thm:4.6}
若 $L/K$ 是 Abel 的, 则跳跃都是整数, 即若 $G_i\neq G_{i+1}$, 则 $\varphi(i)\in\Z$.
\end{theorem}

\begin{proof}
见 \cite{Serre1962}, V.7. (证明相当初等但复杂. 它不要求剩余类域有限, 只要求剩余类域扩张可分.)
\end{proof}

因此, 对有限 Abel 扩张 $L/K$, $G_0=G^0$ 上的滤列形如
\[
G^0\supsetneq G^{i_1}\supsetneq G^{i_2}\ldots\quad i_j\in\N.
\]
此外对充分大的 $n$ 有 $G^n=\{1\}$, 且 $(G^{i_j}:G^{i_j+1})$ 整除 $q-1$ 或 $q$. 对无限 Abel 扩张, 同样的陈述成立, 只是滤列不必终止: 我们只能说
\[
\bigcap G^i=\{1\}.
\]

\begin{example}\label{ex:4.7}
设 $L=K_{\pi,n}$. 若 $G_i\neq G_{i+1}$, 则 $i=0,q-1,\ldots,q^n-1$, 而在那些点 $\varphi$ 取值 $1,2,3,\ldots,n$. 因此我们对所有这些扩张验证了 Hasse-Arf 定理, 并且由 (\ref{prop:4.4}), 对所有子扩张也成立. 特别地, 我们表明了局部 Kronecker-Weber 定理蕴含 Hasse-Arf 定理.
\end{example}

\subsection*{局部 Kronecker-Weber 定理}

如常, $K$ 是局部非阿基米德 域, 且 $K$ 的所有扩张都要求为某个固定的可分代数闭包 $K^{\mathrm{al}}$ 的子域.

\begin{theorem}\label{thm:4.8}
对 $K$ 的每个素元 $\pi$,
\[
K_{\pi}\cdot K^{\mathrm{un}}=K^{\mathrm{ab}}.
\]
\end{theorem}

\begin{lemma}\label{lem:4.9}
设 $L$ 是 $K$ 的 Abel 全分歧扩张. 若 $L\supset K_{\pi}$, 则 $L=K_{\pi}$.
\end{lemma}

\begin{proof}
令 $G=\Gal(L/K)$, $H=\Gal(L/K_{\pi})$, 于是 $G/H=\Gal(K_{\pi}/K)$. 考虑 (Abel 群的) 图表
\[
\begin{tikzcd}[column sep=small]
 & 1\ar[d] & 1\ar[d] & 1\ar[d] & \\
1\ar[r] & G^{n+1}\cap H\ar[r]\ar[d] & G^{n+1}\ar[r]\ar[d] & (G/H)^{n+1}\ar[r]\ar[d] & 1\\
1\ar[r] & G^n\cap H\ar[r]\ar[d] & G^n\ar[r]\ar[d] & (G/H)^n\ar[r]\ar[d] & 1\\
1\ar[r] & \frac{G^n\cap H}{G^{n+1}\cap H}\ar[r]\ar[d] & \frac{G^n}{G^{n+1}}\ar[r]\ar[d] & \frac{(G/H)^n}{(G/H)^{n+1}}\ar[r]\ar[d] & 1\\
 & 1 & 1 & 1 &
\end{tikzcd}
\]
其中各列显然正合, 而命题\ref{prop:4.4} 表明上面两行正合. 因此 (例如由蛇形引理) 第三行也正合, 从而
\[
(G^n:G^{n+1})=((G/H)^n:(G/H)^{n+1})(G^n\cap H:G^{n+1}\cap H).
\]
由于右边 $\le q$ 且 $(G/H)^n:(G/H)^{n+1}$ 等于 $q-1$ 或 $q$, 我们推出对所有 $n$ 有 $G^n\cap H=G^{n+1}\cap H$. 于是
\[
G^{n+1}\cap H=G^n\cap H=\cdots=G^0\cap H=H,
\]
即对所有 $n$ 有 $H\subset G^{n+1}$. 由于 $\bigcap G^n=1$, 这表明 $H=1$.
\end{proof}

\begin{lemma}\label{lem:4.10}
$K_{\pi}$ 的每个有限非分歧扩张都包含在 $K_{\pi}\cdot K^{\mathrm{un}}$ 中.
\end{lemma}

\begin{proof}
设 $L$ 是 $K_{\pi}$ 的一个非分歧扩张. 则对某个 $n$, $L=K_{\pi}\cdot L'$ 对某 $K_{\pi,n}$ 的非分歧扩张 $L'$ 成立.\footnote{$K$ 的有限扩张 $L$ 是可分的, 故由某个元素 $\theta$ 生成, 而 $\theta$ 在 $K$ 上的极小方程只有有限多个系数. 由于 $K_{\pi}$ 是子域链 $K_{\pi,n}$ 的并, 其中必有一个含有该多项式的全部系数。} 现在应用 (ANT, 7.58) 可知 $L'=K_{\pi,n}\cdot L''$ 对某 $K$ 的非分歧扩张 $L''$ 成立.
\end{proof}

\begin{lemma}\label{lem:4.11}
设 $L$ 是 $K_{\pi}$ 的一个指数为 $m$ 的有限 Abel 扩张 (即对所有 $\tau\in\Gal(L/K)$, $\tau^m=1$), 并设 $K_m$ 是 $K_{\pi}$ 的 $m$ 次非分歧扩张. 则存在 $K_{\pi}$ 的一个全分歧 Abel 扩张 $L_t$ 使得
\[
L\subset L_t\cdot K_m=L\cdot K_m.
\]
\end{lemma}

\begin{proof}
对每个 $\tau\in\Gal(LK_m/K_{\pi})$, $\tau^m|_L=1=\tau^m|_{K_m}$, 因此 $\Gal(LK_m/K_{\pi})$ 仍是指数为 $m$ 的 Abel 群. 取 $\tau\in\Gal(LK_m/K_{\pi})$ 使得 $\tau|_{K_m}$ 是 Frobenius 自同构. 则 $\tau$ 的阶为 $m$, 因此
\[
\Gal(LK_m/K)=\langle\tau\rangle\times H\quad(\text{直积})
\]
对某子群 $H$. 令 $L_t=L^{\langle\tau\rangle}$; 则 $L_t$ 在 $K$ 上全分歧, 且 $L\cdot K_m=L_t\cdot K_m$.
\end{proof}

\begin{proof}[定理\ref{thm:4.8} 的证明]
设 $L$ 是 $K$ 的有限 Abel 扩张; 我们要证明 $L\subset K_{\pi}\cdot K^{\mathrm{un}}$. 把引理\ref{lem:4.11} 用于 $L\cdot K_{\pi}$ 表明存在 $K_{\pi}$ 的全分歧扩张 $L_t$ 与 $K_{\pi}$ 的非分歧扩张 $L_u$ 使得
\[
L\cdot K_{\pi}\subset L_t\cdot L_u.
\]
现在 (\ref{lem:4.9}) 推出 $L_t\subset K_{\pi}$, 而 (\ref{lem:4.10}) 推出 $L_u\subset K_{\pi}\cdot K^{\mathrm{un}}$.
\end{proof}

\begin{corollary}\label{cor:4.12}
$\Q_p$ 的每个有限 Abel 扩张都包含在一个分圆扩张中.
\end{corollary}

\begin{example}\label{ex:4.13}
$K$ 的一个有限 Abel 扩张 $L$ 未必能写成 $L=L_t\cdot L_u$, 其中 $L_t$ 在 $K$ 上全分歧, $L_u$ 在 $K$ 上非分歧. 考虑: 域 $\Q_5[\zeta_5]$ 在 $\Q_5$ 上全分歧, 次数为 $4$, Galois 群 $(\Z/5)^{\times}$ 是阶 $4$ 的循环群. 注意 $624=5^4-1$, 因此 $\Q_5[\zeta_{624}]$ 在 $\Q_5$ 上非分歧, 次数为 $4$, 其 Galois 群也是循环的. 显然
\[
\Gal(\Q_5[\zeta_{3120}]/\Q_5)=\langle\sigma\rangle\times\langle\tau\rangle,
\]
其中
\[
\begin{cases}
\sigma(\zeta_{624})=\zeta_{624}^5,\\
\sigma(\zeta_5)=\zeta_5,
\end{cases}\qquad
\begin{cases}
\tau(\zeta_{624})=\zeta_{624},\\
\tau(\zeta_5)=\zeta_5^2.
\end{cases}
\]
令 $L$ 为 $\langle\sigma^2\tau\rangle$ 的固定域. 则 $L$ 是 $\Q_5$ 的循环 Galois 扩张, 次数为 $4$. 它的最大非分歧子域
\[
L_u=L\cap\Q_5[\zeta_{624}]=\Q_5[\zeta_{624}]^{\langle\sigma^2\rangle}=\Q_5[\zeta_{624}^{25}]
\]
在 $\Q_5$ 上次数为 $2$. 若存在域 $L_t$ 使得 $L=L_t\cdot L_u$, 则 $\Gal(L/\Q_5)$ 将是两个阶 $2$ 循环群的乘积, 与它是循环群矛盾.

我们无需命题\ref{prop:3.10} 就重新得到了 $K_{\pi}\cdot K^{\mathrm{un}}$ 与 $\pi$ 无关的结论. 然而, 要证明 $\phi_{\pi}$ 与 $\pi$ 无关仍需要该命题.
\end{example}

\begin{remark}\label{rem:4.14}
原始的 Kronecker-Weber 定理 (由 Hilbert 在 1896 年利用对分歧群的分析证明, 此前 Kronecker 与 Weber 有不完全的证明) 陈述: $\Q$ 的每个有限 Abel 扩张都包含在一个分圆扩张中. 对 $\Q_p$ 同样的陈述称为\emph{局部 Kronecker-Weber 定理}, 而定理\ref{thm:4.8} 通常被称为 $K$ 上的局部 Kronecker-Weber 定理. 事实上可以对 $\Q_p$ 给出局部 Kronecker-Weber 定理的一个初等证明 (见 \cite{Cassels1986}, p.\,151).
\end{remark}

\begin{remark}\label{rem:4.15}
在第 \ref{ch:III} 章, 我们将不借助 Hasse-Arf 定理、直接从定理\ref{thm:1.1} 推出局部 Kronecker-Weber 定理 — 这是 Lubin 与 Tate 的原始思路. 上面的证明遵循 \cite{Gold1981} 与 \cite{Lubin1981} (显然 Leopoldt 更早发现了类似的证明但未发表). 对特征零的局部域, 不用 Hasse-Arf 定理或上同调、但比上述更复杂的局部 Kronecker-Weber 定理证明, 见 \cite{Rosen1981}. 正如 Iwasawa 所指出的 (1986, p.\,115), 一旦命题\ref{prop:4.4} 与 Abel 扩张 $K_{\pi,n}/K$ 的某些性质被当作已知, 那么 $K$ 上的局部 Kronecker-Weber 定理与 Hasse-Arf 定理本质上等价.
\end{remark}

\subsection*{整体 Kronecker-Weber 定理}

既然现在如此容易, 我们不妨也证明原始的 Kronecker-Weber 定理.

\begin{theorem}\label{thm:4.16}
$\Q$ 的每个 Abel 扩张都包含在一个分圆扩张中.
\end{theorem}

\begin{lemma}\label{lem:4.17}
设 $K$ 是 $\Q$ 的有限 Galois 扩张, Galois 群为 $G$. 则 $G$ 由那些在扩张 $K/\Q$ 中分歧的 $K$ 的素理想 $\mathfrak{p}$ 的惯性群生成.
\end{lemma}

\begin{proof}
令 $H$ 为惯性群生成的 $G$ 的子群, 并令 $M$ 为 $H$ 的固定域. 则对所有 $K$ 的素理想 $\mathfrak{p}$ 有 $K^{I(\mathfrak{p})}\supset M$, 因此 $\mathfrak{p}\cap M$ 在扩张 $M/\Q$ 中非分歧. 于是 $M$ 是 $\Q$ 的非分歧扩张, 从而等于 $\Q$ (由 ANT, 4.9).
\end{proof}

\begin{proof}[Kronecker-Weber 定理的证明]\footnote{这里遵循 \cite{Cassels1986}, p.~236。}
设 $K$ 是 $\Q$ 的 Abel 扩张. 设 $p$ 为一个素数, 并设 $\mathfrak{p}$ 是 $\mathcal{O}_K$ 中位于其上的素理想. 由局部 Kronecker-Weber 定理, $K_{\mathfrak{p}}$ 包含在一个 $\Q_p$ 的分圆扩张中, 记为 $K_{\mathfrak{p}}=\Q_p[u_p,v_p]$, 其中 $u_p$ 是一个 $p^{s_p}$ 次单位根, $v_p$ 是一个与 $p$ 互素次数的单位根. 注意 $s_p$ 只依赖于 $p$ (而非 $\mathfrak{p}$), 因为 $\Gal(K_\mathfrak{p}/\Q_p)$ 在位于 $p$ 之上的素理想上传递作用, 从而在域 $K_\mathfrak{p}$ 的集合上传递.

令 $L$ 为 $\Q$ 的由所有在 $K$ 中分歧的素数 $p$ 的 $p^{s_p}$ 次单位根生成的分圆扩张, 并令 $K'=K\cdot L$. 则 $K'$ 在 $\Q$ 上再次 Abel, 且对 $K'$ 的每个素理想 $\mathfrak{p}'$, 我们有
\begin{equation*}
K_{\mathfrak{p}'}'\subset\Q_p[u_p,w_p],
\tag{$*$}
\end{equation*}
其中 $w_p$ 是一个与 $p$ 互素次数的单位根. 显然用 $K'$ 替换 $K$ 来证明定理已足够, 因此我们可以假设 $K\supset L$.

现在我们有
\[
[K:\Q]\ge[L:\Q]=\prod_p\varphi(p^{s_p})\quad(\text{乘积取遍在 }K\text{ 中分歧的素数 }p).
\]
由于 $K/\Q$ 的 Galois 群 $G$ 是交换的, 惯性群 $I(p)$ 只依赖于底层的素数 $p$, 因此记为 $I(p)$. 由 (*) 我们有
\[
(I(p):1)\le\varphi(p^{s_p})
\]
因为 $I(p)$ 是 $\Q_p[u_p,w_p]$ 相应群的商群. 由 (\ref{lem:4.17}), $G$ 由群 $I(p)$ 生成, 因此存在满射 $\prod I(p)\to G$; 于是
\[
(G:1)\le\prod_p(I(p):1)\le\prod\varphi(p^{s_p}).
\]
但 $(G:1)=[K:\Q]$, 因此处处取等号, 且 $K=L$.
\end{proof}

\begin{remark}\label{rem:4.18}
在这一点上, 完成整体类域论主定理在 $K=\Q$ 情形的证明已不太困难 (见引言). 由 $L$ 包含在一个分圆扩张这一事实我们推出 Artin 映射 $\phi_L/\Q$ 有一个模. 现在用 Dirichlet 关于算术级数中素数密度的定理表明 $\phi_L/\Q$ 是满的. 我们知道 $\Nm(I_L^S)$ 包含在核中, 因此此刻所缺的只是
\[
(I_K^S:\Nm I_L^S\cdot i(K_{\mathfrak{m},1}))\le[L:\Q].
\]
这是第一个不等式, 用分析方法不难 (见 \cite{Janusz1996}, IV, 5.6). 一旦有了它, 存在定理就由 $\Q[\zeta_m]$ 有模 $(m)\infty$ 这一事实推出.
\end{remark}

\begin{NOTES}\label{note:4.19}
Lubin-Tate 扩张理论的原始文献是 Lubin 与 Tate \citeyear{LubinTate1965}. 该理论也在 \cite{Serre1967a}, \cite{Iwasawa1986}, 以及 Fesenko 与 Vostokov \citeyear{FesenkoVostokov1993} 中论述.
\end{NOTES}

\subsection*{这一切从何而来?}\label{note:4.20}
我们看到 Lubin-Tate 形式群律对一个表面上非常复杂的问题 — 给出局部域最大 Abel 扩张的显式生成元并描述 Galois 群如何在它们上作用 — 提供了一个异常简单的解. 唯一的神秘之处是有人怎么会想到它们. 下面的揣测或许有帮助.

这种理论可能存在, 是由椭圆曲线的复乘理论暗示的. 这里人们证明, 对二次虚数域 $K$, 存在一条椭圆曲线 $E$, 在同构意义下唯一, 以 $\mathcal{O}_K$ 为其自同态环. 对每个 $n$, $E$ 上阶整除 $n$ 的点构成一个循环 $\mathcal{O}_K$-模, 正是这一事实使得添加它们的坐标得到 Abel 扩张成为可能.

在为非阿基米德局部域寻找类似理论时, 自然要用局部类比 — 即形式群律 — 来替换代数群 $E$. 因此我们寻找一个形式群律, 其自同态环大到足以使它的挠点构成一个循环模. 自同态环显而易见的候选仍然是 $K$ 中的整数环 $\mathcal{O}_K$. 起初自然只要求该形式群承认一个对应于 $\mathcal{O}_K$ 的素元 $\pi$ 的自同态. 对形式群律高度 (height) 的考虑, 加上希望挠点构成循环模, 暗示这个自同态应由一个满足 $f(T)\equiv T^q\bmod\mathfrak{m}$ 的幂级数 $f(T)$ 给出. 此外, 由于我们确实希望该形式群依赖于 $\pi$ 的选取 (因为我们想构造的扩张 $K_{\pi}$ 依赖于它), 自然要求 $f(T)=\pi T+\cdots$. 于是我们被引向集合 $\mathcal{F}_{\pi}$, 而一旦有了它, 理论便自然展开.

实际上, 真实的故事更有趣 — 见下.

\begin{namedblock}{LUBIN MO220796}
由于我曾读过并欣赏 Lazard 关于一维形式群(律)的论文——该文处理的是基域特征为 $p$ 的情形——我决定考察 $p$-进环上的形式群. 出于某种原因, 我想了解这些东西的自同态环, 并逐渐认识到两类情形的相似性: 一方面是椭圆曲线及其模 $p$ 超奇异约化(当这一现象确实发生时), 另一方面是高度大于 $1$ 的 $p$-进整数环上的形式群.

我曾选修或旁听过 Tate 关于椭圆曲线算术的第一门课, 因而对这一切早有准备. 此外, 我知道 Weierstrass 预备定理, 以及它为任何掌握它的人所带来的威力. 在为我的学位论文证明某个结果时, 我曾考虑考察形式群上的挠点, 而且我认为对我而言这很清楚: 它们构成自同态环上的一个模. 请特别注意, 把它们用作 Galois 群的表示模完全不是我的想法.

但 Tate 时时刻刻都在我肩后看着, 无疑看到了许多我未曾考虑的事情. 在我提交学位论文时, 我还没有构造出高度为 $h$、自同态环 $\mathfrak{o}$ 等于 $\Q_p$ 上 $h$ 次局部域 $k$ 的整数环的形式群. 我只对非分歧情形有构造, 而且用的是基于 Lazard 技巧、极其繁琐的逐次次数方法. 在我完成学位论文后不久, 我坐在从 Brunswick 开往 Boston 的巴士上, 发现不仅能在所有情形构造出具有这种极大自同态结构的形式群, 而且其中有一个可以取多项式形式 $\pi x+x^{q}$. Tate 告诉我, 当他看到这一点时, 一切都各就各位了(Everything Fell Into Place). 其结果就是我们论文中那条精彩而优美的引理, 对此我完全不敢邀功. 我的记忆一向不可靠, 但我记得论文其余部分相当迅速地凝聚成形. 别忘了 Tate 当时已是类域论各个方面的大师. 可是, 如果你形式群的自同态环是 $\mathfrak{o}$, 而该形式群的 Tate 模是此自同态环上的秩一模, 那么 $k$ 上 $k(F[p^{\infty}])$ 的 Galois 群与子群 $\mathfrak{o}^{*}\subset k^{*}$ 之间的同构, 难道还不足以让你想到互反映射吗?
\end{namedblock}

\begin{namedblock}{TATE 致 SERRE 的信, 1964 年 1 月 10 日; COLLECTED WORKS, L8}
[由 Lubin 的结果]我们得到一个同态 $G_{K}\to U_{K}=K$ 的单位群. 把它限制到惯性群 $G_{K^{\mathrm{un}}}$ 上, 我们得到一个同态 $G_{K^{\mathrm{un}}}\to U_{K}$; 由上文的唯一性注记, 它很可能是典范的, 而且几乎无法怀疑它实际上就是\emph{互反律同构}(或者其负元!)\ldots 奇迹似乎在于: 一旦放弃代数群而转向形式群, 复乘理论便\emph{普遍地}(局部地)适用.
\end{namedblock}

\phantomsection
\section*{A 附录: 无限 Galois 理论与逆向极限}\label{sec:I.A}%
\addcontentsline{toc}{section}{A 附录: 无限 Galois 理论与逆向极限}
\index{wuxiangaloislilun@无限 Galois 理论(infinite Galois theory)}\index{nixiangjixian@逆向极限(inverse limit)}

我们回顾本章所需的两个专题 (也见 FT 第 7 章).

% 附录中定理按 A.x 编号
\renewcommand{\thetheorem}{A.\arabic{theorem}}
\setcounter{theorem}{0}

\subsection*{Galois 理论用于无限扩张}

固定一个域 $K$.

\begin{definition}\label{def:A.1}
一个域 $\Omega\supset K$ (未必有限次) 称为在 $K$ 上\emph{Galois}\index{galuoiskuozhang@Galois 扩张(Galois extension)!wuxian@无限(infinite)}, 如果

(a) 它在 $K$ 上代数且可分, 即 $\Omega$ 中每个元素都是系数在 $K$ 中的多项式的单根;

(b) 它在 $K$ 上\emph{正规}\index{zhengguikuozhang@正规扩张(normal extension)}, 即每个系数在 $K$ 中、在 $\Omega$ 中有根的不可约多项式在 $\Omega[X]$ 中分裂.
\end{definition}

\begin{proposition}\label{prop:A.2}
一个域 $\Omega$ 在 $K$ 上 Galois, 当且仅当它是 $K$ 的有限 Galois 扩张的并.
\end{proposition}

\begin{proof}
设 $\Omega$ 在 $K$ 上 Galois. 对每个 $\alpha\in\Omega$, $\alpha$ 在 $K$ 上的极小多项式在 $\Omega$ 中的分裂域是 $K$ 的有限 Galois 扩张, 而 $\Omega$ 是这些域的并. 反之, 若 $\Omega$ 是 $K$ 的有限 Galois 扩张的并, 则它在 $K$ 上代数且可分. 此外, 若 $f(X)\in K[X]$ 在 $\Omega$ 中有根, 则它在 $\Omega$ 的某个有限 Galois 子扩张 $E$ 中有根, 因此在 $E[X]$ 中分裂.
\end{proof}

若 $\Omega$ 是 $K$ 的 Galois 扩张, 则 Galois 群 $\Gal(\Omega/K)$ 定义为固定 $K$ 中元素的 $\Omega$ 的自同构群, 赋予如下拓扑: 以集合
\[
\Gal(\Omega/E),\quad \Omega\supset E\supset K,\ [E:K]<\infty
\]
作为 $1$ 处的基本邻域系. 这意味着 $\Gal(\Omega/K)$ 中两个元素接近, 是指它们在某个 “大” 的域 $E$ ($\Omega\supset E\supset K$, $[E:K]<\infty$) 上一致.

\begin{proposition}\label{prop:A.3}
群 $\Gal(\Omega/K)$ 是紧的且 Hausdorff 的.
\end{proposition}

\begin{proof}
考虑映射
\[
\sigma\mapsto\sigma|_E\colon\Gal(\Omega/K)\to\prod\Gal(E/K),
\]
其中乘积取遍所有 $E\subset\Omega$ 且在 $K$ 上有限 Galois 的子域. 命题 \ref{prop:A.2} 表明该映射是单的. 当我们给每个群 $\Gal(E/K)$ 赋予离散拓扑, 给乘积赋予积拓扑时, 在 $\Gal(\Omega/K)$ 上诱导的拓扑就是上面的拓扑. 我断言像是闭的. 像等于满足如下条件的族 $(\sigma_E)_E$ 的集合: 当 $E'\supset E$ 时 $\sigma_{E'}|_E=\sigma_E$. 设 $(\sigma_E)$ 不在像中. 则存在一对域 $E_2\supset E_1$ 使得 $\sigma_{E_2}|_{E_1}\neq\sigma_{E_1}$. 令
\[
U=\prod_{E\neq E_1,E_2}\Gal(E/K)\times\{\sigma_{E_2}\}\times\{\sigma_{E_1}\}.
\]
这是 $(\sigma_E)_E$ 的一个开邻域, 且 $U\cap\mathrm{Im}(\Gal(\Omega/K))=\varnothing$.
\end{proof}

$\Gal(\Omega/K)$ 上的拓扑离散, 当且仅当 $\Omega$ 是 $K$ 的有限扩张.

\begin{theorem}\label{thm:A.4}
设 $\Omega$ 是 $K$ 的 (可能无限的) Galois 扩张, Galois 群为 $G$. 则在 $\Omega$ 的子域与 $G$ 的闭子群之间存在一一对应. 更精确地:

(a) 对 $\Omega$ 的子域 $E$, $H\stackrel{\mathrm{def}}{=}\Gal(\Omega/E)$ 是 $G$ 的闭子群, 且 $E=\Omega^H$.

(b) 若 $H$ 是 $G$ 的子群, 则 $\Gal(\Omega/\Omega^H)$ 是 $H$ 在 $\Gal(\Omega/K)$ 中的闭包.

此外, $G$ 的正规闭子群对应于 $K$ 的 Galois 扩张, 而 $G$ 的开子群对应于 $K$ 的有限扩张.
\end{theorem}

\begin{proof}
设 $E\subset\Omega$ 是 $K$ 的有限扩张. 由于每个 $K$-同态 $E\to\Omega$ 都可扩张到 $\Omega$, 映射 $\sigma\mapsto\sigma|_{\Omega}\colon\Gal(\Omega/K)\to\Hom_K(E,\Omega)$ 是满的, 因此诱导一个双射
\[
\Gal(\Omega/K)/\Gal(\Omega/E)\to\Hom_K(E,\Omega).
\]
这表明 $\Gal(\Omega/E)$ 在 $\Gal(\Omega/E)$ 中的指数为 $[E:K]$. 由于它开 (按定义), 它也闭 (其补集是有限个开陪集的并).

设 $E\subset\Omega$ 是 $K$ 的任意扩张. 则 $E=\bigcup E_i$, 其中 $E_i$ 跑遍 $E$ 中包含的 $K$ 的有限扩张. 相应地, $\Gal(\Omega/E)=\bigcap\Gal(\Omega/E_i)$, 因此闭. 此外, 若 $\alpha\in\Omega$ 不被 $\Gal(\Omega/E)$ 固定, 则它不被任何 $\Gal(\Omega/E_i)$ 固定, 从而不在 $E$ 中. 于是 $E=\Omega^{\Gal(\Omega/E)}$.

设 $H$ 是 $\Gal(\Omega/K)$ 的子群, 并令 $H'=\Gal(\Omega/\Omega^H)$. 由有限扩张的 Galois 理论, 对每个 $\Gal(\Omega/K)$ 的开子群 $U$ 有 $UH=UH'$; 而由拓扑群理论, $H$ 的闭包是 $\bigcap UH$ (交取遍 $\Gal(\Omega/K)$ 的开子群 $U$), 因此 $\bar{H}=\bigcap UH=\bigcap UH'=H'$.
\end{proof}

\begin{example}\label{ex:A.5}
(a) 给 $\Z$ 赋予以有限指数子群为基本邻域系的拓扑, 并令 $\hat{\Z}$ 为完备化. 则 $\hat{\Z}=\prod_{\ell\text{ 素数}}\Z_{\ell}$, 且对每个 $m$ 有 $\hat{\Z}/m\hat{\Z}=\Z/m\Z$. 令 $\F$ 为 $\F_q$ 的代数闭包. 存在典范同构
\[
\hat{\Z}\to\Gal(\F/\F_q)
\]
把 $1\in\hat{\Z}$ 送到自同构 $x\mapsto x^q$. $\F_q$ 的 $m$ 次扩张对应于 $\hat{\Z}$ 的子群 $m\hat{\Z}$. 设 $\sigma$ 是 $\F/\F_q$ 的满足 $\sigma(x)=x^q$ 的自同构. 对 $\alpha\in\hat{\Z}$, 定义 $\sigma^{\alpha}$ 为 $\Gal(\F/\F_q)$ 中如下元素: 对每个 $m$, 当 $a\in\Z$ 充分接近 $\alpha$ (依赖于 $m$) 时 $\sigma^{\alpha}|_{\F_{q^m}}=\sigma^a|_{\F_{q^m}}$. 上述同构把 $\alpha$ 送到 $\sigma^{\alpha}$. (对 $\hat{\Z}$ 与同构 $\hat{\Z}\to\Gal(\F/\F_q)$ 的详细描述, 见 \cite{Artin1951}, 9.2.)

(b) 令 $\Omega_p=\bigcup_r\Q[\zeta_{p^r}]$, 其中 $\zeta_{p^r}$ 是一个本原 $p^r$ 次单位根. 对 $u\in\Z_p^{\times}$, 写 $u=a_0+a_1p+a_2p^2+\cdots$, $0\le a_i\le p-1$, 并定义
\[
\zeta_{p^r}^u=\zeta_{p^r}^{a_0+a_1p+\cdots+a_s p^s},\quad\text{每个 }s\ge r-1.
\]
这定义了 $\Z_p^{\times}$ 在 $\Omega_p$ 上的一个作用, 实际上是一个拓扑群同构
\[
\Z_p^{\times}\to\Gal(\Omega_p/\Q).
\]

(c) 令 $\Omega=\bigcup\Q[\zeta_n]$, 其中 $\zeta_n$ 是一个本原 $n$ 次单位根. 则
\[
\Omega=\Omega_2\cdot\Omega_3\cdot\Omega_5\cdots,\quad \Omega_p\cap\Omega_{p'}=\Q,\ p\neq p'.
\]
正如有限个有限 Galois 扩张的情形, 这推出
\[
\Gal(\Omega/\Q)=\prod\Gal(\Omega_p/\Q)
\]
(闭子群拓扑积). 于是存在一个同构
\[
\hat{\Z}^{\times}\to\Gal(\Omega/\Q).
\]
它可以描述如下: 若 $\zeta$ 是一个 $n$ 次单位根且 $u\in\Z^{\times}$, 则
\[
\zeta^u=\zeta^m\quad\text{对任何满足 }m\equiv u\bmod n\text{ 的 }m\in\Z.
\]

(d) 令 $\Omega_p=\bigcup\Q[\zeta_{p^r}]$, 如 (b). 则
\[
\Z_p^{\times}\simeq\Delta_p\times C_p,
\]
其中对 $p\neq 2$ 有 $\Delta_p\simeq(\Z/(p))^{\times}\approx\Z/(p-1)$, 而 $\Delta_2=\Z/2\Z$, 且 $C_p\approx\Z_p$. 令 $\Omega_p'=\Omega_p^{\Delta_p}$. 则 $\Gal(\Omega_p'/\Q_p)\approx\Z_p$. 令
\[
\Omega'=\Omega_2'\cdot\Omega_3'\cdots
\]
(在 $\Q^{\mathrm{al}}$ 内的复合). 则 $\Gal(\Omega'/\Q)\approx\prod\Z_p\simeq\hat{\Z}$.

(e) 对 $\Q$ 的每个有限扩张 $K$, $K\cdot\Omega'$ 是 $K$ 的 Galois 扩张, 其 Galois 群同构于 $\Z$ 中一个有限指数的子群. 但每个这样的子群又同构于 $\Z$ (因为它又是一个 $\Z$ 的有限指数子群的完备化). 因此 $\Gal(K\cdot\Omega'/K)\approx\Z$. 若我们固定一个同构, 并令 $K_m$ 为对应于 $m\Z$ 的域, 则我们看到:

(a) $K_m$ 是次数为 $m$ 的循环扩张;

(b) $K_m$ 是分圆的, 即包含在 $K$ 添加单位根所得的扩张中.
\end{example}

\begin{remark}\label{rem:A.6}
一般, 即使 $K=\Q$, 也可能存在 $\Gal(\Omega/K)$ 中有限指数但不开的子群 (见 FT, 7.26).
\end{remark}

\subsection*{逆向极限}

一个偏序集 $(I,\preceq)$ 称为\emph{有向的}\index{youxiangji@有向集(directed set)}, 如果对任意 $\alpha,\beta\in I$, 存在 $\gamma\in I$ 使得 $\alpha,\beta\preceq\gamma$. 例如, 以整除关系排序的正整数集
\[
m\preceq n\iff m\mid n
\]
是有向的.

一个\emph{逆向系}\index{nixiangxi@逆向系(inverse system)} (或投射系) 的集合族是指一族由有向集 $I$ 索引的集合 $(S_{\alpha})_{\alpha\in I}$, 连同对每个对 $\alpha\preceq\beta$ 的一个映射 $\varphi_{\alpha,\beta}\colon S_{\beta}\to S_{\alpha}$, 使得

(a) 对所有 $\alpha\in I$, $\varphi_{\alpha,\alpha}$ 是恒等映射;

(b) 对所有 $I$ 中的 $\alpha\preceq\beta\preceq\gamma$, $\varphi_{\alpha,\beta}\circ\varphi_{\beta,\gamma}=\varphi_{\alpha,\gamma}$.

一个集合 $S$ 连同对每个 $\alpha\in I$ 的一个映射 $\varphi_{\alpha}\colon S\to S_{\alpha}$ 使得 $\varphi_{\alpha}=\varphi_{\alpha,\beta}\circ\varphi_{\beta}$, 若满足明显的万有性质, 则称为逆向系 $(S_{\alpha})$, $(\varphi_{\alpha,\beta})$ 的\emph{逆向极限}\index{nixiangjixian@逆向极限(inverse limit)} (或投射极限).

每个集合、群或环的逆向系都有逆向极限. 例如 $\hat{\Z}=\varprojlim\Z/m\Z$.

\begin{example}\label{ex:A.7}
离散赋值环 $R$ 的完备化 $\hat{R}$ 典范同构于 $\varprojlim R/\mathfrak{m}^n$, 其中 $\mathfrak{m}$ 是 $R$ 中的极大理想.
\end{example}

一个群 $G$ 的\emph{profinite 完备化}\index{profinitewanchenghua@profinite 完备化(profinite completion)}定义为 $\varprojlim G/H$, 其中 $H$ 跑遍 $G$ 中有限指数的正规子群.

\begin{proposition}\label{prop:A.8}
有限 Abel 群的正合列组成的逆向系的逆向极限仍正合.
\end{proposition}

\begin{remark}\label{rem:A.9}
Galois 群 $\Gal(\Omega/K)$ 是群 $\Gal(E/K)$ 的投射极限, 其中 $E$ 跑遍 $\Omega$ 中在 $K$ 上有限且 Galois 的子域. 一个拓扑群若是有限群的一个投射极限, 则称为\emph{profinite 群}\index{profinitequn@profinite 群(profinite group)}. 它们恰为紧的全不连通拓扑群. (一个拓扑空间全不连通, 是指它的连通分支都是单点集.) 见 \cite{Serre1964} 或 \cite{Shatz1972}.
\end{remark}

% 恢复定理计数器为 节.序号 格式, 供后续章节使用
\renewcommand{\thetheorem}{\thesection.\arabic{theorem}}
    % I   局部类域论: Lubin-Tate 理论
% ============ 第二章 II: 群的上同调 ============
% 源: J.S. Milne, Class Field Theory (v4.03), Chapter II, pp. 51--71
% 翻译自 OCR 草稿 J.S. Milne - Class Field Theory.pdf_by_PaddleOCR-VL-1.6.md (行 2093--3783)
% 编号约定见 main.tex: 定理类环境按"节.序号"编号, 全部共享 theorem 计数器.
%   §1 中 1.6/1.7/1.8/1.9 为叙述性条目(非定理环境), 故在 rem:1.10 前 \setcounter{theorem}{9}.
%   §2 中 2.1/2.2/2.3/2.4 为叙述性条目, 故在 rem:2.5 前 \setcounter{theorem}{4}.
%   附录 A 改用 \renewcommand{\thetheorem}{A.\arabic{theorem}}.

% 本章局部使用的算子宏 (若已在别处定义则跳过, 不冲突)
\providecommand{\Res}{\operatorname{Res}}
\providecommand{\Cor}{\operatorname{Cor}}
\providecommand{\Inf}{\operatorname{Inf}}
\providecommand{\Ver}{\operatorname{Ver}}
\providecommand{\Coker}{\operatorname{Coker}}
\providecommand{\Ab}{\mathsf{Ab}}      % 交换群范畴
\providecommand{\Mod}{\mathsf{Mod}}    % 模范畴

\chapter{群的上同调}\label{ch:II}\index{qundeshangtongdiao@群的上同调(cohomology of groups)}

我们暂离数论, 来做一点同调代数. 在本章末尾的一个附录中, 我们回顾导出函子的一般理论.

\section{上同调}\label{sec:II.1}\index{shangtongdiao@上同调(cohomology)}

\subsection{\texorpdfstring{$G$}{G}-模的范畴}

设 $G$ 为一个群. 一个 \emph{$G$-模}是一个交换群 $M$ 连同映射
\[
 (g,m)\mapsto g m\colon G\times M\to M
\]
使得对所有 $g,g'\in G$, $m,m'\in M$ 有
\begin{enumerate}[(a)]
\item $g(m+m')=g m+g m'$;
\item $(gg')(m)=g(g'm)$, $1m=m$.
\end{enumerate}
等价地, 一个 $G$-模是一个交换群 $M$ 连同群同态 $G\to\operatorname{Aut}(M)$. 我们也说 $G$ \emph{作用在} $M$ 上, 并说作用是\emph{平凡的}如果对所有 $g\in G$, $m\in M$ 有 $gm=m$.

一个\emph{$G$-模同态}(或 $G$-同态)是一个映射 $\alpha\colon M\to N$ 使得
\begin{enumerate}[(a)]
\item $\alpha(m+m')=\alpha(m)+\alpha(m')$ (即 $\alpha$ 是交换群的同态);
\item $\alpha(gm)=g(\alpha(m))$ 对所有 $g\in G$, $m\in M$.
\end{enumerate}
我们记 $\operatorname{Hom}_G(M,N)$ 为 $G$-同态 $M\to N$ 的集合.

\begin{remark}\label{rem:ii.1.1}
$G$ 的\emph{群代数} $\Z[G]$ 是以 $G$ 的元素为基的自由交换群, 其乘法由 $G$ 上的群律给出. 于是 $\Z[G]$ 的元素是有限和
\[
 \sum_i n_i g_i,\qquad n_i\in\Z,\quad g_i\in G,
\]
并且
\[
 \Bigl(\sum_i n_i g_i\Bigr)\Bigl(\sum_j n_j' g_j'\Bigr)=\sum_{i,j} n_i n_j'(g_i g_j').
\]
一个交换群 $M$ 上的 $G$-模结构唯一地扩张为一个 $\Z[G]$-模结构, 而一个交换群同态是 $G$-模同态当且仅当它是 $\Z[G]$-模同态. 因此 $G$-模范畴 $\Mod_G$ 可以等同于环 $\Z[G]$ 上的模范畴. 特别地, $\Mod_G$ 是一个 Abel 范畴.
\end{remark}

如果 $M$ 和 $N$ 是 $G$-模, 那么同态集合 $\operatorname{Hom}(M,N)$ (把 $M$ 和 $N$ 只看作交换群) 带如下结构成为一个 $G$-模:
\begin{equation*}
 (\varphi+\varphi')(m)=\varphi(m)+\varphi'(m),\qquad (g\varphi)(m)=g(\varphi(g^{-1}m)).
\tag{11}
\end{equation*}

\subsubsection*{诱导模}

对 $H$-模 $M$, 我们定义 $\Ind_H^G(M)$ 为满足 $\varphi(hg)=h\varphi(g)$ 对所有 $h\in H$ 的映射 (不必是同态) $\varphi\colon G\to M$ 的集合. 于是 $\Ind_H^G(M)$ 成为一个 $G$-模, 其运算为
\begin{equation*}
 (\varphi+\varphi')(x)=\varphi(x)+\varphi'(x),\qquad (g\varphi)(x)=\varphi(xg).
\tag{12}
\end{equation*}
一个 $H$-模同态 $\alpha\colon M\to M'$ 定义 $G$-模同态
\[
 \varphi\mapsto \alpha\circ\varphi\colon \Ind_H^G(M)\to \Ind_H^G(M').
\]

\begin{lemma}\label{lem:ii.1.2}
\begin{enumerate}[(a)]
\item 对每个 $G$-模 $M$ 与 $H$-模 $N$,
\[
 \operatorname{Hom}_G(M,\Ind_H^G(N))\simeq \operatorname{Hom}_H(M,N).
\]
\item 函子
\[
 \Ind_H^G\colon \Mod_H\to \Mod_G
\]
是正合的.
\end{enumerate}
\end{lemma}

\begin{proof}
(a) 给定 $G$-同态 $\alpha\colon M\to \Ind_H^G(N)$, 我们定义映射 $\beta\colon M\to N$ 为
\[
 \beta(m)=\alpha(m)(1_G),
\]
其中 $1_G$ 是 $G$ 的单位元. 对每个 $g\in G$,
\[
 \beta(gm)\stackrel{\mathrm{def}}{=}(\alpha(gm))(1_G)=(g(\alpha(m)))(1_G)\stackrel{(12)}{=}\alpha(m)(g).
\]
因为 $\alpha(m)\in \Ind_H^G(N)$, 当 $g\in H$ 时 $\alpha(m)(g)=g(\alpha(m)(1_G))=g(\beta(m))$. 因此 $\beta$ 是一个 $H$-同态 $M\to N$.

反之, 给定 $H$-同态 $\beta\colon M\to N$, 我们定义 $\alpha$ 为映射 $M\to \Ind_H^G(N)$, 满足 $\alpha(m)(g)=\beta(gm)$. 则 $\alpha$ 是一个 $G$-同态.

由于映射 $\alpha\mapsto\beta$ 与 $\beta\mapsto\alpha$ 互逆, 二者都是同构.

(b) 给定正合列
\[
 0\to M\to N\to P\to 0,
\]
我们必须证明
\[
 0\to \Ind_H^G M\to \Ind_H^G N\to \Ind_H^G P\to 0
\]
正合. 除最后一项外这是显然的. 取 $H$ 在 $G$ 中的右陪集代表系 $S$, 即 $G=\coprod_{s\in S}Hs$, 并设 $\varphi\in \Ind_H^G(P)$. 对每个 $s\in S$, 取一个映到 $P$ 中 $\varphi(s)$ 的 $n(s)\in N$, 并定义 $\tilde\varphi(hs)=h\cdot n(s)$. 则 $\tilde\varphi\in \Ind_H^G(N)$ 且映到 $\varphi$.
\end{proof}

设 $\phi$ 为映射 $\Ind_H^G(N)\to N$, $\varphi\mapsto \varphi(1_G)$; 这是一个 $H$-同态, 引理表明 $(\Ind_H^G(N),\phi)$ 具有如下泛性质: 对任何从 $G$-模 $M$ 到 $N$ 的 $H$-同态 $\beta$, 存在唯一的 $G$-同态 $\alpha\colon M\to \Ind_H^G(N)$ 使得 $\phi\circ\alpha=\beta$:
\[
\begin{tikzcd}
	M && {\Ind_H^G(N)} \\
	&& N
	\arrow["\alpha"{description}, from=1-1, to=1-3]
	\arrow["\beta"'{pos=0.4}, curve={height=18pt}, from=1-1, to=2-3]
	\arrow["\phi", from=1-3, to=2-3]
\end{tikzcd}
\]
(其中 $\alpha$ 为 $G$-映射, 而 $\phi$ 与 $\beta$ 为 $H$-映射).
当 $H=\{1\}$ 时, 一个 $H$-模不过是一个交换群. 此时我们从记号 $\Ind_H^G$ 中略去 $H$. 于是
\[
 \Ind^G(M_0)=\{\varphi\colon G\to M_0\}\quad(\text{映射, 不必是同态})=\operatorname{Hom}(\Z[G],M_0)\quad(\text{交换群的同态}).
\]
一个 $G$-模称为\emph{诱导的}, 如果它同构于某个交换群 $M_0$ 的 $\Ind^G(M_0)$.

\begin{remark}\label{rem:ii.1.3}
设 $G$ 为有限群.
\begin{enumerate}[(a)]
\item 一个 $G$-模 $M$ 是诱导的当且仅当存在 $M$ 的子交换群 $M_0$ 使得
\[
 M=\bigoplus_{g\in G} g M_0\qquad(\text{交换群的直接和}),
\]
此时有 $G$-模的同构
\[
 \varphi\mapsto \sum_{g\in G} g\otimes \varphi(g^{-1})\colon \Ind^G(M_0)\to \Z[G]\otimes_\Z M_0.
\]
这里 $\Z[G]\otimes_\Z M_0$ 带如下 $G$-结构: $g(z\otimes m)=gz\otimes m$.
\item 设 $H$ 为 $G$ 的子群. 一个诱导 $G$-模在视为 $H$-模时也是诱导的: 取 $H$ 在 $G$ 中的右陪集代表系 $S$, 即 $G=\coprod_{s\in S}Hs$; 若 $M=\bigoplus_{g\in G}gM_0$, 则 $M=\bigoplus_{h\in H}hM_1$, 其中 $M_1=\bigoplus_{s\in S}sM_0$.
\item 设 $M$ 为 $G$-模, $M_0$ 为视为交换群的 $M$. 则
\[
 \pi\colon \Ind^G(M_0)\to M,\qquad \varphi\mapsto \sum_{g\in G} g\varphi(g^{-1})
\]
是一个满 $G$-模同态. 它对应于映射
\[
 \Z[G]\otimes M_0\to M,\qquad \Bigl(\sum n_g g\Bigr)\otimes m\mapsto \sum n_g g m.
\]
\end{enumerate}
\end{remark}

\begin{remark}\label{rem:ii.1.4}
设 $M$ 和 $N$ 为 $G$-模. 则规则
\[
 g(m\otimes n)=gm\otimes gn
\]
在 $M\otimes_\Z N$ 上定义了一个 $G$-模结构. 设 $M_0$ 为视为交换群的 $M$. 则 $\Z[G]\otimes_\Z M=\Z[G]\otimes_\Z M_0$ 作为交换群, 但它们的 $G$-模结构不对应. 然而容易验证, 把 $g\otimes m$ 映到 $g\otimes gm$ 的映射 $\Z[G]\otimes_\Z M_0\to \Z[G]\otimes_\Z M$ 是一个 $G$-模同构.
\end{remark}

\subsubsection*{内射 $G$-模}

一个 $G$-模 $I$ 称为\emph{内射的}, 如果从一个 $G$-模的子模到它的每个 $G$-同态都扩张到整个模上, 等价地, 如果 $\operatorname{Hom}_G(\cdot,I)$ 是一个正合函子.

\begin{proposition}\label{prop:ii.1.5}
范畴 $\Mod_G$ 有足够多的内射对象, 即每个 $G$-模 $M$ 都可以嵌入一个内射 $G$-模 $I$ 中: $M\hookrightarrow I$.
\end{proposition}

\begin{proof}
当 $G=\{1\}$ 时, $\Mod_G$ 就是交换群范畴, 这由附录 (\ref{prop:ii.A.4}) 证得. 现在设 $M$ 为 $G$-模, $M_0$ 为视为交换群的 $M$. 我们可以把 $M_0$ 嵌入一个内射交换群, 比如 $M_0\hookrightarrow I$. 应用函子 $\Ind^G$, 我们得到 $G$-模的含入 $\Ind^G(M_0)\hookrightarrow \Ind^G(I)$. 又存在 $G$-模的含入 $M\hookrightarrow \Ind^G(M_0)$, 把 $m$ 送到函数 $g\mapsto g m$. 把这些映射复合, 我们得到内射同态 $M\hookrightarrow \Ind^G(I)$, 故只需证明 $\Ind^G(I)$ 是一个内射 $G$-模; 但自然同构
\[
 \operatorname{Hom}_{\Mod_G}(M,\Ind^G(I))\simeq \operatorname{Hom}_{\Ab}(M,I)
\]
表明 $\operatorname{Hom}_{\Mod_G}(-,\Ind^G(I))$ 是正合的.
\end{proof}

\noindent\phantomsection\hypertarget{item:ii.1.6}{}\textbf{1.6} 第零个群 $H^0(G,M)=M^G$, 因为
\[
 0\to M^G\to I^{0G}\xrightarrow{d^0} I^{1G}
\]
正合, 而 $H^0(G,M)\stackrel{\mathrm{def}}{=}\dfrac{\Ker(d^0)}{\Im(d^{-1})}=\Ker(d^0)$.

\noindent\phantomsection\hypertarget{item:ii.1.7}{}\textbf{1.7} 如果 $I$ 是内射 $G$-模, 那么对所有 $r>0$ 有 $H^r(G,I)=0$, 因为
\[
 0\to I\to I\to 0\to\cdots
\]
是 $I$ 的一个内射分解.

\noindent\phantomsection\hypertarget{item:ii.1.8}{}\textbf{1.8} 设 $M\to I^\bullet$ 与 $N\to J^\bullet$ 为 $G$-模 $M$ 与 $N$ 的内射分解. 任何 $G$-模同态 $\alpha\colon M\to N$ 都扩张为一个复形的映射
\[
 \begin{tikzcd}
 M \ar[r] \ar[d, "\alpha"'] & I^\bullet \ar[d, "\alpha^\bullet"] \\
 N \ar[r] & J^\bullet.
 \end{tikzcd}
\]
并且同态
\[
 H^r(\alpha^\bullet)\colon H^r(I^\bullet)\to H^r(J^\bullet)
\]
与 $\alpha^\bullet$ 的选取无关. 将此陈述用于恒等映射 $\id\colon M\to M$, 我们看到群 $H^r(G,M)$ 的定义是良定的 (到一个良定的同构). 一般 $\alpha$ 的情形的陈述则表明 $M\mapsto H^r(G,M)$ 是从 $G$-模范畴到交换群范畴的一个函子.

\phantomsection\hypertarget{item:ii.1.9}{}\subsubsection*{1.9 短正合列}

一个 $G$-模的短正合列
\[
 0\to M'\to M\to M''\to 0
\]
产生一个长正合列
\[
 0\to H^0(G,M')\to\cdots\to H^r(G,M)\to H^r(G,M'')\xrightarrow{\delta^r} H^{r+1}(G,M')\to\cdots
\]
而且对应
\[
 \text{短正合列}\mapsto\text{长正合列}
\]
是函子性的, 即一个短正合列的态射诱导一个长正合列的态射 (见 \ref{prop:ii.A.11}).

% 1.6/1.7/1.8 为叙述性条目, 1.9 为子节标题, 均不占定理计数器; 调至 9 使下一条显示为 1.10
\setcounter{theorem}{9}
\begin{remark}\label{rem:ii.1.10}
函子族 $(H^r(G,\cdot))_{r\ge 0}$ 与边缘同态 $\delta^r$ 由性质 (\hyperlink{item:ii.1.6}{1.6}, \hyperlink{item:ii.1.7}{1.7}, \hyperlink{item:ii.1.9}{1.9}) 唯一确定.
\end{remark}

\subsubsection*{Shapiro 引理}

设 $M$ 为 $G$-模, 并把 $\Z$ 视为带平凡作用的 $G$-模: 对所有 $g\in G$, $m\in\Z$ 有 $gm=m$. 一个同态 $\alpha\colon \Z\to M$ 由 $\alpha(1)$ 唯一确定, 而 $m\in M$ 是某个 $G$-同态 $\Z\to M$ 下 $1$ 的像当且仅当它被 $G$ 固定. 因此
\begin{equation*}
 \operatorname{Hom}_G(\Z,M)\simeq M^G.
\tag{13}
\end{equation*}

\begin{proposition}[Shapiro 引理]\label{prop:ii.1.11}
设 $H$ 为 $G$ 的子群. 对每个 $H$-模 $N$, 存在典范同构
\[
 H^r(G,\Ind_H^G(N))\to H^r(H,N)
\]
对所有 $r\ge 0$.
\end{proposition}

\begin{proof}
当 $r=0$ 时, 该同构是复合
\begin{equation*}
 N^H\overset{(13)}{\simeq}\operatorname{Hom}_H(\Z,N)\overset{(\ref{lem:ii.1.2})}{\simeq}\operatorname{Hom}_G(\Z,\Ind_H^G(N))\overset{(13)}{\simeq}\Ind_H^G(N)^G.
\tag{14}
\end{equation*}
现在取 $N$ 的一个内射分解 $N\to I^\bullet$. 应用函子 $\Ind_H^G$, 我们得到 $G$-模 $\Ind_H^G(N)$ 的一个内射分解 $\Ind_H^G(N)\to \Ind_H^G(I^\bullet)$, 因为函子 $\Ind_H^G$ 正合 (\ref{lem:ii.1.2}) 且保持内射性 (\ref{prop:ii.1.5} 的证明). 于是
\[
 H^r(G,\Ind_H^G(N))=H^r((\Ind_H^G(I^\bullet))^G)\stackrel{(14)}{\simeq} H^r(I^{\bullet H})=H^r(H,N).
\]
\end{proof}

\begin{corollary}\label{cor:ii.1.12}
如果 $M$ 是一个诱导 $G$-模, 那么对所有 $r>0$ 有 $H^r(G,M)=0$.
\end{corollary}

\begin{proof}
如果 $M=\Ind^G(M_0)$, 则
\[
 H^r(G,M)\simeq H^r(\{1\},M_0)=0\quad\text{对 }r>0.
\]
\end{proof}

\begin{remark}\label{rem:ii.1.13}
考虑一个 $G$-模的正合列
\[
 0\to M\to J\to N\to 0.
\]
如果对所有 $r>0$ 有 $H^r(G,J)=0$, 那么上同调列变成正合列
\[
 0\to M^G\to J^G\to N^G\to H^1(G,M)\to 0
\]
以及一族同构
\[
 H^r(G,N)\xrightarrow{\approx} H^{r+1}(G,M),\quad r\ge 1.
\]
例如, 设 $M$ 为 $G$-模, $M_*$ 为诱导模 $\Ind^G(M_0)$, 其中 $M_0$ 是视为交换群的 $M$. 如 (\ref{prop:ii.1.5}) 证明中所注, $M$ 可以等同于 $M_*$ 中形如 $g\mapsto gm$ ($m\in M$) 的映射组成的 $G$-子模. 设 $M_\dagger=M_*/M$. 把上述注记用于序列
\begin{equation*}
 0\to M\to M_*\to M_\dagger\to 0
\tag{15}\label{eqn:ii.1.15}
\end{equation*}
我们得到
\[
 H^r(G,M_\dagger)\simeq H^{r+1}(G,M),\quad\text{所有 }r\ge 1.
\]
更一般地, 一个正合列
\[
 0\to M\to J^1\to\cdots\to J^s\to N\to 0
\]
若对所有 $r,i>0$ 有 $H^r(G,J^i)=0$, 则定义同构
\[
 H^r(G,N)\xrightarrow{\simeq} H^{r+s}(G,M),\quad\text{所有 }r\ge 1.
\]
要证明它, 把该序列拆成短正合列
\[
 \begin{aligned}
 &0\to M\to J^1\to N^1\to 0\\
 &0\to N^1\to J^2\to N^2\to 0\\
 &\cdots\\
 &0\to N^{s-1}\to J^s\to N\to 0
 \end{aligned}
\]
并注意到我们有同构
\[
 H^r(G,N)\simeq H^{r+1}(G,N^{s-1})\simeq H^{r+2}(G,N^{s-2})\simeq\cdots.
\]
\end{remark}

\begin{remark}\label{rem:ii.1.14}
设
\[
 0\to M\xrightarrow{\varepsilon} J^0\xrightarrow{d^0} J^1\xrightarrow{d^1} J^2\to\cdots
\]
为一个正合列, 使得对所有 $s>0$ 与所有 $r$ 有 $H^s(G,J^r)=0$. 那么
\[
 H^r(G,M)=H^r(J^{\bullet G}).
\]
此注记适用于每个由诱导模给出的 $M$ 的分解.
\end{remark}

\subsubsection*{用上链描述上同调群}

设 $P_r$ ($r\ge 0$) 为以 $G$ 的元素组成的 $(r+1)$-元组 $(g_0,\ldots,g_r)$ 为基的自由 $\Z$-模, 带如下 $G$-作用:
\[
 g(g_0,\ldots,g_r)=(gg_0,\ldots,gg_r).
\]
注意 $P_r$ 作为 $\Z[G]$-模也是自由的, 以 $\{(1,g_1,\ldots,g_r)\mid g_i\in G\}$ 为基. 定义同态 $d_r\colon P_r\to P_{r-1}$ 为
\[
 d_r(g_0,\ldots,g_r)=\sum_{i=0}^r (-1)^i (g_0,\ldots,\widehat{g_i},\ldots,g_r),
\]
其中符号 $\hat{\cdot}$ 表示略去 $\cdot$. 令 $P_\bullet$ 为
\[
 \cdots\to P_r\xrightarrow{d_r} P_{r-1}\to\cdots\to P_0,
\]
容易验证 $d_{r-1}\circ d_r=0$, 故这是一个复形. 设 $\varepsilon$ 为把每个基元送到 $1$ 的映射 $P_0\to\Z$.

\begin{lemma}\label{lem:ii.1.15}
复形 $P_\bullet\xrightarrow{\varepsilon}\Z\to 0$ 正合.
\end{lemma}

\begin{proof}
取一个元素 $o\in G$, 定义 $k_r\colon P_r\to P_{r+1}$ 为
\[
 k_r(g_0,\ldots,g_r)=(o,g_0,\ldots,g_r).
\]
容易验证 $d_{r+1}\circ k_r+k_{r-1}\circ d_r=1$. 因此若 $d_r(x)=0$, 则 $x=d_{r+1}(k_r(x))$.
\end{proof}

\begin{proposition}\label{prop:ii.1.16}
对每个 $G$-模 $M$,
\[
 H^r(G,M)\simeq H^r(\operatorname{Hom}_G(P_\bullet,M)).
\]
\end{proposition}

\begin{proof}
这由 $P_\bullet\to\Z$ 是 $\Z$ 的一个投射分解这一事实推出 —— 见例 \ref{ex:ii.A.14}.
\end{proof}

$\operatorname{Hom}(P_r,M)$ 中的一个元素可以等同于一个函数 $\varphi\colon G^{r+1}\to M$, 而 $\varphi$ 被 $G$ 固定当且仅当
\[
 \varphi(gg_0,\ldots,gg_r)=g(\varphi(g_0,\ldots,g_r))\quad\text{对所有 }g,g_0,\ldots,g_r\in G.
\]
因此 $\operatorname{Hom}_G(P_r,M)$ 可以等同于满足此条件的 $\varphi$ 的集合 $\widetilde C^r(G,M)$. 这样的 $\varphi$ 称为 $G$ 在 $M$ 中取值的\emph{齐次 $r$-上链}. 由 $d_{r+1}$ 诱导的边缘映射 $\widetilde d^r\colon \widetilde C^r(G,M)\to \widetilde C^{r+1}(G,M)$ 为
\[
 (\widetilde d^r\varphi)(g_0,\ldots,g_{r+1})=\sum_{i=0}^{r+1}(-1)^i\varphi(g_0,\ldots,\widehat{g_i},\ldots,g_{r+1}).
\]
命题\ref{prop:ii.1.16} 表明
\[
 H^r(G,M)\simeq \frac{\Ker(\widetilde d^r)}{\Im(\widetilde d^{r-1})}.
\]
一个齐次上链 $\varphi\colon G^{r+1}\to M$ 由它在元素 $(1,g_1,g_1g_2,\ldots,g_1\cdots g_r)$ 上的值决定. 因此我们引出由所有映射 $\varphi\colon G^r\to M$ 组成的\emph{非齐次 $r$-上链}群 $C^r(G,M)$. 我们令 $G^0=\{1\}$, 于是 $C^0(G,M)=M$. 定义
\[
 d^r\colon C^r(G,M)\to C^{r+1}(G,M)
\]
为 (对 $(d^r\varphi)(g_1,\cdots,g_{r+1})=$)
\[
 g_1\varphi(g_2,\ldots,g_{r+1})+\sum_{j=1}^r (-1)^j\varphi(g_1,\ldots,g_jg_{j+1},\ldots,g_{r+1})+(-1)^{r+1}\varphi(g_1,\ldots,g_r).
\]
定义
\[
 Z^r(G,M)=\Ker(d^r)\qquad(r\text{-上循环群})
\]
与
\[
 B^r(G,M)=\Im(d^{r-1})\qquad(r\text{-上边界群}).
\]

\begin{proposition}\label{prop:ii.1.17}
映射序列
\[
 C^0(G,M)\xrightarrow{d^0} C^1(G,M)\xrightarrow{d^1}\cdots\xrightarrow{d^{r-1}} C^r(G,M)\xrightarrow{d^r} C^{r+1}(G,M)\to\cdots
\]
是一个复形, 即 $d^r\circ d^{r-1}=0$, 并且存在典范同构
\[
 H^r(G,M)\simeq \frac{Z^r(G,M)}{B^r(G,M)}.
\]
\end{proposition}

\begin{proof}
对 $\varphi\in\widetilde C^r(G,M)$, 定义
\[
 \varphi'(g_1,\ldots,g_r)=\varphi(1,g_1,g_1g_2,\ldots,g_1\cdots g_r).
\]
则 $\varphi\mapsto\varphi'$ 是一个双射 $\widetilde C^r(G,M)\to C^r(G,M)$, 把 $\widetilde C^\bullet(G,M)$ 中的边缘映射变成 $C^\bullet(G,M)$ 中的边缘映射.
\end{proof}

\begin{example}\label{ex:ii.1.18}
\begin{enumerate}[(a)]
\item 一个映射 $\varphi\colon G\to M$ 称为\emph{交叉同态}, 如果
\[
 \varphi(\sigma\tau)=\sigma\varphi(\tau)+\varphi(\sigma),\quad\text{所有 }\sigma,\tau\in G.
\]
注意此条件蕴含对 $G$ 的单位元 $e$ 有
\[
 \varphi(e)=\varphi(ee)=e\varphi(e)+\varphi(e)=2\varphi(e),
\]
从而 $\varphi(e)=0$. 对每个 $m\in M$, 映射 $\sigma\mapsto \sigma m-m$ 是一个交叉同态 —— 称为\emph{主交叉同态}. 由该命题,
\[
 H^1(G,M)=\frac{\{\text{交叉同态 }G\to M\}}{\{\text{主交叉同态}\}}.
\]
如果 $G$ 在 $M$ 上的作用是平凡的, 那么交叉同态就是同态, 而主交叉同态为零. 于是此时
\begin{equation*}
 H^1(G,M)\simeq \operatorname{Hom}(G,M).
\tag{16}
\end{equation*}
\item 设 $M$ 为一个交换群 (其合成律写作乘法). $G$ 被 $M$ 的扩张是一个群的正合列
\[
 1\to M\to E\xrightarrow{\pi} G\to 1.
\]
我们设
\[
 \sigma m=s(\sigma)\cdot m\cdot s(\sigma)^{-1},\quad \sigma\in G,\ m\in M,
\]
其中 $s(\sigma)$ 是 $E$ 中映到 $\sigma$ 的任一元素. 因为 $M$ 交换, $\sigma m$ 只依赖于 $\sigma$, 这在 $G$ 上定义了一个作用. 注意
\[
 s(\sigma)\cdot m=\sigma m\cdot s(\sigma),\quad\text{所有 }\sigma\in G,\ m\in M.
\]
现在取 $\pi$ 的一个截面 $s$, 即一个映射 (不必是同态) $s\colon G\to E$ 使得 $\pi\circ s=\id$. 则 $s(\sigma)s(\sigma')$ 与 $s(\sigma\sigma')$ 都映到 $\sigma\sigma'\in G$, 故它们相差一个元素 $\varphi(\sigma,\sigma')\in M$:
\[
 s(\sigma)s(\sigma')=\varphi(\sigma,\sigma')\cdot s(\sigma\sigma').
\]
由
\[
 s(\sigma)(s(\sigma')s(\sigma''))=(s(\sigma)s(\sigma'))s(\sigma'')
\]
我们推出
\[
 \sigma\varphi(\sigma',\sigma'')\cdot\varphi(\sigma,\sigma'\sigma'')=\varphi(\sigma,\sigma')\cdot\varphi(\sigma\sigma',\sigma''),
\]
即 $\varphi\in Z^2(G,M)$. 如果 $s$ 换成另一个截面, $\varphi$ 换成一个上同调的上循环, 因此 $\varphi$ 在 $H^2(G,M)$ 中的类与 $s$ 的选取无关. 每个这样的 $\varphi$ 都来自一个扩张. 这样, $H^2(G,M)$ 分类带有固定 $G$ 在 $M$ 上作用的、$G$ 被 $M$ 的扩张的同构类.
\end{enumerate}
\end{example}

\begin{remark}\label{rem:ii.1.19}
设 $G$ 为一个群, $M$ 为一个 $G$-模. 对 $m\in M$, 令 $\varphi_m\colon G\to M$ 为常值映射 $\sigma\mapsto m$. 则
\[
 (d^1\varphi_m)(\sigma,\tau)=\sigma m-m+m=\sigma m.
\]
特别地, $(d^1\varphi_m)(1,1)=m$. 因此 $H^2(G,M)$ 中的每个类都由一个满足 $\varphi(1,1)=0$ 的 $2$-上循环 $\varphi$ 表示. 这样的 $2$-上循环称为\emph{规范化}的.
\end{remark}

\begin{example}\label{ex:ii.1.20}
设 $\varphi\colon G\to M$ 为一个交叉同态. 对每个 $\sigma\in G$,
\[
 \begin{aligned}
 \varphi(\sigma^2)&=\sigma\varphi(\sigma)+\varphi(\sigma)\\
 \varphi(\sigma^3)&=\varphi(\sigma\cdot\sigma^2)=\sigma^2\varphi(\sigma)+\sigma\varphi(\sigma)+\varphi(\sigma)\\
 &\cdots
 \end{aligned}
\]
\[
 \varphi(\sigma^f)=\sigma^{f-1}\varphi(\sigma)+\cdots+\sigma\varphi(\sigma)+\varphi(\sigma).
\]
因此, 如果 $G$ 是由 $\sigma$ 生成的 $f$ 阶循环群, 那么一个交叉同态 $\varphi$ 由它在 $\sigma$ 上的值 (记为 $m$) 决定, 而 $m$ 满足方程
\begin{equation*}
 \sigma^{f-1}m+\cdots+\sigma m+m=0.
\tag{17}
\end{equation*}
反之, 如果 $m\in M$ 满足 (17), 那么公式 $\varphi(\sigma^i)=\sigma^{i-1}m+\cdots+\sigma m+m$ 定义了一个交叉同态 $\varphi\colon G\to M$, 并且
\[
 \varphi\text{ 是主交叉同态}\iff m=\sigma x-x\text{ 对某个 }x\in M.
\]
定义映射
\[
 \begin{aligned}
 &\Nm_G\colon M\to M,\quad m\mapsto \sum_{\sigma\in G}\sigma m\\
 &\sigma-1\colon M\to M,\quad m\mapsto \sigma m-m.
 \end{aligned}
\]
于是, 对有限阶循环群 $G$, 我们已证明选取 $G$ 的一个生成元 $\sigma$ 决定了一个同构
\[
 \varphi\mapsto\varphi(\sigma)\colon H^1(G,M)\to \frac{\Ker(\Nm_G)}{(\sigma-1)M}.
\]
\end{example}

\begin{remark}\label{rem:ii.1.21}
设
\[
 0\to M\to N\to P\to 0
\]
为 $G$-模的正合列. 边缘映射
\[
 \delta^r\colon H^r(G,P)\to H^{r+1}(G,M)
\]
有如下描述: 设 $\gamma\in H^r(G,P)$ 由 $r$-上循环 $\varphi\colon G^r\to P$ 表示; 因为 $N$ 满射到 $P$, 存在一个提升 $\varphi$ 的 $r$-上链 $\tilde\varphi\colon G^r\to N$; 因为 $d\varphi=0$, $d\tilde\varphi$ 取值于 $M$ —— 它就是表示 $\delta^r\gamma$ 的上循环.
\end{remark}

\subsubsection*{$L$ 与 $L^\times$ 的上同调}

设 $L$ 为域 $K$ 的有限 Galois 扩张, $G=\Gal(L/K)$. 则 $L$ (视为加法群) 与 $L^\times$ 都是 $G$-模.

\begin{proposition}\label{prop:ii.1.22}
设 $L/K$ 为 Galois 群 $G$ 的有限 Galois 扩张. 则 $H^1(G,L^\times)=0$.
\end{proposition}

\begin{proof}
设 $\varphi\colon G\to L^\times$ 为一个交叉同态. 用乘法记号, 这意味着
\[
 \varphi(\sigma\tau)=\sigma\varphi(\tau)\cdot\varphi(\sigma),\quad \sigma,\tau\in G,
\]
而我们要找一个 $c\in L^\times$ 使得 $\varphi(\sigma)=\sigma c/c$. 对 $a\in L^\times$, 令
\[
 b=\sum_{\sigma\in G}\varphi(\sigma)\cdot\sigma a.
\]
假设 $b\neq 0$. 则
\[
 \tau b=\sum_\sigma\tau\varphi(\sigma)\cdot\tau\sigma a=\sum_\sigma\varphi(\tau)^{-1}\varphi(\tau\sigma)\tau\sigma a=\varphi(\tau)^{-1}b.
\]
于是
\[
 \varphi(\tau)=b/\tau b=\tau(b^{-1})/b^{-1},
\]
这表明 $\varphi$ 是主交叉同态.

余下只需证明存在使 $b\neq 0$ 的 $a$. 回顾 (FT, 5.14, Dedekind 关于特征独立的定理): 设 $L$ 为域, $H$ 为群; 则每个不同的同态 $H\to L^\times$ 的有限族 $\{f_i\}$ 在 $L$ 上线性无关, 即
\[
 \sum a_i\, f_i(\alpha)=0\quad\text{对所有 }\alpha\in H\ \Longrightarrow\ a_1=a_2=\cdots=a_n=0.
\]
当我们把它用于同态 $\sigma\colon L^\times\to L^\times$ ($\sigma\in G$) 时, 我们发现 $\sum_\sigma\varphi(\sigma)\sigma$ 不是零映射, 因此存在使 $b\neq 0$ 的 $a$.
\end{proof}

\begin{corollary}\label{cor:ii.1.23}
设 $L/K$ 为循环扩张, $\sigma$ 生成 $\Gal(L/K)$. 如果 $\Nm_{L/K}a=1$, 那么 $a$ 具有形式 $\dfrac{\sigma b}{b}$.
\end{corollary}

\begin{proof}
我们有 $0=H^1(G,L^\times)=\Ker(\Nm_G)/(\sigma-1)L^\times$.
\end{proof}

\begin{proposition}\label{prop:ii.1.24}
设 $L/K$ 为 Galois 群 $G$ 的有限 Galois 扩张. 则对所有 $r>0$ 有 $H^r(G,L)=0$.
\end{proposition}

\begin{proof}
回顾 (FT, 5.18, 正规基定理): 存在 $\alpha\in L$ 使得 $\{\sigma\alpha\mid \sigma\in G\}$ 是 $L$ 作为 $K$-向量空间的一组基 (这种形式的一组基称为\emph{正规}的).

一个正规基 $(\sigma\alpha)_{\sigma\in G}$ 定义一个 $G$-模同构
\[
 \sum_{\sigma\in G} a_\sigma\sigma\mapsto \sum_{\sigma\in G} a_\sigma\sigma\alpha\colon K[G]\to L.
\]
但 $K[G]=\Ind_{\{1\}}^G K$, 因此 $H^r(G,L)\simeq H^r(\{1\},K)=0$ 对 $r>0$ (Shapiro 引理, \ref{prop:ii.1.11}).
\end{proof}

\subsubsection*{乘积的上同调}

一个 $G$-模的乘积 $M=\prod_i M_i$ 在对角作用下成为 $G$-模:
\[
 \sigma(\ldots,m_i,\ldots)=(\ldots,\sigma m_i,\ldots).
\]

\begin{proposition}\label{prop:ii.1.25}
对所有的 $G$-模族 $(M_i)_i$,
\[
 H^r\Bigl(G,\prod_i M_i\Bigr)\simeq \prod_i H^r(G,M_i).
\]
\end{proposition}

\begin{proof}
交换群正合列的乘积仍正合. 由此推出内射 $G$-模的乘积 $I=\prod I_i$ 仍是内射的, 因为
\[
 \operatorname{Hom}_G(\cdot,I)\simeq \prod_i \operatorname{Hom}_G(\cdot,I_i)
\]
正合. 设 $M_i\to I_i^\bullet$ 为 $M_i$ 的内射分解. 则 $\prod M_i\to\prod I_i^\bullet$ 是 $\prod M_i$ 的内射分解, 且
\[
 \begin{array}{r}
 H^r(G,\prod_i M_i)\simeq H^r((\prod_i I_i^{\bullet})^G)\simeq H^r(\prod_i (I_i^{\bullet G}))\simeq \prod_i H^r(I_i^{\bullet G})\simeq \prod_i H^r(G,M_i).
 \end{array}
\]
特别地, 对任何 $G$-模 $M$, $N$,
\[
 H^r(G,M\oplus N)=H^r(G,M)\oplus H^r(G,N).
\]
这也可以通过注意到说一个模 $P$ 是模 $M$ 与 $N$ 的直和意味着某些映射及这些映射之间的关系存在 (见第 \ref{sec:II.3} 节, 本书 p.\pageref{sec:II.3}) 来证明. 这些映射与关系在我们应用加性函子 $H^r(G,\cdot)$ 时保持.
\end{proof}

\begin{remark}\label{rem:ii.1.26}
任意交换群的逆极限的构造不是正合的. 因此一般地, 不能指望上同调与逆极限交换.
\end{remark}

\subsubsection*{上同调群的函子性质}

设 $M$ 与 $M'$ 分别为 $G$ 与 $G'$ 模. 同态
\[
 \alpha\colon G'\to G,\qquad \beta\colon M\to M'
\]
称为\emph{相容的}, 如果
\[
 \beta(\alpha(g)m)=g(\beta(m)).
\]
它们定义映射
\[
 C^\bullet(G,M)\to C^\bullet(G',M'),\quad \varphi\mapsto \beta\circ\varphi\circ\alpha^r,
\]
从而定义同态
\[
 H^r(G,M)\to H^r(G',M').
\]

\begin{example}\label{ex:ii.1.27}
\begin{enumerate}[(a)]
\item 设 $H$ 为 $G$ 的子群. 对每个 $H$-模 $M$, 映射
\[
 \varphi\mapsto\varphi(1_G)\colon \Ind_H^G(M)\to M
\]
与含入 $H\hookrightarrow G$ 相容, 而它诱导的同态
\[
 H^r(G,\Ind_H^G(M))\to H^r(H,M)
\]
就是 Shapiro 引理\ref{prop:ii.1.11} 中的同构. 要证明这两个同构一致, 先对 $r=0$ 验证, 再用维数移位 (II, \ref{rem:ii.1.28}) 推广到所有 $r$.
\item 设 $H$ 为 $G$ 的子群. 令 $\alpha$ 为含入 $H\hookrightarrow G$, $\beta$ 为在 $G$-模 $M$ 上的恒等映射. 此时我们得到\emph{限制}同态
\[
 \Res\colon H^r(G,M)\to H^r(H,M).
\]
它们也可如下构造: 设 $M\to \Ind_H^G(M)$ 为把 $m$ 送到映射 $g\mapsto gm$ 的同态; 它在上同调上定义的同态与 Shapiro 引理中的同构的复合
\[
 H^r(G,M)\to H^r(G,\Ind_H^G(M))\xrightarrow{\simeq} H^r(H,M)
\]
就是限制映射.
\item 设 $H$ 为 $G$ 的正规子群, 令 $\alpha$ 为商映射 $G\to G/H$, $\beta$ 为含入 $M^H\hookrightarrow M$. 此时我们得到\emph{膨胀}同态
\[
 \Inf\colon H^r(G/H,M^H)\to H^r(G,M).
\]
\item 对任何 $g_0\in G$, 同态 $\alpha\colon G\to G$, $\sigma\mapsto g_0\sigma g_0^{-1}$, 与 $\beta\colon M\to M$, $m\mapsto g_0^{-1}m$, 是相容的. 我断言它们定义的同态
\[
 H^r(G,M)\to H^r(G,M)
\]
都是恒等映射. 对 $r=0$, 该同态是
\[
 m\mapsto g_0^{-1}m\colon M^G\to M^G,
\]
这显然是恒等映射. 设 $r>0$ 且已知对 $r-1$ 成立. 由正合列 \ref{eqn:ii.1.15} (本书 p.\pageref{eqn:ii.1.15})
\[
 0\to M\to M_*\to M_\dagger\to 0,
\]
其中 $M_*=\Ind^G(M_0)$, 我们得到一个有正合行的图表
\[
 \begin{tikzcd}
 H^{r-1}(G,M_*) \ar[r] \ar[d] & H^{r-1}(G,M_\dagger) \ar[r] \ar[d] & H^r(G,M) \ar[r] \ar[d] & 0 \\
 H^{r-1}(G,M_*) \ar[r] & H^{r-1}(G,M_\dagger) \ar[r] & H^r(G,M) \ar[r] & 0.
 \end{tikzcd}
\]
右边的 $0$ 来自 $M_*$ 是诱导模. 竖直映射是由这对 $(\alpha,\beta)$ (对不同模) 定义的. 容易验证图表交换. 由归纳, 中间的竖直映射是恒等映射, 从而第三个也是.
\end{enumerate}
\end{example}

\begin{remark}\label{rem:ii.1.28}
\begin{enumerate}[(a)]
\item (d) 中的证明方法称为\emph{维数移位}.
\item 设 $H$ 为 $G$ 的正规子群. 对每个 $G$-模 $M$, (d) 中的做法给出 $G$ 在 $H^r(H,M)$ 上的一个作用, 上述结果表明它经过 $G/H$.
\end{enumerate}
\end{remark}

\begin{example}\label{ex:ii.1.29}
我们还需要一个上同调群的函子映射. 设 $H$ 为 $G$ 的有限指数子群, $S$ 为 $H$ 在 $G$ 中的左陪集代表系, $G=\coprod_{s\in S}sH$. 设 $M$ 为 $G$-模. 对每个 $m\in M^H$,
\[
 \Nm_{G/H}m\stackrel{\mathrm{def}}{=}\sum_{s\in S}sm
\]
与 $S$ 的选取无关, 且被 $G$ 固定. 因此 $\Nm_{G/H}$ 是一个同态 $M^H\to M^G$.

这可以如下扩张为一个\emph{余限制}同态\footnote{余限制同态也称为转移映射.}\index{yuxianzhitonghui@余限制同态(corestriction homomorphism)}
\[
 \Cor\colon H^r(H,M)\to H^r(G,M)
\]
对所有 $r$: 对每个 $G$-模 $M$, 存在一个典范的 $G$-模同态
\[
 \varphi\mapsto \sum_{s\in S}s\varphi(s^{-1})\colon \Ind_H^G M\to M;
\]
它在上同调上定义的映射与 Shapiro 引理中的同构复合, 给出 $\Cor$,
\[
 H^r(H,M)\xrightarrow{\simeq} H^r(G,\Ind_H^G M)\to H^r(G,M).
\]
\end{example}

\begin{proposition}\label{prop:ii.1.30}
设 $H$ 为 $G$ 的有限指数子群. 复合
\[
 \Cor\circ \Res\colon H^r(G,M)\to H^r(G,M)
\]
是乘以 $(G:H)$.
\end{proposition}

\begin{proof}
对 $m\in M$, 令 $\varphi_m$ 为映射 $g\mapsto gm\colon G\to M$. 则 $\Cor\circ\Res$ 是由如下复合在上同调上定义的映射:
\[
 M\to \Ind_H^G(M)\to M,\quad m\mapsto \varphi_m\mapsto \sum_s s\varphi_m(s^{-1})=\sum_s m=(G:H)m.
\]
\end{proof}

\begin{corollary}\label{cor:ii.1.31}
如果 $(G:1)=m$, 那么对所有 $r>0$ 有 $mH^r(G,M)=0$.
\end{corollary}

\begin{proof}
由该命题, 乘以 $m$ 经过 $H^r(\{1\},M)=0$ 分解:
\[
 H^r(G,M)\xrightarrow{\Res} H^r(\{1\},M)\xrightarrow{\Cor} H^r(G,M).
\]
\end{proof}

\begin{corollary}\label{cor:ii.1.32}
如果 $G$ 有限且 $M$ 作为交换群是有限生成的, 那么 $H^r(G,M)$ 有限.
\end{corollary}

\begin{proof}
从上同调群作为上循环模上边界群的描述清楚地看出 $H^r(G,M)$ 作为交换群是有限生成的, 而我们已经看到它被 $(G:1)$ 零化. 因此它是有限的.
\end{proof}

对交换群 $A$ 与素数 $p$, $A$ 的 $p$-主分量 $A(p)$ 是由所有被 $p$ 的某个幂零化的元素组成的子群.

\begin{corollary}\label{cor:ii.1.33}
设 $G$ 为有限群, $G_p$ 为其 Sylow $p$-子群. 对每个 $G$-模 $M$, 限制映射
\[
 \Res\colon H^r(G,M)\to H^r(G_p,M)
\]
在 $H^r(G,M)$ 的 $p$-主分量上是单射.
\end{corollary}

\begin{proof}
复合
\[
 \Cor\circ\Res\colon H^r(G,M)\to H^r(G_p,M)\to H^r(G,M)
\]
是乘以 $(G:G_p)$, 它不被 $p$ 整除, 因此在 $H^r(G,M)$ 的 $p$-主分量上是单射.
\end{proof}

\subsubsection*{膨胀--限制正合列}

\begin{proposition}\label{prop:ii.1.34}
设 $H$ 为 $G$ 的正规子群, $M$ 为 $G$-模. 设 $r$ 为 $>0$ 的整数. 如果对所有满足 $0<j<r$ 的 $j$ 有 $H^j(H,M)=0$, 那么序列
\[
 0\to H^r(G/H,M^H)\xrightarrow{\Inf} H^r(G,M)\xrightarrow{\Res} H^r(H,M)
\]
正合.
\end{proposition}

\begin{proof}
我们先对 $r=1$ 证明. 此时对 $H^j(H,M)$ 的假设是空真的. 设 $\varphi\colon G\to M$ 为一个交叉同态, 其在 $H$ 上的限制是主交叉同态, 即 $\varphi(h)=hm_0-m_0$. 定义 $\varphi'(g)=\varphi(g)-(gm_0-m_0)$. 则 $\varphi'$ 在 $H^1(G,M)$ 中表示同一类, 但现在对所有 $h\in H$ 有 $\varphi'(h)=0$, 因此 $\varphi'$ 由 $G/H\to M$ 的一个交叉同态"膨胀"而来.

我们证明 $\varphi'$ 取值于 $M^H$. 因为 $\varphi'$ 是交叉同态且对 $h\in H$ 有 $\varphi'(h)=0$, 我们有
\[
 \varphi'(hg)=h\varphi'(g)+\varphi'(h)=h\varphi'(g)
\]
以及
\[
 \varphi'(gh)=g\varphi'(h)+\varphi'(g)=\varphi'(g).
\]
第二个方程说 $\varphi'$ 在 $H$ 的左陪集上为常数. 因为 $H$ 正规, 左陪集与右陪集相同, 故 $\varphi'$ 在右陪集上也为常数. 由于 $hg$ 与 $g$ 在同一个右陪集 $Hg$ 中, 我们有 $\varphi'(hg)=\varphi'(g)$. 与第一个方程结合给出 $h\varphi'(g)=\varphi'(g)$, 即 $\varphi'(g)\in M^H$.

我们已证明 $\varphi'$ 由 $G/H\to M^H$ 的一个交叉同态"膨胀"而来, 这就证明了在 $H^1(G,M)$ 处的正合性. 在 $H^1(G/H,M^H)$ 处的正合性容易证明.

现在假设 $r>1$ 且命题对 $r-1$ 成立. 考虑正合列 \ref{eqn:ii.1.15} (本书 p.\pageref{eqn:ii.1.15}),
\[
 0\to M\to M_*\to M_\dagger\to 0.
\]
则
\[
 H^j(H,M_\dagger)\simeq H^{j+1}(H,M),\quad j>0,
\]
故对 $0<j<r-1$ 有 $H^j(H,M_\dagger)=0$. 由归纳, 序列
\[
 0\to H^{r-1}(G/H,M_\dagger^H)\xrightarrow{\Inf} H^{r-1}(G,M_\dagger)\xrightarrow{\Res} H^{r-1}(H,M_\dagger)
\]
正合, 而它同构于序列
\[
 0\to H^r(G/H,M^H)\xrightarrow{\Inf} H^r(G,M)\xrightarrow{\Res} H^r(H,M).
\]
\end{proof}

\begin{remark}\label{rem:ii.1.35}
\begin{enumerate}[(a)]
\item 稍加努力, 可以把该序列扩张为
\[
 \begin{aligned}
 0\to &H^1(G/H,M^H)\xrightarrow{\Inf} H^1(G,M)\xrightarrow{\Res} H^1(H,M)^{G/H}\to\\
 &H^2(G/H,M^H)\xrightarrow{\Inf} H^2(G,M)^*\to H^1(G/H,H^1(H,M))\to H^3(G/H,M^H),
 \end{aligned}
\]
其中 $H^2(G,M)^*=\Ker(\Res\colon H^2(G,M)\to H^2(H,M))$. 见 Dekimpe, Hartl, and Wauters, J. Algebra 369 (2012), 70--95.
\item 对专家而言, Hochschild--Serre (或 Lyndon) 谱序列形如
\[
 H^i(G/H,H^j(H,M))\Rightarrow H^{i+j}(G,M);
\]
这里 $G/H$ 如 (\ref{rem:ii.1.28}b) 中所述作用在 $H^j(H,M)$ 上. 可以由此序列读出一个正合列
\[
 0\to H^1(G/H,M^H)\xrightarrow{\Inf} H^1(G,M)\xrightarrow{\Res} H^1(H,M).
\]
假设 $H^i(H,M)=0$ 对 $0<i<r$ 迫使谱序列的 $r$ 行消失, 并允许把 $1$ 换成 $r$ 读出同一序列.
\end{enumerate}
\end{remark}

\begin{example}\label{ex:ii.1.36}
设 $\Omega\supset L$ 为 $K$ 的 Galois 扩张. 则 $H\stackrel{\mathrm{def}}{=}\Gal(\Omega/L)$ 是 $G\stackrel{\mathrm{def}}{=}\Gal(\Omega/K)$ 的正规子群. 由命题\ref{prop:ii.1.22}, $H^1(H,\Omega^\times)=0$, 故存在正合列
\[
 0\to H^2(G/H,L^\times)\to H^2(G,\Omega^\times)\to H^2(H,\Omega^\times).
\]
一个不涉及维数移位而直接证明此序列正合的方法可见 \cite{Artin1951}, 6.4.
\end{example}

\subsubsection*{杯积}

\begin{lemma}\label{lem:ii.1.37}
正合列 \ref{eqn:ii.1.15} (本书 p.\pageref{eqn:ii.1.15}),
\[
 0\to M\to M_*\to M_\dagger\to 0
\]
作为交换群序列是分裂的, 即 $M_*\approx M\oplus M_\dagger$ (作为交换群).
\end{lemma}

\begin{proof}
回忆 $M_*$ 由映射 $\varphi\colon G\to M$ 组成, 而该序列中的第一个同态把 $M$ 的元素 $m$ 送到 $\varphi_m$, 其中 $\varphi_m(g)=gm$. 对 $\varphi\in M_*$, 令 $s(\varphi)=\varphi(1_G)$. 则 $s$ 是一个交换群同态 $M_*\to M$ 满足 $s(\varphi_m)=m$, 于是
\[
 M_*\simeq \Ker(1-s)\oplus \Ker(s)\simeq M\oplus M_\dagger
\]
作为交换群.
\end{proof}

设 $G$ 为一个群. 对 $G$-模 $M$ 与 $N$, 我们记 $M\otimes N$ 为 $M\otimes_\Z N$ 视为带如下作用的 $G$-模:
\[
 g(m\otimes n)=gm\otimes gn,\quad g\in G,\ m\in M,\ n\in N.
\]

\begin{proposition}\label{prop:ii.1.38}
存在唯一的一族双加性配对
\[
 (m,n)\mapsto m\cup n\colon H^r(G,M)\times H^s(G,N)\to H^{r+s}(G,M\otimes N),
\]
对所有 $G$-模 $M$, $N$ 与所有整数 $r,s\ge 0$ 定义, 满足如下条件:
\begin{enumerate}[(a)]
\item 当两边视为 $(M,N)$ 上的协变双函子时, 这些映射成为函子态射;
\item 当 $r=s=0$ 时, 该配对是
\[
 (m,n)\mapsto m\otimes n\colon M^G\otimes N^G\to (M\otimes N)^G;
\]
\item 如果 $0\to M'\to M\to M''\to 0$ 是 $G$-模的正合列使得
\[
 0\to M'\otimes N\to M\otimes N\to M''\otimes N\to 0
\]
正合, 那么
\[
 (\delta m'')\cup n=\delta(m''\cup n),\quad m''\in H^r(G,M''),\ n\in H^s(G,N).
\]
这里 $\delta$ 表示连接同态 $H^r(G,M'')\to H^{r+1}(G,M')$ 或 $H^{r+s}(G,M''\otimes N)\to H^{r+s+1}(G,M'\otimes N)$.
\item 如果 $0\to N'\to N\to N''\to 0$ 是 $G$-模的正合列使得
\[
 0\to M\otimes N'\to M\otimes N\to M\otimes N''\to 0
\]
正合, 那么
\[
 m\cup \delta n''=(-1)^r\delta(m\cup n''),\quad m\in H^r(G,M),\ n''\in H^s(G,N'').
\]
\end{enumerate}
\end{proposition}

\begin{proof}
唯一性用维数移位证明, 关键的一点是, 因为序列 (15) 作为交换群序列是分裂正合的 (见 \ref{lem:ii.1.37}), 它在与 $G$-模作 $\Z$ 张量积后仍正合. 对存在性, 我们证明如下定义的配对具有所需性质: 设 $m\in H^r(G,M)$ 与 $n\in H^s(G,N)$ 由 (非齐次) 上循环 $\varphi$ 与 $\psi$ 表示; 则 $m\cup n$ 由如下上循环表示:
\[
 (g_1,\ldots,g_{r+s})\mapsto \varphi(g_1,\ldots,g_r)\otimes g_1\cdots g_r\psi(g_{r+1},\cdots,g_{r+s}).
\]
把细节写清需要好几页, 但证明并不困难.
\end{proof}

\begin{proposition}\label{prop:ii.1.39}
如下公式成立:
\begin{enumerate}[(a)]
\item $(x\cup y)\cup z=x\cup (y\cup z)$ (在 $H^{r+s+t}(G,M\otimes N\otimes P)$ 中);
\item $x\cup y=(-1)^{rs}y\cup x$ 若 $x\in H^r(G,M)$, $y\in H^s(G,N)$ (我们等同 $M\otimes N$ 与 $N\otimes M$);
\item $\Res(x\cup y)=\Res(x)\cup \Res(y)$;
\item $\Cor(x\cup \Res y)=\Cor(x)\cup y$;
\item $\Inf(x\cup y)=\Inf(x)\cup \Inf(y)$.
\end{enumerate}
\end{proposition}

\begin{proof}
在每种情形, 先验证当 $x,y,\ldots$ 次数为 $0$ 时的公式, 再用维数移位推出一般情形. 例如, (d) 的详细证明见 \cite{AtiyahWall1967}, 命题 9. 对 (e), 见 \cite{Weiss1969}, 4-3-9.
\end{proof}

一个配对
\begin{equation*}
 x,y\mapsto (x,y)\colon M\times N\to P
\tag{18}
\end{equation*}
使得
\[
 g(x,y)=(gx,gy),\quad\text{所有 }g\in G,
\]
定义一个 $G$-模同态
\[
 M\otimes N\to P,
\]
从而定义映射
\[
 H^r(G,M\otimes N)\to H^r(G,P).
\]
杯积配对与这些映射的复合, 即配对
\[
 H^r(G,M)\times H^s(G,N)\to H^{r+s}(G,P),
\]
将被称为由 (18) 定义的\emph{杯积配对}.

\section{同调}\label{sec:II.2}\index{tongdaoqu@同调群(homology group)}

\subsubsection*{同调群的定义}

对 $G$-模 $M$, 令 $M_G$ 为 $M$ 上 $G$ 作用平凡的最大商群. 于是 $M_G$ 是 $M$ 模去由
\[
 \{gm-m\mid g\in G,\ m\in M\}
\]
生成的子群所得的商. 注意这是 $M^G$ 的对偶概念, $M^G$ 是 $M$ 上 $G$ 作用平凡的最大子对象. 上同调群的定义对偶地给出同调群.

详言之, 函子
\[
 M\mapsto M_G\colon \Mod_G\to \Ab
\]
是右正合的, 即如果
\[
 0\to M'\to M\to M''\to 0
\]
正合, 那么
\[
 M_G'\to M_G\to M_G''\to 0
\]
正合. 这可以直接验证, 或者由注意到从 $M$ 到 $M_G$ 等于与 $\Z[G]/I_G$ 作张量积看出, 其中 $I_G$ 是增广理想
\[
 \Ker\Bigl(\sum n_g g\mapsto \sum n_g\colon \Z[G]\to\Z\Bigr).
\]
一个 Abel 范畴中的对象 $P$ 称为\emph{投射的}, 如果对任意对象 $N$ 与商对象 $M$, 每个态射 $P\to M$ 都提升为一个态射 $P\to N$. 例如, 每个自由 $\Z[G]$-模作为 $\Z[G]$-模 (等价地, $G$-模) 是投射的.

设 $M$ 为 $G$-模. 设 $(m_i)_{i\in I}$ 是 $M$ 作为 $\Z[G]$-模的一族生成元, $\Z[G]^{(I)}$ 是以 $I$ 为指标集的 $\Z[G]$ 的拷贝的直和. 映射 $\sum_{i\in I}\gamma_i\mapsto \sum \gamma_i m_i$ 是一个满 $G$-同态 $\Z[G]^{(I)}\to M$. 这表明每个 $G$-模都是一个自由 $G$-模的商, 因此范畴 $\Mod_G$ 有足够多的投射对象.

设 $M$ 为 $G$-模. 取 $M$ 的一个投射分解
\[
 \cdots\longrightarrow P_2\xrightarrow{d_2} P_1\xrightarrow{d_1} P_0\xrightarrow{d_0} M\longrightarrow 0.
\]
复形
\[
 \cdots\longrightarrow (P_2)_G\xrightarrow{d_2} (P_1)_G\xrightarrow{d_1} (P_0)_G\longrightarrow 0
\]
未必正合, 我们令
\[
 H_r(G,M)=\frac{\Ker(d_r)}{\Im(d_{r+1})}.
\]
这些群有如下基本性质.

\noindent\phantomsection\hypertarget{item:ii.2.1}{}\textbf{2.1} 第零个群 $H_0(G,M)=M_G$, 因为
\[
 P_{1G}\to P_{0G}\to M_G\to 0
\]
正合.

\noindent\phantomsection\hypertarget{item:ii.2.2}{}\textbf{2.2} 如果 $P$ 是投射 $G$-模, 那么对所有 $r>0$ 有 $H_r(G,P)=0$, 因为
\[
 \cdots\to 0\to P\to P\to 0
\]
是 $P$ 的一个投射分解.

\noindent\phantomsection\hypertarget{item:ii.2.3}{}\textbf{2.3} 设 $P_\bullet\to M$ 与 $Q_\bullet\to N$ 为 $G$-模 $M$ 与 $N$ 的投射分解. 任何 $G$-模同态 $\alpha\colon M\to N$ 都扩张为一个复形的态射
\[
 \begin{tikzcd}
 P_\bullet \ar[r] \ar[d, "\alpha_\bullet"'] & M \ar[d, "\alpha"] \\
 Q_\bullet \ar[r] & N
 \end{tikzcd}
\]
并且同态
\[
 H_r(\alpha_\bullet)\colon H_r(P_\bullet)\to H_r(Q_\bullet)
\]
与 $\alpha_\bullet$ 的选取无关. 当我们把此陈述用于恒等映射 $\id\colon M\to M$ 时, 我们发现群 $H_r(G,M)$ 的定义是良定的 (到一个良定的同构). 一般 $\alpha$ 的情形的陈述则表明 $M\mapsto H_r(G,M)$ 是从 $G$-模范畴到交换群范畴的一个函子.

\phantomsection\hypertarget{item:ii.2.4}{}\subsubsection*{2.4 短正合列}

一个 $G$-模的短正合列
\[
 0\to M'\to M\to M''\to 0
\]
产生一个长正合列
\[
 \cdots\to H_r(G,M)\to H_r(G,M'')\xrightarrow{\delta_r} H_{r-1}(G,M')\to\cdots\to H_0(G,M'')\to 0.
\]
而且对应
\[
 \text{短正合列}\mapsto\text{长正合列}
\]
是函子性的, 即一个短正合列的态射诱导一个长正合列的态射.

% 2.1/2.2/2.3 为叙述性条目, 2.4 为子节标题, 均不占定理计数器; 调至 4 使下一条显示为 2.5
\setcounter{theorem}{4}
\begin{remark}\label{rem:ii.2.5}
函子族 $(H_r(G,\cdot))_{r\ge 0}$ 与边缘映射 $\delta_r$ 由性质 (\hyperlink{item:ii.2.1}{2.1}, \hyperlink{item:ii.2.2}{2.2}, \hyperlink{item:ii.2.4}{2.4}) 唯一确定.
\end{remark}

正如上同调的情形, 也可以把 $H_r(G,M)$ 明确描述为 $r$-循环群模去 $r$-边界子群 —— 见本章后面的参考文献.

\subsubsection*{群 $H_1(G,\Z)$}

只用上面列出的同调群的性质, 我们将计算 $H_1(G,\Z)$.

增广映射是
\[
 \Z[G]\to\Z,\quad \sum n_g g\mapsto \sum n_g.
\]
它的核称为\emph{增广理想} $I_G$. 显然 $I_G$ 是以 $\{g-1\mid g\in G,\ g\neq 1\}$ 为基的自由 $\Z$-子模, 因此
\[
 M/I_G M=M_G\stackrel{(\hyperlink{item:ii.2.1}{2.1})}{\simeq} H_0(G,M).
\]
考虑正合 (增广) 序列:
\begin{equation*}
 0\to I_G\to \Z[G]\to\Z\to 0.
\tag{19}
\end{equation*}
$G$-模 $\Z[G]$ 是投射的 (因为它是自由 $\Z[G]$-模), 故 $H_1(G,\Z[G])=0$. 因此我们得一族同调群的正合列
\[
 0\to H_1(G,\Z)\to I_G/I_G^2\to \Z[G]/I_G\to\Z\to 0.
\]
中间的映射由含入 $I_G\hookrightarrow \Z[G]$ 诱导, 故为零. 因此该序列表明
\begin{equation*}
 H_1(G,\Z)\xrightarrow{\simeq} I_G/I_G^2
\tag{20}
\end{equation*}
以及
\[
 \Z[G]_G\xrightarrow{\simeq}\Z,
\]
即 $\Z$ 是 $\Z[G]$ 上 $G$ 作用平凡的最大商群. 注意 $I_G^2\stackrel{\mathrm{def}}{=}I_G\cdot I_G$ 是 $\Z[G]$ 的由形如
\[
 (g-1)(g'-1),\quad g,g'\in G
\]
的元素生成的 $\Z$-子模.

\begin{lemma}\label{lem:ii.2.6}
设 $G^c$ 为 $G$ 的换位子子群, 于是 $G/G^c$ 是 $G$ 的最大交换商 $G^{\mathrm{ab}}$. 则映射 $g\mapsto (g-1)+I_G^2$ 诱导一个同构
\begin{equation*}
 \frac{G}{G^c}\to \frac{I_G}{I_G^2}.
\tag{21}
\end{equation*}
\end{lemma}

\begin{proof}
考虑映射
\[
 g\mapsto (g-1)+I_G^2\colon G\to I_G/I_G^2.
\]
这是一个同态, 因为
\[
 \begin{aligned}
 gg'-1&=(g-1)(g'-1)+(g-1)+(g'-1)\\
 &\equiv (g-1)+(g'-1)\bmod I_G^2.
 \end{aligned}
\]
因为 $I_G/I_G^2$ 交换, 该映射经过 $G^{\mathrm{ab}}$. 为证明它是同构, 我们构造逆. 回忆 $I_G$ 以元素 $g-1$ 为基自由生成. 考虑同态 $I_G\to G^{\mathrm{ab}}$ 把 $g-1$ 映到 $g$ 的类. 由恒等式
\[
 (g-1)(g'-1)=(gg'-1)-(g-1)-(g'-1)
\]
我们看到 $(g-1)(g'-1)$ 映到 $1$. 因为 $I_G^2$ 由这种形式的元素生成, 这表明该映射经过 $I_G/I_G^2$. 我们构造的两个映射互逆.
\end{proof}

\begin{proposition}\label{prop:ii.2.7}
存在一个典范同构
\[
 H_1(G,\Z)\simeq G^{\mathrm{ab}}.
\]
\end{proposition}

\begin{proof}
把同构 (20) 与 (21) 的逆复合.
\end{proof}

\begin{aside}\label{as:ii.2.8}
对任何群 $G$, 存在一个拓扑空间 $BG$, 称为 $G$ 的\emph{分类空间}, 使得 $G=\pi_1(BG)$ 且对所有 $r$ 有 $H_r(BG,\Z)=H_r(G,\Z)$ (\cite{Rosenberg1994}, 5.1.16, 5.1.27). 用 $BG$ 的术语, 该命题简单地说就是 $H_1(BG,\Z)\simeq \pi_1(BG)^{\mathrm{ab}}$.
\end{aside}

\section{Tate 群}\label{sec:II.3}\index{tateshangtongqun@Tate 上同调群(Tate cohomology group)}

本节中 $G$ 始终为有限群.

对 $G$-模 $M$, \emph{范数映射} $\Nm_G\colon M\to M$ 定义为
\[
 m\mapsto \sum_{g\in G} g m.
\]
设 $g'\in G$. 当 $g$ 跑遍 $G$ 的元素时, $g'g$ 与 $gg'$ 也跑遍 $G$ 的元素, 故 $g'(\Nm_G(m))=\Nm_G(m)=\Nm_G(g'm)$. 因此
\[
 \Im(\Nm_G)\subset M^G,\qquad I_G M\subset \Ker(\Nm_G).
\]
于是我们得到一个交换的正合图表.
\[
 0\longrightarrow \Ker(\Nm_G)/I_G M\longrightarrow M/I_G M\longrightarrow M^G\longrightarrow M^G/\Nm_G(M)\longrightarrow 0.
\]
因为 $H_0(G,M)=M/I_G M$ 且 $H^0(G,M)=M^G$, 此图底行可改写为
\[
 0\to \Ker(\Nm_G)/I_G M\to H_0(G,M)\xrightarrow{\Nm_G} H^0(G,M)\to M^G/\Nm_G(M)\to 0.
\]
Tate 定义
\[
 H_T^r(G,M)=
 \begin{cases}
 H^r(G,M) & r>0\\
 M^G/\Nm_G(M) & r=0\\
 \Ker(\Nm_G)/I_G M & r=-1\\
 H_{-r-1}(G,M) & r<-1.
 \end{cases}
\]
于是该正合列现在变成
\[
 0\to H_T^{-1}(G,M)\to H_0(G,M)\xrightarrow{\Nm_G} H^0(G,M)\to H_T^0(G,M)\to 0.
\]
群 $H_T^r(G,M)$ 称为\emph{Tate 上同调群}. 常把 $H_T^r(G,M)$ 记作 $\hat H^r(G,M)$. 因为不会引起歧义, 当 $r>0$ 时我们有时略去下标 $T$.

对任何 $G$-模的短正合列
\[
 0\to M'\to M\to M''\to 0
\]
我们得到一个图表 (见下页图片). 对图的中部应用扩张蛇形引理 (\ref{lem:ii.A.1}), 我们得到一个 (很) 长的正合列
\begin{equation*}
 \begin{aligned}
 &\cdots\to H_T^r(G,M')\to H_T^r(G,M)\to H_T^r(G,M'')\xrightarrow{\delta} H_T^{r+1}(G,M)\to\cdots\\
 &(-\infty<r<\infty).
 \end{aligned}
\tag{22}
\end{equation*}

\begin{proposition}\label{prop:ii.3.1}
如果 $M$ 是诱导的, 那么对所有 $r\in\Z$ 有 $H_T^r(G,M)=0$.
\end{proposition}

\begin{proof}
对 $r>0$, 这在 (\ref{cor:ii.1.12}) 中已证.

情形 $r=0$. 回忆 (\ref{rem:ii.1.3}) $M\simeq \Z[G]\otimes_\Z X$ 对某个交换群 $X$, 因此 $M$ 的元素 $m$ 可唯一写成 $m=\sum_g(g\otimes x_g)$ 且 $x_g\in X$. 如果 $m$ 被 $g'\in G$ 固定, 则 $\sum_g(g'g\otimes x_g)=\sum_g(g\otimes x_g)$; 特别地,\footnote{这里 $e$ 是 $G$ 的单位元.} $g'\otimes x_e=g'\otimes x_{g'}$, 即 $x_e=x_{g'}$. 因此若 $m\in M^G$, 则对所有 $g\in G$ 有 $x_g=x_e$, 且 $m=(\sum_g g)\otimes x_e=\Nm_G(e\otimes x_e)$.

情形 $r=-1$. 如果 $\Nm_G(\sum_g g\otimes x_g)=0$, 则 $\sum_g x_g=0$, 于是
\[
 \sum_g(g\otimes x_g)=\sum_g((g-1)\otimes x_g)\in I_G M.
\]

情形 $r<-1$. 设 $M=\Z[G]\otimes_\Z X$. 把 $X$ 写成一个自由交换群 $X_0$ 的商 $X_0\to X$. 该映射的核也将是一个自由交换群,\footnote{我们用到了: 自由交换群的每个子群都是自由交换群. Darij Grinberg 指出, 这一事实绝非平凡 (且是非构造性的), 而且可以避免使用它.} 于是我们有一个正合列
\[
 0\to X_1\to X_0\to X\to 0
\]
其中 $X_0$ 与 $X_1$ 为自由交换群. 把它与 $\Z[G]$ 作张量积, 我们得到一个 $\Z[G]$-模的正合列
\begin{equation*}
 0\to M_1\to M_0\to M\to 0.
\tag{23}
\end{equation*}
现在
\[
 H_T^r(G,M_0)=0=H_T^r(G,M_1)
\]
对 $r<-1$ 成立 (因为 $M_0$ 与 $M_1$ 是自由 $\Z[G]$-模), 对 $r\ge -1$ 也成立 (因为 $M_0$ 与 $M_1$ 是诱导模). (23) 的 (很) 长正合列现在表明对所有 $r$ 有 $H_T^r(G,M)=0$.
\end{proof}

我们知道 \ref{eqn:ii.1.15} (本书 p.\pageref{eqn:ii.1.15}), 每个 $G$-模 $M$ 都是一个诱导模的子模. 它也是一个诱导 $G$-模的商, 因为把增广序列 (19) 与 $M$ 作张量积, 我们得到一个 $G$-模的正合列
\[
 0\to M'\to \Z[G]\otimes_\Z M\to M\to 0.
\]
这允许我们在两个方向上都用维数移位证明结果. 利用它, 可以把我们对群 $H^r(G,M)$ ($r\ge 0$) 证明的大部分结果推广到所有 $r$. 例如, Shapiro 引理及其推论都成立, 并且我们在 (\ref{ex:ii.1.27}) 与 (\ref{ex:ii.1.29}) 中定义的限制与余限制映射推广到所有 Tate 上同调群. 具体地说, 存在典范同态.\footnote{然而, 膨胀映射 $\Inf\colon H_T^r(G/H,M^H)\to H_T^r(G,M)$ ($H$ 为 $G$ 的正规子群) 只对 $r\ge 1$ 有定义:}\index{xianzhitonghui@限制同态(restriction homomorphism)}\index{yuxianzhitonghui@余限制同态(corestriction homomorphism)}
\[
 \Res\colon H_T^r(G,M)\to H_T^r(H,M)\quad(H\text{ 为 }G\text{ 的子群});
\]
\[
 \Cor\colon H_T^r(H,M)\to H_T^r(G,M)\quad(H\text{ 为 }G\text{ 的子群}).
\]
复合 $\Res\circ\Cor$ 仍是乘以 $(G:H)$. 因为对所有 $r$ 有 $H_T^r(\{1\},M)=0$, (\ref{prop:ii.1.30}) 证明中的论证表明 $H_T^r(G,M)$ 被 $|G|$ 零化.

注意
\[
 H_T^{-2}(G,\Z)=H_1(G,\Z)\simeq G/G^c.
\]

\begin{proposition}\label{prop:ii.3.2}
设 $H$ 为 $G$ 的子群.
\begin{enumerate}[(a)]
\item 映射 $\Cor\colon H_T^{-2}(H,\Z)\to H_T^{-2}(G,\Z)$ 对应于由含入 $H\hookrightarrow G$ 诱导的映射 $H/H^c\to G/G^c$.
\item 映射 $\Res\colon H_T^{-2}(G,\Z)\to H_T^{-2}(H,\Z)$ 对应于\emph{转移}映射 $G/G^c\to H/H^c$ (V, \ref{prop:v.3.18}).
\end{enumerate}
\end{proposition}

\begin{proof}
见 \cite{Weiss1969}, 3-5-5, 3-5-7. 我将在笔记的未来版本中加入证明.
\end{proof}

对所有群 $H_T^r(G,M)$ 有一个统一的明确描述. 事实上, 存在一个 $G$-模的明确复形 $L_\bullet$ (双向无限), 使得
\[
 H_T^r(G,M)=H^r(\operatorname{Hom}_G(L_\bullet,M)).
\]
对 $r>0$ 这直接导出上面我们用非齐次上循环与上边界给出的 $H^r(G,M)$ 的描述. 见 \cite{Weiss1969}, 第 I 章.

\subsubsection*{杯积}

当群 $G$ 有限时, 杯积以唯一的方式推广到所有上同调群, 并具有同样的属性列表. 见 \cite{Weiss1969}, 第 IV 章. 笔记的未来版本将包含我们需要的全部属性的完整列表 (但不包含证明, 它们冗长但不难).

\bigskip

我们首先计算 $\Q$, $\Z$ 与 $\Q/\Z$ 视为带平凡作用的 $G$-模时的上同调.

\begin{lemma}\label{lem:ii.3.3}
对每个有限群 $G$,
\begin{enumerate}[(a)]
\item 对所有 $r\in\Z$ 有 $H_T^r(G,\Q)=0$;
\item $H_T^0(G,\Z)=\Z/(G:1)\Z$ 且 $H^1(G,\Z)=0$;
\item 存在一个典范同构
\[
 \Hom(G,\Q/\Z)\to H^2(G,\Z).
\]
\end{enumerate}
\end{lemma}

\begin{proof}
(a) 群 $\Q$ 唯一可除, 即对所有非零整数 $m$, 乘以 $m\colon \Q\to\Q$ 是同构. 因此映射 $H_T^r(m)\colon H_T^r(G,\Q)\to H_T^r(G,\Q)$ (也是乘以 $m$) 是同构. 当我们取 $m=(G:1)$ 时, 发现乘以 $m$ 在 $H_T^r(G,\Q)$ 上既为零 (见 \ref{prop:ii.1.30}) 又是同构, 这只有在 $H_T^r(G,\Q)=0$ 时才可能.

(b) 因为 $G$ 平凡地作用在 $\Z$ 上, $\Z^G=\Z$ 且范数映射是乘以 $(G:1)$. 故 $H_T^0(G,\Z)=\Z/(G:1)\Z$. 此外 $H^1(G,\Z)=\Hom(G,\Z)$ (见 \ref{ex:ii.1.18}a), 但因为 $\Z$ 无挠, 所有同态 $G\to\Z$ 都为零.

(c) 序列
\[
 0\to\Z\to\Q\to\Q/\Z\to 0
\]
的上同调列是一个正合列
\[
 \begin{tikzcd}[column sep=small]
 H^1(G,\Q)\ar[r]\ar[d,equal] & H^1(G,\Q/\Z)\ar[r]\ar[d,equal] & H^2(G,\Z)\ar[r] & H^2(G,\Q)\ar[d,equal]\\
 0 & \Hom(G,\Q/\Z) & & 0
 \end{tikzcd}
\]
\end{proof}

设 $G$ 为 $m$ 阶有限循环群, 生成元为 $\sigma$. 那么由定义,
\[
 H_T^0(G,M)=\frac{\Ker(\sigma-1)}{\Im(\Nm_G)},\qquad H_T^{-1}(G,M)=\frac{\Ker(\Nm_G)}{\Im(\sigma-1)}.
\]
回忆在 (\ref{ex:ii.1.20}) 中我们证明了 $H^1(G,M)$ 同构于与 $H_T^{-1}(G,M)$ 相同的群. 事实上, 有限循环群的上同调群以 $2$ 为周期.

\begin{proposition}\label{prop:ii.3.4}
设 $G$ 为有限阶循环群. 选取 $G$ 的一个生成元决定同构
\[
 H_T^r(G,M)\xrightarrow{\simeq} H_T^{r+2}(G,M)
\]
对所有 $G$-模 $M$ 与所有 $r\in\Z$.
\end{proposition}

\begin{proof}
设 $\sigma$ 生成 $G$. 则序列
\[
 0\to \Z\xrightarrow{m\mapsto \sum_{g\in G} g m} \Z[G]\xrightarrow{\sigma-1} \Z[G]\xrightarrow{\sigma^i\mapsto 1} \Z\to 0
\]
正合. 因为该序列中的群与 $\Z[G]\to\Z$ 的核 $I_G$ 都是自由 $\Z$-模, 该序列与 $M$ 作张量积后仍正合. 于是
\[
 0\to M\to \Z[G]\otimes_\Z M\to \Z[G]\otimes_\Z M\to M\to 0
\]
是一个 $G$-模的正合列. 回忆 (\ref{rem:ii.1.4}) $\Z[G]\otimes_\Z M\approx \Z[G]\otimes_\Z M_0$, 其中 $M_0$ 是视为交换群的 $M$, 因此由 (\ref{prop:ii.3.1}) 对所有 $r$ 有 $H_T^r(G,\Z[G]\otimes_\Z M)=0$. 于是 (见 \ref{rem:ii.1.13}) 该序列定义同构
\[
 H_T^r(G,M)\xrightarrow{\simeq} H_T^{r+2}(G,M)
\]
对所有 $r$.
\end{proof}

\begin{remark}\label{rem:ii.3.5}
设 $\gamma$ 为 $H^2(G,\Z)$ 中对应于同构 $H^2(G,\Z)\simeq \Hom(G,\Q/\Z)$ 下把所选生成元 $\sigma$ 送到 $1/m$ 的映射的元素. 则映射 $H_T^r(G,M)\to H_T^{r+2}(G,M)$ 为 $x\mapsto x\cup\gamma$. 在未来版本中我将证明此结论. Serre, Local Fields, VIII, \S 4 给出如下提示: "周期同构由与 $\gamma$ 的杯积给出: 例如这由 Cartan--Eilenberg (Homological Algebra) p. 252 给出的杯积公式推出." 另见 \cite{Weiss1969}, 4-5-10.
\end{remark}

设 $G$ 为有限循环群, $M$ 为 $G$-模. 当上同调群 $H^r(G,M)$ 有限时, 我们定义 $M$ 的\emph{Herbrand 商}\index{herbandshang@Herbrand 商(Herbrand quotient)}为
\[
 h(M)=\frac{|H_T^0(G,M)|}{|H_T^1(G,M)|}.
\]

\begin{proposition}\label{prop:ii.3.6}
设 $0\to M'\to M\to M''\to 0$ 为 $G$-模的正合列. 如果 Herbrand 商 $h(M')$, $h(M)$, $h(M'')$ 中有两个有定义, 那么第三个也有定义, 且
\[
 h(M)=h(M')h(M'').
\]
\end{proposition}

\begin{proof}
我们可以如下截断 (很) 长的正合上同调列:
\[
 0\to K\to H_T^0(M')\to H_T^0(M)\to H_T^0(M'')\to H^1(M')\to H^1(M)\to H^1(M'')\to K'\to 0,
\]
其中
\[
 K=\Coker(H_T^{-1}(M)\to H_T^{-1}(M''))\approx \Coker(H^1(M)\to H^1(M''))=K'.
\]
第一个陈述现在显然, 第二个由下一个引理推出.
\end{proof}

\begin{lemma}\label{lem:ii.3.7}
设
\[
 0\to A_0\to A_1\to\cdots\to A_r\to 0
\]
为有限群的一个正合列. 则
\[
 \frac{|A_0|\cdot |A_2|\cdots\cdot}{|A_1|\cdot |A_3|\cdots\cdot}=1.
\]
\end{lemma}

\begin{proof}
对短正合列, 即 $r=2$, 这是显然的, 但每个正合列都可拆成短正合列,
\[
 \begin{aligned}
 0\to A_0\to A_1\to C_1\to 0\\
 0\to C_1\to A_2\to C_2\to 0\\
 \cdots\\
 0\to C_{r-1}\to A_{r-1}\to A_r\to 0.
 \end{aligned}
\]
这里 $C_i=\Coker(A_{i-1}\to A_i)=\Ker(A_{i+1}\to A_{i+2})$. 从这些序列我们发现
\[
 1=\frac{|A_0|\cdot |C_1|}{|A_1|}=\frac{|A_0|\cdot |A_2|}{|A_1|\cdot |C_2|}=\cdots.
\]
\end{proof}

\begin{proposition}\label{prop:ii.3.8}
如果 $M$ 是有限模, 那么 $h(M)=1$.
\end{proposition}

\begin{proof}
考虑正合列
\[
 0\to M^G\to M\xrightarrow{g-1} M\to M_G\to 0
\]
与
\[
 0\to H_T^{-1}(M)\to M_G\xrightarrow{\Nm_G} M^G\to H_T^0(M)\to 0,
\]
其中 $g$ 是 $G$ 的任意生成元. 由第一个序列我们发现 $M^G$ 与 $M_G$ 有相同的阶, 再由第二个序列得 $H_T^{-1}(M)$ 与 $H_T^0(M)$ 有相同的阶.
\end{proof}

\begin{corollary}\label{cor:ii.3.9}
设 $\alpha\colon M\to N$ 为一个有有限核与有限余核的 $G$-模同态. 如果 $h(M)$ 或 $h(N)$ 有定义, 那么另一个也有定义, 且它们相等.
\end{corollary}

\begin{proof}
假设 $h(N)$ 有定义, 考虑正合列:
\[
 0\to \alpha(M)\to N\to \Coker(\alpha)\to 0,
\]
\[
 0\to \Ker(\alpha)\to M\to \alpha(M)\to 0.
\]
由第一个序列我们发现 $h(\alpha M)$ 有定义且等于 $h(N)$, 由第二个序列发现 $h(M)$ 有定义且等于 $h(\alpha M)$.
\end{proof}

\subsubsection*{Tate 定理}

在本节的余下部分, 所有上同调群都是 Tate 群, 因此除了主要陈述外我们略去下标 $T$.

\begin{theorem}\label{thm:ii.3.10}
设 $G$ 为有限群, $M$ 为 $G$-模. 如果
\[
 H^1(H,M)=0=H^2(H,M)
\]
对所有 $G$ 的子群 $H$ 成立, 那么对所有 $r\in\Z$ 有 $H_T^r(G,M)=0$.
\end{theorem}

\begin{proof}
如果 $G$ 循环, 这由上同调的周期性推出.

现在假设 $G$ 可解. 我们将用对 $G$ 的阶的归纳来证明此情形.

因为 $G$ 可解, 它包含一个真正规子群 $H$ 使得 $G/H$ 循环. 因为 $H$ 的阶小于 $G$ 的阶, 且对 $(H,M)$ 满足定理的假设, 对所有 $r$ 有 $H^r(H,M)=0$. 因此 (见 \ref{prop:ii.1.34}) 我们拥有正合列
\begin{equation*}
 0\to H^r(G/H,M^H)\to H^r(G,M)\to H^r(H,M)
\tag{$*$}
\end{equation*}
对所有 $r\ge 1$. 因为 $H^1(G,M)=0=H^2(G,M)$, 有 $H^1(G/H,M^H)=H^2(G/H,M^H)=0$, 而因为 $G/H$ 循环, 这表明对所有 $r$ 有 $H^r(G/H,M^H)=0$. 因此序列 (*) 表明对所有 $r>0$ 有 $H^r(G,M)=0$. 我们接着证明 $H^0(G,M)=0$. 设 $x\in M^G$. 因为 $H^0(G/H,M^H)=0$, 存在 $y\in M^H$ 使得 $\Nm_{G/H}(y)=x$, 而因为 $H^0(H,M)=0$, 存在 $z\in M$ 使得 $\Nm_H(z)=y$. 现在
\[
 \Nm_G(z)=(\Nm_{G/H}\circ \Nm_H)(z)=x.
\]
于是我们现在知道对所有 $r\ge 0$ 有 $H^r(G,M)=0$.

为了进一步推进, 我们使用由增广序列 (19) 与 $M$ 作张量积得到的精确序列
\[
 0\to M'\to \Z[G]\otimes_\Z M\to M\to 0.
\]
回忆
\[
 \Z[G]\otimes_\Z M\overset{1.4}{\simeq}\Z[G]\otimes_\Z M_0\overset{1.3}{\simeq}\Ind^G(M_0)
\]
是诱导的, 因此对所有 $r$ 与所有 $G$ 的子群 $H$ 有 $H^r(H,\Z[G]\otimes_\Z M)=0$. 于是
\[
 H^r(H,M)\simeq H^{r+1}(H,M')
\]
对所有 $r$ 与所有 $H$. 特别地, $M'$ 满足定理的假设, 故 (由已证部分) 对所有 $r\ge 0$ 有 $H^r(G,M')=0$. 特别地,
\[
 0=H^0(G,M')=H^{-1}(G,M).
\]
重复此论证给出 $H^{-2}(G,M)=0$, 等等. 这证明了 $G$ 可解时的定理.

现在考虑任意有限群 $G$ 的情形. 如果 $G$ 与 $M$ 满足定理的假设, 那么 $G_p$ 与 $M$ 也满足, 其中 $G_p$ 是一个 Sylow $p$-子群; 由于 $G_p$ 是 $p$-群 (从而可解), 前面的论证给出对所有 $r$ 有 $H_T^r(G_p,M)=0$. 对 $r>0$, 这表明 (见 \ref{cor:ii.1.33}) $H^r(G,M)$ 的 $p$-主分量对所有 $p$ 为零, 从而 $H^r(G,M)=0$. 对 $r=0$: 设 $x\in M^G$. 由 $H_T^0(G_p,M)=M^{G_p}/\Nm_{G_p}(M)=0$ 知 $x=\Nm_{G_p}(z)$ (某个 $z\in M$), 于是 $[G:G_p]\cdot x=\Nm_G(z)\in\Nm_G(M)$, 即 $[G:G_p]$ 零化 $H_T^0(G,M)$; 对所有素数 $p$ 取遍, 各 $[G:G_p]$ 无公共素因子, 故 $H_T^0(G,M)=0$. 对 $r<0$, 用维数移位: 由正合列 $0\to M'\to\Z[G]\otimes_\Z M\to M\to 0$ 及 $\Z[G]\otimes_\Z M$ 的 Tate 群在所有次数为零, 得 $H_T^r(G,M)\simeq H_T^{r+1}(G,M')$ 对所有 $r$ 成立, 且 $M'$ 及其逐次移位得到的模仍满足定理假设, 于是 $r<0$ 的情形归结为已证的 $r=0$ 情形. 这表明对所有 $r$ 有 $H_T^r(G,M)=0$.\erratum{原书此段对一切 $r$ 直接引用 1.33; 但该 Sylow 归约仅对 $r>0$ 有效, 而维数移位论证还需要 $r=0$ 的情形. 依 Milne errata (K. Kedlaya) 补正如上.}
\end{proof}

\begin{theorem}[Tate 定理]\label{thm:ii.3.11}
设 $G$ 为有限群, $C$ 为 $G$-模. 假设对所有 $G$ 的子群 $H$ (包括 $H=G$)，
\begin{enumerate}[(a)]
\item $H^1(H,C)=0$, 且
\item $H^2(H,C)$ 是一个阶等于 $(H:1)$ 的循环群.
\end{enumerate}
那么对所有 $r$, 存在同构
\[
 H_T^r(G,\Z)\to H_T^{r+2}(G,C)
\]
只依赖于选取 $H^2(G,C)$ 的一个生成元.
\end{theorem}

\begin{proof}
选取 $H^2(G,C)$ 的一个生成元 $\gamma$. 因为 $\Cor\circ\Res=(G:H)$, 对任何子群 $H$ 有 $\Res(\gamma)$ 生成 $H^2(H,C)$.

设 $\varphi$ 为表示 $\gamma$ 的一个上循环. 定义 $C(\varphi)$ 为 $C$ 与以符号 $x_\sigma$ 为基的自由交换群的直接和, 每个 $\sigma\in G$, $\sigma\neq 1$, 并把 $G$ 在 $C$ 上的作用扩张为在 $C(\varphi)$ 上的作用, 规定
\[
 \sigma x_\tau=x_{\sigma\tau}-x_\sigma+\varphi(\sigma,\tau).
\]
符号 "$x_1$" 应解释为 $\varphi(1,1)$. 这确实定义了 $G$ 在 $C(\varphi)$ 上的一个作用, 因为
\[
 \rho\sigma x_\tau=x_{\rho\sigma\tau}-x_{\rho\sigma}+\varphi(\rho\sigma,\tau),
\]
而
\[
 \begin{aligned}
 \rho(\sigma x_\tau)&=\rho(x_{\sigma\tau}-x_\sigma+\varphi(\sigma,\tau))\\
 &=x_{\rho\sigma\tau}-x_\rho+\varphi(\rho,\sigma\tau)-(x_{\rho\sigma}-x_\rho+\varphi(\rho,\sigma))+\rho\varphi(\sigma,\tau).
 \end{aligned}
\]
它们一致是因为 $\varphi$ 满足上循环条件
\[
 \rho\varphi(\sigma,\tau)+\varphi(\rho,\sigma\tau)=\varphi(\rho\sigma,\tau)+\varphi(\rho,\sigma).
\]
注意 $\varphi$ 是 $1$-上链 $\sigma\mapsto x_\sigma$ 的上边界, 因此 $\gamma$ 在 $H^2(G,C(\varphi))$ 中映到零. 为此原因, $C(\varphi)$ 称为 $\gamma$ 的\emph{分裂模}.

我们先证明这些假设蕴含
\[
 H^1(H,C(\varphi))=0=H^2(H,C(\varphi))
\]
对所有 $G$ 的子群 $H$.

回忆我们有一个正合列
\[
 0\to I_G\to \Z[G]\to\Z\to 0,
\]
其中 $I_G$ 是以元素 $\sigma-1$ ($\sigma\in G$, $\sigma\neq 1$) 为基的自由交换群. 因为 $\Z[G]$ 是诱导的, 对所有 $r$ 有 $H^r(H,\Z[G])=0$, 于是
\[
 \begin{array}{r l}
 H^1(H,I_G) & \simeq H^0(H,\Z)\simeq \Z/(H:1)\Z,\\
 H^2(H,I_G) & \simeq H^1(H,\Z)=0.
 \end{array}
\]
定义 $\alpha$ 为加性映射 $C(\varphi)\to \Z[G]$, 使得
\[
 \begin{aligned}
 \alpha(c)&=0\quad\text{对所有 }c\in C,\\
 \alpha(x_\sigma)&=\sigma-1.
 \end{aligned}
\]
显然,
\[
 0\to C\to C(\varphi)\xrightarrow{\alpha} I_G\to 0
\]
是一个 $G$-模的正合列. 它的上同调列读为
\[
 0\to H^1(H,C(\varphi))\to H^1(H,I_G)\to H^2(H,C)\xrightarrow{0} H^2(H,C(\varphi))\to 0.
\]
两端的零来自 $H^1(H,C)=0$ 与 $H^2(H,I_G)=0$. 映射 $H^2(H,C)\to H^2(H,C(\varphi))$ 为零, 因为 $H^2(H,C)$ 由 $\Res(\gamma)$ 生成, 而它映到 $\gamma$ 在 $H^2(G,C(\varphi))$ 中像的限制, 即零. 因此 $H^1(H,I_G)\to H^2(H,C)$ 是满射, 从而是同构 (因为两个群有相同的阶). 它的核与余核, 即 $H^1(H,C(\varphi))$ 与 $H^2(H,C(\varphi))$, 因此都为零.

我们从定理\ref{thm:ii.3.10} 推出对所有 $r$ 有 $H^r(H,C(\varphi))=0$. 把两个短正合列拼接起来, 我们得到一个正合列
\begin{equation*}
 0\to C\to C(\varphi)\to \Z[G]\to\Z\to 0
\tag{24}
\end{equation*}
具有性质: 对所有 $r$ 有 $H^r(G,C(\varphi))=0=H^r(G,\Z[G])$. 因此双边缘映射是一个同构 (\ref{rem:ii.1.13})
\[
 H^r(G,\Z)\to H^{r+2}(G,C).
\]
\end{proof}

\begin{remark}\label{rem:ii.3.12}
如果 $M$ 是一个满足 $\Tor_1^\Z(M,C)=0$ 的 $G$-模, 例如如果 $M$ 或 $C$ 作为 $\Z$-模无挠, 那么可以把上面的四项序列与 $M$ 作张量积并得到同构
\[
 H_T^r(G,M)\to H_T^{r+2}(G,M\otimes C).
\]
\end{remark}

\begin{remark}\label{rem:ii.3.13}
映射 $H_T^r(G,\Z)\to H_T^{r+2}(G,C)$ 是与所选元素 $\gamma\in H^2(G,C)$ 的杯积.
\end{remark}

\begin{example}\label{ex:ii.3.14}
设 $K$ 为一个局部域. 我们将证明对任何有限 Galois 扩张 $L/K$ (Galois 群 $G$), $H^2(G,L^\times)$ 是一个阶为 $[L:K]$ 的循环群, 带一个典范生成元 $u_{L/K}$. 因为已知 $H^1(G,L^\times)=0$ (Hilbert 定理 90), Tate 定理表明与 $u_{L/K}$ 的杯积是一个同构
\[
 G^{\mathrm{ab}}=H_T^{-2}(G,\Z)\to H_T^0(G,L^\times)=K^\times/\Nm\, L^\times.
\]
逆同构 $K^\times/\Nm L^\times\xrightarrow{\simeq} G^{\mathrm{ab}}$ 是\emph{局部 Artin 映射}. 整体 Artin 映射可由类似论证得到.
\end{example}

\section{Profinite 群的上同调}\label{sec:II.4}\index{profinitequn@profinite 群(profinite group)}\index{galoisqun@Galois 群(Galois group)}

\subsubsection*{正向极限}

一个偏序集 $(I,\le)$ 称为\emph{有向的}, 如果对 $I$ 中任意两个元素 $i$ 与 $j$, 存在 $k$ 使得 $i,j\le k$. 假设对有向集 $(I,\le)$ 中每个元素 $i$ 我们有一个集合 $A_i$, 且对每条不等式 $i\le j$ 我们有一个映射 $\alpha_{ji}\colon A_i\to A_j$. 如果
\begin{enumerate}[(a)]
\item 对所有 $i\in I$ 有 $\alpha_{ii}=\id$;
\item 每当 $i\le j\le k$ 有 $\alpha_{kj}\circ\alpha_{ji}=\alpha_{ki}$,
\end{enumerate}
那么该族 $(A_i,\alpha_{ji})$ 称为一个\emph{正向系统}. 在 $A_i$ 的无交并 $\coprod A_i$ 上引入等价关系: $a_i\in A_i$ 与 $a_j\in A_j$ 相关当且仅当对某个 $k\ge i,j$ 有 $\alpha_{ki}(a_i)=\alpha_{kj}(a_j)$. 相应的商集称为 $A_i$ 的\emph{正向极限} (相对于 $\alpha_{ji}$):
\[
 A=\varinjlim A_i.
\]
对每个 $i$ 有一个典范映射
\[
 \alpha_i\colon A_i\to A,
\]
具有如下性质:
\begin{enumerate}[(a)]
\item 对 $j\ge i$ 有 $\alpha_i=\alpha_j\circ\alpha_{ji}$;
\item $\alpha_i(a_i)=\alpha_j(a_j)\iff$ 对某个 $k\ge i,j$ 有 $\alpha_{ki}(a_i)=\alpha_{kj}(a_j)$;
\item $A=\bigcup \alpha_i(A_i)$.
\end{enumerate}
系统 $(A,\alpha_i)$ 有如下泛性质: 设 $T$ 为集合, $(\beta_i)$ ($\beta_i\colon A_i\to T$) 为一族映射使得对 $i\le j$ 有 $\beta_i=\beta_j\circ\alpha_{ji}$; 则存在唯一映射 $\beta\colon A\to T$ 使得对所有 $i$ 有 $\beta_i=\beta\circ\alpha_i$.

如果 $A_i$ 是交换群且 $\alpha_{ij}$ 是同态, 那么 $A$ 有唯一的交换群结构使得 $\alpha_i$ 是同态.

\begin{lemma}\label{lem:ii.4.1}
对任何正合列的正向系统
\[
 A_i\to B_i\to C_i,
\]
序列
\[
 \varinjlim A_i\to \varinjlim B_i\to \varinjlim C_i
\]
仍正合. 因此正向极限的构造与上同调在复合形中的计算交换.
\end{lemma}

\begin{proof}
留作读者练习.
\end{proof}

\subsubsection*{Profinite 群}

设 $G$ 为一个 profinite 群. 这意味着 $G$ 是一个紧拓扑群\footnote{按照 Bourbaki, 我要求紧空间是 Hausdorff 的.}, 其开正规子群构成 $1$ 处的邻域的基本系. 注意每个开子群指数有限 (因为它的陪集覆盖 $G$). 例如, 带离散拓扑的有限群是 profinite 的, 且每个离散 profinite 群有限. 一个 Galois 群 $G=\Gal(L/K)$ 是一个 profinite 群 —— 开子群恰好是固定 $L$ 中 $K$ 的某个有限扩张的那些 —— 且每个 profinite 群都作为某个 Galois 群出现. 对 profinite 群, 我们用拓扑来修正上同调群的概念.

首先, 我们只考虑那些使得映射
\[
 G\times M\to M
\]
连续的 $G$-模, 其中 $M$ 赋予离散拓扑, 即每个子集都开的拓扑. 等价条件:
\begin{itemize}
\item $M=\bigcup M^H$, $H$ 跑遍 $G$ 的开子群;
\item $M$ 中任何元素的稳定子群在 $G$ 中开.
\end{itemize}
满足这些条件的模称为\emph{离散 $G$-模}.

离散 $G$-模构成一个有足够多内射对象的 Abel 范畴, 因此我们可以像前面一样通过取内射分解来定义上同调群 $H_{\mathrm{cts}}^r(G,M)$ ($r\ge 0$). 此外, 这些群可以用\emph{连续上链}计算: 对 $r\ge 0$, 设 $C_{\mathrm{cts}}^r(G,M)$ 为连续映射 $G^r\to M$ 的群, 并如前面一样定义 $d^r\colon C_{\mathrm{cts}}^r(G,M)\to C_{\mathrm{cts}}^{r+1}(G,M)$; 则
\[
 H_{\mathrm{cts}}^r(G,M)=\frac{Z_{\mathrm{cts}}^r(G,M)}{B_{\mathrm{cts}}^r(G,M)},
\]
其中 $Z_{\mathrm{cts}}^r(G,M)=\Ker(d^r)$, $B_{\mathrm{cts}}^r(G,M)=\Im(d^{r-1})$.

设 $\varphi\colon G^r\to M$ 为一个连续 $r$-上链. 则 $\varphi(G^r)$ 紧 (因为 $G^r$ 紧) 且离散 (因为 $M$ 离散), 因此有限, 从而包含于某个开正规子群 $H_0$ 的 $M^{H_0}$ 中. $\varphi(G^r)$ 中每点 $m$ 的逆像 $\varphi^{-1}(m)$ 开, 因此包含某个开正规子群 $H(m)$ 的 $H(m)^r$ 的一个平移. 令 $H_1=\bigcap_{m\in\varphi(G^r)}H(m)$.\footnote{这个群可能不够小, 不足以使 $\varphi$ 经过 $(G/H_1)^r$ 分解. 取而代之, 对 $G^r$ 的每个点, 取包含该点且完全含于 $\varphi$ 的一个纤维中的某个开正规子群的一个平移. 这样我们便得到 $G^r$ 的一个开覆盖, 由于 $G^r$ 紧, 它可以替换为一个有限子覆盖. 现在令 $H_1$ 为所有其平移在该子覆盖中出现的子群的交.} 这又是 $G$ 的一个开子群, 且 $\varphi$ 经过 $(G/H_1)^r$. 令 $H=H_0\cap H_1$. 则 $\varphi$ 由 $G/H$ 上一个在 $M^H$ 中取值的 $r$-上循环膨胀而来. 换言之, 映射
\[
 \varinjlim C^r(G/H,M^H)\to C_{\mathrm{cts}}^r(G,M)
\]
是满射. 不难看到它也是单射. 因为过渡到正向极限与 kernel 和 cokernel 的构造交换, 从而与上同调的构造交换, 我们得到如下命题:

\begin{proposition}\label{prop:ii.4.2}
映射 $\Inf\colon H^r(G/H,M^H)\to H_{\mathrm{cts}}^r(G,M)$ 把 $H_{\mathrm{cts}}^r(G,M)$ 实现为群 $H^r(G/H,M^H)$ 当 $H$ 跑遍 $G$ 的开正规子群时的正向极限:
\[
 \varinjlim H^r(G/H,M^H)=H_{\mathrm{cts}}^r(G,M).
\]
\end{proposition}

确切地说, 该陈述意味着 $H_{\mathrm{cts}}^r(G,M)$ 的每个元素都由一个群 $H^r(G/H,M^H)$ 中的元素膨胀而来, 且如果 $a\in H^r(G/H,M^H)$ 与 $a'\in H^r(G/H',M^{H'})$ 在 $H_{\mathrm{cts}}^r(G,M)$ 中映到同一元素, 那么对某个包含于 $H\cap H'$ 的开子群 $H''$ 它们在 $H^r(G/H'',M^{H''})$ 中映到同一元素.

\begin{corollary}\label{cor:ii.4.3}
对每个 profinite 群 $G$ 与离散 $G$-模 $M$, 对所有 $r>0$ 有 $H^r(G,M)$ 是一个挠群.
\end{corollary}

\begin{proof}
每个群 $H^r(G/H,M^H)$ ($r>0$) 是挠群 (见 \ref{cor:ii.1.31}).
\end{proof}

关于 $r\ge 0$ 的上同调群 $H^r(G,M)$ 的大部分理论对由连续上链定义的群仍成立. 例如, 如果 $H$ 是 $G$ 的闭子群, 则存在映射 $\Inf$, $\Res$ 与 $\Cor$ (最后一个要求 $H$ 指数有限), 并且存在杯积映射.

今后, 所有上同调群都将用连续上链定义 (并略去下标 cts). 实际上, 这意味着要么 $G$ 是无限 profinite 群, 此时取连续上链是有关系的; 要么 $G$ 有限, 此时则无关 (且群对所有 $r\in\Z$ 有定义).

\begin{proposition}\label{prop:ii.4.4}
设 $G$ 为 profinite 群, $M$ 为离散 $G$-模. 如果 $M=\varinjlim M_i$, 其中 $M_i\subset M$, 那么 $H^r(G,M)=\varinjlim H^r(G,M_i)$.
\end{proposition}

\begin{proof}
因为 $G$ 紧而 $M$ 离散, 任何 $r$-上链 $f\colon G^r\to M$ 的像有限. 因为 $M_i$ 形成 $M$ 的子模的一个有向系 (即给定 $M_i$ 与 $M_j$, 存在同时包含二者的 $M_k$) 且 $M=\bigcup M_i$, $M$ 的每个有限子集都包含于某个 $M_i$. 由此可知 $C^r(G,M)=\varinjlim C^r(G,M_i)$, 于是引理\ref{lem:ii.4.1} 表明
\[
 H^r(C^\bullet(G,M))=\varinjlim H^r(C^\bullet(G,M_i)).
\]
\end{proof}

\begin{aside}\label{as:ii.4.5}
(给专家的.) 设 $G$ 为 profinite 群, 令 $M(G)$ 为所有 $G$-模的范畴, $C(G)$ 为离散 $G$-模的范畴. 则 $C(G)$ 是 $M(G)$ 的全子范畴. 此外, 存在一个函子
\[
 M\mapsto M^*\stackrel{\mathrm{def}}{=}\bigcup_{H\text{ 在 }G\text{ 中开}} M^H\colon \mathsf{M}(G)\to \mathsf{C}(G).
\]
显然,
\[
 \operatorname{Hom}_G(M,N^*)=\operatorname{Hom}_G(M,N)
\]
对离散 $G$-模 $M$. 含入函子 $i\colon \mathsf{C}(G)\to \mathsf{M}(G)$ 正合, 但不保持内射性 —— 因此一般地 $H^r(G,M)\neq H^r(G,iM)$. 另一方面, $M\mapsto M^*$ 保持内射性, 但只是左正合的 —— 因此 $\mathsf{C}(M)$ 有足够多的内射对象. 再次有 $H^r(G,M)\neq H^r(G,M^*)$ (然而存在一个谱序列……).
\end{aside}

\begin{NOTES}
在 1930 年代中期, Hurewicz 证明了"非球面空间" $X$ 的同调群只依赖于该空间的基本群 $\pi$. 因此可以把空间的同调群 $H_r(X,\Z)$ 视为群 $\pi$ 的同调群 $H_r(\pi,\Z)$. 直到 1940 年代中期, Hopf, Eckmann, Eilenberg, MacLane, Freudenthal 及其他人才给出群 $G$ 的同调群与上同调群的纯代数定义. 这时人们发现 $H^1$ 与交叉同态模主交叉同态一致, $H^2$ 与"因子集"的同构类群一致, 后者早就被引入 (例如 I. Schur, Über die Darstellung der endlichen…, 1904; O. Schreier, Über die Erweiterungen von Gruppen, 1926; R. Brauer, Über Zusammenhänge…, 1926). 更多历史见 \cite{MacLane1978}.

我们对 Tate 定理的证明遵循 Tate 的原始证明 (\cite{Tate1952}). 那时还没有 Tate 群的公开说明, 因此 Tate 只对 $r\ge 0$ 证明他的定理, 但以一个神秘的承诺结尾: 要把结果推广到负的 $r$ 并重新得到 Artin 映射. Tate 上同调群的构造在 \cite{CartanEilenberg1956} 中首次发表, 遵循 Tate 的想法.

前三节内容的优秀来源是 \cite{Serre1962}, 第 3 部分. \cite{AtiyahWall1967} 中有同一内容的稍简版本. 另见 \cite{Iyanaga1975}, 第 I 章, 与 \cite{Weiss1969}. 对 profinite 群的上同调, 见 \cite{Serre1964} 与 \cite{Shatz1972}.

对上同调在类域论中起源的讨论, 见 \sx{857593}.
\end{NOTES}

\phantomsection
\section*{A 附录: 若干同调代数}\label{sec:II.A}\addcontentsline{toc}{section}{A 附录: 若干同调代数}\index{tongdaodaishu@同调代数(homological algebra)}\index{daochuhanshu@导出函子(derived functor)}

Ça me semble extrêmement plaisant de ficher comme ça beaucoup de choses, pas drôles quand on les prend séparément, sous le grand chapeau des foncteurs dérivés.

Grothendieck, 给 Serre 的信, 18.2.1955.\footnote{依我看, 把形形色色单独拿出来看并无多少趣味的东西, 统统塞进导出函子这顶大帽子底下, 是件极为愉快的事.}

% 附录中定理按 A.x 编号
\renewcommand{\thetheorem}{A.\arabic{theorem}}
\setcounter{theorem}{0}

\subsubsection*{若干正合列}\footnote{基于 F. Lemmermeyer, \emph{The Snake Lemma}.}

\begin{lemma}[扩张蛇形引理]\label{lem:ii.A.1}
下述交换图导出长正合列 (蛇形):
\[
\begin{tikzcd}[column sep=small, row sep=small]
 0 \ar[r] & \Ker f \ar[r] & \Ker a \ar[r] \ar[d, hook] & \Ker b \ar[r] \ar[d, hook] & \Ker c \ar[d, hook] \ar[dddll, bend left=55] & & \\
 & & A \ar[r, "f"] \ar[d, "a"'] & B \ar[r] \ar[d, "b"'] & C \ar[r] \ar[d, "c"'] & 0 & \\
 & 0 \ar[r] & A' \ar[r] \ar[d] & B' \ar[r, "g'"] \ar[d] & C' \ar[d] & & \\
 & & \Coker a \ar[r] & \Coker b \ar[r] & \Coker c \ar[r] & \Coker g' \ar[r] & 0
\end{tikzcd}
\]
即长正合列 (蛇形)
\[
 0\to \Ker f\to \Ker a\to \Ker b\to \Ker c\to \Coker a\to \Coker b\to \Coker c\to \Coker g'\to 0.
\]
\end{lemma}

\begin{proof}
除首尾两项外这是标准的. 一个小图表追踪表明 $\Ker f\subset \Ker a$, 由此得到在 $\Ker a$ 处的正合性. 在 $\Coker g'$ 处的正合性的证明同样直接.
\end{proof}

\begin{lemma}[核--余核引理]\label{lem:ii.A.2}
每对同态
\[
 A\xrightarrow{f} B\xrightarrow{g} C
\]
的交换群产生一个正合列
\[
 0\to \Ker f\to \Ker(g\circ f)\xrightarrow{f} \Ker g\to \Coker f\to \Coker(g\circ f)\to \Coker g\to 0.
\]
\end{lemma}

\begin{proof}
把扩张蛇形引理应用于
\[
 \begin{tikzcd}
 & A \ar[r, "f"] \ar[d, "g\circ f"'] & B \ar[r] \ar[d, "g"'] & \Coker(f) \ar[r] \ar[d] & 0 \\
 0 \ar[r] & C \ar[r, "\mathrm{id}"'] & C \ar[r] & 0 &
 \end{tikzcd}
\]
\end{proof}

\subsubsection*{范畴论的语言}

一个范畴 $\mathsf{C}$ 由一个非空的对象类 $\operatorname{ob}(\mathsf{C})$, 对每对对象 $A,B$ 一个集合 $\Hom(A,B)$ (称为从 $A$ 到 $B$ 的态射集合), 以及对每个三元组对象 $A,B,C$ 一个映射
\[
 (\alpha,\beta)\mapsto \beta\circ\alpha\colon \Hom(A,B)\times \Hom(B,C)\to \Hom(A,C)
\]
组成, 满足如下条件:
\begin{enumerate}[(a)]
\item 态射的复合结合;
\item 对每个对象 $A$, $\Hom(A,A)$ 有一个元素 $\id_A$, 它是复合的左, 右单位元.
\end{enumerate}
要理解集合 $\Hom(A,B)$ 互不相交, 因此一个态射决定它的源与目标.

一个\emph{协变函子} $F\colon \mathsf{C}\to \mathsf{D}$ 是一个"映射", 它把 $\mathsf{C}$ 的每个对象 $A$ 附到一个 $\mathsf{D}$ 的对象 $F(A)$, 把每个态射 $\alpha\colon A\to B$ 附到一个态射 $F(\alpha)\colon F(A)\to F(B)$, 并保持态射的复合: 若 $\alpha\colon A\to B$, $\beta\colon B\to C$, 则 $F(\alpha\circ\beta)=F(\alpha)\circ F(\beta)$; 此外 $F(\id_A)=\id_{F(A)}$.

一个函子 $F\colon \mathsf{C}\to \mathsf{D}$ 是函子 $G\colon \mathsf{D}\to \mathsf{C}$ 的\emph{左伴随}, 如果存在自然同构
\[
 \operatorname{Hom}_{\mathsf{D}}(F(A),B)\simeq \operatorname{Hom}_{\mathsf{C}}(A,G(B)).
\]
如果集合 $\Hom(A,B)$ 被赋予交换群的结构使得复合映射是双加性的, 且 $\mathsf{C}$ 中每个有限对象族都有直和, 那么 $\mathsf{C}$ (连同这些结构) 称为一个\emph{加性范畴}. 说 $A$ 与 $B$ 容许一个直和意味着在 $\mathsf{C}$ 中存在一个对象 $A\oplus B$ 与映射 $i_A\colon A\to A\oplus B$, $i_B\colon B\to A\oplus B$, $p_A\colon A\oplus B\to A$, $p_B\colon A\oplus B\to B$ 使得:
\[
 \begin{aligned}
 p_A\circ i_A&=\id_A & p_B\circ i_B&=\id_B\\
 p_A\circ i_B&=0 & p_B\circ i_A&=0
 \end{aligned}
\]
\[
 i_A p_A+i_B p_B=1_{A\oplus B}.
\]
设 $\mathsf{C}$ 为一个加性范畴. 一个序列
\[
 0\to A\to B\xrightarrow{\alpha} C
\]
正合, 如果对所有对象 $T$ 交换群序列
\[
 0\to \Hom(T,A)\to \Hom(T,B)\to \Hom(T,C)
\]
正合. 一个序列
\[
 A\xrightarrow{\beta} B\to C\to 0
\]
正合, 如果对所有对象 $T$ 交换群序列
\[
 0\to \Hom(C,T)\to \Hom(B,T)\to \Hom(A,T)
\]
正合. 当第一个序列正合时, $A$ 称为 $\alpha$ 的\emph{核}; 当第二个序列正合时, $C$ 称为 $\beta$ 的\emph{余核}.

一个态射 $A\to B$ 称为\emph{单的} (或单态射), 如果 $0\to A\to B$ 正合. 类似地, 它称为\emph{满的} (或满态射), 如果 $A\to B\to 0$ 正合.

设 $\mathsf{C}$ 为一个加性范畴, 其中每个态射都有核与余核. 设 $\alpha\colon A\to B$ 为一个态射. $\alpha$ 的余核 $B\to C$ 的核称为 $\alpha$ 的\emph{像}, 而 $\alpha$ 的核的余核称为 $\alpha$ 的\emph{余像}. 存在一个从 $\alpha$ 的余像到像的典范映射, 如果它总是同构, 那么 $\mathsf{C}$ 称为一个\emph{Abel 范畴}.\footnote{如果你对这一段感到困惑, sx2263197 或许会有所帮助.}

加性范畴之间的函子将假定是加性的, 即映射 $\Hom(A,B)\to \Hom(F(A),F(B))$ 是交换群同态. 这样的函子称为\emph{正合的}, 如果它把正合列映到正合列.

例如, 对每个环 $R$, $R$-模的范畴是一个 Abel 范畴; 且对每个拓扑空间 $X$, $X$ 上交换群层范畴是一个 Abel 范畴.

在本节余下部分, $\mathsf{C}$ 将是一个 Abel 范畴. 就本课程而言, 读者取 $\mathsf{C}$ 为环上的模范畴 (例如群环 $\Z[G]$ 上的模范畴) 不会损失什么.

\subsubsection*{内射对象}

设 $\mathsf{C}$ 为一个 Abel 范畴. $\mathsf{C}$ 的一个对象 $I$ 是\emph{内射的}, 如果 $\operatorname{Hom}(\cdot,I)$ 是一个正合函子, 即如果
\[
 0\to A\to B\to C\to 0
\]
在 $\mathsf{C}$ 中正合蕴含
\[
 0\to \Hom(C,I)\to \Hom(B,I)\to \Hom(A,I)\to 0
\]
正合. 最后一个序列除在 $\Hom(A,I)$ 处外自动正合, 因此说 $I$ 内射意味着: 每当 $A$ 是 $B$ 的子对象, 每个同态 $A\to I$ 都扩张到 $B$.

一个 Abel 范畴 $\mathsf{C}$ 称为\emph{有足够多的内射对象}, 如果每个对象都容许一个嵌入内射对象的单同态.

\begin{proposition}\label{prop:ii.A.3}
主理想整环 $R$ 上的模 $M$ 是内射的当且仅当它是可除的, 即映射 $x\mapsto rx$ 对所有 $r\in R$, $r\neq 0$ 是满射.
\end{proposition}

\begin{proof}
内射 $\Longrightarrow$ 可除: 设 $m\in M$, $r$ 为 $R$ 的非零元素. 映射 $x\mapsto rx\colon R\to R$ 是单的, 而 $x\mapsto xm\colon R\to M$ 到 $R$ 的每个扩张都将 $1$ 送到满足 $rm'=m$ 的元素 $m'$.

可除 $\Longrightarrow$ 内射: 假设 $M$ 可除, 考虑同态 $\alpha\colon A\to M$, 其中 $A$ 是 $B$ 的子模. 把 Zorn 引理用于满足如下条件的对 $(A',\alpha')$ 的集合: $A'$ 是包含 $A$ 的 $B$ 的子模且 $\alpha'$ 把 $\alpha$ 扩张到 $A'$, 我们得到一个极大的这样的对 $(A_1,\alpha_1)$. 如果 $A_1\neq B$, 则存在 $b\in B\setminus A_1$, 我们定义 $I=\{r\in R\mid rb\in A_1\}$. 因为 $M$ 可除, 映射 $r\mapsto \alpha_1(rb)\colon I\to M$ 扩张到 $R$\footnote{目前这一点留作读者的练习.}, 但这蕴含 $\alpha_1$ 扩张到 $A_1+Rb$, 与 $(A_1,\alpha_1)$ 的极大性矛盾.
\end{proof}

因此, 当 $R$ 为主理想整环时, 内射模的每个商都是内射的, 例如 $K/R$, 其中 $K$ 是 $R$ 的分式域. 特别地, $\Q/\Z$ 是一个内射交换群.

\begin{proposition}\label{prop:ii.A.4}
主理想整环 $R$ 上的模范畴有足够多的内射对象.
\end{proposition}

\begin{proof}
把 $M$ 写成自由 $R$-模 $F$ 的商 $M\approx F/N$, 并考虑正合交换图表,
\[
\begin{tikzcd}[column sep=small, row sep=small]
	& & 0 & 0 \\
	0 & N & F & M & 0 \\
	0 & N & {F\otimes_R K} & {(F\otimes_R K)/N} & 0 \\
	& 0 & 0 & 0
	\arrow[from=1-3, to=2-3]
	\arrow[from=1-4, to=2-4]
	\arrow[from=2-1, to=2-2]
	\arrow[from=2-2, to=2-3]
	\arrow[from=2-3, to=2-4]
	\arrow[from=2-4, to=2-5]
	\arrow[equal, from=2-2, to=3-2]
	\arrow[from=2-3, to=3-3]
	\arrow[from=2-4, to=3-4]
	\arrow[from=3-1, to=3-2]
	\arrow[from=3-2, to=3-3]
	\arrow[from=3-3, to=3-4]
	\arrow[from=3-4, to=3-5]
	\arrow[from=3-2, to=4-2]
	\arrow[from=3-3, to=4-3]
	\arrow[from=3-4, to=4-4]
\end{tikzcd}
\]
因为 $F\to F\otimes_R K$ 是单射且 $F\otimes_R K$ 是内射对象, 映射 $M\to (F\otimes_R K)/N$ 是 $M$ 嵌入一个内射 $R$-模的单同态.
\end{proof}

对交换群, 该命题也可如下证明: 对交换群 $M$, 令 $M^\vee=\Hom(M,\Q/\Z)$; 取一个自由交换群 $F$ 满射到 $M^\vee$, $F\to M^\vee$; 则 $M\hookrightarrow M^{\vee\vee}\hookrightarrow F^\vee$, 且 $F^\vee$ 内射.

\begin{proposition}\label{prop:ii.A.5}
一个容许正合左伴随的函子保持内射对象.
\end{proposition}

\begin{proof}
设 $F'$ 为函子 $F\colon \mathsf{C}\to \mathsf{D}$ 的正合左伴随. 对 $\mathsf{C}$ 中每个内射对象 $I$, 函子 $\operatorname{Hom}_{\mathsf{D}}(\cdot,F(I))$ 同构于函子 $\operatorname{Hom}_{\mathsf{C}}(F'(\cdot),I)$, 后者正合, 因为它是两个正合函子 $F'$ 与 $\operatorname{Hom}_{\mathsf{C}}(\cdot,I)$ 的复合.
\end{proof}

\subsubsection*{右导出函子}

设 $\mathsf{C}$ 为一个有足够多内射对象的 Abel 范畴, $F\colon \mathsf{C}\to \mathsf{D}$ 为从一个 Abel 范畴到另一个 Abel 范畴的左正合函子. 于是 $\mathsf{C}$ 中的一个短正合列
\[
 0\to M'\to M\to M''\to 0
\]
产生 $\mathsf{D}$ 中的一个正合列
\[
 0\to F(M')\to F(M)\to F(M'').
\]
导出函子理论提供把这个最后的序列自然扩张为一个长正合列.

设 $M$ 为 $\mathsf{C}$ 的一个对象. $M$ 的一个\emph{分解}是一个长正合列
\[
 0\to M\to I^0\xrightarrow{d^0} I^1\to\cdots\to I^r\xrightarrow{d^r} I^{r+1}\to\cdots.
\]
如果 $I^r$ 是 $\mathsf{C}$ 的内射对象, 则称为\emph{内射分解}. 我们有时把这个复形记为 $M\to I^\bullet$.

\begin{lemma}\label{lem:ii.A.6}
一个对象 $M$ 的内射分解 $M\to I^\bullet$ 存在; 且如果 $M\to J^\bullet$ 是第二个内射分解, 则存在从 $M\to I^\bullet$ 到 $M\to J^\bullet$ 的一个态射, 即存在一个交换图表,
\[
 \begin{tikzcd}
 0 \ar[r] & M \ar[r] \ar[d, no head, double] & I^0 \ar[r] \ar[d] & I^1 \ar[r] \ar[d] & \cdots \\
 0 \ar[r] & M \ar[r] & J^0 \ar[r] & J^1 \ar[r] & \cdots
 \end{tikzcd}
\]
\end{lemma}

\begin{proof}
由假设, 存在一个单同态
\[
 0\to M\to I^0
\]
其中 $I^0$ 内射. 令 $B^1$ 为该映射的余核. 再次, 存在一个单同态
\[
 0\to B^1\to I^1
\]
其中 $I^1$ 内射. 现在
\[
 0\to M\to I^0\to I^1
\]
正合. 令 $B^2=\Coker(B^1\to I^1)$, 并如此继续.

类似地, 可以用内射对象的定义一步步构造分解的态射.
\end{proof}

对一个对象 $M$ 应用左正合函子 $F$ 到它的内射分解 $M\to I^\bullet$, 我们得到一个复形
\[
 F(I^\bullet)\colon F(I^0)\to F(I^1)\to\cdots\to F(I^r)\xrightarrow{F(d^r)} F(I^{r+1})\to\cdots
\]
它未必仍正合, 因此上同调群
\[
 H^r(F(I^\bullet))\stackrel{\mathrm{def}}{=}\frac{\Ker(F(d^r))}{\Im(F(d^{r-1}))}
\]
可能非零.

\begin{remark}\label{rem:ii.A.7}
因为 $F$ 左正合, 序列
\[
 0\to F(M)\to F(I^0)\xrightarrow{d^0} F(I^1)
\]
正合. 因此
\[
 H^0(F(I^\bullet))\stackrel{\mathrm{def}}{=}\Ker(d^0)=F(M).
\]
\end{remark}

\begin{proposition}\label{prop:ii.A.8}
设 $M\to I^\bullet$ 与 $N\to J^\bullet$ 为 $\mathsf{C}$ 中对象 $M$ 与 $N$ 的内射分解. 每个态射 $\alpha\colon M\to N$ 都扩张为一个复形的映射
\[
 \begin{tikzcd}
 M \ar[r] \ar[d, "\alpha"'] & I^\bullet \ar[d, "\alpha^\bullet"] \\
 N \ar[r] & J^\bullet.
 \end{tikzcd}
\]
对每个左正合函子 $F$, 态射
\[
 H^r(F(\alpha^\bullet))\colon H^r(F(I^\bullet))\to H^r(F(J^\bullet))
\]
与 $\alpha^\bullet$ 的选取无关.
\end{proposition}

我们在下面讨论证明.

该命题用于恒等态射 $M\to M$ 表明对象 $H^r(F(I^\bullet))$ 的定义到一个良定的同构是良定的: 给定 $M$ 的两个内射分解 $M\to I^\bullet$ 与 $M\to J^\bullet$, 存在扩张恒等映射的态射 $\alpha^\bullet\colon I^\bullet\to J^\bullet$, 它定义的映射 $H^r(F(I^\bullet))\to H^r(F(J^\bullet))$ 是同构且与 $\alpha^\bullet$ 的选取无关. 现在对每个 $\mathsf{C}$ 的对象 $M$, 我们选取一个内射分解 $M\to I^\bullet$, 并定义
\[
 (R^rF)(M)=H^r(F(I^\bullet)).
\]
一个态射 $\alpha\colon M\to N$ 产生一个良定的态射 $(R^rF)(M)\to (R^rF)(N)$, 这些映射使 $R^rF$ 成为一个函子. 它们称为 $F$ 的\emph{右导出函子}.

下面两个引理证明了比该命题稍精确一些的结论.

\begin{lemma}\label{lem:ii.A.9}
设 $M\to I^\bullet$ 与 $N\to J^\bullet$ 为 $\mathsf{C}$ 中对象 $M$ 与 $N$ 的分解. 如果 $N\to J^\bullet$ 是一个内射分解, 那么每个态射 $\alpha\colon M\to N$ 都扩张为一个复形的态射
\[
 \begin{tikzcd}
 M \ar[r] \ar[d] & I^\bullet \ar[d] \\
 N \ar[r] & J^\bullet
 \end{tikzcd}
\]
\end{lemma}

\begin{proof}
\cite{BucurDeleanu1968}, 7.5.
\end{proof}

两个复形态射 $\alpha^\bullet,\beta^\bullet\colon I^\bullet\to J^\bullet$ 称为\emph{同伦的}, 如果存在一族态射 $k^r\colon I^r\to J^{r-1}$ (一个同伦) 使得对所有 $r$,
\[
 \alpha^r-\beta^r=d^{r-1}\circ k^r+k^{r+1}\circ d^r.
\]
注意对每个 $x\in Z^r(I^\bullet)\stackrel{\mathrm{def}}{=}\Ker(d^r)$,
\[
 \alpha^r(x)-\beta^r(x)=d^{r-1}(k^r(x))\in \Im(d^{r-1})\stackrel{\mathrm{def}}{=}B^r(J^\bullet).
\]
因此 $\alpha^r(x)$ 与 $\beta^r(x)$ 在 $H^r(J^\bullet)$ 中有相同的像, 故同伦的态射在上同调上定义相同的态射.

\begin{lemma}\label{lem:ii.A.10}
设 $M\to I^\bullet$ 为 $M$ 的一个分解, $N\to J^\bullet$ 为 $N$ 的一个内射分解. 把态射 $M\to N$ 扩张到 $I^\bullet\to J^\bullet$ 的任何两个扩张 $\alpha^\bullet$ 与 $\beta^\bullet$ 都同伦.
\end{lemma}

\noindent 证明. 见 \cite{BucurDeleanu1968}, 7.5.

这推出命题 \ref{prop:ii.A.8} 的第二个陈述, 因为族 $(F(k^r))$ 是从 $F(\alpha^\bullet)$ 到 $F(\beta^\bullet)$ 的一个同伦.

\begin{proposition}\label{prop:ii.A.11}
一个短正合列
\[
 0\to A\to B\to C\to 0
\]
在 $\mathsf{C}$ 中产生一个长正合列
\begin{align*}
 0\to F(A)\to F(B)\to F(C)\to R^1F(A)\to\cdots\\
 \cdots\to R^rF(A)\to R^rF(B)\to R^rF(C)\to\cdots
 \end{align*}
且长正合列与短正合列的对应是函子性的.
\end{proposition}

第二个条件意味着一个交换图表
\[
 \begin{tikzcd}
 0 \ar[r] & A \ar[r] \ar[d] & B \ar[r] \ar[d] & C \ar[r] \ar[d] & 0 \\
 0 \ar[r] & A' \ar[r] & B' \ar[r] & C' \ar[r] & 0
 \end{tikzcd}
\]
产生一个交换图表
\[
 \begin{tikzcd}
 \cdots \ar[r] & R^{r-1}F(C) \ar[r] \ar[d] & R^rF(A) \ar[r] \ar[d] & R^rF(B) \ar[r] \ar[d] & R^rF(C) \ar[r] \ar[d] & \cdots \\
 \cdots \ar[r] & R^{r-1}F(C') \ar[r] & R^rF(A') \ar[r] & R^rF(B') \ar[r] & R^rF(C') \ar[r] & \cdots
 \end{tikzcd}
\]

对该命题的证明, 见 \cite{BucurDeleanu1968}, 7.6.

\begin{remark}\label{rem:ii.A.12}
$F$ 的右导出函子由如下三个性质唯一确定 (到一个唯一的函子同构):
\begin{enumerate}[(a)]
\item $R^0F=F$;
\item 当 $I$ 内射且 $r>0$ 时 $R^rF(I)=0$;
\item (\ref{prop:ii.A.11}) 中的性质.
\end{enumerate}
\end{remark}

\subsubsection*{变体}

通过反转某些箭头的方向, 可以得到上面某些定义的变体, 例如投射对象, 左导出函子, 等等.

\subsubsection*{Ext 群}

设 $\mathsf{C}$ 为一个 Abel 范畴.

设 $A\in\mathsf{C}$. 如果 $\mathsf{C}$ 有足够多的内射对象, 那么我们可以定义左正合函子 $\Hom(A,\cdot)$ 的右导出函子. 记第 $r$ 个右导出函子为 $\operatorname{Ext}^r(A,\cdot)$. 要计算 $\operatorname{Ext}^r(A,B)$, 我们取 $B$ 的一个内射分解 $B\to I^\bullet$, 并令
\[
 \operatorname{Ext}^r(A,B)=H^r(\Hom(A,I^\bullet)).
\]
设 $B\in\mathsf{C}$. 如果 $\mathsf{C}$ 有足够多的投射对象, 那么我们可以定义左正合反变函子 $\Hom(\cdot,B)$ 的右导出函子. 记第 $r$ 个右导出函子为 $\operatorname{Ext}^r(\cdot,B)$. 要计算 $\operatorname{Ext}^r(A,B)$, 我们取 $A$ 的一个投射分解 $P_\bullet\to A$, 并令
\[
 \operatorname{Ext}^r(A,B)=H^r(\Hom(P_\bullet,B)).
\]

\begin{proposition}\label{prop:ii.A.13}
如果 $\mathsf{C}$ 有足够多的内射对象与足够多的投射对象, 那么 $\operatorname{Ext}^r(A,B)$ 的两个定义一致.
\end{proposition}

\begin{proof}
我们用投射定义 $\operatorname{Ext}^r$, 并证明它们具有刻画 $\Hom(A,\cdot)$ 的右导出函子的性质.

首先, 显然 $\operatorname{Ext}^0(A,B)=\Hom(A,B)$.

说 $I$ 内射意味着 $\Hom(\cdot,I)$ 正合. 因此 $\Hom(P_\bullet,I)$ 正合, 于是
\[
 \operatorname{Ext}^r(A,I)\stackrel{\mathrm{def}}{=}H^r(\Hom(P_\bullet,I))=0.
\]
最后, 如果
\[
 0\to B'\to B\to B''\to 0
\]
正合, 那么因为 $P_\bullet\to A$ 是由投射对象给出的 $A$ 的分解, 复形序列
\[
 0\to \Hom(P_\bullet,B')\to \Hom(P_\bullet,B)\to \Hom(P_\bullet,B'')\to 0
\]
正合. 用一个标准程序, 我们从中得到一个长正合列
\[
 \cdots\to H^r(\Hom(P_\bullet,B'))\to H^r(\Hom(P_\bullet,B))\to H^r(\Hom(P_\bullet,B''))\to\cdots.
\]
\end{proof}

\begin{example}\label{ex:ii.A.14}
设 $G$ 为一个群. 则 $\Mod_G$ 既有足够多的内射对象也有足够多的投射对象. 对每个 $G$-模 $M$, $\Hom_G(\Z,M)=M^G$, 因此函子 $\Hom(\Z,\cdot)$ 与 $H^0(G,\cdot)$ 一致. 于是它们的右导出函子也一致:
\[
 \operatorname{Ext}_G^r(\Z,M)\simeq H^r(G,M).
\]
最后一个命题允许我们通过取 $\Z$ 的投射分解 $P_\bullet\to\Z$ 并令
\begin{equation*}
 H^r(G,M)=H^r(\Hom_G(P_\bullet,M))
\tag{25}
\end{equation*}
来计算这些群.
\end{example}

\begin{remark}\label{rem:ii.A.15}
采用公式 (25) 作为 $H^r(G,M)$ 的定义会稍微缩短本章的叙述. 这是 Atiyah and Wall (1967) 采取的方法, 但它不是自然的定义.
\end{remark}

\subsubsection*{参考文献}

对导出函子的一般概念, 见 \cite{BucurDeleanu1968} 的第 7 章.

% 恢复定理计数器为 节.序号 格式, 供后续章节使用
\renewcommand{\thetheorem}{\thesection.\arabic{theorem}}
   % II  群的上同调
% ============ 第三章 III: 局部类域论: 上同调 ============
% 源: J.S. Milne, Class Field Theory (v4.03), Chapter III, pp. 72--106
% 翻译自 OCR 草稿 J.S. Milne - Class Field Theory.pdf_by_PaddleOCR-VL-1.6.md (行 3787--4584)
% 编号约定见 main.tex: 定理类环境按"节.序号"编号, 全部共享 theorem 计数器.
%   本章各节的编号条目均为定理类环境或叙述性条目, 序号连续, 无需 \setcounter 干预;
%   唯一例外是 §3 中的 3.3, 它是不带环境名的编号条目, 故在其后 \setcounter{theorem}{3}.
%   源书中的公式编号 (26)--(34) 用 \tag 如实保留, 使正文中的交叉引用得以落实.

\chapter{局部类域论: 上同调}\label{ch:III}\index{jubuleiyulun@局部类域论(local class field theory)}\index{shangtongdiao@上同调(cohomology)}

在本章中, 我们展开局部类域论的上同调进路. 我们首先证明局部 Artin 映射的存在性 (定理\ref{thm:iii.3.4}). 结合第 \ref{ch:I} 章 \S\S 1--3 的定理, 这就足以完成局部类域论主定理的证明. 对于那些跳过了 Lubin-Tate 理论的读者, 我们随后给出一个不依赖第 \ref{ch:I} 章 (\S 1 除外) 的主定理证明.

贯穿本章, ``局部域'' 指 ``非阿基米德局部域''. 一如既往, $K^{\mathrm{al}}$ 表示 $K$ 的一个代数闭包 (在 $K$ 的特征为 $p\neq 0$ 时指可分代数闭包), 而 ``$K$ 的扩张'' 指 ``$K^{\mathrm{al}}$ 中包含 $K$ 的子域''. 所有上同调群都用连续上链来计算 (见 II, \S 4). 对域的 Galois 扩张 $L/K$ (可能无限), 记
\[
 H^2(L/K)=H^2(\Gal(L/K),L^\times).
\]
眼下 $H^2(L/K)$ 只是一个上同调群, 但在下一章我们将会看到它作为 $L/K$ 的相对 Brauer 群有明确的解释.

% ==================================================================
\section{非分歧扩张的上同调}\label{sec:III.1}\index{feifenqikuozhang@非分歧扩张(unramified extension)}

\subsubsection*{单位群的上同调}

设 $K$ 为局部域.

\begin{proposition}\label{prop:iii.1.1}
设 $L/K$ 为 Galois 群为 $G$ 的有限非分歧扩张, 并设 $U_L$ 为 $L$ 中的单位群. 则
\[
 H_T^r(G,U_L)=0,\quad\text{所有 }r.
\]
\end{proposition}

\begin{proof}
若 $\pi$ 为 $L$ 的一个素元, 那么 $L^\times$ 的每个元素都可唯一写成 $\alpha=u\pi^m$ 的形式, 其中 $u\in U_L$, $m\in\Z$. 于是
\begin{equation}
 L^\times=U_L\cdot\pi^{\Z}\simeq U_L\times\Z.\tag{26}
\end{equation}
因为 $L$ 在 $K$ 上非分歧, 我们可以取 $\pi\in K$. 于是对 $\tau\in\Gal(L/K)$ 有 $\tau\alpha=\tau(u\pi^m)=(\tau u)\pi^m$, 因此当我们让 $G$ 平凡地作用在 $\pi^{\Z}\simeq\Z$ 上时, (26) 成为 $G$-模的一个分解. 所以 $H^r(G,U_L)$ 是 $H^r(G,L^\times)$ 的一个直和因子 (见 II, \ref{prop:ii.1.25}). 由于依 Hilbert 定理 90 (II, \ref{prop:ii.1.22}) 有 $H^1(G,L^\times)=0$, 这表明 $H^1(G,U_L)=0$. 因为 $G$ 是循环群, 为完成定理的证明, 只需 (依 II, \ref{prop:ii.3.4}) 证明 $H_T^0(G,U_L)=0$. 这一点由下一个命题完成.
\end{proof}

\begin{proposition}\label{prop:iii.1.2}
设 $L/K$ 为有限非分歧扩张. 则范映射 $\Nm_{L/K}\colon U_L\to U_K$ 是满射.
\end{proposition}

我们先需要一些引理. 设 $l$ 与 $k$ 为 $L$ 与 $K$ 的剩余类域. 因为 $L/K$ 非分歧, $G$ 在 $\mathcal{O}_L$ 上的作用定义了一个同构 $G\simeq\Gal(l/k)$.

\begin{lemma}\label{lem:iii.1.3}
对 $m>0$, 令 $U_L^{(m)}=1+\mathfrak{m}_L^m$. 则作为 $G$-模有
\[
 U_L/U_L^{(1)}\xrightarrow{\simeq} l^\times
\]
\[
 U_L^{(m)}/U_L^{(m+1)}\xrightarrow{\simeq} l .
\]
\end{lemma}

\begin{proof}
设 $\pi$ 为 $K$ 的一个素元. 它在 $L$ 中仍是素元, 且
\[
 U_L^{(m)}=\{1+a\pi^m\mid a\in\mathcal{O}_L\}.
\]
映射
\[
 u\mapsto u\mod\mathfrak{m}_L\colon U_L\to l^\times
\]
\[
 1+a\pi^m\mapsto a\mod\mathfrak{m}_L\colon U_L^{(m)}\to l
\]
诱导出所需的同构.
\end{proof}

\begin{lemma}\label{lem:iii.1.4}
对所有 $r$ 有 $H_T^r(G,l^\times)=0$. 特别地, 范映射 $l^\times\to k^\times$ 是满射.
\end{lemma}

\begin{proof}
依 Hilbert 定理 90 (II, \ref{prop:ii.1.22}), $H^1(G,l^\times)=0$, 又因为 $l^\times$ 有限, 其 Herbrand 商 $h(l^\times)=1$ (见 II, \ref{prop:ii.3.8}). 因此 $H^2(G,l^\times)=0$, 依 (II, \ref{prop:ii.3.4}) 这推出所有的群都是零.
\end{proof}

\begin{lemma}\label{lem:iii.1.5}
群 $H_T^r(G,l)=0$ 对所有 $r$ 成立. 特别地, 迹映射 $l\to k$ 是满射.
\end{lemma}

\begin{proof}
在 (II, \ref{prop:ii.1.24}) 中对 $r>0$ 证明了这一点, 依 (II, \ref{prop:ii.3.4}) 这蕴含它对所有 $r$ 成立.
\end{proof}

\begin{proof}[命题\ref{prop:iii.1.2} 的证明]
有交换图表 ($m>0$)
\[
\begin{tikzcd}
	{U_L} & {l^\times} \\
	{U_K} & {k^\times}
	\arrow[from=1-1, to=1-2]
	\arrow["\Nm"', from=1-1, to=2-1]
	\arrow["\Nm", from=1-2, to=2-2]
	\arrow[from=2-1, to=2-2]
\end{tikzcd}
\qquad
\begin{tikzcd}
	{U_L^{(m)}} & l \\
	{U_K^{(m)}} & k
	\arrow[from=1-1, to=1-2]
	\arrow["\Nm"', from=1-1, to=2-1]
	\arrow["\Tr", from=1-2, to=2-2]
	\arrow[from=2-1, to=2-2]
\end{tikzcd}
\]
取 $u\in U_K$. 因为范映射 $l^\times\to k^\times$ 是满射, 存在 $v_0\in U_L$ 使得 $\Nm(v_0)$ 与 $u$ 在 $k^\times$ 中有相同的像, 即 $u/\Nm(v_0)\in U_K^{(1)}$. 因为迹映射 $l\to k$ 是满射, 存在 $v_1\in U_L^{(1)}$ 使得 $\Nm(v_1)\equiv u/\Nm(v_0)\mod U_K^{(2)}$. 如此继续下去, 我们得到一列元素 $v_0,v_1,v_2,v_3,\ldots$, $v_i\in U_L^{(i)}$, 使得 $u/\Nm(v_0\cdots v_i)\in U_K^{(i+1)}$. 令 $v=\lim_{m\to\infty}\prod_{j=0}^{m}v_j$. 则 $u/\Nm(v)\in\bigcap U_K^{(i)}=\{1\}$.
\end{proof}

\begin{corollary}\label{cor:iii.1.6}
设 $L/K$ 为 Galois 群为 $G$ 的无限非分歧扩张. 则对 $r>0$ 有 $H^r(G,U_L)=0$ (连续上链).
\end{corollary}

\begin{proof}
域 $L$ 是 $K$ 的有限扩张 $K'$ 的并, 于是 (见 II, \ref{prop:ii.4.2})
\[
 H^r(\Gal(L/K),U_L)\simeq\varinjlim_{K'}H^r(\Gal(K'/K),U_{K'})=0.
\]
\end{proof}

\subsubsection*{不变映射}

设 $L$ 为 $K$ 的一个非分歧扩张 (可能无限), 并设 $G=\Gal(L/K)$.

由于 $H^2(G,U_L)=0=H^3(G,U_L)$, 短正合列
\[
 0\to U_L\to L^\times\xrightarrow{\ord_L}\Z\to 0
\]
的上同调序列给出一个同构
\[
 H^2(G,L^\times)\xrightarrow[\simeq]{\quad H^2(\ord_L)\quad}H^2(G,\Z).
\]
群 $H^r(G,\Q)$ 对 $r>0$ 是挠群 (见 II, \ref{cor:ii.4.3}) 并且是唯一可除的 (因为 $\Q$ 如此), 故为零. 因此短正合列
\[
 0\to\Z\to\Q\to\Q/\Z\to 0
\]
(带平凡的 $G$-作用) 的上同调序列给出一个同构
\[
 H^1(G,\Q/\Z)\xrightarrow{\delta}H^2(G,\Z).
\]
回顾
\[
 H^1(G,\Q/\Z)\simeq\Hom_{\mathrm{cts}}(G,\Q/\Z),
\]
并且 $G$ 有一个典范的拓扑生成元, 即 Frobenius 元素 $\sigma=\Frob_{L/K}$. 复合映射
\[
 H^2(L/K)\xrightarrow[\simeq]{\ord_L}H^2(G,\Z)\xleftarrow[\simeq]{\delta}H^1(G,\Q/\Z)\simeq\Hom_{\mathrm{cts}}(G,\Q/\Z)\xrightarrow{f\mapsto f(\sigma)}\Q/\Z
\]
称为\emph{不变映射}
\[
 \inv_{L/K}\colon H^2(L/K)\to\Q/\Z.
\]

考虑一个域扩张塔
\[
 E\supset L\supset K,
\]
其中 $E$ 与 $L$ 都在 $K$ 上非分歧 (从而 Galois). 那么
\[
\begin{tikzcd}
	{H^2(L/K)} & {\Q/\Z} \\
	{H^2(E/K)} & {\Q/\Z}
	\arrow["{\inv_{L/K}}", from=1-1, to=1-2]
	\arrow["\Inf"', from=1-1, to=2-1]
	\arrow[equal, from=1-2, to=2-2]
	\arrow["{\inv_{E/K}}"', from=2-1, to=2-2]
\end{tikzcd}
\]
交换, 因为 $\inv$ 的定义中所有的映射都与 $\Inf$ 相容.

\begin{theorem}\label{thm:iii.1.7}
存在唯一一个同构
\[
 \inv_K\colon H^2(K^{\mathrm{un}}/K)\to\Q/\Z,
\]
具有以下性质: 对每个在 $K$ 上 $n$ 次有限的 $L\subset K^{\mathrm{un}}$, $\inv_K$ 诱导出同构
\[
 \inv_{L/K}\colon H^2(L/K)\to\tfrac{1}{[L:K]}\Z/\Z.
\]
\end{theorem}

\begin{proof}
这是上面讨论的直接推论.
\end{proof}

\begin{proposition}\label{prop:iii.1.8}
设 $L$ 为 $K$ 的 $n$ 次有限扩张, 并设 $K^{\mathrm{un}}$ 与 $L^{\mathrm{un}}$ 为 $K$ 与 $L$ 的最大非分歧扩张. 则下列图表交换:
\begin{equation}
\begin{tikzcd}[column sep=small]
	{H^2(K^{\mathrm{un}}/K)} & {H^2(L^{\mathrm{un}}/L)} \\
	{\Q/\Z} & {\Q/\Z}
	\arrow["\Res", from=1-1, to=1-2]
	\arrow["{\inv_K}"', from=1-1, to=2-1]
	\arrow["{\inv_L}", from=1-2, to=2-2]
	\arrow["n"', from=2-1, to=2-2]
\end{tikzcd}\tag{27}
\end{equation}
\end{proposition}

\begin{proof}
局部域的最大非分歧扩张可由添加所有 $m$ 次单位根得到, 其中 $m$ 不被剩余特征整除. 因此 $L^{\mathrm{un}}=L\cdot K^{\mathrm{un}}$, 从而映射
\[
 \tau\mapsto\tau|K^{\mathrm{un}}\colon \Gal(L^{\mathrm{un}}/L)\to\Gal(K^{\mathrm{un}}/K)
\]
是单射. (27) 中记作 $\Res$ 的映射由相容的同态
\[
\begin{tikzcd}
	{\Gal(K^{\mathrm{un}}/K)} & {\Gal(L^{\mathrm{un}}/L)} \\
	{K^{\mathrm{un}\times}} & {L^{\mathrm{un}\times}}
	\arrow[from=2-1, to=2-2]
\end{tikzcd}
\]
所定义 (II, \ref{ex:ii.1.27}b); 其中第一行由 $\tau\mapsto\tau|K^{\mathrm{un}}$ 的反向给出.

令 $\Gamma_K=\Gal(K^{\mathrm{un}}/K)$ 与 $\Gamma_L=\Gal(L^{\mathrm{un}}/L)$, 并考虑图表
\[
\begin{tikzcd}[column sep=small, row sep=large]
	{H^2(K^{\mathrm{un}}/K)} & {H^2(\Gamma_K,\Z)} & {H^1(\Gamma_K,\Q/\Z)} & {\Q/\Z} \\
	{H^2(L^{\mathrm{un}}/L)} & {H^2(\Gamma_L,\Z)} & {H^1(\Gamma_L,\Q/\Z)} & {\Q/\Z}
	\arrow["{\ord_K}"{description}, "\simeq"{above}, from=1-1, to=1-2]
	\arrow["\delta"{description}, "\simeq"{above}, from=1-3, to=1-2]
	\arrow["{g\mapsto g(\sigma_K)}", from=1-3, to=1-4]
	\arrow["{\ord_L}"{description}, "\simeq"{above}, from=2-1, to=2-2]
	\arrow["\delta"{description}, "\simeq"{above}, from=2-3, to=2-2]
	\arrow["{g\mapsto g(\sigma_L)}"', from=2-3, to=2-4]
	\arrow["\Res"', from=1-1, to=2-1]
	\arrow["{e\,\Res}", from=1-2, to=2-2]
	\arrow["\Res"', from=1-3, to=2-3]
	\arrow["f", from=1-4, to=2-4]
\end{tikzcd}
\]
这里 $e$ 是 $L/K$ 的分歧指数, $f$ 是剩余类次数. 第一个方块来自交换方块
\[
\begin{tikzcd}
	{K^{\mathrm{un}\times}} & \Z \\
	{L^{\mathrm{un}\times}} & \Z
	\arrow["{\ord_K}", from=1-1, to=1-2]
	\arrow[from=1-1, to=2-1]
	\arrow["e", from=1-2, to=2-2]
	\arrow["{\ord_L}"', from=2-1, to=2-2]
\end{tikzcd}
\]
第二个方块表达的事实是限制映射与边缘映射可交换. 除因子 ``$e$'' 外, 第三个方块是
\[
\begin{tikzcd}
	{\Hom(\Gamma_K,\Q/\Z)} & {\Q/\Z} \\
	{\Hom(\Gamma_L,\Q/\Z)} & {\Q/\Z}
	\arrow["{g\mapsto g(\sigma_K)}", from=1-1, to=1-2]
	\arrow["{g\mapsto g|\Gamma_L}"', from=1-1, to=2-1]
	\arrow["f", from=1-2, to=2-2]
	\arrow["{g\mapsto g(\sigma_L)}"', from=2-1, to=2-2]
\end{tikzcd}
\]
Frobenius 元素 $\sigma_K$ 与 $\sigma_L$ 由它们在剩余类域上分别诱导 $x\mapsto x^q$ 与 $x\mapsto x^{q^f}$ 这一事实所确定, 其中 $q=|k|$, $q^f=|l|$, 因此 $\sigma_L|K^{\mathrm{un}}=\sigma_K^f$. 现在该方块交换是清楚的, 又因为 $n=ef$, 这就证明了命题.
\end{proof}

\subsubsection*{局部 Artin 映射}

设 $L$ 为 $K$ 的 Galois 群为 $G$ 的有限非分歧扩张, 并令 $n=[L:K]$. \emph{局部基本类} $u_{L/K}$ 是 $H^2(L/K)$ 中被不变映射
\[
 \inv_{L/K}\colon H^2(L/K)\xrightarrow{\simeq}\tfrac{1}{[L:K]}\Z/\Z
\]
映到生成元 $\tfrac{1}{[L:K]}$ 的那个元素. 偶对 $(G,L^\times)$ 满足 Tate 定理 (II, \ref{thm:ii.3.11}) 的假设,\footnote{对 $H^1$ 的条件, 用 Hilbert 定理 90; 对 $H^2$ 的条件, 用不变映射.} 于是与基本类 $u_{L/K}$ 作杯积定义了同构
\[
 H_T^r(G,\Z)\to H_T^{r+2}(G,L^\times)
\]
(对所有 $r\in\Z$). 当 $r=-2$ 时, 这成为
\[
\begin{tikzcd}
	{H^{-2}(G,\Z)} & {H^0(G,L^\times)} \\
	G & {K^\times/\Nm(L^\times)}
	\arrow["\simeq", from=1-1, to=1-2]
	\arrow[equal, from=1-1, to=2-1]
	\arrow[equal, from=1-2, to=2-2]
\end{tikzcd}
\]
我们现在明确地计算这个映射.

$K$ 的一个素元 $\pi$ 也是 $L$ 的素元, 并给出一个 $G$-模的分解
\[
 L^\times=U_L\cdot\pi^{\Z}\simeq U_L\times\Z.
\]
于是
\[
 H^r(G,L^\times)\simeq H^r(G,U_L)\oplus H^r(G,\pi^{\Z}).
\]
取 $G$ 的一个生成元 $\sigma$ (例如 Frobenius 生成元), 并令
\[
 f\in H^1(G,\Q/\Z)\simeq\Hom(G,\Q/\Z)
\]
为满足对所有 $i$ 有 $f(\sigma^i)=\tfrac{i}{n}\mod\Z$ 的元素. 它生成 $H^1(G,\Q/\Z)$.

从正合列
\[
 0\to\Z\to\Q\to\Q/\Z\to 0
\]
以及 $H^r(G,\Q)=0$ (对所有 $r$) 这一事实, 我们得到一个同构
\[
 \delta\colon H^1(G,\Q/\Z)\to H^2(G,\Z).
\]
按照 (II, \ref{rem:ii.1.21}) 中对 $\delta$ 的描述, 为构造 $\delta f$, 我们首先把 $f$ 提升为一个 $1$-上链 $\tilde f\colon G\to\Q$. 我们取 $\tilde f$ 为映射 $\sigma^i\mapsto\tfrac{i}{n}$, 其中 $0\le i<n$. 那么
\[
 d\tilde f(\sigma^i,\sigma^j)=\sigma^i\tilde f(\sigma^j)-\tilde f(\sigma^{i+j})+\tilde f(\sigma^i)=
 \begin{cases}
 0 & \text{若 } i+j\le n-1,\\
 1 & \text{若 } i+j>n-1.
 \end{cases}
\]
当把 $\Z$ 与 $L^\times$ 的子群 $\pi^{\Z}$ 等同起来时, 我们发现基本类 $u_{L/K}\in H^2(G,L^\times)$ 由上循环
\[
 \varphi(\sigma^i,\sigma^j)=
 \begin{cases}
 1 & \text{若 } i+j\le n-1,\\
 \pi & \text{若 } i+j>n-1
 \end{cases}
\]
表示.

从正合列
\[
 0\to I\to\Z[G]\to\Z\to 0
\]
\[
 0\to L^\times\to L^\times(\varphi)\to I\to 0
\]
(见 II, \ref{thm:ii.3.11} 的证明) 我们得到边缘映射
\[
 H^{-2}(G,\Z)\to H^{-1}(G,I)
\]
\[
 H^{-1}(G,I)\to H^0(G,L^\times),
\]
因为 $\Z[G]$ 与 $L^\times(\varphi)$ 有平凡的上同调, 它们是同构. 这里 $L^\times(\varphi)$ 是 $\varphi$ 的分裂模 $L^\times\oplus\bigoplus_{\sigma\in G,\sigma\neq 1}\Z x_\sigma$.

最后, $H^{-2}(G,\Z)\overset{\mathrm{def}}{=}H_1(G,\Z)\simeq G$ (见 II, \ref{prop:ii.2.7}).

\begin{proposition}\label{prop:iii.1.9}
在复合映射
\[
\begin{tikzcd}
	{H^{-2}(G,\Z)} & {H^0(G,L^\times)} \\
	G & {K^\times/\Nm(L^\times)}
	\arrow["\simeq", from=1-1, to=1-2]
	\arrow[equal, from=1-1, to=2-1]
	\arrow[equal, from=1-2, to=2-2]
\end{tikzcd}
\]
下, Frobenius 元素 $\sigma\in G$ 映到 $\pi$ 在 $K^\times/\Nm(L^\times)$ 中的类.
\end{proposition}

注意, 因为 $K$ 中所有单位都是来自 $L^\times$ 的范 (见 \ref{prop:iii.1.2}), $\pi\mod\Nm(L^\times)$ 的类与素元 $\pi$ 的选取无关. 另一方面, $G$-模 $L^\times(\varphi)$ 以及该映射依赖于 $G$ 的生成元 $\sigma$ 的选取.

\begin{proof}
从同构 $H^{-2}(G,\Z)\simeq G$ 的构造我们看到, $\sigma$ 在边缘映射 $H^{-2}(G,\Z)\to H^{-1}(G,I_G)\subset I_G/I_G^2$ 下的像由 $\sigma-1$ 表示.

边缘映射 $H^{-1}(G,I_G)\to H^0(G,L^\times)$ 是由下列图表经蛇形引理给出的那个映射 (我们把 $I_G$ 记作 $I$):
\[
\begin{tikzcd}[column sep=small]
	0 & {L^\times} & {L^\times(\varphi)} & I & 0 \\
	0 & {L^\times} & {L^\times(\varphi)} & I & 0
	\arrow[from=1-1, to=1-2]
	\arrow[from=1-2, to=1-3]
	\arrow[from=1-3, to=1-4]
	\arrow[from=1-4, to=1-5]
	\arrow["{\Nm_G}"', from=1-2, to=2-2]
	\arrow["{\Nm_G}"', from=1-3, to=2-3]
	\arrow["{\Nm_G}"', from=1-4, to=2-4]
	\arrow[from=2-1, to=2-2]
	\arrow[from=2-2, to=2-3]
	\arrow[from=2-3, to=2-4]
	\arrow[from=2-4, to=2-5]
\end{tikzcd}
\]
连接两行的竖直映射是 $\Nm_G=\sum_{i=0}^{n-1}\sigma^i$. 元素 $(\sigma-1)+I^2$ 是 $x_\sigma+I\cdot L^\times(\varphi)$ 在 $L^\times(\varphi)_G$ 中的像, 而 $\Nm_G(x_\sigma+I\cdot L^\times(\varphi))$ 是下列元素之和:
\[
\begin{aligned}
 x_\sigma &= x_\sigma\\
 \sigma x_\sigma &= x_{\sigma^2}-x_\sigma+\varphi(\sigma,\sigma)\\
 \sigma^2 x_\sigma &= x_{\sigma^3}-x_{\sigma^2}+\varphi(\sigma,\sigma^2)\\
 &\;\;\vdots\\
 \sigma^{n-1}x_\sigma &= x_1-x_{\sigma^{n-1}}+\varphi(\sigma,\sigma^{n-1}).
\end{aligned}
\]
把这些加起来, 记住在 $L^\times(\varphi)$ 的因子 $L^\times$ 上 ``$+$'' 实际上是 ``$\cdot$'', 我们得到
\[
 \Nm_G(x_\sigma)=\prod_{i=1}^{n-1}\varphi(\sigma,\sigma^i)=\pi.
\]
证明完成.
\end{proof}

\begin{remark}\label{rem:iii.1.10}
上面这个命题\ref{prop:iii.1.9} 的证明使用了 Tate 关于同构 $H^r(G,\Z)\to H^{r+2}(G,C)$ 的原始定义, 它比其他文献中用杯积描述该映射的证明更简单.
\end{remark}

% ==================================================================
\section{分歧扩张的上同调}\label{sec:III.2}\index{fenqikuozhang@分歧扩张(ramified extension)}

由于 Hilbert 定理 90 (II, \ref{prop:ii.1.22}), 对任意 Galois 扩张塔 $E\supset L\supset K$ 都有一个正合列
\[
 0\to H^2(L/K)\xrightarrow{\Inf}H^2(E/K)\xrightarrow{\Res}H^2(E/L)
\]
(见 II, \ref{ex:ii.1.36}).

下一个定理把定理\ref{thm:iii.1.7} 推广到分歧扩张.

\begin{theorem}\label{thm:iii.2.1}
对每个局部域 $K$, 存在一个典范同构
\[
 \inv_K\colon H^2(K^{\mathrm{al}}/K)\to\Q/\Z.
\]
设 $L$ 为 $K$ 的 $n<\infty$ 次 Galois 扩张. 那么图表
\begin{equation}
\begin{tikzcd}[column sep=small]
	0 & {H^2(L/K)} & {H^2(K^{\mathrm{al}}/K)} & {H^2(K^{\mathrm{al}}/L)} \\
	0 & {\frac{1}{n}\Z/\Z} & {\Q/\Z} & {\Q/\Z}
	\arrow[from=1-1, to=1-2]
	\arrow[from=1-2, to=1-3]
	\arrow["\Res", from=1-3, to=1-4]
	\arrow["{\inv_{L/K}}"', from=1-2, to=2-2]
	\arrow["{\inv_K}"', from=1-3, to=2-3]
	\arrow["{\inv_L}", from=1-4, to=2-4]
	\arrow[from=2-1, to=2-2]
	\arrow[from=2-2, to=2-3]
	\arrow["n"', from=2-3, to=2-4]
\end{tikzcd}\tag{28}
\end{equation}
交换, 从而定义了一个同构
\[
 \inv_{L/K}\colon H^2(L/K)\to\tfrac{1}{n}\Z/\Z.
\]
\end{theorem}

证明将占据本节余下的大部分篇幅.

\begin{lemma}\label{lem:iii.2.2}
若 $L/K$ 是 $n$ 次 Galois 扩张, 则 $H^2(L/K)$ 包含一个典范同构于 $\tfrac{1}{n}\Z/\Z$ 的子群.
\end{lemma}

\begin{proof}
考虑图表
\[
\begin{tikzcd}[column sep=small]
	0 & {\Ker(\Res)} & {H^2(K^{\mathrm{un}}/K)} & {H^2(L^{\mathrm{un}}/L)} \\
	0 & {H^2(L/K)} & {H^2(K^{\mathrm{al}}/K)} & {H^2(K^{\mathrm{al}}/L)}
	\arrow[from=1-1, to=1-2]
	\arrow[from=1-2, to=1-3]
	\arrow["\Res", from=1-3, to=1-4]
	\arrow[from=1-2, to=2-2]
	\arrow["\Inf"', from=1-3, to=2-3]
	\arrow["\Inf", from=1-4, to=2-4]
	\arrow[from=2-1, to=2-2]
	\arrow[from=2-2, to=2-3]
	\arrow["\Res"', from=2-3, to=2-4]
\end{tikzcd}
\]
由于两个膨胀映射都是单射, 第一个竖直映射也是单射, 而 \ref{prop:iii.1.8} 表明顶行限制映射的核典范同构于 $\tfrac{1}{n}\Z/\Z$.
\end{proof}

为完成定理\ref{thm:iii.2.1} 的证明, 只需证明映射 $\tfrac{1}{n}\Z/\Z\hookrightarrow H^2(L/K)$ 是同构. 有两种不同的途径来证明这一点. 在下一章关于 Brauer 群的内容中, 我们将证明 $H^2(K^{\mathrm{un}}/K)\simeq H^2(K^{\mathrm{al}}/K)$ ($K$ 上每个中心单代数都被一个非分歧扩张分裂; 见 IV, \S 4), 于是上面图表中的竖直映射都是同构. 第二种证明, 也就是我们现在给出的证明, 则表明 $|H^2(L/K)|\le n$.

\begin{lemma}\label{lem:iii.2.3}
设 $L$ 为 $K$ 的 Galois 群为 $G$ 的有限 Galois 扩张. 则存在 $\mathcal{O}_L$ 的一个在 $G$ 下稳定的开子群 $V$, 使得对所有 $r>0$ 有 $H^r(G,V)=0$.
\end{lemma}

\begin{proof}
设 $\{x_\tau\mid\tau\in G\}$ 为 $L$ 在 $K$ 上的一个正规基 (见 FT, 5.18). 诸 $x_\tau$ 在 $\mathcal{O}_K$ 中有一个公分母 $d$ (见 ANT, 2.6). 把每个 $x_\tau$ 换成 $d\cdot x_\tau$ 之后, 可设它们都落在 $\mathcal{O}_L$ 中. 取 $V=\sum\mathcal{O}_K x_\tau$. 因为正规基在 $G$ 下稳定, $V$ 也在 $G$ 下稳定. 设 $\pi$ 为 $\mathcal{O}_L$ 的一个素元. 那么对某个 $m>0$ 有 $V\supset\pi^m\mathcal{O}_L$, 这表明 $V$ 是开的 (它是 $\pi^m\mathcal{O}_L$ 的陪集之并). 最后, 作为 $G$-模
\[
 V\simeq\mathcal{O}_K[G]\simeq\Ind^G\mathcal{O}_K,
\]
于是对所有 $r>0$ 有 $H^r(G,V)=0$ (II, \ref{cor:ii.1.12}).
\end{proof}

\begin{lemma}\label{lem:iii.2.4}
设 $L$, $K$, $G$ 如上个引理. 则存在 $U_L$ 的一个在 $G$ 下稳定的开子群 $V$, 使得对所有 $r>0$ 有 $H^r(G,V)=0$.
\end{lemma}

\begin{proof}
我只对特征零的 $K$ 给出证明.\footnote{特征 $p$ 情形的证明见 \cite{Serre1967b}, p.~134.} 幂级数
\[
 e^x=1+x+\cdots+x^n/n!+\cdots
\]
在 $\ord(x)>\ord(p)/(p-1)$ 时收敛 (ANT, 7.29). 它定义了 $L$ 中 $0$ 的一个开邻域到 $L^\times$ 中 $1$ 的一个开邻域上的同构, 其逆映射为
\[
 \log(x)=(x-1)-(x-1)^2/2+(x-1)^3/3-\cdots.
\]
两个映射都与 $G$ 的作用可交换 (因为 $G$ 作用是连续的). 若 $V'$ 是 (\ref{lem:iii.2.3}) 中那样的 $0$ 的开邻域, 则 $\pi^M V'$ 有同样的性质, 于是我们可以取 $V=\exp(\pi^M V')$, 其中 $M$ 取得足够大, 使得指数函数在 $\pi^M V'$ 上有定义并且是同构.
\end{proof}

\begin{lemma}\label{lem:iii.2.5}
设 $L/K$ 为 $n$ 次循环扩张; 则 $h(U_L)=1$ 且 $h(L^\times)=n$.
\end{lemma}

\begin{proof}
设 $V$ 为 $U_L$ 的一个开子群, 满足对所有 $r$ 有 $H^r(G,V)=0$. 因为 $U_L$ 是紧的, 商群 $U_L/V$ 有限, 于是依 (II, \ref{cor:ii.3.9}) 有 $h(U_L)=h(V)=1$. 现在 $h(L^\times)=h(U_L)\cdot h(\Z)=h(\Z)$, 且
\[
 h(\Z)=\bigl|H_T^0(G,\Z)\bigr|/\bigl|H^1(G,\Z)\bigr|=\bigl|H_T^0(G,\Z)\bigr|=\bigl|\Z/n\Z\bigr|=n.
\]
\end{proof}

\begin{lemma}\label{lem:iii.2.6}
设 $L$ 为 $K$ 的 $n$ 阶有限 Galois 扩张, 则 $H^2(L/K)$ 的阶为 $n$.
\end{lemma}

\begin{proof}
我们知道 $H^2(L/K)$ 的阶被 $n$ 整除, 并且当 $L/K$ 循环时它等于 $n$. 我们对 $[L:K]$ 作归纳来证明引理. 因为群 $\Gal(L/K)$ 是可解的 (ANT, 7.59), 存在 Galois 扩张 $K'/K$ 满足 $L\supsetneq K'\supsetneq K$. 从正合列
\[
 0\to H^2(K'/K)\to H^2(L/K)\to H^2(L/K')
\]
我们看出
\[
 \bigl|H^2(L/K)\bigr|\le\bigl|H^2(K'/K)\bigr|\cdot\bigl|H^2(L/K')\bigr|=n.
\]
\end{proof}

\begin{proof}[定理\ref{thm:iii.2.1} 的证明]
从 (\ref{lem:iii.2.2}) 证明中的图表我们看到, 对 $K$ 的每个有限 Galois 扩张 $L$, $H^2(K^{\mathrm{al}}/K)$ 的子群 $H^2(L/K)$ 包含在 $H^2(K^{\mathrm{un}}/K)$ 中. 由于 $H^2(K^{\mathrm{al}}/K)=\bigcup H^2(L/K)$, 这证明了膨胀映射 $H^2(K^{\mathrm{un}}/K)\to H^2(K^{\mathrm{al}}/K)$ 是同构. 把它求逆再与定理\ref{thm:iii.1.7} 的不变映射 $\inv_K\colon H^2(K^{\mathrm{un}}/K)\to\Q/\Z$ 复合即可.
\end{proof}

事实上, 定理\ref{thm:iii.2.1} 中的图表即使当 $L/K$ 不是 Galois 扩张时也交换:
\begin{equation}
 \inv_L\circ\Res=[L:K]\inv_K.\tag{29}
\end{equation}
(在这种情形下, 我们把 $H^2(L/K)$ 定义为限制映射的核.)

我断言
\begin{equation}
 \inv_K\circ\Cor=\inv_L.\tag{30}
\end{equation}
因为 $\Res$ 是满射 (见图表 (28)), 为证明这一点只需证明
\[
 \inv_K\circ\Cor\circ\Res=\inv_L\circ\Res.
\]
但依 (II, \ref{prop:ii.1.30}) 有 $\Cor\circ\Res=[L:K]$, 于是这就是前面的等式 (29).

\subsubsection*{基本类}

设 $L$ 为 $K$ 的 Galois 群为 $G$ 的有限 Galois 扩张. 如同非分歧的情形, 我们定义 $L/K$ 的\emph{基本类} $u_{L/K}$ 为 $H^2(L/K)$ 中满足
\[
 \inv_{L/K}(u_{L/K})=\frac{1}{[L:K]}\bmod\Z
\]
的元素, 等价地, 满足
\[
 \inv_K(u_{L/K})=\frac{1}{[L:K]}\bmod\Z
\]
的元素.

\begin{lemma}\label{lem:iii.2.7}
设 $L\supset E\supset K$, 其中 $L/K$ 有限 Galois. 则
\begin{equation}
 \Res(u_{L/K})=u_{L/E},\tag{31}
\end{equation}
\begin{equation}
 \Cor(u_{L/E})=[E:K]u_{L/K}.\tag{32}
\end{equation}
此外, 若 $E/K$ 也是 Galois 扩张, 则
\begin{equation}
 \Inf(u_{E/K})=[L:E]u_{L/K}.\tag{33}
\end{equation}
\end{lemma}

\begin{proof}
考虑
\[
\begin{tikzcd}[column sep=small]
	{H^2(K^{\mathrm{al}}/K)} & {H^2(K^{\mathrm{al}}/E)} & {H^2(K^{\mathrm{al}}/L)} \\
	{\Q/\Z} & {\Q/\Z} & {\Q/\Z}
	\arrow["\Res", from=1-1, to=1-2]
	\arrow["\Res", from=1-2, to=1-3]
	\arrow["{\inv_K}"', "\simeq"', from=1-1, to=2-1]
	\arrow["{\inv_E}"', "\simeq"', from=1-2, to=2-2]
	\arrow["{\inv_L}", "\simeq"', from=1-3, to=2-3]
	\arrow["{[E:K]}"', from=2-1, to=2-2]
	\arrow["{[L:E]}"', from=2-2, to=2-3]
\end{tikzcd}
\]
把核--余核引理 (II, \ref{lem:ii.A.2}) 用于两行 (并利用该引理的函子性), 我们得到一个交换图表
\[
\begin{tikzcd}[column sep=small]
	0 & {H^2(E/K)} & {H^2(L/K)} & {H^2(L/E)} \\
	0 & {\frac{1}{[E:K]}\Z/\Z} & {\frac{1}{[L:K]}\Z/\Z} & {\frac{1}{[L:E]}\Z/\Z}
	\arrow[from=1-1, to=1-2]
	\arrow["\Inf", from=1-2, to=1-3]
	\arrow["\Res", from=1-3, to=1-4]
	\arrow["{\inv_{E/K}}"', from=1-2, to=2-2]
	\arrow["{\inv_{L/K}}"', from=1-3, to=2-3]
	\arrow["{\inv_{L/E}}", from=1-4, to=2-4]
	\arrow[from=2-1, to=2-2]
	\arrow["\id"{description}, from=2-2, to=2-3]
	\arrow["{[E:K]}"', from=2-3, to=2-4]
\end{tikzcd}
\]
这里 $H^2(E/K)$ 定义为 $\Res\colon H^2(K^{\mathrm{al}}/K)\to H^2(K^{\mathrm{al}}/E)$ 的核. 等式 (31) 表达的是第二个方块交换. 当 $E/K$ 为 Galois 扩张时, 等式 (33) 表达的是第一个方块交换.

等式 (32) 可由比较
\[
\begin{aligned}
 \inv_K(\Cor_{E/K}(u_{L/E}))&\overset{(30)}{=}\inv_E(u_{L/E})=\frac{1}{[L:E]}\bmod\Z,\\
 \inv_K(u_{L/K})&=\frac{1}{[L:K]}\bmod\Z
\end{aligned}
\]
得到.
\end{proof}

\begin{remark}\label{rem:iii.2.8}
对每个局部域的有限扩张 $L/K$ 都附有群 $K^\times/\Nm(L^\times)$ 与 $H^2(L/K)$. 当 $L/K$ 循环时它们同构, 否则不同构. 第一个群总是同构于 $\Gal(M/K)$, 其中 $M$ 是 $L/K$ 中最大的 Abel 子扩张 (见 \ref{thm:iii.3.4}, \ref{thm:iii.3.5}), 而第二个群总是 $[L:K]$ 阶循环群. 于是, 当 $L/K$ Abel 但不循环时, 这两个群有相同的阶但不同构; 而当 $L/K$ 不 Abel 时, 它们的阶不同.
\end{remark}

% ==================================================================
\section{局部 Artin 映射}\label{sec:III.3}\index{jubuartinying@局部 Artin 映射(local Artin map)}\index{jubuhuxiang@局部互反映射(local reciprocity map)}

偶对 $(G,L^\times)$ 满足 Tate 定理 (II, \ref{thm:ii.3.11}) 的假设, 于是我们证明了下述结果.

\begin{theorem}\label{thm:iii.3.1}
对局部域的每个有限 Galois 扩张 $L/K$ 以及 $r\in\Z$, 由 $x\mapsto x\cup u_{L/K}$ 定义的同态
\[
 H_T^r(\Gal(L/K),\Z)\to H_T^{r+2}(\Gal(L/K),L^\times)
\]
是同构. 当 $r=-2$ 时, 这成为同构
\[
 G^{\mathrm{ab}}\simeq K^\times/\Nm_{L/K}(L^\times).
\]

若 $L\supset E\supset K$ 且 $L/K$ Galois, 则定理中的同态与 $\Res$ 及 $\Cor$ 都可交换. 例如, 为证明关于 $\Res$ 的陈述, 我们必须证明
\[
 \Res(x\cup u_{L/K})=\Res(x)\cup u_{L/E}
\]
对所有 $x\in H_T^r(\Gal(L/K),L^\times)$ 成立. 但杯积的一个标准性质 (II, \ref{prop:ii.1.39}c) 是
\[
 \Res(x\cup u_{L/K})=\Res(x)\cup\Res(u_{L/K}),
\]
于是这由 (31) 推出. 关于 $\Cor$ 的证明类似.

若 $L\supset E\supset K$, 其中 $L/K$ 与 $E/K$ 都 Galois, 且 $x\in H^r(\Gal(E/K),E^\times)$, $r\ge 1$, 则
\[
 \Inf(x\cup u_{E/K})=[L:E]\Inf(x)\cup u_{L/K}.
\]
同样, 这由杯积的一个标准性质 (\ref{prop:ii.1.39}e) 以及公式 (33) 推出.
\end{theorem}

当 $r=-2$ 时, 定理中的映射成为
\[
 \Gal(L/K)^{\mathrm{ab}}\xrightarrow{\simeq}K^\times/\Nm_{L/K}L^\times.
\]
我们把它的逆映射记作
\[
 \phi_{L/K}\colon K^\times/\Nm_{L/K}L^\times\to\Gal(L/K)^{\mathrm{ab}}
\]
并称之为\emph{局部 Artin 映射}, 或者记作 $\mathrm{rec}_{L/K}$ 并称之为\emph{局部互反映射}. (后一个名字更常见, 但很难看出这个映射里有什么互反性.)

\begin{lemma}\label{lem:iii.3.2}
设 $L\supset E\supset K$ 为局部域, 且 $L/K$ Galois. 则下列图表交换:
\[
\begin{tikzcd}
	{E^\times} & {\Gal(L/E)^{\mathrm{ab}}} \\
	{K^\times} & {\Gal(L/K)^{\mathrm{ab}}}
	\arrow["{\phi_{L/E}}", from=1-1, to=1-2]
	\arrow["{\Nm_{E/K}}"', from=1-1, to=2-1]
	\arrow[from=1-2, to=2-2]
	\arrow["{\phi_{L/K}}"', from=2-1, to=2-2]
\end{tikzcd}\qquad
\begin{tikzcd}
	{E^\times} & {\Gal(L/E)^{\mathrm{ab}}} \\
	{K^\times} & {\Gal(L/K)^{\mathrm{ab}}}
	\arrow["{\phi_{L/E}}", from=1-1, to=1-2]
	\arrow[from=2-1, to=1-1]
	\arrow["{\phi_{L/K}}"', from=2-1, to=2-2]
	\arrow["{\operatorname{Ver}}", from=2-2, to=1-2]
\end{tikzcd}
\]
两图中未标名的竖直箭头分别由包含映射 $\Gal(L/E)\hookrightarrow\Gal(L/K)$ 与 $K^\times\hookrightarrow E^\times$ 诱导; 第二个图右侧的竖直箭头为转移映射 $\operatorname{Ver}$.
\end{lemma}

\begin{proof}
待补 (第二个图表其实并不需要).
\end{proof}

\noindent\textbf{3.3} 设 $L\supset E\supset K$ 为局部域, 且 $L$ 与 $E$ 都在 $K$ 上 Galois. 则下列图表交换:
\[
\begin{tikzcd}[column sep=small]
	{K^\times} & {\Gal(L/K)^{\mathrm{ab}}} \\
	& {\Gal(E/K)^{\mathrm{ab}}}
	\arrow["{\phi_{L/K}}", from=1-1, to=1-2]
	\arrow["{\phi_{E/K}}"', from=1-1, to=2-2]
	\arrow[from=1-2, to=2-2]
\end{tikzcd}
\]
未标名的箭头由满射 $\sigma\mapsto\sigma|E\colon\Gal(L/K)\to\Gal(E/K)$ 诱导.

\begin{proof}
这直接由局部 Artin 映射的定义并利用 (33) 推出.\footnote{Hervé Jacquet 指出, 这个陈述似乎并不正确, 但利用命题\ref{prop:iii.3.6} 可以证明它; 我不知道不用 \ref{prop:iii.3.6} 能否证明.}
\end{proof}

% 3.3 是不带环境名的编号条目, 不占定理计数器; 调至 3 使下一条显示为 3.4
\setcounter{theorem}{3}

特别地, 若 $L\supset E\supset K$ 是 $K$ 的有限 Abel 扩张塔, 则对所有 $a\in K^\times$ 有 $\phi_{L/K}(a)|E=\phi_{E/K}(a)$, 于是我们可以定义 $\varphi_K\colon K^\times\to\Gal(K^{\mathrm{ab}}/K)$ 为这样的同态: 对每个有限 Abel 扩张 $L/K$, $\varphi_K(a)|L=\varphi_{L/K}(a)$.

\begin{theorem}\label{thm:iii.3.4}
对每个局部域 $K$, 存在一个同态 (局部 Artin 映射)
\[
 \phi_K\colon K^\times\to\Gal(K^{\mathrm{ab}}/K),
\]
具有以下性质:
\begin{enumerate}[(a)]
\item 对 $K$ 的每个素元 $\pi$, $\phi_K(\pi)|K^{\mathrm{un}}=\Frob_K$;
\item 对 $K$ 的每个有限 Abel 扩张 $L$, $\Nm_{L/K}(L^\times)$ 包含在映射 $a\mapsto\phi_K(a)|L$ 的核中, 且 $\phi_K$ 诱导出同构
\[
 \phi_{L/K}\colon K^\times/\Nm_{L/K}(L^\times)\to\Gal(L/K).
\]
\end{enumerate}
\end{theorem}

\begin{proof}
除了 (a) 之外一切都由上述内容显然得出; 而 (a) 可由以下事实推出: 对 $K$ 的非分歧扩张 $L$, $\phi_{L/K}$ 与 \S 1 中定义的那个一致, 于是我们可应用命题\ref{prop:iii.1.9}.
\end{proof}

对读过第 \ref{ch:I} 章 \S 3 的读者来说, 这就完成了局部类域论主定理的证明 —— 如果愿意, 他们现在可以跳到第 \ref{ch:V} 章. 其他读者则需要继续读至本章结束.

\begin{theorem}[范限制定理]\label{thm:iii.3.5}
设 $L$ 为 $K$ 的有限扩张, 并设 $E$ 为 $L$ 中包含的 $K$ 的最大 Abel 扩张; 则
\[
 \Nm_{L/K}(L^\times)=\Nm_{E/K}(E^\times).
\]
\end{theorem}

\begin{proof}
由范映射的传递性, $\Nm_{L/K}(L^\times)$ 是 $\Nm_{E/K}(E^\times)$ 的子群. 若 $L/K$ 是 Galois 扩张, 则 $\Gal(E/K)=\Gal(L/K)^{\mathrm{ab}}$, 于是定理\ref{thm:iii.3.1} 表明两个范群在 $K^\times$ 中都有指数 $[E:K]$. 这意味着它们相等.

在一般情形下, 我们取 $K$ 的一个包含 $L$ 的有限 Galois 扩张 $L'$. 令 $G=\Gal(L'/K)$, $H=\Gal(L'/L)$. 于是
\[
\begin{tikzcd}
	{L'} \\
	L \\
	E \\
	K
	\arrow[from=1-1, to=2-1]
	\arrow[from=2-1, to=3-1]
	\arrow[from=3-1, to=4-1]
\end{tikzcd}
\]
那么 $E$ 是 $L'$ 中在 $K$ 上 Abel 且包含于 $L$ 的最大子域. 因此固定它的 $G$ 的子群是 $G'\cdot H$, 其中 $G'$ 是 $G$ 的派生群. 设 $a\in\Nm(E^\times)$. 我们必须证明 $a\in\Nm(L^\times)$. 考虑交换图表
\[
\begin{tikzcd}
	{L^\times} & {H/H'} \\
	{K^\times} & {G/G'} \\
	{K^\times} & {G/G'H}
	\arrow["{\phi_{L'/L}}", from=1-1, to=1-2]
	\arrow["\Nm"', from=1-1, to=2-1]
	\arrow[from=1-2, to=2-2]
	\arrow["{\phi_{L'/K}}", from=2-1, to=2-2]
	\arrow[equal, from=2-1, to=3-1]
	\arrow[from=2-2, to=3-2]
	\arrow["{\phi_{E/K}}"', from=3-1, to=3-2]
\end{tikzcd}
\]
$G/G'$ 中的元素 $\phi_{L'/K}(a)$ 在 $G/G'H$ 中映到 $1$. 由于 $\phi_{L'/L}$ 是满射, 我们从图表中看出存在 $b\in L^\times$ 使得 $\phi_{L'/K}(a)=\phi_{L'/K}(\Nm(b))$, 从而 $a/\Nm(b)\in\Nm(L^\times)$, 设 $a/\Nm(b)=\Nm(c)$. 现在
\[
 a=\Nm_{L/K}(b\cdot\Nm_{L'/L}(c))\in\Nm_{L/K}(L^\times).
\]
\end{proof}

定理\ref{thm:iii.3.5} 表明, 想用范群来分类局部域的非 Abel 扩张是没有希望的.

\subsubsection*{局部 Artin 映射的另一描述}

设 $L/K$ 为 Galois 群为 $G$ 的有限 Abel 扩张, 并设 $u_{L/K}\in H^2(G,L^\times)$ 为基本类. 局部 Artin 映射 $\phi_{L/K}$ 是同构
\[
 x\mapsto x\cup u_{L/K}\colon H_T^{-2}(G,\Z)\to H_T^0(G,L^\times)
\]
的逆. 这个定义难以操作, 因为同时涉及同调群与上同调群的杯积没有很方便的描述. 我们改为纯粹用上同调群来重新解释这个映射. 考虑杯积配对
\[
 H^0(G,L^\times)\times H^2(G,\Z)\ \longrightarrow\ H^2(G,L^\times)\xrightarrow{\inv_{L/K}}\Q/\Z.
\]
给定元素 $a\in H^0(G,L^\times)=K^\times$ 以及由上循环 $f\colon G\times G\to\Z$ 表示的类 $c\in H^2(G,\Z)$, 杯积类 $a\cup c$ 由上循环 $(\sigma,\tau)\mapsto a^{f(\sigma,\tau)}$ 表示. 还要回顾我们有一个同构
\[
 \Hom(G,\Q/\Z)=H^1(G,\Q/\Z)\xrightarrow{\delta}H^2(G,\Z).
\]

\begin{proposition}\label{prop:iii.3.6}
对每个 $\chi\in\Hom_{\mathrm{cts}}(G,\Q/\Z)$ 与 $a\in K^\times$,
\[
 \chi(\phi_{L/K}(a))=\inv_K(a\cup\delta\chi).
\]
\end{proposition}

\begin{proof}
见 \cite{Serre1962}, 第 XI 章的 ``附录'', 以及 \cite{Serre1967a}, p. 140.
\end{proof}

利用这一点, 我们可以给出命题\ref{thm:iii.3.4} 的另一个证明.

\begin{lemma}\label{lem:iii.3.7}
若 $L/K$ 非分歧, 则 $\phi_{L/K}$ 把 $a\in K^\times$ 送到 $\Frob^{\ord_K(a)}$.
\end{lemma}

\begin{proof}
回顾 $\inv_K$ 定义为复合映射
\[
 H^2(G,L^\times)\xrightarrow{\ord}H^2(G,\Z)\xleftarrow{\delta}\Hom(G,\Q/\Z)\xrightarrow{\chi\mapsto\chi(\sigma)}\Q/\Z.
\]
由杯积的函子性
\[
 \ord(a\cup\delta\chi)=\ord(a)\cup\delta\chi,\quad a\in K^\times,\quad \chi\in\Hom(G,\Q/\Z),
\]
这里左边的 $\ord$ 表示由 $\ord_L\colon L^\times\to\Z$ 在 $H^2$ 上诱导的映射, 右边则是该映射本身. 设 $a\in H^0(G,L^\times)=K^\times$, 并令 $m=\ord_L(a)$. 对每个 $\chi\in\Hom(G,\Q/\Z)$, 在上面的图表中
\[
 a\cup\delta\chi\mapsto\ord(a)\cup\delta\chi\mapsto m\chi\mapsto\chi(\sigma^m),\quad \sigma=\Frob,
\]
即 $\inv_K(a\cup\delta\chi)=\chi(\sigma^m)$. 把它与 (\ref{prop:iii.3.6}) 中的公式结合起来, 我们得到
\[
 \chi(\phi_{L/K}(a))=\chi(\sigma^{\ord_K(a)})
\]
对所有 $\chi\in\Hom(G,\Q/\Z)$ 成立, 从而 $\phi_{L/K}(a)=\sigma^{\ord_K(a)}$.
\end{proof}

对 $G$ 的每个特征标 $\chi$, 我们得到 $K^\times$ 的一个特征标 $a\mapsto\inv_K(a\cup\delta\chi)$. 由对偶性我们得到一个映射 $K^\times\to G$. 命题\ref{prop:iii.3.6} 表明这个映射就是 $\phi_{L/K}\colon K^\times\to\Gal(L/K)$.

% ==================================================================
\section{Hilbert 符号}\label{sec:III.4}\index{hilbertfuhao@Hilbert 符号(Hilbert symbol)}

由于 $K^\times$ 中有限指数子群的交是 $\{1\}$, 存在定理蕴含范群的交为 $1$. 下述关键结果是我们证明存在定理的出发点.

\begin{proposition}\label{prop:iii.4.1}
设 $K$ 为一个含有 $n$ 次本原单位根的局部域. $K^\times$ 中任何元素, 若它是 $K$ 的每个次数整除 $n$ 的循环扩张的范, 则它是一个 $n$ 次幂.
\end{proposition}

这将由 Hilbert 符号的理论推出.

\subsection*{一个非常特殊的情形}

在这一小节中, 我们说明如何只用 \cite{Serre1970} 第 III 章中的初等结果来得到 $K=\Q_p$, $n=2$ 这一特殊情形下的命题.

对 $a,b\in\Q_p^\times$, 定义
\[
 (a,b)_p=
 \begin{cases}
 \hphantom{-}1 & \text{若 } z^2=ax^2+by^2 \text{ 在 } \Q_p \text{ 中有非平凡解},\\
 -1 & \text{其他情形}.
 \end{cases}
\]
显然 $(a,b)_p$ 只依赖于 $a$ 与 $b$ 模平方的类, 于是这是一个配对
\[
 a,b\mapsto(a,b)_p\colon \Q_p^\times/\Q_p^{\times 2}\times\Q_p^\times/\Q_p^{\times 2}\to\{\pm 1\}.
\]
Serre 证明了 (同上, 1.1, 1.2) 这个配对是双乘性的
\[
 (aa',b)_p=(a,b)_p(a',b)_p,\quad (a,bb')_p=(a,b)_p(a,b')_p,
\]
对称的 (在这种情形下与斜对称是一回事)
\[
 (b,a)_p=(a,b)_p^{-1}=(a,b)_p,
\]
并且是非退化的 (即左核与右核都是平凡的).

设 $a$ 为 $\Q_p^\times$ 中一个非平方元. 那么
\begin{align*}
 &b \text{ 是来自 } \Q_p[\sqrt{a}] \text{ 的范}\\
 \iff{}& b=z^2-ax^2 \text{ 在 } \Q_p \text{ 中有解}\\
 \iff{}& z^2=ax^2+by^2 \text{ 在 } \Q_p \text{ 中有非平凡解}\\
 \iff{}& (a,b)_p=1.
\end{align*}
于是, 由配对的退化性,
\[
 b \text{ 对所有 } a \text{ 都是来自 } \Q_p[\sqrt{a}] \text{ 的范}\implies (a,b)_p=1 \text{ 对所有 } a \implies b\in\Q_p^{\times 2}.
\]
这就在此情形下证明了命题.

\begin{aside}\label{aside:iii.4.2}
映射 $a,b\mapsto(a,b)_p$ 称为 Hilbert 符号. 它对 $\R$ 也有定义, 并且 Serre 证明了 (Hilbert 乘积公式)
\begin{equation}
 \prod_{p\le\infty}(a,b)_p=1.\tag{34}
\end{equation}
\end{aside}

\subsubsection*{一般情形}

在这一小节中, $\Hom(G,\Z/n\Z)$ 表示连续同态 $G\to\Z/n\Z$ 构成的群.

\paragraph*{步骤 1: 回顾.} 回顾对每个循环群 $\bar G$ 与 $\bar G$-模 $M$, 都有同构
\[
 H_T^r(\bar G,M)\to H_T^{r+2}(\bar G,M),
 \]
一旦选定 $\bar G$ 的一个生成元 $\sigma$, 它们就成为典范的. 更确切地说, 令 $n=(\bar G:1)$, 并设 $\chi$ 为把 $\sigma$ 送到 $1\mod n$ 的同构 $\bar G\to\Z/n\Z$. 从
\[
 0\to\Z\xrightarrow{n}\Z\to\Z/n\Z\to 0
\]
的上同调序列, 我们得到一个边缘映射 $\delta\colon H^1(\bar G,\Z/n\Z)\to H^2(\bar G,\Z)$. 同构
\[
 H_T^r(\bar G,M)\to H_T^{r+2}(\bar G,M)
\]
就是 $b\mapsto\delta\chi\cup b$.

现在设 $K$ 为一个域 (不必是局部域), 并设 $L/K$ 为 Galois 群为 $\bar G$ (阶为 $n$) 的循环扩张. 当 $r=0$ 且 $M=L^\times$ 时, 该同构成为
\[
 b\mapsto\delta\chi\cup b\colon K^\times/\Nm L^\times\to H^2(\bar G,L^\times).
\]
对 $\chi\in\Hom(\bar G,\Z/n\Z)$ 与 $b\in K^\times$, 定义
\[
 (\chi,b)'=\delta\chi\cup b\in H^2(\bar G,L^\times).
\]
上面的讨论表明, 若 $\chi$ 是同构 $\bar G\to\Z/n\Z$, 则
\[
 (\chi,b)'=0\iff b \text{ 是来自 } L^\times \text{ 的范}.
\]

\paragraph*{步骤 2.} 设 $K$ 为一个域 (不必是局部域). 若 $K$ 的特征为 $p\neq 0$, 则我们要求 $n$ 与 $p$ 互素. 设 $K^{\mathrm{al}}$ 为 $K$ 的一个可分代数闭包 (于是当 $K$ 完美时它就是一个代数闭包), 并设 $G=\Gal(K^{\mathrm{al}}/K)$. 设 $\mu_n$ 为 $K^{\mathrm{al}}$ 中 $n$ 次单位根构成的群, 视为 $G$-模. 在我们对 $n$ 的假设下, $\mu_n$ 是一个 $n$ 阶循环群.

从
\[
 0\to\mu_n\to K^{\mathrm{al}\times}\xrightarrow{x\mapsto x^n}K^{\mathrm{al}\times}\to 0
\]
的上同调序列以及 Hilbert 定理 90 (II, \ref{prop:ii.1.22}), 我们得到
\[
 \delta\colon K^\times/K^{\times n}\xrightarrow{\simeq}H^1(G,\mu_n),\qquad H^2(G,\mu_n)\xrightarrow{\simeq}H^2(G,K^{\mathrm{al}\times})_n.
\]
这里 $H^2(G,K^{\mathrm{al}\times})_n$ 表示 $H^2(G,K^{\mathrm{al}\times})$ 中被 $n$ 杀死的元素构成的子群.

考虑杯积配对
\[
\begin{tikzcd}[column sep=small]
	{H^1(G,\Z/n\Z)} & {\times} & {H^1(G,\mu_n)} & {\longrightarrow} & {H^2(G,\mu_n)} \\
	{\Hom(G,\Z/n\Z)} & & {K^\times/K^{\times n}} & & {H^2(G,K^{\mathrm{al}\times})_n}
	\arrow[equal, no head, from=1-1, to=2-1]
	\arrow[equal, no head, from=1-3, to=2-3]
	\arrow[equal, no head, from=1-5, to=2-5]
\end{tikzcd}
\]
对 $\chi\in\Hom(G,\Z/n\Z)$ 与 $b\in K^\times/K^{\times n}$, 我们把 $(\chi,b)$ 记作该偶对在上述配对下的像: $(\chi,b)=\chi\cup\delta b$ (作为 $H^2(G,K^{\mathrm{al}\times})$ 中的元素).

设 $\chi\in\Hom(G,\Z/n\Z)$, 并 (为简单起见) 假设 $\chi$ 的阶为 $n$. 设 $L_\chi$ 为 $K^{\mathrm{al}}$ 中被 $\Ker(\chi)$ 固定的子域. 于是 $L_\chi$ 是 $K$ 的 $n$ 次循环扩张, 且 $\chi$ 诱导出其 Galois 群 $\bar G$ 与 $\Z/n\Z$ 的一个同构. 设 $b\in K^\times$, 并令
\[
 (\chi,b)'=\delta\chi\cup b\in H^2(\bar G,L^\times),
\]
如同步骤 1. 那么
\[
 \Inf((\chi,b)')=(\chi,b).
\]
由于膨胀映射 $H^2(\bar G,L_\chi)\to H^2(G,K^{\mathrm{al}\times})$ 是单射 (见 II, \ref{ex:ii.1.36}), 我们推出
\[
 (\chi,b)=0\iff b \text{ 是来自 } L_\chi^\times \text{ 的范}.
\]

\paragraph*{步骤 3.} 现在假设 $K$ 是局部域, 于是我们有一个典范同构
\[
 \inv_K\colon H^2(G,K^{\mathrm{al}\times})\to\Q/\Z.
\]
对 $\chi\in\Hom(G,\Z/n\Z)\simeq H^1(G,\Z/n\Z)$ 与 $b\in K^\times/K^{\times n}\simeq H^1(G,\mu_n)$, 我们现在定义
\[
 (\chi,b)=\inv_K(\chi\cup\delta b)\in\tfrac{1}{n}\Z/\Z.
\]
于是我们有一个典范配对
\[
 \chi,b\mapsto(\chi,b)\colon H^1(G,\Z/n\Z)\times H^1(G,\mu_n)\to H^2(G,\mu_n)\simeq\tfrac{1}{n}\Z/\Z.
\]

\begin{proposition}\label{prop:iii.4.3}
这个配对的左核为零.
\end{proposition}

\begin{proof}
设 $\chi\in\Hom(G,\Z/n\Z)$; 我们必须证明: 若对所有 $b\in K^\times/K^{\times n}$ 都有 $(\chi,b)=0$, 则 $\chi=0$.

回顾局部 Artin 映射 $\phi_K\colon K^\times\to\Gal(K^{\mathrm{ab}}/K)$ 与
\[
 \tau\mapsto\tau|L\colon \Gal(K^{\mathrm{ab}}/K)\to\Gal(L/K)
\]
的复合对每个包含在 $K^{\mathrm{ab}}$ 中的 $K$ 的有限扩张 $L$ 都是满射. 这蕴含 $\phi_K$ 的像在 $\Gal(K^{\mathrm{ab}}/K)$ 中稠密. 回顾 (\ref{prop:iii.3.6}): 对 $b\in K^\times$,
\[
 (\chi,b\mod K^{\times n})=\chi(\phi_K(b)).
\]
因为 $\chi$ 连续, 若它在 $\phi_K(K^\times)$ 上为零, 则它在整个 $\Gal(K^{\mathrm{ab}}/K)$ 上为零.
\end{proof}

\paragraph*{步骤 4.} 现在假设 $K$ 含有一个 $n$ 次本原单位根. 于是作为 $G$-模有 $\mu_n\approx\Z/n\Z$ (该同构依赖于一个 $n$ 次本原单位根的选取). 特别地, $\mu_n\otimes\mu_n\approx\mu_n$, 从而 $H^2(G,\mu_n\otimes\mu_n)\approx H^2(G,\mu_n)$ (非典范). 存在一个典范同构
\[
 H^2(G,\mu_n)\otimes\mu_n\xrightarrow{\simeq}H^2(G,\mu_n\otimes\mu_n),
\]
它可以定义为 $x,y\mapsto x\cup y$. 把它与由不变映射给出的同构 $H^2(G,\mu_n)\to\tfrac{1}{n}\Z/\Z$ 复合, 我们得到一个同构
\[
 H^2(G,\mu_n\otimes\mu_n)\simeq\mu_n.
\]
于是我们有一个杯积配对
\[
\begin{tikzcd}[column sep=small]
	{H^1(G,\mu_n)} & {\times} & {H^1(G,\mu_n)} & {\longrightarrow} & {H^2(G,\mu_n\otimes\mu_n)} \\
	{K^\times/K^{\times n}} & & {K^\times/K^{\times n}} & & {\mu_n}
	\arrow[equal, no head, from=1-1, to=2-1]
	\arrow[equal, no head, from=1-3, to=2-3]
	\arrow[equal, no head, from=1-5, to=2-5]
\end{tikzcd}
\]
对 $a,b\in K^\times$, 它们在这个配对下于 $\mu_n$ 中的像记作 $(a,b)$. 这个配对称为 \emph{Hilbert 符号}.

\begin{theorem}\label{thm:iii.4.4}
Hilbert 符号具有下列性质.
\begin{enumerate}[(a)]
\item 它是双乘性的, 即
\[
 (aa',b)=(a,b)(a',b),\quad (a,bb')=(a,b)(a,b').
\]
\item 它是斜对称的, 即
\[
 (b,a)=(a,b)^{-1}.
\]
\item 它是非退化的, 即
\[
 (a,b)=1 \text{ 对所有 } b\in K^\times/K^{\times n}\implies a\in K^{\times n},
\]
\[
 (a,b)=1 \text{ 对所有 } a\in K^\times/K^{\times n}\implies b\in K^{\times n}.
\]
\item $(a,b)=1$ 当且仅当 $b$ 是来自 $K[\sqrt[n]{a}]$ 的范.
\end{enumerate}
\end{theorem}

\begin{proof}
陈述 (a) 与 (b) 可由 Hilbert 符号用杯积的定义直接得出.

接下来证明 (c). 选取一个 $n$ 次单位根, 并用它来定义同构 $\Z/n\Z\to\mu_n$ 与 $\mu_n\to\mu_n\otimes\mu_n$. 由这些映射定义的图表
\[
\begin{tikzcd}[column sep=small, row sep=large]
	{H^1(G,\mu_n)} & {\times} & {H^1(G,\mu_n)} & {\longrightarrow} & {H^2(G,\mu_n\otimes\mu_n)} & {\simeq} & {H^2(G,\mu_n)\otimes\mu_n} \\
	{H^1(G,\Z/n\Z)} & {\times} & {H^1(G,\mu_n)} & {\longrightarrow} & {H^2(G,\mu_n)}
	\arrow[from=2-1, to=1-1]
	\arrow[equal, no head, from=2-3, to=1-3]
	\arrow[from=2-5, to=1-5]
\end{tikzcd}
\]
交换. 由命题\ref{prop:iii.4.3}, 下方配对的左核为零. 因此上方配对的左核也为零, 而由斜对称性, 这蕴含它的右核为零.

最后证明 (d). 假设 (为简单起见) $K[\sqrt[n]{a}]$ 在 $K$ 上的次数为 $n$. 取一个 $n$ 次本原单位根 $\zeta$. 那么 $K[\sqrt[n]{a}]$ 在 $K$ 上 $n$ 次循环, 其 Galois 群由把 $\sqrt[n]{a}$ 送到 $\zeta\sqrt[n]{a}$ 的映射 $\sigma$ 生成. 取定这个 $\zeta$ 之后, $a\in K^\times/K^{\times n}$ 对应于把 $\sigma$ 送到 $1\mod n\Z$ 的元素 $\chi_a\in\Hom(G,\Z/n\Z)$. 由于 $K[\sqrt[n]{a}]=L_{\chi_a}$, (d) 由上面的公式推出.
\end{proof}

\begin{proof}[命题\ref{prop:iii.4.1} 的证明]
我们可以像 $K=\Q_p$, $n=2$ 的情形那样, 从定理\ref{thm:iii.4.4} 推出命题\ref{prop:iii.4.1}: 若 $b$ 对所有 $a$ 都是来自 $K[\sqrt[n]{a}]$ 的范, 则 (由 \ref{thm:iii.4.4}d) 对所有 $a$ 有 $(a,b)=1$, 从而 (由 \ref{thm:iii.4.4}c) $b\in K^{\times n}$.
\end{proof}

\begin{remark}\label{rem:iii.4.5}
Hilbert 符号通过公式
\[
 \phi_K(b)(a^{\frac{1}{n}})=(a,b)a^{\frac{1}{n}}
\]
与局部 Artin 映射相联系. 注意 Galois 理论告诉我们, 对每个 $\tau\in\Gal(K[a^{\frac{1}{n}}]/K)$, 都有 $\tau a^{\frac{1}{n}}=\zeta a^{\frac{1}{n}}$, 其中 $\zeta$ 是某个 $n$ 次单位根 (记住, 我们假设 $K$ 含有 $n$ 次单位根), 所以这个公式的要点在于单位根是相同的. 该公式的证明是一个从命题\ref{prop:iii.3.6} 出发的杯积练习.
\end{remark}

\begin{remark}\label{rem:iii.4.6}
存在定理将表明 $\phi_K$ 定义了一个同构 $K^\times/K^{\times n}\to\Gal(L/K)$, 其中 $L/K$ 是 $K$ 的指数为 $n$ 的最大 Abel 扩张 (没有存在定理时我们只知道这个映射是满射). 显然 $\Hom(G,\Z/n\Z)=\Hom(\Gal(L/K),\Z/n\Z)$, 后者是 $\Gal(L/K)$ 的对偶. 于是 $K^\times/K^{\times n}$ 与 $\Hom(G,\Z/n\Z)$ 有相同的阶. 这表明步骤 2 中的配对
\[
 H^1(G,\Z/n\Z)\times H^1(G,\mu_n)\to H^2(G,\mu_n)\simeq\tfrac{1}{n}\Z/\Z
\]
是非退化的. 更一般地, 对每个有限 $G$-模 $M$, 杯积配对
\[
 H^r(G,M)\times H^{2-r}(G,\check M)\to H^2(G,\mu_n)\simeq\tfrac{1}{n}\Z/\Z
\]
都是非退化的 —— 这里 $\check M=\Hom(M,\mu_n)$. 这个对偶定理由 Poitou 与 Tate 在 1950 年代早期证明, 它是我的书 ``Arithmetic Duality Theorems'' 中那些定理的一个出发点. 注意它蕴含 $H^2(G,\mu_n\otimes\mu_n)\simeq(\mu_n)_G$.
\end{remark}

\begin{remark}\label{rem:iii.4.7}
Hilbert 符号可以用中心单代数来解释 (见下一章). 对 $a,b\in K^\times$ 以及 $n$ 次本原单位根 $\zeta$, 定义 $A(a,b;\zeta)$ 为以 $i$, $j$ 为生成元、满足关系 $i^n=a$, $j^n=b$, $ij=\zeta ji$ 的 $K$-代数. 那么 $(a,b)$ 就是 $A(a,b;\zeta)$ 所代表的 $\mathrm{Br}(K)_n$ 中的类. 注意当 $n=2$ 时我们重新得到四元数代数 $H(a,b)$.
\end{remark}

\begin{remark}\label{rem:iii.4.8}
Hilbert 符号
\[
 (\ ,\ )\colon K^\times/K^{\times n}\times K^\times/K^{\times n}\to\mu_n
\]
对 $K=\R$ 与 $n=2$ 也有定义. 此外, 乘积公式 (34) 对所有数域与所有 $n$ 都成立 (见下文 V, \ref{ex:v.5.4}).
\end{remark}

% ==================================================================
\section{存在定理}\label{sec:III.5}\index{cunzaidingli@存在定理(existence theorem)}

设 $K$ 为局部域. 回顾 $K^\times$ 的子群 $N$ 称为一个\emph{范群}, 如果存在有限 Abel 扩张 $L/K$ 使得 $\Nm_{L/K}(L^\times)=N$. 由于 $K^\times/N\xrightarrow{\simeq}\Gal(L/K)$, 这样的群 $N$ 在 $K^\times$ 中指数有限, 从而是开的 (I, \ref{lem:1.3}).

\begin{theorem}[存在定理]\label{thm:iii.5.1}
$K^\times$ 的每个有限指数的开子群都是范群.
\end{theorem}

在这份草稿中, 我只对特征零的 $K$ 证明这一点. 此时有限指数的子群是开的, 于是我们可以使用 Hilbert 符号的理论.

\begin{remark}\label{rem:iii.5.2}
在开始定理\ref{thm:iii.5.1} 的证明之前, 我们应当考虑一下抽象的情形: 设 $Z$ 为一个无限 Abel 群, 并设 $\mathcal{N}$ 为一族有限指数的子群; 为证明 $\mathcal{N}$ 包含 $Z$ 的所有有限指数子群, 我们需要知道些什么? 显然, 一个起点是要知道:

(a) 若 $I$ 包含 $\mathcal{N}$ 的一个元素, 则 $I$ 在 $\mathcal{N}$ 中;

(b) 若 $N_1$ 与 $N_2$ 在 $\mathcal{N}$ 中, 则 $N_1\cap N_2$ 包含 $\mathcal{N}$ 的一个元素 (从而在 $\mathcal{N}$ 中). 我们知道 $K^\times$ 的范子群族具有这些性质 (推论\ref{prop:iii.1.2}), 但是例如 $\Z$ 中包含 $n\Z$ 的子群构成的集合也具有这些性质. 显然我们需要知道更多. 设 $D$ 为 $\mathcal{N}$ 中所有群的交. 若 $Z$ 的每个有限指数子群 $I$ 都包含 $\mathcal{N}$ 的一个元素, 则它包含 $D$. 反过来, 若 $D$ 是可除的, 即对所有整数 $n$ 有 $D=nD$, 那么 $D$ 包含在每个有限指数的子群中 (指数为 $n$ 的 $Z$ 的子群包含 $nZ$, 而 $nZ$ 包含 $nD=D$). 因此我们需要 $D$ 可除, 但即使这样也还不够. 例如, 取 $Z=\Z$ 与 $\mathcal{N}=\{p^n\Z\mid n\in\N\}$. 那么 $D=\bigcap p^n\Z=0$ 是可除的, 但并非 $\Z$ 的每个有限指数子群都包含 $\mathcal{N}$ 的某个元素. [也许我应当列出一族 $\mathcal{N}$ 为包含所有有限指数的开子群所应满足的完整 ``公理'' 清单 —— 见 \cite{Serre1962}, XI.5, 或 Artin 与 Tate \citeyear{ArtinTate1961}, 第 14 章. 这对我们处理整体存在定理时或许有用.]
\end{remark}

\subsubsection*{定理\ref{thm:iii.5.1} 的证明}

\paragraph*{步骤 1.} 对所有有限扩张 $L/K$, 范映射 $L^\times\to K^\times$ 有闭的像与紧的核.

\begin{proof}
由于像的指数有限, 它是开的 (见 I, \ref{lem:1.3}), 从而是闭的. 核是闭的, 且等式
\[
 \ord_K(\Nm(a))=f\cdot\ord_L(a)
\]
表明它包含在 $U_L$ 中, 而 $U_L$ 是紧的.
\end{proof}

令
\[
 D_K=\bigcap \Nm_{L/K}(L^\times),
\]
其中 $L$ 跑遍 $K$ 的有限扩张 (或者只跑遍有限 Abel 扩张 —— 由定理\ref{thm:iii.3.5}, 这是同一个群).

\paragraph*{步骤 2.} 对任意有限扩张 $K'/K$, 有 $\Nm_{K'/K}D_{K'}=D_K$.

\begin{proof}
设 $a\in D_K$, 并考虑诸集合
\[
 \Nm_{L/K'}L^\times\cap\Nm_{K'/K}^{-1}(a)
\]
其中 $L/K'$ 有限. 它们是紧的且非空的, 并且任意两个都包含第三个. 因此它们的交非空 (点集拓扑). 交中的一个元素落在 $D_{K'}$ 中且范为 $a$.
\end{proof}

\paragraph*{步骤 3.} 群 $D_K$ 是可除的.

\begin{proof}
设 $n>1$ 为一个整数 —— 我们必须证明 $nD_K=D_K$. 设 $a\in D_K$. 对每个含有 $n$ 次本原单位根的 $K$ 的有限扩张 $L$, 考虑集合
\[
 E(L)=\{b\in K^\times\mid b^n=a,\ b\in\Nm_{L/K}L^\times\}.
\]
它是非空的: 对某个 $a'\in D_L$ 有 $a=\Nm_{L/K}a'$; 依命题\ref{prop:iii.4.1}, 对某个 $c\in L$ 有 $a'=c^n$, 且
\[
 \Nm_{L/K}(c)^n=\Nm_{L/K}(a')=a;
\]
因此 $b\stackrel{\mathrm{def}}{=}\Nm_{L/K}(c)\in E(L)$. 此外 $E(L\cdot L')\subset E(L)\cap E(L')$. 由于每个集合都是有限的, 它们的交非空. 交中的一个元素落在 $D_K$ 中且 $n$ 次幂为 $a$.
\end{proof}

\paragraph*{步骤 4.} $D_K=\{1\}$. [实际上, 我们这里不用这一点.]

\begin{proof}
取 $K$ 的一个素元. 令 $V_{m,n}=U^{(m)}\times\pi^{n\Z}$. 那么 $V_{m,n}$ 是 $K^\times$ 的一个 (开的) 有限指数子群, 因此包含 $D_K$. 现在 $D_K\subset\bigcap_{m,n}V_{m,n}=\{1\}$.
\end{proof}

\paragraph*{步骤 5.} $K^\times$ 中每个包含 $U_K$ 的有限指数子群 $I$ 都是范子群.

\begin{proof}
由于 $\ord_K\colon K^\times\to\Z$ 是满射且核为 $U_K$, 所说的那些子群 $I$ 就是 $K^\times$ 中形如 $\ord_K^{-1}(n\Z)$ ($n\ge 1$) 的子群. 设 $K_n$ 为 $K$ 的 $n$ 次非分歧扩张. 那么 $\Nm_{K_n/K}(K_n^\times)$ 是 $K^\times$ 中包含 $U_K$ (见 \ref{prop:iii.1.2}) 且在 $\Z$ 中的像为 $n\Z$ 的子群. 结论成立.
\end{proof}

我们现在来证明定理. 设 $\mathcal{N}$ 为 $K^\times$ 中范群构成的集合, 于是 $D_K=\bigcap_{N\in\mathcal{N}}N$. 设 $I$ 为 $K^\times$ 的一个有限指数子群. 因为 $D_K$ 可除, 有 $I\supset D_K$, 于是 $I\supset\bigcap_{N\in\mathcal{N}}N\supset\bigcap_{N\in\mathcal{N}}(N\cap U_K)$. 因此诸集合 $(N\cap U_K)\setminus I$ 的交为空. 由于它们是紧的, 某个有限子族的交已经为空, 即存在 $\mathcal{N}$ 的某个有限子集 $\mathcal{S}$ 使得 $I\supset\bigcap_{N\in\mathcal{S}}(N\cap U_K)$. 由于诸集合 $N\cap U_K$ 中任意两个都包含第三个, 这蕴含对某个 $N$ 有 $I\supset N\cap U_K$.

取定一个满足 $N\cap U_K\subset I$ 的范群 $N$. 我断言 $I$ 包含
\[
 N\cap\bigl(U_K\cdot(N\cap I)\bigr).
\]
事实上, 交中的每个元素都有形式 $ab$, 其中 $a\in U_K$, $b\in N\cap I$, $ab\in N$. 最后两条陈述蕴含 $a\in N$. 因此 $a\in N\cap U_K\subset I$, 从而 $ab\in I$.

现在 $N\cap I$ 在 $K^\times$ 中指数有限, 因为 $N$ 与 $I$ 都是如此, 且 $K^\times/(N\cap I)$ 单射入 $(K^\times/N)\times(K^\times/I)$. 因此 $U_K\cdot(N\cap I)$ 是 $K^\times$ 中包含 $U_K$ 的有限指数子群, 从而是一个范群 (步骤 5). 现在 $N\cap\bigl(U_K\cdot(N\cap I)\bigr)$ 作为两个范群的交, 包含一个范群. 因此 $I$ 包含一个范群, 而这蕴含它本身是一个范群.

\begin{NOTES}
由 Krasner 引理 (ANT, 7.60) 推出, 局部域的每个有限 Abel 扩张都来自整体域的某个有限 Abel 扩张的完备化. 在 1930 年代, Hasse 与 F. K. Schmidt\footnote{H. Hasse, \emph{Die Normenresttheorie relative-Abelscher Zahlkörper als Klassenkörper im Kleinen}, J. für Mathematik (Crelle) 162 (1930), 145–154. F. K. Schmidt, \emph{Zur Klassenkörpertheorie im Kleinen}, 同上, 155–168.} 能够从整体类域论的主定理推出局部类域论的主定理.

从现代的视角看, 这似乎是一种奇怪的做法. 在 1940 年代, Chevalley 在他的类域论的代数进路中直接发展了局部类域论, 并用它来构造整体类域论. F. K. Schmidt 也证明了局部类域论可以独立于整体类域论来建立.

当时, 局部 Artin 映射还没有好的描述, 也没有构造局部域的最大 Abel 扩张的明确方法 (当然 $\Q_p$ 除外).

1958 年 Dwork 给出了局部 Artin 映射的一个明确描述, 它复现在 \cite{Serre1962} 中, 但不太令人满意.

1965 年 Lubin 与 Tate 引入了 Lubin-Tate 形式群律, 给出了 $K^{\mathrm{ab}}$ 的一个明确构造以及局部 Artin 映射的一个明确描述. 不过, 他们在证明中用到了局部 Artin 映射的存在性 (我们的定理 I, \ref{thm:1.1}).

1980 年代早期, Lubin, Gold 与 Rosen 各自独立地给出了 $K^{\mathrm{ab}}=K_\pi\cdot K^{\mathrm{un}}$ 的 ``初等'' 证明. Iwasawa 在他的书 (\cite{Iwasawa1986}) 中从 Lubin-Tate 的视角展开了整个局部类域论, 也给出了 Hilbert 符号等的明确公式 (由 de Shalit 与 Wiles 得到).

其他的非上同调进路可以在 \cite{Hazewinkel1975}, \cite{Neukirch1986}, 以及 Fesenko 与 Vostokov \citeyear{FesenkoVostokov1993} 中找到. 非上同调进路的缺点自然在于它们不提供关于局部域上同调群的信息, 而这些群在其他课题 (例如椭圆曲线) 中有重要应用.

关于顺分歧扩张情形中局部 Artin 映射的明确描述, 见 Newton, Rachel, Explicit local reciprocity for tame extensions. Math. Proc. Cambridge Philos. Soc. 152 (2012), no. 3, 425--454.

在本章中, 我在很大程度上遵循了 \cite{Serre1962} 与 \cite{Serre1967a}.
\end{NOTES}
  % III 局部类域论: 上同调
% ============================================================
% chIV.tex — 第四章: Brauer 群
% 原书 Chapter IV, Brauer Groups (pp.111--138)
% 约定:
%   - 定理类环境共享 theorem 计数器, 按 节.序号 编号;
%   - 叙述性编号条目 (1.8, 5.2) 不占定理计数器, 手排加粗编号,
%     其后的 \setcounter{theorem}{N} 使下一条编号与原书一致;
%   - 原书公式编号 (35)--(42) 跨节连续, 用 \tag 硬编码;
%   - 标签前缀 iv.;
%   - 勘误: 例 1.11 末句改引 GT 7.19(a); 例 2.1 删 "in its centre"
%     (依 Milne errata, 正文取修正内容, 脚注注明原书错误).
% ============================================================
\chapter{Brauer 群}\label{ch:IV}\index{bulajunqun@Brauer 群(Brauer group)}
\providecommand{\Br}{\operatorname{Br}}

在本章中, 我定义一个域的 Brauer 群, 并证明它给出上同调群 $H^2(K^{\mathrm{al}}/K)$ 的一个具体解释. 除了澄清类域论之外, Brauer 群还有许多应用, 例如应用于有限群的表示论以及非代数闭域上半单代数群的分类.

纵观全章, $k$ 是一个域, 并且所有 $k$ 上的向量空间都取有限维. \emph{$k$-代数}\index{kdaishu@k-代数(k-algebra)}是指一个环 $A$, 它在其中心中包含 $k$, 且作为 $k$-向量空间有限维. 我们不假定 $A$ 交换; 例如 $A$ 可以是 $k$ 上 $n\times n$ 矩阵环 $M_n(k)$. $k$-代数的 \emph{$k$-子代数}是指包含 $k$ 的子环. $k$-代数的同态 $\varphi\colon A\to B$ 是一个环同态, 且满足 $\varphi(a)=a$ 对所有 $a\in k$. $k$-代数 $A$ 的\emph{反代数}\index{fandaishu@反代数(opposite algebra)} $A^{\mathrm{opp}}$ 是与 $A$ 有相同底集与加法的代数, 但其乘法 $\cdot$ 定义为 $\alpha\cdot\beta=\beta\alpha$. 设 $e_1,\ldots,e_n$ 是 $A$ 作为 $k$-向量空间的一组基, 则
\[
e_{i}e_{j}=\sum_{l}a_{ij}^{l}e_{l}
\]
对某些 $a_{ij}^{l}\in k$ 成立, 这些元素称为 $A$ 相对于基 $(e_i)_i$ 的\emph{结构常数}\index{jiegouchangshu@结构常数(structure constants)}. 选定基之后, 代数 $A$ 由其结构常数唯一确定.

\section{单代数; 半单模}\label{sec:IV.1}

\subsection*{半单模}

本节中 $A$ 是一个 $k$-代数.

所谓 $A$-\emph{模}\index{mo@模(module)}, 是指有限生成的左 $A$-模 $V$. 特别地, 这意味着 $1v=v$ 对所有 $v\in V$. 这样的 $V$ 作为 $k$-向量空间也是有限维的, 因此给出一个 $A$-模就相当于给出一个 (有限维) $k$-向量空间连同一个 $k$-代数同态 $A\to\End_k(V)$, 即 $A$ 在 $V$ 上的一个\emph{表示}\index{biaoshi@表示(representation)}. 若该同态是单射, 即若 $ax=0$ 对所有 $x\in V$ 成立便蕴含 $a=0$, 则称该模是\emph{忠实的}\index{zhongshi@忠实(faithful)}.

$A$-模 $V$ 称为\emph{单的}\index{danmu@单模(simple module)}, 如果它非零且除 $0$ 外没有真 $A$-子模; 称为\emph{半单的}\index{bandanmu@半单模(semisimple module)}, 如果它同构于若干单 $A$-模的直和; 称为\emph{不可分解的}\index{bukefenjiemu@不可分解模(indecomposable module)}, 如果它不能写成两个非零 $A$-模的直和. 于是单模是半单的, 而不可分解模是半单的当且仅当它是单的.\footnote{有些作者用「不可约」与「完全可约」分别指「单」与「半单」; 对他们而言, 「不可约」蕴含「完全可约」.}

\begin{example}\label{ex:iv.1.1}
设 $V=k^2$, 并设 $A=k[\alpha]$, 其中 $\alpha\in M_2(k)$. $V$ 的 $A$-子模就是在 $\alpha$ 下稳定的 $k$-子空间.

若 $\alpha=\begin{pmatrix}a&1\\0&a\end{pmatrix}$, 则 $\left\{\begin{pmatrix}*\\0\end{pmatrix}\right\}$ 是 $V$ 的一个 $A$-子模. 事实上, 它是唯一的非平凡子模, 故 $V$ 不可分解但不半单.

若 $\alpha=\begin{pmatrix}a&0\\0&b\end{pmatrix}$, $a\neq b$, 则在 $\alpha$ 下稳定的直线只有 $L_1\stackrel{\mathrm{def}}{=}\left\{\begin{pmatrix}*\\0\end{pmatrix}\right\}$ 与 $L_2\stackrel{\mathrm{def}}{=}\left\{\begin{pmatrix}0\\*\end{pmatrix}\right\}$. 由于 $V=L_1\oplus L_2$ (作为 $A$-模), 它是半单的.

若 $\alpha=\begin{pmatrix}a&0\\0&a\end{pmatrix}$, 则 $V$ 仍分解为两条直线的直和, 但分解不再唯一.

最后, 若 $A=M_2(k)$, 则 $V$ 是单 $A$-模: 对每个非零的 $v\in V$ 有 $V=M_2(k)\cdot v$.
\end{example}

\begin{theorem}\label{thm:iv.1.2}
每个 $A$-模 $V$ 都有滤列
\[
V=V_0\supset\cdots\supset V_r=0,
\]
其商 $V_i/V_{i+1}$ 是单 $A$-模. 若
\[
V=W_0\supset\cdots\supset W_s=0
\]
是另一个这样的滤列, 则 $r=s$, 且存在 $\{0,\ldots,r-1\}$ 的一个置换 $\sigma$, 使得 $V_i/V_{i+1}\simeq W_{\sigma(i)}/W_{\sigma(i)+1}$ 对所有 $i$ 成立.
\end{theorem}

\begin{proof}
若 $V$ 单, 则 $V\supset 0$ 就是这样一个滤列. 否则 $V$ 包含一个子模 $W$, $V\neq W\neq 0$, 于是可对 $V/W$ 与 $W$ 作同样的论证. 由于 $V$ 是有限维 $k$-向量空间, 此过程在有限步后终止.

唯一性部分可以完全照 Jordan-H\"older 定理 (GT, 6.2) 的证明进行.
\end{proof}

\begin{corollary}\label{cor:iv.1.3}
假设
\[
\begin{aligned}
V&\simeq V_1\oplus\cdots\oplus V_r\\
&\simeq W_1\oplus\cdots\oplus W_s,
\end{aligned}
\]
其中所有 $A$-模 $V_i$ 与 $W_j$ 都是单模. 则 $r=s$, 且存在 $\{1,\ldots,r\}$ 的一个置换 $\sigma$, 使得 $V_i\simeq W_{\sigma(i)}$ 对所有 $i$ 成立.
\end{corollary}

\begin{proof}
这些分解定义了 $A$-模的滤列, 其商分别为 $\{V_1,\ldots,V_r\}$ 与 $\{W_1,\ldots,W_s\}$.
\end{proof}

\begin{proposition}\label{prop:iv.1.4}
设 $V$ 是 $A$-模. 若 $V$ 是单子模之和 (不必是直和), 比如 $V=\sum_{i\in I}S_i$, 则对 $V$ 的任何子模 $W$, 存在 $I$ 的子集 $J$, 使得
\[
V=W\oplus\bigoplus_{i\in J}S_i.
\]
\end{proposition}

\begin{proof}
在「和 $S_J=\sum_{j\in J}S_j$ 是直和且 $W\cap S_J=0$」的 $I$ 的子集中取一个极大的 $J$. 我断言 $W+S_J=V$ (从而 $V$ 是 $W$ 与诸 $S_j$ ($j\in J$) 的直和). 为证此, 只需证每个 $S_i$ 都含于 $W+S_J$. 由于 $S_i$ 单, $S_i\cap(W+S_J)$ 等于 $0$ 或 $S_i$. 在前一种情形, $W\cap(S_i+S_J)=0$, 因为否则存在向量 $w\in W$, $s_i\in S_i$, $s_J\in S_J$, 使 $w=s_i+s_J\neq 0$, 于是 $s_i=w-s_J\neq 0$, 与 $S_i\cap(W+S_J)=0$ 矛盾. 因此必是后一种情形, 即 $S_i\subset W+S_J$.
\end{proof}

\begin{corollary}\label{cor:iv.1.5}
$A$-模 $V$ 的下列条件等价:

\begin{enumerate}[(a)]
\item $V$ 半单;
\item $V$ 是单子模之和;
\item $V$ 的每个子模都有补模.
\end{enumerate}
\end{corollary}

\begin{proof}
显然 (a) 蕴含 (b), 而命题表明 (b) 蕴含 (c). 设 (c) 成立. 若 $V$ 单, 则它半单. 否则它包含一个非零真子模 $W$, 它有补模 $W'$, 即 $V=W\oplus W'$. 若 $W$ 与 $W'$ 单, 则 $V$ 半单. 否则可继续此论证, 由于 $V$ 作为 $k$-向量空间是有限维的, 它在有限步后终止.
\end{proof}

\begin{corollary}\label{cor:iv.1.6}
半单模的和、子模与商模都半单.
\end{corollary}

\begin{proof}
每一个都是单模之和: 对和, 这由 (\ref{cor:iv.1.5}) 推出; 对商, 由 (\ref{prop:iv.1.4}) 推出; 对子模, 由每个子模有补模、从而也是一个商模这一事实推出.
\end{proof}

每个半单 $A$-模 $V$ 可以写成直和
\begin{equation}
V\simeq m_{1}S_{1}\oplus\cdots\oplus m_{r}S_{r}\tag{35}
\end{equation}
其中每个 $S_i$ 单, 且互不同构. 若一个 $A$-模同构于某个单模 $S$ 的若干副本的直和, 则称它为\emph{等型的}\index{dengxingmu@等型模(isotypic module)} (以 $S$ 的同构类为型). 分解 (35) 表明每个半单模 $V$ 是若干互不同型的等型模的直和, 这些直和项称为 $V$ 的\emph{等型分支}\index{dengxingfenzhi@等型分支(isotypic component)}. $V$ 对应于单模 $S$ 的等型分支是 $V$ 中同构于 $S$ 的所有单子模之和. 由这一描述可见, 半单 $A$-模的同态 $V\to V'$ 把 $V$ 的每个等型分支映入 $V'$ 的同型的等型分支.

\begin{proposition}\label{prop:iv.1.7}
设 $V$ 是半单 $A$-模. $V$ 的一个子模在 $V$ 的所有自同态下稳定, 当且仅当它是 $V$ 的若干等型分支之和.
\end{proposition}

\begin{proof}
$V$ 的每个自同态都保持其等型分支, 这证明了充分性. 反之, 设 $W$ 是 $V$ 的一个在 $V$ 的所有自同态下稳定的子模, 并设 $S$ 是 $W$ 的一个单子模. 若 $S'$ 是 $V$ 的同构于 $S$ 的子模, 则 $V$ 的自同态
\[
V\xrightarrow{\ \text{投射}\ }S\xrightarrow{\ \simeq\ }S'\hookrightarrow V
\]
把 $W$ 映入 $W$, 从而 $S'\subset W$. 因此 $W$ 包含 $V$ 的以 $S$ 的同构类为型的等型分支.
\end{proof}

\noindent\phantomsection\hypertarget{item:iv.1.8}{}\textbf{1.8} 令 ${}_A A$ 表示视为左 $A$-模的 $A$. 用 $A$ 的元素 $a$ 作右乘 $x\mapsto xa$ 是 ${}_A A$ 的一个 $A$-线性自同态; 反之, 每个 $A$-线性映射 $\varphi\colon {}_A A\to {}_A A$ 都是这种形式, 其中 $a=\varphi(1)$. 于是
\[
\varphi\mapsto\varphi(1)\colon\End_A({}_A A)\xrightarrow{\ \simeq\ }A\qquad(\text{作为 $k$-向量空间}).
\]
令 $\varphi_a$ 为映射 $x\mapsto xa$, 则
\[
(\varphi_{a}\circ\varphi_{b})(1)\stackrel{\mathrm{def}}{=}\varphi_{a}(\varphi_{b}(1))=\varphi_{a}(b)=ba=\varphi_{ba}(1),
\]
从而
\[
\End_A({}_A A)\simeq A^{\mathrm{opp}}\qquad(\text{作为 $k$-代数}).
\]
更一般地, 若 $V$ 是秩 $n$ 的自由 $A$-模, 则 $V$ 的一组基的选择确定一个同构
\[
\End_A(V)\to M_n(A^{\mathrm{opp}}).
\]

$k$-代数 $A$ 称为\emph{半单的}\index{bandandaishu@半单代数(semisimple algebra)}, 若每个 $A$-模都半单. 由于每个 $A$-模都是若干个 ${}_A A$ 的直和的商, 只需验证 $A$-模 ${}_A A$ 半单.

\setcounter{theorem}{8} % 1.8 为叙述性编号条目, 不占定理计数器; 调至 8 使下一条显示为 1.9

\begin{proposition}\label{prop:iv.1.9}
设 $A$ 是半单 $k$-代数. $A$-模 ${}_A A$ 的等型分支恰是 $A$ 的极小双侧理想. $A$ 的每个双侧理想都是极小双侧理想的直和.
\end{proposition}

\begin{proof}
$A$ 的双侧理想就是 ${}_A A$ 的在用 $A$ 的元素作右乘下稳定的子模, 即 ${}_A A$ 的自同态 (见 \hyperlink{item:iv.1.8}{1.8}) 作用下稳定的子模, 因此它们是 ${}_A A$ 的等型分支之和 (见 \ref{prop:iv.1.7}). 特别地, 极小双侧理想恰是 ${}_A A$ 的等型分支. 第二个陈述由上述讨论显然.
\end{proof}

\subsection*{单 $k$-代数}

$k$-代数 $A$ 称为\emph{单的}\index{dandaishu@单代数(simple algebra)}, 若它除 $0$ 外没有真双侧理想. 我们将频繁使用下述观察:

$k$-代数同态 $f\colon A\to B$ 的核是 $A$ 中不含 $1$ 的理想; 因此若 $A$ 单, 则 $f$ 是单射.

\begin{example}\label{ex:iv.1.10}
若 $k$-代数 $A$ 的每个非零元素 $a$ 都有逆, 即存在 $b$ 使 $ab=1=ba$, 则称 $A$ 为\emph{可除代数}\index{kechudaishu@可除代数(division algebra)}. 于是可除代数满足成为域的全部公理, 只是不交换 (因此有时称为\emph{斜域}). 显然可除代数没有非零真理想 — 左的、右的或双侧的都没有 — 从而是单的.

线性代数的许多内容并不要求域交换. 例如, 通常的论证表明可除代数 $D$ 上有限生成模 $V$ 有基, 且所有基的元素个数相同, 记为 $n$; $n$ 称为 $V$ 的\emph{维数}. 特别地, 所有有限生成 $D$-模都是自由的.
\end{example}

\begin{example}\label{ex:iv.1.11}
设 $D$ 是 $k$ 上的可除代数, 考虑矩阵代数 $M_n(D)$. 对 $A,B\in M_n(D)$, $AB$ 的第 $j$ 列 $(AB)_j$ 等于 $AB_j$, 其中 $B_j$ 是 $B$ 的第 $j$ 列. 因此, 给定矩阵 $B$,
\[
B_{j}=0\Rightarrow(AB)_{j}=0,\qquad
B_{j}\neq0\Rightarrow(AB)_{j}\ \text{任意}.
\]
由此可见, 形如 $L(I)$ 的集合是 $M_n(D)$ 中的左理想, 这里 $I$ 是 $\{1,2,\ldots,n\}$ 的子集, $L(I)$ 是第 $j$ 列 ($j\notin I$) 为零的矩阵之集.\footnote{并非 $M_n(D)$ 中所有左理想都形如 $L(I)$ — 例如, 对列施加一个线性关系即可定义 $M_n(D)$ 的一个左理想.} 此外, 每个 $L(\{j\})$ 是 $M_n(D)$ 中的极小左理想. 例如当 $n=4$ 时,
\[
L(\{1,3\})=\left\{\begin{pmatrix}*&0&*&0\\ *&0&*&0\\ *&0&*&0\\ *&0&*&0\end{pmatrix}\right\},\qquad
L(\{3\})=\left\{\begin{pmatrix}0&0&*&0\\ 0&0&*&0\\ 0&0&*&0\\ 0&0&*&0\end{pmatrix}\right\}
\]
分别为一个左理想与一个极小左理想. 对右理想也有类似的陈述. 由此可知 $M_n(D)$ 的每个非零双侧理想都是整个环, 故 $M_n(D)$ 是单 $k$-代数 (GT, 7.19a).\erratum{原书末句给出一个直接论证; 依 Milne errata 改为引用 GT 7.19(a).}
\end{example}

\begin{example}\label{ex:iv.1.12}
对 $a,b\in k^{\times}$, 设 $H(a,b)$ 是以 $1,i,j,ij$ 为基 (作为 $k$-向量空间)、且乘法由
\[
i^{2}=a,\qquad j^{2}=b,\qquad ij=-ji
\]
确定的 $k$-代数, 称为 $k$ 上的\emph{四元数代数}\index{siyuanshudaishu@四元数代数(quaternion algebra)}. 例如当 $k=\R$ 时, $H(-1,-1)$ 就是通常的四元数代数. 可以证明 (见习题\ref{ex:iv.5.7}) $H(a,b)$ 或是可除代数, 或是同构于 $M_2(k)$. 特别地, 它是单的.
\end{example}

\subsection*{中心化子}

设 $A$ 是 $k$-代数 $B$ 的一个 $k$-子代数. $A$ 在 $B$ 中的\emph{中心化子}\index{zhongxinhuazi@中心化子(centralizer)}是
\[
C_{B}(A)=\{b\in B\mid ba=ab\ \text{对所有 }a\in A\};
\]
它仍是 $B$ 的 $k$-子代数.

\begin{example}\label{ex:iv.1.13}
以下例子中, 中心化子都取在 $M_n(k)$ 中.

(a) 设 $A$ 为 $M_n(k)$ 中纯量矩阵之集, 即 $A=kI_n$. 显然 $C(A)=M_n(k)$.

(b) 设 $A=M_n(k)$. 则 $C(A)$ 是 $M_n(k)$ 的中心, 我们现在来计算它. 设 $e_{ij}$ 是在 $(i,j)$ 位置为 $1$、其余位置为 $0$ 的矩阵, 则
\[
e_{i j}e_{l m}=\left\{\begin{array}{l l}{e_{i m}}&{\text{若 }j=l}\\ {0}&{\text{若 }j\neq l.}\end{array}\right.
\]
设 $\alpha=(a_{ij})\in M_n(k)$, 则 $\alpha=\sum_{i,j}a_{ij}e_{ij}$, 于是 $\alpha e_{lm}=\sum_i a_{il}e_{im}$, $e_{lm}\alpha=\sum_j a_{mj}e_{lj}$. 若 $\alpha$ 在 $M_n(k)$ 的中心中, 则 $\alpha e_{lm}=e_{lm}\alpha$, 从而 $a_{il}=0$ ($i\neq l$), $a_{mj}=0$ ($j\neq m$), $a_{ll}=a_{mm}$. 由此可知 $M_n(k)$ 的中心是 $k$ (等同于纯量矩阵之集).

(c) 设 $A$ 为 $M_n(k)$ 中对角矩阵之集. 此时 $C(A)=A$.

注意在所有三种情形中都有 $C(C(A))=A$.
\end{example}

\begin{theorem}[双中心化子定理]\label{thm:iv.1.14}
设 $A$ 是 $k$-代数, $V$ 是忠实的半单 $A$-模. 则 $C(C(A))=A$ (中心化子取在 $\End_k(V)$ 中).
\end{theorem}

\begin{proof}
令 $D=C(A)$, $B=C(D)$. 显然 $A\subset B$; 取 $v_1,\ldots,v_n$ 为 $V$ 作为 $k$-向量空间的生成元, 由下述引理即得反包含.
\end{proof}

\begin{lemma}\label{lem:iv.1.15}
对任何 $v_1,\ldots,v_n\in V$ 与 $b\in B$, 存在 $a\in A$, 使得
\[
av_{1}=bv_{1},\quad av_{2}=bv_{2},\quad\ldots,\quad av_{n}=bv_{n}.
\]
\end{lemma}

\begin{proof}
先证 $n=1$ 的情形. 注意 $Av_1$ 是 $V$ 的 $A$-子模, 于是 (见 \ref{cor:iv.1.5}) 存在 $V$ 的 $A$-子模 $W$ 使 $V=Av_1\oplus W$. 设 $\pi\colon V\to V$ 为映射 $(av_1,w)\mapsto(av_1,0)$ (到 $Av_1$ 上的投射). 它是 $A$-线性的, 故属于 $D$, 且具有性质: $\pi(v)=v$ 当且仅当 $v\in Av_1$. 现在
\[
\pi(bv_{1})=b(\pi v_{1})=bv_{1},
\]
故 $bv_1\in Av_1$, 即为所求.

再证一般情形. 设 $W$ 是 $V$ 的 $n$ 个副本的直和, $A$ 按对角方式作用, 即
\[
a(v_{1},\ldots,v_{n})=(av_{1},\ldots,av_{n}),\qquad a\in A,\ v_{i}\in V.
\]
则 $W$ 仍是半单 $A$-模 (\ref{cor:iv.1.6}). $A$ 在 $\End_k(W)$ 中的中心化子由矩阵 $(\gamma_{ij})_{1\le i,j\le n}$ 组成, 其中 $\gamma_{ij}\in\End_k(V)$ 且 $\gamma_{ij}\alpha=\alpha\gamma_{ij}$ 对所有 $\alpha\in A$ 成立, 即 $\gamma_{ij}\in D$. 换言之, $A$ 在 $\End_k(W)$ 中的中心化子是 $M_n(D)$. 与例\ref{ex:iv.1.13}(b) 中同样的论证 (利用在第 $ij$ 位置为 $\delta$、其余为 $0$ 的矩阵 $e_{ij}(\delta)$) 表明, $M_n(D)$ 在 $\End_k(W)$ 中的中心化子由对角矩阵
\[
\begin{pmatrix}\beta&0&\cdots&0\\ 0&\beta&\cdots&0\\ \vdots&\vdots&\ddots&\vdots\\ 0&0&\cdots&\beta\end{pmatrix}
\qquad(\beta\in B)
\]
组成. 现在把引理的 $n=1$ 情形应用于 $A$、$W$、$b$ 与向量 $(v_1,\ldots,v_n)$ 即完成证明.
\end{proof}

\subsection*{单 $k$-代数的分类}

\begin{lemma}[Schur 引理]\label{lem:iv.1.16}
单 $A$-模的自同态代数是可除代数.
\end{lemma}

\begin{proof}
设 $\gamma$ 是 $A$-线性映射 $S\to S$. 则 $\Ker(\gamma)$ 是 $S$ 的 $A$-子模, 从而是 $S$ 或 $0$. 前一种情形 $\gamma$ 为零; 后一种情形它是同构, 即它有逆, 且其逆也 $A$-线性.
\end{proof}

\begin{theorem}\label{thm:iv.1.17}
每个单 $k$-代数 $A$ 都同构于 $M_n(D)$, 其中 $n$ 为某正整数, $D$ 为某个可除 $k$-代数.
\end{theorem}

\begin{proof}
选取一个单 $A$-模 $S$, 例如取 $A$ 的任一极小左理想. 则 $A$ 在 $S$ 上忠实作用, 因为 $A\to\End_k(S)$ 的核将是 $A$ 中不含 $1$ 的双侧理想, 从而为 $0$.

设 $D$ 是 $A$ 在 $k$-代数 $\End_k(S)$ ($S\to S$ 的 $k$-线性映射代数) 中的中心化子. 由双中心化子定理 (\ref{thm:iv.1.14}), $D$ 在 $\End_k(S)$ 中的中心化子是 $A$, 即 $A=\End_D(S)$. Schur 引理表明 $D$ 是可除代数. 因此 $S$ 是自由 $D$-模 (\ref{ex:iv.1.10}), 设 $S\simeq D^n$, 于是 $\End_D(S)\simeq M_n(D^{\mathrm{opp}})$ (见 \hyperlink{item:iv.1.8}{1.8}).
\end{proof}

\subsection*{单 $k$-代数上的模}

设 $A$ 是 $k$-代数. ${}_A A$ 的子模就是 $A$ 中的左理想, 而 ${}_A A$ 的单子模就是极小左理想.

\begin{proposition}\label{prop:iv.1.18}
单 $k$-代数是半单的.
\end{proposition}

\begin{proof}
设 $A$ 是单 $k$-代数. 只需证 $A$-模 ${}_A A$ 半单. 由定理\ref{thm:iv.1.17}, 可设 $A=M_n(D)$, $D$ 为可除代数. 我们在 \ref{ex:iv.1.11} 中看到, 集合 $L(\{i\})$ 是 $M_n(D)$ 中的极小左理想, 且 $M_n(D)=L(\{1\})\oplus\cdots\oplus L(\{n\})$ (作为 $M_n(D)$-模). 这表明 $A$ 半单.
\end{proof}

\begin{theorem}\label{thm:iv.1.19}
设 $A$ 是半单 $k$-代数. 下列条件等价:

\begin{enumerate}[(a)]
\item $A$ 单;
\item $A$-模 ${}_A A$ 是等型的;
\item 任何两个单 $A$-模同构.
\end{enumerate}
\end{theorem}

\begin{proof}
(a) 与 (b) 的等价性由命题\ref{prop:iv.1.9} 推出, 而 (c) 显然蕴含 (b). 最后 (b) 蕴含 (c): 若 ${}_A A$ 等型, 则若干个 ${}_A A$ 的直和也等型, 而每个 $A$-模都是这种直和的商.
\end{proof}

\begin{corollary}\label{cor:iv.1.20}
设 $A$ 是单 $k$-代数. $A$ 的任何两个极小左理想作为左 $A$-模同构, 且 $A$ 是它的极小左理想之直和.
\end{corollary}

\begin{proof}
极小左理想是单 $A$-模, 故第一个陈述由定理 (c) 推出; 第二个陈述在 \ref{prop:iv.1.18} 的证明中已证.
\end{proof}

\begin{corollary}\label{cor:iv.1.21}
设 $A$ 是单 $k$-代数, $S$ 是单 $A$-模. 每个 $A$-模都同构于若干个 $S$ 的直和. 在 $k$ 上维数相同的任何两个 $A$-模同构.
\end{corollary}

\begin{proof}
由于 $A$ 半单, 第一个断言由定理 (c) 推出, 第二个断言由第一个推出.
\end{proof}

\begin{corollary}\label{cor:iv.1.22}
定理\ref{thm:iv.1.17} 中的整数 $n$ 由 $A$ 唯一确定, 而 $D$ 在同构意义下唯一确定.
\end{corollary}

\begin{proof}
若 $A\simeq M_n(D)$, 则对任何单 $A$-模 $S$ 有 $D^{\mathrm{opp}}\simeq\End_A(S)$, 而 $n$ 是 $S$ 作为 $D$-向量空间的维数. 由于任何两个单 $A$-模同构 (\ref{thm:iv.1.19}), 这就证明了断言.
\end{proof}

\section{Brauer 群的定义}\label{sec:IV.2}

\subsection*{代数的张量积}

设 $A$ 与 $B$ 为 $k$-代数, $A\otimes_k B$ 是 $A$ 与 $B$ 作为 $k$-向量空间的张量积. $A\otimes_k B$ 上有唯一的 $k$-双线性乘法使得
\[
(a\otimes b)(a^{\prime}\otimes b^{\prime})=aa^{\prime}\otimes bb^{\prime},\qquad a,a^{\prime}\in A,\ b,b^{\prime}\in B.
\]
当把 $k$ 等同于 $k\cdot(1\otimes 1)\subset A\otimes_k B$ 时, $A\otimes_k B$ 成为 $k$-代数. 若 $(e_i)_i$ 与 $(f_j)_j$ 分别是 $A$ 与 $B$ 作为 $k$-向量空间的基, 则 $(e_i\otimes f_j)_{i,j}$ 是 $A\otimes_k B$ 的一组基, 且 $A\otimes_k B$ 的结构常数可由 $A$ 与 $B$ 的结构常数经一个显然的公式得到. 我们将用到张量积的下述性质: 它按意义是交换的, 即对任意两个 $k$-代数 $A$ 与 $B$, 映射 $a\otimes b\mapsto b\otimes a$ 扩充为 $k$-代数同构
\[
A\otimes_{k}B\to B\otimes_{k}A;
\]
且按意义是结合的, 即对任意三个 $k$-代数 $A,B,C$, 映射 $a\otimes(b\otimes c)\mapsto(a\otimes b)\otimes c$ 扩充为 $k$-代数同构
\[
A\otimes_{k}(B\otimes_{k}C)\to(A\otimes_{k}B)\otimes_{k}C.
\]

\begin{example}\label{ex:iv.2.1}
对每个 $k$-代数 $A$,
\[
A\otimes_{k}M_{n}(k)\simeq M_{n}(A).
\]
为看出这一点, 注意: 一个包含子环 $R$ 的环 $B$ 同构于 $M_n(R)$, 当且仅当它承认一组作为左 $R$-模的基 $(e_{ij})_{1\le i,j\le n}$, 满足
\[
e_{i j}e_{l m}=\left\{\begin{array}{l l}{e_{i m}}&{\text{若 }j=l}\\ {0}&{\text{若 }j\neq l,}\end{array}\right.
\]
且每个 $e_{ij}$ 与 $R$ 的元素交换.\erratum{原书还要求 $R$ 含于 $B$ 的中心; 这一条件过强 (它蕴含 $R$ 交换), 故在涉及 Brauer 群的应用中不适用. 依 Milne errata (M. Malcic) 改为: 只要求 $B$ 包含子环 $R$, 且诸 $e_{ij}$ 与 $R$ 的元素交换.}
若 $(e_{ij})$ 是 $M_n(k)$ 的标准基, 则 $(1\otimes e_{ij})$ 是 $A\otimes M_n(k)$ 的具有上述性质的 $A$-基.

更一般地,
\begin{equation}
A\otimes_{k}M_{n}(A^{\prime})\simeq M_{n}(A\otimes_{k}A^{\prime})\tag{36}
\end{equation}
对任意 $k$-代数 $A$ 与 $A^{\prime}$ 成立.
\end{example}

\begin{example}\label{ex:iv.2.2}
对任意 $m,n$, $M_m(k)\otimes M_n(k)\simeq M_{mn}(k)$. 为看出这一点, 注意由前例, $M_m(k)\otimes_k M_n(k)\simeq M_m(M_n(k))$, 而元素为 $n\times n$ 矩阵的 $m\times m$ 矩阵就是一个 $mn\times mn$ 矩阵 (去掉内层括号). 或者, 设 $(e_{ij})$ 与 $(f_{i'j'})$ 分别是 $M_m(k)$ 与 $M_n(k)$ 的标准基, 直接验证 $(e_{ij}\otimes f_{i'j'})$ 具有正确的乘法性质.
\end{example}

\subsection*{张量积中的中心化子}

下一命题表明, 子代数张量积的中心化子是它们中心化子的张量积.

\begin{proposition}\label{prop:iv.2.3}
设 $A$ 与 $A'$ 是 $k$-代数, 各有子代数 $B$ 与 $B'$, 并设 $C(B)$ 与 $C(B')$ 分别是 $B$ 在 $A$ 中、$B'$ 在 $A'$ 中的中心化子. 则 $B\otimes_k B'$ 在 $A\otimes_k A'$ 中的中心化子是 $C(B)\otimes_k C(B')$, 即
\[
C_{A\otimes_{k}A^{\prime}}(B\otimes_{k}B^{\prime})=C_{A}(B)\otimes_{k}C_{A^{\prime}}(B^{\prime}).
\]
\end{proposition}

\begin{proof}
显然 $C(B\otimes_k B')\supset C(B)\otimes_k C(B')$.

设 $(f_i)_i$ 是 $A'$ 作为 $k$-向量空间的一组基, 则 $(1\otimes f_i)_i$ 是 $A\otimes_k A'$ 作为 $A$-模的一组基, 于是 $\alpha\in A\otimes_k A'$ 可唯一地写成 $\alpha=\sum_i\alpha_i\otimes f_i$, $\alpha_i\in A$. 设 $\beta\in B$, 则 $\alpha$ 与 $\beta\otimes 1$ 交换当且仅当 $\beta\alpha_i=\alpha_i\beta$ 对所有 $i$ 成立. 因此 $B\otimes 1$ 在 $A\otimes A'$ 中的中心化子是 $C(B)\otimes A'$. 类似地, $1\otimes B'$ 在 $C(B)\otimes A'$ 中的中心化子是 $C(B)\otimes C(B')$.

显然 $C(B\otimes B')\subset C(B\otimes 1)$, 故 $C(B\otimes B')$ 包含于 $C(B)\otimes A'$, 事实上包含于 $1\otimes B'$ 在 $C(B)\otimes A'$ 中的中心化子, 即 $C(B)\otimes C(B')$ 中. 证毕.
\end{proof}

特别地, 两个 $k$-代数张量积的中心是它们中心的张量积: $Z(A\otimes_k B)=Z(A)\otimes_k Z(B)$.

\begin{corollary}\label{cor:iv.2.4}
单 $k$-代数的中心是域.
\end{corollary}

\begin{proof}
显然可除代数的中心是域, 而 Wedderburn 定理 (\ref{thm:iv.1.17}) 表明每个单 $k$-代数都同构于 $M_n(D)$, $D$ 为某可除代数. 现在 $M_n(D)\simeq M_n(k)\otimes_k D$, 故 $Z(M_n(D))\simeq k\otimes_k Z(D)\simeq Z(D)$.
\end{proof}

若 $k$-代数 $A$ 的中心是 $k$, 则称它为\emph{中心的}\index{zhongxindaishu@中心代数(central algebra)}; 既中心又单的 $k$-代数称为\emph{中心单的}\index{zhongxindandaishu@中心单代数(central simple algebra)}. 该推论表明每个单 $k$-代数都是 $k$ 的某个有限扩张上的中心单代数.

\subsection*{原始元素}

在继续之前, 先回顾一点 Bourbaki《代数》第二版中的线性代数.

设 $V$ 是 $k$-向量空间 (可以无穷维), $(e_i)_{i\in I}$ 是 $V$ 的一组基. 每个 $v\in V$ 可唯一写成 $v=\sum a_ie_i$, 定义
\[
J(v)=\{i\in I\mid a_{i}\neq0\};
\]
它是 $I$ 的有限子集, 且当且仅当 $v=0$ 时为空.

设 $W$ 是 $V$ 的子空间. $W$ 的非零元素 $w$ 称为\emph{原始的}\index{yuanshiyuansu@原始元素(primitive element)} (相对于 $V$ 与基 $(e_i)_{i\in I}$), 如果 $J(w)$ 在诸集合 $J(w')$ ($w'$ 为 $W$ 的非零元素) 中极小, 且在表达式 $w=\sum a_ie_i$ 中至少有一个 $a_i=1$.

\begin{proposition}\label{prop:iv.2.5}
(a) 设 $w$ 是 $W$ 的非零元素, $J(w)$ 极小, 并设 $w'$ 是 $W$ 的另一个非零元素. 则 $J(w')\subset J(w)$ 当且仅当 $w'=cw$ 对某个非零 $c\in k$ 成立, 此时 $J(w')=J(w)$.

(b) $W$ 的原始元素之集张成 $W$.
\end{proposition}

\begin{proof}
(a) 若 $w'=cw$, $c\neq 0$, 则显然 $J(w')=J(w)$. 反之, 设 $w=\sum_{i\in J(w)}a_ie_i$ 与 $w'=\sum_{i\in J(w')}b_ie_i$ 是 $W$ 的非零元素, $J(w')\subset J(w)$. 取 $j\in J(w')$, 则 $J(w-a_jb_j^{-1}w')\subset J(w)\setminus\{j\}$, 故由 $J(w)$ 的极小性得 $w-a_jb_j^{-1}w'=0$.

(b) 设 $w=\sum_{i\in J(w)}a_ie_i\in W$, 需证 $w$ 落在原始元素张成的子空间中. 不妨设 $w\neq 0$. 在满足 $J(w')\subset J(w)$ 的 $W$ 的非零元素 $w'$ 中, 取一个使 $J(w')$ 元素个数最少的; 其某个纯量倍数 $w_0$ 将是原始的, 设 $w_0=\sum_{i\in I(w_0)}b_ie_i$, 其中某个 $b_i=1$.

现在 $w=a_jw_0+(w-a_jw_0)$, 其中 $w-a_jw_0\in W$ 且 $|J(w-a_jw_0)|<|J(w)|$. 若 $w-a_jw_0$ 是原始元素的 $k$-线性组合, 则 $w$ 亦然; 否则重复此过程. 由于 $|J(w)|$ 有限, 过程在有限步后终止.
\end{proof}

本节的结果 (与证明) 不要求 $k$ 交换, 即 $k$ 可以是可除代数.

\subsection*{张量积的单性}

\begin{proposition}\label{prop:iv.2.6}
两个单 $k$-代数的张量积仍是单的, 只要其中至少一个是中心的.
\end{proposition}

\begin{proof}
由 \ref{thm:iv.1.17}, 可设其中一个代数是 $M_n(D)$, 其中 $D$ 是中心为 $k$ 的可除代数. 设 $A$ 是另一个单 $k$-代数. 在下面的引理中我们证明 $A\otimes_k D$ 单. 因此 $A\otimes_k D\simeq M_m(D')$, $D'$ 为某可除代数, 从而
\[
A\otimes_{k}M_{n}(D)\simeq M_{n}(A\otimes_{k}D)\simeq M_{n}(M_{m}(D^{\prime}))\simeq M_{mn}(D^{\prime}),
\]
它是单的.
\end{proof}

\begin{lemma}\label{lem:iv.2.7}
设 $A$ 是 $k$-代数, $D$ 是中心为 $k$ 的可除代数. 则 $A\otimes D$ 中任何双侧理想 $\mathfrak{A}$ 作为左 $D$-模由 $\mathfrak{a}\stackrel{\mathrm{def}}{=}\mathfrak{A}\cap(A\otimes 1)$ 生成.
\end{lemma}

\begin{proof}
用规则
\[
\delta(\alpha\otimes\delta^{\prime})=\alpha\otimes\delta\delta^{\prime},\qquad \alpha\in A,\ \delta,\delta^{\prime}\in D
\]
使 $A\otimes_k D$ 成为左 $D$-模. 则 $A\otimes_k D$ 的理想 $\mathfrak{A}$ 特别地是 $A\otimes_k D$ 的 $D$-子模.

设 $(e_i)$ 是 $A$ 作为 $k$-向量空间的基, 则 $(e_i\otimes 1)$ 是 $A\otimes_k D$ 作为左 $D$-向量空间的基. 设 $\alpha\in\mathfrak{A}$ 是相对于这组基的原始元素, 设
\[
\alpha=\sum_{i\in J(\alpha)}\delta_{i}(e_{i}\otimes1)=\sum_{i\in J(\alpha)}e_{i}\otimes\delta_{i}.
\]
对任何非零 $\delta\in D$, $\alpha\delta\in\mathfrak{A}$, 且 $\alpha\delta=\sum e_i\otimes\delta_i\delta=\sum_{i\in J(\alpha)}(\delta_i\delta)(e_i\otimes 1)$. 特别地 $J(\alpha\delta)=J(\alpha)$, 故 $\alpha\delta=\delta'\alpha$ 对某个 $\delta'\in D$ 成立 (见 \ref{prop:iv.2.5}a). 由于某个 $\delta_i=1$, 这蕴含 $\delta=\delta'$, 从而每个 $\delta_i$ 与每个 $\delta\in D$ 交换. 于是 $\delta_i$ 落在 $D$ 的中心 $k$ 中, 而 $\alpha\in A\otimes 1$. 我们已证明 $\mathfrak{A}$ 的每个原始元素都属于 $A\otimes 1$, 这就完成了证明, 因为 $\mathfrak{A}$ 作为 $D$-模由其原始元素生成 (见 \ref{prop:iv.2.5}b).
\end{proof}

\begin{corollary}\label{cor:iv.2.8}
两个中心单 $k$-代数的张量积仍是中心单的.
\end{corollary}

\begin{proof}
结合命题\ref{prop:iv.2.3} 与命题\ref{prop:iv.2.6}.
\end{proof}

设 $A$ 是 $k$ 上的中心单代数, $V$ 表示视为 $k$-向量空间的 $A$. 左乘使 $V$ 成为左 $A$-模, 右乘使它成为右 $A$-模, 或者说左 $A^{\mathrm{opp}}$-模. 这两个作用把 $A$ 与 $A^{\mathrm{opp}}$ 等同于 $\End_k(V)$ 中相互交换的子代数. 由张量积的泛性质, 得到一个同态
\[
a\otimes a^{\prime}\mapsto(v\mapsto ava^{\prime})\colon A\otimes_{k}A^{\mathrm{opp}}\to\End_{k}(V).
\]
由于 $A\otimes_k A^{\mathrm{opp}}$ 单, 且该同态的核不含 $1$, 它是单射. 比较维数:
\[
[A\otimes_{k}A^{\mathrm{opp}}:k]=[A:k]^{2}=(\dim V)^{2}=[\End_{k}(V):k],
\]
故该同态是同构. 我们证明了下述结果.

\begin{corollary}\label{cor:iv.2.9}
对中心单 $k$-代数 $A$,
\[
A\otimes_{k}A^{\mathrm{opp}}\simeq\End_{k}(A)\simeq M_{n}(k),\qquad n=[A:k].
\]
\end{corollary}

\subsection*{Noether-Skolem 定理}

\begin{theorem}[Skolem--Noether]\label{thm:iv.2.10}
设 $f,g\colon A\to B$ 是从 $k$-代数 $A$ 到 $k$-代数 $B$ 的同态. 若 $A$ 单而 $B$ 中心单, 则存在可逆元素 $b\in B$, 使得 $f(a)=b\cdot g(a)\cdot b^{-1}$ 对所有 $a\in A$ 成立, 即 $f$ 与 $g$ 相差 $B$ 的一个内自同构.
\end{theorem}

\begin{proof}
若 $B=M_n(k)=\End_k(k^n)$, 则这两个同态定义了 $A$ 在 $k^n$ 上的作用 — 用 $V_f$ 与 $V_g$ 分别表示带 $f$ 与 $g$ 所定义的作用的 $k^n$. 据 (\ref{cor:iv.1.21}), 具有相同维数的两个 $A$-模同构, 而同构 $b\colon V_g\to V_f$ 是满足 $f(a)\cdot b=b\cdot g(a)$ (对所有 $a\in A$) 的元素 $b\in M_n(k)$.

一般情形, 考虑同态
\[
f\otimes1,\ g\otimes1\colon A\otimes_{k}B^{\mathrm{opp}}\to B\otimes_{k}B^{\mathrm{opp}}.
\]
由于 $B\otimes_k B^{\mathrm{opp}}$ 是 $k$ 上的矩阵代数, 证明的第一部分表明存在 $b\in B\otimes_k B^{\mathrm{opp}}$, 使得
\[
(f\otimes1)(a\otimes b^{\prime})=b\cdot(g\otimes1)(a\otimes b^{\prime})\cdot b^{-1}
\]
对所有 $a\in A$, $b'\in B^{\mathrm{opp}}$ 成立. 在此式中取 $a=1$, 得 $(1\otimes b^{\prime})=b\cdot(1\otimes b^{\prime})\cdot b^{-1}$ 对所有 $b'\in B^{\mathrm{opp}}$ 成立. 因此 (见 \ref{prop:iv.2.3}), $b\in C_{B\otimes_k B^{\mathrm{opp}}}(k\otimes B^{\mathrm{opp}})=B\otimes_k k$, 即 $b=b_0\otimes 1$, $b_0\in B$. 在式中取 $b'=1$, 得
\[
f(a)\otimes1=(b_{0}\cdot g(a)\cdot b_{0}^{-1})\otimes1
\]
对所有 $a\in A$ 成立, 故 $b_0$ 即所求元素.
\end{proof}

\begin{corollary}\label{cor:iv.2.11}
设 $A$ 是 $k$ 上的中心单代数, $B_1$ 与 $B_2$ 是 $A$ 的单 $k$-子代数. 任何同构 $f\colon B_1\to B_2$ 都由 $A$ 的一个内自同构诱导, 即存在可逆元 $a\in A$, 使 $f(b)=aba^{-1}$ 对所有 $b\in B_1$ 成立.
\end{corollary}

\begin{proof}
这是定理中两个映射为 $B_1\xrightarrow{J}B_2\hookrightarrow A$ 与 $B_1\xrightarrow{\mathrm{id}}B_1\hookrightarrow A$ 的特殊情形.
\end{proof}

\begin{corollary}\label{cor:iv.2.12}
中心单 $k$-代数的所有自同构都是内自同构.

例如, $M_n(k)$ 的自同构群是 $\mathrm{PGL}_n(k)\stackrel{\mathrm{def}}{=}\mathrm{GL}_n(k)/k^{\times}I_n$.
\end{corollary}

\subsection*{Brauer 群的定义}

设 $A$ 与 $B$ 是 $k$ 上的中心单代数. 若对某些 $m,n$ 有 $A\otimes_k M_n(k)\simeq B\otimes_k M_m(k)$, 则称 $A$ 与 $B$ \emph{相似}\index{xiangsi@相似(similar)}, 记作 $A\sim B$. 这是一个等价关系: 它显然自反且对称, 而 \ref{ex:iv.2.2} 表明它传递. 定义 $\Br(k)$ 为 $k$ 上中心单代数的相似类之集, 并记 $[A]$ 为 $A$ 的相似类. 对类 $[A]$ 与 $[B]$, 定义
\[
[A][B]=[A\otimes_{k}B].
\]
这是良定义的 (即若 $A\sim A'$ 且 $B\sim B'$, 则 $A\otimes_k B\sim A'\otimes_k B'$), 而张量积的结合律与交换律表明它结合且交换. 对每个 $n$, $[M_n(k)]$ 是单位元, 又因 $A\otimes_k A^{\mathrm{opp}}\simeq M_n(k)$ (见 \ref{cor:iv.2.9}), $[A]$ 以 $[A^{\mathrm{opp}}]$ 为逆. 因此 $\Br(k)$ 是交换群, 称为 $k$ 的 \emph{Brauer 群}\index{bulajunqun@Brauer 群(Brauer group)}.

\begin{remark}\label{rem:iv.2.13}
(a) Wedderburn 定理 (\ref{thm:iv.1.17}, \ref{cor:iv.1.22}) 表明, $k$ 上每个中心单代数都同构于 $M_n(D)$, 其中 $D$ 是某个中心可除代数, 且 $D$ 由 $A$ (甚至由 $A$ 的相似类) 在同构意义下唯一确定. 因此每个相似类都由一个中心可除代数代表, 而两个中心可除代数代表同一相似类当且仅当它们同构.

(b) 我们应当验证\footnote{我曾亲见 Brauer — 他通常态度温和 — 对无视集合与类之区别的「现代」数学家大发雷霆. 正如他指出, 若无视这一区别, 便会得到矛盾 (Russell 悖论); 而一旦系统中有了矛盾, 便能证明任何命题.}相似类构成一个集合, 而不仅仅是一个类. 对每个 $n>0$, 考虑族 $(a_{ij}^{l})_{1\le i,j,l\le n}$, $a_{ij}^{l}\in k$, 它们是 $k$ 上中心可除代数的结构常数. 显然这些族构成一个集合, 每个族定义一个 $k$ 上的中心可除代数, 而这些可除代数包含 $k$ 的 Brauer 群的一个代表集.
\end{remark}

\begin{example}\label{ex:iv.2.14}
(a) 若 $k$ 代数闭, 则 $\Br(k)=0$. 为证此, 需证 $k$ 上任何中心可除代数 $D$ 都等于 $k$. 设 $\alpha\in D$, $k[\alpha]$ 是 $D$ 中由 $k$ 与 $\alpha$ 生成的子代数, 则 $k[\alpha]$ 是 $k$ 上有限次的交换域 (因为它是 $k$ 上有限次的整环). 于是 $k[\alpha]=k$, $\alpha\in k$. 由 $\alpha$ 任意, 得 $D=k$.

(b) Frobenius 证明了 Hamilton 四元数代数是 $\R$ 上唯一的中心可除代数. 因此 $\R$ 的 Brauer 群是 $2$ 阶循环群, 等于 $\{[\R],[\mathbb{H}]\}$ (见下文 \S 4).

(c) Wedderburn 证明了所有有限可除代数都交换. 因此有限域的 Brauer 群为零 (见下文定理\ref{thm:iv.4.1}).

(d) Hasse 证明了非阿基米德局部域的 Brauer 群典范地同构于 $\Q/\Z$ (见下文 本书 p.\pageref{sec:IV.4}).

(e) Albert、Brauer、Hasse 与 Noether 证明了, 对数域 $K$ 有正合列
\[
0\to\Br(K)\to\bigoplus_{v}\Br(K_{v})\xrightarrow{\ \sum\ }\Q/\Z\to0.
\]
和取遍 $K$ 的全部素点 (包括无穷素点). 这一陈述的深度与类域论主要定理相当 (见第 \ref{ch:VII}、\ref{ch:VIII} 章).
\end{example}

\subsection*{基域扩张}

\begin{proposition}\label{prop:iv.2.15}
设 $A$ 是 $k$ 上的中心单代数, $K$ 是包含 $k$ 的域 (不必在 $k$ 上有限次). 则 $A\otimes_k K$ 是 $K$ 上的中心单代数.
\end{proposition}

\begin{proof}
与命题\ref{prop:iv.2.3} 证明中相同的论证表明 $A\otimes_k K$ 的中心是 $k\otimes_k K=K$ (该论证不要求 $K$ 在 $k$ 上有限次). 此外, 引理\ref{lem:iv.2.7} 的证明未用到 $D$ 在 $k$ 上有限维. 因此当 $A$ 是可除代数时, $A\otimes_k K$ 中每个双侧理想作为 $A$-模由它与 $K$ 的交生成, 从而为 $0$ 或 $A\otimes_k K$. 一般情形, $A\simeq M_n(D)$, 于是
\[
\begin{aligned}
A\otimes_{k}K\simeq M_{n}(D)\otimes_{k}K&\simeq(M_{n}(k)\otimes_{k}D)\otimes_{k}K\\
&\simeq M_{n}(k)\otimes_{k}(D\otimes_{k}K)\\
&\simeq M_{n}(D\otimes_{k}K)\\
&\simeq M_{n}(K)\otimes_{K}(D\otimes_{k}K),
\end{aligned}
\]
它是单的.
\end{proof}

\begin{corollary}\label{cor:iv.2.16}
对 $k$ 上的中心单代数 $A$, $[A:k]$ 是平方数.
\end{corollary}

\begin{proof}
显然 $[A:k]=[A\otimes_k k^{\mathrm{al}}:k^{\mathrm{al}}]$, 而 $A\otimes_k k^{\mathrm{al}}\simeq M_n(k^{\mathrm{al}})$ 对某个 $n$ 成立.
\end{proof}

设 $L$ 是包含 $k$ 的域 (不必有限次), 则
\[
M_{n}(k)\otimes L\simeq M_{n}(L),
\]
且
\[
(A\otimes_{k}L)\otimes_{L}(A^{\prime}\otimes_{k}L)=A\otimes_{k}\bigl(L\otimes_{L}(A^{\prime}\otimes_{k}L)\bigr)=(A\otimes_{k}A^{\prime})\otimes_{k}L.
\]
因此映射 $A\mapsto A\otimes_k L$ 定义一个同态
\[
\Br(k)\to\Br(L).
\]
我们把这个同态的核记作 $\Br(L/k)$ — 它由这样的中心单 $k$-代数 $A$ 的相似类组成: $L$-代数 $A\otimes_k L$ 是矩阵代数.

若中心单代数 $A$ (或其在 $\Br(k)$ 中的类) 满足 $A\otimes_k L$ 是 $L$ 上的矩阵代数, 则称 $A$ 被 $L$ \emph{分裂}\index{fenlieyu@分裂域(splitting field)}, 并称 $L$ 为 $A$ 的\emph{分裂域}. 于是 $\Br(L/k)$ 由被 $L$ 分裂的 $\Br(k)$ 的元素组成.

\begin{proposition}\label{prop:iv.2.17}
对每个域 $k$, $\Br(k)=\bigcup\Br(K/k)$, 其中 $K$ 遍历固定代数闭包 $k^{\mathrm{al}}$ 内所有包含 $k$ 的有限扩张.
\end{proposition}

\begin{proof}
需证 $k$ 上每个中心单代数 $A$ 都被 $k$ 的一个有限扩张分裂. 我们知道 $A\otimes_k k^{\mathrm{al}}\simeq M_n(k^{\mathrm{al}})$, 即存在 $A\otimes_k k^{\mathrm{al}}$ 的基 $(e_{ij})_{1\le i,j\le n}$ 使 $e_{ij}e_{lm}=\delta_{jl}e_{im}$ 对所有 $i,j,l,m$ 成立. 由于 $A\otimes_k k^{\mathrm{al}}=\bigcup_{[K:k]<\infty}A\otimes_k K$, 诸 $e_{ij}$ 都落在某个 $A\otimes_k K$ 中, 从而 $A\otimes_k K\simeq M_n(K)$.
\end{proof}

\section{Brauer 群与上同调}\label{sec:IV.3}

对域的 Galois 扩张 $L/k$, 令 $H^2(L/k)=H^2(\Gal(L/k),L^{\times})$. 我们将证明存在自然同构 $H^2(L/k)\simeq\Br(L/k)$, 但先需考察中心单代数的极大子域.

\subsection*{极大子域}

我们需要双中心化子定理的一个变体, 其中 $M_n(k)$ 被换成中心单代数.

\begin{theorem}\label{thm:iv.3.1}
设 $B$ 是中心单 $k$-代数 $A$ 的单 $k$-子代数. 则 $B$ 在 $A$ 中的中心化子 $C=C(B)$ 单, 且 $B$ 是 $C$ 的中心化子. 此外
\[
[B:k][C:k]=[A:k].
\]
\end{theorem}

\begin{proof}
设 $V$ 表示视为 $k$-向量空间的 $B$. 则 $B$ 与 $B^{\mathrm{opp}}$ 分别通过右乘与左乘作用在 $V$ 上, 且互为中心化子 (见 \hyperlink{item:iv.1.8}{1.8}).

考虑中心单代数 $A\otimes_k\End_k(V)$. 命题\ref{prop:iv.2.3} 表明 $B\otimes 1$ 在此代数中的中心化子是 $C\otimes\End_k(V)$, 而 $1\otimes B$ 的中心化子是 $A\otimes B^{\mathrm{opp}}$. 对 $B$ 嵌入 $A\otimes_k\End_k(V)$ 的两个嵌入 $b\mapsto b\otimes 1$ 与 $b\mapsto 1\otimes b$ 应用 Noether-Skolem 定理 (\ref{thm:iv.2.10}, \ref{cor:iv.2.11}), 得到此 $k$-代数的一个可逆元 $u$, 使 $b\otimes 1=u(1\otimes b)u^{-1}$ 对所有 $b\in B$ 成立. 于是显然
\[
C(B\otimes1)=u\cdot C(1\otimes B)\cdot u^{-1}
\]
(中心化子取在 $A\otimes_k\End_k(V)$ 中), 这表明这些中心化子同构. 因此 $C\otimes_k\End_k(V)$ 单, 因为 $A\otimes_k B^{\mathrm{opp}}$ 单 (见 \ref{prop:iv.2.6}); 而这又蕴含 $C$ 本身单, 因为对 $C$ 的每个理想 $\mathfrak{a}$, $\mathfrak{a}\otimes_k\End_k(V)$ 是 $C\otimes_k\End_k(V)$ 的理想. 由于 $\End_k(V)$ 在 $k$ 上次数为 $[B:k]^2$,
\[
[C\otimes_{k}\End(V):k]=[C:k][B:k]^{2},
\]
且显然
\[
[A\otimes B^{\mathrm{opp}}:k]=[A:k][B:k].
\]
比较这些等式, 得
\[
[A:k]=[B:k][C:k].
\]
若 $B'$ 表示 $C$ 在 $A$ 中的中心化子, 则 $B'\supset B$. 但由上述, $[A:k]=[C:k][B':k]$, 故 $[B:k]=[B':k]$, 从而 $B=B'$.
\end{proof}

\begin{remark}\label{rem:iv.3.2}
当 $A=\End_k(V)$ ($V$ 为 $k$-向量空间) 时, 定理\ref{thm:iv.3.1} 可由定理\ref{thm:iv.1.14} 推出, 因为此时 $V$ 是忠实的半单 $B$-模. 这一观察可用来给出定理的另一个证明, 因为基域经有限扩张后 $A$ 就具有这种形式 (见 \ref{prop:iv.2.17}).
\end{remark}

\begin{corollary}\label{cor:iv.3.3}
若在定理的陈述中 $B$ 的中心是 $k$, 则 $C$ 的中心也是 $k$, 且典范同态 $B\otimes_k C\to A$ 是同构.
\end{corollary}

\begin{proof}
$B$ 与 $C$ 的中心都等于 $B\cap C$, 故 $B$ 中心蕴含 $C$ 中心. 于是 $k$-代数 $B\otimes_k C$ 中心单, 这蕴含 $B\otimes_k C\to A$ 单射. 它是满射, 因为两代数在 $k$ 上维数相同.
\end{proof}

\begin{corollary}\label{cor:iv.3.4}
设 $A$ 是 $k$ 上的中心单代数, $L$ 是 $A$ 中包含 $k$ 的子域. 下列条件等价:

\begin{enumerate}[(a)]
\item $L$ 等于它在 $A$ 中的中心化子;
\item $[A:k]=[L:k]^2$;
\item $L$ 是 $A$ 的极大交换 $k$-子代数.
\end{enumerate}
\end{corollary}

\begin{proof}
(a) $\iff$ (b). 由于 $L$ 交换, 它包含于其中心化子 $C(L)$ 中, 而
\[
[A:k]=[L:k][C(L):k],
\]
故 $C(L)=L$ 当且仅当 $[A:k]=[L:k]^2$.

(b) $\implies$ (c). 设 $L'$ 是包含 $L$ 的极大交换子代数, 则 $L'\subset C(L)$, 于是
\[
[A:k]\ge[L:k][L^{\prime}:k]\ge[L:k]^{2}.
\]
故 (b) 蕴含 $[L':k]=[L:k]$.

(c) $\implies$ (a). 若 $L\neq C(L)$, 则存在 $\gamma\in C(L)\setminus L$, 此时 $L[\gamma]$ 是 $A$ 的交换子代数, 与 $L$ 的极大性矛盾.
\end{proof}

\begin{corollary}\label{cor:iv.3.5}
中心可除 $k$-代数 $D$ 的包含 $k$ 的极大子域恰是那些在 $k$ 上次数为 $\sqrt{[D:k]}$ 的子域.
\end{corollary}

\begin{proof}
$D$ 的任何交换 $k$-子代数是 $k$ 上有限次的整环, 从而是域.
\end{proof}

\begin{corollary}\label{cor:iv.3.6}
设 $A$ 是 $k$ 上的中心单代数. $k$ 上有限次的域 $L$ 分裂 $A$, 当且仅当存在与 $A$ 相似的代数 $B$ 包含 $L$ 且
\begin{equation}
[B:k]=[L:k]^{2}.\tag{37}
\end{equation}
特别地, $A$ 的每个在 $k$ 上次数为 $[A:k]^{1/2}$ 的子域 $L$ 都分裂 $A$.
\end{corollary}

\begin{proof}
设 $L$ 分裂 $A$, 则 $L$ 也分裂 $A^{\mathrm{opp}}$, 设 $A^{\mathrm{opp}}\otimes_k L=\End_L(V)$. 该等式表明 $A^{\mathrm{opp}}\otimes_k L$ 是 $L$ 在 $\End_k(V)$ 中的中心化子, 故 $L$ 是 $A^{\mathrm{opp}}\otimes_k L$ 在 $\End_k(V)$ 中的中心化子 (见 \ref{thm:iv.3.1}). 设 $B$ 是 $A^{\mathrm{opp}}$ 在 $\End_k(V)$ 中的中心化子. 我断言 $B$ 满足所要求的条件. 显然 $B\supset L$, 且推论\ref{cor:iv.3.3} 表明 $B$ 中心单且 $B\otimes_k A^{\mathrm{opp}}\simeq\End_k(V)$. 用 $A$ 对两边作张量积并利用 $A\otimes_k A^{\mathrm{opp}}$ 是矩阵代数, 得 $B\sim A$.

反之, 只需证 $L$ 分裂 $B$. 由于 $L$ 交换, $L=L^{\mathrm{opp}}\subset B^{\mathrm{opp}}$, 又因 $[L:k]=\sqrt{[B:k]}$, $L$ 等于它在 $B^{\mathrm{opp}}$ 中的中心化子. 因此 $1\otimes L$ 在 $B\otimes_k B^{\mathrm{opp}}$ 中的中心化子是 $B\otimes_k L$. 当我们把 $B\otimes_k B^{\mathrm{opp}}$ 等同于 $\End_k(B)$ ($B$ 作为 $k$-向量空间的自同态 — 见 \ref{cor:iv.2.9}) 时, $L$ 的中心化子便等同于 $\End_L(B)$ (作为 $L$-向量空间的自同态). 证毕.
\end{proof}

\begin{corollary}\label{cor:iv.3.7}
设 $D$ 是 $k$ 上次数为 $n^2$ 的中心可除代数, $L$ 是 $k$ 上次数为 $n$ 的域. 则 $L$ 分裂 $D$ 当且仅当 $L$ 能嵌入 $D$ (即存在 $k$-代数同态 $L\to D$).
\end{corollary}

\begin{proof}
若 $L$ 嵌入 $D$, 则上一推论表明它分裂 $D$. 反之, 若 $L$ 分裂 $D$, 则存在包含 $L$、与 $D$ 相似且次数为 $[L:k]^2$ 的 $k$ 上中心单代数 $B$. 但 $B\sim D$ 蕴含 $B\simeq M_m(D)$ 对某个 $m$ 成立 (见注记 \ref{rem:iv.2.13}a), 而次数条件蕴含 $m=1$.
\end{proof}

\begin{proposition}\label{prop:iv.3.8}
$k$ 上的每个中心可除代数都包含一个在 $k$ 上可分的极大子域.
\end{proposition}

\begin{proof}
设 $D$ 是 $k$ 上的中心可除代数, $L$ 是在 $D$ 的可分子域中极大的子域. 若 $L$ 不是 $D$ 的极大子域, 则其中心化子 $D'$ 是 $L$ 上不等于 $L$ 的中心可除代数 (见 \ref{thm:iv.3.1}). 下一个引理表明 $D'$ 包含一个真包含 $L$ 的可分子域, 与 $L$ 的定义矛盾.
\end{proof}

\begin{lemma}\label{lem:iv.3.9}
$k$ 上每个不等于 $k$ 的中心可除代数都包含一个在 $k$ 上可分且不等于 $k$ 的子域.
\end{lemma}

\begin{proof}
设 $D$ 是 $k$ 上的中心可除代数. 取 $D$ 的一组基 $(e_i)_{1\le i\le n^2}$, $e_1=1$, 并设 $x=\sum_i a_ie_i$ 是 $D$ 的元素. 对每个整数 $m$, $x^m=\sum_i P_i(m;a_1,\ldots,a_{n^2})e_i$, 其中诸 $P_i$ 是系数在 $k$ 中、只依赖 $D$ 的结构常数的多项式. 因此, 若 $K$ 是包含 $k$ 的域, 而
\[
x=\sum a_{i}(e_{i}\otimes1),\qquad a_{i}\in K
\]
是 $D\otimes_k K$ 的元素, 则 $x^m=\sum_i P_i(m;a_1,\ldots,a_{n^2})\,e_i\otimes 1$, 具有同样的多项式 $P_i$.

当 $k$ 是完美域时引理已知, 故可设 $k$ 特征 $p\neq 0$ 且无限. 假设对每个 $a\in D$, $k[a]$ 都是 $k$ 上的纯不可分扩张, 则对某个 $r>0$ 有 $a^{p^r}\in k$, 且容易看出存在一个对 $D$ 的所有元素都适用的 $r$. 对这个特定的 $r$ 与 $i\neq 1$, $P_i(p^r;a_1,\ldots,a_{n^2})=0$ 对所有 $a_1,\ldots,a_{n^2}\in k$ 成立, 这蕴含 $P_i(p^r;X_1,\ldots,X_{n^2})$ 是零多项式 (见 FT, 5.18 的证明). 因此对 $D\otimes_k k^{\mathrm{al}}$ 的每个元素 $a$ 都有 $a^{p^r}\in k^{\mathrm{al}}$. 但 $D\otimes_k k^{\mathrm{al}}\simeq M_n(k^{\mathrm{al}})$, 于是
\[
\operatorname{diag}(1,0,\ldots,0)=\operatorname{diag}(1,0,\ldots,0)^{p^{r}}\in k^{\mathrm{al}},
\]
这蕴含 $n=1$.
\end{proof}

\begin{corollary}\label{cor:iv.3.10}
Brauer 群 $\Br(k)=\bigcup\Br(L/k)$, 其中 $L/k$ 取遍 $k$ 的某个固定可分代数闭包中的有限 Galois 扩张.
\end{corollary}

\begin{proof}
该命题表明 $\Br(k)$ 的每个元素都被一个有限可分扩张分裂, 从而被一个有限 Galois 扩张分裂.
\end{proof}

\subsection*{中心单代数与 2-上循环}

取定 $k$ 的一个有限 Galois 扩张 $L$, 令 $G=\Gal(L/k)$. 定义 $\mathcal{A}(L/k)$ 为 $k$ 上包含 $L$ 且次数 $[A:k]=[L:k]^2$ (从而 $L$ 等于它在 $A$ 中的中心化子) 的中心单代数 $A$ 之类.

取定 $A\in\mathcal{A}(L/k)$. 对每个 $\sigma\in G$, Noether-Skolem 定理的推论\ref{cor:iv.2.11} 表明存在元素 $e_{\sigma}\in A$, 使得
\begin{equation}
\sigma a=e_{\sigma}ae_{\sigma}^{-1}\qquad\text{对所有 }a\in L.\tag{38}
\end{equation}
此外, $e_\sigma$ 在 (38) 之下在相差乘以 $L^{\times}$ 中一个元素的意义下确定: 若 $f_\sigma$ 也具有该性质, 则 $f_\sigma^{-1}e_\sigma$ 中心化 $L$. 注意 (38) 可写成
\begin{equation}
e_{\sigma}\cdot a=\sigma a\cdot e_{\sigma}\qquad\text{对所有 }a\in L,\tag{39}
\end{equation}
这表明把 $e_\sigma$ 移过 $a\in L$ 时, $a$ 被 $\sigma a$ 替换. 显然 $e_\sigma e_\tau$ 具有 (38) 对 $\sigma\tau$ 的性质, 于是
\begin{equation}
e_{\sigma}e_{\tau}=\varphi(\sigma,\tau)e_{\sigma\tau}\tag{40}
\end{equation}
对某个 $\varphi(\sigma,\tau)\in L^{\times}$ 成立. 注意
\[
e_{\rho}(e_{\sigma}e_{\tau})=e_{\rho}\bigl(\varphi(\sigma,\tau)e_{\sigma\tau}\bigr)=\rho\varphi(\sigma,\tau)\cdot\varphi(\rho,\sigma\tau)\cdot e_{\rho\sigma\tau}
\]
而
\[
(e_{\rho}e_{\sigma})e_{\tau}=\varphi(\rho,\sigma)e_{\rho\sigma}e_{\tau}=\varphi(\rho,\sigma)\varphi(\rho\sigma,\tau)\cdot e_{\rho\sigma\tau}.
\]
因此结合律蕴含 $\varphi$ 是一个 2-上循环; 若取 $e_1=1$, 它甚至是规范的 2-上循环. 换一组 $e_\sigma$ 的选取得到上同调的 2-上循环, 于是我们有良定义的映射 $A\mapsto\gamma(A)\colon\mathcal{A}(L/k)\to H^2(L/k)$.

\begin{theorem}\label{thm:iv.3.11}
映射 $\gamma\colon\mathcal{A}(L/k)\to H^2(L/k)$ 是满射, 其纤维是同构类.
\end{theorem}

先需要一个引理.

\begin{lemma}\label{lem:iv.3.12}
设 $A\in\mathcal{A}(L/k)$, 并取 $e_\sigma$ 满足 (38). 则集合 $(e_\sigma)_{\sigma\in G}$ 是 $A$ 作为 $L$ 上左向量空间的基.
\end{lemma}

\begin{proof}
注意
\[
\dim_{L}(A)=\frac{\dim_{k}(A)}{\dim_{k}(L)}=n,
\]
故只需证诸 $e_\sigma$ 线性无关. 假设不然, 设 $(e_\sigma)_{\sigma\in J}$ 是一个极大线性无关组. 若 $\tau\notin J$, 则
\[
e_{\tau}=\sum_{\sigma}a_{\sigma}e_{\sigma}
\]
对某些 $a_\sigma\in L$ 成立. 取 $a\in L$. 用两种方式计算 $e_\tau a$:
\[
e_{\tau}a=\tau a\cdot e_{\tau}=\sum_{\sigma\in J}\tau a\cdot a_{\sigma}e_{\sigma},
\]
\[
e_{\tau}a=\sum_{\sigma\in J}a_{\sigma}e_{\sigma}\cdot a=\sum_{\sigma\in J}a_{\sigma}\cdot\sigma a\cdot e_{\sigma},
\]
比较得 $\tau a\cdot a_\sigma=\sigma a\cdot a_\sigma$ 对所有 $\sigma\in J$ 成立. 对至少一个 $\sigma\in J$ 有 $a_\sigma\neq 0$, 此时该方程表明 $\tau=\sigma$, 与 $\tau\notin J$ 矛盾. 因此 $J=G$.
\end{proof}

现在, $A$ 由下列性质唯一确定: $A\supset L$; $(e_\sigma)_{\sigma\in G}$ 是 $A$ 作为 $L$-向量空间的基; $A$ 中的乘法满足方程 (39) 与 (40).

设 $A^{\prime}\in\mathcal{A}(L/k)$ 且 $\gamma(A)=\gamma(A')$. 该条件表明可选取 $A$ 与 $A'$ 的基 $(e_\sigma)$ 与 $(e'_\sigma)$, 使其满足 (39) 与 (40) 且具有同一个 2-上循环 $\varphi$. 映射 $\sum a_\sigma e_\sigma\mapsto\sum a_\sigma e'_\sigma$ 是 $k$-代数同构 $A\to A'$.

反之, 设 $A$ 与 $A'$ 是 $\mathcal{A}(L/k)$ 中同构的元素. Noether-Skolem 定理允许我们选取同构 $f\colon A\to A'$, 使 $f(L)=L$ 且 $f|L$ 为恒等映射. 若 $e_\sigma$ 对 $A$ 满足条件 (39), 则 $f(e_\sigma)$ 对 $A'$ 满足 (39). 在选取 $(e_\sigma)$ 与 $(f(e_\sigma))$ 之下, $A$ 与 $A'$ 定义同一个上循环.

这些论述表明映射 $A\mapsto\gamma(A)$ 定义一个单射
\[
\mathcal{A}(L/k)/\simeq\ \to H^{2}(L/k).
\]
为证该映射满, 我们构造其逆.

设 $\varphi\colon G\times G\to L^{\times}$ 是规范 2-上循环. 定义 $A(\varphi)$ 为以 $(e_\sigma)_{\sigma\in G}$ 为基的 $L$-向量空间, 赋予由 (39) 与 (40) 给出的乘法. 则 $e_1$ 是乘法的单位元, 且上循环条件恰好表明
\[
e_{\rho}(e_{\sigma}e_{\tau})=(e_{\rho}e_{\sigma})e_{\tau}.
\]
由此 $A(\varphi)$ 是 $k$-代数. 我们把 $L$ 等同于 $A(\varphi)$ 的子域 $Le_1$.

\begin{lemma}\label{lem:iv.3.13}
代数 $A(\varphi)$ 在 $k$ 上中心单.
\end{lemma}

\begin{proof}
设 $\alpha=\sum a_\sigma e_\sigma$ 中心化 $L$, 并设 $a\in L$. 比较 $a\alpha=\sum aa_\sigma\cdot e_\sigma$ 与 $\alpha a=\sum a_\sigma(\sigma a)\cdot e_\sigma$, 得 $a_\sigma=0$ ($\sigma\neq 1$), 故 $\alpha=a_1e_1\in L$. 因此 $L$ 在 $A(\varphi)$ 中的中心化子是 $L$.

设 $\alpha$ 在 $A(\varphi)$ 的中心中. 则 $\alpha$ 中心化 $L$, 故 $\alpha\in L$, 设 $\alpha=ae_1$, $a\in L$. 比较 $e_\sigma\cdot\alpha=(\sigma a)e_\sigma$ 与 $\alpha\cdot e_\sigma=ae_\sigma$, 得 $\alpha\in k$. 因此 $A(\varphi)$ 中心.

设 $\mathfrak{A}$ 是 $A(\varphi)$ 中的一个双侧理想; 特别地, 它是 $A(\varphi)$ 的 $L$-子空间. 若 $\mathfrak{A}$ 包含某个 $e_\sigma$, 则 (40) 表明它包含所有 $e_\sigma$, 从而等于 $A(\varphi)$. 设 $\mathfrak{A}\neq 0$, 并设 $\alpha=\sum a_\sigma e_\sigma$ 是 $\mathfrak{A}$ 的一个原始元素, 设 $a_{\sigma_0}=1$. 若 $a_{\sigma_1}\neq 0$, $\sigma_1\neq\sigma_0$, 则对每个 $a\in L$,
\[
(\sigma_{1}a)\cdot\alpha-\alpha\cdot a=\sum_{\sigma}a_{\sigma}(\sigma_{1}a-\sigma a)e_{\sigma}\in\mathfrak{A}.
\]
若选取 $a$ 使 $\sigma_1 a\neq\sigma_0 a$, 则该元素非零但非零系数比 $\alpha$ 少, 与 $\alpha$ 的原始性矛盾. 因此 $\alpha=e_{\sigma_0}$, 从而我们证明了 $\mathfrak{A}=A(\varphi)$.
\end{proof}

设 $\varphi$ 与 $\varphi'$ 是上同调的 2-上循环, 即
\[
a(\sigma)\cdot\sigma a(\tau)\cdot\varphi^{\prime}(\sigma,\tau)=a(\sigma\tau)\cdot\varphi(\sigma,\tau)
\]
对某个映射 $a\colon G\to L^{\times}$ 成立. 立即可验证, 把 $e_\sigma$ 送到 $a(\sigma)e'_\sigma$ 的 $L$-线性映射 $A(\varphi)\to A(\varphi')$ 是 $k$-代数同构. 因此 $\varphi\mapsto A(\varphi)$ 定义一个映射 $H^2(L/k)\to\mathcal{A}(L/k)/\simeq$, 它显然是 $A\mapsto\gamma(A)$ 的逆. 这完成了定理\ref{thm:iv.3.11} 的证明.

代数 $A(\varphi)$ 称为\emph{交叉乘积代数}\index{jiaochachengjidaishu@交叉乘积代数(crossed-product algebra)}. 在群上同调出现之前, 2-上循环 $\varphi\colon G\times G\to L^{\times}$ 曾被称为\emph{因子集} (有时至今仍这样称呼).

\begin{theorem}\label{thm:iv.3.14}
对每个有限 Galois 扩张 $L/k$, 映射 $\varphi\mapsto[A(\varphi)]$ 定义交换群的同构 $H^2(L/k)\to\Br(L/k)$.
\end{theorem}

\begin{proof}
为证该映射双射, 只需 (由定理\ref{thm:iv.3.11} 之后) 证明映射 $A\mapsto[A]\colon\mathcal{A}(L/k)/\simeq\to\Br(L/k)$ 双射.

若 $A$ 与 $A'$ 是 $k$ 上相似的中心单代数, 则 (见 \ref{rem:iv.2.13}) 存在中心可除代数 $D$ 使 $A\sim D\sim A'$, 设 $A\simeq M_n(D)$, $A'\simeq M_{n'}(D)$. 若 $[A:k]=[A':k]$, 则 $n=n'$, 从而 $A\simeq A'$. 这证明了映射 $\mathcal{A}(L/k)/\simeq\to\Br(L/k)$ 单, 而 \ref{cor:iv.3.6} 证明它满.
\end{proof}

\begin{lemma}\label{lem:iv.3.15}
对任意两个 2-上循环 $\varphi$ 与 $\varphi'$, 有 $A(\varphi+\varphi')\sim A(\varphi)\otimes_k A(\varphi')$.
\end{lemma}

\begin{proof}
证明略显繁琐, 因为我们必须把 $A(\varphi)\otimes_k A(\varphi')$ 认作交叉乘积代数上的矩阵代数, 而非交叉乘积代数本身. 这里只给出证明梗概.\footnote{细节见 \cite{Blanchard1972}, pp.\,94--95, 或 \cite{FarbDennis1993}, p.\,126--128.}

设 $A=A(\varphi)$, $B=A(\varphi')$, $C=A(\varphi+\varphi')$. 把 $A$ 与 $B$ 视为左 $L$-模 (用左乘), 定义
\[
V=A\otimes_{L}B.
\]
具体地, $V$ 是 $A\otimes_k B$ 的最大商空间, 使得
\[
la\otimes_{L}b=a\otimes_{L}lb
\]
对所有 $a\in A$, $b\in B$, $l\in L$ 成立.

$k$-向量空间 $V$ 有唯一的右 $A\otimes_k B$-模结构使
\[
(a^{\prime}\otimes_{L}b^{\prime})(a\otimes_{k}b)=a^{\prime}a\otimes_{L}b^{\prime}b\qquad(a^{\prime},a\in A,\ b^{\prime},b\in B),
\]
以及唯一的左 $C$-模结构使
\[
(le_{\sigma}^{\prime\prime})(a\otimes_{L}b)=le_{\sigma}a\otimes_{L}e_{\sigma}^{\prime}b\qquad(l\in L,\ \sigma\in G,\ a\in A,\ b\in B).
\]
这里 $(e_\sigma)$, $(e'_\sigma)$, $(e''_\sigma)$ 分别是 $A=A(\varphi)$, $B=A(\varphi')$, $C=A(\varphi+\varphi')$ 的标准基.

这两个作用交换, 于是 $A\otimes_k B$ 在 $V$ 上的右作用定义 $k$-代数同态
\[
f\colon(A\otimes_{k}B)^{\mathrm{opp}}\to\End_{C}(V).
\]
该同态单, 因为 $A\otimes_k B$ (从而其反代数) 单. 由于 $(A\otimes_k B)^{\mathrm{opp}}$ 与 $\End_C(V)$ 在 $k$ 上的次数都是 $n^4$, 其中 $n=[L:k]$, $f$ 是同构. 如 (\ref{cor:iv.1.21}) 中所证, 同一 $k$-维数的单环上的任意两模同构, 故作为 $C$-模有 $V\simeq C^n$. 从而
\[
\End_{C}(V)\simeq\End_{C}(C^{n})=M_{n}(C^{\mathrm{opp}}),
\]
把此同构与 $f$ 复合, 得 $k$-代数同构
\[
(A\otimes_{k}B)^{\mathrm{opp}}\to M_{n}(C)^{\mathrm{opp}}.
\]
同一映射又可解释为同构
\[
A\otimes_{k}B\to M_{n}(C).
\]
\end{proof}

\begin{corollary}\label{cor:iv.3.16}
对每个可分代数闭包 $k^{\mathrm{al}}$, 存在典范同构 $\Br(k)\to H^2(k^{\mathrm{al}}/k)$.
\end{corollary}

\begin{proof}
对每个域塔 $E\supset L\supset k$ ($E$ 与 $L$ 均在 $k$ 上有限 Galois), 图表
\[
\begin{tikzcd}
H^{2}(L/k)\ar[r,"\Inf"]\ar[d] & H^{2}(E/k)\ar[d]\\
\Br(L/k)\ar[r,hook] & \Br(E/k)
\end{tikzcd}
\]
交换 (竖直映射把 $\varphi$ 送到 $[A(\varphi)]$).\footnote{应当补一个该图表交换的证明. (原书作者自注)} 现在利用
\[
\Br(k)=\bigcup\Br(L/k),\qquad
H^{2}(k^{\mathrm{al}}/k)=\bigcup H^{2}(L/k),
\]
其中两个并都取遍 $k^{\mathrm{al}}$ 中 $k$ 的有限 Galois 扩张 $L$.
\end{proof}

\begin{corollary}\label{cor:iv.3.17}
对每个域 $k$, $\Br(k)$ 是挠群; 且对每个有限扩张 $L/k$, $\Br(L/k)$ 被 $[L:k]$ 零化.
\end{corollary}

\begin{proof}
同样的陈述对上同调群成立.
\end{proof}

\section{特殊域的 Brauer 群}\label{sec:IV.4}

上一节的结果使我们能把第 \ref{ch:III} 章的某些结果解释为关于一个域的 Brauer 群的陈述. 本节将独立于第 \ref{ch:III} 章 (但并非完全独立于第 \ref{ch:II} 章) 导出同样的结果.

\subsection*{有限域}

设 $k$ 是有限域. 我们在 III, \ref{lem:iii.1.4} 中看到, 对 $k$ 的每个有限扩张 $L$ 有 $H^{2}(L/k)=0$, 从而 $\Br(k)=0$. 以下是这一事实更直接的证明.

\begin{theorem}[Wedderburn]\label{thm:iv.4.1}
每个有限可除代数都交换.
\end{theorem}

\begin{proof}
设 $D$ 是中心为 $k$ 的有限可除代数, $[D:k]=n^2$. $D$ 的每个元素都属于某个形如 $k[\alpha]$ 的子域, 从而属于某个极大子域. $D$ 的每个极大子域都有 $q^n$ 个元素, 因此彼此同构, 从而 (由 Noether-Skolem) 互相共轭. 于是对每个极大子域 $L$,
\[
D^{\times}=\bigcup_{\alpha}\alpha L^{\times}\alpha^{-1},
\]
但有限群不能等于一个真子群的全部共轭之并 (因为并的元素个数太少), 故 $D=L$.
\end{proof}

\subsection*{实数域}

设 $G=\Gal(\CC/\R)=\{1,\sigma\}$, 则
\[
H^{2}(\CC/\R)\simeq H_{T}^{0}(G,\CC^{\times})=\R^{\times}/\Nm_{G}(\CC^{\times})=\{\pm1\},
\]
故 $\Br(\CC/\R)$ 是 $2$ 阶循环群. $H^2(\CC/\R)$ 的非零元由 2-上循环 $\varphi\colon G\times G\to\CC^{\times}$ 代表:
\[
\varphi(\rho,\tau)=\left\{\begin{array}{c l}-1&\text{若 }\rho=\sigma=\tau\\ 1&\text{其他情形}\end{array}\right.
\]
设 $\mathbb{H}$ 是 $\R$ 上通常的四元数代数, 则把 $x_\sigma$ 送到 $j$ 的 $\CC$-线性映射 $A(\varphi)\to\mathbb{H}$ 是 $\R$-代数同构. 由此, $\R$ 上每个中心单代数都同构于 $\R$ 上的矩阵代数或 $\mathbb{H}$ 上的矩阵代数.

\subsection*{非阿基米德局部域}

设 $K$ 是非阿基米德局部域. 在第 \ref{ch:III} 章定理\ref{thm:iii.2.1} 中我们定义了同构 $H^2(K^{\mathrm{al}}/K)\simeq\Q/\Z$, 从而有同构 $\Br(K)\simeq\Q/\Z$. 本小节说明如何更直接地构造这个同构.

设 $D$ 是 $K$ 上的中心可除代数, $n^2=[D:K]$. 对 $D$ 中每个包含 $K$ 的子域 $L$, 绝对值 $|\cdot|$ 有到 $L$ 的唯一扩张. 由于 $D$ 的每个元素 $\alpha$ 都属于这样一个子域 (例如 $K[\alpha]$), 绝对值 $|\cdot|$ 有到 $D$ 的唯一扩张. 可以验证 $|\cdot|$ 是通常意义下 $D$ 上的非阿基米德绝对值, 即

(a) $|\alpha|=0\iff\alpha=0$;

(b) 对所有 $\alpha,\beta\in D$, $|\alpha\beta|=|\alpha||\beta|$;

(c) 对所有 $\alpha,\beta\in D$, $|\alpha+\beta|\le\max\{|\alpha|,|\beta|\}$.

设 $q$ 为 $K$ 的剩余类域 $k$ 的元素个数, 并用公式
\[
|\alpha|=(1/q)^{\mathrm{ord}(\alpha)}
\]
对 $\alpha\in D$ 定义 $\mathrm{ord}(\alpha)$. 则 $\mathrm{ord}$ 把 $K$ 上的加性赋值 $\mathrm{ord}_K$ (标准化为映 $K^{\times}$ 到 $\Z$ 上) 扩张到 $D$. 对 $D$ 中每个包含 $K$ 的子域 $L$, $[L:K]\le n$, 故 $\mathrm{ord}(L^{\times})\subset n^{-1}\Z$, 从而也有 $\mathrm{ord}(D^{\times})\subset n^{-1}\Z$.

令
\[
\mathcal{O}_{D}=\{\alpha\in D\mid\mathrm{ord}(\alpha)\ge0\},\qquad
\mathfrak{P}=\{\alpha\in D\mid\mathrm{ord}(\alpha)>0\}.
\]
则 $\mathcal{O}_D$ 是 $D$ 的子环, 称为\emph{整数环}. 对 $D$ 中每个包含 $K$ 的子域 $L$, $\mathcal{O}_D\cap L=\mathcal{O}_L$, 故 $\mathcal{O}_D$ 恰由 $D$ 中在 $\mathcal{O}_K$ 上整的元素组成. 此外 $\mathfrak{P}$ 是 $D$ 的极大双侧理想 (显然), 且 $\mathcal{O}_D$ 的非零双侧理想恰为 $\mathfrak{P}$ 的幂 (证明与交换情形相同). 于是对某个 $e$ 有 $\mathfrak{P}^e=\mathfrak{p}\mathcal{O}_D$. 此时 $\mathrm{ord}(D^{\times})=e^{-1}\Z$, 故 $e\le n$.

显然, $\mathcal{O}_D$ 中不在 $\mathfrak{P}$ 内的元素都是单位. 因此 $d\stackrel{\mathrm{def}}{=}\mathcal{O}_D/\mathfrak{P}$ 仍是可除代数, 从而是域 (\ref{thm:iv.4.1}). 设 $f$ 为其在 $k$ 上的次数, 写 $d=k[a]$. 可把 $a$ 提升为 $\mathcal{O}_D$ 的元素 $\alpha$. 由于 $[K[\alpha]:K]\le n$, 得 $f\le n$.

与交换情形相同的论证给出 $n^2=ef$: 即 $\mathcal{O}_D$ 是秩 $m$ 的自由 $\mathcal{O}_K$-模; 因 $\mathcal{O}_D\otimes_{\mathcal{O}_K}K=D$, 故 $m=n^2$; 又因 $\mathcal{O}_D\otimes_{\mathcal{O}_K}k=\mathcal{O}_D/\mathfrak{p}\mathcal{O}_D$, 它在 $k$ 上也是 $n^2$ 维自由模. 现在考虑 $k$-向量空间的滤列
\[
\mathcal{O}_{D}\supset\mathfrak{P}\supset\mathfrak{P}^{2}\supset\cdots\supset\mathfrak{P}^{e}=\mathfrak{p}\mathcal{O}_{D}.
\]
由 $f$ 的定义, $\mathcal{O}_D/\mathfrak{P}=d$ 作为 $k$-向量空间维数为 $f$, 而相邻商是 $d$ 上的一维向量空间. 故 $\mathcal{O}_D/\mathfrak{p}\mathcal{O}_D$ 在 $k$ 上维数为 $ef$, 从而 $ef=n^2$.

由于 $e\le n$, $f\le n$, 等式 $ef=n^2$ 蕴含 $e=f=n$. 特别地, 每个不等于 $K$ 的中心可除代数都是分歧的. 再写 $d=k[a]$, 并把 $a$ 提升为 $D$ 中的元素 $\alpha$, 则 $K[\alpha]$ 是剩余类域为 $d$ 的域, 故 $[K[\alpha]:K]\ge[d:k]=n$. 因此 $K[\alpha]$ 在 $K$ 上次数为 $n$ 且非分歧. 它是极大子域, 从而分裂 $D$. 我们证明了 $\Br(K)$ 的每个元素都被一个非分歧扩张分裂, 即 $\Br(K)$ 等于其子群 $\Br(K^{\mathrm{un}}/K)$.

下面定义映射
\[
\operatorname{inv}_{K}\colon\Br(K)\to\Q/\Z.
\]
$\Br(K)$ 的一个元素由 $K$ 上的一个中心可除代数 $D$ 代表 (在同构意义下唯一). 由刚才所证, $D$ 有一个在 $K$ 上非分歧的极大子域 $L$. 设 $\sigma$ 是 $L$ 的 Frobenius 自同构. 据 Noether-Skolem 定理, 存在 $\alpha\in D$ 使 $\sigma x=\alpha x\alpha^{-1}$ 对所有 $x\in L$ 成立. 若 $\alpha'$ 也具有此性质, 则 $\alpha'=c\alpha$ 对某个 $c\in L$ 成立, 于是
\[
\mathrm{ord}(\alpha^{\prime})=\mathrm{ord}(c)+\mathrm{ord}(\alpha)\equiv\mathrm{ord}(\alpha)\ \bmod\ \Z.
\]
我们定义
\[
\operatorname{inv}_{K}(D)=\mathrm{ord}(\alpha)\ \bmod\ \Z.
\]
它只依赖 $D$ 的同构类.

\begin{example}\label{ex:iv.4.2}
设 $L$ 是 $K$ 的 $n$ 次非分歧扩张, $\sigma$ 是 $L/K$ 的 Frobenius 自同构, 于是 $G\stackrel{\mathrm{def}}{=}\Gal(L/K)=\{\sigma^i\mid 0\le i\le n-1\}$. 设 $\varphi$ 为 2-上循环
\[
\varphi(\sigma^{i},\sigma^{j})=\left\{\begin{array}{l l}{1}&{\text{若 }i+j\le n-1}\\ {\pi}&{\text{若 }i+j>n-1,}\end{array}\right.
\]
其中 $\pi$ 是 $K$ 的素元 (见 III, \ref{prop:iii.1.9} 之前的讨论). 交叉乘积代数 $A(\varphi)$ 等于 $\bigoplus_{0\le i\le n-1}Le_i$, 其乘法由
\[
e_{i}\cdot a=\sigma^{i}a\cdot e_{i}\qquad(\text{所有 }a\in L)
\]
与
\[
e_{i}e_{j}=\left\{\begin{array}{l l}{e_{i+j}}&{\text{若 }i+j\le n-1}\\ {\pi e_{i+j-n}}&{\text{若 }i+j>n-1}\end{array}\right.
\]
确定. 我们通过把 $e_0$ 等同于 $1$ 把 $L$ 等同于 $A(\varphi)$ 的子域. 由于 $e_1ae_1^{-1}=\sigma a$ 对 $a\in L$ 成立, 可用 $e_1$ 计算 $A(\varphi)$ 的不变量. 由上述规则, $e_1^n=e_{n-1}e_1=\pi e_0=\pi$, 于是
\[
\operatorname{inv}_{K}(A(\varphi))=\mathrm{ord}(e_{1})=\tfrac{1}{n}\,\mathrm{ord}(e_{1}^{n})=\tfrac{1}{n}\,\mathrm{ord}(\pi)=\tfrac{1}{n},
\]
正如所料.
\end{example}

\begin{proposition}\label{prop:iv.4.3}
刚定义的映射 $\operatorname{inv}_K\colon\Br(K)\to\Q/\Z$ 是双射.
\end{proposition}

\begin{proof}
设 $L$ 是 $K$ 的 $n$ 次非分歧扩张 (包含在 $K$ 的某个固定代数闭包 $K^{\mathrm{al}}$ 中), $l/k$ 是相应的剩余类域扩张. 由于范映射 $l\to k$ 与 $l^{\times}\to k^{\times}$ 都是满射, 且 $U_L$ 有一个滤列, 其商为 $l^{\times}$ 或 $l$, 可知范映射 $U_L\to U_K$ 满 (见 III, \ref{prop:iii.1.2}). 因此 $H_T^0(G,U_L)=0$, 又 (由循环群上同调的周期性) 这蕴含 $H^2(G,U_L)=0$. 由于对 $K$ 的任意素元 $\pi$ 有 $L^{\times}=U_L\times\pi^{\Z}$,
\[
H^{2}(L/K)=H^{2}(G,\pi^{\Z}).
\]
考察
\[
0\to\Z\to\Q\to\Q/\Z\to0
\]
的上同调序列可知, $H^2(G,\pi^{\Z})$ 是 $n$ 阶循环群, 由上一例中考虑的上循环 $\varphi$ 的类生成 (见 III, \ref{prop:iii.1.9} 之前的讨论). 因此 $\Br(L/K)$ 是 $n$ 阶循环群, 由 $[A(\varphi)]$ 生成. 由此可知 $\operatorname{inv}_K\colon\Br(K^{\mathrm{un}}/K)\to\Q/\Z$ 是同构, 而前面已看到 $\Br(K^{\mathrm{un}}/K)=\Br(K^{\mathrm{al}}/K)$.
\end{proof}

\begin{remark}\label{rem:iv.4.4}
(a) 例\ref{ex:iv.4.2} 中的计算表明, 本章定义的不变量映射与上一章用上同调定义的一致. 特别地, 这表明本章定义的映射是同态.

(b) 与例\ref{ex:iv.4.2} 同样的计算给出 $\operatorname{inv}_K(A(\varphi^i))=\frac{i}{n}\bmod\Z$. 我断言: 若 $i$ 与 $n$ 互素, 则 $A(\varphi^i)$ 是中心可除代数. 若不然, 则 $A(\varphi^i)\simeq M_r(D)$, $D$ 为某个次数 $m^2$ ($m<n$) 的中心可除代数, 而 $\operatorname{inv}_K(A(\varphi^i))=\operatorname{inv}_K(D)\in\frac{1}{m}\Z/\Z$, 矛盾. 由此可知, $K$ 上每个中心可除代数都同构于恰好一个形如 $A(\varphi^i)$ 的可除代数 (对某个 $n\ge 1$ 及某个与 $n$ 互素的 $i$). 特别地, 对中心可除代数 $D$, $[D]$ 在 $\Br(K)$ 中的阶是 $\sqrt{[D:K]}$.

(c) 设 $D$ 是 $K$ 上次数为 $n^2$ 的中心可除代数. 由于映射 $\Br(K)\to\Br(L)$ 把不变量乘以 $[L:K]$ (见 III, \ref{thm:iii.2.1}), $D$ 被 $K$ 的每个 $n$ 次扩张 $L$ 分裂. 因此每个这样的 $L$ 都能嵌入 $D$ (由 \ref{cor:iv.3.7}). 换言之, $K[X]$ 中每个 $n$ 次不可约多项式在 $D$ 中都有根.
\end{remark}

\begin{NOTES}\label{note:iv.4}
关于有限域、实数域与非阿基米德局部域的结果分别归功于 Wedderburn、Frobenius 与 Hasse. 非阿基米德局部域的 Brauer 群的计算本质上就是 Hasse 原来的 (前上同调的) 证明.
\end{NOTES}

\section{补充}\label{sec:IV.5}

\subsection*{半单代数}

回忆: 若每个 $A$-模都半单, 则称 $k$-代数 $A$ 半单. 单 $k$-代数半单 (\ref{prop:iv.1.18}), 而 Maschke 定理 (GT, 7.4) 表明, 当 $G$ 的阶不被 $k$ 的特征整除时, 群代数 $k[G]$ 半单.

\begin{example}\label{ex:iv.5.1}
设 $A$ 是若干单 $k$-代数的有限乘积. 单因子的每个极小左理想都是 ${}_A A$ 的单 $A$-子模, 故 $A$ 是单 $A$-模之和, 从而半单. 由于每个 $A$-模都是若干个 ${}_A A$ 的直和的商, 这表明 $A$ 半单.
\end{example}

在陈述本节主要结果之前, 先回顾一些初等的模论.

\noindent\textbf{5.2} 设 $A$ 是 $k$-代数, 考虑模
\[
M=M_{1}\oplus\cdots\oplus M_{n},\qquad
N=N_{1}\oplus\cdots\oplus N_{m}.
\]
设 $\alpha$ 是 $A$-线性映射 $M\to N$. 对 $x_j\in M_j$, 设
\[
\alpha(0,\ldots,0,x_{j},0,\ldots,0)=(y_{1},\ldots,y_{m}),
\]
则 $x_j\mapsto y_i$ 是 $A$-线性映射 $M_j\to N_i$, 记作 $\alpha_{ij}$. 于是 $\alpha$ 定义一个 $m\times n$ 矩阵, 其 $ij$ 系数是 $A$-线性映射 $M_j\to N_i$. 反之, 每个这样的矩阵 $(\alpha_{ij})$ 定义一个 $A$-线性映射 $M\to N$, 即
\[
\begin{pmatrix}x_{1}\\ \vdots\\ x_{j}\\ \vdots\\ x_{n}\end{pmatrix}\mapsto
\begin{pmatrix}\alpha_{11}&\cdots&\alpha_{1j}&\cdots&\alpha_{1n}\\ \vdots&&\vdots&&\vdots\\ \alpha_{i1}&\cdots&\alpha_{ij}&\cdots&\alpha_{jn}\\ \vdots&&\vdots&&\vdots\\ \alpha_{m1}&\cdots&\alpha_{mj}&\cdots&\alpha_{mn}\end{pmatrix}
\begin{pmatrix}x_{1}\\ \vdots\\ x_{j}\\ \vdots\\ x_{n}\end{pmatrix}
\stackrel{\mathrm{def}}{=}
\begin{pmatrix}\alpha_{11}(x_{1})+\cdots+\alpha_{1n}(x_{n})\\ \vdots\\ \alpha_{i1}(x_{1})+\cdots+\alpha_{in}(x_{n})\\ \vdots\\ \alpha_{m1}(x_{1})+\cdots+\alpha_{mn}(x_{n})\end{pmatrix}.
\]
于是可见
\begin{equation}
\Hom_{A}(M,N)\simeq\bigl(\Hom_{A}(M_{j},N_{i})\bigr)_{1\le j\le n,\,1\le i\le m}\tag{41}
\end{equation}
($k$-向量空间的同构). 当 $M=N$ 时, 它成为 $k$-代数同构. 例如, 若 $M$ 是 $m$ 个 $M_0$ 副本的直和, 则
\begin{equation}
\End_{A}(M)\simeq M_{m}\bigl(\End_{A}(M_{0})\bigr)\tag{42}
\end{equation}
(系数在环 $\End_A(M_0)$ 中的 $m\times m$ 矩阵).

\setcounter{theorem}{2} % 5.2 为叙述性编号条目, 不占定理计数器; 调至 2 使下一条显示为 5.3

\begin{theorem}\label{thm:iv.5.3}
设 $V$ 是有限维 $k$-向量空间, $A$ 是 $\End_k(V)$ 的 $k$-子代数. 若 $V$ 作为 $A$-模半单, 则 $A$ 在 $\End_k(V)$ 中的中心化子是若干单 $k$-代数的乘积 (从而它是半单 $k$-代数).
\end{theorem}

\begin{proof}
由假设可写 $V\simeq\bigoplus_i r_iS_i$, 其中诸 $S_i$ 是单 $A$-模, 互不同构. $A$ 在 $\End_k(V)$ 中的中心化子是 $\End_A(V)$, 而 $\End_A(V)\simeq\End_A\bigl(\bigoplus_i r_iS_i\bigr)$. 由于 $\Hom_A(S_j,S_i)=0$ ($i\neq j$),
\[
\begin{aligned}
\End_{A}\Bigl(\bigoplus_i r_{i}S_{i}\Bigr)&\simeq\prod_{i}\End_{A}(r_{i}S_{i})\quad\text{由 (41)}\\
&\simeq\prod_{i}M_{r_{i}}(D_{i})\qquad\ \ \text{由 (42)},
\end{aligned}
\]
其中 $D_i=\End_A(S_i)$. 据 Schur 引理 (\ref{lem:iv.1.16}), $D_i$ 是可除代数, 故 $M_{r_i}(D_i)$ 是单 $k$-代数 (\ref{ex:iv.1.11}).
\end{proof}

\begin{theorem}\label{thm:iv.5.4}
每个半单 $k$-代数都同构于若干单 $k$-代数的乘积.
\end{theorem}

\begin{proof}
选取一个 $A$ 忠实作用的 $A$-模 $V$, 例如 $V={}_A A$. 则 $A$ 等于它在 $\End_k(V)$ 中的双中心化子 $C(C(A))$ (见 \ref{thm:iv.1.14}). 据定理\ref{thm:iv.5.3}, $C(A)$ 半单, 故 $C(C(A))$ 是若干单代数的乘积 (\ref{thm:iv.5.3}).
\end{proof}

\subsection*{半单 $k$-代数上的模}

设 $A=B\times C$ 是若干 $k$-代数的乘积. $B$-模 $M$ 通过作用 $(b,c)m=bm$ 成为 $A$-模.

\begin{theorem}\label{thm:iv.5.5}
设 $A$ 是半单 $k$-代数, 设 $A=A_1\times\cdots\times A_t$, 诸 $A_i$ 单. 对每个 $A_i$, 设 $S_i$ 是一个单 $A_i$-模.

(a) 每个 $S_i$ 都是单 $A$-模, 且每个单 $A$-模都恰好同构于其中一个 $S_i$.

(b) 每个 $A$-模都同构于 $\bigoplus r_iS_i$ (某些 $r_i\in\N$), 而两个模 $\bigoplus r_iS_i$ 与 $\bigoplus r'_iS_i$ 同构当且仅当 $r_i=r'_i$ 对所有 $i$ 成立.
\end{theorem}

\begin{proof}
(a) 显然每个 $S_i$ 视为 $A$-模时单, 且互不同构. 由 \ref{cor:iv.1.20}, ${}_A A\simeq\bigoplus r_iS_i$ (某些 $r_i\in\N$). 设 $S$ 是单 $A$-模, $x$ 是 $S$ 的非零元素, 则映射 $a\mapsto ax\colon {}_A A\to S$ 满射, 于是它限制到 ${}_A A$ 的某个 $S_i$ 上非零, 从而是同构.

(b) 第一部分由 (a) 与半单环的定义推出, 第二部分由 \ref{cor:iv.1.3} 推出.
\end{proof}

\begin{proposition}\label{prop:iv.5.6}
设 $A$ 是半单 $k$-代数, $k'$ 是 $k$ 的可分扩张, 则 $A\otimes_k k'$ 半单.
\end{proposition}

\begin{proof}
不妨设 $A$ 单, 则 $A$ 的中心 $F$ 是域. 由于 $k'/k$ 可分, $F\otimes_k k'$ 是若干域 $F_i$ 的乘积. 现在
\[
A\otimes_{k}k^{\prime}\simeq A\otimes_{F}(F\otimes_{k}k^{\prime})\simeq A\otimes_{F}\Bigl(\prod F_{i}\Bigr)\simeq\prod A\otimes_{F}F_{i},
\]
而每个代数 $A\otimes_F F_i$ 单 (\ref{prop:iv.2.15}).
\end{proof}

\subsection*{代数、上同调与群扩张}

设 $A$ 是 $k$ 上次数为 $n^2$ 的中心单代数, 并设 $A$ 包含一个在 $k$ 上次数为 $n$ 的 Galois 域 $L$. 令 $E$ 为满足 $\alpha L\alpha^{-1}=L$ 的可逆元素 $\alpha\in A$ 之集. 每个 $\alpha\in E$ 定义 $\Gal(L/k)$ 的元素 $x\mapsto\alpha x\alpha^{-1}$, 而 Noether-Skolem 定理蕴含 $\Gal(L/k)$ 的每个元素都来自某个 $\alpha\in E$. 由于 $[L:k]=\sqrt{[A:k]}$, $L$ 的中心化子是 $L$ 本身, 于是序列
\[
1\to L^{\times}\to E^{\times}\to\Gal(L/k)\to1
\]
正合. 不难证明, 把 $A$ 送到这个序列的映射定义了从 $\mathcal{A}(L/k)$ 到 $\Gal(L/k)$ 被 $L^{\times}$ 扩张的同构类之集的同构, 从而到 $H^2(L/k)$ 的同构 (见 II, \ref{ex:ii.1.18}). 参见 Serre 1950/51.

\subsection*{Brauer 群与 K-理论}

设 $k$ 是包含 $n$ 次本原单位根 $\zeta$ 的域. 对元素 $a,b\in k^{\times}$, 可以配一个具有生成元 $i,j$ 与关系
\[
i^{n}=a,\qquad j^{n}=b,\qquad ij=\zeta ji
\]
的 $k$-代数 $A(a,b;\zeta)$. 它是 $k$ 上的中心单代数.

域 $F$ 的 \emph{Milnor K-群}\index{milnorkqun@Milnor K-群(Milnor K-group)} $K_2F$ 定义为 $F^{\times}\otimes_{\Z}F^{\times}$ 除以由形如 $u\otimes(1-u)$ 的元素生成的子群, 其中 $u\in F^{\times}$ 且 $1-u\in F^{\times}$. 于是 $K_2F$ 以元素对 $\{a,b\}$ (每对 $F^{\times}$ 中的元素一个) 为生成元, 具有关系
\[
\{ab,c\}=\{a,c\}\{b,c\},\qquad
\{a,bc\}=\{a,b\}\{a,c\},\qquad
\{u,1-u\}=1.
\]
已知这些关系蕴含
\[
\{u,v\}=\{v,u\}^{-1},\qquad
\{u,-u\}=1
\]
(见 \cite{Rosenberg1994}, p.\,214).

不难证明, 视为 $\Br(k)$ 中元素的 $A(a,b;\zeta)$ 满足这些关系, 于是有良定义的同态
\[
K_{2}k\to\Br(k).
\]
引人注目的是, 已被证明 (Merkurjev--Suslin 定理, 1980 年代初) 这个映射定义了从 $K_2k/nK_2k$ 到 $\Br(k)$ 中被 $n$ 零化的元素所成子群的同构, 从而我们有了 $\Br(k)_n$ 的一个生成元与关系式的显式描述.\footnote{参见 Merkurjev, A.~S.; Suslin, A.~A., \emph{K-cohomology of Severi-Brauer varieties and the norm residue homomorphism} (俄文), Izv.\ Akad.\ Nauk SSSR Ser.\ Mat.\ 46 (1982), no.~5, 1011--1046, 1135--1136. 该定理在 Kersten, Ina, \emph{Brauergruppen von K\"orpern}, Aspects of Mathematics, D6, Friedr.\ Vieweg \& Sohn, Braunschweig, 1990 一书中有讨论.}

\begin{exercise}\label{ex:iv.5.7}
设 $F$ 是特征 $\neq 2$ 的域, 并按 (\ref{ex:iv.1.12}) 定义四元数代数 $H(a,b)$. 于是 $H(a,b)$ 有基 $1,i,j,k$, 且 $i^2=a$, $j^2=b$, $k=ij=-ji$. 它是 $F$ 上的中心单代数.

(a) 证明 $k$ 上每个 $4$ 维中心单代数都同构于某个 $H(a,b)$ ($a,b\in F^{\times}$).

(b) 据 Wedderburn 定理, $H(a,b)\simeq M_2(F)$ 或 $H(a,b)$ 是可除代数. 证明第一种情形出现当且仅当 $w^2-ax^2-by^2+abz^2$ 在 $F$ 中有非平凡零点. (提示: 对 $\alpha=w+xi+yj+zk$, 令 $\bar{\alpha}=w-xi-yj-zk$, 注意 $\alpha\bar{\alpha}=w^2-ax^2-\cdots$.)

(c) 证明 $H(1,1)\simeq M_2(F)$. (提示: 考虑矩阵 $e_{12}+e_{21}$ 与 $e_{11}-e_{22}$.)

(d) 证明对任意 $x,y\in F^{\times}$ 有 $H(a,b)\simeq H(ax^2,by^2)$.

(e) 证明 $H(a,b)\otimes_F L$ 就是由视 $a,b$ 为 $L$ 中元素所定义的 $L$ 上四元数代数.

(f) 验证 $H(a,b)$ 事实上在 $F$ 上中心单.

(g) 证明当 $a,\,1-a\in F^{\times}$ 时, $H(a,1-a)\simeq M_2(F)$.

(h) 证明 $H(1,b)\simeq H(a,-a)\simeq M_2(F)$. (提示: 考虑 $j+k$ 与 $i+j$.)

(i) 证明 $H(a,b)\simeq H(a,b)^{\mathrm{opp}}$.

(j) 证明 $H(a,b)\simeq M_2(F)$ 当且仅当 $a\in\Nm(F[\sqrt{b}])$.

(k) 证明映射 $\{a,b\}\mapsto[H(a,b)]\colon K_2F\to\Br(F)$ 良定义.
\end{exercise}

\begin{exercise}\label{ex:iv.5.8}
对给定的 $\Q$ 上中心可除代数 $D$, 确定能嵌入 $D$ 的数域.
\end{exercise}

\begin{exercise}\label{ex:iv.5.9}
确定能嵌入 $\Q$ 上某个中心可除代数的数域.

习题\ref{ex:iv.5.8} 与 \ref{ex:iv.5.9} 的解见 \sx{178999}.
\end{exercise}

\begin{NOTES}\label{note:iv.5}
Brauer 群由 R.~Brauer 引入并研究, 随后 Albert、Brauer、Noether、Hasse 等人自二十世纪二十年代起继续研究. 经典的论述是:

Artin, E., Nesbitt, C., and Thrall, R., \emph{Rings with Minimum Condition}, University of Michigan Press, 1944.

除了古怪的术语 (例如把张量积称为 Kronecker 积) 之外, 后者至今仍是一本优秀的书.

其他著作包括 \cite{Blanchard1972}、\cite{FarbDennis1993} 与 \cite{Herstein1968}. 这些书都论及 Brauer 群, 但也涵盖更多内容 (不过没有数论). 其中第二本有大量习题.

在撰写本章第一版时, 我发现下述文献特别有用:

Serre, J-P., \emph{Applications alg\'ebriques de la cohomologie des groupes}, I, II, S\'eminaire Henri Cartan, 1950/51 (现可在 www.numdam.org 获取).
\end{NOTES}
   % IV  Brauer 群
% ============================================================
% 第五章  整体类域论: 主要定理的陈述
% 源: Milne CFT v4.03, Chapter V (书页 p.147--177)
% OCR: 行 5537--6795; 脚注 1--13 从原书 PDF 恢复
% 勘误: p.154-155 三处 1-s -> s-1 (Yiming Tang); 定理 3.6 补 "L/K 有模 m" (Zhuo Zhang)
% 标签前缀: v.  公式 (43)(44) 用 \tag 硬编码
% ============================================================
\chapter{整体类域论: 主要定理的陈述}\label{ch:V}
\index{zhengtileiyulun@整体类域论(global class field theory)}

\begin{quote}
\itshape La théorie du corps de classes a une réputation de difficulté qui est en partie justifiée. Mais il faut faire une distinction: il n'est peut-être pas en effet dans la science de théorie où tout à la fois les démonstrations soient aussi ardues, et les résultats d'une aussi parfaite simplicité et d'une aussi grande puissance.
\end{quote}
\hspace*{\fill}J.~Herbrand, 1936, p.~2\footnote{类域论以困难著称, 这在一定程度上是有道理的. 但必须作出区分: 科学中也许再没有别的理论, 其证明如此艰深, 而其结果却又如此完美地简洁、如此威力巨大. (原书脚注给出的英译, 此处译出.)}

本章陈述并解释整体类域论的主要定理. 这些定理将在第 \ref{ch:VII} 章中证明.

本章始终设 $K$ 为数域, 尽管大多数结果对 $\F_p(T)$ 的有限扩张同样成立.

回顾: 对数域 $K$, 我们把 $K$ 的\textbf{素} (prime) 定义为 $K$ 的非平凡绝对值的等价类. 素有两类: \textbf{有限素} (finite prime), 可与 $\mathcal{O}_K$ 的素理想等同; 以及\textbf{无限素} (infinite prime). 实无限素可与一个嵌入 $K\hookrightarrow\R$ 等同, 复无限素可与一对共轭嵌入 $K\hookrightarrow\CC$ 等同. 我们用 $\mathfrak{p}$ 或 $v$ 表示素 (有限或无限), 用 $S$ 表示 $K$ 的一个有限素集, 同时也表示 $K$ 的有限扩张 $L$ 中位于其上的素集. 无限素集记作 $S_\infty$.

$K$ 在素 $\mathfrak{p}$ (相应地 $v$) 处的完备化记作 $K_{\mathfrak{p}}$ (相应地 $K_v$), 包含映射 $K\hookrightarrow K_{\mathfrak{p}}$ (相应地 $K\hookrightarrow K_v$) 记作 $a\mapsto a_{\mathfrak{p}}$ (相应地 $a\mapsto a_v$).

\section{射线类群}\label{sec:V.1}

\subsection*{与 $S$ 互素的理想}

设 $I=I_K$ 为 $K$ 的分式理想群. 对 $K$ 的素的有限集 $S$, 定义 $I^S$ 为由不在 $S$ 中的素理想生成的 $I$ 的子群. 每个 $\mathfrak{a}\in I^S$ 有唯一的分解
\[
\mathfrak{a}=\mathfrak{p}_1^{n_1}\cdots\mathfrak{p}_s^{n_s},\quad\mathfrak{p}_i\notin S,\quad n_i\in\Z,
\]
因此 $I^S$ 可与由不在 $S$ 中的素理想生成的自由 Abel 群等同. 定义
\[
K^S=\{a\in K^{\times}\mid(a)\in I^S\}=\{a\in K^{\times}\mid\mathrm{ord}_{\mathfrak{p}}(a)=0\ \text{对所有有限}\ \mathfrak{p}\in S\}.
\]
设 $i\colon K^S\to I^S$ 为把 $a\in K^S$ 送到理想 $a\mathcal{O}_K$ 的映射.

例如, 当 $K=\Q$ 且 $S$ 为整除整数 $n$ 的素数集时, $I^S$ 是分式理想
\[
\{(r/s)\mid r,s\in\Z,\ \gcd(r,n)=1=\gcd(s,n)\}
\]
的集合, 而
\[
\Q^S=\{r/s\mid r,s\in\Z,\ \gcd(r,n)=1=\gcd(s,n)\}.
\]
此时自然映射 $\Q^S\to I^S$ 是满射, 核为 $\{\pm1\}$.

\begin{lemma}\label{lem:v.1.1}
对 $\mathcal{O}_K$ 中素理想的每个有限集 $S$, 序列
\[
0\to U_K\to K^S\to I^S\to C\to0
\]
正合. (此处 $U_K=\mathcal{O}_K^{\times}$, 而 $C$ 是完整理想类群 $I/i(K^{\times})$.)
\end{lemma}

\begin{proof}
为证 $I^S\to C$ 满, 须证每个理想类 $\mathcal{C}$ 都由 $I^S$ 中的理想代表. 设 $\mathfrak{a}$ 代表 $\mathcal{C}$. 则 $\mathfrak{a}=\mathfrak{b}\mathfrak{c}^{-1}$, 其中 $\mathfrak{b},\mathfrak{c}$ 为整理想, 且对任意 $c\in\mathfrak{c}$, $\mathfrak{a}\cdot(c)=\mathfrak{b}\mathfrak{c}^{-1}\cdot(c)$ 是整理想, 故可设 $\mathfrak{a}$ 本身是整理想. 写 $\mathfrak{a}=\prod_{\mathfrak{p}\in S}\mathfrak{p}^{n(\mathfrak{p})}\mathfrak{b}$, 其中 $\mathfrak{b}\in I^S$. 对每个 $\mathfrak{p}\in S$, 取 $\pi_{\mathfrak{p}}\in\mathfrak{p}\setminus\mathfrak{p}^2$, 使 $\mathrm{ord}_{\mathfrak{p}}(\pi_{\mathfrak{p}})=1$. 由中国剩余定理, 存在 $a\in\mathcal{O}_K$ 使得
\[
a\equiv\pi_{\mathfrak{p}}^{n(\mathfrak{p})}\bmod\mathfrak{p}^{n(\mathfrak{p})+1}\qquad(\text{对所有 }\mathfrak{p}\in S).
\]
这些同余式蕴含 $\mathrm{ord}_{\mathfrak{p}}(a)=n(\mathfrak{p})$ 对所有 $\mathfrak{p}\in S$ 成立, 于是 $(a)=\prod_{\mathfrak{p}\in S}\mathfrak{p}^{n(\mathfrak{p})}\mathfrak{b}'$, 其中 $\mathfrak{b}'\in I^S$. 现在 $a^{-1}\mathfrak{a}\in I^S$ 且与 $\mathfrak{a}$ 在 $C$ 中代表同一个类.

其次, 若 $\mathfrak{a}\in I^S$ 在 $C$ 中映为零, 则 $\mathfrak{a}=(\alpha)$ 对某个 $\alpha\in K^S$ 成立, 且 $\alpha$ 在相差一个单位的意义下唯一确定. 这就证明了其余各处的正合性.
\end{proof}

\begin{remark}\label{rem:v.1.2}
事实上, $C$ 中每个类都可由 $I^S$ 中的一个整理想 $\mathfrak{a}$ 代表: 设该类由 $\mathfrak{a}\in I^S$ 代表; 写 $\mathfrak{a}=\mathfrak{b}\mathfrak{c}^{-1}$, 其中 $\mathfrak{b},\mathfrak{c}$ 为 $I^S$ 中的整理想, 取非零的 $c\in\mathfrak{c}\cap K^S$ (由中国剩余定理存在), 并注意 $c\mathfrak{a}$ 是整理想.
\end{remark}

\subsection*{模}

\begin{definition}\label{def:v.1.3}
$K$ 的一个\emph{模} (modulus) 是一个函数
\[
m\colon\{K\ \text{的素}\}\to\Z
\]
满足:
\begin{enumerate}[(a)]
\item $m(\mathfrak{p})\ge 0$ 对所有素 $\mathfrak{p}$, 且除有限多个 $\mathfrak{p}$ 外 $m(\mathfrak{p})=0$;
\item 若 $\mathfrak{p}$ 为实素, 则 $m(\mathfrak{p})=0$ 或 $1$;
\item 若 $\mathfrak{p}$ 为复素, 则 $m(\mathfrak{p})=0$.
\end{enumerate}
\end{definition}

习惯上记
\[
\mathfrak{m}=\prod_{\mathfrak{p}}\mathfrak{p}^{m(\mathfrak{p})}.
\]
称模 $\mathfrak{m}=\prod\mathfrak{p}^{m(\mathfrak{p})}$ \textbf{整除}模 $\mathfrak{n}=\prod_{\mathfrak{p}}\mathfrak{p}^{n(\mathfrak{p})}$, 若 $m(\mathfrak{p})\le n(\mathfrak{p})$ 对所有 $\mathfrak{p}$ 成立. 特别地, 素 $\mathfrak{p}$ 整除模 $\mathfrak{m}$ 当且仅当 $m(\mathfrak{p})>0$.

模 $\mathfrak{m}$ 可写成
\[
\mathfrak{m}=\mathfrak{m}_\infty\mathfrak{m}_0,
\]
其中 $\mathfrak{m}_\infty$ 是一些实素之积, $\mathfrak{m}_0$ 是素理想的正数次幂之积, 故可与 $\mathcal{O}_K$ 中的一个理想等同.

\subsection*{射线类群}

对模 $\mathfrak{m}$, 定义 $K_{\mathfrak{m},1}$ 为满足如下条件的 $a\in K^{\times}$ 的集合:
\[
\begin{cases}
\mathrm{ord}_{\mathfrak{p}}(a-1)\ge m(\mathfrak{p}) & \text{对所有整除 }\mathfrak{m}\text{ 的有限 }\mathfrak{p},\\[2pt]
a_{\mathfrak{p}}>0 & \text{对所有整除 }\mathfrak{m}\text{ 的实 }\mathfrak{p}.
\end{cases}
\]
注意
\[
\mathrm{ord}_{\mathfrak{p}}(a-1)\ge m(\mathfrak{p})
\iff \pi^{m(\mathfrak{p})}\mid(a_{\mathfrak{p}}-1)
\iff a\mapsto 1\ \text{于}\ (\mathcal{O}_{\mathfrak{p}}/\mathfrak{p}^{m(\mathfrak{p})})^{\times}\simeq(\hat{\mathcal{O}}_{\mathfrak{p}}/\hat{\mathfrak{p}}^{m(\mathfrak{p})})^{\times},
\]
其中 $\pi$ 是 $K$ 在 $\mathfrak{p}$ 处的完备化 $K_{\mathfrak{p}}$ 中的素元. 令
\[
S(\mathfrak{m})=\{\text{整除 }\mathfrak{m}\text{ 的素}\}.
\]
对任意 $a\in K_{\mathfrak{m},1}$ 及整除 $\mathfrak{m}$ 的素理想 $\mathfrak{p}$, 有 $\mathrm{ord}_{\mathfrak{p}}(a-1)>0=\mathrm{ord}_{\mathfrak{p}}(1)$, 故
\[
\mathrm{ord}_{\mathfrak{p}}(a)=\mathrm{ord}_{\mathfrak{p}}\bigl((a-1)+1\bigr)=0.
\]
因此对任意 $a\in K_{\mathfrak{m},1}$, 理想 $(a)$ 属于 $I^{S(\mathfrak{m})}$. 设 $i$ 表示映射 $a\mapsto(a)\colon K_{\mathfrak{m},1}\to I^{S(\mathfrak{m})}$. 商群
\[
C_{\mathfrak{m}}=I^{S(\mathfrak{m})}/i(K_{\mathfrak{m},1})
\]
称为\textbf{模 $\mathfrak{m}$ 的 (射线) 类群}\index{shexianleiqun@射线类群(ray class group)}\index{leiqun@类群(class group)!shexianleiqun@射线类群(ray class group)}.

\begin{example}\label{ex:v.1.4}
表达式 $\mathfrak{m}=(2)^3\cdot(17)^2\cdot(19)\cdot\infty$ 是 $\Q$ 的一个模, 其中 $\mathfrak{m}_0=(2)^3\cdot(17)^2\cdot(19)$, $\mathfrak{m}_\infty=\infty$ (此处 $\infty$ 表示 $\Q$ 的唯一无限素). 此外, $\Q_{\mathfrak{m},1}$ 由满足
\[
\begin{cases}
\mathrm{ord}_2(a-1)\ge 3,\\
\mathrm{ord}_{17}(a-1)\ge 2,\\
\mathrm{ord}_{19}(a-1)\ge 1
\end{cases}
\]
的正有理数 $a$ 组成. 条件 ``$2$ 处'' 是说 $a$ 是两个奇整数之商 $a=b/c$, 且 $bc^{-1}$ 在 $(\Z/8\Z)^{\times}$ 中的像为 $1$. 其余条件可类似表述.
\end{example}

\begin{lemma}\label{lem:v.1.5}
设 $S$ 为 $K$ 的素理想的有限集. 则每个 $\alpha\in K^S$ 可写成 $\alpha=a/b$, 其中 $a,b\in\mathcal{O}_K\cap K^S$.
\end{lemma}

\begin{proof}
因 $\alpha\in K^S$, 有 $(\alpha)=\mathfrak{a}/\mathfrak{b}$, 其中 $\mathfrak{a},\mathfrak{b}$ 为 $I^S$ 中的整理想. 显然 $\mathfrak{a}$ 与 $\mathfrak{b}$ 代表理想类群中同一个元 $C$, 由注记 \ref{rem:v.1.2} 可取 $I^S$ 中的整理想 $\mathfrak{c}$ 代表 $C^{-1}$. 于是 $(\alpha)=\mathfrak{ac}/\mathfrak{bc}=(a)/(b)$, 对某个 $a,b\in\mathcal{O}_K\cap K^S$ 成立.
\end{proof}

\begin{proposition}\label{prop:v.1.6}
$C_{\mathfrak{m}}$ 中每个类都由一个整理想 $\mathfrak{a}$ 代表, 且两个整理想 $\mathfrak{a}$ 与 $\mathfrak{b}$ 代表 $C_{\mathfrak{m}}$ 中同一个类, 当且仅当存在非零的 $a,b\in\mathcal{O}_K$ 使得 $a\mathfrak{a}=b\mathfrak{b}$,
\[
a\equiv b\equiv 1\bmod\mathfrak{m}_0
\]
且 $a$ 与 $b$ 在每个整除 $\mathfrak{m}$ 的实素处有相同的符号.
\end{proposition}

\begin{proof}
设该类由 $\mathfrak{a}\in I^S$ 代表, 写 $\mathfrak{a}=\mathfrak{b}\mathfrak{c}^{-1}$, 其中 $\mathfrak{b},\mathfrak{c}$ 为 $I^S$ 中的整理想. 中国剩余定理表明存在非零的 $c\in\mathfrak{c}\cap K_{\mathfrak{m}_0,1}$, 而强逼近定理表明可取 $c$ 在所有实素处均为正. 此时 $c\mathfrak{a}$ 是整理想且与 $\mathfrak{a}$ 在 $C_{\mathfrak{m}}$ 中代表同一个类. 结论的第二部分由引理 \ref{lem:v.1.5} 推出.
\end{proof}

于是, 例如, 通常的理想类群可与下述等价关系下整理想的集合等同: $\mathfrak{a}\sim\mathfrak{b}$ 当且仅当 $a\mathfrak{a}=b\mathfrak{b}$ 对某对非零的 $a,b\in\mathcal{O}_K$ 成立.

\begin{theorem}\label{thm:v.1.7}
对 $K$ 的每个模 $\mathfrak{m}$, 有正合列
\[
0\to U/U_{\mathfrak{m},1}\to K_{\mathfrak{m}}/K_{\mathfrak{m},1}\to C_{\mathfrak{m}}\to C\to0
\]
与典范同构
\[
K_{\mathfrak{m}}/K_{\mathfrak{m},1}\simeq\prod_{\substack{\mathfrak{p}\ \mathrm{real}\\ \mathfrak{p}|\mathfrak{m}}}\{\pm\}\times\prod_{\substack{\mathfrak{p}\ \mathrm{finite}\\ \mathfrak{p}|\mathfrak{m}}}(\mathcal{O}_K/\mathfrak{p}^{m(\mathfrak{p})})^{\times}
\simeq\prod_{\substack{\mathfrak{p}\ \mathrm{real}\\ \mathfrak{p}|\mathfrak{m}}}\{\pm\}\times(\mathcal{O}_K/\mathfrak{m}_0)^{\times},
\]
其中
\begin{align*}
K_{\mathfrak{m}}&=K^{S(\mathfrak{m})}=\{\alpha\in K^{\times}\mid \mathrm{ord}_{\mathfrak{p}}(\alpha)=0\ \text{对所有 }\mathfrak{p}|\mathfrak{m}_0\},\\
U&=\mathcal{O}_K^{\times}\quad(\text{$K$ 的单位群}),\\
U_{\mathfrak{m},1}&=U\cap K_{\mathfrak{m},1}.
\end{align*}
因此 $C_{\mathfrak{m}}$ 是有限群, 其阶为
\[
h_{\mathfrak{m}}=h\cdot(U:U_{\mathfrak{m},1})^{-1}\cdot 2^{r_0}\cdot\Nm(\mathfrak{m}_0)\cdot\prod_{\mathfrak{p}|\mathfrak{m}_0}\Bigl(1-\frac{1}{\Nm\mathfrak{p}}\Bigr),
\]
其中 $r_0$ 是整除 $\mathfrak{m}$ 的实素个数, $h$ 是 $K$ 的类数 ($C$ 的阶).
\end{theorem}

\begin{proof}
包含映射 $I^{S(\mathfrak{m})}\to I$ 定义了一个同态 $C_{\mathfrak{m}}\to C$. 考虑映射对
\[
K_{\mathfrak{m},1}\xrightarrow{\ f\ }K_{\mathfrak{m}}\xrightarrow{\ g\ }I^{S(\mathfrak{m})}.
\]
由引理 \ref{lem:v.1.1}, $g$ 的核与余核分别为 $U$ 与 $C$. $g\circ f$ 的余核是 $C_{\mathfrak{m}}$ (按定义), 其核是 $K_{\mathfrak{m},1}\cap U=U_{\mathfrak{m},1}$. 最后 $f$ 是单射. 因此该映射对的核--余核序列 (见 II, \ref{lem:ii.A.2}) 为
\[
0\to U_{\mathfrak{m},1}\to U\to K_{\mathfrak{m}}/K_{\mathfrak{m},1}\to C_{\mathfrak{m}}\to C\to0.
\]
下面证明 $K_{\mathfrak{m}}$ 与所给群典范同构. 设 $\mathfrak{p}$ 是整除 $\mathfrak{m}$ 的素. 若 $\mathfrak{p}$ 为实素, 则把 $\alpha\in K_{\mathfrak{m}}$ 映到 $\alpha_{\mathfrak{p}}$ 的符号 (回顾: 实素是一个嵌入 $K\hookrightarrow\R$, 而 $\alpha_{\mathfrak{p}}$ 表示 $\alpha$ 在该嵌入下的像). 若 $\mathfrak{p}$ 为有限素, 即它是 $\mathcal{O}_K$ 中的素理想, 则把 $\alpha\in K_{\mathfrak{m}}$ 映到 $[a][b]^{-1}\in(\mathcal{O}_K/\mathfrak{p}^{m(\mathfrak{p})})^{\times}$, 其中 $a,b$ 如引理 \ref{lem:v.1.5} 所述. 由于 $a$ 与 $b$ 都和 $\mathfrak{p}$ 互素, 它们在 $\mathcal{O}_K/\mathfrak{p}^{m(\mathfrak{p})}$ 中的类 $[a]$ 与 $[b]$ 可逆, 故此映射有意义. 弱逼近定理 (ANT, 7.20) 表明映射 $K_{\mathfrak{m}}\to\prod\{\pm\}\times\prod(\mathcal{O}_K/\mathfrak{p}^{m(\mathfrak{p})})^{\times}$ 是满射, 其核显然是 $K_{\mathfrak{m},1}$.

中国剩余定理表明有环的同构
\[
\mathcal{O}_K/\mathfrak{m}_0\simeq\prod_{\mathfrak{p}|\mathfrak{m}}\mathcal{O}_K/\mathfrak{p}^{m(\mathfrak{p})},
\]
从而有群的同构
\[
(\mathcal{O}_K/\mathfrak{m}_0)^{\times}\simeq\prod(\mathcal{O}_K/\mathfrak{p}^{m(\mathfrak{p})})^{\times}.
\]
这就完成了同构部分的证明.

剩下计算群的阶. 注意 $\mathcal{O}_K/\mathfrak{p}^m$ 是局部环, 极大理想为 $\mathfrak{p}/\mathfrak{p}^m$ (因为其理想对应于 $\mathcal{O}_K$ 中包含 $\mathfrak{p}^m$ 的理想), 故其单位是不属于 $\mathfrak{p}/\mathfrak{p}^m$ 的元素. 滤链
\[
(\mathcal{O}_K/\mathfrak{p}^m)^{\times}\supset(1+\mathfrak{p})/\mathfrak{p}^m\supset\cdots\supset(1+\mathfrak{p}^{m-1})/\mathfrak{p}^m\supset 0
\]
的商群同构于
\[
k^{\times},\,k,\,\ldots,\,k,\qquad k\overset{\mathrm{def}}{=}\mathcal{O}_K/\mathfrak{p},
\]
故 $(\mathcal{O}_K/\mathfrak{p}^m)^{\times}$ 的阶为 $(q-1)q^{m-1}$, $q=(\mathcal{O}_K:\mathfrak{p})\overset{\mathrm{def}}{=}\Nm\mathfrak{p}$. 这表明
\[
(C_{\mathfrak{m}}:1)=(C:1)\cdot(K_{\mathfrak{m}}:K_{\mathfrak{m},1})\cdot(U:U_{\mathfrak{m},1})^{-1}
\]
等于定理陈述中的表达式.
\end{proof}

\begin{example}\label{ex:v.1.8}
\begin{enumerate}[(a)]
\item 若 $\mathfrak{m}=1$, 则 $C_{\mathfrak{m}}=C$.
\item 当 $\mathfrak{m}$ 是所有实素之积时, $C_{\mathfrak{m}}$ 是\textbf{狭窄类群} (narrow class group)\index{leiqun@类群(class group)!xiazhai@狭窄类群(narrow class group)}, 且有正合列
\[
0\to U/U_+\to K^{\times}/K_+\to C_{\mathfrak{m}}\to C\to1,
\]
其中 $K_+$ 是全正元素群 (即在所有实嵌入下均为正), $U_+$ 是全正单位群. 此外 $K^{\times}/K_+\simeq\prod_{\mathfrak{p}\,\mathrm{real}}\{\pm\}$, 故 $C_{\mathfrak{m}}\to C$ 的核是所有可能的符号模去由单位产生的那些.

对 $\Q$, 狭窄类群平凡. 对 $\Q[\sqrt{d}]$, $d>0$, 有两个实素, 且 $U=\{\pm\varepsilon^m\mid m\in\Z\}\approx(\Z/2\Z)\times\Z$, 其中 $\varepsilon$ 是基本单位. 设 $\bar{\varepsilon}$ 为 $\varepsilon$ 的共轭. 则依 $\varepsilon$ 与 $\bar{\varepsilon}$ 符号相同或不同, 有 $h_m=2h$ 或 $h$. 注意 $\Nm(\varepsilon)=+1$ 当符号相同, $-1$ 当符号相反. 对小的 $d$ 值:
\[
\begin{array}{cccc}
d & h & \varepsilon & \Nm(\varepsilon)\\
2 & 1 & 1+\sqrt{2} & -1\\
3 & 1 & 2+\sqrt{3} & 1\\
5 & 1 & (1+\sqrt{5})/2 & -1\\
6 & 1 & 5+2\sqrt{6} & 1
\end{array}
\]
因此 $\Q[\sqrt{3}]$ 与 $\Q[\sqrt{6}]$ 的类数为 $1$ 而狭窄类数为 $2$, 而对 $\Q[\sqrt{2}]$ 与 $\Q[\sqrt{5}]$, 两个类数都是 $1$.
\item 对域 $\Q$ 与模 $(m)$, 该正合列成为
\[
0\to\{\pm1\}\to(\Z/m\Z)^{\times}\to C_{\mathfrak{m}}\to 0.
\]
对模 $\infty(m)$, 正合列成为
\[
0\to\{\pm1\}\to\{\pm\}\times(\Z/m\Z)^{\times}\to C_{\mathfrak{m}}\to 0.
\]
此处 $-1$ 映到 $(-,[-1])$, 而积的子群 $(\Z/m\Z)^{\times}$ 同构地映满商群 $C_{\mathfrak{m}}$.
\end{enumerate}
\end{example}

\subsection*{Frobenius 元}

我们回顾 Frobenius 元的理论 (ANT, 第 8 章). 设 $K$ 为数域, $L$ 为 $K$ 的有限 Galois 扩张, 群为 $G$. 设 $\mathfrak{p}$ 是 $K$ 的一个素理想, $\mathfrak{P}$ 是 $L$ 中位于其上的素理想. 分解群 $D(\mathfrak{P})$ (或 $G(\mathfrak{P})$) 定义为
\[
\{\tau\in G\mid\tau\mathfrak{P}=\mathfrak{P}\}.
\]
等价地, 它是 $G$ 中对 $\mathfrak{P}$-adic 拓扑连续, 从而延拓到完备化 $L_{\mathfrak{P}}$ 的元素之集. 由此得到同构
\[
D(\mathfrak{P})\to\Gal(L_{\mathfrak{P}}/K_{\mathfrak{p}}).
\]
设 $\mathfrak{P}$ 在 $\mathfrak{p}$ 上不分歧. 则 $\Gal(L_{\mathfrak{P}}/K_{\mathfrak{p}})$ 在 $\mathcal{O}_L$ 上的作用诱导同构
\[
\Gal(L_{\mathfrak{P}}/K_{\mathfrak{p}})\to\Gal(l/k),
\]
其中 $l$ 与 $k$ 是剩余类域. 图示如下:
\[
\begin{tikzcd}[column sep=small, row sep=small]
\mathfrak{P} \ar[dd, no head] & L \ar[r, no head] \ar[dd, "f", "D(\mathfrak{P})"'] & L_{\mathfrak{P}} \ar[r, no head] \ar[dd, "f", "D(\mathfrak{P})"'] & l \ar[dd, no head, "D(\mathfrak{P})"'] \\
& & & \\
\mathfrak{P}_D \ar[dd, no head] & L^{D(\mathfrak{P})} \ar[r, no head] \ar[dd, "g"'] & K_{\mathfrak{p}} \ar[r, no head] \ar[ddl, no head] & k \\
& & & \\
\mathfrak{p} & K \ar[r, no head] & \mathfrak{p} &
\end{tikzcd}
\]
群 $\Gal(l/k)$ 是循环群, 有典范生成元, 即 Frobenius 元 $x\mapsto x^q$, 其中 $q$ 是 $k$ 的元素个数. 从而 $D(\mathfrak{P})$ 循环, $D(\mathfrak{P})$ 中对应于 $\Gal(l/k)$ 里 Frobenius 元的生成元称为在 $\mathfrak{P}$ 处的 \textbf{Frobenius 元}\index{fuluobeniusiyuan@Frobenius 元(Frobenius element)} $(\mathfrak{P},L/K)$. 它是 $\Gal(L/K)$ 中唯一满足下列两个条件的元素 $\sigma$:

(a) $\sigma\in D(\mathfrak{P})$, 即 $\sigma\mathfrak{P}=\mathfrak{P}$;

(b) 对所有 $\alpha\in\mathcal{O}_L$, $\sigma\alpha\equiv\alpha^q\bmod\mathfrak{P}$, 其中 $q$ 是剩余类域 $\mathcal{O}_K/\mathfrak{p}$, $\mathfrak{p}=\mathfrak{P}\cap K$ 的元素个数.

现在列举 $(\mathfrak{P},L/K)$ 的基本性质.

\textbf{1.9.} 设 $\tau\mathfrak{P}$ 是整除 $\mathfrak{p}$ 的另一个素. 则 $D(\tau\mathfrak{P})=\tau D(\mathfrak{P})\tau^{-1}$, 且
\[
(\tau\mathfrak{P},L/K)=\tau(\mathfrak{P},L/K)\tau^{-1}.
\]

\begin{proof}
若 $\rho\in D(\mathfrak{P})$, 则
\[
\tau\rho\tau^{-1}(\tau\mathfrak{P})=\tau\rho\mathfrak{P}=\tau\mathfrak{P},
\]
故 $\tau\rho\tau^{-1}\in D(\tau\mathfrak{P})$. 于是 $\tau D(\mathfrak{P})\tau^{-1}\subset D(\tau\mathfrak{P})$, 又因两者阶相同, 必相等.

设 $\alpha\in\mathcal{O}_L$, $\sigma=(\mathfrak{P},L/K)$; 则
\[
\tau\sigma\tau^{-1}(\alpha)=\tau\bigl((\tau^{-1}\alpha)^q+a\bigr),\quad\text{某个 }a\in\mathfrak{P},
\]
而
\[
\tau\bigl((\tau^{-1}\alpha)^q+a\bigr)=\alpha^q+\tau a\equiv\alpha^q\bmod\tau\mathfrak{P}.
\]
\end{proof}

由于 $G$ 在整除 $\mathfrak{p}$ 的素上传递作用, 这表明
\[
\{(\mathfrak{P},L/K)\mid\mathfrak{P}|\mathfrak{p}\}
\]
是 $G$ 中的一个共轭类, 我们记之为 $(\mathfrak{p},L/K)$. 当 $L/K$ 为 Abel 扩张时, $(\mathfrak{p},L/K)$ 只含一个元素, 我们把它视为 $\Gal(L/K)$ 的一个元素 (而非由单个元素构成的集合).

\phantomsection\hypertarget{item:v.1.10}{}\textbf{1.10.} 考虑域塔
\[
\begin{tikzcd}
M \ar[d, no head] & \Omega \ar[d, no head] \\
L \ar[d, no head] & \mathfrak{P} \ar[d, no head] \\
K & \mathfrak{p}
\end{tikzcd}
\]
其中 $\Omega$ 是 $M$ 中位于 $\mathfrak{p}$ 之上的素理想, $\mathfrak{P}=\Omega\cap\mathcal{O}_L$. 设 $\Omega$ 在 $\mathfrak{p}$ 上不分歧; 则
\[
(\Omega,M/L)=(\Omega,M/K)^{f(\mathfrak{P}/\mathfrak{p})}.
\]

\begin{proof}
设 $k(\Omega)\supset k(\mathfrak{P})\supset k(\mathfrak{p})$ 为剩余类域塔. 则 $f(\mathfrak{P}/\mathfrak{p})\overset{\mathrm{def}}{=}[k(\mathfrak{P}):k(\mathfrak{p})]$, 且 $\Gal(k(\Omega)/k(\mathfrak{P}))$ 中的 Frobenius 元是 $\Gal(k(\Omega)/k(\mathfrak{p}))$ 中 Frobenius 元的 $f(\mathfrak{P}/\mathfrak{p})$ 次幂. 其余是直接的.
\end{proof}

\phantomsection\hypertarget{item:v.1.11}{}\textbf{1.11.} 在 \hyperlink{item:v.1.10}{1.10} 中设 $L$ 在 $K$ 上 Galois; 则
\[
(\Omega,M/K)\big|_L=(\mathfrak{P},L/K).
\]

\begin{proof}
显然 $(\Omega,M/K)|_L$ 满足刻画 $(\mathfrak{P},L/K)$ 的条件.
\end{proof}

设 $L_1$ 与 $L_2$ 是含于某个域 $\Omega$ 的 $K$ 的 Galois 扩张, 并令 $M=L_1\cdot L_2$. 则 $M$ 在 $K$ 上 Galois, 且有单同态
\[
\sigma\mapsto(\sigma|L_1,\sigma|L_2)\colon\Gal(M/K)\to\Gal(L_1/K)\times\Gal(L_2/K).
\]

\phantomsection\hypertarget{item:v.1.12}{}\textbf{1.12.} 设 $\mathfrak{Q}$ 是 $\mathcal{O}_M$ 的素理想, $\mathfrak{P}_i=\mathfrak{Q}\cap\mathcal{O}_{L_i}$. 在上述映射下,
\[
(\mathfrak{Q},M/K)\mapsto(\mathfrak{P}_1,L_1/K)\times(\mathfrak{P}_2,L_2/K).
\]

\begin{proof}
应用 \hyperlink{item:v.1.11}{1.11}.
\end{proof}

注意 $\mathfrak{p}$ 在 $L$ 中完全分裂当且仅当 $(\mathfrak{P},L/K)=1$ 对某个 (从而所有) 位于其上的 $\mathfrak{P}$ 成立. 因此在 \hyperlink{item:v.1.12}{1.12} 的情形, $\mathfrak{p}$ 在 $M$ 中完全分裂当且仅当它在 $L_1$ 与 $L_2$ 中都完全分裂.

\section{L-级数与算术级数中素数的密度}\label{sec:V.2}

我们先简要回顾 Dirichlet $L$-级数的初等理论 (例如见 \cite{Serre1970}, 第 VI 章).

设 $m$ 为整数. \textbf{模 $m$ 的 Dirichlet 特征}\index{dirichlettezheng@Dirichlet 特征(Dirichlet character)}是同态 $\chi\colon(\Z/m\Z)^{\times}\to\CC^{\times}$. 因 $(\Z/m\Z)^{\times}$ 有限, $\chi([n])$ 对所有 $n$ 是 $1$ 的根. 模 $m$ 的 Dirichlet 特征可视为与 $m$ 互素的整数集上的乘性函数, 其在 $n$ 处的值只依赖 $n$ 模 $m$. 常把 $\chi$ 延拓到全体整数上, 规定当 $\gcd(m,n)\ne 1$ 时 $\chi(n)=0$. 模 $m$ 的平凡 Dirichlet 特征, 即在与 $m$ 互素的所有整数上取值 $1$, 称为主 Dirichlet 特征 $\chi_0$.

对模 $m$ 的 Dirichlet 特征 $\chi$, 附带一个 Dirichlet 级数
\[
L(s,\chi)=\prod_{p\nmid m}\frac{1}{1-\chi(p)p^{-s}}=\sum_{n>0}\chi(n)/n^s.
\]
两个表达式都对 $\Re(s)>1$ 的复数 $s$ 收敛, 二者的相等正是唯一分解的解析表述. 注意 $L(s,\chi_0)$ 与 Riemann zeta 函数 $\zeta(s)$ 仅相差被积的因子 $\frac{1}{1-p^{-s}}$ (对整除 $m$ 的 $p$).

\begin{theorem}\label{thm:v.2.1}
\begin{enumerate}[(a)]
\item zeta 函数 $\zeta(s)$ 可延拓为半平面 $\Re(s)>0$ 上的亚纯函数, 且
\[
\zeta(s)=\frac{1}{s-1}+\varphi(s),
\]
其中 $\varphi(s)$ 在 $\Re(s)>0$ 上全纯.
\item 若 $\chi\ne\chi_0$, 则 $L(s,\chi)$ 的级数对 $\Re(s)>0$ 收敛, 且 $L(1,\chi)\ne 0$.
\end{enumerate}
\end{theorem}

\begin{proof}
见 \cite{Serre1970}, 第 VI 章, 命题 10, 12, 定理 1.
\end{proof}

对 (a) 中的等式取对数, 利用近似 $-\log(1-x)\sim x$ 可得
\[
\sum 1/p^s\sim\log\frac{1}{s-1}\quad(s\downarrow 1).\erratum{原书此处及下文密度定义、定理 \ref{thm:v.2.2} 的证明中共三处误作 $\log\frac{1}{1-s}$, 依 Milne errata (Yiming Tang) 改正.}
\]
我们的意思是: 商 $\dfrac{\sum 1/p^s}{-\log(s-1)}$ 当 $s$ 从大于 $1$ 的实数一侧趋近 $1$ 时收敛到 $1$. 这一结果使如下定义显得合理: 若
\[
\sum_{p\in T}1/p^s\sim\delta\log\frac{1}{s-1}\quad(s\downarrow 1),
\]
则称素数集 $T$ 具有 \textbf{Dirichlet (或解析) 密度} $\delta$.

定义 $f_{\chi}(s)=\sum_{p\nmid m}\chi(p)/p^s$. 则 (b) 表明, 对 $\chi\ne\chi_0$, $f_{\chi}(s)$ 在 $s=1$ 附近有界. 一个初等论证 (\cite{Serre1970}, 第 VI 章, 引理 9) 表明, 对每个与 $m$ 互素的 $a$,
\[
\sum_{p\equiv a\bmod m}1/p^s=\frac{1}{\varphi(m)}\sum_{\chi}\chi(a)^{-1}f_{\chi}(s),
\]
其中 $\varphi(m)\overset{\mathrm{def}}{=}|(\Z/m\Z)^{\times}|$, 求和遍历模 $m$ 的所有 Dirichlet 特征.

\begin{theorem}\label{thm:v.2.2}
对每个与 $m$ 互素的 $a$, 算术级数
\[
\ldots,\ a-2m,\ a-m,\ a,\ a+m,\ a+2m,\ \ldots
\]
中的素数具有 Dirichlet 密度 $1/\varphi(m)$.
\end{theorem}

\begin{proof}
对 $\chi\ne\chi_0$, $f_{\chi}(s)$ 在 $s=1$ 附近有界, 故
\[
\sum_{p\equiv a\bmod m}1/p^s\sim\frac{1}{\varphi(m)}\chi_0(a)^{-1}f_{\chi_0}(s)\sim\frac{1}{\varphi(m)}\log\frac{1}{s-1}\quad(s\downarrow 1).
\]
\end{proof}

\begin{corollary}\label{cor:v.2.3}
对每个 $m$, 在分圆域 $\Q[\zeta_m]$ 中分裂的素数集具有 Dirichlet 密度 $1/\varphi(m)$.
\end{corollary}

\begin{proof}
素理想 $(p)$ 在 $\Q[\zeta_m]$ 中分裂当且仅当 $p\equiv 1\bmod m$.
\end{proof}

现在说明上述结果如何从 $\Q$ 推广到任意数域. 证明将在第 \ref{ch:VI} 章给出.

设 $K$ 为数域, $\mathfrak{m}$ 为 $K$ 的模. \textbf{模 $\mathfrak{m}$ 的 (Dirichlet 或 Weber) 特征}是同态 $\chi\colon C_{\mathfrak{m}}\to\CC^{\times}$, 其值同样是 $1$ 的根. 换一种说法, 这样的特征是一个乘性函数 $I^S\to\CC^{\times}$, 它在 $i(K_{\mathfrak{m},1})$ 上取零值, 对某个满足 $S(\mathfrak{m})=S$ 的模 $\mathfrak{m}$. 模 $\mathfrak{m}$ 的主特征是只取值 $1$ 的函数 $\chi_0\colon C_{\mathfrak{m}}\to\CC^{\times}$.

对模 $\mathfrak{m}$ 的特征 $\chi$, 附带一个 Dirichlet 级数
\[
\mathcal{L}(s,\chi)=\prod_{\mathfrak{p}\nmid\mathfrak{m}}\frac{1}{1-\chi(\mathfrak{p})\Nm\mathfrak{p}^{-s}}=\sum_{(\mathfrak{a},\mathfrak{m}_0)=\mathcal{O}_K}\chi(\mathfrak{a})/\Nm\mathfrak{a}^s.
\]
乘积遍历与 $\mathfrak{m}_0$ 互素的素理想, 求和遍历与 $\mathfrak{m}_0$ 互素的整理想. 同样, 两个表达式都对 $\Re(s)>1$ 收敛, 其相等正是理想唯一分解的解析表述. $L$-级数 $\mathcal{L}(s,\chi_0)$ 与 Dedekind zeta 函数
\[
\zeta_K(s)\overset{\mathrm{def}}{=}\prod_{\mathfrak{p}}\frac{1}{1-\Nm\mathfrak{p}^{-s}}=\sum_{\mathfrak{a}\subset\mathcal{O}_K}\Nm\mathfrak{a}^{-s}
\]
仅相差有限多个因子.

\begin{theorem}\label{thm:v.2.4}
\begin{enumerate}[(a)]
\item zeta 函数 $\zeta_K(s)$ 可延拓为半平面 $\Re(s)>0$ 上的亚纯函数, 且当 $s\downarrow 1$ 时
\[
\zeta_K(s)\sim\frac{2^{r_1}(2\pi)^{r_2}\,\mathrm{Reg}(K)}{w_K|\Delta_{K/\Q}|^{1/2}}\,h_K\cdot\frac{1}{s-1},
\]
其中 $r_1$ 与 $r_2$ 分别是 $K$ 的实素与复素个数, $\mathrm{Reg}(K)$ 是 $K$ 的 regulator (见 ANT, 第 5 章), $w_K$ 是 $K$ 中 $1$ 的根的个数, $\Delta_{K/\Q}$ 是 $K/\Q$ 的判别式, $h_K$ 是类数.
\item 若 $\chi\ne\chi_0$, 则 $L(s,\chi)$ 的级数对 $\Re(s)>0$ 收敛, 且 $L(1,\chi)\ne 0$.
\end{enumerate}
\end{theorem}

(b) 的证明用到存在定理 (见 \ref{thm:v.3.6}).

同样, 对 (a) 中的等式取对数可得
\[
\sum_{\mathfrak{p}}1/\Nm\mathfrak{p}^s\sim\log\frac{1}{s-1}\quad(s\downarrow 1),
\]
且若
\[
\sum_{\mathfrak{p}\in T}1/\Nm\mathfrak{p}^s\sim\delta\log\frac{1}{s-1}\quad(s\downarrow 1),
\]
则称 $K$ 中素理想集 $T$ 具有 \textbf{Dirichlet (或自然) 密度} $\delta$.

与前一情形类似的论证证明:

\begin{theorem}\label{thm:v.2.5}
对每个与 $\mathfrak{m}_0$ 互素的理想 $\mathfrak{a}$, $\mathcal{O}_K$ 中在 $C_{\mathfrak{m}}$ 里类为 $[\mathfrak{a}]$ 的素理想具有 Dirichlet 密度 $1/h_{\mathfrak{m}}$.
\end{theorem}

定理 \ref{thm:v.2.4} 与 \ref{thm:v.2.5} 证明中的分析与 $K=\Q$ 的情形相同, 但数论部分要困难得多.

\section{理想语言的主要定理}\label{sec:V.3}

\subsection*{Artin 映射}

设 $L/K$ 为有限 Abel 扩张, Galois 群为 $G$. 回顾: 对在 $L$ 中不分歧的素理想 $\mathfrak{p}$, 有唯一确定的 Frobenius 自同构 $\sigma=(\mathfrak{p},L/K)\in\Gal(L/K)$, 由下述条件刻画: 对每个位于 $\mathfrak{p}$ 之上的素理想 $\mathfrak{P}$, $\sigma\mathfrak{P}=\mathfrak{P}$, 且 $\sigma\alpha\equiv\alpha^{\Nm\mathfrak{p}}\bmod\mathfrak{P}$.

对包含所有在 $L$ 中分歧的素的 $K$ 的素的每个有限集 $S$, 有同态
\[
\psi_{L/K}\colon I^S\to\Gal(L/K),\qquad\mathfrak{p}_1^{n_1}\cdots\mathfrak{p}_t^{n_t}\mapsto\prod(\mathfrak{p}_i,L/K)^{n_i},
\]
称为\textbf{整体 Artin 映射} (或互反映射)\index{artinyingshe@Artin 映射(Artin map)!zhengti@整体 Artin 映射(global Artin map)}\index{huxiangyingshe@互反映射(reciprocity map)}.

\begin{example}\label{ex:v.3.1}
设 $K=\Q[\sqrt{m}]$, $m$ 为无平方因子整数. 在 $K$ 中分歧的有限素集 $S$: 当 $m\equiv 1\bmod 4$ 时由整除 $m$ 的素组成, 否则由整除 $m$ 的素连同 $2$ 组成. 把 $\Gal(K/\Q)$ 与 $\{\pm1\}$ 等同. Artin 映射是由
\[
p\mapsto\Bigl(\frac{m}{p}\Bigr)\colon I^S\to\Gal(K/\Q)
\]
确定的同态, 其中 $\bigl(\frac{m}{p}\bigr)$ 是二次剩余 (Legendre) 符号.
\end{example}

\begin{example}\label{ex:v.3.2}
设 $L=\Q[\zeta_m]$, $\zeta_m$ 为本原 $m$ 次单位根. 设 $m$ 为奇数或被 $4$ 整除 (此时在 $L$ 中分歧的素恰为整除 $m$ 的素). 把与 $m$ 互素的整数 $n$ 送到 $L$ 的自同构 $\zeta\mapsto\zeta^n$ 的映射定义了一个同构 $(\Z/m\Z)^{\times}\to\Gal(L/\Q)$. 对不整除 $n$ 的 $p$, $(p,L/K)=[p]$. 若 $r$ 与 $s$ 是与 $m$ 互素的正整数, 则 $r/s$ 定义了一个类 $[r/s]=[r][s]^{-1}\in(\Z/m\Z)^{\times}$, 而 Artin 映射是复合
\[
I^S\xrightarrow{\ (r/s)\mapsto[r/s]\ }(\Z/m\Z)^{\times}\xrightarrow{\ [n]\mapsto(\xi\mapsto\xi^n)\ }\Gal(L/\Q).
\]
\end{example}

回顾 (ANT, p.~68): 对数域的任意有限扩张 $L/K$, 范理想映射 $\Nm_{L/K}\colon I_L\to I_K$ 是从 $L$ 的分式理想群到 $K$ 的相应群的唯一同态, 使得对 $L$ 的任意素理想 $\mathfrak{P}$, $\Nm_{L/K}(\mathfrak{P})=\mathfrak{p}^{f(\mathfrak{P}/\mathfrak{p})}$, 其中 $\mathfrak{p}=\mathfrak{P}\cap\mathcal{O}_K$. 对任意 $\alpha\in L$, $\Nm_{L/K}((\alpha))=(\Nm_{L/K}\alpha)$.

\begin{proposition}\label{prop:v.3.3}
设 $L$ 是 $K$ 的 Abel 扩张, $K'$ 是中间域: $L\supset K'\supset K$. 则下图交换:
\[
\begin{tikzcd}
I_{K'}^S\ar[r, "\psi_{L/K'}"]\ar[d, "\Nm"'] & \Gal(L/K')\ar[d, "\text{限制}\ |\!\ _K"]\\
I_K^S\ar[r, "\psi_{L/K}"'] & \Gal(L/K)
\end{tikzcd}
\]
此处 $S$ 是 $K$ 的素理想的任意有限集, 包含所有在 $L$ 中分歧者, 以及 $K'$ 中位于 $S$ 的素之上的素集.
\end{proposition}

\begin{proof}
设 $\mathfrak{p}'$ 是 $K'$ 的任意素理想, 位于 $K$ 的不在 $S$ 中的素理想 $\mathfrak{p}$ 之上. 则 $\Nm_{K'/K}(\mathfrak{p}')=\mathfrak{p}^{f(\mathfrak{p}'/\mathfrak{p})}$, 须证 $\psi_{L/K'}(\mathfrak{p}')=\psi_{L/K}(\mathfrak{p}^{f(\mathfrak{p}'/\mathfrak{p})})$, 即对每个位于 $\mathfrak{p}$ 之上的 $L$ 的素理想 $\mathfrak{P}$, $(\mathfrak{P},L/K')=(\mathfrak{P},L/K)^{f(\mathfrak{p}'/\mathfrak{p})}$. 但这正是 \hyperlink{item:v.1.10}{1.10} 中证明过的.
\end{proof}

\begin{corollary}\label{cor:v.3.4}
对 $K$ 的每个有限 Abel 扩张 $L$, $\Nm_{L/K}(I_L^S)$ 包含在 $\psi_{L/K}\colon I^S\to\Gal(L/K)$ 的核中.
\end{corollary}

\begin{proof}
在上图中取 $K'=L$.
\end{proof}

于是, 当 $L/K$ 为有限 Abel 扩张时, Artin 映射诱导同态
\[
\psi_{L/K}\colon I_K^S/\Nm(I_L^S)\to\Gal(L/K).
\]
群 $I^S/\Nm(I_L^S)$ 是无限的 (因为有无穷多个素理想不分裂), 故 $\psi_{L/K}$ 不可能单.

\subsection*{整体类域论的主要定理}

设 $S$ 为 $K$ 的素的有限集. 称同态 $\psi\colon I^S\to G$ \textbf{容许一个模}, 若存在模 $\mathfrak{m}$, $S(\mathfrak{m})\supset S$, 使得 $\psi(i(K_{\mathfrak{m},1}))=0$. 换言之, $\psi$ 容许一个模当且仅当它经过某个满足 $S(\mathfrak{m})\supset S$ 的 $C_{\mathfrak{m}}$ 分解.

\begin{theorem}[互反律]\label{thm:v.3.5}
设 $L$ 是 $K$ 的有限 Abel 扩张, $S$ 是 $K$ 中在 $L$ 分歧的素集. 则 Artin 映射 $\psi\colon I^S\to\Gal(L/K)$ 容许一个满足 $S(\mathfrak{m})=S$ 的模 $\mathfrak{m}$, 且它定义同构
\[
I_K^{S(\mathfrak{m})}/i(K_{\mathfrak{m},1})\cdot\Nm(I_L^{S(\mathfrak{m})})\to\Gal(L/K).
\]
\end{theorem}

如定理陈述中的模称为 $L$ 的\textbf{定义模} (defining modulus).

注意该定理并未蕴含 $K$ 有哪怕一个非平凡 Abel 扩张. 记 $I_K^{\mathfrak{m}}$ 为 $K$ 中 $S(\mathfrak{m})$-理想群, $I_L^{\mathfrak{m}}$ 为 $L$ 中 $S(\mathfrak{m})'$-理想群, 其中 $S(\mathfrak{m})'$ 包含位于 $S(\mathfrak{m})$ 的素之上的 $L$ 的素. 称 $I_K^{\mathfrak{m}}$ 的子群 $H$ 为\textbf{模 $\mathfrak{m}$ 的同余子群}, 若
\[
I_K^{\mathfrak{m}}\supset H\supset i(K_{\mathfrak{m},1}).
\]

\begin{theorem}[存在定理]\label{thm:v.3.6}
对模 $\mathfrak{m}$ 的每个同余子群 $H$, 存在有限 Abel 扩张 $L/K$, 以 $\mathfrak{m}$ 为模 (即在不整除 $\mathfrak{m}$ 的素处不分歧), 使得
\[
H=i(K_{\mathfrak{m},1})\cdot\Nm_{L/K}(I_L^{\mathfrak{m}}).\erratum{原书未指明 $L/K$ 以 $\mathfrak{m}$ 为模; 若无此条件, 存在定理并不能 (立即) 推出下文关于 Artin 映射诱导同构的断言. 依 Milne errata (Zhuo Zhang) 补正.}
\]
\end{theorem}

注意, 对如定理中的 $H$ 与 $L$, Artin 映射 $\psi_{L/K}$ 诱导同构
\[
I^{S(\mathfrak{m})}/H\to\Gal(L/K).
\]

特别地, 对每个模 $\mathfrak{m}$ 存在一个域 $L_{\mathfrak{m}}$, 称为\textbf{模 $\mathfrak{m}$ 的射线类域}\index{shexianleiyu@射线类域(ray class field)}, 使得 Artin 映射定义同构 $C_{\mathfrak{m}}\to\Gal(L_{\mathfrak{m}}/K)$. 对域 $L\subset L_{\mathfrak{m}}$, 令
\[
\Nm(C_{L,\mathfrak{m}})=i(K_{\mathfrak{m},1})\cdot\Nm(I_L^{\mathfrak{m}})\ \mathrm{mod}\ i(K_{\mathfrak{m},1}).
\]

\begin{corollary}\label{cor:v.3.7}
固定模 $\mathfrak{m}$. 则映射 $L\mapsto\Nm(C_{L,\mathfrak{m}})$ 是从含于 $L_{\mathfrak{m}}$ 的 $K$ 的 Abel 扩张之集到 $C_{\mathfrak{m}}$ 的子群之集的双射. 此外,
\begin{align*}
L_1\subset L_2&\iff\Nm(C_{L_1,\mathfrak{m}})\supset\Nm(C_{L_2,\mathfrak{m}});\\
\Nm(C_{L_1\cdot L_2,\mathfrak{m}})&=\Nm(C_{L_1,\mathfrak{m}})\cap\Nm(C_{L_2,\mathfrak{m}});\\
\Nm(C_{L_1\cap L_2,\mathfrak{m}})&=\Nm(C_{L_1,\mathfrak{m}})\cdot\Nm(C_{L_2,\mathfrak{m}}).
\end{align*}
\end{corollary}

\begin{remark}\label{rem:v.3.8}
设 $L/K$ 是 Galois 群为 $G$ 的 Abel 扩张. 据互反律, 存在模 $\mathfrak{m}$, 其支撑为在 $L$ 分歧的素集, 使 Artin 映射 $\psi_{L/K}\colon I^{S(\mathfrak{m})}\to G$ 在 $i(K_{\mathfrak{m},1})$ 上取值 $1$. 考虑定理 \ref{thm:v.1.7} 中的映射
\[
(\mathcal{O}_K/\mathfrak{p}^{m(\mathfrak{p})})^{\times}\hookrightarrow K_{\mathfrak{m}}/K_{\mathfrak{m},1}\xrightarrow{\ i\ }C_{\mathfrak{m}}\xrightarrow{\ \psi_{L/K}\ }G.
\]
显然存在最小整数 $f(\mathfrak{p})\le m(\mathfrak{p})$ 使该映射经过 $(\mathcal{O}_K/\mathfrak{p}^{f(\mathfrak{p})})^{\times}$ 分解. 模 $\mathfrak{f}(L/K)=\mathfrak{m}_\infty\prod\mathfrak{p}^{f(\mathfrak{p})}$ 便是使 $\psi_{L/K}$ 经过 $C_{\mathfrak{f}}$ 分解的最小模, 称为 $L/K$ 的\textbf{导子}\index{daozi@导子(conductor)}\footnote{德语为 F\"uhrer\,---\,在 1930 年代的德国, 公开谈论类域论可能有风险.}. 导子 $\mathfrak{f}(L/K)$ 恰被在 $L$ 中分歧的素整除.

含 $K$ 的射线类域 $L_{\mathfrak{m}}$ 的子域, 恰是导子整除 $\mathfrak{m}$ 的那些域.\footnote{设 $L$ 是导子整除 $\mathfrak{m}$ 的 Abel 扩张. 则在 $L$ 中分裂的素也在 $L_{\mathfrak{m}}$ 中分裂. 再由 Chebotarev 密度定理即得 $L\subset L_{\mathfrak{m}}$ (见下文 \ref{thm:v.3.25}).} $K$ 的每个 Abel 扩张都含于某个 $L_{\mathfrak{m}}$.
\end{remark}

下文第五节将用 idèle 重述定理\ref{thm:v.3.5} 与 \ref{thm:v.3.6}, 并在第 \ref{ch:VII} 章证明重述后的定理.

如下文所讨论的, Artin 映射是满射有一个相当简单的解析证明. 因此证明互反律的难点在于证明 Artin 映射容许一个导子, 且
\[
\bigl(I^{S(\mathfrak{m})}:i(K_{\mathfrak{m},1})\cdot\Nm(I_L^{S(\mathfrak{m})})\bigr)\le[L:K].
\]
为证明存在定理, 须对每个模构造射线类域. 遗憾的是, 我们不知道如何直接构造射线类域; 实际做法是构造足够多的扩张, 迫使定理成立.

\begin{example}\label{ex:v.3.9}
模 $\mathfrak{m}=1$ 的射线类群是理想类群, 相应的射线类域是 \textbf{Hilbert 类域}\index{hilbertleiyu@Hilbert 类域(Hilbert class field)}; 它是 $K$ 的极大 Abel 扩张, 在包括实素的所有素处不分歧 (即实素保持实). 例如, $\Q$ 的 Hilbert 类域是 $\Q$ 自身 (因 $\Q$ 类数为 $1$). $\Q[\sqrt{-5}]$ 的 Hilbert 类域是 $\Q[\sqrt{-1},\sqrt{5}]$\,---\,$2$ 与 $5$ 都在 $\Q[\sqrt{-5}]$ 中分歧, 但在 $\Q[\sqrt{-1}]$ 中只有 $2$ 分歧, 在 $\Q[\sqrt{5}]$ 中只有 $5$ 分歧, 由此推出 $\Q[\sqrt{-5}]$ 中整除 $2$ 与 $5$ 的素不在 $\Q[\sqrt{-1},\sqrt{5}]$ 中分歧.
\end{example}

\begin{example}\label{ex:v.3.10}
设 $m$ 为奇数或被 $4$ 整除的正整数. 模 $(m)$ 的射线类域是 $\Q[\zeta_m+\bar{\zeta}_m]$, 模 $\infty(m)$ 的射线类域是 $\Q[\zeta_m]$.\footnote{这由 \ref{ex:v.1.8}(c) 中对射线类群的描述以及 ANT, 8.18 中对 Frobenius 元的描述推出.} 于是互反律蕴含 Kronecker--Weber 定理: $\Q$ 的每个 Abel 扩张的导子整除某个 $\infty(m)$, 从而含于某个分圆域.
\end{example}

\begin{example}\label{ex:v.3.11}
设 $d$ 为无平方因子整数. 我们通过求使 $\Q[\sqrt{d}]\subset\Q[\zeta_m]$ 的最小整数 $m$ 来计算 $K=\Q[\sqrt{d}]$ 在 $\Q$ 上的导子.

先考虑奇素数 $p$. 此时 $\Gal(\Q[\zeta_p]/\Q)\simeq(\Z/p\Z)^{\times}$ 是 $p-1$ 阶循环群, 故有唯一的 $2$ 阶商群. 因此 $\Q[\zeta_p]$ 含唯一的二次域, 由于它只能在 $p$ 处分歧, 必为 $\Q[\sqrt{p^*}]$, 其中 $p^*=(-1)^{\frac{p-1}{2}}p$ (符号的选取使 $p^*\equiv 1\bmod 4$).

其次注意 $\zeta_8=(1+i)/\sqrt{2}$, 故 $\zeta_8+\bar{\zeta}_8=\sqrt{2}$. 因此 $\Q[\sqrt{2}]\subset\Q[\zeta_8]$ (事实上, $\Q[\sqrt{2}]$ 是 $\Q[\zeta_8]$ 的极大实子域, 且 $\Q[\zeta_8]=\Q[i,\sqrt{2}]$).

设 $n$ 为整除 $d$ 的奇素数之积 (故 $d=\pm n$ 或 $\pm 2n$). 我断言
\[
\Q[\sqrt{d}]\subset\Q[\zeta_n]\quad\text{若 } d\equiv 1\bmod 4,
\]
\[
\Q[\sqrt{d}]\subset\Q[\zeta_{4n}]\quad\text{若 } d\equiv 3\bmod 4,
\]
\[
\Q[\sqrt{d}]\subset\Q[\zeta_{8n}]\quad\text{若 } d\equiv 2\bmod 4,
\]
且每种情形这都是包含 $\Q[\sqrt{d}]$ 的最小分圆域. 例如, 注意 $d=p_1\cdots p_r$, $d\equiv 1\bmod 4$, 蕴含 $d=p_1^*\cdots p_r^*$, 故 $\Q[\sqrt{d}]\subset\Q[\zeta_n]$. 又注意若 $d$ 为偶数, 则 $\Q[\sqrt{d}]$ 不含于 $\Q[\zeta_{4n}]$, 否则 $\Q[\zeta_{4n}]$ 将包含 $i,\sqrt{d},\sqrt{d/2}$, 从而包含 $i,\sqrt{2}$ 与 $\zeta_8$.

结论: $\Q[\sqrt{d}]$ 的导子是 $|\Delta_{K/\Q}|$ 或 $\infty|\Delta_{K/\Q}|$, 视 $d>0$ 或 $d<0$ 而定, 此处 $\Delta_{K/\Q}$ 是 $K/\Q$ 的判别式.
\end{example}

\begin{remark}\label{rem:v.3.12}
对有限 Abel 扩张 $L/K$ 及满足 $S(\mathfrak{m})$ 等于 $K$ 中在 $L$ 分歧的素集的模 $\mathfrak{m}$, 令
\[
T(L/K,\mathfrak{m})=i(K_{\mathfrak{m},1})\cdot\Nm_{L/K}(I_L^{S(\mathfrak{m})})\subset I_K.
\]
Takagi 证明了: 当 $\mathfrak{m}$ 充分可除时 (事实上, 对所有被 $L/K$ 的导子 $\mathfrak{f}$ 整除的 $\mathfrak{m}$), 群 $T(L/K,\mathfrak{m})$ 不依赖于 $\mathfrak{m}$, 且 $I_K^{S(\mathfrak{m})}/T(L/K,\mathfrak{m})\approx\Gal(L/K)$. 因此当 $\mathfrak{f}|\mathfrak{m}$ 时, $T(L/K,\mathfrak{m})$ 常被称为 $L/K$ 的 \textbf{Takagi 群}. Artin 证明了 Takagi 群正是映射 $\mathfrak{a}\mapsto(\mathfrak{a},L/K)\colon I_K^{S(\mathfrak{m})}\to\Gal(L/K)$ 的核.
\end{remark}

\begin{exercise}\label{exc:v.3.13}
利用二次互反律求使 $i(\Q_{m,1})$ 含于 Artin 映射之核的最小 $m$, 以此计算 $\Q[\sqrt{d}]/\Q$ 的导子 (对照习题 \ref{ex:0.11}).
\end{exercise}

对应于同余子群 $H$ 的域 $L$ 称为 $H$ 的\textbf{类域}, 这也是本学科名称的由来. 注意: 对不整除 $L/K$ 的导子的 $K$ 的素 $\mathfrak{p}$, 位于 $\mathfrak{p}$ 之上的素的剩余类次数 $f(\mathfrak{P}/\mathfrak{p})$ 等于 $\mathfrak{p}$ 在 $I^{\mathfrak{m}}/H$ 中的阶 (因为这是 $\Gal(k(\mathfrak{P})/k(\mathfrak{p}))$ 中 Frobenius 元的阶). 于是我们得到了用 $K$ 的理想结构对 $K$ 的 Abel 扩张的分类, 并且对每个 Abel 扩张都知道 $K$ 中素的分解规律.

\begin{exercise}\label{exc:v.3.14}
验证下表的最后一行:
\begin{center}
\begin{tabular}{|c|c|c|c|c|c|}
\hline
判别式 & $-15$ & $-20$ & $-23$ & $-24$ & $-31$\\
\hline
类数 & $2$ & $2$ & $3$ & $2$ & $3$\\
\hline
Hilbert 类域 & $X^2+3$ & $X^2+1$ & $X^3-X-1$ & $X^2+3$ & $X^3+X-1$\\
\hline
\end{tabular}
\end{center}
第一行列出前五个类数不为 $1$ 的虚二次域的判别式, 第二行是它们的类数, 最后一行是 Hilbert 类域生成元的极小多项式. (注意对全虚域, 类数与狭窄类数相等.)

[例如, 判别式为 $-24$ 的二次域是 $\Q[\sqrt{-6}]$, 其 Hilbert 类域是 $L=\Q[\sqrt{-6},\sqrt{-3}]=\Q[\sqrt{2},\sqrt{-3}]$. 于是 $L=\Q[\alpha]$, $\alpha=\sqrt{2}+\sqrt{-3}$, 其极小多项式为 $X^4+2X^2+5$. 然而 $\Z[\alpha]$ 并不等于 $L$ 的整数环; 对照 \mo{122765}.]
\end{exercise}

\begin{exercise}\label{exc:v.3.15}
本习题解释当我们忽略 $K$ 的素理想的有限集 $S$ 时会发生什么. 设 $\mathfrak{m}$ 是 $K$ 的满足 $S(\mathfrak{m})\cap S=\emptyset$ 的模, $H$ 是 $I^{S\cup S(\mathfrak{m})}$ 的包含 $i(K_{\mathfrak{m},1})$ 的子群. 称 $K$ 的扩张 $L$ 为 $H$ 的 \textbf{$S$-类域}, 若

(a) $L$ 是 $K$ 的有限 Abel 扩张, 且 $S$ 中的素理想在 $L$ 中完全分裂;

(b) $m(\mathfrak{p})=0\Rightarrow\mathfrak{p}$ 不在 $L$ 中分歧;

(c) 不在 $S\cup S(\mathfrak{m})$ 中且在 $L$ 中分裂的素理想恰是 $H$ 中的那些.

证明: 对如上每个群 $H$ 都存在 $S$-类域 $L$, 它是唯一的, 且 $I^{S\cup S(\mathfrak{m})}/H\simeq\Gal(L/K)$; 此外, 每个满足 (a) 的域 $L$ 都是某个 $H$ 的 $S$-类域.

提示: 证明 $I^{S\cup S(\mathfrak{m})}/i(K_{\mathfrak{m},1})\simeq I^{S(\mathfrak{m})}/\langle S\rangle\cdot i(K_{\mathfrak{m},1})$, 其中 $\langle S\rangle$ 是由 $S$ 中的素生成的 $I^{S(\mathfrak{m})}$ 的子群.
\end{exercise}

\subsection*{范限制定理}

在对 $K$ 的 Abel 扩张的分类中, 我们对 $L$ 附上群 $H=i(K_{\mathfrak{m},1})\cdot\Nm(I_L^{S(\mathfrak{m})})$, 其中 $\mathfrak{m}$ 是充分大到成为定义模的模 (此时 $(I^{S(\mathfrak{m})}:H)=[L:K]$). 人们或许希望类似的办法对非 Abel 扩张也有效, 但下面的定理表明并非如此.

\begin{theorem}[范限制定理]\label{thm:v.3.16}
设 $L$ 是 $K$ 的有限扩张, $L'/K$ 是 $L/K$ 的极大 Abel 子扩张. 对 $L'/K$ 的每个定义模 $\mathfrak{m}$,
\[
i(K_{\mathfrak{m},1})\cdot\Nm_{L/K}(I_L^{S(\mathfrak{m})})=i(K_{\mathfrak{m},1})\cdot\Nm_{L'/K}(I_{L'}^{S(\mathfrak{m})}).
\]
\end{theorem}

例如, 若 $L$ 是 $K$ 的在 $K$ 上非 Galois 的三次扩张, 则 $L'=K$, 故
\[
i(K_{\mathfrak{m},1})\cdot\Nm_{L/K}(I_L^{S(\mathfrak{m})})=I_K^{S(\mathfrak{m})}.
\]

范限制定理表明: 对非 Abel 扩张 $L/K$, $\mathrm{Spl}(L/K)$ 不能用同余条件描述. 该定理的证明见例如 \cite{GrantLeitzel1969}.

\subsection*{主理想定理}

\begin{quote}
\itshape ... \`a ma honte, je ne suis par arriv\'e \`a retrouver ce \guillemotleft corollaire\guillemotright\ que tous les id\'eaux de $K$ deviennent principaux dans la plus grande extension ab\'elienne non ramifi\'ee dans le fini. Si \c ca s'explique en deux mots, je t'en serais fort reconnaissant.
\end{quote}
\hspace*{\fill}Grothendieck, 致 Serre 的信, 1956.9.19\footnote{\guillemotleft\,\ldots\ 令我羞愧的是, 我一直没能重新找出这个\guillemotleft 推论\guillemotright: $K$ 的所有理想在有限素处不分歧的极大 Abel 扩张中都变为主理想. 若能用两句话解释, 我将不胜感激.\guillemotright\ Serre 回信道: \guillemotleft 随信附上一篇关于 \guillemotleft Hauptidealsatz\guillemotright\ 的小文章, 解释该定理如何化归为一个 (实际上非常神秘的) 群论定理. 事实上, 这正是 Artin 本人在其论文 (Abh.~Hamburg, 大约第 7--10 卷) 中给出的化归; 如果你能找到该定理的一个漂亮的上同调证明, 那就再好不过了, 但迄今为止人人都被卡住了.\guillemotright\ 另见 Serre 信件的尾注 (\cite{GrothendieckSerre2001}).}

下面的定理由 Hilbert 在 1900 年前后猜想.

\begin{theorem}[主理想定理]\label{thm:v.3.17}
$K$ 的每个理想在 $K$ 的 Hilbert 类域中都变为主理想.
\end{theorem}

我解释一下证明的想法. 回顾: 对群 $G$, $G$ 的\textbf{换位子群} (或\textbf{导群}) $G'$ 是由换位子 $ghg^{-1}h^{-1}$, $g,h\in G$, 生成的子群. 商群 $G^{\mathrm{ab}}=G/G'$ 是 Abel 群, 且是 $G$ 的极大 Abel 商. 若 $L$ 是 $K$ 的 Galois 扩张, Galois 群为 $G$, 则 $L^{G'}$ 是 $K$ 的 Galois 群为 $G^{\mathrm{ab}}$ 的 Abel 扩张, 且是 $L$ 中包含 $K$ 的极大 Abel 扩张.

设有域
\[
L\supset K'\supset K,
\]
其中 $L$ 在 $K$ 上 Galois (不必 Abel). 对 $K$ 的素的每个有限集 $S$, $\mathfrak{a}\mapsto\mathfrak{a}\mathcal{O}_{K'}$ 是同态 $I_K^S\to I_{K'}^S$. 考虑:
\[
\begin{tikzcd}
I_K^S\ar[r, "\psi_{L/K}"]\ar[d, "\text{典范}"'] & \Gal(L/K)^{\mathrm{ab}}\ar[d, "?"]\\
I_{K'}^S\ar[r, "\psi_{L/K'}"'] & \Gal(L/K')^{\mathrm{ab}}
\end{tikzcd}
\]
使图交换的映射 ``?'' 是什么?

在描述它之前, 需要解释群论中的一个构造. 设 $H$ 是群 $G$ 的有限指数子群, 把 $G$ 写成陪集的不交并
\[
G=Hg_1\cup Hg_2\cup\cdots\cup Hg_n.
\]
对 $g\in G$, 当 $g\in Hg_i$ 时令 $\varphi(g)=g_i$, 并定义
\[
V(g)=\prod_{i=1}^{n}g_i g\,\varphi(g_i g)^{-1}\ \mathrm{mod}\ H',
\]
其中 $H'$ 是 $H$ 的换位子群.

\begin{proposition}\label{prop:v.3.18}
映射 $g\mapsto V(g)$ 是同态 $G\to H/H'$, 且不依赖陪集代表元 $g_i$ 的选取.
\end{proposition}

\begin{proof}
验证是直接的\,---\,例如见 \cite{Hall1959}, 14.2.1.
\end{proof}

于是, 只要有了群 $G$ 及其有限指数子群 $H$, 就有确定的同态
\[
V\colon G^{\mathrm{ab}}\to H^{\mathrm{ab}},
\]
称为 \textbf{Verlagerung} (或\textbf{转移}) 映射\index{zhuanyiying@转移映射(transfer homomorphism)}\index{verlagerung@Verlagerung(Verlagerung)}. 粗略地说, 我们想定义一个从 $G$ 到 $H$ 的\guillemotleft 范\guillemotright, 由于群不交换这行不通, 但过渡到相应的 Abel 群后就行得通了.

在上图的情形, $\Gal(L/K')$ 是 $\Gal(L/K)$ 的子群, 从而 Verlagerung 是同态
\[
V\colon\Gal(L/K)^{\mathrm{ab}}\to\Gal(L/K')^{\mathrm{ab}}.
\]
Emil Artin 证明了这正是使上图交换的映射 (对照 II, \ref{prop:ii.3.2}b). 考虑域
\[
K''\supset K'\supset K,
\]
其中 $K'$ 是 $K$ 的 Hilbert 类域, $K''$ 是 $K'$ 的 Hilbert 类域. 则

(a) $K''$ 在 $K$ 上正规 (因为 $K''$ 的任何共轭仍是 $K'$ 的 Abel 非分歧扩张, 从而含于 $K''$);

(b) $K'$ 是 $K''$ 中包含 $K$ 的极大 Abel 扩张 (因为 $K''$ 中包含 $K$ 的每个 Abel 扩张在 $K$ 上非分歧, 从而含于 $K'$).

由 (b) 得 $\Gal(K''/K)^{\mathrm{ab}}=\Gal(K'/K)$. 因此取 $L=K''$ 时, 前图成为
\[
\begin{tikzcd}
C_K\ar[r, "\simeq"]\ar[d, "\text{典范}"'] & \Gal(K'/K)\ar[d, "\mathrm{Ver}"]\\
C_{K'}\ar[r, "\simeq"'] & \Gal(K''/K')
\end{tikzcd}
\]
其中 $C_K$ 与 $C_{K'}$ 是 $K$ 与 $K'$ 的类群. 设 $G=\Gal(K''/K)$, $H=\Gal(K''/K')$. 由 (b), $H$ 是 $G$ 的换位子群. 下面的定理 (由 Emil Artin 猜想) 表明此情形下 $V$ 为零, 从而典范映射 $C_K\to C_{K'}$ 为零, 即 $K$ 的每个理想在 $K'$ 中变为主理想.

\begin{theorem}\label{thm:v.3.19}
设 $G$ 是有限群, $H$ 是其换位子群; 则
\[
V\colon G^{\mathrm{ab}}\to H^{\mathrm{ab}}
\]
为零.
\end{theorem}

\begin{proof}
这是群论的定理, 也称为主理想定理. 由 Furtw\"angler 于 1930 年证明. 简单证明见: Witt, Proc.~International Congress of Mathematicians, Amsterdam, 1954, Vol.~2, pp.~71--73. 另见 Zassenhaus, Theory of Groups, V, 定理 12, 或 \cite{ArtinTate1961}, 第 XII 章.
\end{proof}

\begin{remark}\label{rem:v.3.20}
事实上容易看出: 存在 $K$ 的一个次数整除类数 $h$ 的扩张 $L$, 使 $K$ 的每个理想在 $L$ 中变为主理想: 把 $K$ 的类群写成循环群直和; 对每个直和项取生成元 $a_i$, 设 $a_i$ 在类群中的阶为 $h_i$; 写 $a_i^{h_i}=(a_i)$, 并定义 $L$ 为对每个 $i$ 向 $K$ 添加 $a_i$ 的 $h_i$ 次根所得的域.

然而如此构造的域通常并非 $K$ 的 Hilbert 类域\,---\,它甚至不必在 $K$ 上 Galois. 可能存在次数 $<h$ 的 $K$ 的域, 使 $K$ 的每个理想在其中变为主理想.
\end{remark}

\begin{remark}\label{rem:v.3.21}
注意主理想定理说的是: $K$ 中元素的\guillemotleft 理想\guillemotright\ 分解在 Hilbert 类域 $L$ 中成为实际的分解: 设 $a\in K$, $(a)=\mathfrak{p}_1^{r_1}\cdots\mathfrak{p}_s^{r_s}$; 则 $\mathfrak{p}_i\mathcal{O}_L=(\pi_i)$ 对某个 $\pi_i\in\mathcal{O}_L$ 成立, 于是在 $\mathcal{O}_L$ 中 $a=u\pi_1^{r_1}\cdots\pi_s^{r_s}$.
\end{remark}

\begin{remark}\label{rem:v.3.22}
主理想定理当然并不蕴含 $K$ 的 Hilbert 类域 $K'$ 的每个理想都是主理想, 因为并非 $K'$ 的每个理想都属于同态 $I_K\to I_{K'}$ 的像. 可以继续作 $K'$ 的 Hilbert 类域 $K''$, 等等, 得到塔
\[
K\subset K'\subset K''\subset\cdots\subset K^{(n)}\subset\cdots
\]
其中 $K^{(n+1)}$ 是 $K^{(n)}$ 的 Hilbert 类域. 对 $K^{\mathrm{al}}$ 的每个自同构 $\tau$, $\tau K^{(n+1)}$ 显然是 $\tau K^{(n)}$ 的极大 Abel 非分歧扩张, 这蕴含 $\tau K'=K'$, $\tau K''=K''$, 等等. 特别地, $K^{(n)}$ 在 $K$ 上 Galois. \textbf{类域塔问题} (Hasse 于 1925 年提出) 问: 该塔是否总是有限, 从而终止于一个类数为一的域? Golod 与 Shafarevich 于 1964 年证明答案是否定的 (见 \cite{Roquette1967}). 例如 $\Q[\sqrt{-2\cdot 3\cdot 5\cdot 7\cdot 11\cdot 13}]$ 有无限类域塔. 事实上, 只要 $d$ 的素因子多于 $8$ 个, $\Q[\sqrt{d}]$ 就有无限类域塔.
\end{remark}

\subsection*{Chebotarev 密度定理}

设 $L$ 是 $K$ 的 Galois 扩张, Galois 群为 $G$. 回顾: 对在 $L$ 中不分歧的 $K$ 的每个素理想 $\mathfrak{p}$,
\[
(\mathfrak{p},L/K)\overset{\mathrm{def}}{=}\{(\mathfrak{P},L/K)\mid\mathfrak{P}|\mathfrak{p}\}
\]
是 $G$ 中的一个共轭类.

\begin{theorem}[Chebotarev 密度定理]\label{thm:v.3.23}
设 $L/K$ 是 Galois 群为 $G$ 的数域有限扩张, $C$ 是 $G$ 中的一个共轭类. 则满足 $(\mathfrak{p},L/K)=C$ 的 $K$ 的素理想之集在 $K$ 的所有素理想之集中的密度为 $|C|/|G|$. 特别地, 若 $G$ 为 Abel 群, 则对固定的 $\tau\in G$, 满足 $(\mathfrak{p},L/K)=\tau$ 的 $K$ 的素理想 $\mathfrak{p}$ 之集的密度为 $1/(G:1)$.
\end{theorem}

\begin{proof}
对 Abel 扩张 $L/K$, 这由定理 \ref{thm:v.2.5} 与定理 \ref{thm:v.3.5} 推出: 后者说映射 $\mathfrak{p}\mapsto(\mathfrak{p},L/K)$ 对某个模 $\mathfrak{m}$ 诱导满同态 $C_{\mathfrak{m}}\to\Gal(L/K)$, 前者说素理想在 $C_{\mathfrak{m}}$ 的类中均匀分布. 非 Abel 情形由 Abel 情形经一个巧妙的论证导出\,---\,见第 \ref{ch:VIII} 章.
\end{proof}

\begin{corollary}\label{cor:v.3.24}
若多项式 $f(X)\in K[X]$ 对 $K$ 中除有限多个素理想外的所有素理想 $\mathfrak{p}$ 都模 $\mathfrak{p}$ 分解为一次因子, 则它在 $K[X]$ 中分裂.
\end{corollary}

\begin{proof}
对 $f(X)$ 的分裂域应用定理.
\end{proof}

对数域的有限扩张 $L/K$ 及 $K$ 的素的有限集 $S$, 设 $\mathrm{Spl}_S(L/K)$ 为 $K$ 中不在 $S$ 内且在 $L$ 中分裂的素之集.

\begin{theorem}\label{thm:v.3.25}
若 $L$ 与 $M$ 都是 $K$ 的 Galois 扩张, 则
\[
L\subset M\iff\mathrm{Spl}_S(L/K)\supset\mathrm{Spl}_S(M/K).
\]
从而
\[
L=M\iff\mathrm{Spl}_S(L/K)=\mathrm{Spl}_S(M/K).
\]
\end{theorem}

\begin{proof}
作为 (\hyperlink{item:v.1.12}{1.12}) 的推论,
\[
\mathrm{Spl}_S(LM/K)=\mathrm{Spl}_S(L/K)\cap\mathrm{Spl}_S(M/K).
\]
于是
\[
\mathrm{Spl}_S(L/K)\supset\mathrm{Spl}_S(M/K)\Rightarrow\mathrm{Spl}_S(LM/K)=\mathrm{Spl}_S(M/K),
\]
由 Chebotarev 密度定理, 这蕴含
\[
[LM:K]=[M:K],
\]
故 $L\subset M$. 反向蕴含是显然的.
\end{proof}

\begin{remark}\label{rem:v.3.26}
\begin{enumerate}[(a)]
\item 没有 Galois 假设时定理 \ref{thm:v.3.25} 不成立 (见 Cassels and Fr\"ohlich 1967, p.~363).
\item 在定理 \ref{thm:v.3.25} 的陈述中, $S$ 可换成密度为 $0$ 的任意素集.
\item 设 $f(X)$ 是 $K[X]$ 中的不可约多项式. 若 $f(X)$ 对几乎所有素理想 $\mathfrak{p}$ 都模 $\mathfrak{p}$ 有根, 则 $f(X)$ 在 $K$ 中有根 (Cassels and Fr\"ohlich 1967, p.~363, 6.2).
\item Chebotarev 用 Dirichlet (解析) 密度证明了其定理. Artin 指出它对自然密度 (见下文) 也应成立, 而事实上确实如此.
\item 给定数域间的有限 Galois 扩张 $K/\Q$ 及共轭类 $C$, Chebotarev 定理说存在 $K$ 的素理想 $\mathfrak{p}$, 它在 $\Q$ 上不分歧, 使 $(\mathfrak{p},L/K)=C$. 有趣的是 $\Nm\mathfrak{p}$ 需要有多大.  关于此问题的近作, 见 Ahn, Jeoung-Hwan; Kwon, Soun-Hi.~An explicit upper bound for the least prime ideal in the Chebotarev density theorem.~Ann.~Inst.~Fourier (Grenoble) 69 (2019), no.~3, 1411--1458.
\end{enumerate}
\end{remark}

\subsection*{导子--判别式公式}

两个 Dirichlet 特征 $\chi\colon I^S\to\CC^{\times}$ 与 $\chi'\colon I^{S'}\to\CC^{\times}$ 称为\textbf{相合的} (cotrained), 若它们在某个 $I^{S''}$, $S''\supset S\cup S'$, 上一致. 这是一个等价关系. 每个等价类中存在唯一的 $\chi$, 其 $S$ 最小\,---\,这样的 $\chi$ 称为\textbf{本原的} (primitive).

设 $\chi_1$ 是与 $\chi$ 等价的本原特征. 使 $\chi_1$ 在 $i(K_{\mathfrak{m},1})$ 上取零值的最小模 $\mathfrak{m}$ 称为 $\chi$ 的\textbf{导子} $\mathfrak{f}(\chi)$. 令 $\mathfrak{f}(\chi)=\mathfrak{f}_\infty(\chi)\mathfrak{f}_0(\chi)$, 其中 $\mathfrak{f}_\infty$ 与 $\mathfrak{f}_0$ 分别只被无限素与有限素整除.

\begin{theorem}[F\"uhrerdiskriminantenproduktformel]\label{thm:v.3.27}
对 Galois 群为 $G$ 的数域的每个有限 Abel 扩张 $L/K$,
\begin{align*}
\mathrm{disc}(L/K)&=\prod_{\chi\in G^{\vee}}\mathfrak{f}_0(\chi\circ\psi_{L/K}),\qquad G^{\vee}\overset{\mathrm{def}}{=}\mathrm{Hom}(G,\CC^{\times});\\
\mathfrak{f}(L/K)&=\mathrm{lcm}_{\chi}\bigl(\mathfrak{f}(\chi\circ\psi_{L/K})\bigr).
\end{align*}
\end{theorem}

显然 $\bigcap\mathrm{Ker}(\chi\colon G\to\CC^{\times})=0$, 由此得第二个等式. 第一个等式的证明从略\,---\,它实际上是关于局部域的陈述.

\begin{remark}\label{rem:v.3.28}
设 $H$ 是特征 $\chi$ 的核. 则 $\mathfrak{f}(\chi)=\mathfrak{f}(L^H/K)$. 例如, 若 $\chi$ 单, 则 $\mathfrak{f}(\chi)=\mathfrak{f}(L/K)$.
\end{remark}

\begin{example}\label{ex:v.3.29}
设 $L=\Q[\sqrt{d}]$. 则 $G\simeq\Z/2\Z$, 只有两个特征 $G\to\CC^{\times}$: 平凡特征 $\chi_0$ 与单射特征 $\chi_1$. 因此定理说 $\Delta_{K/\Q}=\mathfrak{f}_0(\chi_1)=\mathfrak{f}_0(K/\Q)$, 与例 \ref{ex:v.3.11} 中所证一致.
\end{example}

\begin{example}\label{ex:v.3.30}
设 $L=\Q[\zeta_p]$, $p$ 为奇素数. 则 $G\simeq(\Z/p\Z)^{\times}$ 是 $p-1$ 阶循环群, 故有 $p-2$ 个非平凡特征. 若 $\chi$ 非平凡, 则 $\mathfrak{f}(\chi)|\infty(p)$, 但 $\mathfrak{f}(\chi)=1$ 或 $\infty$ 不可能 (因为那将蕴含 $\chi$ 平凡). 因此当 $\chi$ 非平凡时 $\mathfrak{f}(\chi)=(p)$ 或 $\infty(p)$, 于是导子--判别式公式 (正确地) 给出 $\Delta_{L/\Q}=\pm p^{p-2}$.
\end{example}

\begin{exercise}\label{exc:v.3.31}
对扩张 $\Q[\zeta_{p^2}]/\Q$ 验证导子--判别式公式.
\end{exercise}

\begin{remark}\label{rem:v.3.32}
设 $L/K$ 是 Galois 群为 $G$ (不必 Abel) 的数域有限扩张. 对 $G$ 在有限维向量空间 $V$ 上的表示 $\rho\colon G\to\mathrm{GL}(V)$, Artin 附带一个 Dirichlet $L$-级数 $L(s,\rho)$ (见引言). 这些 Artin $L$-级数的解析性质至今仍未完全弄清.

当 $G$ 交换时, Artin 映射 $I^S\to G$ 把 $G$ 的特征与 Dirichlet/Weber 特征等同, 从而把 Artin $L$-级数与 Dirichlet $L$-级数等同. 这正是 Artin 寻求该映射的动机 (而不像今天看起来那么自然的那样, 是为了在早已知道抽象同构的群 $C_{\mathfrak{m}}$ 与 $G$ 之间构造典范同构).

Langlands 纲领的一部分正是 Dirichlet $L$-级数与 Abel Artin $L$-级数之间这一对应的巨大推广.
\end{remark}

\subsection*{互反律与幂互反律}

设 $K$ 含本原 $n$ 次单位根, $a\in K$. 若 $\sqrt[n]{a}$ 是 $X^n-a$ 的一个根, 则其余根形如 $\zeta\sqrt[n]{a}$, $\zeta$ 为 $n$ 次单位根. 因此 $L\overset{\mathrm{def}}{=}K[\sqrt[n]{a}]$ 在 $K$ 上 Galois, 且 $\sigma\sqrt[n]{a}=\zeta\sqrt[n]{a}$ 对某个 $n$ 次单位根 $\zeta$ 成立.

若 $\mathfrak{p}$ 是 $K$ 的与 $n$ 和 $a$ 互素的素理想, 则 $\mathfrak{p}$ 在 $L$ 中不分歧, 且可由公式
\[
(\mathfrak{p},L/K)(\sqrt[n]{a})=\Bigl(\frac{a}{\mathfrak{p}}\Bigr)_{\!n}\sqrt[n]{a}
\]
定义一个 $n$ 次单位根 $\bigl(\frac{a}{\mathfrak{p}}\bigr)_n$.

可以证明
\[
\Bigl(\frac{a}{\mathfrak{p}}\Bigr)_{\!n}=1\iff a\ \text{模}\ \mathfrak{p}\ \text{是}\ n\ \text{次幂},
\]
故 $\bigl(\frac{a}{\mathfrak{p}}\bigr)_n$ 推广了二次剩余符号. 因此 $\bigl(\frac{a}{\mathfrak{p}}\bigr)_n$ 称为\textbf{幂剩余符号}\index{mishengyufuhao@幂剩余符号(power residue symbol)}. Artin 互反律蕴含这些符号的所有已知的互反律, 因而正如 Artin 指出的, 可视为它们到没有单位根的域的推广. 我们将在第 \ref{ch:VIII} 章解释这一点.

\subsection*{一个初等的未解决问题}

设 $K$ 为数域, $S$ 为 $K$ 的素理想的非空有限集, $p$ 为不被 $S$ 中任何素整除的素数. 是否存在域列 $\ldots,L_n,L_{n+1},\ldots$ 使得

(a) $L_n$ 在 $S$ 之外的素处不分歧;

(b) $p^n|[L_n:K]$?

一个关键情形 (答案未知) 是 $K=\Q$, $S=\{l\}$.

更明确地 (也稍难些): 固定 $\Q$ 中互异的素数 $p\ne l$. 是否存在首一不可约多项式列 $f_n(X)\in\Z[X]$ 使得

(a) $\mathrm{disc}(f_n(X))$ 不被 $l$ 以外的任何素数整除;

(b) $p^n|\deg f_n(X)$?

见 \cite{Milne2006}, pp.~48--49, 了解这个问题出现于何处. 它当然非常困难.\footnote{设 $K$ 为数域, $S$ 为 $K$ 的素集, $K^S$ 为 $K$ 在 $S$ 外不分歧的极大扩张. 设 $P$ 是使 $\ell^\infty\mid [K_S:K]$ 的素数 $\ell$ 之集. 现已知: 只要 $S$ 包含 $K$ 中位于 $\Q$ 的至少两个素数之上的全部素理想, $P$ 就包含所有素数. 见 Chenevier and Clozel, JAMS 22 (2009), 467--519, Cor.~5.2.}

\subsection*{显式整体类域论: Kronecker 的青春之梦与 Hilbert 第十二问题}

与局部类域论不同, 整体类域论 (一般而言) 并不给出数域 $K$ 的 Abel 扩张的显式构造.

Gauss 知道 $\Q$ 的分圆扩张是 Abel 的. 1840 年代末, Kronecker 萌生了一个想法: 分圆域及其子域穷尽了 $\Q$ 的 Abel 扩张; 而且虚二次数域 $E$ 的每个 Abel 扩张都含于向 $E$ 添加单位根与模函数 $j$ 的某些特殊值所得的扩张中. 许多年后他称这一想法为\guillemotleft 我青春时期最珍爱的梦\guillemotright\ (mein liebster Jugendtraum).

更精确地说, Kronecker 之梦\footnote{关于 Kronecker 的想法及其工作的细致梳理, 见 \cite{Schappacher1998}.} 是: $\Q$ 的每个 Abel 扩张都含于向 $\Q$ 添加函数 $e^{2\pi iz}$ 对所有 $z\in\Q^{\times}$ 的值所得的域; 虚二次域 $K$ 的每个 Abel 扩张都含于向 $K$ 添加函数 $j(z)$ 对所有 $z\in K^{\times}$ 的值所得的域.

后来 Hilbert 把它列入其著名问题的第十二个: 对任意数域 $K$, 找出对 $K$ 起着指数函数对 $\Q$、模函数 $j$ 对虚二次域所起作用的函数.\footnote{见 Mathematical Developments arising from Hilbert's Problems, Proc.~Symp.~Pure Math.~XXVIII, Part 1, AMS, 1976, pp.~18--20.} Kronecker 之梦的第一部分\,---\,$\Q$ 的每个 Abel 扩张都是分圆扩张的子域\,---\,由 Weber (1886, 1899, 1907, 1911) 与 Hilbert (1896) 证明.

上述\guillemotleft Kronecker 之梦\guillemotright\ 的表述并不完全正确. 设 $K$ 为虚二次域. 可写 $\mathcal{O}_K=\Z+\Z\tau$, $\Im(\tau)>0$. 已知 $K[j(\tau)]$ 是 $K$ 的 Hilbert 类域. 现向 $K$ 添加所有单位根及所有值 $j(\tau)$, $\tau\in K$, $\Im(\tau)>0$. 所得的域 $K'$ 在 $K$ 上 Abel, 且 $[K^{\mathrm{ab}}:K']$ 是若干 $2$ 阶群之积. 要得到整个 $K^{\mathrm{ab}}$, 还须添加其他椭圆函数的特殊值. 这些结论由 Weber (1908) 与 Feuter (1914) 部分证明, 由 Takagi (1920) 完全证明.

从现代观点看, 椭圆模函数的特殊值与带复乘的椭圆曲线的算术相关联, 正是关于后者的结果使人得以证明前者生成虚二次域的 Abel 扩张.

自 1950 年代末 Taniyama、Shimura 与 Weil 的工作开始, 椭圆曲线与椭圆模曲线的理论已被成功地推广到高维. 在该理论中, 由虚二次域给出复乘的椭圆曲线被由 CM 域给出复乘的 Abel 簇所取代\,---\,CM 域即全实域 $F$ 的二次全虚扩张 $K$, 而椭圆模函数被自守函数所取代.

在哲学层面, 人们预期: 除 $\Q$ 外, 不能通过添加自守函数的特殊值得到全实域的 Abel 扩张. 然而, 粗略地说, 已知对 CM 域 $K$ 确能得到与该限制相容的最大可能的 Abel 扩张.

更精确地, 设 $K$ 为 CM 域, $F$ 为 $K$ 的极大全实子域. 则 $G\overset{\mathrm{def}}{=}\Gal(\Q^{\mathrm{al}}/K)$ 是 $G'\overset{\mathrm{def}}{=}\Gal(\Q^{\mathrm{al}}/F)$ 中指数为 $2$ 的子群, 相应的 Verlagerung 是同态 $V\colon G'^{\mathrm{ab}}\to G^{\mathrm{ab}}$. 此情形下 $V$ 有非常简单的描述.

\begin{theorem}\label{thm:v.3.33}
设 $K$ 为 CM 域, $F$ 为 $K$ 的全实子域, $[K:F]=2$. 设 $H$ 是 Verlagerung 映射
\[
\Gal(F^{\mathrm{al}}/F)^{\mathrm{ab}}\to\Gal(K^{\mathrm{al}}/K)^{\mathrm{ab}}
\]
的像. 则向 $K$ 添加所有定义在有理权 Shimura 簇的典范模型上的自守函数的特殊值所得的 $K$ 的扩张是 $(K^{\mathrm{ab}})^H\cdot\Q^{\mathrm{ab}}$.
\end{theorem}

\begin{proof}
见 \cite{Wei1993}, \cite{Wei1994}.
\end{proof}

\subsection*{注}

同余群与 $K$ 的 Abel 扩张之间的联系在 Artin 定义其映射之前就已被知晓. 它在约 1870--1920 年间缓慢浮现, 主要贡献者是 Kronecker, Weber, Hilbert 与 Takagi. Chebotarev 于 (1926) 年证明了其定理 (Frobenius 更早证明了一个较不精确的结果), 而 Artin 于 1927 年定义了他的映射并证明它给出同构. (更早, 人们已知 $I^m/H\approx\Gal(L/K)$, 但并无典范同构.) 证明这些形式上纯代数的主要定理竟需要分析\,---\,以 Chebotarev (或 Frobenius) 定理的形态出现\,---\,被视为一个缺陷; 经过大量努力, Chevalley 于 1940 年成功地给出了主要定理的纯代数证明. (难点在于证明: 若 $L/K$ 是素次数 $p$ 的数域 Abel 扩张, 则 $K$ 中至少有一个素在 $L$ 中既不分裂也不分歧.\footnote{类域论对 $k$ 有限的 $k(X)$ 的有限扩张也成立; 局部类域论对剩余类域 $k$ 仅为拟有限的局部域, 即满足 $\Gal(k^{\mathrm{al}}/k)\simeq\hat{\Z}$ 的完全域, 也成立. 当人们试图对 $k$ 拟有限的 $k(X)$ 的有限扩张建立整体类域论时, 发现一切都行得通, 只是对某些 $k$ 确实存在每个素都分裂的非平凡 Abel 扩张. 见本书 Arithmetic Duality Theorems, 第 I 章附录.}) 他还引入了 idèle, 使直接对无限扩张陈述类域论成为可能. 群上同调 (至少 2-上循环等) 从 1920 年代起就在类域论中被隐含地使用, 但在 1950 年代被 Nakayama, Hochschild 与 Tate 系统地使用. 1951/52 年, Artin 与 Tate 在一个影响深远的研究班中给出了类域论的纯代数且极富上同调色彩的处理. 此后, 我们对局部类域论的理解有了重要改进 (主要归功于 Lubin 与 Tate). 非 Abel 类域论是 Langlands 纲领的一部分, 该纲领是一系列相互交织的庞大猜想, 目前已取得一些进展, 尤其是在局部情形与函数域情形. 对更一般的域 ($\Q$ 或 $\F_p$ 上有限超越次数的域), Bloch, Kato, Saito 等人建立了相当令人满意的 Abel 类域论. 它使用了代数 $K$-理论 (综述见 \cite{Raskind1995}, 近期结果见 \cite{Wiesend2007}).

\section{Idèle}\label{sec:V.4}

定理\ref{thm:v.3.5} 与 \ref{thm:v.3.6} 表明: 对任意数域 $K$, 存在典范同构 $\varprojlim_{\mathfrak{m}}C_{\mathfrak{m}}\to\Gal(K^{\mathrm{ab}}/K)$. 与其直接研究 $\varprojlim_{\mathfrak{m}}C_{\mathfrak{m}}$, 不如引入另一个以它为商的群更为自然\,---\,这就是 idèle 类群.

\phantomsection\label{pg:orig169}%
\subsection*{拓扑群}

带拓扑的群 $G$ 称为\textbf{拓扑群}, 若映射
\[
g,g'\mapsto gg'\colon G\times G\to G,\qquad g\mapsto g^{-1}\colon G\to G
\]
连续. 平移映射
\[
g\mapsto ag\colon G\to G
\]
此时是同胚.

一般地, 要在一个集合上确定拓扑, 须给出每点的基本邻域组, 即该点的一组邻域, 使每个邻域都包含组中的一个. 由于平移是同胚, 拓扑群的拓扑由 $1$ 的基本邻域组确定.

我们将用到拓扑群的若干一般性事实, 它们可在许多书中找到. 幸运的是, 我们只需要十分初等的内容.

设 $(X_i)_i$ 是一族 (可以无限) 拓扑空间. $\prod X_i$ 上的\textbf{乘积拓扑}是使形如 $\prod U_i$ ($U_i$ 在 $X_i$ 中开, 且除有限多个 $i$ 外 $U_i=X_i$) 的集合构成基的拓扑. Tychonoff 定理说紧空间的乘积是紧的. 然而, 局部紧空间的无穷乘积一般不再局部紧: 若 $V_i$ 是 $x_i$ 在 $X_i$ 中的紧邻域, 则 $\prod V_i$ 是紧的, 但除非除有限多个 $i$ 外 $V_i=X_i$, 它不是 $(x_i)$ 的邻域.

\subsection*{Idèle}

今后常用 $v$ 记 $K$ 的一个素. 则:

$|\cdot|_v$ = $v$ 的正规化绝对值 (使乘积公式成立的那个);

$K_v$ = $K$ 在 $v$ 处的完备化;

$\mathfrak{p}_v$ = $\mathcal{O}_K$ 中相应的素理想 (当 $v$ 有限时);

$\mathcal{O}_v$ = $K_v$ 的整数环;

$U_v=\mathcal{O}_v^{\times}$;

$\hat{\mathfrak{p}}_v$ = $\mathfrak{p}_v$ 的完备化 = $\mathcal{O}_v$ 的极大理想.

回顾: 对所有 $v$, $K_v$ 局部紧\,---\,事实上 $\mathcal{O}_v$ 是 $0$ 的紧邻域. 类似地 $K_v^{\times}$ 局部紧; 事实上
\[
1+\hat{\mathfrak{p}}_v\supset 1+\hat{\mathfrak{p}}_v^2\supset 1+\hat{\mathfrak{p}}_v^3\supset\cdots
\]
是由开紧子群构成的 $1$ 的基本邻域组.

我们想把所有的群 $K_v^{\times}$ 合成一个大的拓扑空间, 但 $\prod K_v^{\times}$ 不是局部紧的. 于是定义 \textbf{idèle} 群\index{idele@idèle(idèle)}为\footnote{idèle 的其他记号: $\mathbb{A}_K^{\times}$, $\mathrm{GL}_1(\mathbb{A}_K)$.}
\[
\mathbb{I}_K=\Bigl\{(a_v)\in\prod K_v^{\times}\ \Big|\ a_v\in\mathcal{O}_v^{\times}\ \text{对除有限多个外的所有}\ v\Bigr\}.
\]
下文中不带下标的 $\mathbb{I}$ 表示 $\mathbb{I}_K$.

对包含所有无限素的素的每个有限集 $S$, 令
\[
\mathbb{I}_S=\prod_{v\in S}K_v^{\times}\times\prod_{v\notin S}\mathcal{O}_v^{\times},
\]
带乘积拓扑. 第一个因子是局部紧空间的有限乘积, 故局部紧; 第二个因子是紧空间的乘积, 故紧 (由 Tychonoff). 因此 $\mathbb{I}_S$ 局部紧. 注意
\[
\mathbb{I}=\bigcup\mathbb{I}_S.
\]
我们想在 $\mathbb{I}$ 上赋予拓扑, 使每个 $\mathbb{I}_S$ 在 $\mathbb{I}$ 中开且继承乘积拓扑. 做法: 规定开集基由形如 $\prod_v V_v$ 的集合组成, 其中 $V_v$ 在 $K_v^{\times}$ 中开, 且对\textbf{几乎所有}\footnote{此处\guillemotleft 几乎所有\guillemotright\ 意为\guillemotleft 除可能有限多个外的所有\guillemotright.} $v$ 有 $V_v=\mathcal{O}_v^{\times}$. 两个这种形状的集合的交仍包含一个这种形状的集合, 故它们确构成某拓扑的基. 显然该拓扑具有我们想要的性质, 而且它赋予 $\mathbb{I}$ 拓扑群的结构. 下列集合构成 $1$ 的基本邻域组: 对每个包含 $S_\infty$ 的有限素集 $S$ 及 $\varepsilon>0$, 定义
\[
U(S,\varepsilon)=\{(a_v)\mid |a_v-1|_v<\varepsilon,\ v\in S,\quad |a_v|_v=1,\ \text{所有 }v\notin S\}.
\]

\textbf{4.1.} 存在典范的满同态
\[
(a_v)\mapsto\prod_{v\ \mathrm{finite}}\mathfrak{p}_v^{\mathrm{ord}_{\mathfrak{p}}(a_v)}\colon\mathbb{I}_K\to I_K,
\]
其核为 $\mathbb{I}_{S_\infty}$.

可以把 idèle 想成理想的加厚: 它包含无限素的因子, 也包含有限素处的单位. 注意 $\mathbb{I}/\mathbb{I}_{S_\infty}$ 是可数多个 $\Z$ 的带离散拓扑的直和, 而 $\prod K_v^{\times}/\mathbb{I}_{S_\infty}$ 是可数多个 $\Z$ 的直积, 后者本身不可数.

\textbf{4.2.} 存在典范的单射 (对角) 同态
\[
a\mapsto(a,a,a,\ldots)\colon K^{\times}\to\mathbb{I}_K.
\]
我断言其像是离散的. 由于是群, 只须证 $1\in K^{\times}$ 在诱导拓扑下是开的. 取 $U=U(S,\varepsilon)$, $S$ 为任意包含 $S_\infty$ 的有限集, $1>\varepsilon>0$. 对每个 $a\in K^{\times}\cap U$,
\[
\begin{cases}
|a-1|_v<\varepsilon & \text{对所有 } v\in S,\\
|a|_v=1 & \text{对所有 } v\notin S.
\end{cases}
\]
第二个条件蕴含
\[
|a-1|_v\le\max\bigl(|a|_v,|-1|_v\bigr)\le 1.
\]
因此若 $a\in K^{\times}\cap U$, 则 $\prod_v|a-1|_v<\varepsilon^{|S|}<1$, 与乘积公式矛盾, 除非 $a=1$.

商群 $C=\mathbb{I}/K^{\times}$ 称为 $K$ 的 \textbf{idèle 类群}\index{ideleleiqun@idèle 类群(idèle class group)}. 它映满 $K$ 的理想类群. 它不紧 (见下文 (\hyperlink{item:v.4.4}{4.4})).

\textbf{4.3.} 存在典范的单同态
\[
a\mapsto(1,\ldots,1,a,1,\ldots,1)\colon K_v^{\times}\to\mathbb{I}_K
\]
($a$ 在第 $v$ 个位置). $K_v^{\times}$ 上诱导的拓扑是其自然拓扑, 因为
\[
U(S,\varepsilon)\cap K_v^{\times}=
\begin{cases}
\{a\mid|a-1|_v<\varepsilon\} & v\in S,\\
\{a\mid|a|_v=1\} & v\notin S,
\end{cases}
\]
且这些集合构成 $K_v^{\times}$ 中 $1$ 的基本邻域组.

\phantomsection\hypertarget{item:v.4.4}{}\textbf{4.4.} 存在典范的满同态
\[
\mathbf{a}=(a_v)\mapsto c(\mathbf{a})=\prod|a_v|_v\colon\mathbb{I}\to\R_{>0}.
\]
$\mathbf{a}$ 的像称为 $\mathbf{a}$ 的\textbf{容量} (content). 定义
\[
\mathbb{I}^1=\mathrm{Ker}(c)=\{\mathbf{a}\in\mathbb{I}\mid c(\mathbf{a})=1\}.
\]
注意由乘积公式, $K^{\times}\subset\mathbb{I}^1$. 商群 $\mathbb{I}/K^{\times}$ 不可能紧, 因为它满射到 $\R_{>0}$; 但可以证明 $\mathbb{I}^1/K^{\times}$ 是紧的.

\setcounter{theorem}{4} % 4.1--4.4 为叙述性编号条目, 不占定理计数器; 调至 4 使下一条显示为 4.5
\begin{aside}\label{asi:v.4.5}
按同样方式仅用有限素定义的 $\mathbb{I}_f$ 称为\textbf{有限 idèle} 群. 我们有
\[
\prod_{v\ \mathrm{finite}}\mathcal{O}_v^{\times}\subset\mathbb{I}_f\subset\prod_{v\ \mathrm{finite}}K_v^{\times}.
\]
子群 $\prod\mathcal{O}_v^{\times}$ 在 $\mathbb{I}_f$ 中开且紧, 且 $\mathbb{I}_f/\prod\mathcal{O}_v^{\times}=I$ ($K$ 的理想群).

仍有 $K^{\times}$ 到 $\mathbb{I}_f$ 的对角嵌入, 但此时 $K^{\times}$ 上的诱导拓扑可描述如下: $U_K\overset{\mathrm{def}}{=}\mathcal{O}_K^{\times}$ 是开的, 而 $1$ 的基本邻域组由 $U_K$ 的有限指数子群构成 (非平凡定理).\footnote{取自 \sx{140729}. 第一个断言: 有限 idèle 中由 $\prod\mathcal{O}_v^{\times}$ 给出的子集按限制直积拓扑的定义是开的, 而 $K^{\times}$ 中的元素属于该子集当且仅当它在每个完备化中都是单位, 即当且仅当它属于 $\mathcal{O}_K^{\times}$. 第二个断言的证明只是第一个论证的推广: 有限 idèle 中 $1$ 的开邻域基可以这样给出\,---\,在有限多个素 $v$ 处任取 $K_v^{\times}$ 的开子群, 在其余 $v$ 处取 $\mathcal{O}_v^{\times}$. 若只想要开集基, 不妨只考虑非常小的那些, 故可设在所选的有限多个 $v$ 处取 $\mathcal{O}_v^{\times}$ 的小开子集. 注意 $\mathcal{O}_v^{\times}$ 的子群都是有限指数的, 且由素元的幂生成. 把这样的邻域与 $K^{\times}$ 的对角像相交, 得到的是 $\mathcal{O}_K^{\times}$ 中落在所选的那些 $\mathcal{O}_v^{\times}$ 的开子群里的元素构成的子群, 其指数恰为诸指数之积, 故为有限指数子群; 反之, $\mathcal{O}_K^{\times}$ 的任何有限指数子群由其在各完备化中的像确定.} 特别地, $K^{\times}$ 是 $\mathbb{I}_f$ 的离散子群 $\iff U_K$ 有限 $\iff K=\Q$ 或虚二次域.
\end{aside}

\begin{NOTES}
adèle 拓扑并不在 idèle 上诱导 idèle 拓扑 (因为 $\prod_{v|\infty}K_v^{\times}\times\prod_{v\,\mathrm{finite}}\mathcal{O}_v^{\times}$ 在 adèle 中不开). 作为练习, 证明 idèle 在诱导拓扑下甚至不是拓扑群. (先考虑 $\Q$ 的情形, 忘掉无限素. (a) 确保理解 adèle 中 $0$ 的标准邻域基. (b) 把 (a) 中的邻域平移加 $1$, 与 idèle 相交, 得到 adèle 拓扑下 idèle 中 $1$ 的邻域基. (c) 验证 (b) 中某个邻域在映射 $x\mapsto x^{-1}$ 下的原像不开.)
\end{NOTES}

\begin{NOTES}
设 $(V_v,U_v)_v$ 是一族拓扑空间 $V_v$, 每个带有子空间 $U_v$. 该族的\textbf{限制直积} $\prod(V_v,U_v)$ 是 $\prod V_v$ 中由满足\guillemotleft 对几乎所有 $v$, $a_v\in U_v$\guillemotright\ 的族 $(a_v)$ 组成的子集, 其拓扑使 $\prod U_v$ 开且带乘积拓扑. 限制直积有时记作 $\prod' V_v$. 于是 $\mathbb{A}_K=\prod_v(K_v,U_v)$, 其中 $v$ 有限时 $U_v=\mathcal{O}_v$, 否则 $U_v=K_v$. 类似地 $\mathbb{A}_K^{\times}=\prod_v(K_v^{\times},U_v^{\times})$, 其中 $v$ 有限时 $U_v=\mathcal{O}_v^{\times}$, 否则 $U_v=K_v^{\times}$.
\end{NOTES}

\begin{NOTES}
再举些例子. 设 $K$ 是类数为 $1$ 的域. 设 $x=(x_v)$ 是一个 idèle. 则 $x$ 的理想是主理想, 设等于 $(a)$. 此时 $x/a\in\prod_{\mathrm{finite}}\mathcal{O}_v^{\times}\times K_\infty^{\times}$. 故 $\mathbb{A}_K^{\times}=K^{\times}\cdot\prod_{\mathrm{finite}}\mathcal{O}_v^{\times}\times K_\infty^{\times}$. 然而 $K^{\times}\cap(\prod_{\mathrm{finite}}\mathcal{O}_v^{\times}\times K_\infty^{\times})=U_K$.
\end{NOTES}

\begin{NOTES}
取自 \sx{80763}. 普林斯顿资格考试的问题: 关于 adèle 与 idèle 的主要紧性定理给出了哪些结果? 答: 理想类群的有限性与单位定理 (用经典方法证明) 可证范数一 idèle 类群的紧性; 而用测度论方法证明的范数一 idèle 类群的紧性又可反推理想类群的有限性与单位定理. (见 \cite{Cassels1967}.) 另一方面, 紧性结果确实可推广到可除代数, 那里没有明显的理想类群等概念. (见 Weil 的 Basic Number Theory.)
\end{NOTES}

\subsection*{把射线类群实现为 $\mathbb{I}$ 的商}

我们已经看到类群 $C_K=I/i(K^{\times})$ 可以实现为 $\mathbb{I}$ 的商. 现在要对 $C_{\mathfrak{m}}$ 证明同样的事.

设 $\mathfrak{m}$ 为模. 对 $\mathfrak{p}|\mathfrak{m}$, 令
\[
W_{\mathfrak{m}}(\mathfrak{p})=
\begin{cases}
\R_{>0} & \mathfrak{p}\ \text{实},\\
1+\hat{\mathfrak{p}}^{m(\mathfrak{p})} & \mathfrak{p}\ \text{有限}.
\end{cases}
\]
于是在每种情形, $W_{\mathfrak{m}}(\mathfrak{p})$ 都是 $K_{\mathfrak{p}}^{\times}$ 中 $1$ 的邻域.

定义 $\mathbb{I}_{\mathfrak{m}}$ 为满足 $a_{\mathfrak{p}}\in W_{\mathfrak{m}}(\mathfrak{p})$ (对所有 $\mathfrak{p}|\mathfrak{m}$) 的 idèle $(a_{\mathfrak{p}})_{\mathfrak{p}}$ 之集:
\[
\mathbb{I}_{\mathfrak{m}}=\Bigl(\prod_{\mathfrak{p}\nmid\mathfrak{m}}K_{\mathfrak{p}}^{\times}\times\prod_{\mathfrak{p}|\mathfrak{m}}W_{\mathfrak{m}}(\mathfrak{p})\Bigr)\cap\mathbb{I}.
\]
换言之, $\mathbb{I}_{\mathfrak{m}}$ 由 $K$ 的素标号的族 $(a_{\mathfrak{p}})_{\mathfrak{p}}$ 中满足
\[
\begin{cases}
a_{\mathfrak{p}}\in K_{\mathfrak{p}}^{\times} & \text{对所有 }\mathfrak{p},\\
a_{\mathfrak{p}}\in\mathcal{O}_{\mathfrak{p}}^{\times} & \text{对几乎所有 }\mathfrak{p},\\
a_{\mathfrak{p}}\in W_{\mathfrak{m}}(\mathfrak{p}) & \text{对所有 }\mathfrak{p}|\mathfrak{m}
\end{cases}
\]
者组成. 定义 $W_{\mathfrak{m}}$ 为 $\mathbb{I}_{\mathfrak{m}}$ 中使 $a_{\mathfrak{p}}$ 在所有不整除 $\mathfrak{m}$ 的有限 $\mathfrak{p}$ 处为单位的 idèle $(a_{\mathfrak{p}})_{\mathfrak{p}}$ 之集:
\[
W_{\mathfrak{m}}=\prod_{\substack{\mathfrak{p}\nmid\mathfrak{m}\\ \mathfrak{p}\ \mathrm{infinite}}}K_{\mathfrak{p}}^{\times}\times\prod_{\mathfrak{p}|\mathfrak{m}}W_{\mathfrak{m}}(\mathfrak{p})\times\prod_{\substack{\mathfrak{p}\nmid\mathfrak{m}\\ \mathfrak{p}\ \mathrm{finite}}}U_{\mathfrak{p}}.
\]
换言之, $W_{\mathfrak{m}}$ 由满足
\[
\begin{cases}
a_{\mathfrak{p}}\in K_{\mathfrak{p}}^{\times} & \text{对所有无限 }\mathfrak{p},\\
a_{\mathfrak{p}}\in\mathcal{O}_{\mathfrak{p}}^{\times} & \text{对所有有限 }\mathfrak{p},\\
a_{\mathfrak{p}}\in W_{\mathfrak{m}}(\mathfrak{p}) & \text{对所有 }\mathfrak{p}|\mathfrak{m}
\end{cases}
\]
的族 $(a_{\mathfrak{p}})_{\mathfrak{p}}$ 组成. 注意
\[
K_{\mathfrak{m},1}=K^{\times}\cap\prod_{\mathfrak{p}|\mathfrak{m}}W_{\mathfrak{m}}(\mathfrak{p})\qquad(\text{相交在 }\prod_{\mathfrak{p}|\mathfrak{m}}K_{\mathfrak{p}}^{\times}\ \text{内}),
\]
且
\[
K_{\mathfrak{m},1}=K^{\times}\cap\mathbb{I}_{\mathfrak{m}}\qquad(\text{相交在 }\mathbb{I}\ \text{内}).
\]

\begin{proposition}\label{prop:v.4.6}
设 $\mathfrak{m}$ 是 $K$ 的模.

(a) 映射 $\mathrm{id}\colon\mathbb{I}_{\mathfrak{m}}\to I^{S(\mathfrak{m})}$ 定义同构
\[
\mathbb{I}_{\mathfrak{m}}/K_{\mathfrak{m},1}\cdot W_{\mathfrak{m}}\xrightarrow{\ \simeq\ }C_{\mathfrak{m}}.
\]

(b) 包含 $\mathbb{I}_{\mathfrak{m}}\hookrightarrow\mathbb{I}$ 定义同构
\[
\mathbb{I}_{\mathfrak{m}}/K_{\mathfrak{m},1}\to\mathbb{I}/K^{\times}.
\]
\end{proposition}

\begin{proof}
(a) 考虑映射对
\[
K_{\mathfrak{m},1}\to\mathbb{I}_{\mathfrak{m}}\xrightarrow{\ \mathrm{id}\ }I^{S(\mathfrak{m})}.
\]
第一个映射单, 第二个满且核为 $W_{\mathfrak{m}}$, 故该映射对的核--余核序列 (II, \ref{lem:ii.A.2}) 为
\[
W_{\mathfrak{m}}\to\mathbb{I}_{\mathfrak{m}}/K_{\mathfrak{m},1}\to C_{\mathfrak{m}}\to 1.
\]
这证明了命题的 (a).

(b) $\mathbb{I}_{\mathfrak{m}}\to\mathbb{I}/K^{\times}$ 的核是 $K^{\times}\cap\mathbb{I}_{\mathfrak{m}}$ (相交在 $\mathbb{I}$ 内), 如刚才所见它是 $K_{\mathfrak{m},1}$. 于是包含映射定义单射
\[
\mathbb{I}_{\mathfrak{m}}/K_{\mathfrak{m},1}\hookrightarrow\mathbb{I}/K^{\times}.
\]
至于满性, 应用弱逼近定理 (ANT, 7.20). 设 $S=S(\mathfrak{m})$, $\mathbf{a}=(a_v)\in\mathbb{I}$. 若取 $b\in K$ 使其在 $K_v^{\times}$ 中对所有 $v\in S$ 都非常接近 $a_v$, 则 $a_v/b$ 在 $K_v^{\times}$ 中对所有 $v\in S$ 接近 $1$; 事实上可取 $b$ 使 $a_v/b\in W_{\mathfrak{m}}(\mathfrak{p})$ 对所有 $v\in S$ 成立. 例如对 $S$ 中的实素 $v$, 只须取 $b$ 与 $a_v$ 在 $K_v$ 中同号. 此时 $\mathbf{a}/b\in\mathbb{I}_{\mathfrak{m}}$, 且它在 $\mathbb{I}/K^{\times}$ 中的像为 $\mathbf{a}$.
\end{proof}

\subsection*{理想的特征与 idèle 的特征}

设 $S\supset S_\infty$ 是 $K$ 的素的有限集, $G$ 是有限 Abel 群. 同态
\[
\psi\colon I^S\to G
\]
称为\textbf{容许一个模}, 若存在支撑含于 $S$ 的模 $\mathfrak{m}$ 使 $\psi(i(K_{\mathfrak{m},1}))=1$. 例如, 对每个 Abel 扩张 $L/K$, Artin 证明了 Artin 映射
\[
I^S\to\Gal(L/K)
\]
容许一个模.

\begin{proposition}\label{prop:v.4.7}
若 $\psi\colon I^S\to G$ 容许一个模, 则存在唯一的同态 $\phi\colon\mathbb{I}\to G$ 使得

(a) $\phi$ 连续 ($G$ 取离散拓扑);

(b) $\phi(K^{\times})=1$;

(c) $\phi(\mathbf{a})=\psi(\mathrm{id}(\mathbf{a}))$, 对所有 $\mathbf{a}\in\mathbb{I}^S\overset{\mathrm{def}}{=}\{\mathbf{a}\mid a_v=1\ \text{对所有}\ v\in S\}$.

此外, 每个满足 (b) 的连续同态 $\phi\colon\mathbb{I}\to G$ 都由某个 $\psi$ 产生.
\end{proposition}

\begin{proof}
因 $\psi$ 容许模 $\mathfrak{m}$, 它经过 $I^{\mathfrak{m}}/i(K_{\mathfrak{m},1})=C_{\mathfrak{m}}$ 分解. 于是有图:
\begin{equation}
\begin{tikzcd}[column sep=small]
I^{\mathfrak{m}}\ar[r] & C_{\mathfrak{m}}\ar[r, "\psi"] & G\\
\mathbb{I}_{\mathfrak{m}}/K_{\mathfrak{m},1}\ar[r]\ar[u, "\simeq"] & \mathbb{I}_{\mathfrak{m}}/K_{\mathfrak{m},1}W_{\mathfrak{m}}\ar[u, "\simeq"'] & \\
\mathbb{I}\ar[r]\ar[u, hookrightarrow] & \mathbb{I}/K^{\times}\ar[u, two heads] &
\end{tikzcd}\tag{43}
\end{equation}
其中的同构即命题 \ref{prop:v.4.6} 中的那些, 其余未命名的映射是商映射. 定义 $\phi$ 为复合 $\mathbb{I}\to G$. 它显然具有性质 (a) 与 (b), 且
\[
\phi(\mathbf{a})=\psi(\mathrm{id}(\mathbf{a}))\quad\text{对所有}\ \mathbf{a}\in\mathbb{I}_{\mathfrak{m}},
\]
从而更是满足性质 (c).

为证该映射由 (a), (b), (c) 唯一确定, 只须证 $\mathbb{I}^S K^{\times}$ 在 $\mathbb{I}$ 中稠密, 而这由弱逼近定理 (ANT, 7.20) 推出: 设 $\mathbf{a}\in\mathbb{I}$; 取 $b\in K^{\times}$ 使其对 $v\in S$ 非常接近 $a_v$, 并令 $\mathbf{a}'$ 为满足 $a_v'b=a_v$ (对所有 $v\notin S$) 的 $\mathbb{I}^S$ 中的元素. 则 $\mathbf{a}'b\in\mathbb{I}^S\cdot K$ 且在 $\mathbb{I}$ 中接近 $\mathbf{a}$.

反之, 设 $\phi\colon\mathbb{I}\to G$ 连续. 其核包含 $1$ 的一个开邻域, 故对某个 $S$ 与 $\varepsilon$ 有 $U(S,\varepsilon)\subset\mathrm{Ker}(\phi)$. 考虑无限素 $v$. $\phi$ 在 $K_v^{\times}$ 上的限制是连续映射 $\R^{\times}\to G$ 或 $\CC^{\times}\to G$. 显然 $K_v^{\times}$ 中含 $1$ 的连通分支, 即 $\R_{>0}$ 或 $\CC^{\times}$, 映到 $1$, 故在核中. 综合这些观察, 可见 $\phi$ 的核包含某个 $\mathfrak{m}$ 的 $W_{\mathfrak{m}}$.

现在可以再次使用证明开头的图. 我们有同态 $\phi\colon\mathbb{I}/K^{\times}\to G$, 可把它\guillemotleft 限制\guillemotright\ 为同态 $\mathbb{I}_{\mathfrak{m}}/K_{\mathfrak{m},1}\to G$. 该同态在 $W_{\mathfrak{m}}$ 上平凡, 故经过 $\mathbb{I}_{\mathfrak{m}}/K_{\mathfrak{m},1}W_{\mathfrak{m}}$ 分解. 此同态可转移到 $C_{\mathfrak{m}}$ 上, 再与 $\mathbb{I}\to C_{\mathfrak{m}}$ 复合. 这便是要找的 $\psi$.
\end{proof}

\begin{remark}\label{rem:v.4.8}
设 $G$ 为交换拓扑群. 称同态 $\psi\colon I^S\to G$ \textbf{可容许的} (admissible), 若对 $G$ 中 $1$ 的每个邻域 $N$, 存在模 $\mathfrak{m}$ 使 $\psi(i(K_{\mathfrak{m},1}))\subset N$. 则每个可容许同态 $\psi$ 定义满足命题 \ref{prop:v.4.7} 条件 (a), (b), (c) 的同态 $\phi\colon\mathbb{I}\to G$. 此外, 若 $G$ 完备且\guillemotleft 无小子群\guillemotright, 即存在 $1$ 的不含非平凡子群的邻域, 则每个满足 (b) 的连续同态 $\phi\colon\mathbb{I}\to G$ 都由某个可容许的 $\psi$ 产生. 证明与命题 \ref{prop:v.4.7} 相同 (见 \cite{Tate1967}, 命题 4.1).

圆群 $G=\{z\in\CC\mid|z|=1\}$ 完备且无小子群. 可容许的 $\psi\colon I^S\to G$ 及相应的 $\phi$ 称为 \textbf{Hecke 特征}\index{heckelezheng@Hecke 特征(Hecke character)}.
\end{remark}

\begin{remark}\label{rem:v.4.9}
给定 $\psi$ 后我们选取了一个 $\mathfrak{m}$, 然后说明如何构造 $\phi$. 实践中更常见的做法是: 在已知 $\phi$ 满足条件 (a), (b), (c) 的情况下直接把它认出来. 为此, 下列观察有用.

(a) 设 $\mathbf{a}=(a_v)$ 是满足 $a_v=1$ 对所有有限素、$a_v>0$ 对所有实素的 idèle; 则 $\phi(\mathbf{a})=1$. 为看出这一点, 注意 $\prod_{v|\infty}K_v^{\times}$ 作为 $\mathbb{I}$ 的子群上的诱导拓扑是其自然拓扑, 故 $\phi$ 在其上的限制在含 $1$ 的连通分支上平凡.

(b) 设 $\mathbf{a}=(a_v)$ 是满足 $a_v=1$ 对所有 $v\in S$、$a_v$ 为单位对所有 $v\notin S$ 的 idèle; 则 $\phi(\mathbf{a})=1$. 事实上这直接由条件 (c) 推出.

(c) 若 $\mathbf{a}$ \guillemotleft 接近 $1$\guillemotright, 则 $\phi(\mathbf{a})=1$. 事实上, 鉴于 $G$ 取离散拓扑, 这直接由条件 (a) 推出.

(d) 结合 (a), (b), (c) 可知: 若 $\mathbf{a}=(a_v)$ 满足
\[
\begin{cases}
a_v>0 & \text{当 } v\ \text{实};\\
a_v\ \text{接近}\ 1 & \text{当 } v\in S\ \text{有限};\\
a_v\ \text{为单位} & \text{当 } v\notin S,
\end{cases}
\]
则 $\phi(\mathbf{a})=1$. 事实上, (a) 与 (b) 表明可以用某些类型的 idèle 乘 $\mathbf{a}$ 而不改变 $\phi(\mathbf{a})$ 的值. 显然, 若 $a_v$ 对 $S$ 中的有限 $v$ 接近 $1$, 可以用这样的 idèle 乘之使其更接近 $1$.
\end{remark}

\begin{example}\label{ex:v.4.10}
设 $L=\Q[\zeta_p]$, $\psi$ 为 Artin 映射
\[
I^S\to(\Z/p\Z)^{\times}\to\Gal(L/\Q),\qquad S=\{p,\infty\}.
\]
回顾: 第一个映射把 (唯一) 由 $(r/s)$, $r,s>0$, $(p,r)=1=(p,s)$, 代表的理想送到 $[r][s]^{-1}$, 第二个映射把 $[m]$ 送到自同构 $\zeta\mapsto\zeta^m$. 总的说来, 对任意素数 $l\ne p$, 该映射把 $(l)$ 送到 $l$ 处的 Frobenius 自同构 $\zeta\mapsto\zeta^l$. 设 $\phi\colon\mathbb{I}\to\Gal(L/\Q)$ 为定理中对应于 $\psi$ 的同态. 我们要显式确定 $\phi$.

设 $\mathbf{a}=(a_\infty,a_2,\ldots,a_p,\ldots,a_l,\ldots)$ 是 $\Q$ 的 idèle. 若 $a_\infty=1=a_p$, 则 $\phi(\mathbf{a})=\psi(\mathrm{id}(\mathbf{a}))$. 于是 $\phi(\mathbf{a})=\zeta_p^m$, 其中 $m=\prod l^{\mathrm{ord}_l(a_l)}$.

考虑 $\mathbf{p}=(1,\ldots,1,p,1,\ldots)$ ($p$ 在第 $p$ 个位置). 则
\[
\mathbf{p}/p=(p^{-1},\ldots,p^{-1},1,p^{-1},\ldots).
\]
据上述注记的 (d), $\phi(\mathbf{p}/p)=1$, 故
\[
\phi(\mathbf{p})=\phi(\mathbf{p}/p)\phi(p)=1.
\]
此处 $p$ 既表示元素 $p\in\Q_p$, 也表示主 idèle $(p,p,\ldots)$.

考虑 $\mathbf{a}=(1,\ldots,1,u,1,\ldots)$, $u\in\Z_p^{\times}$, $u$ 在第 $p$ 个位置. 写
\[
u^{-1}=a_0+a_1p+\cdots+a_sp^s+\cdots,\qquad 0\le a_i<p,\quad a_i\in\Z,
\]
并令 $c=a_0+\cdots+a_sp^s\in\Z$. 注意 $c>0$. 则 $uc\in 1+p^{s+1}\Z_p$, 即对大的 $s$ 它\guillemotleft 接近 $1$\guillemotright. 写
\[
\mathbf{a}c=(c,c,\ldots,c,\overset{l|c}{1},c,\ldots,\overset{p}{uc},c,\ldots)\,(1,\ldots,1,\overset{l|c}{c},1,\ldots).
\]
第一个因子是 $\mathbf{a}c$, 只是把整除 $c$ 的素 $l$ 处的分量挪到了第二个因子. 对大的 $s$, 由上述注记的 (d) 有 $\phi(\text{第一个因子})=1$. 第二个因子属于 $\mathbb{I}^S$, 而我们对 $\phi|\mathbb{I}^S$ 的描述表明 $\phi(\text{第二个因子})$ 把 $\zeta$ 送到 $\zeta^c$. 结论:
\begin{equation}
\phi(\mathbf{a})(\zeta)=\zeta^c=\zeta^{u^{-1}}.\tag{44}
\end{equation}

考虑 $\mathbf{a}=(-1,1,\ldots,1)$. 则
\[
-\mathbf{a}=(1,-1,\ldots,-1,\overset{p}{1},-1,\ldots)\,(1,\ldots,1,\overset{p}{-1},1,\ldots),
\]
而
\[
\phi(\mathbf{a})=\phi(-\mathbf{a})=\phi(1,\ldots,1,-1,1,\ldots).
\]
据 (44), $\phi(\mathbf{a})(\zeta)=\zeta^{-1}$.

由于 $\phi$ 是同态, 这就完成了对它的显式描述.
\end{example}

\begin{remark}\label{rem:v.4.11}
由 (43) 得到典范同态 $\pi_{\mathfrak{m}}\colon\mathbb{I}\to C_{\mathfrak{m}}$. 这是唯一的连续同态 $\pi_{\mathfrak{m}}\colon\mathbb{I}\to C_{\mathfrak{m}}$ 使得

(a) $\pi_{\mathfrak{m}}(K^{\times})=1$;

(b) $\pi_{\mathfrak{m}}(\mathbf{a})=\mathrm{id}(\mathbf{a})$ 对所有 $\mathbf{a}\in\mathbb{I}^{S(\mathfrak{m})}$.

(在命题 \ref{prop:v.4.7} 中取 $\psi$ 为 $I^S\to C_{\mathfrak{m}}$.)

若 $\mathfrak{m}|\mathfrak{m}'$, 则 $\pi_{\mathfrak{m}'}$ 与典范同态 $C_{\mathfrak{m}'}\to C_{\mathfrak{m}}$ 的复合满足刻画 $\pi_{\mathfrak{m}}$ 的条件. 因此诸 $\pi_{\mathfrak{m}}$ 合成一个连续同态 $\pi\colon\mathbb{I}\to\varprojlim C_{\mathfrak{m}}$. 我们要确定该映射的核与像.

由于每个 $\pi_{\mathfrak{m}}\colon\mathbb{I}\to C_{\mathfrak{m}}$ 都是满射, 像是稠密的. 事实上 $\mathbb{I}^1\to C_{\mathfrak{m}}$ 是满射, 故 $\pi_{\mathfrak{m}}(\mathbb{I}^1)$ 稠密. 但 $\pi_{\mathfrak{m}}(\mathbb{I}^1)$ 是紧的, 因为 $\pi_{\mathfrak{m}}$ 经过紧群 $\mathbb{I}^1/K^{\times}$ 分解, 因而是完备的. 这表明 $\pi$ 是满射.

设 $\mathbb{I}_\infty^+$ 为满足 $a_v=1$ ($v$ 有限)、$a_v>0$ ($v$ 实) 的 idèle $\mathbf{a}$ 之集. 则 $\mathbb{I}_\infty^+$ 同构于 $(K\otimes_{\Q}\R)^{\times}=\prod_{v|\infty}K_v^{\times}$ 的单位分支. $\mathbb{I}\to\varprojlim C_{\mathfrak{m}}$ 的核包含 $\mathbb{I}_\infty^+\cdot K^{\times}$, 从而包含其闭包. 事实上它就等于闭包. [在将来的版本中, 将更详细地考察这些内容.]
\end{remark}

\subsection*{idèle 的范}

设 $L$ 是数域 $K$ 的有限扩张, $v$ 是 $K$ 的素. 回顾 (ANT, 8.2): 存在典范同构
\[
L\otimes_K K_v\to\prod_{w|v}L_w.
\]
由此 (同书 8.3) 对任意 $\alpha\in L$,
\[
\Nm_{L/K}\alpha=\prod_{w|v}\Nm_{L_w/K_v}\alpha\qquad(\text{等式在}\ K_v\ \text{内}).
\]
对 idèle $\mathbf{a}=(a_w)\in\mathbb{I}_L$, 定义 $\Nm_{L/K}(\mathbf{a})$ 为 $\mathbb{I}_K$ 中满足 $b_v=\prod_{w|v}\Nm_{L_w/K_v}a_w$ 的 idèle $\mathbf{b}$. 前面的注记表明下述图中左侧的正方形交换, 右侧的正方形显然也交换:
\[
\begin{tikzcd}
L^{\times}\ar[r]\ar[d, "\Nm_{L/K}"'] & \mathbb{I}_L\ar[r, "\mathrm{id}"]\ar[d, "\Nm_{L/K}"'] & I_L\ar[d, "\Nm_{L/K}"']\\
K^{\times}\ar[r] & \mathbb{I}_K\ar[r, "\mathrm{id}"'] & I_K
\end{tikzcd}
\]
于是得到交换图:
\[
\begin{tikzcd}
C_L\ar[r]\ar[d, "\Nm_{L/K}"'] & \mathbf{C}_L\ar[d, "\Nm_{L/K}"']\\
C_K\ar[r] & \mathbf{C}_K
\end{tikzcd}
\]
($C_K=$ 理想类群 $I/i(K^{\times})$; $\mathbf{C}_K=$ idèle 类群 $\mathbb{I}/K^{\times}$.)

\begin{proposition}\label{prop:v.4.12}
若 $L/K$ 是特征零的局部域的有限扩张, 则

(a) $\Nm_{L/K}(L^{\times})=\R_{>0}$ (情形 $K=\R$, $L=\CC$);

(b) $\Nm_{L/K}(L^{\times})\supset 1+\mathfrak{p}_K^m$ 对某个 $m$ (情形 $K$ 非阿基米德);

(c) $\Nm_{L/K}(L^{\times})\supset\mathcal{O}_K^{\times}$ (情形 $K$ 非阿基米德且 $L/K$ 非分歧).
\end{proposition}

\begin{proof}
(a) 显然. (b) 见 (I \ref{lem:1.3}), (c) 见 (III, \ref{prop:iii.1.2}).
\end{proof}

\begin{corollary}\label{cor:v.4.13}
设 $L/K$ 是数域的有限扩张. 则 $\Nm_{L/K}\mathbb{I}_L\supset W_{\mathfrak{m}}$ 对某个模 $\mathfrak{m}$.
\end{corollary}

\section{idèle 语言的主要定理}\label{sec:V.5}

用理想陈述的类域论主要定理非常显式, 且对许多目的而言是最有用的表述. 但它有一些缺点: 必须先固定一个模 $\mathfrak{m}$, 而该理论只描述导子整除 $\mathfrak{m}$ 的那些 Abel 扩张; 特别地, 它对 $K$ 的无限 Abel 扩张不作任何描述. 用 idèle 陈述的主要定理使人能同时考虑无限 Abel 扩张, 或等价地, 一并考虑所有有限 Abel 扩张. 它还使整体与局部 Artin 映射之间的关系变得透明.

设 $L$ 是 $K$ 的有限 Abel 扩张. 设 $v$ 是 $K$ 的素, $w$ 是 $L$ 中位于 $v$ 之上的素. 回顾 $w$ 的分解群 $D(w)$ 是子群
\[
D(w)=\{\sigma\in\Gal(L/K)\mid\sigma w=w\}.
\]
其元素唯一地延拓为 $L_w/K_v$ 的自同构, 且 $D(w)\simeq\Gal(L_w/K_v)$. 局部类域论 (I, \ref{thm:1.1}) 给出同态 (局部 Artin 映射)
\[
\phi_v\colon K_v^{\times}\to D(w)\subset G.
\]

\begin{lemma}\label{lem:v.5.1}
$G$ 的子群 $D(w)$ 与映射 $\phi_v$ 不依赖于位于 $v$ 之上的素 $w$ 的选取.
\end{lemma}

\begin{proof}
位于 $v$ 之上的其他素形如 $\sigma w$, $\sigma\in G$, 且 $\sigma\colon L\to L$ 由连续性延拓为保持 $K_v$ 的同态 $\sigma\colon L_w\to L_{\sigma w}$. 我们有
\[
D(\sigma w)=\sigma D(w)\sigma^{-1},
\]
而它等于 $D(w)$, 因为 $G$ 交换.

设 $\Omega$ 与 $\Omega'$ 分别是 $K_v$ 的包含 $L_w$ 与 $L_{\sigma w}$ 的极大 Abel 扩张. 由第 \ref{ch:III} 章, 我们有局部 Artin 映射 $\phi_v\colon K^{\times}\to\Gal(\Omega/K_v)$ 与 $\phi_v'\colon K^{\times}\to\Gal(\Omega'/K_v)$ 同构 $\tilde{\sigma}\colon\Omega\to\Omega'$ 的选取确定同构
\[
\rho\mapsto\tilde{\sigma}\circ\rho\circ\tilde{\sigma}^{-1}\colon\Gal(\Omega/K_v)\to\Gal(\Omega'/K_v),
\]
它不依赖于 $\tilde{\sigma}$. 此外它与 $\phi_v$ 的复合是 $\phi_v'$ (因为它满足刻画 $\phi_v'$ 的条件).
\end{proof}

\begin{proposition}\label{prop:v.5.2}
存在唯一的连续同态 $\phi_K\colon\mathbb{I}\to\Gal(K^{\mathrm{ab}}/K)$ 具有如下性质: 对任何含于 $K^{\mathrm{ab}}$ 的 $K$ 的有限扩张 $L$ 及 $L$ 中位于 $K$ 的素 $v$ 之上的素 $w$, 图
\[
\begin{tikzcd}
K_v^{\times}\ar[r, "\phi_v"]\ar[d] & \Gal(L_w/K_v)\ar[d, "\text{限制}\ |\!\ _L"]\\
\mathbb{I}_K\ar[r, "\mathbf{a}\mapsto\phi_K(\mathbf{a})|L"'] & \Gal(L/K)
\end{tikzcd}
\]
交换.
\end{proposition}

\begin{proof}
设 $\mathbf{a}\in\mathbb{I}$, $L\subset K^{\mathrm{ab}}$ 在 $K$ 上有限. 若 $a_v\in U_v$ 且 $L_w/K_v$ 非分歧, 则 $\phi_v(a_v)=1$ (见 III, \ref{sec:III.1}). 因此除有限多个 $v$ 外 $\phi_v(a_v)=1$, 故可定义
\[
\phi_{L/K}(\mathbf{a})=\prod_v\phi_v(a_v)
\]
(乘积在 $\Gal(L/K)$ 内). 显然 $\phi_{L/K}$ 是使上图交换的唯一同态.

若 $L'\supset L$, 则局部 Artin 映射的性质表明 $\phi_{L'/K}(\mathbf{a})|L=\phi_{L/K}(\mathbf{a})$. 因此存在唯一的同态 $\phi\colon\mathbb{I}\to\Gal(K^{\mathrm{ab}}/K)$ 使得 $\phi(\mathbf{a})|L=\phi_{L/K}(\mathbf{a})$ 对所有含于 $K^{\mathrm{ab}}$ 的 $K$ 的有限扩张 $L$ 成立.

再由局部 Artin 映射的性质, 对任何 $K\subset K'\subset L\subset K^{\mathrm{ab}}$ ($L$ 在 $K$ 上有限), 图
\[
\begin{tikzcd}
\mathbb{I}_{K'}^S\ar[r, "\phi_{L/K'}"]\ar[d, "\Nm"'] & \Gal(L/K')\ar[d]\\
\mathbb{I}_K^S\ar[r, "\phi_{L/K}"'] & \Gal(L/K)
\end{tikzcd}
\]
交换. 取 $K'=L$ 可知 $\Nm_{L/K}(\mathbb{I}_L^S)$ 含于 $\phi_{L/K}$ 的核. 特别地, $\phi_{L/K}$ 的核包含 $\mathbb{I}_K^S$ 的一个开子群 (推论 \ref{cor:v.4.13}), 这蕴含 $\phi_K$ 连续.
\end{proof}

\begin{theorem}[互反律]\label{thm:v.5.3}
同态 $\phi_K\colon\mathbb{I}_K\to\Gal(K^{\mathrm{ab}}/K)$ 具有如下性质:

(a) $\phi_K(K^{\times})=1$;

(b) 对 $K$ 的每个有限 Abel 扩张 $L$, $\phi_K$ 定义同构
\[
\phi_{L/K}\colon\mathbb{I}_K/(K^{\times}\cdot\Nm(\mathbb{I}_L))\to\Gal(L/K).
\]
\end{theorem}

我们在命题的证明中已看到 $\phi_{L/K}(\Nm(\mathbb{I}_L))=1$, 故 (假设定理的 (a)) $\phi_{L/K}$ 确实经过 $\mathbb{I}_K/K^{\times}\cdot\Nm(\mathbb{I}_L)$ 分解. (b) 也可表述为: $\phi$ 定义同构
\[
\phi_{L/K}\colon\mathbf{C}_K/\Nm(\mathbf{C}_L)\to\Gal(L/K).
\]

\begin{example}\label{ex:v.5.4}
定理的 (a) 说: 对每个 $b\in K^{\times}$, $\prod\phi_v(b)=1$. 把它应用于扩张 $K[a^{\frac{1}{n}}]/K$ (设 $K$ 含本原 $n$ 次单位根), 便得到 Hilbert 符号的乘积公式:
\[
\prod_v(a,b)_v=1.
\]
见 (III, \ref{rem:iii.4.8}).
\end{example}

\begin{theorem}[存在定理]\label{thm:v.5.5}
固定 $K$ 的一个代数闭包 $K^{\mathrm{al}}$; 对 $C_K$ 的每个有限指数开子群 $N$, 存在唯一的含于 $K^{\mathrm{al}}$ 的 $K$ 的 Abel 扩张 $L$ 使得 $\Nm_{L/K}C_L=N$.
\end{theorem}

$C_K$ 的子群称为\textbf{范群}, 若它形如 $\Nm(C_L)$, $L$ 为 $K$ 的某个有限 Abel 扩张. 存在定理表明: 范群恰好是 $C_K$ 的有限指数开子群. 若 $N$ 是这样的群, 则满足 $\Nm(C_L)=N$, 即 $N=\mathrm{Ker}(\phi_{L/K})$, 的 $K$ 的有限 Abel 扩张 $L$ 称为属于 $N$ 的 $K$ 的\textbf{类域}.

如上陈述的存在定理对所有整体域成立.

\begin{corollary}\label{cor:v.5.6}
映射 $L\mapsto\Nm(C_L)$ 是从 $K$ 的有限 Abel 扩张之集到 $C_K$ 的有限指数开子群之集的双射. 此外,
\begin{align*}
L_1\subset L_2&\iff\Nm(\mathbf{C}_{L_1})\supset\Nm(\mathbf{C}_{L_2});\\
\Nm(\mathbf{C}_{L_1\cdot L_2})&=\Nm(\mathbf{C}_{L_1})\cap\Nm(\mathbf{C}_{L_2});\\
\Nm(\mathbf{C}_{L_1\cap L_2})&=\Nm(\mathbf{C}_{L_1})\cdot\Nm(\mathbf{C}_{L_2}).
\end{align*}
\end{corollary}

\begin{remark}\label{rem:v.5.7}
\begin{enumerate}[(a)]
\item 在数域情形, 映射
\[
\phi_K\colon\mathbb{I}_K\to\Gal(K^{\mathrm{ab}}/K)
\]
是满射. 对 $K$ 的无限素 $v$, 记 $K_v^+$ 为 $K_v^{\times}$ 中含 $1$ 的连通分支; 于是依 $v$ 为复或实, $K_v^+$ 同构于 $\CC^{\times}$ 或 $\R_{>0}$. 显然 $\prod_{v|\infty}K_v^+\subset\mathrm{Ker}(\phi_K)$. 按定义 $K^{\times}\subset\mathrm{Ker}(\phi_K)$, 故 $K^{\times}\cdot(\prod_{v|\infty}K_v^+)\subset\mathrm{Ker}(\phi_K)$. 但 $\phi_K$ 是连续同态而 $\Gal(K^{\mathrm{ab}}/K)$ 是 Hausdorff 的, 故核是闭子群. 于是 $\mathrm{Ker}(\phi_K)$ 包含 $K^{\times}\cdot(\prod_{v|\infty}K_v^+)$ 的闭包. 这恰好就是核是一条定理. $K^{\times}\cdot(\prod_{v|\infty}K_v^+)$ 的闭包在 $C_K$ 中的像是 $C_K$ 中含 $1$ 的连通分支.

对 $K$ 的每个有限 Abel 扩张 $L$, Artin 映射定义同构 $C_K/\Nm(C_L)\to\Gal(L/K)$. 对 $L$ 取逆向极限, 右端得到与 $\Gal(K^{\mathrm{ab}}/K)$ 的同构, 于是问题化为计算体系 $C_K/\Nm(C_L)$ 的逆向极限. 存在定理表明群 $\Nm(C_L)$ 恰为有限指数开子群, 故这是一个拓扑学问题.

\item 在函数域情形, Artin 映射 $\phi_K\colon\mathbb{I}_K/K^{\times}\to\Gal(K^{\mathrm{ab}}/K)$ 是单射, 但不满 (其像稠密).
\end{enumerate}
\end{remark}

\begin{remark}\label{rem:v.5.8}
设整体 Artin 映射 $\phi\colon\mathbb{I}\to\Gal(K^{\mathrm{ab}}/K)$ 的核包含 $K^{\times}$. 则对 $K$ 的每个有限 Abel 扩张 $L/K$, $\phi_{L/K}\colon\mathbb{I}\to\Gal(L/K)$ 由一个容许模的同态 $\psi\colon I^S\to\Gal(L/K)$ 产生 (如命题 \ref{prop:v.4.7}). 此外, 由于 $\phi$ 是局部 Artin 映射之积, $\psi$ 必为理想语言的整体 Artin 映射 (因而它容许一个模). 由 idèle 版本的定理\ref{thm:v.5.3} 与 \ref{thm:v.5.5} 推出理想版本的定理\ref{thm:v.3.5} 与 \ref{thm:v.3.6} 是一个直接的练习. 我们将在第 \ref{ch:VII} 章证明定理\ref{thm:v.5.3} 与 \ref{thm:v.5.5}.
\end{remark}

\phantomsection\label{pg:orig180}%
\begin{NOTES}
在其 1951 年的报告 Sur la théorie du corps de classes 中, Weil 写道:\footnote{寻找 $C_k$ 的一个解释 ($k$ 为数域时), 使其在某种意义上类似于 $k$ 为函数域时由 Galois 群给出的解释, 在我看来是当前数论的基本问题之一; 也许这样的解释蕴含着 Riemann 假设的钥匙\ldots\ (原书脚注给出的英译, 此处译出.)}
\begin{quote}
\itshape La recherche d'une interprétation de $C_k$ si $k$ est un corps de nombres, analogue en quelque manière \`a l'interprétation par un groupe de Galois quand $k$ est un corps de fonctions, me semble constituer l'un des problèmes fondamentaux de la théorie des nombres \`a l'heure actuelle; il se peut qu'une telle interprétation renferme la clef de l'hypothèse de Riemann\ldots
\end{quote}
\end{NOTES}
关于此段的讨论见 \mo{41296}.

\subsection*{例}

\begin{lemma}\label{lem:v.5.9}
映射
\[
(r,t,(u_p))\mapsto(rt,ru_2,ru_3,ru_5,\ldots)\colon\Q^{\times}\times\R_{>0}\times\prod\Z_p^{\times}\to\mathbb{I}_{\Q}
\]
是拓扑群的同构 ($\Q^{\times}$ 取离散拓扑).
\end{lemma}

\begin{proof}
任何 idèle $\mathbf{a}=(a_\infty,a_2,\ldots,a_p,\ldots)$ 可写成
\[
\mathbf{a}=a(t,u_2,u_3,u_5,\ldots),\qquad a\in\Q^{\times},\ t\in\R_{>0},\ u_p\in\Z_p^{\times},
\]
\,---\,取 $a=(\mathrm{sign}(a_\infty))\prod p^{\mathrm{ord}_p(a_p)}$, $t=a_\infty/a$, $u_p=a_p/a$. 而且表达式唯一, 因为对所有的 $p$ 都是 $p$-adic 单位的正有理数只有 $1$.

形如
\[
\{1\}\times U\times\prod_{p\ \mathrm{finite}}U_p
\]
的集合 ($U$, $U_p$ 分别为 $1$ 在 $\R^{\times}$, $\Q_p^{\times}$ 中的开邻域, 且除有限多个 $p$ 外 $U_p=\Z_p^{\times}$) 在左边构成 $1$ 的基本邻域组, 在右边亦然.
\end{proof}

于是有拓扑群的典范同构
\[
C_{\Q}\to\R_{>0}\times\prod_{p\ \mathrm{finite}}\Z_p^{\times}=\R_{>0}\times\hat{\Z}^{\times}.
\]
令
\[
\Q^{\mathrm{cyc}}=\bigcup\Q[\zeta_n].
\]
此情形下, 整体互反映射是
\[
\phi\colon\mathbb{I}_{\Q}\to\hat{\Z}^{\times}\to\Gal(\Q^{\mathrm{cyc}}/\Q)
\]
的\guillemotleft 逆\guillemotright, 其中 $\mathbb{I}_{\Q}\to\hat{\Z}^{\times}$ 是上述投影映射, 而 $\hat{\Z}^{\times}\to\Gal(\Q^{\mathrm{cyc}}/\Q)$ 是典范同构 (见 I A.5c).

\begin{summary}\label{note:v.5.10}
\phantomsection\hypertarget{item:v.5.10}{}
设 $K$ 是代数数域. 存在连续满同态 (互反映射或 Artin 映射)
\[
\phi_K\colon\mathbb{I}_K\to\Gal(K^{\mathrm{ab}}/K)
\]
使得对每个含于 $K^{\mathrm{ab}}$ 的 $K$ 的有限扩张 $L$, $\phi_K$ 给出交换图
\[
\begin{tikzcd}
\mathbb{I}_K/K^{\times}\ar[r, "\phi_K"]\ar[d] & \Gal(K^{\mathrm{ab}}/K)\ar[d, "\tau\mapsto\tau|L"]\\
\mathbb{I}_K/(K^{\times}\cdot\Nm(\mathbb{I}_L))\ar[r, "\phi_{L/K}"', "\simeq"] & \Gal(L/K)
\end{tikzcd}
\]
它由下列两条性质确定:

(a) 对每个满足如下条件的 $u=(u_v)\in\mathbb{I}_K$, $\phi_{L/K}(u)=1$:

\quad i) 若 $v$ 在 $L$ 中不分歧, 则 $u_v$ 是单位;

\quad ii) 若 $v$ 在 $L$ 中分歧, 则 $u_v$ 充分接近 $1$ (只依赖于 $L/K$);

\quad iii) 若 $v$ 是实素但在 $L$ 中变复, 则 $u_v>0$.

(b) 对 $K$ 中在 $L$ 不分歧的每个素 $v$, idèle
\[
\alpha=(1,\ldots,1,\overset{v}{\pi},1,\ldots),\qquad \pi\ \text{为}\ \mathcal{O}_v\ \text{的素元},
\]
映到 $\Gal(L/K)$ 中的 Frobenius 元 $(\mathfrak{p}_v,L/K)$.

为看出满足这些条件的映射至多一个, 设 $\alpha\in\mathbb{I}_K$, 用弱逼近定理取 $a\in K^{\times}$, 使其对所有在 $L$ 分歧或变复的素 $v$ 接近 $\alpha_v$. 则 $\alpha=au\beta$, 其中 $u$ 是 (a) 中的 idèle, $\beta$ 是形如 (b) 的 idèle 的有限乘积. 此时 $\phi_{L/K}(\alpha)=\phi_{L/K}(\beta)$, 后者可用 (b) 计算.

对 $K=\Q$, Artin 映射经过 $\{\pm\}\times\mathbb{I}_f/\Q^{\times}$ 分解, 而该商群的每个元素都被 $\hat{\Z}^{\times}\subset\mathbb{I}_f$ 中的元素唯一代表. 此情形下得到图
\[
\begin{tikzcd}
\hat{\Z}^{\times}\ar[r, "\phi_{\Q}", "\simeq"'] & \Gal(\Q^{\mathrm{ab}}/\Q) & =\bigcup_m\Q[\zeta_m]\\
(\Z/m\Z)^{\times}\ar[r, "\simeq"']\ar[u] & \Gal(\Q[\zeta_m]/\Q)\ar[u, hookrightarrow]
\end{tikzcd}
\]
(下方的同构把 $[n]$ 送到自同构 $\zeta_m\mapsto\zeta_m^n$.)
它交换, 且有逆. 这可以通过把 idèle $\alpha$ 写成上述形状 $au\beta$ 来验证, 但看一个例子更有启发. 设 $p$ 是不整除 $m$ 的素数, 并令
\[
\alpha=p\cdot(1,\ldots,1,p_p^{-1},1,\ldots)\in\Z\cdot\mathbb{I}_f=\mathbb{I}_f.
\]
则
\[
\alpha=(p,\ldots,p,\overset{p}{1},p,\ldots)\in\hat{\Z}^{\times}
\]
在 $(\Z/m\Z)^{\times}$ 中的像为 $[p]$, 而 $[p]$ 作为 $(p,\Q[\zeta_m]/\Q)$ 作用在 $\Q[\zeta_m]$ 上. 另一方面, $\phi_{\Q}(\alpha)=\phi_{\Q}((1,\ldots,1,p_p^{-1},1,\ldots))$, 它作为 $(p,\Q[\zeta_m]/\Q)^{-1}$ 作用.

\end{summary}
\begin{exercise}\label{exc:v.5.11}
证明: 即使 $K$ 是数域, $K$ 的 idèle 类群也有有限指数但不开的子群.
\end{exercise}

    % V   整体类域论: 陈述
% ============================================================
% 第六章  L-级数与素数的密度
% 源: Milne CFT v4.03, Chapter VI (书页 p.183--197)
% OCR: 行 6797--7441; 脚注 1--3 从原书 PDF 恢复
% 勘误: CFT 官网勘误表无本章条目.
% 标签前缀: vi.  本章无跨节编号公式
% ============================================================
\chapter{L-级数与素数的密度}\label{ch:VI}
\index{ljishu@L-级数(L-series)}\index{sulixiang@素理想(prime ideal)!midu@密度(density)}

Euler 用初等形式的 Riemann zeta 函数证明了素数有无穷多个. 为了证明素数在模 $m$ 的各个算术级数中均匀分布, Dirichlet 给 $(\Z/m\Z)^{\times}$ 的特征附加了 $L$-级数 (视作实变量的函数). Riemann 开创了把 Riemann zeta 函数当作复变量函数来研究的方法. 本章中, 我们将 (遵循 Weber) 把 Dirichlet 的方法推广到研究素理想在射线类群各个类中的分布. 除射线类群的定义外, 本章独立于前面各章.

本章需要用到一点复分析. 回顾: 幂级数 $1+z+\frac{z^2}{2!}+\cdots$ 对所有 $z\in\CC$ 收敛到一个全纯函数, 该函数记作 $e^z$. 对每个正实数 $n$ 与复数 $z$, 定义 $n^z=e^{(\log n)z}$, 其中 $\log$ 是自然对数 (函数 $\R_{>0}\to\R$, 为 $e^r$ 的反函数).

\section{Dirichlet 级数与 Euler 乘积}\label{sec:VI.1}

\textbf{Dirichlet 级数} (Dirichlet series) 是形如
\[
f(s)=\sum_{n\ge1}\frac{a(n)}{n^{s}}\qquad a(n)\in\CC,\quad s=\sigma+it\in\CC
\]
的级数. 属于数域 $K$ 的 \textbf{Euler 乘积} (Euler product) 是形如
\[
g(s)=\prod_{\mathfrak{p}}\frac{1}{(1-\theta_{1}(\mathfrak{p})\Nm\mathfrak{p}^{-s})\cdots(1-\theta_{d}(\mathfrak{p})\Nm\mathfrak{p}^{-s})},\quad\theta_{i}(\mathfrak{p})\in\CC,\quad s\in\CC
\]
的乘积, 其中 $\mathfrak{p}$ 跑遍 $\mathcal{O}_K$ 的除有限多个之外的全部素理想.

\begin{example}\label{ex:vi.1.1}
(a) Riemann zeta 函数是
\[
\zeta(s)=\sum_{n\ge1}\frac{1}{n^{s}}=\prod_{p}\frac{1}{1-p^{-s}}.
\]
众所周知, $\zeta(s)$ 的性态, 特别是在临界带 $0\le\Re(s)\le1$ 中的性态, 与素数的分布有关.

(b) Dedekind zeta 函数. 对每个数域 $K$,
\[
\zeta_K(s)=\sum_{\mathfrak{a}\ge0}\frac{1}{\Nm\mathfrak{a}^{s}}=\prod_{\mathfrak{p}}\frac{1}{1-\Nm\mathfrak{p}^{-s}}.
\]
这里 $\Nm\mathfrak{a}=(\mathcal{O}_K:\mathfrak{a})$. 和式跑遍 $\mathcal{O}_K$ 中的整理想, 乘积跑遍 $\mathcal{O}_K$ 中的素理想.

(c) \textbf{Dirichlet 特征} (Dirichlet character) 是同态\footnote{当 $K=\Q$ 且 $\mathfrak{m}=\infty(m)$, 从而 $C_{\mathfrak{m}}=(\Z/m\Z)^{\times}$ 时, 这些特征与 $L$-级数由 Dirichlet 引入; 对一般的射线类群, 它们由 Weber 引入. 一些作者 (有时也包括本书作者本人) 把术语 ``Dirichlet 特征''与``Dirichlet $L$-级数''限制于 $K=\Q$ 的情形, 而把更一般的对象称为``Weber 特征''与``Weber $L$-级数''. Dirichlet 用字母 $L$ 表示他的 $L$-函数, 这个字母沿用至今.}
\[
\chi\colon I^{\mathfrak{m}}\to\CC^{\times}
\]
其核包含 $i(K_{\mathfrak{m},1})$ (对某个模 $\mathfrak{m}$), 即 $\chi$ 是射线类群 $C_{\mathfrak{m}}$ 的特征. 对这样的特征, 相应的 Dirichlet $L$-级数是
\[
L(s,\chi)=\sum_{\substack{\mathfrak{a}\subset\mathcal{O}_{K},\\(\mathfrak{m},\mathfrak{a})=1}}\frac{\chi(\mathfrak{a})}{\Nm\mathfrak{a}^{s}}=\prod_{(\mathfrak{m},\mathfrak{p})=1}\frac{1}{1-\chi(\mathfrak{p})\Nm\mathfrak{p}^{-s}}.
\]

(d) \textbf{Hecke 特征} (Hecke character, 或 Gr\"ossen character) 是连续同态
\[
\psi:\mathbb{I}_K/K^{\times}\to\CC^{\times}
\]
其像包含于单位圆周. 若它在 $\mathbb{I}_K$ 于无限素处的单位分支上取值 $1$, 则它经过某个 $C_{\mathfrak{m}}$ 分解, 从而是 Dirichlet 特征; 反之, Dirichlet 特征定义一个具有离散像的 Hecke 特征. Hecke 特征在某个 $\prod_{v\notin S}U_v$ ($S$ 为包含无限素的素的有限集) 上取值 $1$, 相应的 Hecke $L$-级数是
\[
{\cal L}_{S}(s,\psi)=\prod_{v\not\in S}\frac{1}{1-\psi(\pi_{v})\Nm\mathfrak{p}_{v}^{-s}},
\]
其中 $\pi_v$ 是一个 idèle, 其 $v$-分量为一个素元, 其余分量为 $1$.

(e) 设 $L$ 为 $K$ 的有限 Galois 扩张, Galois 群为 $G$. 设 $V$ 是 $\CC$ 上的有限维向量空间,
\[
\rho:G\to\mathrm{GL}(V)
\]
是 $G$ 到 $V$ 的线性自同构群的同态. 我们称 $\rho$ 为 $G$ 的一个 (有限维)\textbf{表示} (representation). $\rho$ 的迹 (trace) 是把 $\sigma$ 映到自同构 $\rho(\sigma)$ 的迹的映射. 对 $\sigma\in G$, 令
\[
P_{\sigma}(T)\overset{\mathrm{def}}{=}\det(1-\rho(\sigma)T\mid V)=\prod_{i=1}^{\dim V}(1-a_{i}(\sigma)T),\qquad a_{i}\in\CC,
\]
这是 $\rho(\sigma)$ 的特征多项式. 由于 $P_{\sigma}(T)$ 只依赖于 $\sigma$ 的共轭类, 对每个在 $L$ 中不分歧的 $K$ 的素 $\mathfrak{p}$, 可定义 $P_{\mathfrak{p}}(T)$ 为 $(\mathfrak{P},L/K)$ 的特征多项式, 其中 $\mathfrak{P}$ 是 $L$ 中位于 $\mathfrak{p}$ 之上的任一素理想. 附加到 $\rho$ 上的 \textbf{Artin $L$-级数} (Artin $L$-series) 是
\[
L(s,\rho)=\prod_{\mathfrak{p}}\frac{1}{P_{\mathfrak{p}}(\Nm\mathfrak{p}^{-s})}=\prod_{\mathfrak{p}}\frac{1}{(1-a_{1}(\mathfrak{p})\Nm\mathfrak{p}^{-s})\cdots(1-a_{\dim V}(\mathfrak{p})\Nm\mathfrak{p}^{-s})}
\]
(乘积跑遍 $K$ 的全部不分歧素; $a_i(\mathfrak{p})=a_i((\mathfrak{P},L/K))$).
\end{example}

\section{收敛性结果}\label{sec:VI.2}

我们来研究 Dirichlet 级数与 Euler 乘积的初等解析性质.

\subsection*{Dirichlet 级数}

\begin{proposition}\label{prop:vi.2.1}
设
\[
f(s)=\sum_{n\ge1}\frac{a(n)}{n^{s}}.
\]
令 $S(x)=\sum_{n\le x}a(n)$, 并设存在正常数 $a$ 与 $b$, 使得对所有充分大的 $x$ 有 $|S(x)|\le ax^{b}$. 则级数 $f(s)$ 在
\[
D(b,\delta,\varepsilon)=\{\Re(s)\ge b+\delta,\quad|\arg(s-b)|\le\frac{\pi}{2}-\varepsilon\}
\]
上对一切 $\delta,\varepsilon>0$ 一致收敛, 且它收敛到半平面 $\Re(s)>b$ 上的一个解析函数.
\end{proposition}

\begin{proof}
由于每个满足 $\Re(s)>b$ 的点 $s$ 都有形如 $D(b,\delta,\varepsilon)$ 的邻域, 断言的第二部分由第一部分推出. 为证第一部分, 我们使用一致收敛的 Cauchy 准则. 对充分大的整数 $n_1<n_2$,
\begin{align*}
\left|\sum_{n=n_{1}}^{n_{2}}\frac{a(n)}{n^{s}}\right|
&=\left|\sum_{n_{1}}^{n_{2}}\frac{S(n)-S(n-1)}{n^{s}}\right|\\
&=\left|\sum_{n_{1}}^{n_{2}}\frac{S(n)}{n^{s}}-\sum_{n_{1}-1}^{n_{2}-1}\frac{S(n)}{(n+1)^{s}}\right|\\
&=\left|\frac{S(n_{2})}{n_{2}^{s}}-\frac{S(n_{1}-1)}{n_{1}^{s}}+\sum_{n_{1}}^{n_{2}-1}S(n)\left(\frac{1}{n^{s}}-\frac{1}{(n+1)^{s}}\right)\right|\\
&\le\frac{|S(n_{2})|}{n_{2}^{\sigma}}+\frac{|S(n_{1}-1)|}{n_{1}^{\sigma}}+\sum_{n_{1}}^{n_{2}-1}|S(n)|\left|s\int_{n}^{n+1}\frac{dt}{t^{s+1}}\right|\\
&\le\frac{a}{n_{2}^{\sigma-b}}+\frac{a}{n_{1}^{\sigma-b}}+\sum_{n_{1}}^{n_{2}-1}|s|an^{b}\left|\int_{n}^{n+1}\frac{dt}{t^{s+1}}\right|\\
&\le\frac{2a}{n_{1}^{\sigma-b}}+|s|a\sum_{n_{1}}^{n_{2}-1}\left|\int_{n}^{n+1}\frac{t^{b}\,dt}{t^{s+1}}\right|\\
&\le\frac{2a}{n_{1}^{\sigma-b}}+|s|a\int_{n_{1}}^{\infty}\frac{dt}{t^{\sigma+1-b}}\\
&\le\frac{2a}{n_{1}^{\sigma-b}}-\frac{|s|a}{\sigma-b}\cdot\frac{1}{t^{\sigma-b}}\bigg|_{n_{1}}^{\infty}\\
&\le\frac{2a}{n_{1}^{\sigma-b}}+\frac{|s|a}{(\sigma-b)n_{1}^{\sigma-b}}.
\end{align*}
但对 $s\in D(b,\delta,\varepsilon)$,
\[
\frac{|s|}{\sigma-b}=\frac{|s-b+b|}{\sigma-b}\le\frac{|s-b|}{\sigma-b}+\frac{b}{\sigma-b}=\frac{1}{\cos\theta}+\frac{b}{\sigma-b}\le\frac{1}{\cos\theta}+\frac{b}{\delta},
\]
其中 $\theta=\arg(s-b)$. 由于 $|\theta|\le\frac{\pi}{2}-\varepsilon$, $\frac{1}{\cos\theta}$ 被某个数 $M$ 界住, 于是
\[
\left|\sum_{n=n_{1}}^{n_{2}}\frac{a(n)}{n^{s}}\right|\le\frac{2a}{n_{1}^{\sigma-b}}+\frac{(M+\frac{b}{\delta})a}{n_{1}^{\sigma-b}}.
\]
上式右端当 $n_1\to\infty$ 时趋于零, 故可应用 Cauchy 准则推知 $f(s)$ 一致收敛.
\end{proof}

\begin{remark}\label{rem:vi.2.2}
(a) 对 Dirichlet 级数 $\zeta(s)$, $S(x)$ 为 $[x]$, 故 $\zeta(s)$ 的级数对 $\Re(s)>1$ 收敛. 对 $\zeta_K(s)$, $S(x)$ 是数值范数 $\le x$ 的整理想的个数. 显然 $S(x)$ 有限, 而实际上 (见下文 \ref{prop:vi.2.8}) $S(x)\le Cx$. 因此 $\zeta_K$ (以及 $L(s,\chi)$) 的级数对 $\Re(s)>1$ 收敛. (也可以直接证明 Euler 乘积对 $\Re(s)>1$ 收敛, 从而推出 Dirichlet 级数收敛; 见 \cite{FrohlichTaylor1991}, VIII, 2.2.)

(b) 设 $f(s)=\sum\frac{a(n)}{n^{s}}$ 是满足 $a(n)\ge0$ 的 Dirichlet 级数. 若 $f(s)$ 对所有 $\Re(s)>b$ 的 $s$ 收敛, 但对任何 $\varepsilon>0$ 都不在半平面 $\{s\mid\Re(s)>b-\varepsilon\}$ 上收敛, 则当实数 $s$ 从右侧趋于 $1$ 时 $f(s)\to\infty$, 也就是说, $f(s)$ 的收敛域受到实轴上一个奇点的限制. (见 \cite{Serre1970}, III, 2.3.) 例如, $\zeta(s)$ 的级数在任何半平面 $\Re(s)>1-\varepsilon$ ($\varepsilon>0$) 上都不收敛, 而我们将会看到, $\zeta(s)$ 在 $s=1$ 处确实有极点.
\end{remark}

\begin{lemma}\label{lem:vi.2.3}
zeta 函数 $\zeta(s)$ 可解析延拓为 $\Re(s)>0$ 上的亚纯函数, 其唯一 (可能的) 极点在 $s=1$ 处.
\end{lemma}

\begin{proof}
定义
\[
\zeta_{2}(s)=1-\frac{1}{2^{s}}+\frac{1}{3^{s}}-\frac{1}{4^{s}}+\cdots
\]
对这一 Dirichlet 级数, $S(x)$ 为 $0$ 或 $1$, 故 $\zeta_2(s)$ 对 $\Re(s)>0$ 解析. 注意
\[
\left(1+\frac{1}{2^{s}}+\frac{1}{3^{s}}+\frac{1}{4^{s}}+\cdots\right)-2\left(\frac{1}{2^{s}}+\frac{1}{4^{s}}+\frac{1}{6^{s}}+\cdots\right)=1-\frac{1}{2^{s}}+\frac{1}{3^{s}}-\frac{1}{4^{s}}+\cdots,
\]
即
\[
\zeta(s)-\frac{2}{2^{s}}\zeta(s)=\zeta_{2}(s),
\]
亦即
\[
\zeta(s)=\frac{\zeta_{2}(s)}{1-2^{1-s}}.
\]
于是 $\zeta(s)$ 在 $\Re(s)>0$ 上解析, 可能的极点只在 $2^{s-1}=1$ 处. 但
\[
2^{s-1}=1\iff e^{(\log 2)(s-1)}=1\iff(\log 2)(s-1)=2k\pi i,
\]
故 $\zeta(s)$ 除可能在
\[
s=1+\frac{2k\pi i}{\log 2},\quad k\in\Z
\]
处外解析. 实际上, 唯一可能的极点是 $s=1$. 为此, 定义
\[
\zeta_{3}(s)=1+\frac{1}{2^{s}}-\frac{2}{3^{s}}+\frac{1}{4^{s}}+\frac{1}{5^{s}}-\frac{2}{6^{s}}+\cdots
\]
并观察 (如同对 $\zeta_2(s)$ 那样) $\zeta_3(s)$ 对 $\Re(s)>0$ 解析, 且
\[
\zeta(s)=\frac{\zeta_{3}(s)}{1-3^{1-s}}.
\]
因此 $\zeta(s)$ 对 $\Re(s)>0$ 解析, 可能的极点只在
\[
s=1+\frac{2k\pi i}{\log 3},\quad k\in\Z
\]
处. 于是, 在 $\zeta(s)$ 的极点处必须有
\[
\frac{2k\pi i}{\log 2}=\frac{2k'\pi i}{\log 3},
\]
即
\[
2^{k'}=3^{k},\qquad k,k'\in\Z.
\]
由唯一分解, 这仅当 $k=0=k'$ 时才可能.
\end{proof}

\begin{lemma}\label{lem:vi.2.4}
对实数 $s>1$,
\[
\frac{1}{s-1}\le\zeta(s)\le 1+\frac{1}{s-1}.
\]
因此 $\zeta(s)$ 在 $s=1$ 处有留数为 $1$ 的单极点, 即
\[
\zeta(s)=\frac{1}{s-1}+\text{在 $1$ 附近全纯的函数}.
\]
\end{lemma}

\begin{proof}
固定实数 $s>1$. 考察 $y=x^{-s}$ 的图象即知
\[
\int_{1}^{\infty}x^{-s}\,dx\le\zeta(s)\le 1+\int_{1}^{\infty}x^{-s}\,dx.
\]
而
\[
\int_{1}^{\infty}x^{-s}\,dx=\left.\frac{x^{1-s}}{1-s}\right|_{1}^{\infty}=\frac{1}{s-1},
\]
这就给出了不等式. 由于 $\zeta(s)$ 在 $s=1$ 附近亚纯, 在 $s=1$ 附近有
\[
\zeta(s)=\frac{c}{(s-1)^{m}}+\frac{g(s)}{(s-1)^{m-1}},
\]
其中 $m\in\N$, $c\in\CC$, $g(s)$ 在 $s=1$ 附近全纯. 上述不等式蕴含 $m=1$ 且 $c=1$.
\end{proof}

\begin{proposition}\label{prop:vi.2.5}
设 $f(s)$ 为一 Dirichlet 级数, 且存在实常数 $C$ 与 $b<1$, 使得
\[
|S(n)-a_{0}n|\le Cn^{b}.
\]
则 $f(s)$ 可延拓为 $\Re(s)>b$ 上的亚纯函数, 在 $s=1$ 处有留数为 $a_0$ 的单极点, 即在 $s=1$ 附近
\[
f(s)=\frac{a_{0}}{s-1}+\text{全纯函数}.
\]
\end{proposition}

\begin{proof}
对 Dirichlet 级数 $f(s)-a_0\zeta(s)$, 有 $|S(n)|\le Cn^{b}$, 故 $f(s)-a_0\zeta(s)$ 对 $\Re(s)>b$ 收敛.
\end{proof}

\subsection*{Euler 乘积}

回顾: 无穷乘积 $\prod_{n=1}^{\infty}(1+b_n)$ ($b_n\in\CC$, $b_n\ne-1$) 称为\textbf{收敛}的 (convergent), 如果部分乘积序列
\[
\Pi_{m}=\prod_{n=1}^{m}(1+b_{n})
\]
收敛到一个非零值. 此外, 若 $\prod_{n=1}^{\infty}1+|b_{n}|$ 收敛, 则称该乘积\textbf{绝对收敛} (absolutely convergent). 绝对收敛的乘积必收敛, 且此时乘积的任何重排都 (绝对) 收敛到同一个值.

\begin{lemma}\label{lem:vi.2.6}
乘积 $\prod_{1}^{\infty}1+b_n$ 绝对收敛当且仅当级数 $\sum b_n$ 绝对收敛.
\end{lemma}

\begin{proof}
不妨设 $b_n\ge0$ 对所有 $n$ 成立. 此时 $\Pi_m\overset{\text{def}}{=}\prod_1^m 1+b_n$ 与 $\Sigma_m\overset{\text{def}}{=}\sum_1^m b_n$ 都是单调递增序列. 由于 $\Pi_m\ge\Sigma_m$, 显然若 $\Pi_m$ 收敛则 $\Sigma_m$ 也收敛. 反之, 注意
\[
e^{\Sigma_{m}}=\prod_{i=1}^{m}e^{b_{i}}\ge\prod_{i=1}^{m}(1+b_{i})=\Pi_{m},
\]
故若序列 $\Sigma_m$ 收敛, 则序列 $\Pi_m$ 有上界, 从而也收敛.
\end{proof}

回顾: 有限和的乘积, 比如
\[
\left(\sum_{i=1}^{l}a_{i}\right)\left(\sum_{i=1}^{m}b_{i}\right)\left(\sum_{i=1}^{n}c_{i}\right)
\]
是一个和
\[
\sum_{\substack{1\le i\le l\\1\le j\le m\\1\le k\le n}}a_{i}b_{j}c_{k},
\]
其中每一项恰好由从每个和式中各取一项相乘而成. 又回顾
\[
\frac{1}{1-t}=1+t+t^{2}+\cdots,\qquad|t|<1.
\]
于是 (至少形式上)
\begin{align*}
\prod_{p}\frac{1}{1-p^{-s}}&=(1+2^{-s}+(2^{2})^{-s}+\cdots)(1+3^{-s}+(3^{2})^{-s}+\cdots)(1+5^{-s}+(5^{2})^{-s}+\cdots)\cdots\\
&=\sum n^{-s},
\end{align*}
因为每个正整数都能以唯一的方式写成素数幂的乘积. 这个恒等式有时被称为唯一分解的解析形式. 现在我们证明一个更一般的结果.

\begin{proposition}\label{prop:vi.2.7}
设 $\chi$ 是数域 $K$ 的 Dirichlet 特征. 对所有满足 $\Re(s)>1$ 的 $s$, Euler 乘积 $\prod_{\mathfrak{p}\nmid\mathfrak{m}}\frac{1}{1-\chi(\mathfrak{p})\Nm\mathfrak{p}^{-s}}$ 收敛到 $L(s,\chi)$.
\end{proposition}

\begin{proof}
对 $\Re(s)>1$,
\[
\frac{1}{1-\chi(\mathfrak{p})\Nm\mathfrak{p}^{-s}}=1+\frac{\chi(\mathfrak{p})}{\Nm\mathfrak{p}^{s}}+\frac{\chi(\mathfrak{p}^{2})}{(\Nm\mathfrak{p}^{2})^{s}}+\cdots.
\]
现在
\[
\prod_{\substack{(\mathfrak{p},\mathfrak{m})=1\\ \Nm\mathfrak{p}\le t_{0}}}\frac{1}{1-\chi(\mathfrak{p})\Nm\mathfrak{p}^{-s}}=\sum\frac{\chi(\mathfrak{a})}{\Nm(\mathfrak{a})^{s}},
\]
其中右端的和跑遍所有能表示为数值范数 $\le t_0$ 的素理想之乘积的整理想. 当 $t_0\to\infty$ 时, 右端 (绝对) 收敛到 $L(s,\chi)$. 因此该无穷乘积收敛, 且其值为 $L(s,\chi)$.
\end{proof}

\subsection*{部分 zeta 函数; 留数公式}

设 $K$ 为数域, $\mathfrak{m}$ 为一个模. 对 $C_{\mathfrak{m}}\overset{\mathrm{def}}{=}I^{\mathfrak{m}}/i(K_{\mathfrak{m},1})$ 中的每个类 $\mathfrak{k}$, 定义\textbf{部分 zeta 函数} (partial zeta function)
\[
\zeta(s,\mathfrak{k})=\sum_{\mathfrak{a}\ge0,\ \mathfrak{a}\in\mathfrak{k}}\frac{1}{\Nm\mathfrak{a}^{s}}\qquad(\text{对 $\mathfrak{k}$ 中的整理想求和}).
\]
注意: 对 $C_{\mathfrak{m}}$ 的每个特征 $\chi$,
\[
L(s,\chi)=\sum_{\mathfrak{k}\in C_{\mathfrak{m}}}\chi(\mathfrak{k})\zeta(s,\mathfrak{k}).
\]
特别地,
\[
\zeta_K(s)=\sum_{\mathfrak{k}\in C_{\mathfrak{m}}}\zeta(s,\mathfrak{k}).
\]
因此, 知道 $\zeta(s,\mathfrak{k})$ 就会为我们提供关于 $L(s,\chi)$ 与 $\zeta_K(s)$ 的信息.

令
\[
S(x,\mathfrak{k})=\left|\left\{\mathfrak{a}\in\mathfrak{k}\mid\mathfrak{a}\ \text{为整理想},\ \Nm\mathfrak{a}\le x\right\}\right|,
\]
即它是 Dirichlet 级数 $\zeta(s,\mathfrak{k})$ 的 $S(x)$. 回顾 (ANT, 第 5 章, p.~87): 存在同态
\[
l\colon U\to\R^{r+s},\quad u\mapsto(\log|\sigma_{1}(u)|,\ldots,2\log|\sigma_{r+s}(u)|),
\]
其核为 $U$ 的挠子群, 其像是一个 $r+s-1$ 维的格. \textbf{regulator} (调整子) $\mathrm{reg}(K)$ 定义为该格的基本平行多面体的体积. 令 $U_{\mathfrak{m},1}=U\cap K_{\mathfrak{m},1}$. 则 $U_{\mathfrak{m},1}$ 在 $U$ 中指数有限, 我们定义 $\mathrm{reg}(\mathfrak{m})$ 为 $l(U_{\mathfrak{m},1})$ 的基本平行多面体的体积. 于是
\[
\mathrm{reg}(\mathfrak{m})=\mathrm{reg}(K)\,(U:U_{\mathfrak{m},1}).
\]

\begin{proposition}\label{prop:vi.2.8}
对所有 $x\ge1$,
\[
|S(x,\mathfrak{k})-g_{\mathfrak{m}}x|\le Cx^{1-\frac{1}{d}},\quad g_{\mathfrak{m}}=\frac{2^{r}(2\pi)^{s}\,\mathrm{reg}(\mathfrak{m})}{w_{\mathfrak{m}}\Nm(\mathfrak{m})|\Delta_{K/\Q}|^{\frac{1}{2}}},\quad d=[K:\Q],
\]
其中
\begin{align*}
r&=\text{实素的个数},\\
s&=\text{复素的个数},\\
w_{\mathfrak{m}}&=\text{$K_{\mathfrak{m},1}$ 中 $1$ 的根的个数},\\
\Nm(\mathfrak{m})&=\Nm(\mathfrak{m}_{0})\,2^{r_{0}},
\end{align*}
$r_0$ 为 $\mathfrak{m}$ 中实素的个数, 而 $\Delta_{K/\Q}$ 为 $K/\Q$ 的判别式.
\end{proposition}

\begin{proof}
首先证明存在整理想 $\mathfrak{b}_0\in\mathfrak{t}^{-1}$. 则对每个整的 $\mathfrak{a}\in\mathfrak{t}$, 有 $\mathfrak{a}\mathfrak{b}_0=(\alpha)$, 其中 $\alpha\in\mathcal{O}_K$. 此时 $S(x,\mathfrak{t})$ 就是满足 $\alpha\in\mathfrak{b}_0\cap K_{\mathfrak{m},1}$ 且 $|\Nm(\alpha)|\le x\Nm(\mathfrak{b}_0)$ 的主理想 $(\alpha)$ 的个数. 下面只需计数. 所用技巧与单位定理证明中的类似; 细节见 \cite{Lang1970}, VI.3, 定理 3. (\cite{Janusz1996}, IV.2.11 中证明了一个稍弱的结果.)
\end{proof}

\begin{corollary}\label{cor:vi.2.9}
部分 zeta 函数 $\zeta(s,\mathfrak{k})$ 在 $\Re(s)>1-\frac{1}{d}$ 上解析, 除 $s=1$ 处一个留数为 $g_{\mathfrak{m}}$ 的单极点外.
\end{corollary}

\begin{proof}
应用命题 \ref{prop:vi.2.5}.
\end{proof}

注意 $g_{\mathfrak{m}}$ 不依赖于类 $\mathfrak{k}$.

\begin{lemma}\label{lem:vi.2.10}
设 $A$ 为有限 Abel 群, $\chi:A\to\CC^{\times}$ 为非平凡特征 (即不把每个元素都映到 $1$ 的同态). 则
\[
\sum_{a\in A}\chi(a)=0.
\]
\end{lemma}

\begin{proof}
由于 $\chi$ 非平凡, 存在 $b\in A$ 使得 $\chi(b)\ne1$. 但
\[
\sum_{a\in A}\chi(a)=\sum_{a\in A}\chi(ab)=\Bigl(\sum_{a}\chi(a)\Bigr)\chi(b),
\]
故
\[
\bigl(\chi(b)-1\bigr)\sum_{a}\chi(a)=0,
\]
从而 $\sum_{a}\chi(a)=0$.
\end{proof}

\begin{corollary}\label{cor:vi.2.11}
若 $\chi$ 不是平凡特征, 则 $L(s,\chi)$ 在 $\Re(s)>1-\frac{1}{d}$ 上解析.
\end{corollary}

\begin{proof}
在 $s=1$ 附近,
\[
L(s,\chi)=\sum_{\mathfrak{k}\in C_{\mathfrak{m}}}\chi(\mathfrak{k})\cdot\zeta(s,\mathfrak{k})=\frac{\sum_{\mathfrak{k}\in C_{\mathfrak{m}}}\chi(\mathfrak{k})\,g_{\mathfrak{m}}}{s-1}+\text{全纯函数},
\]
而引理表明第一项为零.
\end{proof}

以后我们将看到 $L(1,\chi)\ne0$.

\begin{corollary}\label{cor:vi.2.12}
Dedekind zeta 函数 $\zeta_K(s)$ 在 $\Re(s)>1-\frac{1}{d}$ 上解析, 除 $s=1$ 处一个单极点外, 该极点的留数为
\[
\frac{2^{r}(2\pi)^{s}\,\mathrm{reg}(K)}{w_{K}|\Delta|^{\frac{1}{2}}}\,h_{K}.
\]
\end{corollary}

\begin{proof}
回顾 $\zeta_K(s)=\sum_{\mathfrak{k}\in C_K}\zeta(s,\mathfrak{k})$.
\end{proof}

\begin{example}\label{ex:vi.2.13}
(a) 对 $K=\Q$, 上面的公式化为 $1=\frac{2}{2}$.

(b) 对 $K=\Q[\sqrt{d}]$, 公式化为
\[
\lim_{s\to1}(s-1)\zeta_{K}(s)=
\begin{cases}
\dfrac{2\log(u)}{\Delta^{\frac{1}{2}}}\,h_{K}, & u>1\ \text{为一个基本单位},\ d>0,\\[2ex]
\dfrac{2\pi}{w_{K}|\Delta|^{\frac{1}{2}}}\,h_{K}, & d<0.
\end{cases}
\]
可以求出左端表达式的闭公式, 由此得到类数的一个非常简单的表达式. 回顾: $K/\Q$ 的 Artin 映射可以视为一个特征 $\chi\colon I^S\to\{\pm1\}$, 其中 $S$ 是在 $K$ 中分歧的素的集合. 我们不把它看作理想上的映射, 而是看作正整数上的映射, 并规定: 若 $m$ 被某个在 $K$ 中分歧的素数整除, 则 $\chi(m)=0$, 以此把它延拓到所有正整数上. 于是 $\chi$ 现在是正整数集上的乘性映射, 取值
\[
\chi(p)=
\begin{cases}
1 & \text{若 $p$ 在 $K$ 中分裂},\\
-1 & \text{若 $p$ 在 $K$ 中保持素},\\
0 & \text{若 $p$ 在 $K$ 中分歧}.
\end{cases}
\]
对判别式 $<-4$ 的虚二次域, 公式化为
\[
h_{K}=\frac{1}{2-\chi(2)}\sum_{\substack{(x,\Delta)=1\\ 0<x<|\Delta|/2}}\chi(x).
\]
例如, 若 $K=\Q[\sqrt{-5}]$, 则 $|\Delta|=20$, 且
\[
h=\frac{1}{2-0}\bigl(\chi(1)+\chi(3)+\chi(7)+\chi(9)\bigr)=\frac{4}{2}=2,
\]
因为 $2$ 分歧, 而
\[
-5\equiv1\equiv1^{2}\bmod 3,\qquad -5\equiv2\equiv3^{2}\bmod 7.
\]
更多细节见 \cite{BorevichShafarevich1966}, 第 5 章, 第 4 节.
\end{example}

\section{在扩张中分裂的素理想的密度}\label{sec:VI.3}

对 $K$ 的素理想集 $T$, 定义 $\zeta_{K,T}(s)=\prod_{\mathfrak{p}\in T}\frac{1}{1-\Nm\mathfrak{p}^{-s}}$. 若 $\zeta_{K,T}(s)$ 的某个正整数次幂 $\zeta_{K,T}(s)^{n}$ 可延拓为 $1$ 的一个邻域上的亚纯函数, 且在 $1$ 处有 $m$ 阶极点, 则称\footnote{仿照 \cite{Marcus1977}, p.~188; 或第 2 版 p.~134.} $T$ 具有\textbf{极密度} (polar density) $\delta(T)=m/n$.

\begin{proposition}\label{prop:vi.3.1}
(a) $K$ 的全部素理想之集具有极密度 $1$.

(b) 任何 (具有极密度的) 集合的极密度都 $\ge0$.

(c) 设 $T$ 是 $T_1$ 与 $T_2$ 的不交并. 若 $T$, $T_1$, $T_2$ 中任意两个具有极密度, 则第三个也具有, 且 $\delta(T)=\delta(T_1)+\delta(T_2)$.

(d) 若 $T\subset T'$, 则 (当两者都有定义时) $\delta(T)\le\delta(T')$.

(e) 有限集的密度为零.
\end{proposition}

\begin{proof}
(a) 我们知道 $\zeta_K(s)$ 本身可延拓到 $1$ 的邻域, 且在 $1$ 处有单极点.

(b) 称 $T$ 具有负密度, 即是说 $\zeta_{K,T}(s)$ 在 $s=1$ 的邻域内全纯且在该点为零. 但 $\zeta_{K,T}(1)=\prod_{\mathfrak{p}\in T}\frac{1}{1-\mathfrak{p}^{-1}}>0$.

(c) 显然
\[
\zeta_{K,T}(s)=\zeta_{K,T_{1}}(s)\cdot\zeta_{K,T_{2}}(s).
\]
比方说, 设 $\zeta_{K,T}(s)^{m}$ 与 $\zeta_{K,T_{1}}(s)^{m_{1}}$ 可延拓为 $1$ 的邻域上的亚纯函数, 分别在 $1$ 处有 $n$ 阶与 $n_1$ 阶极点. 则
\[
\zeta_{K,T_{2}}(s)^{mm_{1}}=\zeta_{K,T}(s)^{mm_{1}}/\zeta_{K,T_{1}}(s)^{mm_{1}}
\]
可延拓为 $1$ 的邻域上的亚纯函数, 且在 $1$ 处有 $m_1n-mn_1$ 阶极点. 因此
\[
\delta(T_{2})=\frac{m_{1}n}{mm_{1}}-\frac{mn_{1}}{mm_{1}}=\delta(T)-\delta(T_{2}).
\]

(d) 结合 (c) 与 (b).

(e) 显然.
\end{proof}

\begin{proposition}\label{prop:vi.3.2}
若 $T$ 不含数值范数为素数 (属于 $\Z$) 的素理想, 则 $\delta(T)=0$.
\end{proposition}

\begin{proof}
对 $\mathfrak{p}\in T$, $\Nm\mathfrak{p}=p^{f}$, 其中 $f\ge2$. 此外, 对给定的 $p$, 位于其上的 $K$ 的素理想至多 $[K:\Q]$ 个. 因此 $\zeta_{K,T}(s)$ 可分解为 $\prod_{i=1}^{d}g_{i}(s)$, 即 $d$ 个关于素数的无穷乘积之积, 其中每个 $g_i(s)$ 的因子或为 $1$, 或为 $\frac{1}{1-p^{-fs}}$ ($f\ge2$) 的形状. 对每个 $i$, $g_{i}(1)\le\sum_{n>0}n^{-2}=\zeta(2)$. 因此 $g_i(s)$ 在 $1$ 处全纯.
\end{proof}

\begin{corollary}\label{cor:vi.3.3}
设 $T_1$ 与 $T_2$ 为 $K$ 中素理想的集合. 若二者仅相差数值范数不是素数的素理想, 且其中一者具有极密度, 则另一者亦然, 且两者的密度相等.
\end{corollary}

\begin{theorem}\label{thm:vi.3.4}
设 $L$ 为 $K$ 的有限扩张, $M$ 为其 Galois 闭包. 则在 $L$ 中完全分裂的 $K$ 的素理想之集具有密度 $1/[M:K]$.
\end{theorem}

\begin{proof}
$K$ 的素理想 $\mathfrak{p}$ 在 $L$ 中完全分裂当且仅当它在 $M$ 中完全分裂.\footnote{这里解释一下``素在扩张 $L$ 中完全分裂当且仅当在其 Galois 闭包中完全分裂''这一说法. 若素在 $L$ 中完全分裂, 则它在 $L$ 的每个共轭 $L'$ 中也完全分裂, 于是问题化为: 若素在两个域 $L$ 与 $L'$ 中都完全分裂, 则它在它们的复合中完全分裂. 这一点容易得到 (例如) 由下述判别法: $K$ 的素 $\mathfrak{p}$ 在 $L$ 中完全分裂当且仅当 $L\otimes_{K}K_{\mathfrak{p}}$ 是 $K_{\mathfrak{p}}$ 的若干拷贝之积 (复合 $LL'$ 是 $L\otimes_{K}L'$ 的直因子).} 因此, 只需在 $L$ 为 $K$ 的 Galois 扩张的假设下证明本定理. 设 $S$ 为在 $L$ 中完全分裂的 $K$ 的素理想之集, $T$ 为 $L$ 中位于 $S$ 的素理想之上的素理想之集. 对应于每个 $\mathfrak{p}\in S$, 恰有 $[L:K]$ 个素理想 $\mathfrak{P}\in T$, 且对其中每个都有 $\Nm_{L/K}\mathfrak{P}=\mathfrak{p}$, 从而 $\Nm\mathfrak{P}=\Nm\mathfrak{p}$. 于是 $\zeta_{L,T}(s)=\zeta_{K,S}(s)^{[L:K]}$. 但 $T$ 包含 $L$ 中每个在 $L/K$ 内不分歧且 $\Nm\mathfrak{P}=p$ 为有理素数的素理想. 因此 $T$ 与 $L$ 的全部素理想之集仅相差一个极密度为 $0$ 的集合, 从而 $T$ 具有密度 $1$. 这蕴含 $\zeta_{K,S}(s)$ 具有标志着 $S$ 的密度为 $1/[L:K]$ 的性质.
\end{proof}

\begin{corollary}\label{cor:vi.3.5}
若 $f(X)\in K[X]$ 对除有限多个之外的所有素理想 $\mathfrak{p}$ 都模 $\mathfrak{p}$ 分解为一次因式之积, 则 $f$ 在 $K$ 中分解为一次因式之积.
\end{corollary}

\begin{proof}
对 $f$ 的分裂域应用该定理.
\end{proof}

\begin{corollary}[Bauer 1916]\label{cor:vi.3.6}
对数域 $K$ 的 Galois 扩张 $L$ 与 $M$,
\[
L\subset M\iff\operatorname{Spl}(L)\supset\operatorname{Spl}(M).
\]
因此
\[
L=M\iff\operatorname{Spl}(L)=\operatorname{Spl}(M),
\]
且
\[
L\mapsto\operatorname{Spl}(L)
\]
是从 $K$ 的有限 Galois 扩张之集 (含于某个固定的代数闭包) 到 $\{\mathfrak{p}\subset\mathcal{O}_K\}$ 的子集之集的单射.
\end{corollary}

\begin{proof}
证明见定理 \ref{thm:v.3.25} 的证明.
\end{proof}

\begin{example}\label{ex:vi.3.7}
设 $f(X)$ 为三次不可约多项式. 使 $f(X)$ 模 $\mathfrak{p}$ 分裂的素 $\mathfrak{p}$ 之集的密度为 $1/3$ 或 $1/6$, 视 $f(X)$ 的 Galois 群为 $C_3$ 还是 $S_3$ 而定.
\end{example}

\begin{corollary}\label{cor:vi.3.8}
对每个 Abel 扩张 $L/K$ 与每个包含在 $L$ 中分歧的素的有限素集 $S\supset S_{\infty}$, Artin 映射 $\psi_{L/K}\colon I^{S}\to\Gal(L/K)$ 是满射.
\end{corollary}

\begin{proof}
设 $H$ 为 $\psi_{L/K}$ 的像. 对所有 $\mathfrak{p}\notin S$, $(\mathfrak{p},L^{H}/K)=(\mathfrak{p},L/K)|_{L^{H}}=1$, 这蕴含 $\mathfrak{p}$ 在 $L^{H}$ 中分裂. 于是除有限多个之外的全部素理想都在 $L^{H}$ 中分裂, 从而 $[L^{H}:K]=1$.
\end{proof}

\section{算术级数中素理想的密度}\label{sec:VI.4}

设 $f(s)$ 与 $g(s)$ 是两个 (至少) 对实数 $s>1$ 有定义的函数. 若 $f(s)-g(s)$ 在
\[
1<s<1+\varepsilon\qquad(\text{$s$ 为实数, 某个 }\varepsilon>0)
\]
上有界, 则记
\[
f(s)\sim g(s)\qquad\text{当 }s\downarrow 1.
\]
注意
\[
f(s)\sim\delta\log\frac{1}{s-1}\qquad\text{当 }s\downarrow 1
\]
蕴含
\[
\lim_{s\downarrow1}\frac{f(s)}{\log\frac{1}{s-1}}=\delta.
\]
当 $f(s)$ 与 $g(s)$ 都在 $s=1$ 的邻域内全纯 (可能除 $s=1$ 处有极点) 时,
\[
f(s)\sim g(s)\qquad\text{当 }s\downarrow 1
\]
当且仅当 $f(s)$ 与 $g(s)$ 相差一个在 $1$ 的邻域内全纯的函数.

设 $T$ 为 $K$ 的一个素集. 若存在 $\delta$ 使得
\[
\sum_{\mathfrak{p}\in T}\frac{1}{\Nm\mathfrak{p}^{s}}\sim\delta\log\frac{1}{s-1}\qquad\text{当 }s\downarrow 1,
\]
则称 $T$ 具有 \textbf{Dirichlet 密度} (Dirichlet density) $\delta$.

若极限
\[
\lim_{x\to\infty}\frac{\#\{\mathfrak{p}\in T\mid\Nm\mathfrak{p}\le x\}}{\#\{\mathfrak{p}\mid\Nm\mathfrak{p}\le x\}}
\]
存在, 则称之为 $T$ 的\textbf{自然密度} (natural density).

\begin{proposition}\label{prop:vi.4.1}
(a) 若极密度存在, 则 Dirichlet 密度也存在, 且两者相等.

(b) 若自然密度存在, 则 Dirichlet 密度也存在, 且两者相等.
\end{proposition}

\begin{proof}
(a) 若 $T$ 有极密度 $m/n$, 则
\[
\zeta_{K,T}(s)^{n}=\frac{a}{(s-1)^{m}}+\frac{g(s)}{(s-1)^{m-1}},
\]
其中 $g(s)$ 在 $s=1$ 附近全纯. 此外 $a>0$, 因为对实数 $s>1$ 有 $\zeta_{K,T}(s)>0$. 取对数 (并应用引理\ref{lem:vi.4.3}) 得
\[
n\sum_{T}\frac{1}{\Nm\mathfrak{p}^{s}}\sim m\log\frac{1}{s-1}\qquad\text{当 }s\downarrow 1,
\]
这表明 $T$ 有 Dirichlet 密度 $m/n$.

(b) 见 \cite{Goldstein1971}, p.~252.
\end{proof}

\begin{remark}\label{rem:vi.4.2}
(a) 一个集合 $T$ 可以有 Dirichlet 密度而没有自然密度. 例如, 设 $T$ 为 (十进制下) 首位数字是 $1$ 的素数之集. 则 $T$ 没有自然密度, 但其 Dirichlet 密度为 $\log_{10}(2)=0.3010300\ldots$ (结论见 \cite{Serre1970}, VI, 4.5). 因此, 断言一个素数集有自然密度 $\delta$ 比断言它有 Dirichlet 密度 $\delta$ 更强. 本讲义中计算密度的所有集合其实都具有自然密度, 但我们不对此证明.

(b) 由定义, 极密度是有理数. 因此, 每个具有非有理自然密度的集合都没有极密度.
\end{remark}

回顾: 指数函数
\[
e^{z}=\sum\frac{z^{n}}{n!}=e^{x}(\cos y+i\sin y),\quad z=x+iy
\]
定义了一个从
\[
\{z\in\CC\mid-\pi<\Im(z)<\pi\}
\]
到 $\CC$ 中负实轴的补集
\[
\{z\in\R\mid z\le0\}\text{ 在 $\CC$ 中的补集}
\]
的同构, 按定义其逆就是对数函数 $\log$ 的 (主支). 用这个定义,
\[
\log z=\log|z|+i\arg z,
\]
其中右端的 $\log$ 是微积分课程中定义的函数, 且 $-\pi<\arg z<\pi$. 用这个定义,
\[
\log\frac{1}{1-z}=z+\frac{z^{2}}{2}+\frac{z^{3}}{3}+\cdots,\qquad|z|<1.
\]

\begin{lemma}\label{lem:vi.4.3}
设 $u_1,u_2,\ldots$ 是一个不小于 $2$ 的实数序列, 使得
\[
f(s)\overset{\mathrm{def}}{=}\prod_{j=1}^{\infty}\frac{1}{1-u_{j}^{-s}}
\]
在每个区域 $D(1,\delta,\varepsilon)$ ($\delta,\varepsilon>0$) 上一致收敛. 则
\[
\log f(s)\sim\sum\frac{1}{u_{j}^{s}}\qquad\text{当 }s\downarrow 1.
\]
\end{lemma}

\begin{proof}
我们有
\begin{align*}
\log f(s)&=\sum_{j=1}^{\infty}\log\frac{1}{1-u_{j}^{-s}}\\
&=\sum_{j}\sum_{m=1}^{\infty}\frac{1}{mu_{j}^{sm}}\\
&=\sum_{j}\frac{1}{u_{j}^{s}}+\sum_{j}\sum_{m=2}^{\infty}\frac{1}{mu_{j}^{sm}}\\
&=\sum_{j}\frac{1}{u_{j}^{s}}+g(s),
\end{align*}
其中
\[
|g(s)|\le\sum_{j=1}^{\infty}\sum_{m=2}^{\infty}\left|\frac{1}{mu_{j}^{sm}}\right|=\sum_{j=1}^{\infty}\sum_{m=2}^{\infty}\frac{1}{mu_{j}^{m\sigma}},\qquad\sigma=\Re(s).
\]
利用 ($u\ge2$, $\sigma>1$)
\[
\sum_{m=2}^{\infty}\frac{1}{mu^{m\sigma}}\le\sum_{m=2}^{\infty}\frac{1}{2}\left(\frac{1}{u^{\sigma}}\right)^{m}=\frac{1}{2}\left(\frac{1}{1-u^{-\sigma}}-u^{-\sigma}-1\right)=\frac{1}{2}\cdot\frac{u^{-2\sigma}}{1-u^{-\sigma}}<\frac{1}{u^{2\sigma}}
\]
来估计内层和. 于是
\[
|g(s)|\le f(2\sigma).
\]
由于 $f(s)$ 对 $\Re(s)>1$ 全纯, $f(2s)$ 对 $\Re(s)>\frac{1}{2}$ 全纯, 故当 $\sigma\downarrow 1$ 时 $g(\sigma)$ 有界.
\end{proof}

\begin{proposition}\label{prop:vi.4.4}
(a) $K$ 的全部素理想之集的 Dirichlet 密度为 $1$.

(b) 任何 (具有 Dirichlet 密度的) 集合的 Dirichlet 密度都 $\ge0$.

(c) 若 $T$ 有限, 则 $\delta(T)=0$.

(d) 设 $T$ 是 $T_1$ 与 $T_2$ 的不交并. 若 $\delta(T_1)$, $\delta(T_2)$, $\delta(T)$ 中任意两个有定义, 则第三个也有定义, 且 $\delta(T)=\delta(T_1)+\delta(T_2)$.

(e) 若 $T\subset T'$, 则 (设两者都有定义) $\delta(T)\le\delta(T')$.
\end{proposition}

\begin{proof}
(a) 全部素理想之集甚至具有极密度 $1$.

(b) 对实数 $s>0$, $\frac{1}{\Nm\mathfrak{p}^{s}}>0$; 而对 $s=1+\varepsilon$, $\log\frac{1}{s-1}=-\log\varepsilon$ 在 $0<\varepsilon<1$ 时为正.

(c) 当 $T$ 有限时, $\sum_{\mathfrak{p}\in T}\frac{1}{\Nm\mathfrak{p}^{s}}$ 对所有 $s$ 全纯, 从而在任一点附近有界.

(d) 显然, 对 $\Re(s)>1$,
\[
\sum_{\mathfrak{p}\in T}\frac{1}{\Nm\mathfrak{p}^{s}}=\sum_{\mathfrak{p}\in T_{1}}\frac{1}{\Nm\mathfrak{p}^{s}}+\sum_{\mathfrak{p}\in T_{2}}\frac{1}{\Nm\mathfrak{p}^{s}}.
\]
因此, 比方说若
\[
\sum_{\mathfrak{p}\in T_{1}}\frac{1}{\Nm\mathfrak{p}^{s}}\sim\delta_{1}\log\frac{1}{s-1},\qquad\sum_{\mathfrak{p}\in T_{2}}\frac{1}{\Nm\mathfrak{p}^{s}}\sim\delta_{2}\log\frac{1}{s-1},
\]
则
\[
\sum_{\mathfrak{p}\in T}\frac{1}{\Nm\mathfrak{p}^{s}}\sim(\delta_{1}+\delta_{2})\log\frac{1}{s-1}.
\]

(e) 若 $\delta(T)$ 与 $\delta(T')$ 都存在, 则 $\delta(T'\setminus T)$ 也存在, 且
\[
\delta(T')-\delta(T)\overset{(c)}{=}\delta(T'\setminus T)\overset{(a)}{\ge}0.
\]
\end{proof}

\begin{proposition}\label{prop:vi.4.5}
设 $T$ 为 $K$ 中在 $\Q$ 上次数为 $1$ 的素理想之集, 即满足剩余类次数 $f(\mathfrak{p}/p)=1$ 的素理想. 则 $\delta(T)=1$.
\end{proposition}

\begin{proof}
$T$ 的补集具有极密度 $1$ (命题 \ref{prop:vi.3.2}).
\end{proof}

\begin{corollary}\label{cor:vi.4.6}
设 $T$ 如命题 \ref{prop:vi.4.5} 所设. 对每个具有 Dirichlet 密度的 $K$ 的素集 $S$,
\[
\delta(T\cap S)=\delta(S).
\]
\end{corollary}

\begin{proof}
$T$ 的补集 $T'$ 的密度为 $0$, 由此易得 $\delta(S\cap T')=0$. 由于 $S$ 是 $S\cap T$ 与 $S\cap T'$ 的不交并, 这蕴含 $\delta(S\cap T)$ 有定义且等于 $\delta(S)$.
\end{proof}

\begin{lemma}\label{lem:vi.4.7}
设 $A$ 为有限 Abel 群, $a\in A$. 则
\[
\sum_{\chi\in A^{\vee}}\chi(a)=0
\]
这里 $A^{\vee}$ 是 $A$ 的特征群, 即 $A^{\vee}=\Hom(A,\CC^{\times})$.
\end{lemma}

\begin{proof}
若 $a=1$, 则对所有 $\chi$ 有 $\chi(a)=1$, 故由 $A^{\vee}$ 与 $A$ 有相同个数的元素 (它事实上 (非典范地) 同构于 $A$) 即得结论. 若 $a\ne1$, 则存在特征 $\chi_1$ 使得 $\chi_1(a)\ne1$. 于是
\[
\sum_{\chi\in A^{\vee}}\chi(a)=\sum_{\chi\in A^{\vee}}(\chi_{1}\chi)(a)=\sum_{\chi\in A^{\vee}}\chi_{1}(a)\chi(a)=\chi_{1}(a)\sum_{\chi\in A^{\vee}}\chi(a).
\]
由于 $\chi_1(a)\ne1$, 这蕴含 $\sum_{\chi\in A^{\vee}}\chi(a)=0$.

另一种做法: 借助同构
\[
a\mapsto(\chi\mapsto\chi(a))\colon A\to(A^{\vee})^{\vee}
\]
把 $A$ 与 $A^{\vee\vee}$ 等同, 然后应用引理\ref{lem:vi.2.10}.
\end{proof}

\begin{theorem}\label{thm:vi.4.8}
设 $\mathfrak{m}$ 为 $K$ 的模, $H$ 为 $\mathfrak{m}$ 的一个同余子群:
\[
I^{\mathfrak{m}}\supset H\supset i(K_{\mathfrak{m},1}).
\]
则
\[
\delta(\{\mathfrak{p}\in H\})=
\begin{cases}
1/(I^{S(\mathfrak{m})}:H) & \text{若 $L(1,\chi)$ 对 $I^{S(\mathfrak{m})}/H$ 的所有特征 $\chi\ne\chi_0$ 都非零};\\
0 & \text{否则}.
\end{cases}
\]
\end{theorem}

\begin{proof}
设 $\chi$ 是 $I^{\mathfrak{m}}$ 的在 $H$ 上平凡的特征, 并令
\[
L(s,\chi)=\prod_{\mathfrak{p}\nmid\mathfrak{m}}\frac{1}{1-\chi(\mathfrak{p})\Nm\mathfrak{p}^{-s}}.
\]
则引理\ref{lem:vi.4.3} 证明中的论证表明
\[
\log L(s,\chi)-\sum_{\mathfrak{p}\nmid\mathfrak{m}}\frac{\chi(\mathfrak{p})}{\Nm\mathfrak{p}^{s}}
\]
在 $\Re(s)>\frac{1}{2}$ 上全纯. 特别地,
\[
\log(L(s,\chi))\sim\sum_{\mathfrak{p}\nmid\mathfrak{m}}\frac{\chi(\mathfrak{p})}{\Nm\mathfrak{p}^{s}}\qquad\text{当 }s\downarrow 1.
\]
但 (见引理\ref{lem:vi.4.7})
\[
\sum_{\chi}\chi(\mathfrak{p})=
\begin{cases}
h & \mathfrak{p}\in H,\\
0 & \mathfrak{p}\notin H,
\end{cases}
\]
故对全体 $\chi$ 求和得
\[
\sum_{\chi}\log L(s,\chi)\sim h\sum_{\mathfrak{p}\in H}\frac{1}{\Nm\mathfrak{p}^{s}}\qquad\text{当 }s\downarrow 1.
\]
若 $\chi\ne\chi_0$, 则 $L(s,\chi)$ 在 $s=1$ 附近全纯, 设 $L(s,\chi)=(s-1)^{m(\chi)}g(s)$, 其中 $m(\chi)\ge0$, $g(1)\ne0$. 于是
\[
\log L(s,\chi)\sim m(\chi)\log(s-1)=-m(\chi)\log\frac{1}{s-1}.
\]
若 $\chi=\chi_0$, 则 $L(s,\chi)=\zeta_K(s)/\prod_{\mathfrak{p}\mid\mathfrak{m}}\frac{1}{1-\Nm\mathfrak{p}^{-s}}$, 故
\[
\log L(s,\chi_{0})\sim\log\zeta_{K}(s)\sim\log\frac{1}{s-1}\qquad\text{当 }s\downarrow 1.
\]
合并这些结果, 得
\[
h\sum_{\mathfrak{p}\in H}\Nm\mathfrak{p}^{-s}\sim\Bigl(1-\sum_{\chi\ne\chi_{0}}m(\chi)\Bigr)\log\frac{1}{s-1},
\]
从而
\[
\delta(\{\mathfrak{p}\in H\})=\frac{1-\sum_{\chi\ne\chi_{0}}m(\chi)}{h}.
\]
这表明: 若 $L(1,\chi)\ne0$ 对所有 $\chi\ne\chi_0$ 成立, 则 $\delta(\{\mathfrak{p}\in H\})=\frac{1}{h}$; 否则 $\delta(\{\mathfrak{p}\in H\})=0$ (且至多有一个 $L(s,\chi)$ 在 $s=1$ 处有零点, 并且只能是单零点).
\end{proof}

\subsection*{第二个不等式}

\begin{theorem}\label{thm:vi.4.9}
对 $K$ 的每个 Galois 扩张 $L$ 与每个模 $\mathfrak{m}$,
\[
\bigl(I^{S(\mathfrak{m})}:i(K_{\mathfrak{m},1})\cdot\Nm(I_{L}^{S(\mathfrak{m})})\bigr)\le[L:K].
\]
\end{theorem}

\begin{proof}
令 $H=\Nm_{L/K}I_{L}^{\mathfrak{m}}\cdot i(K_{\mathfrak{m},1})$. 由定理 \ref{thm:vi.4.8} 知, $\delta(\{\mathfrak{p}\in H\})=1/(I^{S(\mathfrak{m})}:H)$ 或 $0$, 且前一种情形恰好在对 $I^{S(\mathfrak{m})}/H$ 的所有非平凡特征 $\chi$ 都有 $L(1,\chi)\ne0$ 时发生.

若 $\mathfrak{p}$ 在 $L$ 中分解, 即对所有 $\mathfrak{P}|\mathfrak{p}$ 有 $f(\mathfrak{P}/\mathfrak{p})=1$, 则 $\mathfrak{p}$ 是 $\mathcal{O}_L$ 中位于其上的任一素理想的范, 从而 $\{\mathfrak{p}\in H\}$ 包含在 $L$ 中分裂的素理想之集. 于是定理 \ref{thm:vi.3.4} 表明
\[
\delta(\{\mathfrak{p}\in H\})\ge[L:K]^{-1}\ne0.
\]
合并两方面结论, 我们得到:
(a) $\delta(\{\mathfrak{p}\in H\})\ne0$;
(b) 对 $I^{S(\mathfrak{m})}/H$ 的所有非平凡特征 $\chi$, $L(1,\chi)\ne0$;
(c) $(I^{S(\mathfrak{m})}:H)=\delta(\{\mathfrak{p}\in H\})^{-1}\le[L:K]$.
\end{proof}

\begin{corollary}\label{cor:vi.4.10}
设 $\chi$ 为 $C_{\mathfrak{m}}$ 的非平凡特征, 并设存在 $K$ 的 Galois 扩张 $L$ 使得 $\Nm_{L/K}C_{\mathfrak{m},L}\subset\Ker(\chi)$. 则 $L(1,\chi)\ne0$.
\end{corollary}

\begin{proof}
这在定理的证明过程中已经得到.
\end{proof}

存在定理 (第 \ref{ch:V} 章, 定理 \ref{thm:v.3.6}) 蕴含: 该推论的条件对所有 $\chi$ 都满足. 不用类域论也能证明 $L(1,\chi)\ne0$, 但此处我们更愿意回到类域论. Chebotarev 定理的证明将在第 \ref{ch:VIII} 章完成.
   % VI  L-级数与素数密度
% ============================================================
% 第七章  整体类域论: 主要定理的证明
% 源: Milne CFT v4.03, Chapter VII (书页 p.199--226)
% OCR: 行 7441--8470; 脚注 1--7 从原书 PDF 恢复
% 勘误: CFT 勘误表无本章条目; 原书无编号公式, 但 (*) 标记照排;
%       4.3 证明缺失公式 h(C_L)=h(I_{L,T})/h(U(T)) 及两张图
%       (引理 6.1 域图, 引理 6.5 蛇形图) 均从原书恢复.
%       原书 "H^2(-/K)" 的连字符在 pdftotext 中丢失, 记号已还原.
% 标签前缀: vii.
% ============================================================
\chapter{整体类域论: 主要定理的证明}\label{ch:VII}
\index{zhengtileiyulun@整体类域论(global class field theory)}

\begin{quote}
\itshape J'ai revu un peu la théorie du corps de classes, dont j'ai enfin l'impression d'avoir compris les énoncés essentiels (bien entendu, pas les démonstrations!)
\end{quote}
\hspace*{\fill}Grothendieck, 致 Serre 的信, 1956.9.19\footnote{``我近来稍稍重温了类域论, 终于觉得自己理解了它的几个主要结论 (当然, 不包括证明!)'' (原书脚注给出的英译, 此处译出.)}

本章中我们遵循 \cite{Tate1967} 的方法 (另见 \cite{ArtinTate1961}) 证明整体类域论的主要定理, 即互反律与存在定理 (第 \ref{ch:V} 章的定理 \ref{thm:v.5.3} 与 \ref{thm:v.5.5}). 全章均使用 idèle 而非理想.

本章独立于第 \ref{ch:VI} 章, 只是用第 \ref{ch:VI} 章的定理 \ref{thm:vi.4.9} 可以替代第 \ref{sec:VII.6} 节. 我们需要引用第 \ref{ch:V} 章的地方仅有: idèle 类群的定义, 以及整体 Artin 映射 $\phi\colon\mathbb{I}\to\Gal(L/K)$ 作为局部 Artin 映射之``积''的定义 (第 \ref{sec:VII.5} 节). 另一方面, 我们将频繁使用第 \ref{ch:II}、\ref{ch:III} 两章的结果.

我们使用第 \ref{ch:V} 章的记号 (特别见本书 p.\pageref{pg:orig169}). 特别地, $|\cdot|_v$ 表示 $v$ 的规范绝对值 (使乘积公式成立的那个).

\section{概述}\label{sec:VII.1}

设 $L/K$ 是数域的有限 Galois 扩张, Galois 群为 $G$. Idèle 类群 $\mathbf{C}_L\overset{\text{def}}{=}\mathbb{I}_L/L^{\times}$ 在整体类域论中所扮演的角色, 恰如乘法群 $L^{\times}$ 在局部类域论中的角色. 特别地, 当 $L/K$ 为 Abel 扩张时, 我们将证明存在同构
\[
\phi\colon\mathbf{C}_{K}/\Nm_{L/K}(\mathbf{C}_{L})\to G,
\]
其局部分量即局部 Artin 映射, 也就是说: 对 $K$ 的任意素 $v$ 与其上的 $L$ 的素 $w$, 下图交换:
\[
\begin{tikzcd}
K_{v}^{\times}\ar[r, "\phi_{v}"]\ar[d] & \Gal(L_{w}/K_{v})\ar[d]\\
\mathbb{I}_{K}\ar[r, "\phi"'] & \Gal(L/K),
\end{tikzcd}
\]
其中 $\phi_v$ 是第 \ref{ch:I}、\ref{ch:III} 两章的局部 Artin 映射.

根据 Tate 定理 (第 \ref{ch:II} 章, 3.11), 为得到这样的同构, 只需证明: 对每个具有 Galois 群 $G$ 的有限 Galois 扩张 $L/K$,

(a) $H^{1}(G,\mathbf{C}_{L})=0$;

(b) $H^{2}(G,\mathbf{C}_{L})$ 是 $[L:K]$ 阶循环群, 带有一个典范生成元 $u_{L/K}$;

(c) 若 $E\supset L\supset K$ 是 $K$ 的两个有限 Galois 扩张, 则 $\mathrm{Res}(u_{E/K})=u_{E/L}$.

此时同构 $\phi_{L/K}$ 便是 $u_{L/K}$ 所定义之映射的逆:
\[
H_{T}^{-2}(G,\mathbb{Z})\to H_{T}^{0}(G,\mathbf{C}_{L}).
\]
一旦证明了基本类 $u_{L/K}$ 与局部基本类相容, $\phi_{L/K}$ 就会是局部 Artin 映射的乘积.

实际上, 我们采用一条略为不同的路线. 我们将把整体 Artin 映射 $\mathbf{C}_K\to\Gal(L/K)$ 定义为局部 Artin 映射的``积'', 并用比 (a)(b) 稍弱的结果推出它具有应有的性质.

在第 \ref{sec:VII.2} 节中, 我们用局部域的上同调表示 idèle 的上同调:
\[
\begin{aligned}
&H^{0}(G,\mathbb{I}_{L})=\mathbb{I}_{K};\\
&H^{r}_{T}(G,\mathbb{I}_{L})\simeq\bigoplus_{v}H^{r}_{T}(G^{v},L^{v\times})
\end{aligned}
\]
(对 $K$ 的素 $v$ 求和; 取定 $w|v$, 则 $G^v$ 是 $w$ 的分解群, $L^v=L_w$). 在第 \ref{sec:VII.3} 节算出 $S$-单位群的 Herbrand 商之后, 我们在第 \ref{sec:VII.4} 节证明第一个不等式:
\[
\text{对任意循环扩张 }L/K,\quad(\mathbf{C}_K:\Nm_{L/K}\mathbf{C}_L)\ge[L:K].
\]
在第 \ref{sec:VII.4} 节我们还将证明: 对任何 Abel 扩张, Galois 群由 Frobenius 元生成. 第 \ref{sec:VII.5} 节陈述定理:
\[
\begin{gathered}
\text{对数域的每个 Galois 扩张 }L/K:\ \text{(a)}\ (\mathbf{C}_K:\Nm_{L/K}\mathbf{C}_L)\le[L:K]\ \text{(第二个不等式)};\\
\text{(b)}\ H^{1}(G,\mathbf{C}_{L})=1;\quad \text{(c)}\ H^{2}(G,\mathbf{C}_{L})\ \text{的阶}\le[L:K],
\end{gathered}
\]
并利用第 \ref{ch:VI} 章的定理 \ref{thm:vi.4.9} 证明之. 下一节给出该定理的另一个证明, 避免使用复分析.

在作一些关于 Brauer 群的预备之后, 第 \ref{sec:VII.7} 节通过证明: 对每个 Abel 扩张 $L/K$, $K^{\times}$ 包含在 $\phi_{L/K}:\mathbb{I}_K\to\Gal(L/K)$ 的核中, 来完成互反律的证明. 由于我们已经知道 $\Nm_{L/K}(\mathbb{I}_L)$ 包含在 $\phi_{L/K}$ 的核中, 且 $\phi_{L/K}$ 是满射 (因为 $\Gal(L/K)$ 由 Frobenius 元生成), 第二个不等式便蕴含 $\phi_{L/K}$ 是同构.

我们在第 \ref{sec:VII.9} 节证明存在定理, 证明途径是: $C_K$ 的每个 (开) 有限指数子群都包含如下扩张的某个子扩张的范群, 该扩张由先添加一个单位根、再作 Kummer 扩张而得到.

在某种程度上, $K$ 的循环分圆扩张扮演着与局部域的无分歧扩张相同的角色. 例如, 互反律证明的最后一步的关键在于: $\mathrm{Br}(K)$ 的每个元素都被一个循环分圆扩张分裂.

\section{Idèle 的上同调}\label{sec:VII.2}

设 $L/K$ 是数域的有限 Galois 扩张, Galois 群为 $G$. 回顾: $\sigma\in G$ 按规则 $|\sigma a|_{\sigma w}=|a|_w$ 作用在 $L$ 的位于 $K$ 的固定素 $v$ 之上的素 $w$ 上. 因此 $\sigma$ 是赋值域的同构
\[
(L,|\mid_{w})\to(L,|\mid_{\sigma w}),
\]
从而可延拓到完备化上: 有交换图
\[
\begin{tikzcd}
L_{w}\ar[r, "\sigma"] & L_{\sigma w}\\
L\ar[u, "i_{w}"']\ar[r, "\sigma"'] & L\ar[u, "i_{\sigma w}"']
\end{tikzcd}
\]
取定 $K$ 的素 $v$, 设 $w_0$ 是 $L$ 的位于 $v$ 之上的素. 映射 $\sigma\mapsto\sigma w_0$ 定义双射
\[
G/G_{w_{0}}\to\{w\mid v\},
\]
其中 $G_{w_0}$ 是 $w_0$ 的分解群.

我们希望把 $G$ 在 $L$ 上的作用延拓为 $G$ 在 $\prod_{w|v}L_w$ 上的作用. 回顾 (ANT, 8.2): 映射
\[
a\otimes b\mapsto(i_{w}(a)b)_{w}\colon L\otimes_{K}K_{v}\to\prod_{w|v}L_{w}
\]
是同构. 群 $G$ 通过其在 $L$ 上的作用作用于 $L\otimes_K K_v$, 我们用这个同构把该作用转移到 $\prod_{w|v}L_w$ 上. 于是:

(a) $G$ 的元素连续地作用于 $\prod_{w|v}L_w$;

(b) 所有形如 $(a,\ldots,a)$ ($a\in K_v$) 的元素都被 $G$ 固定;

(c) 对每个 $a\in L$, $\sigma(\ldots,i_w(a),\ldots)=(\ldots,i_w(\sigma a),\ldots)$.

这些条件唯一确定该作用 (实际上, (a) 与 (c) 已确定它, 因为 $K$ 在 $K_v$ 中稠密).

\begin{lemma}\label{lem:vii.2.1}
对 $\sigma\in G$ 与 $\alpha=(\alpha(w))\in\prod_{w\mid v}L_w$,
\begin{equation*}
(\sigma\alpha)(w)=\sigma\bigl(\alpha(\sigma^{-1}w)\bigr).
\tag{$*$}
\end{equation*}
\end{lemma}

\begin{proof}
规则 $(*)$ 确实定义了 $G$ 在 $\prod L_w$ 上的连续作用, 故只需验证它满足 (b) 与 (c). 条件 (b) 显然. 对 (c), 设 $\alpha(w)=i_w(a)$, $a\in L$. 则 (由 $(*)$)
\[
(\sigma\alpha)(w)=\sigma\bigl(i_{\sigma^{-1}w}(a)\bigr).
\]
在上面的交换图中把 $w$ 换成 $\sigma^{-1}w$, 便得公式 $\sigma\circ i_{\sigma^{-1}w}=i_{w}\circ\sigma$. 因此
\[
(\sigma\alpha)(w)=i_{w}(\sigma a),
\]
这正是所要证的.
\end{proof}

用更朴素的话说: $(\sigma\alpha)(\sigma w)=\sigma(\alpha(w))$, 即若 $\alpha$ 在 $w$-分量处是 $a$, 则 $\sigma\alpha$ 在 $\sigma w$-分量处是 $\sigma a$.

注意 $G$ 在 $\prod_{w|v}L_w$ 上的作用诱导出 $G$ 在 $\prod_{w|v}L_w$ 的子集 $\prod_{w|v}L_w^{\times}$ 与 $\prod_{w|v}U_w$ 上的作用.

\begin{proposition}\label{prop:vii.2.2}
取定 $w_0|v$, 设 $G_{w_0}$ 为其分解群. 对 $\alpha\in\prod_{w|v}L_w$ 与 $\sigma\in G$, 定义 $f_{\alpha}(\sigma)=\sigma(\alpha(\sigma^{-1}w_{0}))$. 则 $f_{\alpha}\in\operatorname{Ind}_{G_{w_{0}}}^{G}(L_{w_{0}})$, 且映射
\[
\alpha\mapsto f_{\alpha}\colon\prod_{w\mid v}L_{w}\to\operatorname{Ind}_{G_{w_{0}}}^{G}(L_{w_{0}})
\]
是 $G$-模的同构. 把 $L_w$ 换成 $L_w^{\times}$ 或 $U_w$ 后同样的结论成立.
\end{proposition}

\begin{proof}
回顾 (II, \S \ref{sec:II.1}):
\[
\operatorname{Ind}_{G_{w_{0}}}^{G}(L_{w_{0}})=\{f:G\to L_{w_{0}}\mid f(\rho\sigma)=\rho f(\sigma),\ \text{一切 }\rho\in G_{w_{0}}\},
\]
且 $\tau\in G$ 按 $(\tau f)(\sigma)=f(\sigma\tau)$ 作用于 $f\in\operatorname{Ind}_{G_{w_{0}}}^{G}(L_{w_{0}})$. 对 $\rho\in G_{w_{0}}$,
\[
f_{\alpha}(\rho\sigma)=\rho\sigma\bigl(\alpha(\sigma^{-1}\rho^{-1}w_{0})\bigr)=\rho\sigma\bigl(\alpha(\sigma^{-1}w_{0})\bigr)=\rho f_{\alpha}(\sigma),
\]
故 $f_{\alpha}\in\operatorname{Ind}_{G_{w_{0}}}^{G}(L_{w_{0}})$. 此外,
\[
(\tau f_{\alpha})(\sigma)=f_{\alpha}(\sigma\tau)=\sigma\tau\bigl(\alpha(\tau^{-1}\sigma^{-1}w_{0})\bigr)=\sigma(\tau\alpha)(\sigma^{-1}w_{0})=f_{\tau\alpha}(\sigma),
\]
故 $\alpha\mapsto f_\alpha$ 是 $G$-模同态 $\prod_{w|v}L_w\to\operatorname{Ind}_{G_{w_{0}}}^{G}(L_{w_{0}})$. 给定 $f\in\operatorname{Ind}_{G_{w_{0}}}^{G}(L_{w_{0}})$, 令
\[
\alpha_{f}(w)=\sigma\bigl(f(\sigma^{-1})\bigr),\quad w=\sigma w_{0}.
\]
则 $f\mapsto\alpha_f$ 是 $\alpha\mapsto f_\alpha$ 的逆.
\end{proof}

\begin{proposition}\label{prop:vii.2.3}
对一切 $r$,
\[
H^{r}\Bigl(G,\prod_{w\mid v}L_{w}^{\times}\Bigr)\simeq H^{r}(G_{w_{0}},L_{w_{0}}^{\times}).
\]
特别地,
\[
H^{0}\Bigl(G,\prod_{w|v}L_{w}^{\times}\Bigr)\simeq K_{v}^{\times}.
\]
把 $L_w^{\times}$ 换成 $U_w$ 后同样的结论成立.
\end{proposition}

\begin{proof}
由 Shapiro 引理 (II, \ref{prop:ii.1.11}),
\[
H^{r}\Bigl(G,\prod_{w|v}L_{w}^{\times}\Bigr)=H^{r}\bigl(G,\operatorname{Ind}_{G_{w_{0}}}^{G}L_{w_{0}}^{\times}\bigr)=H^{r}(G_{w_{0}},L_{w_{0}}^{\times}).
\]
\end{proof}

\begin{remark}\label{rem:vii.2.4}
群 $H^r(G_{w_0},L_{w_0}^{\times})$ 在典范同构的意义下不依赖于整除 $v$ 的素 $w_0$: 设 $w$ 是另一个这样的素, 则可写 $w=\sigma w_0$, 且我们有相容的同构对
\[
\tau\mapsto\sigma\tau\sigma^{-1}\colon G_{w_{0}}\to G_{w},\qquad x\mapsto\sigma^{-1}x\colon L_{w}\to L_{w_{0}},
\]
从而对每个 $r$ 有同构
\[
H^{r}(G_{w},L_{w}^{\times})\to H^{r}(G_{w_{0}},L_{w_{0}}^{\times})
\]
(见 II, \S \ref{sec:II.1}). 若 $w=\sigma'w_0$, 则 $\sigma'=\sigma\tau$, $\tau\in G_{w_0}$. 由 $\sigma$ 与 $\sigma'$ 定义的两个映射相差 $H^r(G_{w_0},L_{w_0}^{\times})$ 的由 $\tau$ 给出的自同构, 而它是恒等映射 (II, \ref{ex:ii.1.27}d). 因此 $H^r(G_{w_0},L_{w_0}^{\times})$ 与 $H^r(G_w,L_w^{\times})$ 典范同构. 这提示如下记号: 取定 $w|v$, 令
\[
G^{v}=G_{w},\quad L^{v}=L_{w},\quad U^{v}=U_{w}.
\]
这些对象只在相差非典范同构的意义下有定义, 但 $H^r(G^v,L^{v\times})$ 与 $H^r(G^v,U^v)$ 在典范同构的意义下有定义.
\end{remark}

我们赋予 $\mathbb{I}_L$ 如下的 $G$-作用, 使包含
\[
\prod_{w\mid v}L_{w}^{\times}\to\mathbb{I}_{L}
\]
都是 $G$-同态. 于是若 $\alpha$ 在 $w$-分量处是 $a_w$, 则 $\sigma\alpha$ 在 $\sigma w$-分量处是 $\sigma a_w$.

\begin{proposition}\label{prop:vii.2.5}
(a) 映射 $\mathbb{I}_K\hookrightarrow\mathbb{I}_L$ 诱导同构 $\mathbb{I}_K\to\mathbb{I}_L^{G}$.

(b) 对所有 $r\ge0$,
\[
H_{T}^{r}(G,\mathbb{I}_{L})=\bigoplus_{v}H_{T}^{r}(G^{v},L^{v\times}).
\]
\end{proposition}

\begin{proof}
(a) 显然, $\alpha=(a_w)$ 被 $G$ 固定当且仅当每个子族 $(a_w)_{w|v}$ 被 $G$ 固定. 而 $(a_w)_{w|v}$ 被 $G$ 固定仅当 $a_w$ 不依赖于 $w$ 且落在 $K_v^{\times}$ 中.

(b) 对 $K$ 的素的每个有限集 $S$, 令
\[
\mathbb{I}_{L,S}=\prod_{v\in S}\Bigl(\prod_{w|v}L_{w}^{\times}\Bigr)\times\prod_{v\notin S}\Bigl(\prod_{w|v}U_{w}\Bigr).
\]
则 $\mathbb{I}_{L,S}$ 在 $G$ 的作用下稳定, 且 $\mathbb{I}_L$ 是 $\mathbb{I}_{L,S}$ 的归纳极限, 其中 $S$ 跑遍 $K$ 的包含全部无限素与全部在 $L$ 中分歧的素的有限素集. 于是 (见 II, \ref{prop:ii.4.4}),
\[
H^{r}(G,\mathbb{I}_{L})=\varinjlim H^{r}(G,\mathbb{I}_{L,S}).
\]
应用 (II, \ref{prop:ii.1.25}) 与 (\ref{prop:vii.2.3}) 得
\begin{align*}
H^{r}(G,\mathbb{I}_{L,S})&=\prod_{v\in S}H^{r}\Bigl(G,\prod_{w|v}L_{w}^{\times}\Bigr)\times\prod_{v\notin S}H^{r}\Bigl(G,\prod_{w|v}U_{w}\Bigr)\\
&=\prod_{v\in S}H^{r}(G^{v},L^{v\times})\times\prod_{v\notin S}H^{r}(G^{v},U^{v}).
\end{align*}
由 (III, \ref{prop:iii.1.1}), 当 $r>0$ 时第二个因子为零, 故
\begin{align*}
H^{r}(G,\mathbb{I}_{L})&=\varinjlim_{S}H^{r}(G,\mathbb{I}_{L,S})\\
&=\varinjlim_{S}\bigoplus_{v\in S}H^{r}(G^{v},L^{v\times})\\
&=\bigoplus_{\text{一切 }v}H^{r}(G^{v},L^{v\times}).
\end{align*}
对 $r\le0$ 使用 Tate 群后同样的论证成立.
\end{proof}

\begin{corollary}\label{cor:vii.2.6}
(a) $H^{1}(G,\mathbb{I}_{L})=0$;

(b) $H^{2}(G,\mathbb{I}_{L})\simeq\bigoplus_{v}\bigl(\tfrac{1}{n_{v}}\mathbb{Z}/\mathbb{Z}\bigr)$, 其中 $n_v=[L^{v}:K_v]$.
\end{corollary}

\begin{proof}
(a) 应用 Hilbert 定理 90 (II, \ref{prop:ii.1.22}).

(b) 由第 \ref{ch:III} 章的定理\ref{thm:iii.2.1}, 不变映射给出同构
\[
H^{2}(G^{v},L^{v\times})\to\frac{1}{n_{v}}\mathbb{Z}/\mathbb{Z}.
\]
\end{proof}

\begin{proposition}\label{prop:vii.2.7}
设 $S$ 为 $K$ 的素的有限集, $T$ 为 $L$ 中位于 $S$ 的素之上的素集. 若 $L/K$ 循环, 则 Herbrand 商
\[
h(\mathbb{I}_{L,T})=\prod_{v\in S}n_{v},\quad n_{v}=[L^{v}:K_{v}].
\]
\end{proposition}

\begin{proof}
我们有
\[
\mathbb{I}_{L,T}=\Biggl(\prod_{v\in S}\Bigl(\prod_{w|v}L_{w}^{\times}\Bigr)\Biggr)\times\Biggl(\prod_{v\notin S}\Bigl(\prod_{w|v}U_{w}\Bigr)\Biggr).
\]
第二个因子的 Herbrand 商为 $1$ (应用 III, \ref{lem:iii.2.5}), 故
\begin{align*}
h(G,\mathbb{I}_{L,T})&=\prod_{v\in S}h\Bigl(G,\prod_{w|v}L_{w}^{\times}\Bigr)\\
&=\prod_{v\in S}h(G^{v},L^{v\times})\\
&=\prod_{v\in S}\bigl|H^{2}(G^{v},L^{v\times})\bigr|\\
&=\prod_{v\in S}n_{v}.
\end{align*}
\end{proof}

\subsection*{Idèle 上的范映射}

设 $L/K$ 是数域的有限 Galois 扩张. 与对任何 $G$-模一样, 存在范映射
\[
x\mapsto\prod_{\sigma\in G}\sigma x\colon\mathbb{I}_{L}\to\mathbb{I}_{L}^{G}=\mathbb{I}_{K}.
\]
我们需要考察这个映射. 回顾 (ANT, 8.4): 有交换图
\[
\begin{tikzcd}
L^{\times}\ar[r]\ar[d, "\Nm_{L/K}"'] & \prod_{w\mid v}L_{w}^{\times}\ar[d, "(a_{w})\mapsto\prod\Nm_{Lw/K_{v}}a_{w}"]\\
K^{\times}\ar[r] & K_{v}^{\times}.
\end{tikzcd}
\]
对每个 $w$, $\Nm L_w^{\times}$ 在 $K_v^{\times}$ 中开 (因为它指数有限; 见 I, \ref{lem:1.3}), 且对任何不分歧的 $w$, 范映射把 $U_w$ 映到 $U_v$ 上 (见 III, \ref{prop:iii.1.2}). 图中右侧竖直映射的像等于 $\Nm L_w^{\times}$ (对任意 $w|v$; 因为任意两个 $L_w$ 都 $K_v$-同构). 我们把它记作 $\Nm L^{v\times}$.

设 $S\supset S_{\infty}$ 为 $K$ 的素的有限集, 包含全部分歧素, 并设 $T$ 为位于 $S$ 的素之上的素集. 上面的论述表明
\[
\Nm_{L/K}\mathbb{I}_{L,T}=\prod_{v\in S}V_{v}\times\prod_{v\notin S}U_{v},
\]
其中 $V_v$ 是 $K_v^{\times}$ 的指数有限的开子群. 它是 $\mathbb{I}_{K,S}$ 中的开子群, 而 $\mathbb{I}_{K,S}$ 在 $\mathbb{I}_K$ 中开. 于是我们证明了下述结果.

\begin{proposition}\label{prop:vii.2.8}
对数域的每个有限 Galois 扩张 $L/K$, $\Nm_{L/K}\mathbb{I}_L$ 包含 $\mathbb{I}_K$ 的一个开子群, 从而它本身是开的.
\end{proposition}

事实上, 我们之前 (V, \ref{cor:v.4.13}) 已经证明: 每个范群都对某个模 $\mathfrak{m}$ 包含 $W_{\mathfrak{m}}$, 而 $W_{\mathfrak{m}}$ 是开的.

接下来考虑 idèle 类上的范映射. 考虑
\[
\begin{tikzcd}
0\ar[r] & L^{\times}\ar[r]\ar[d, "\Nm_{L/K}"'] & \mathbb{I}_{L}\ar[r]\ar[d, "\Nm_{L/K}"'] & \mathbf{C}_{L}\ar[r]\ar[d] & 0\\
0\ar[r] & K^{\times}\ar[r] & \mathbb{I}_{K}\ar[r] & \mathbf{C}_{K}\ar[r] & 0.
\end{tikzcd}
\]
左侧方块交换, 故范映射 $\mathbb{I}_L\to\mathbb{I}_K$ 诱导范映射 $\Nm_{L/K}\colon\mathbf{C}_L\to\mathbf{C}_K$. 由蛇形引理, 商映射 $\mathbb{I}_K\to\mathbf{C}_K$ 诱导同构
\[
\mathbb{I}_{K}/K^{\times}\cdot\Nm(\mathbb{I}_{L})\to\mathbf{C}_{K}/\Nm(\mathbf{C}_{L}).
\]

\section{单位群的上同调}\label{sec:VII.3}

设 $L/K$ 是数域的有限扩张, Galois 群为 $G$. 设 $S\supset S_{\infty}$ 为 $K$ 的素的有限集, 并设 $T$ 为 $L$ 中位于 $S$ 的素之上的素集. 由于 $T$ 在 $G$ 的作用下稳定, $T$-单位群
\[
U(T)\overset{\mathrm{def}}{=}\{\alpha\in L\mid\operatorname{ord}_{w}(\alpha)=0,\ \text{一切 }w\notin T\}
\]
也在 $G$ 的作用下稳定.

\begin{proposition}\label{prop:vii.3.1}
在上述情况下, 设 $G$ 循环. 则 Herbrand 商 $h(U(T))$ 有定义, 且满足
\[
n\cdot h(U(T))=\prod_{v\in S}n_{v},
\]
其中 $n=[L:K]$, $n_v=[L^{v}:K_v]$.
\end{proposition}

证明分两步: 我们先证明同一实向量空间中任意两个 $G$-稳定的满格\footnote{回顾: 实向量空间 $V$ 的子群 $M$ 称为\textbf{满格} (full lattice), 如果 $M$ 由 $V$ 的一组基作为 $\Z$-模生成; 等价地, $M$ 有限生成且典范映射 $\R\otimes_{\Z}M\to V$ 是同构. $\Q$-向量空间中满格的定义类似.}有相同的 Herbrand 商; 然后构造两个这样的格, 一个的 Herbrand 商为 $n\cdot h(U(T))$, 另一个的 Herbrand 商为 $\prod n_v$.

\begin{lemma}\label{lem:vii.3.2}
设 $G$ 为有限群, $k$ 为无限域. 设 $M$ 与 $N$ 是 $k[G]$-模, 视为 $k$-向量空间时有限维. 若对某个域 $\Omega\supset k$, $M\otimes_k\Omega$ 与 $N\otimes_k\Omega$ 作为 $\Omega[G]$-模同构, 则它们作为 $k[G]$-模已经同构.
\end{lemma}

\begin{proof}
首先注意: 若 $V$ 是 $k$ 上一个齐次线性方程组的解空间, 则同一方程组在 $\Omega$ 上的解空间有一组坐标在 $k$ 中的基. 实际上, 求解空间的标准算法在 $k$ 上与在 $\Omega$ 上执行给出相同的结果.

$k$-线性映射 $\alpha\colon M\to N$ 是 $G$-同态, 当且仅当 $\alpha(\sigma m)=\sigma\alpha(m)$ 对一切 $m\in M$, $\sigma\in G$ 成立. 给 $M$ 与 $N$ 取定基后, 给出 $k$-线性映射 $\alpha\colon M\to N$ 等价于给出一个矩阵 $A$, 而 ``$\alpha$ 是 $G$-同态''这一条件形如 $A\cdot B(\sigma)=C(\sigma)\cdot A$ (对某些矩阵 $B(\sigma)$, $C(\sigma)$). 这是 $A$ 的系数的线性条件, 故由前述观察, 存在 $k[G]$-同态 $\alpha_1,\ldots,\alpha_r\colon M\to N$, 构成 $\Omega[G]$-同态空间 $M\otimes_k\Omega\to N\otimes_k\Omega$ 的一组 $\Omega$-基.

由于 $M\otimes_k\Omega$ 与 $N\otimes_k\Omega$ 作为 $\Omega[G]$-模同构, 存在 $a_1,\ldots,a_r\in\Omega$ 使 $\sum a_i\alpha_i$ 是同构, 从而行列式非零. 但 $\det(\sum a_i\alpha_i)$ 是系数属于 $k$ 的关于 $a_i$ 的多项式, 而上一句话表明其系数不全为零. 由于 $k$ 无限, 可取 $k$ 中的 $a_i$ 使 $\sum a_i\alpha_i$ 的行列式非零\footnote{这里用到: 系数取自无限域 $k$ 的多项式 $f$ 若对所有 $(a_1,\ldots,a_n)\in k^{n}$ 都有 $f(a_1,\ldots,a_n)=0$, 则 $f=0$ (FT, 5.18 的证明).}, 从而是 $k[G]$-同构 $M\to N$.
\end{proof}

\begin{remark}\label{rem:vii.3.3}
若嫌上述证明过于初等, 这里另给一个. 设 $k$ 特征为零. $M$ 作为 $k[G]$-模的自同构群 $H$ 是若干形如 $\mathrm{GL}_d(D)$ ($D$ 为 $k$ 上的除法代数) 的群之积. 同构 $M\to N$ 的函子是 $H$ 的一个主齐性空间 (非空, 因为存在某个包含 $k$ 的域上的同构), 从而定义了 $H^{1}(k,H)$ 的一个元素. 而 Hilbert 定理 90 的推广表明 $H^{1}(k,H)=1$.
\end{remark}

\begin{lemma}\label{lem:vii.3.4}
设 $G$ 为有限循环群, $M$ 与 $N$ 是 $G$-模, 作为 $\Z$-模有限生成, 且 $M\otimes_{\Z}\Q$ 与 $N\otimes_{\Z}\Q$ 作为 $G$-模同构. 若 $h(M)$ 与 $h(N)$ 之一有定义, 则另一个也有定义, 且两者相等.
\end{lemma}

\begin{proof}
由 (II, \ref{cor:ii.3.9}), 可设 $M$ 与 $N$ 无挠. 取同构
\[
\alpha\colon M\otimes\Q\to N\otimes\Q.
\]
则 $\alpha(M)$ 与 $N$ 是同一 $\Q$-向量空间中的格, 故 $\alpha(M)\subset n^{-1}N$ 对某个 $n\in\N$ 成立 (把 $\alpha(M)$ 的一组基的元素用 $N$ 的基表出, 取系数的公分母为 $n$). 把 $\alpha$ 换成 $n\alpha$ 后, 有 $\alpha(M)\subset N$. 现在有正合列
\[
0\to M\xrightarrow{\alpha}N\to N/\alpha(M)\to 0,
\]
其中 $N/\alpha(M)$ 有限, 再一次应用 (II, \ref{cor:ii.3.9}) 便得 $h(M)=h(N)$.
\end{proof}

\begin{lemma}\label{lem:vii.3.5}
设 $G$ 为有限循环群, $V$ 是一个实向量空间, $G$ 线性地作用于其上 (即 $V$ 是 $\R[G]$-模). 设 $M$ 与 $N$ 是 $V$ 中两个 $G$-稳定的满格. 若 $h(M)$ 与 $h(N)$ 之一有定义, 则另一个也有定义, 且两者相等.
\end{lemma}

\begin{proof}
由于 $M$ 与 $N$ 是 $V$ 中的满格, 典范映射
\[
M\otimes_{\Z}\R\to V,\quad N\otimes_{\Z}\R\to V
\]
是同构. 这些映射是 $G$-同态, 因此由 (\ref{lem:vii.3.2}), 作为 $\Q[G]$-模有 $M\otimes_{\Z}\Q\approx N\otimes_{\Z}\Q$, 从而可应用引理\ref{lem:vii.3.4}.
\end{proof}

现在完成命题\ref{prop:vii.3.1} 的证明. 设 $V$ 是以 $T$ 的元素为指标的若干份 $\R$ 的乘积, 即
\[
V=\operatorname{Hom}(T,\R).
\]
我们让 $G$ 按如下规则作用于 $V$:
\[
(\sigma f)(w)=f(\sigma^{-1}w),\quad\sigma\in G,\quad w\in T.
\]
$V$ 中的第一个格是 $N\overset{\text{def}}{=}\operatorname{Hom}(T,\Z)$. 对每个 $v\in S$, 取位于 $v$ 之上的一个 $w$, 并设 $G^v$ 为 $w$ 的分解群. 集合 $G^{v}\cdot w$ ($v\in S$) 恰是 $G$ 作用在 $T$ 上的轨道. 特别地, $T$ 是这些集合的不交并, 故
\[
\operatorname{Hom}(T,\Z)=\bigoplus_{v}\operatorname{Hom}(G/G^{v},\Z).
\]
但 $\operatorname{Hom}(G/G^{v},\Z)=\operatorname{Ind}_{G^{v}}^{G}(\Z)$ ($\Z$ 视为平凡的 $G^v$-模). 因此
\[
h(G,N)=\prod_{v}h\bigl(G,\operatorname{Ind}_{G^{v}}^{G}(\Z)\bigr)=\prod h(G^{v},\Z)=\prod n_{v}.
\]
现在定义第二个格. 设 $\lambda\colon U(T)\to V$ 为映射 $a\mapsto(\ldots,\log|a|_{w},\ldots)$, 并令 $M^{0}$ 为 $\lambda$ 的像. 注意 $\lambda$ 与 $G$ 的作用交换. $\lambda$ 的核由 $L^{\times}$ 中满足 $|a|_{w}=1$ (对所有 $w$, 包括无限素) 的元素 $a$ 组成, 这些正是 $L$ 中 $1$ 的根, 故 $h(U(T))=h(M^{0})$. 乘积公式表明 $M^{0}$ 包含在 $V$ 的子空间
\[
V^{0}\colon\sum x_{w}=0
\]
之中, 而 $T$-单位定理的证明表明 $M^{0}$ 是 $V^0$ 中的格 (参见 ANT, 5.8). 向量 $e=(1,1,\ldots,1)$ 在 $G$ 下稳定, 我们令 $M=M^{0}+\Z e$. 则 $M\otimes_{\Z}\R=V^{0}+\R e=V$, 故 $M$ 是 $V$ 中的格. 此外,
\[
h(M)=h(M^{0})\cdot h(\Z)=h(M^{0})\cdot n.
\]
这就完成了命题\ref{prop:vii.3.1} 的证明.

\section{Idèle 类的上同调 I: 第一个不等式}\label{sec:VII.4}

设 $L/K$ 是数域的有限 Galois 扩张, Galois 群为 $G$. 我们有 $G$-模的行正合交换图
\[
\begin{tikzcd}
0\ar[r] & K^{\times}\ar[r]\ar[d] & \mathbb{I}_{K}\ar[r]\ar[d] & \mathbf{C}_{K}\ar[r]\ar[d] & 0\\
0\ar[r] & L^{\times}\ar[r] & \mathbb{I}_{L}\ar[r] & \mathbf{C}_{L}\ar[r] & 0.
\end{tikzcd}
\]
各行正合正是 idèle 类群的定义. 左侧方块的竖直箭头是自然包含, 故该方块交换, 这表明右侧的竖直箭头存在.

\begin{lemma}\label{lem:vii.4.1}
典范映射 $\mathbf{C}_K\to\mathbf{C}_L$ 诱导同构
\[
\mathbf{C}_{K}\to\mathbf{C}_{L}^{G}=H^{0}(G,\mathbf{C}_{L}).
\]
\end{lemma}

\begin{proof}
由上图的下一行得到上同调序列
\[
\begin{tikzcd}
0\ar[r] & H^{0}(G,L^{\times})\ar[r]\ar[d, equal] & H^{0}(G,\mathbb{I}_{L})\ar[r]\ar[d, equal] & H^{0}(G,\mathbf{C}_{L})\ar[r]\ar[d, equal] & H^{1}(G,L^{\times})\ar[d, equal]\\
 & L^{\times G} & \mathbb{I}_{L}^{G} & \mathbf{C}_{L}^{G} & 0
\end{tikzcd}
\]
它可以与
\[
0\to K^{\times}\to\mathbb{I}_{K}\to\mathbf{C}_{K}\to 0
\]
等同.
\end{proof}

数域的理想类群有限, 且由素理想的类生成, 因此它由有限多个素理想生成.

\begin{lemma}\label{lem:vii.4.2}
设 $K$ 为数域, $S\supset S_{\infty}$ 为 $K$ 的素的有限集, 包含 $K$ 的理想类群的一组生成元. 则
\[
\mathbb{I}_{K}=K^{\times}\cdot\mathbb{I}_{S}.
\]
\end{lemma}

\begin{proof}
``$S$ 包含 $K$ 的理想类群的一组生成元''这一条件意味着: 每个分式理想 $\mathfrak{a}$ 都能写成
\[
\mathfrak{a}=\mathfrak{b}\cdot(c),
\]
其中 $\mathfrak{b}$ 落在由 $S$ 中的素理想生成的群内, $c\in K^{\times}$. 因此在商群 $I^{S}=I/\langle S\rangle$ 中 $\mathfrak{a}=(c)$, 从而 $I^{S}/i(K^{\times})=0$.
\end{proof}

对 $K$ 的每个有限素集 $S\supset S_{\infty}$, 自然映射 $\mathbb{I}\to I^{S}$ 定义同构 $\mathbb{I}/\mathbb{I}_{S}\to I^{S}$. 两边对 $K^{\times}$ 作商, 得 $\mathbb{I}/K^{\times}\cdot\mathbb{I}_{S}\simeq I^{S}/i(K^{\times})\simeq 0$.

回顾: 我们要证明, 对每个 Abel 扩张 $L/K$ 有 $\mathbf{C}_K/\Nm_{L/K}\mathbf{C}_L\simeq\Gal(L/K)$, 且对每个 Galois 扩张有 $H^{1}(G,\mathbf{C}_L)=1$. 对循环扩张, 这两个断言蕴含 Herbrand 商
\[
h(\mathbf{C}_{L})\overset{\mathrm{def}}{=}\frac{(\mathbf{C}_{K}:\Nm\mathbf{C}_{L})}{\bigl|H^{1}(G,\mathbf{C}_{L})\bigr|}=[L:K].
\]
作为第一步, 我们验证这个等式.

\begin{theorem}\label{thm:vii.4.3}
对数域的每个有限循环扩张 $L/K$, $h(\mathbf{C}_L)=[L:K]$.
\end{theorem}

\begin{proof}
设 $S$ 为 $K$ 的任意一个满足下列条件的有限素集:

(a) $S\supset S_{\infty}$, 即 $K$ 的无限素之集;

(b) $S$ 包含在 $L$ 中分歧的全部素;

(c) $S$ 包含 $\mathfrak{P}\cap\mathcal{O}_K$, 其中 $\mathfrak{P}$ 取遍 $L$ 的理想类群的一组生成元.

设 $T$ 为 $L$ 中位于 $S$ 的素之上的素集. 条件 (c) 蕴含 $\mathbb{I}_L=\mathbb{I}_{L,T}\cdot L^{\times}$, 故
\[
\mathbf{C}_{L}\overset{\mathrm{def}}{=}\mathbb{I}_{L}/L^{\times}=L^{\times}\cdot\mathbb{I}_{L,T}/L^{\times}\simeq\mathbb{I}_{L,T}/L^{\times}\cap\mathbb{I}_{T}.
\]
注意
\[
L^{\times}\cap\mathbb{I}_{T}=\{\alpha\in L\mid\operatorname{ord}_{w}(\alpha)=0,\ \forall w\notin T\}=U(T),
\]
于是
\[
h(\mathbf{C}_{L})=h(\mathbb{I}_{L,T})/h(U(T)).
\]
定理由命题 \ref{prop:vii.2.7} 与命题 \ref{prop:vii.3.1} 推出.
\end{proof}

\begin{corollary}[第一个不等式]\label{cor:vii.4.4}
若 $L/K$ 是 $n$ 次循环扩张, 则
\[
(\mathbb{I}_{K}:K^{\times}\cdot\Nm(\mathbb{I}_{L}))\ge n.
\]
\end{corollary}

\begin{proof}
由于 $h(\mathbf{C}_L)=n$, 其分子必 $\ge n$.
\end{proof}

下面给出第一个不等式的一些应用.

\begin{lemma}\label{lem:vii.4.5}
设 $L$ 是 $K$ 的有限可解扩张 (即 Galois 群可解的有限 Galois 扩张). 若存在 $\mathbb{I}_K$ 的子群 $D$ 使得

(a) $D\subset\Nm_{L/K}\mathbb{I}_L$, 且

(b) $K^{\times}\cdot D$ 在 $\mathbb{I}_K$ 中稠密,

则 $L=K$.
\end{lemma}

\begin{proof}
若 $L\ne K$, 则存在 $L$ 的子域 $K'$, 它在 $K$ 上循环且 $\ne K$. 则
\[
D\subset\Nm_{L/K}(\mathbb{I}_{L})=\Nm_{K'/K}\bigl(\Nm_{L/K}(\mathbb{I}_{L})\bigr)\subset\Nm_{K'/K}(\mathbb{I}_{K'}).
\]
因此 $K^{\times}\cdot\Nm_{K'/K}\mathbb{I}_{K'}$ 在 $\mathbb{I}_K$ 中稠密. 由于它是 $\Nm_{K'/K}\mathbb{I}_{K'}$ 的平移之并, 它是开的 (\ref{prop:vii.2.8}); 又由于它是 $\mathbb{I}_K$ 的子群, 它也是闭的. 于是 $K^{\times}\cdot\Nm_{K'/K}\mathbb{I}_{K'}=\mathbb{I}_K$, 而第一个不等式蕴含 $K'=K$.
\end{proof}

\begin{proposition}\label{prop:vii.4.6}
设 $L$ 是 $K$ 的有限可解扩张. 若 $L\ne K$, 则 $K$ 中有无穷多个素在 $L$ 中不完全分裂.
\end{proposition}

\begin{proof}
设只有有限多个, 并设 $S\supset S_{\infty}$ 为 $K$ 的有限素集, 包含所有不完全分裂的素. 我们对
\[
D=\mathbb{I}^{S}\overset{\mathrm{def}}{=}\{(a_{v})\mid a_{v}=1,\ \text{一切 }v\in S\}
\]
应用引理. 对 $w|v\notin S$, 有 $L_w=K_v$, 故显然 $D\subset\Nm(\mathbb{I}_L)$. 设 $\mathbf{a}=(a_v)\in\mathbb{I}$. 由弱逼近定理 (ANT, 7.20), 存在 $b\in K^{\times}$, 它对所有 $v\in S$ 都在 $K_v$ 中非常接近 $a_v$. 取 $\mathbf{a}'$ 为 $\mathbb{I}^{S}$ 中使 $\mathbf{b}\mathbf{a}'$ 的 $v$-分量对所有 $v\notin S$ 都等于 $a_v$ 的元素. 则 $\mathbf{b}\mathbf{a}'$ 在 $\mathbb{I}_K$ 中接近 $\mathbf{a}$. 故 $K^{\times}\cdot D$ 在 $\mathbb{I}_K$ 中稠密.
\end{proof}

\begin{proposition}\label{prop:vii.4.7}
设 $L/K$ 是具有 Galois 群 $G$ 的有限可解扩张, $T$ 为包含在 $L$ 中分歧者的素理想有限集. 则 $\mathfrak{P}\notin T$ 的 Frobenius 元 $(\mathfrak{P},L/K)$ 生成 $G$.
\end{proposition}

\begin{proof}
必要时把 $T$ 换大, 可设它在 $G$ 下稳定. 于是由 $\mathfrak{P}\notin T$ 处的 Frobenius 元生成的子群 $H$ 是正规的. 令 $E=L^{H}$. 回顾 (\hyperlink{item:v.1.11}{V, 1.11}):
\[
(\mathfrak{P}\cap E,E/K)=(\mathfrak{P},L/K)\big|_{E},
\]
而后者是恒等映射. 因此 $K$ 中不位于 $T$ 的素之下的所有素 $\mathfrak{p}$ 都在 $E$ 中分裂, 这表明 $E=K$. 由 Galois 理论主定理, 这蕴含 $H=G$.
\end{proof}

\begin{corollary}\label{cor:vii.4.8}
对每个 Abel 扩张 $L/K$ 与 $K$ 的每个包含在 $L$ 中分歧者的有限素集 $S\supset S_{\infty}$, 映射
\[
\mathfrak{p}\mapsto(\mathfrak{p},L/K)\colon I^{S}\to\Gal(L/K)
\]
是满射. (回顾: $I^{S}$ 是由不在 $S$ 中的素理想生成的分式理想群.)
\end{corollary}

\begin{proof}
像包含 $\mathfrak{P}$ 整除某个不在 $S$ 中的素 $p$ 的全部 Frobenius 元 $(\mathfrak{P},L/K)$, 而这些元生成 $\Gal(L/K)$.
\end{proof}

\begin{remark}\label{rem:vii.4.9}
当然, 命题 \ref{prop:vii.4.6} 比用复分析可得的结果 (见定理 \ref{thm:vi.3.4}) 弱得多, 但它足以证明互反律.
\end{remark}

\section{Idèle 类的上同调 II: 第二个不等式}\label{sec:VII.5}

\begin{theorem}\label{thm:vii.5.1}
设 $L/K$ 是数域的 Galois 扩张, Galois 群为 $G$. 则

(a) (第二个不等式) 指数 $(\mathbb{I}_K:K^{\times}\cdot\Nm(\mathbb{I}_L))$ 有限, 且整除 $[L:K]$;

(b) 群 $H^{1}(G,\mathbf{C}_{L})=0$;

(c) 群 $H^{2}(G,\mathbf{C}_{L})$ 有限, 且其阶整除 $[L:K]$.
\end{theorem}

\begin{lemma}\label{lem:vii.5.2}
若 $G$ 循环, 则定理的断言 (a)(b)(c) 相互等价 (且 $(\mathbb{I}_K:K^{\times}\cdot\Nm(\mathbb{I}_L))=(H^{2}(G,\mathbf{C}_L):1)=[L:K]$).
\end{lemma}

\begin{proof}
不加限制于 $G$, 我们有
\[
\mathbb{I}_{K}/K^{\times}\cdot\Nm_{L/K}(\mathbb{I}_{L})\simeq\mathbf{C}_{K}/\Nm_{L/K}(\mathbf{C}_{L})=H_{T}^{0}(G,\mathbf{C}_{L}).
\]
若 $G$ 循环, 其上同调是周期的, 故 $H_T^{0}(G,\mathbf{C}_L)\approx H^{2}(G,\mathbf{C}_L)$. 这证明 (a) 与 (c) 等价. 定理\ref{thm:vii.4.3} 断言 Herbrand 商 $h(\mathbf{C}_L)=[L:K]$, 故 (a) 与 (c) 各自等价于 (b).
\end{proof}

\begin{lemma}\label{lem:vii.5.3}
只需对 $G$ 是 $p$-群 ($p$ 为素数) 的情形证明该定理.
\end{lemma}

\begin{proof}
回顾 (II, \ref{cor:ii.1.33}): 若 $H$ 是 $G$ 的 Sylow $p$-子群, 则对每个 $G$-模 $M$, 映射
\[
\operatorname{Res}\colon H_{T}^{r}(G,M)\to H_{T}^{r}(H,M)
\]
在 $p$-primary 分量上是单射. 因此, 若定理对 $L/L^{H}$ 成立, 则 $p$ 不整除 $H_T^{1}(G,\mathbf{C}_L)$ 的阶, 且整除 $H_T^{0}(G,\mathbf{C}_L)$ 与 $H_T^{2}(G,\mathbf{C}_L)$ 之阶的 $p$ 的幂次小于整除 $[L:K]$ 的 $p$ 的幂次. 对所有 $p$ 应用此结论, 便得引理.
\end{proof}

\begin{lemma}\label{lem:vii.5.4}
只需对 $G$ 是素数 $p$ 阶循环群的情形证明该定理.
\end{lemma}

\begin{proof}
由上一引理, 可设 $G$ 是 $p$-群. 我们对 $(G:1)$ 作归纳来证明定理. 由于 $G$ 是 $p$-群, 它有指数为 $p$ 的正规子群 $H$ (见 GT 4.17). 考虑正合列 (II, \ref{prop:ii.1.34})
\[
0\to H^{1}(G/H,\mathbf{C}_{K'})\xrightarrow{\operatorname{Inf}}H^{1}(G,\mathbf{C}_{L})\xrightarrow{\operatorname{Res}}H^{1}(H,\mathbf{C}_{L}),
\]
其中 $K'=L^{H}$. 由假设 $H^{1}(G/H,\mathbf{C}_{K'})=0$, 由归纳假设 $H^{1}(H,\mathbf{C}_{L})=0$. 故 $H^{1}(G,\mathbf{C}_{L})=0$, 即断言 (b) 成立.

由于 $H^{1}(H,\mathbf{C}_{L})=0$, 序列
\[
0\to H^{2}(G/H,\mathbf{C}_{K'})\to H^{2}(G,\mathbf{C}_{L})\to H^{2}(H,\mathbf{C}_{L})
\]
正合, 由此推出断言 (c) 成立.

最后注意
\[
(\mathbf{C}_{K}:\Nm_{L/K}(\mathbf{C}_{L}))=(\mathbf{C}_{K}:\Nm_{K'/K}(\mathbf{C}_{K'}))\cdot(\Nm_{K'/K}(\mathbf{C}_{K'}):\Nm_{L/K}(\mathbf{C}_{L})),
\]
它整除 $p[L:K']$, 因为 $\Nm_{K'/K}$ 诱导满射
\[
\mathbf{C}_{K'}/\Nm_{L/K'}(\mathbf{C}_{L})\to\Nm_{K'/K}(\mathbf{C}_{K'})/\Nm_{L/K}(\mathbf{C}_{L}).
\]
要借助引理\ref{lem:vii.5.2} 完成定理\ref{thm:vii.5.1} 的证明, 剩下的只需证明第二个不等式对 $p$ 次循环扩张成立; 但在 (VI \ref{thm:vi.4.9}) 中我们已对所有有限 Galois 扩张证明了第二个不等式. (关于该命题的理想语言与 idèle 语言两种形式之间的转换, 见命题 \ref{prop:v.4.6}.) 下一节我们给出第二个不等式的一个不依赖于第 \ref{ch:VI} 章的代数证明.
\end{proof}

\begin{remark}\label{rem:vii.5.5}
对数域的有限 Galois 扩张 $L/K$, 我们给它附加了群 $\mathbf{C}_K/\Nm(\mathbf{C}_L)$ 与 $H^{2}(G,\mathbf{C}_L)$. 当 $L/K$ 循环时, 它们 (在取定 $G$ 的一个生成元后) 典范同构, 否则不然. 第一个群总同构于 $G^{\mathrm{ab}}$, 第二个总为 $[L:K]$ 阶循环群. 于是, 当 $G$ Abel 但不循环时, 两群阶相同但不互相同构; 当 $G$ 非 Abel 时, 两群阶不同.
\end{remark}

\section{第二个不等式的代数证明}\label{sec:VII.6}

我们将对 $L/K$ 为 $p$ (素数) 次循环扩张的情形证明第二个不等式.

\begin{lemma}\label{lem:vii.6.1}
只需在 $K$ 包含 $p$ 次单位根的情形下证明第二个不等式.
\end{lemma}

\begin{proof}
设 $\zeta$ 为一个本原 $p$ 次单位根 (取在包含 $L$ 的某个固定的 $K$ 的代数闭包中), 令 $K'=K[\zeta]$, $L'=K'\cdot L=L[\zeta]$. 则 $[K':K]=m|p-1$, 从而与 $p$ 互素. 于是 $K'\cap L=K$, 且有如下图形:
\[
\begin{tikzcd}[row sep=2.2em, column sep=2.2em]
 & & L' \\
 & K' \ar[ur, "p"] & & L \ar[ul, "m"'] \\
 & & K \ar[ur, "p"] \ar[ul, "m"']
\end{tikzcd}
\]
映射
\[
\Gal(L'/K)\to\Gal(L/K)\times\Gal(K'/K)
\]
是同构. 考虑图:
\[
\begin{tikzcd}[column sep=1.4em]
\mathbf{C}_{L} \ar[r, "\Nm_{L/K}"] \ar[d, "i_{L}"'] & \mathbf{C}_{K} \ar[r] \ar[d, "i_{K}"'] & \mathbf{C}_{K}/\Nm\mathbf{C}_{L} \ar[r] \ar[d] & 0 \\
\mathbf{C}_{L'} \ar[r, "\Nm_{L'/K'}"] \ar[d, "\Nm_{L'/L}"'] & \mathbf{C}_{K'} \ar[r] \ar[d, "\Nm_{K'/K}"'] & \mathbf{C}_{K'}/\Nm\mathbf{C}_{L'} \ar[r] \ar[d] & 0 \\
\mathbf{C}_{L} \ar[r, "\Nm_{L/K}"] & \mathbf{C}_{K} \ar[r] & \mathbf{C}_{K}/\Nm\mathbf{C}_{L} \ar[r] & 0
\end{tikzcd}
\]
这里 $i_L$ 与 $i_K$ 是由包含 $\mathbb{I}_L\hookrightarrow\mathbb{I}_{L'}$ 与 $\mathbb{I}_K\hookrightarrow\mathbb{I}_{K'}$ 诱导的映射; $\Nm_{L/K}$ 与 $\Nm_{L'/K'}$ 是映射
\[
x\mapsto\sum\sigma x,\quad\sigma\in\Gal(L/K)\simeq\Gal(L'/K'),
\]
而 $\Nm_{L'/L}$ 与 $\Nm_{K'/K}$ 是映射
\[
x\mapsto\sum\sigma x,\quad\sigma\in\Gal(L'/L)\simeq\Gal(K'/K).
\]
显然左侧的方块交换, 这蕴含图的其余部分存在. 复合
\[
\Nm_{L'/L}\circ i_{L}\quad\text{与}\quad\Nm_{K'/K}\circ i_{K}
\]
都是乘以 $m$. 因此复合
\[
\mathbf{C}_{K}/\Nm\mathbf{C}_{L}\to\mathbf{C}_{K'}/\Nm\mathbf{C}_{L'}\to\mathbf{C}_{K}/\Nm\mathbf{C}_{L}
\]
也是乘以 $m$, 从而是同构 (显然, $\mathbf{C}_K$ 中的 $p$ 次幂都是范, 故 $\mathbf{C}_K/\Nm\mathbf{C}_L$ 被 $p$ 杀死). 特别地, 第二个映射是满射, 故
\[
(\mathbf{C}_{K}:\Nm\mathbf{C}_{L})\ \text{整除}\ (\mathbf{C}_{K'}:\Nm\mathbf{C}_{L'}),
\]
而由假设, 后者整除 $p$.
\end{proof}

我们将对如下情形证明第二个不等式: $K$ 包含本原 $p$ 次单位根, 且 $L$ 是 $K$ 的指数为 $p$ 的有限 Abel 扩张, Galois 群为 $G$. 设 $[L:K]=p^{r}$, 则 $G\approx(\Z/p\Z)^{r}$. 由 Kummer 理论 (见本章附录), 存在 $a_1,\ldots,a_r\in K$ 使
\[
L=K\bigl[a_{1}^{\frac{1}{p}},\ldots,a_{r}^{\frac{1}{p}}\bigr].
\]
设 $S$ 为 $K$ 的素的有限集, 满足:

(a) $S$ 包含无限素;

(b) $S$ 包含 $p$ 的全部因子;

(c) $S$ 充分大, 使得所有 $a_i$ 都是 $S$-单位;

(d) $S$ 包含 $K$ 的理想类群的一组生成元, 从而 $\mathbb{I}_K=\mathbb{I}_{K,S}\cdot K^{\times}$ (见 \ref{lem:vii.4.2}).

注意, 根据 (\ref{prop:vii.A.5}), 条件 (b) 与 (c) 蕴含 $S$ 包含在 $L$ 中分歧的全部素.

照例, 我们以 $U(S)$ 记 $S$-单位群, 即在 $S$ 之外的所有素处都是单位的 $K^{\times}$ 中元素之群. 回顾单位定理:
\[
U(S)\approx\Z^{s-1}\times U(S)_{\mathrm{tors}},\qquad s=|S|.
\]
群 $U(S)_{\mathrm{tors}}$ 循环, 且在当前情形其阶被 $p$ 整除 (因为它包含 $\mu_p$), 故
\[
U(S)/U(S)^{p}\approx(\Z/p\Z)^{s}.
\]
令 $M=K[U(S)^{\frac{1}{p}}]$. 这是与群
\[
U(S)\cdot K^{\times p}/K^{\times p}\simeq U(S)/U(S)\cap K^{\times p}=U(S)/U(S)^{p}\approx(\Z/p\Z)^{s}
\]
相对应的 Kummer 扩张. 于是我们有扩张
\[
M\overset{p^{t}}{\supset}L\overset{p^{r}}{\supset}K,\qquad r+t=s.
\]

\begin{lemma}\label{lem:vii.6.2}
存在与 $S$ 不交的 $K$ 的素的有限集 $T$, 使得 $\{(\mathfrak{p}_{v},M/K)\mid v\in T\}$ 构成 $\Gal(M/L)$ (视为 $\F_p$-向量空间) 的一组基.
\end{lemma}

\begin{proof}
注意 $S$ 包含 $K$ 的在 $M$ 中分歧的全部素 (由 \ref{prop:vii.A.5}). 因此, 对位于不在 $S$ 中的素 $v$ 之上的 $M$ 的素 $w$, 扩张 $M_w/K_v$ 不分歧, 从而是指数为 $p$ 的循环扩张. 于是它要么是 $p$ 阶循环的, 要么平凡. 这样, 若 $M_w\ne L_w$, 则 $L_w=K_v$ (这里我用 $w$ 同时表示 $M$ 的素与其整除的 $L$ 的素).

由 (\ref{prop:vii.4.7}), 存在 $M$ 的素的有限集 $T'$, 其元素均不整除 $S$ 中的素, 使得 $w\in T'$ 的 Frobenius 元 $(\mathfrak{p}_w,M/L)$ 生成 $\Gal(M/L)$. 必要时把 $T'$ 换成其子集, 可设这些 Frobenius 元构成 $\Gal(M/L)$ 的一组基. 设 $w\in T'$ 整除 $L$ 的素 $w_L$ 与 $K$ 的素 $w_K$. 由于 $M_w\ne L_{w_L}$, 有 $L_{w_L}=K_{w_K}$, 从而 $(\mathfrak{p}_w,M/L)=(\mathfrak{p}_w,M/K)$. 于是 $T=\{w_K\mid w\in T'\}$ 具有所需性质.
\end{proof}

注意 $T$ 的阶为 $t$, 其中 $p^{t}=[M:L]$, 且若 $w$ 是 $M$ 的满足 $w|v\in T$ 的素, 则 $L_w=K_v$.

\begin{lemma}\label{lem:vii.6.3}
采用上述记号: 元素 $a\in U(S)$ 在 $L$ 中成为 $p$ 次幂, 当且仅当它对所有 $v\in T$ 都在 $K_v$ 中成为 $p$ 次幂.
\end{lemma}

\begin{proof}
$\implies$: 设 $a\in U(S)$. 若 $a$ 在 $L$ 中成为 $p$ 次幂, 则它对所有 $w$ 都在 $L_w^{\times}$ 中成为 $p$ 次幂. 但若 $w|v\in T$, 则 $L_w=K_v$, 故它在 $K_v$ 中成为 $p$ 次幂.

$\Longleftarrow$: 设 $a\in U(S)$. 则 $a^{\frac{1}{p}}\in M$. 若 $a^{\frac{1}{p}}$ 对所有 $v\in T$ 都落在 $K_v$ 中, 则 $a^{\frac{1}{p}}$ 被所有 $v\in T$ 的 $(\mathfrak{p}_v,M/K)$ 固定, 从而被 $\Gal(M/L)$ 固定. 于是 $a^{\frac{1}{p}}$ 落在 $L$ 中, 故 $a\in L^{p}$.
\end{proof}

\begin{lemma}\label{lem:vii.6.4}
$\mathbb{I}_K$ 的子群
\[
E=\prod_{v\in S}K_{v}^{\times p}\times\prod_{v\in T}K_{v}^{\times}\times\prod_{v\notin S\cup T}U_{v}
\]
包含在 $\Nm_{L/K}(\mathbb{I}_L)$ 中.
\end{lemma}

\begin{proof}
设 $\mathbf{a}=(a_v)\in E$. 须证 $\mathbf{a}$ 的每个分量 $a_v$ 都是范.

$v\in S$: 由局部类域论,
\[
K_{v}^{\times}/\Nm L_{w}^{\times}\xrightarrow{\ \simeq\ }\Gal(L_{w}/K_{v}).
\]
由于第二个群被 $p$ 杀死, 第一个群也必被 $p$ 杀死, 即 $K_v^{\times p}\subset\Nm L_w^{\times}$.

$v\in T$: 此处 $L_w=K_v$, 故 $K_v$ 的每个元素都是来自 $L_w$ 的范.

$v\notin S\cup T$: 由于 $L_w$ 在 $K_v$ 上不分歧, 范映射 $U_w\to U_v$ 是满射 (见 III, \ref{prop:iii.1.2}).
\end{proof}

现在
\[
(\mathbf{C}_{K}:\Nm\mathbf{C}_{L})=(\mathbb{I}_{K}:K^{\times}\cdot\Nm\mathbb{I}_{L}),
\]
它整除 $(\mathbb{I}_K:K^{\times}E)$, 故剩下只需证明
\[
\bigl(\mathbb{I}_{K}:K^{\times}E\bigr)\ \text{整除}\ p^{r}.
\]
但 $\mathbb{I}_K=K^{\times}\cdot\mathbb{I}_S=K^{\times}\cdot\mathbb{I}_{S\cup T}$, 故
\[
(\mathbb{I}_{K}:K^{\times}E)=(K^{\times}\mathbb{I}_{S\cup T}:K^{\times}E).
\]

\begin{lemma}\label{lem:vii.6.5}
设 $A$, $B$, $C$ 是某个 Abel 群的子群, 且 $A\supset B$. 则
\[
(AC:BC)(A\cap C:B\cap C)=(A:B),
\]
其含义是: 若其中两个指数有限, 则第三个也有限, 且等式成立.
\end{lemma}

\begin{proof}
在下述交换图中, 列与前两行显然正合, 由此 (例如用蛇形引理) 推得底行
\[
\begin{tikzcd}[column sep=1.2em, row sep=1.2em]
 & 0 \ar[d] & 0 \ar[d] & 0 \ar[d] & \\
0 \ar[r] & B\cap C \ar[r] \ar[d] & B \ar[r] \ar[d] & BC/C \ar[d] \ar[r] & 0 \\
0 \ar[r] & A\cap C \ar[r] \ar[d] & A \ar[r] \ar[d] & AC/C \ar[d] \ar[r] & 0 \\
0 \ar[r] & A\cap C/B\cap C \ar[r] \ar[d] & A/B \ar[r] \ar[d] & AC/BC \ar[r] \ar[d] & 0 \\
 & 0 & 0 & 0 &
\end{tikzcd}
\]
正合. 这就证明了断言.
\end{proof}

对 $A=\mathbb{I}_{S\cup T}$, $B=E$, $C=K^{\times}$ 应用该引理, 得
\[
(\mathbb{I}_{K}:K^{\times}E)=\frac{(\mathbb{I}_{S\cup T}:E)}{(U(S\cup T):K^{\times}\cap E)}.
\]

\begin{lemma}\label{lem:vii.6.6}
采用上述记号:
\[
(\mathbb{I}_{S\cup T}:E)=p^{2s}.
\]
\end{lemma}

\begin{lemma}\label{lem:vii.6.7}
采用上述记号:
\[
(U(S\cup T):K^{\times}\cap E)=p^{s+t}.
\]
\end{lemma}

由于 $r+t=s$, 这将证明待证的公式.

(\ref{lem:vii.6.6} 的) 证明: 显然 $(\mathbb{I}_{S\cup T}:E)=\prod_{v\in S}(K_{v}^{\times}:K_{v}^{\times p})$. 由于 $S$ 中有 $s$ 个素, 而 $K$ 包含 $p$ 个互异的 $p$ 次单位根, 下述命题表明
\[
\bigl(\mathbb{I}_{S\cup T}:E\bigr)=\frac{p^{2s}}{\prod_{v\in S}|p|_{v}}.
\]
由假设, $S$ 包含所有使 $|p|_{v}\ne1$ 的素, 故
\[
\prod_{v\in S}|p|_{v}=\prod_{\text{一切 }v}|p|_{v},
\]
由乘积公式, 后者等于 $1$.

\begin{proposition}\label{prop:vii.6.8}
设 $K$ 为特征零的局部域, $\mu_n$ 为 $K^{\times}$ 中 $n$ 次单位根之群. 则
\[
(K^{\times}:K^{\times n})=n\cdot\frac{|\mu_{n}|}{|n|}.
\]
若 $K$ 非阿基米德, $U$ 为 $K$ 的单位群, 则
\[
(U:U^{n})=\frac{|\mu_{n}|}{|n|}.
\]
\end{proposition}

\begin{proof}
若 $K=\CC$, 第一式化为 $1=n\cdot\frac{n}{n^{2}}$; 若 $K=\R$, 当 $n$ 为奇数时化为 $1=n\cdot\frac{1}{n}$, 当 $n$ 为偶数时化为 $2=n\cdot\frac{2}{n}$. 若 $K$ 非阿基米德, 则 $K^{\times}\approx U\times\Z$, 故
\[
(K^{\times}:K^{\times n})=(U:U^{n})(\Z:n\Z)=(U:U^{n})\,n.
\]
因此剩下只需证明第二式.

对 Abel 群 $M$, 记
\[
h_{n}(M)=(M:nM)/(M_{n}:1),\quad M_{n}=\{x\in M\mid nx=0\}.
\]
则 $h_n(M)$ 是把 $M$ 视为具有平凡作用的 $\Z/n\Z$-模时的 Herbrand 商, 故可应用 II, \S\ref{sec:II.2} 的结果.

正如 (III, \ref{lem:iii.2.4}) 的证明中所见, 指数映射给出从 $\mathcal{O}_K$ 的一个有限指数子群到 $U$ 的一个有限指数子群的同构. 因此
\[
h_{n}(U)=h_{n}(\mathcal{O}_{K})=(\mathcal{O}_{K}:n\mathcal{O}_{K})\overset{\mathrm{def}}{=}|n|^{-1}.
\]
故
\[
(U:U^{n})=\frac{(U_{n}:1)}{|n|}.
\]
由于 $U_{n}=\mu_{n}$, 第二式得证.
\end{proof}

(\ref{lem:vii.6.7} 的) 证明: 显然 $K^{\times}\cap E\supset U(S\cup T)^{p}$. 由单位定理 (如前) 得 $(U(S\cup T):U(S\cup T)^{p})=p^{s+t}$, 故剩下只需证明
\[
K^{\times}\cap E\subset U(S\cup T)^{p}.
\]
这一点由下面两个引理完成 (第一个表明第二个引理可用来证明该包含关系).

\begin{lemma}\label{lem:vii.6.9}
在上述假设下, 显然映射
\[
U(S)\to\prod_{v\in T}U_{v}/U_{v}^{p}
\]
是满射.
\end{lemma}

\begin{proof}
设 $H$ 为该映射的核. 为证映射满射, 我们将证明
\[
(U(S):H)=\prod_{v\in T}(U_{v}:U_{v}^{p}).
\]
由于 $T$ 与 $S$ 不交, 对所有 $v\in T$ 有 $|p|_{v}=1$, 故由 (\ref{prop:vii.6.8}) 右端为 $p^{t}$. 另一方面, 由引理\ref{lem:vii.6.3}, $H=U(S)\cap L^{\times p}$, 故
\[
U(S)/H=U(S)/U(S)\cap L^{\times p}\simeq U(S)\cdot L^{\times p}/L^{\times p}.
\]
最后一个群由 Kummer 理论 (见 \ref{thm:vii.A.3}) 对应于扩张 $M/L$, 从而阶为 $[M:L]=p^{t}$.
\end{proof}

\begin{proposition}\label{prop:vii.6.10}
设 $K$ 为包含本原 $n$ 次单位根的数域. 设 $S$ 为包含无限素、$n$ 的因子以及 $K$ 的理想类群的一组代表元的素集. 设 $T$ 为与 $S$ 不交的素集, 使得
\[
U(S)\to\prod_{v\in T}U_{v}/U_{v}^{n}
\]
是满射. 设 $b\in K^{\times}$ 对所有 $v\in S$ 都是 $K_v$ 中的 $n$ 次幂, 且在 $S\cup T$ 之外是单位. 则 $b\in K^{\times n}$.
\end{proposition}

\begin{proof}
令 $L=K[b^{\frac{1}{n}}]$, 须证 $L=K$. 令
\[
D=\prod_{v\in S}K_{v}^{\times}\times\prod_{v\in T}U_{v}^{n}\times\prod_{v\notin S\cup T}U_{v}.
\]
由引理\ref{lem:vii.4.5}, 为证 $L=K$, 只需证明:

(a) $D\subset\Nm_{L/K}\mathbb{I}_L$, 且

(b) $D\cdot K^{\times}=\mathbb{I}_K$.

(a) 设 $\mathbf{d}=(d_v)\in D$. 须验证 $d_v$ 对所有 $v$ 都是来自 $K_v[b^{\frac{1}{n}}]$ 的范. $v\in S$: 此时 $K_v[b^{\frac{1}{n}}]=K_v$, 故 $K_v$ 的每个元素都是范. $v\in T$: 由局部类域论, 指数 $(K_{v}^{\times}:\Nm K_{v}[b^{\frac{1}{n}}]^{\times})$ 等于次数 $[K_v[b^{\frac{1}{n}}]:K_v]$, 后者整除 $n$. 故 $K_v$ 中的每个 $n$ 次幂都是范. $v\notin S\cup T$: 因为 $nb$ 在 $v$ 处是单位, 域 $K_v[b^{\frac{1}{n}}]$ 在 $K_v$ 上不分歧, 故每个单位都是范.

(b) 显然 $\mathbb{I}_S/D=\prod_{v\in T}U_v/U_v^{n}$, 且由假设 $U(S)\to\prod_{v\in T}U_v/U_v^{n}$ 是满射. 故 $\mathbb{I}_S=D\cdot U(S)$, 从而
\[
\mathbb{I}_{K}=\mathbb{I}_{S}\cdot K^{\times}=D\cdot U(S)\cdot K^{\times}=D\cdot K^{\times}.
\]
\end{proof}

这就完成了定理\ref{thm:vii.5.1} (第二个不等式) 的证明.

\section{对 Brauer 群的应用}\label{sec:VII.7}

我们令 $H^{2}(L/K)=H^{2}(\Gal(L/K),L^{\times})$ 且 $H^{2}(-/K)=H^{2}(\Gal(K^{\mathrm{al}}/K),K^{\mathrm{al}\times})$. 读过第 \ref{ch:IV} 章的读者会认出: $H^{2}(L/K)\simeq\mathrm{Br}(L/K)$, $H^{2}(-/K)=\mathrm{Br}(K)$.

\begin{theorem}\label{thm:vii.7.1}
对数域的每个 Galois 扩张 $L/K$ (可以是无限的), 典范映射
\[
H^{2}(L/K)\to\bigoplus_{v}H^{2}(L^{v}/K_{v})
\]
是单射.
\end{theorem}

\begin{proof}
先设 $L/K$ 是具有 Galois 群 $G$ 的有限 Galois 扩张. 由于 $H^{1}(G,\mathbf{C}_{L})=0$,
\[
0\to L^{\times}\to\mathbb{I}_{L}\to\mathbf{C}_{L}\to 0
\]
的上同调序列为
\[
0\to H^{2}(G,L^{\times})\to H^{2}(G,\mathbb{I}_{L})\to\cdots.
\]
但
\[
H^{2}(G,L^{\times})=H^{2}(L/K),
\]
且 (见 \ref{prop:vii.2.5})
\[
H^{2}(G,\mathbb{I}_{L})=\bigoplus H^{2}(G^{v},L^{v\times})=\bigoplus H^{2}(L^{v}/K_{v}),
\]
这就证明了该情形下的定理. 对无限扩张的情形, 在有限 Galois 子扩张上取极限即可.
\end{proof}

称域的扩张 $L/K$ 为\textbf{分圆}的 (cyclotomic), 若 $L\subset K[\zeta]$, 其中 $\zeta$ 是某个单位根. 下述命题将在整体互反律的证明中发挥作用.

\begin{proposition}\label{prop:vii.7.2}
对每个 $\beta\in H^{2}(-/K)$, 存在 $K$ 的循环分圆扩张 $L$, 使 $\beta$ 在 $H^{2}(-/L)$ 中的像为零.
\end{proposition}

\begin{proof}
定理表明 $\beta$ 由其在 $H^{2}(K_v)$ 中的像决定, 从而由不变量 $\mathrm{inv}_{v}(\beta_{v})\in\Q/\Z$ 决定 (见第 \ref{ch:III} 章的定理\ref{thm:iii.2.1}). 对 $K$ 的每个有限扩张 $L$ 与 $L$ 的素 $w|v$, 有 $\mathrm{inv}_{w}(\beta|L)=[L_{w}:K_{v}]\cdot\mathrm{inv}_{v}(\beta)$ (同前), 故只需求一个循环分圆扩张 $L/K$, 使得对所有 $v$,
\[
[L^{v}:K_{v}]\cdot\operatorname{inv}_{v}(\beta_{v})=0\ \bmod\ \Z.
\]
注意: 由于 $H^{2}(-/K)$ (由定理\ref{thm:vii.7.1}) 单射地映入局部群的直和, $\mathrm{inv}_{v}(\beta_{v})=0$ 对几乎所有 $v$ 成立. 故存在整数 $m>0$ 使 $m\,\mathrm{inv}_{v}(\beta_{v})=0$ 对所有 $v$ 成立. 具有所需性质的 $L$ 的存在性由下一个引理保证.
\end{proof}

\begin{lemma}\label{lem:vii.7.3}
给定数域 $K$、$K$ 的有限素的有限集 $S$ 与整数 $m>0$, 存在 $K$ 的全复循环分圆扩张 $L$, 使得 $m|[L^{v}:K_{v}]$ 对所有 $v\in S$ 成立.
\end{lemma}

\begin{proof}
只需对 $\Q$ 与 $m\cdot[K:\Q]$ 证明. 故不妨设 $K=\Q$.

设 $l$ 为素数, $\zeta$ 为本原 $l^{r}$ 次单位根, $r>2$. 则 $\Gal(\Q[\zeta]/\Q)\simeq(\Z/l^{r}\Z)^{\times}$, 且
\[
(\Z/l^{r}\Z)^{\times}\approx
\begin{cases}
\Delta\times C(l^{r-1}) & l\ \text{为奇数},\\
\Delta\times C(2^{r-2}) & l=2,
\end{cases}
\]
其中 $\Delta$ 在两种情形下分别为 $l-1$ 阶与 $2$ 阶, $C(t)$ 表示 $t$ 阶循环群 (\cite{Serre1970}, II, 3.2). 因此 $L(l^{r})\overset{\text{def}}{=}\Q[\zeta]^{\Delta}$ 是 $\Q$ 的 $l^{r-2}$ 次或 $l^{r-3}$ 次循环分圆扩张.

再看 $\Q_p[\zeta]$. 若 $p=l$, 则 $\Q[\zeta]$ 在 $p$ 上全分歧, 故 $[\Q_p[\zeta]:\Q_p]=[\Q[\zeta]:\Q]=\varphi(l^{r})$. 若 $p\ne l$, 则 $\Q[\zeta]$ 在 $p$ 上全不分歧, 而 $[\Q_p[\zeta]:\Q_p]$ 是使 $l^{r}|p^{t}-1$ 的最小整数 $t$. 无论哪种情形, 当 $r\to\infty$ 时都有 $[\Q_p[\zeta]:\Q_p]\to\infty$. 于是对每个 $p$, $[L(l^{r})_p:\Q_p]$ 是随 $r\to\infty$ 趋于 $\infty$ 的 $l$ 的幂.

互异素数幂阶的循环群之积仍是循环群. 因此, 对互异的素数 $l_1,\ldots,l_s$, $L=L(l_1^{r_1})\cdots L(l_s^{r_s})$ 是循环的, 并且显然, 取 $l_1^{r_1}\ldots l_s^{r_s}$ 充分大, 即可保证局部次数 $[L_p:\Q_p]$ 对所有 $p\in S$ 都被 $m$ 整除.
\end{proof}

用更具体的话说, 上述两个结果表明:

若 $K$ 上的中心单代数在所有 $K_v$ 上分裂, 则它在 $K$ 上分裂;

$K$ 上的每个中心单代数都在 $K$ 的某个循环分圆扩张上分裂.

\section{互反律证明的完成}\label{sec:VII.8}

回顾: 对数域的有限 Abel 扩张 $L/K$ (Galois 群为 $G$), 我们定义了同态 $\phi_{L/K}\colon\mathbb{I}_K\to G$, $\phi_{L/K}(\mathbf{a})=\prod_v\phi_v(a_v)$.

\begin{theorem}\label{thm:vii.8.1}
(a) 设 $L/K$ 是数域的有限 Abel 扩张. 则 $\phi_{L/K}$ 在主 idèle 子群 $K^{\times}\subset\mathbb{I}_K$ 上取值 $1$.

(b) 设 $L/K$ 是数域的有限 Galois 扩张. 则对一切 $\alpha\in H^{2}(L/K)$ 有 $\sum\mathrm{inv}_{v}(\alpha)=0$.
\end{theorem}

在证明该定理之前, 我们解释为什么 (a) 蕴含 $L/K$ 的互反律. 断言 (a) 说 $\phi_{L/K}\colon\mathbb{I}_K\to\Gal(L/K)$ 的核包含 $K^{\times}$. 我们已经知道它的核包含 $\Nm_{L/K}(\mathbb{I}_L)$\footnote{因为这在局部成立.}, 故它定义同态
\begin{equation*}
\mathbb{I}_{K}/K^{\times}\cdot\Nm_{L/K}\mathbb{I}_{L}\to\Gal(L/K)\tag{$*$}
\end{equation*}
对在 $L$ 中不分歧的 $K$ 的每个有限素 $v$, $\phi_{L/K}$ 把在 $v$-分量处为素元的 idèle 映到 Frobenius 元 $(\mathfrak{p}_v,L/K)$, 故 (\ref{prop:vii.4.7}) 表明 $\phi_{L/K}$ 是满射. 另一方面, 第二个不等式 (\ref{thm:vii.5.1}) 断言
\[
(\mathbb{I}_{K}:K^{\times}\cdot\Nm_{L/K}\mathbb{I}_{L})\le[L:K],
\]
故同态 $(*)$ 是同构.

\begin{example}\label{ex:vii.8.2}
我们对扩张 $\Q[\zeta_m]/\Q$ 验证 (\ref{thm:vii.8.1}a), 其中 $\zeta_m$ 是本原 $m$ 次单位根. 把 $\Gal(\Q[\zeta_m]/\Q)$ 与 $(\Z/m\Z)^{\times}$ 等同: 对与 $m$ 互素的整数 $n$, $[n]$ 表示把 $\zeta_m$ 映到 $\zeta_m^{n}$ 的 $\Q[\zeta_m]$ 的自同构. 只需证: 对所有 $l|m$ 都有 $\phi(a)|_{\Q[\zeta_{l^{r}}]}=1$. 于是可设 $m=l^{r}$, $m\ne 2$.

同态 $\phi_{\infty}\colon\R^{\times}/\Nm(\CC^{\times})\to\Gal(\Q[\zeta_m]/\Q)$ 把每个负实数映为共轭运算. 故 $\phi_{\infty}(a)=[\operatorname{sign}(a)]$.

设 $a=up^{s}\in\Q_p^{\times}$. 若 $p\ne l$, 则 $p$ 在 $\Q[\zeta_m]$ 中不分歧, 且 $\phi_p(a)$ 的作用是 $p$ 处 Frobenius 元的 $s$ 次幂:
\[
\phi_{p}(up^{s})=[p^{s}].
\]
素数 $l$ 在 $\Q[\zeta_m]$ 中全分歧, 且
\[
\phi_{l}(a)=[u^{-1}]
\]
(见 I, \ref{ex:3.13}).\footnote{正如 Catherine O'Neil 向我指出的, 这里用到了第 \ref{ch:I} 章 \S 3 的一个结果, 与我所说的``可以跳过第 \ref{ch:I} 章 \S 2--4''相悖. 不过, 对 $\Q$ 的分圆扩张, 有可能用初等方法证明定理\ref{thm:vii.8.1} 的 (a). 例如, 用 \hyperlink{item:v.5.10}{V, 5.10} 中的图定义 $\phi_{\Q}$, 再验证其局部分量正确. 这一点将在未来版本中澄清. (另见注记\ref{rem:vii.8.3}.)}

只需在三种情形 $a=-1$, $a=l$, $a=$ 素数 $q\ne l$ 下证明 $\phi(a)=1$. 我们有:
\[
\phi_{p}(-1)=
\begin{cases}
[-1] & \text{若 } p=\infty,\\
[-1] & \text{若 } p=l,\\
[1] & \text{若 } p\ne l,\infty,
\end{cases}
\qquad
\phi_{p}(l)=
\begin{cases}
[1] & \text{若 } p=l,\\
[1] & \text{若 } p\ne l,
\end{cases}
\]
\[
\phi_{p}(q)=
\begin{cases}
[q] & \text{若 } p=q,\\
[q^{-1}] & \text{若 } p=l,\\
[1] & \text{若 } p\ne l,q.
\end{cases}
\]
每种情形下都有 $\prod\phi_p(a)=1$.
\end{example}

\begin{remark}\label{rem:vii.8.3}
在例 V \ref{ex:v.4.10} 中, 我们证明了: 附加到 Artin 映射 $\phi\colon C_{\infty(m)}\to\Gal(\Q[\zeta_m]/\Q)$ 上的同态 $\phi\colon\mathbb{I}_{\Q}\to\Gal(\Q[\zeta_m]/\Q)$ 的局部分量等于局部 Artin 映射. 由于按定义 $\phi(\Q^{\times})=1$, 这给出了 (\ref{thm:vii.8.1}a) 在 $\Q[\zeta_m]/\Q$ 情形的另一个证明.
\end{remark}

\begin{lemma}\label{lem:vii.8.4}
(a) 若 \ref{thm:vii.8.1}(a) 对 $L/K$ 成立, 则它对每个子扩张也成立.

(b) 若 \ref{thm:vii.8.1}(a) 对 $L/K$ 成立, 则它对 $L\cdot K'/K'$ 也成立, 其中 $K'\supset K$ 为任意数域.
\end{lemma}

\begin{proof}
(a) 设 $L\supset K'\supset K$. 则 $\phi_{K'/K}$ 是 $\phi_{L/K}$ 与限制映射 $\Gal(L/K)\to\Gal(K'/K)$ 的复合 (因为这对局部 Artin 映射成立).

(b) 令 $L'=L\cdot K'$. 对 $K'$ 的每个素 $w$, 有交换图
\[
\begin{tikzcd}
K_{w}^{\prime\times}\ar[r, "\phi_{w}"]\ar[d, "\Nm"'] & \Gal(L^{\prime w}/K'_{w})\ar[d, "\operatorname{res}"]\\
K_{v}^{\times}\ar[r, "\phi_{v}"'] & \Gal(L^{v}/K_{v}).
\end{tikzcd}
\]
把它们合并, 得交换图:
\[
\begin{tikzcd}
\mathbb{I}_{K'}\ar[r, "\phi_{L'/K'}"]\ar[d, "\Nm"'] & \Gal(L'/K')\ar[d, "\operatorname{res}"]\\
\mathbb{I}_{K}\ar[r, "\phi_{L/K}"'] & \Gal(L/K).
\end{tikzcd}
\]
由于 idèle 上的范映射把 $K'^{\times}$ 映入 $K^{\times}$, 由先前的结果即得本引理.
\end{proof}

由例\ref{ex:vii.8.2} 与引理\ref{lem:vii.8.4}, 我们得到: (\ref{thm:vii.8.1}a) 对数域 $K$ 的所有分圆扩张\footnote{称扩张 $L/K$ 为\textbf{分圆}的 (cyclotomic), 若 $L\subset K[\zeta]$, $\zeta$ 为某个单位根.}成立.

接下来需要在定理\ref{thm:vii.8.1} 的两个断言之间建立联系.

\begin{lemma}\label{lem:vii.8.5}
设 $L/K$ 是数域的 Abel 扩张. 若 \ref{thm:vii.8.1}(b) 对 $L/K$ 成立, 则 \ref{thm:vii.8.1}(a) 也成立. 反之, 若 $L/K$ 循环且 \ref{thm:vii.8.1}(a) 对 $L/K$ 成立, 则 \ref{thm:vii.8.1}(b) 也成立.
\end{lemma}

\begin{proof}
设 $\chi\in\operatorname{Hom}(G,\Q/\Z)$. 可把 $\chi$ 视为 $H^{1}(G,\Q/\Z)$ 的元素, 其在由序列
\[
0\to\Z\to\Q\to\Q/\Z\to 0
\]
产生的边缘映射下的像是一个元素 $\delta\chi\in H^{2}(G,\Z)$. 考虑图
\[
\begin{tikzcd}
K^{\times}\ar[r]\ar[d] & \mathbb{I}_{K}\ar[r, "\phi_{L/K}"]\ar[d] & G\ar[d, "\chi"]\\
H^{2}(G,L^{\times})\ar[r] & H^{2}(G,\mathbb{I}_{L})\ar[r] & \Q/\Z.
\end{tikzcd}
\]
前两个竖直箭头是与 $\delta\chi$ 的杯积: 若 $\delta\chi$ 由 2-上循环 $n_{\sigma,\tau}$ 表示, 则 $x$ 的像由 2-上循环 $\sigma,\tau\mapsto x^{n_{\sigma,\tau}}$ 表示. 显然左侧方块交换. 右侧竖直映射就是 $\chi$ 本身. 右侧方块交换由 (III, \ref{prop:iii.3.6}) 推出. 现设 \ref{thm:vii.8.1}(b) 对 $L/K$ 成立. 则 $\chi(\phi_{L/K}(a))=0$ 对 $G$ 的所有特征 $\chi$ 成立, 故 $\phi_{L/K}(a)$ 本身为零. 反之, 当 $G$ 循环时, 可取 $\chi$ 为单射, 于是 8.1(a) 蕴含 8.1(b), 因为此时 $\cup\delta\chi$ 是同构 (II, \ref{rem:ii.3.5}).
\end{proof}

既然我们已对分圆扩张知道 \ref{thm:vii.8.1}(a), 也就对循环分圆扩张知道 \ref{thm:vii.8.1}(b). 此外, 一旦证明了定理的 (b), 整个定理便得证. 于是, 下述结果完成了定理的证明.

\begin{lemma}\label{lem:vii.8.6}
若 \ref{thm:vii.8.1}(b) 对循环分圆扩张成立, 则它对所有有限 Galois 扩张成立.
\end{lemma}

\begin{proof}
设 $\beta\in H^{2}(-/K)$. 已知: 若 $\beta\in H^{2}(L/K)$, 其中 $L/K$ 为某个循环分圆扩张, 则 $\sum\mathrm{inv}_{v}(\beta)=0$, 其中 $\beta_v$ 是 $\beta$ 在 $H^{2}(K_v)$ 中的像; 而命题\ref{prop:vii.7.2} 说 $H^{2}(-/K)$ 中的每个 $\beta$ 都落在某个循环分圆扩张 $L/K$ 的 $H^{2}(L/K)$ 中.

[附注: 这里用到 inflation-restriction 序列
\[
0\to H^{2}(\Gal(L/K),L^{\times})\to H^{2}(G_{K},\bar{K}^{\times})\to H^{2}(G_{L},\bar{L}^{\times}),
\]
且如上的 $\beta$ 落在中间的群中, 并在右侧的群中为零.]
\end{proof}

\section{存在定理}\label{sec:VII.9}

本节证明存在定理: idèle 类群的每个有限指数开子群都是范群. 证明的大部分可以从第 \ref{sec:VII.6} 节中抽取出来. 但为了免去一些重复, 我给出一个独立于第 \ref{sec:VII.6} 节的证明 (只依赖几个初等断言).

\begin{lemma}\label{lem:vii.9.1}
若 $U$ 是范群, $V\supset U$, 则 $V$ 也是范群.
\end{lemma}

\begin{proof}
设 $U=\Nm\mathbf{C}_L$. 根据互反律, Artin 映射定义同构
\[
\phi\colon\mathbf{C}_{K}/U\to\Gal(L/K).
\]
设 $M$ 为 $\phi(V)$ 的固定域, 则 $\phi$ 定义同构
\[
\mathbf{C}_{K}/V\to\Gal(M/K),
\]
而根据互反律, $\phi\colon\mathbf{C}_K\to\Gal(M/K)$ 的核是 $\Nm_{M/K}\mathbf{C}_M$.
\end{proof}

从素数分解即可看出: 有理数 $a$ 是 $n$ 次幂, 当且仅当它在 $\R$ 中和在所有素数 $p|a$ 的 $\Q_p$ 中都是 $n$ 次幂. 对数域证明类似的命题则需要互反律 (或复分析).

\begin{proposition}\label{prop:vii.9.2}
设 $K$ 为包含本原 $n$ 次单位根的数域, $S\supset S_{\infty}$ 为 $K$ 的素的有限集, 包含所有整除 $n$ 的素以及足以生成 $K$ 的类群的素. 若 $a\in K^{\times}$ 满足:

-- $a$ 对所有 $v\in S$ 都是 $K_v$ 中的 $n$ 次幂;

-- $a$ 对所有 $v\notin S$ 都是 $K_v$ 中的单位.

则 $a$ 是 $K$ 中的 $n$ 次幂.
\end{proposition}

\begin{proof}
令 $L=K[a^{1/n}]$. 由于 $\zeta_n\in K$, 这是 $K$ 的 Abel 扩张.\footnote{来自 Wojtek Wawr\'ow: 由于证明中的 $L=K[a^{1/n}]$ 是循环的, 此证明似乎只用到第一个不等式 (确切地说是引理\ref{lem:vii.4.5}), 而非完整的互反律. 此外, 这个命题似乎在书中任何地方都没有被用到.} 对每个素 $v\in S$, $X^{n}-a$ 在 $K_v[X]$ 中完全分裂, 故对所有 $w|v$ 有 $L_w=K_v$, 从而范映射 $L_{w}^{\times}\to K_{v}^{\times}$ 是满射. 另一方面, $L$ 在每个 $v\notin S$ 处于 $K$ 不分歧, 故范映射 $\Nm_{L/K}\colon U_w\to U_v$ 是满射. 因此 $\Nm_{L/K}(\mathbb{I}_L)\supset\mathbb{I}_S$, 于是
\[
K^{\times}\cdot\Nm_{L/K}(\mathbb{I}_{L})\supset K^{\times}\cdot\mathbb{I}_{S}=\mathbb{I}_{K}.
\]
互反律现在表明 $L=K$, 故 $a$ 是 $K$ 中的 $n$ 次幂.
\end{proof}

\begin{lemma}[存在定理的关键情形]\label{lem:vii.9.3}
设 $K$ 为包含本原 $p$ 次单位根的数域 ($p$ 为素数). 则 $\mathbf{C}_K$ 的每个使 $\mathbf{C}_K/V$ 为被 $p$ 杀死的有限群的开子群 $V$ 都是范群.
\end{lemma}

\begin{proof}
设 $S\supset S_{\infty}$ 为 $K$ 的素的有限集, 包含无限素、整除 $p$ 的素, 以及足以使 $\mathbb{I}_K=K^{\times}\cdot\mathbb{I}_S$ 的素. 设 $L$ 为由 Kummer 理论对应于群 $U(S)\cdot K^{\times p}$ 的 $K$ 的扩张, 即 $L=K[U(S)^{\frac{1}{p}}]$, 并令
\[
E=\prod_{v\in S}K_{v}^{\times p}\times\prod_{v\notin S}U_{v}.
\]
我们将通过验证
\[
\text{(a)}\ E\subset\Nm(\mathbb{I}_{L});\qquad \text{(b)}\ (\mathbb{I}_{K}:K^{\times}\cdot E)=p^{s}=(\mathbb{I}_{K}:K^{\times}\cdot\Nm_{L/K}(\mathbb{I}_{L}))
\]
来证明 $K^{\times}\cdot E=K^{\times}\cdot\Nm(\mathbb{I}_L)$.

对 $K$ 的任意素 $v$ 与其上的 $L$ 的素 $w$, 局部 Artin 映射是同构
\[
K_{v}^{\times}/\Nm(L_{w}^{\times})\to\Gal(L_{w}/K_{v}).
\]
由于 $L/K$ 具有指数 $p$, $\Nm(L_{w}^{\times})\supset K_{v}^{\times p}$.

对任意 $v\notin S$, $L$ 在 $v$ 处于 $K$ 不分歧, 故范映射 $U_w\to U_v$ 是满射.

合并上述两段的结论, 便得 (a).

由互反律,
\[
(\mathbb{I}_{K}:K^{\times}\cdot\Nm(\mathbb{I}_{L}))=[L:K],
\]
而由 Kummer 理论,
\[
[L:K]=(U(S)\cdot K^{\times p}:K^{\times p}).
\]
但
\[
U(S)\cdot K^{\times p}/K^{\times p}\approx U(S)/U(S)\cap K^{\times p}.
\]
若 $a^{p}\in U(S)$, 则 $a\in U(S)$, 故 $U(S)\cap K^{\times p}=U(S)^{p}$. 于是, 由 Dirichlet 单位定理 (ANT, 5.9),
\[
U(S)\approx U(S)_{\mathrm{torsion}}\oplus\Z^{s-1},
\]
且 $U(S)_{\mathrm{torsion}}$ 是 $K$ 中单位根之群, 是阶被 $p$ 整除的循环群. 故 $(U(S):U(S)^{p})=p^{s}$.

另一方面,
\[
(\mathbb{I}_{K}:K^{\times}\cdot E)=(\mathbb{I}_{S}\cdot K^{\times}:E\cdot K^{\times}),
\]
由 (\ref{lem:vii.6.5}), 它等于
\[
(\mathbb{I}_{S}:E)/(\mathbb{I}_{S}\cap K^{\times}:E\cap K^{\times}).
\]
因此 (见 \ref{prop:vii.6.8}),
\[
\bigl(\mathbb{I}_{S}:E\bigr)=\prod_{v\in S}(K_{v}^{\times}:K_{v}^{\times p})=\prod_{v\in S}\frac{p}{|p|_{v}}\cdot p=p^{2s}.
\]
这里用到 $K$ 包含本原 $p$ 次单位根, 且 $S$ 包含所有使 $|p|_{v}\ne1$ 的 $v$, 故由乘积公式 $\prod_{v\in S}|p|_{v}=\prod_{\text{一切 }v}|p|_{v}=1$. 由此推出 $K^{\times}\cdot E=K^{\times}\cdot\Nm\mathbb{I}_{L}$.

现在设 $V$ 为 $\mathbf{C}_K$ 的使 $\mathbf{C}_K/V$ 被 $p$ 杀死的开子群, $U$ 为 $V$ 在 $\mathbb{I}_K$ 中的原像. 则 $U$ 在 $\mathbb{I}_K$ 中开, 故存在素的有限集 $S$ 使 $U\supset\prod_{v\in S}1\times\prod_{v\notin S}U_v$. 此外, $\mathbb{I}_K/U$ 具有指数 $p$, 故 $U\supset\mathbb{I}_K^{p}$. 于是 $U\supset E\cdot K^{\times}$, 而 $E\cdot K^{\times}/K^{\times}$ 是范群, 故 $U/K^{\times}=V$ 也是范群.
\end{proof}

为简便起见, 在下一个引理的证明中, 我们假定范式限制定理 (该定理要到下一章才证明). 关于避免使用该定理的证明, 见 \cite{CasselsFrohlich1967} 中 Tate 的文章 p.~202.

\begin{lemma}\label{lem:vii.9.4}
设 $U$ 为 $\mathbf{C}_K$ 的有限指数开子群. 若存在有限扩张 $K'/K$ 使得 $\Nm_{K'/K}^{-1}(U)$ 是范群, 则 $U$ 也是范群.
\end{lemma}

\begin{proof}
记 $U'$ 为 $\Nm_{K'/K}^{-1}(U)$, 并设 $L$ 是满足 $\Nm\mathbf{C}_L=U'$ 的 $K'$ 的 Abel 扩张. 若 $M$ 是 $L/K$ 的极大 Abel 子扩张, 则
\[
\Nm_{M/K}\mathbf{C}_{M}=\Nm_{L/K}\mathbf{C}_{L}=\Nm_{K'/K}U'\subset U,
\]
于是可应用引理\ref{lem:vii.9.1}.
\end{proof}

\begin{theorem}\label{thm:vii.9.5}
$\mathbf{C}_K$ 的每个有限指数子群都是范群.
\end{theorem}

\begin{proof}
对 $U$ 的指数作归纳. 若指数为 $1$, 则无须证明. 否则, 存在素数 $p$ 整除 $(\mathbf{C}_K:U)$. 由 (\ref{lem:vii.9.4}), 可设 $K$ 包含 $p$ 次单位根. 取 $\mathbf{C}_K$ 的包含 $U$ 且在 $\mathbf{C}_K$ 中指数为 $p$ 的子群 $U_1$. 由 (\ref{lem:vii.9.3}), 存在 $K$ 的 Abel 扩张 $K'$ 使 $\Nm_{K'/K}\mathbf{C}_{K'}=U_1$; 并且 $K'$ 是 $K$ 上 $p$ 次循环扩张. 令 $U'=\Nm_{K'/K}^{-1}U$. 映射
\[
\Nm_{K'/K}\colon\mathbf{C}_{K'}\to\mathbf{C}_{K}/U
\]
的像为 $U_1/U$, 核为 $U'$. 故
\[
(\mathbf{C}_{K'}:U')=(\mathbf{C}_{K}:U)/p,
\]
由归纳假设 $U'$ 是范群. 现在可应用 (\ref{lem:vii.9.4}) 推出 $U$ 是范群.
\end{proof}

\phantomsection
\section*{A 附录: Kummer 理论}
\addcontentsline{toc}{section}{A 附录: Kummer 理论}
\renewcommand{\thetheorem}{A.\arabic{theorem}}
\setcounter{theorem}{0}

本附录中, $K$ 始终是包含本原 $n$ 次单位根 $\zeta$ 的域. 特别地, $K$ 的特征为 $0$, 或特征为不整除 $n$ 的 $p$. 以 $\mu_n$ 记 $K$ 中 $n$ 次单位根之群. 则 $\mu_n$ 是 $K^{\times}$ 的 $n$ 阶循环子群, 以 $\zeta$ 为生成元.

\begin{proposition}\label{prop:vii.A.1}
设 $L=K[\alpha]$, 其中 $\alpha^{n}\in K$ 且 $\alpha$ 的更小的幂都不在 $K$ 中. 则 $L$ 是 $K$ 的 Galois 扩张, Galois 群为 $n$ 阶循环群. 反之, 若 $L$ 是 $K$ 的 $n$ 次循环扩张, 则 $L=K[\alpha]$, 其中 $\alpha^{n}\in K$.
\end{proposition}

\begin{proof}
见 FT 5.26.
\end{proof}

\begin{proposition}\label{prop:vii.A.2}
$K$ 的两个 $n$ 次循环扩张 $K[a^{\frac{1}{n}}]$ 与 $K[b^{\frac{1}{n}}]$ 相等, 当且仅当 $a=b^{r}c^{n}$, 其中 $r\in\Z$ 与 $n$ 互素, $c\in K^{\times}$; 也就是说, 当且仅当 $a$ 与 $b$ 在 $K^{\times}/K^{\times n}$ 中生成同一个子群.
\end{proposition}

\begin{proof}
见 FT 5.28.
\end{proof}

最后两个结果给出了 $K$ 包含本原 $n$ 次单位根时其 $n$ 次循环扩张的完全分类. 不难把它推广到指数为 $n$ 的全部 Abel 扩张的分类.

\begin{theorem}\label{thm:vii.A.3}
映射
\[
L\mapsto K^{\times}\cap L^{\times n}/K^{\times n}
\]
给出从 $K$ 的含于某个固定代数闭包 $\Omega$ 中的指数为 $n$ 的有限 Abel 扩张到 $K^{\times}/K^{\times n}$ 的有限子群之间的一一对应. 对应于 $B$ 的扩张是 $K[B^{\frac{1}{n}}]$, 即 $\Omega$ 中包含 $K$ 与 $B$ 的每个元素的某个 $n$ 次根的最小子域. 若 $L\leftrightarrow B$, 则 $[L:K]=(B:K^{\times n})$.
\end{theorem}

\begin{proof}
见 FT 5.29.
\end{proof}

\begin{example}\label{ex:vii.A.4}
(a) $\R$ 的二次扩张与 $\R^{\times}/\R^{\times 2}=\{\pm1\}$ 的子群一一对应.

(b) $\Q$ 的指数为 $2$ 的有限 Abel 扩张与
\[
\Q^{\times}/\Q^{\times 2}\approx\{\pm1\}\times\Z/2\Z\times\Z/2\Z\times\cdots
\]
($\Z/2\Z$ 的拷贝以素数为指标) 的有限子群一一对应.
\end{example}

复习完这些代数之后, 我们回到一些数论.

\begin{proposition}\label{prop:vii.A.5}
设 $K$ 为数域, $L=K[a_{1}^{\frac{1}{n}},\cdots,a_{m}^{\frac{1}{n}}]$. 若 $na_i$ 在 $K_v$ 中是单位 (对一切 $i$), 则 $L$ 在 $K$ 的有限素 $v$ 处不分歧.
\end{proposition}

\begin{proof}
只需考虑循环扩张 $K[a^{\frac{1}{n}}]$, 由 \ref{prop:vii.A.1} 可设其次数为 $n$. 令 $\alpha=a^{\frac{1}{n}}$, $f=X^{n}-a$. 则
\[
\operatorname{disc}f=\pm\Nm_{L/K}f'(\alpha)=\pm\Nm_{L/K}n\alpha^{n-1}=\pm n^{n}a^{n-1}.
\]
于是若 $\mathfrak{p}_v$ 不整除 $na$, 它就不整除 $\operatorname{disc}f$, 更不整除 $\operatorname{disc}(\mathcal{O}_L/\mathcal{O}_K)$. (见 ANT, 第 2 章.)
\end{proof}

\begin{remark}\label{rem:vii.A.6}
确定 $K[a^{\frac{1}{n}}]/K$ 在整除 $n$ 的素 $p$ 处的分歧可能相当复杂. 下述概括了 $n=p$ 为素数且 $K$ 包含本原 $p$ 次单位根 $\zeta$ 时的情形. 设 $\pi=\zeta-1$, $a\in K^{\times}$ 与 $p$ 互素且不是 $K$ 中的 $p$ 次幂.

(a) $\mathfrak{p}\mid p$ 在扩张 $K[a^{\frac{1}{p}}]/K$ 中完全分裂, 当且仅当 $X^{p}\equiv a\bmod\mathfrak{p}\pi\mathfrak{p}$ 在 $\mathcal{O}_K$ 中有非零解 (此时称 $a$ 在 $\mathfrak{p}$ 处是\textbf{超 primary} 的, hyperprimary).

(b) $\mathfrak{p}\mid p$ 在扩张 $K[a^{\frac{1}{p}}]/K$ 中不分歧, 当且仅当 $X^{p}\equiv a\bmod\mathfrak{p}\pi$ 在 $\mathcal{O}_K$ 中有非零解 (此时称 $a$ 在 $\mathfrak{p}$ 处是 \textbf{primary} 的).

(c) 扩张 $K[a^{\frac{1}{p}}]/K$ 在 $K$ 的所有有限素处不分歧, 当且仅当 $a$ 是 primary 的且 $(a)=\mathfrak{a}^{p}$ 对某个理想 $\mathfrak{a}$ 成立.

这些结论的证明提示见 \cite{Washington1997}, 习题 9.3, 与 \cite{CasselsFrohlich1967}, 习题 2.12, p.~353. 在 $K=\Q[\zeta]$ 的情形, 另见 \cite{Cassels1986}, pp.~139--140, 与 \cite{FrohlichTaylor1991}, III, 3.11.
\end{remark}

% 恢复定理编号
\renewcommand{\thetheorem}{\thesection.\arabic{theorem}}
  % VII 整体类域论: 证明
% 第八章  补遗
\chapter{补遗}\label{ch:VIII}\index{buyi@补遗(complements)}

在本章中, 我们对类域论作一些补充, 并给出一些应用. 特别地, 我们把 \cite{Serre1970} 的第 I--IV、VI 章中证明的结果推广到任意数域上.

\section{局部的 \texorpdfstring{$n$}{n} 次幂何时是整体的 \texorpdfstring{$n$}{n} 次幂?}\label{sec:VIII.1}

显然, 若数域 $K$ 中的元素 $a$ 是 $K$ 中的 $n$ 次幂, 则它对所有素 $v$ 都是 $K_v$ 中的 $n$ 次幂. 本节研究其逆命题在何种意义上成立.

\begin{theorem}\label{thm:viii.1.1}
设 $K$ 是包含本原 $n$ 次单位根的数域. $K$ 的非零元素是 $K$ 中的 $n$ 次幂, 只要它对除可能有限多个之外的每个素 $v$ 都是 $K_v$ 中的 $n$ 次幂.
\end{theorem}

\begin{proof}
回顾: 对包含本原 $n$ 次单位根的域 $k$, 多项式 $X^{n}-a$ 只要在 $k$ 中有一个根, 就在 $k[X]$ 中完全分裂. 设 $a$ 是 $K$ 的非零元素, 并设 $\beta$ 是 $a$ 在某个扩域中的 $n$ 次根. 若 $a$ 对几乎所有 $v$ 都是 $K_v$ 中的 $n$ 次幂, 则 $X^{n}-a$ 对几乎所有 $v$ 在 $K_v[X]$ 中完全分裂, 从而 $v$ 在 $K[\beta]$ 中完全分裂. 由于这对几乎所有 $v$ 成立, (VII, \ref{prop:vii.4.6}) 表明 $K[\beta]=K$.
\end{proof}

\begin{remark}\label{rem:viii.1.2}
如果不用 (VII, \ref{prop:vii.4.6}) 而用 (VI, \ref{thm:vi.3.4}), 则得到更强的结果: 在该定理的假设下, $K$ 的非零元素是 $K$ 中的 $n$ 次幂, 只要它在密度 $>1/2$ 的一组素 $v$ 处成为 $K_v$ 中的 $n$ 次幂.
\end{remark}

人们或许希望该定理在 $K$ 不含本原 $n$ 次单位根时仍然成立, 但下述习题表明并非如此, 甚至对 $K=\Q$ 与 $n=8$ 都不成立.

\begin{exercise}\label{exc:viii.1.3}
证明 $16$ 在 $\R$ 与所有奇素数 $p$ 的 $\Q_p$ 中都是 $8$ 次幂, 但 (古希腊人早已知道) 它在 $\Q$ 中不是. (提示: 证明 $\Q[\zeta_8]=\Q[i,\sqrt{2}]$ 在所有奇素数处不分歧, 并由此推出: 对奇数 $p$, $\Q_p$ 至少包含 $1+i$, $\sqrt{2}$, $\sqrt{-2}$ 三者之一.)
\end{exercise}

域 $\Q[\zeta_8]$ 在 $\Q$ 上不是循环的, 故该习题与下述定理并不矛盾.

\begin{theorem}\label{thm:viii.1.4}
设 $K$ 是数域, 并设 $n$ 为整数, 使得 $K[\zeta_{2^{t}}]$ 在 $K$ 上是循环的, 其中 $2^{t}$ 是整除 $n$ 的 2 的幂. $K$ 的非零元素是 $K$ 中的 $n$ 次幂, 只要它对除可能有限多个之外的每个素 $v$ 都是 $K_v$ 中的 $n$ 次幂.
\end{theorem}

\begin{proof}
第一步. 不妨设 $n$ 为素数幂来证明该定理. 设 $n=n_1n_2$, 其中 $n_1$ 与 $n_2$ 互素, 并设定理对 $n_1$ 与 $n_2$ 成立. 于是对几乎所有 $v$ 都是 $K_v$ 中 $n$ 次幂的非零元素 $a\in K$ 既是 $K$ 中的 $n_1$ 次幂又是 $K$ 中的 $n_2$ 次幂, 比如 $a=b^{n_1}$, $a=c^{n_2}$. 现在存在整数 $r,s$ 使得 $rn_1+sn_2=1$, 故
\[
a=a^{rn_1}a^{sn_2}=c^{rn_1n_2}b^{sn_1n_2}\in K^{\times n}.
\]

第二步. 当 $n$ 是素数 $p$ 的幂且 $K[\zeta_n]/K$ 是 $p$ 幂阶的循环扩张时, 定理成立. 令 $K'=K[\zeta_n]$, 并设 $a$ 是对几乎所有 $v$ 成为 $K_v$ 中 $n$ 次幂的非零元素. 根据定理\ref{thm:viii.1.1}, $a$ 在 $K'$ 中成为 $n$ 次幂, 比如 $a=\beta^{n}$, 于是
\[
X^{n}-a=\prod_{j=0}^{n-1}(X-\beta\zeta_n^{j})\qquad\text{在 } K'[X] \text{ 中}.
\]
设
\[
X^{n}-a=\prod_{i}f_{i}(X)
\]
是 $X^{n}-a$ 在 $K[X]$ 中分解为不可约因子的分解. 对每个 $i$, 取 $f_i(X)$ 的一个根 $\beta_i=\beta\zeta_n^{j(i)}$. 则 $K[\beta_i]\subset K'$, 从而 $K[\beta_i]$ 是 $K$ 的 Abel 扩张, 特别地, 它在 $K$ 上是 Galois 的, 因而是 $f_i(X)$ 的分裂域. 设 $v$ 是使 $a$ 为 $K_v$ 中 $n$ 次幂的 $K$ 的素. 由假设, $X^{n}-a$ 在 $K_v$ 中有根, 故至少某个 $f_i(X)$ 在 $K_v$ 中有根. 对那个特定的 $i$, $v$ 在 $K[\beta_i]$ 中完全分裂. 于是我们看到: 每个 $v\notin S$ 至少在诸域 $K[\beta_i]$ 之一中完全分裂, 但不同的 $v$ 可能在不同的域中分裂, 由此我们得不出任何结论. 要看出所有的 $v$ 都在同一个 $K[\beta_i]$ 中分裂, 必须用到 $K'$ 在 $K$ 上是素数幂阶循环的这一假设. 该假设蕴含中间域是线性有序的:
\[
K'\supset\cdots\supset K_{3}\supset K_{2}\supset K_{1}\supset K.
\]
取 $i_0$ 使 $K[\beta_{i_0}]$ 是诸 $K[\beta_i]$ 中最小的一个. 则每个 $v\notin S$ 在一个包含 $K[\beta_{i_0}]$ 的域中完全分裂, 从而在 $K[\beta_{i_0}]$ 中完全分裂. 故 $K[\beta_{i_0}]=K$, 即 $X^{n}-a$ 在 $K$ 中至少有一个根.

第三步. 定理成立. 经第一步, 可设 $n=p^{r}$; 经第二步, 可设 $p$ 为奇数. 设 $a$ 对几乎所有 $v$ 是 $K_v$ 中的 $n$ 次幂. 由于 $K[\zeta_{p^r}]$ 在 $K[\zeta_p]$ 上是 $p$ 幂阶的循环扩张, 第二步表明 $a$ 在 $K[\zeta_p]$ 中成为 $n$ 次幂, 比如 $a=b^{n}$. 取范数, 得知 $a^{d}$ 是 $K^{\times}$ 中的 $n$ 次幂, 其中 $d=[K[\zeta_p]:K]<p$. 但 $d$ 与 $p$ 互素, 故这蕴含 $a$ 是 $K^{\times}$ 中的 $n$ 次幂 (我们知道 $a^{d}=1$ 属于 $K^{\times}/K^{\times p^{r}}$, 而这蕴含 $a=1$ 属于 $K^{\times}/K^{\times p^{r}}$).
\end{proof}

\begin{remark}\label{rem:viii.1.5}
(a) 对 $K$ 的素的有限集 $S$ 与整数 $n$, 令 $P(n,S)$ 为所有满足如下条件的 $a\in K^{\times}$ 构成的群: 对所有 $v\notin S$, $a\in K_v^{\times n}$. 对 $2^{t}$ 情形作进一步分析可以发现: 要么 $P(n,S)=K^{\times n}$, 要么存在 $a_0\in K^{\times}$ 使 $P(n,S)=K^{\times n}\sqcup a_0K^{\times n}$. 例如, 对 $K=\Q$, $P(8,S)=\Q^{\times 8}\sqcup 16\,\Q^{\times 8}$ (\cite{ArtinTate1961}, p.~96; 2009 年版 p.~75). 注意 $P(n,S)/K^{\times n}$ 是使
\[
0\to P(n,S)/K^{\times n}\to K^{\times}/K^{\times n}\to\mathbb{I}_{K}/\mathbb{I}_{K}^{n}\to\mathbf{C}_{K}/\mathbf{C}_{K}^{n}\to0
\]
正合的群.

(b) 该定理的表述并非对所有 $K,S,n$ 都成立, 这一事实可视为小数域的一种病态. 例如, 若 $\sqrt{-1}\in K$, 则 $K[\zeta_{2^t}]$ 在 $K$ 上是循环的.

(c) 证明中的关键一步是第二步, 它有点微妙. Whaples 在证明 Grunwald ``定理''时正是在这里出错的: 他用 $a^{\frac{1}{n}}$ 表示 $X^{n}-a$ 的一个根, 却忘了不同的根相对于 $K$ 可能有十分不同的性质 (例如, 它们的极小多项式是诸 $f_i$, 甚至可能有不同的次数). Artin 与 Tate 明确警告不要试图缩短上述论证.
\end{remark}

\section{Grunwald--Wang 定理}\label{sec:VIII.2}

类域论通过把关于 Abel 扩张的问题转化为关于 idèle 类群中有限指数开子群的问题, ``解决''了所有这类问题. 遗憾的是, 关于 idèle 类群中有限指数开子群的问题也可能很难. 本节考虑如下问题.

\begin{question}\label{qu:viii.2.1}
设 $K$ 是代数数域, 并设对有限集 $S$ 中的每个素 $v$ 给定了 $K_v$ 的 $n_v$ 次有限循环扩张 $K^{v}$. 是否存在 $K$ 的 $n=\mathrm{lcm}(n_v)$ 次有限循环扩张 $L/K$, 使得 $L^{v}\approx K^{v}$? (此条件是指: 存在 $L$ 的整除 $v$ 的素 $w$ 以及 $K_v$-同构 $L_w\to K^{v}$.)
\end{question}

下面的例子表明答案并不总是肯定的, 然而我们将会看到, 答案通常是肯定的.

\begin{example}\label{ex:viii.2.2}
设 $\Q^{2}$ 是 $\Q_2$ 的 $8$ 次非分歧扩张. 容易构造 $\Q$ 的 $8$ 次扩张 $L$ 使得 $(2)$ 在 $L$ 中保持素, 从而对唯一的 $w|2$ 有 $L_w\approx\Q^{2}$: 取 $L=\Q[X]/(g(X))$, 其中 $g(X)$ 是 $\Z[X]$ 中任意模 $2$ 不可约的首一 $8$ 次多项式. 然而, 我断言不可能构造出这样的循环扩张. 更精确地说:

设 $L$ 是 $\Q$ 的 $8$ 次循环扩张, 则 $(2)$ 不可能在 $L$ 中保持素.

这由如下三个断言推出:

(i) $16$ 在所有 $p\ne 2$ 的 $\Q_p$ 中都是 $8$ 次幂 (包括 $p=\infty$), 但在 $\Q_2$ 中不是.

(ii) 设 $L/\Q$ 是有限 Abel 扩张; 若 $a\in\Q$ 对除一个之外的所有 $p$ 都成为 $\Q_p$ 中的范数, 则它对所有 $p$ 都是范数.

(iii) $16$ 不是来自 $\Q^{2}$ 的范数.

设 $L$ 是 $\Q$ 的 $8$ 次循环扩张. 由于 $L_v$ 在 $\Q_p$ 上的次数整除 $8$, (i) 表明 $16$ 对所有 $p\ne 2$ 都成为 $\Q_p$ 中的范数. 根据 (ii), 这蕴含 $16$ 是 $\Q_2$ 中的范数. 现在 (iii) 表明 $L^{2}$ 不同构于 $\Q^{2}$.

关于断言 (i) 的证明, 见习题\ref{exc:viii.1.3} (及其解答). 对 (ii), 考虑图表:
\begin{equation*}
\begin{tikzcd}[column sep=small, row sep=scriptsize]
& \Q_p^{\times} \ar[dr, "\phi_p"] & \\
\Q^{\times} \ar[ur, hook] \ar[r] & \mathbb{I}_{\Q} \ar[r, "\phi"] \ar[d, two heads] & \Gal(L^{p}/\Q_p) \\
& \Q_p^{\times} \ar[ur, two heads] &
\end{tikzcd}
\end{equation*}
其中 $\phi_p$ 是到整除 $p$ 的素处的分解群中的局部 Artin 映射, $\phi$ 是整体 Artin 映射 (我们允许 $p=\infty$, 此时 $\Q_p=\R$). 图表交换 (这正是我们定义整体 Artin 映射的方式). 对每个 $a\in\Q^{\times}$, 有乘积公式
\[
\prod_{p=2,3,5,\ldots,\infty}\phi_{p}(a)=1
\]
(VII, \ref{cor:vii.4.8}). 现在 $\phi_p(a)=1$ 当且仅当 $a$ 成为 $\Q_p$ 中的范数, 故乘积公式蕴含 (ii). 对 (iii), 我们利用: 由于 $\Q^{2}$ 在 $\Q_2$ 上非分歧, $\Q_2$ 上的正规赋值延拓为 $\Q^{2}$ 上的正规赋值. 若 $16=\Nm(\alpha)$, 则
\[
4=\mathrm{ord}_{2}(16)=\mathrm{ord}_{2}(\mathrm{Nm}(\alpha))=8\,\mathrm{ord}_{2}(\alpha),
\]
这不可能.
\end{example}

我现在简要解释一下正面的结果 (\cite{ArtinTate1961}, 第 9、10 章). 设 $S$ 是 $K$ 的素的有限集. 复合映射
\[
\prod_{v\in S}K_{v}^{\times}\to\mathbb{I}_{K}\to\mathbf{C}_{K}
\]
是单射, 并且当我们赋予 $P\overset{\text{def}}{=}\prod_{v\in S}K_v^{\times}$ 积拓扑时它是连续的. 当 $S$ 含有多于一个素时, 积拓扑与由 $\mathbf{C}_K$ 诱导的拓扑并不一致 (同上, p.~98). 尽管如此, $P$ 的一个有限指数子群对一种拓扑是开的当且仅当它对另一种拓扑是开的, 并且 $P$ 的有限指数开子群是 $P$ 与 $\mathbf{C}_K$ 中某个有限指数开子群的交 (同上, p.~99). 利用互反律, 这蕴含下述结果.

\begin{theorem}\label{thm:viii.2.3}
对每个 $v\in S$, 设 $N_v$ 是 $K_v^{\times}$ 中有限指数的开子群. 则存在 $K$ 的 Abel 扩张 $L$, 使得对所有 $v\in S$, $L^{v}$ 是 (由局部类域论) 对应于 $N_v$ 的 $K_v$ 的扩张.
\end{theorem}

更细致的分析证明下述结果 (其中 ``特征''指``连续特征'').

\begin{theorem}[GRUNWALD--WANG]\label{thm:viii.2.4}
对每个 $v\in S$, 设 $\chi_v$ 是 $K_v^{\times}$ 的特征. 则存在 $\mathbf{C}_K$ 的特征 $\chi$, 它在 $K_v^{\times}$ 上的限制为 $\chi_v$ 对所有 $v\in S$ 成立. 设 $n_v$ 为 $\chi_v$ 的阶, $n=\mathrm{lcm}(n_v)$. 则可以选取 $\chi$ 使其阶为 $n$, 但当 $2^{t}|n$ 对某个使 $K[\zeta_{2^{t}}]$ 在 $K$ 上不循环的 $t$ 成立时可能例外.
\end{theorem}

\begin{proof}
见 \cite{ArtinTate1961}, 第 X 章, 定理 5. 或者, 也可以通过比较 $S$ 与全体素集的 Poitou--Tate 定理来证明该定理的一个更精确的版本——见这些讲义的未来版本.
\end{proof}

例\ref{ex:viii.2.2} 表明, 若去掉对 $n$ 的条件, 第二个断言可能失效.

\begin{corollary}\label{cor:viii.2.5}
对每个 $v\in S$, 设 $n_v$ 为正整数. 若 $v$ 是无限素, 则 $n_v$ 必须是 $K_v$ 的某个扩张的可能次数. 则存在 $K$ 的 $n=\mathrm{lcm}(n_v)$ 次循环扩张 $L$, 使得局部次数对所有 $v\in S$ 为 $n_v$.
\end{corollary}

\begin{proof}
对无限素 $v\in S$, 以唯一可能的方式选取 $K_v^{\times}$ 上 $n_v$ 阶的特征 $\chi_v$; 对不整除 $2$ 的有限素 $v\in S$, 取 $\chi_v$ 为在 $U_v$ 上取值 $1$ 的 $n_v$ 阶典范特征; 对整除 $2$ 的 $v\in S$, 选取 $n_v$ 阶的 $\chi_v$ 以避开定理中的例外情形. 见同上 p.~105, 定理 5 (不要与\footnote{Serge Lang 记录了本讲义. 在 2009 年版中, 它已被正确编号为定理 6.}同页码 p.~103 上编号相同的定理混淆).
\end{proof}

\paragraph{应用}
设 $D$ 是代数数域 $K$ 上次数为 $d^{2}$ 的中心可除代数, 并设 $i_v$ 是 $D\otimes_K K_v$ 的类在不变映射 $\mathrm{inv}_v:\mathrm{Br}(K_v)\to\Q/\Z$ 下的像. 写 $i_v=m_v/n_v \bmod \Z$, 其中 $m_v,n_v$ 为互素整数且 $n_v>0$.

由于 $\mathrm{Br}(K_v)\to\Q/\Z$ 是单射 (当 $v$ 有限时甚至是同构), $D\otimes_K K_v$ 的类在 $\mathrm{Br}(K_v)$ 中的阶为 $n_v$. 在同构意义下, 恰好存在一个 $K_v$ 上不变量为 $i_v=m_v/n_v$ 的可除代数 $D_{i_v}$, 其次数为 $n_v^{2}$. 于是 $D\otimes_K K_v\approx M_{d/n_v}(D_{n_v})$. 特别地, $n_v|d$, 从而 $n|d$.

由于 $\mathrm{Br}(K)\to\bigoplus_v\mathrm{Br}(K_v)$ (对所有 $v$ 求和) 是单射, $D$ 的类在 $\mathrm{Br}(K)$ 中的阶为 $n=\mathrm{lcm}(n_v)$.

根据推论, 存在 $K$ 的 $n$ 次循环扩张 $L$, 其局部次数对每个满足 $i_v\notin\Z$ 的 $v$ 为 $n_v$. 于是 $L$ 分裂 $D$. 由 (IV, \ref{cor:iv.3.6}), 存在 $K$ 上的中心单代数 $B$ 使得

(a) $B\supset L$;

(b) $B$ 相似于 $D$, 即 $B\approx M_{r}(D)$ 对某个 $r$; 且

(c) $[B:K]=[L:K]^{2}$.

由后两条可见 $r^{2}d^{2}=n^{2}$. 由于 $n|d$, 故实际上 $r=1$, $n=d$. 我们证明了:

\begin{theorem}\label{thm:viii.2.6}
对数域 $K$ 上的每个可除代数 $D$, $D$ 在 $\mathrm{Br}(K)$ 中的阶为 $\sqrt{[D:K]}$, 并且 $D$ 包含 $K$ 的一个循环扩张作为极大子域 (从而被它分裂).
\end{theorem}

由该定理与 Noether--Skolem 定理可以推出: $D$ 是 Dickson 于 1906 年所定义意义下的循环代数, 即 $D$ 作为 $K$-代数由 $L$ 的元素与单个元素 $\beta$ 生成, 满足关系
\[
\beta\cdot\alpha=\sigma\alpha\cdot\beta,\quad \alpha\in L,\qquad \beta^{d}=\gamma,
\]
其中 $\sigma$ 生成 $\Gal(L/K)$, $\gamma\in K^{\times}$.

这些是非常著名的结果\footnote{``1930 年代代数结构理论的最高成就, 无疑是如下定理: 数域上每个有限维中心可除代数都是循环的, 以及用一组数值不变量对这些代数的分类. 后一结果等价于确定数域 Brauer 群的结构. 除了我们已概述的中心单代数的一般理论之外, 这些基本结果的证明还需要 Hasse 关于 $p$-adic 域上中心单代数的结构理论、Hasse 范数定理 (``Hasse 原理''), 以及关于某些循环扩张的 Grünwald 存在定理……数域上中心可除代数循环性的第一个证明由 Brauer、Hasse 与 Noether 给出 (1931).'' Nathan Jacobson, A.A.~Albert 讣告, \emph{Bull.~AMS}, 80 (1974), p.~1081.}, 有许多应用, 例如在 Abel 簇理论中的应用 (见我的讲义 AV, 定理 16.8).

\begin{NOTES}
1933 年, Grunwald 在不带对 $n$ 的条件的情况下证明了定理\ref{thm:viii.2.4}; 1942 年, Whaples 给出了第二个证明. 1948 年, Wang 发现了由 (\ref{exc:viii.1.3}, \ref{ex:viii.2.2}) 提供的反例, 并修正了表述. 在 1933 年之后的十五年间, Grunwald 的``定理''被相当广泛地使用, 例如用来证明数域上的每个可除代数都是循环的. 见 \cite{Roquette2005}, 5.3.
\end{NOTES}

事实上, 正如 Brian Conrad 向我指出的, Grunwald 定理的一个反例实际上早在 1934 年, 即比 Wang 早 14 年就已发表, 但 (显然) 没有人——包括作者本人——注意到这一点. 具体地说, 设 $a,n$ 为整数; 对不整除 $an$ 的素数 $p$, $a$ 是 $\Q_p$ 中的 $n$ 次幂当且仅当它模 $p$ 是 $n$ 次幂. Trost (1934)\footnote{E.~Trost, Zur Theorie der Potenzreste, Nieuw Arch. Wisk. 18 (1934), 58--61.}用初等方法证明了下述定理:

若 $a$ 对几乎所有素数都是 $n$ 次幂且 $8$ 不整除 $n$, 则 $a$ 是 $n$ 次幂; 当 $8$ 整除 $n$ 时, $a$ 要么是 $n$ 次幂, 要么是 $2^{n/2}$ 乘以一个 $n$ 次幂 (两种情形都会出现).

例如, 当 $n=8$ 时, 有可能 $2^{4}a$ 是 $8$ 次幂 (从而 $a$ 不是 $8$ 次幂). 这在 $a=16$ 时发生. 换言之, Trost 已经知道例\ref{ex:viii.2.2} 的断言 (i), 而由它通过标准论证即可导出 Grunwald 定理的反例.

顺带一提, Olga Taussky 审阅过 Grunwald 的论文. Chevalley 为 Mathematical Reviews 审阅过 Whaples 的论文, 甚至仔细到注意到``在阅读基本引理的证明时须提防若干排版错误'', 却没有仔细到发现主要结果是错误的.
\section{范数与二次型的局部--整体原则}\label{sec:VIII.3}

局部--整体 (或 Hasse) 原则问的是: 一个命题若在数域 $K$ 的每个完备化上都成立, 它是否就在 $K$ 上成立. 本节给出类域论使我们能够证明该原则成立的两个情形.

\subsection*{范数}

\begin{theorem}[HASSE 范数定理]\label{thm:viii.3.1}
设 $L/K$ 是数域的循环扩张, $a\in K^{\times}$. 则 $a$ 在 $K_v$ 中的像对除有限多个之外的所有 $v$ 都是来自 $L^{v}$ 的范数; 并且若它对所有 $v$ 都是范数, 则它是 $K$ 中的范数.
\end{theorem}

\begin{proof}
根据定理 (VII, \ref{thm:vii.5.1}), $H^{1}(G,\mathbf{C}_{L})=0$. 由于循环群上同调的周期性, 这蕴含 $H_{T}^{-1}(G,\mathbf{C}_{L})=0$. 因此, 从
\[
1\to L^{\times}\to\mathbb{I}_{L}\to\mathbf{C}_{L}\to0
\]
的上同调列得到
\[
H_{T}^{0}(G,L^{\times})\to H_{T}^{0}(G,\mathbb{I}_{L})
\]
是单射. 但 (见 (VII, \ref{prop:vii.2.5})) 这正是映射
\[
K^{\times}/\Nm(L^{\times})\to\bigoplus_{v}K_{v}^{\times}/\Nm(L^{v\times}).
\]
\end{proof}

\begin{remark}\label{rem:viii.3.2}
该证明对非循环扩张失效, 并且事实上该命题对非循环扩张并不成立. 例如, $5^{2}$ 在所有素处都是来自 $\Q[\sqrt{13},\sqrt{17}]$ 的局部范数, 但不是整体范数 (见 \cite{CasselsFrohlich1967}, 习题 5.3, p.~360). 对给定的非循环 Abel 群 $G$, 下述论文研究以 $G$ 为 Galois 群且满足 Hasse 范数原理的扩张的密度: Frei, C.; Loughran, D.; Newton, R.: The Hasse norm principle for Abelian extensions. \emph{Amer.~J.~Math.} 140 (2018), no.~6, 1639--1685.
\end{remark}

\subsection*{二次型}

回顾: 域 $k$ 上向量空间 $V$ 上的二次型是指映射 $Q\colon V\to k$ 使得

(a) $Q(av)=a^{2}Q(v)$;

(b) $B(v,w)\overset{\text{def}}{=}Q(v+w)-Q(v)-Q(w)$ 是 $V$ 上的双线性型.\footnote{当 $k$ 的特征 $\ne 2$ 时, 通常令 $B(v,w)=\frac{1}{2}(Q(v+w)-Q(v)-Q(w))$.}

二次型 $Q$ 称为\textbf{非退化的} (nondegenerate), 若其伴随双线性型 $B$ 非退化. 设 $c\in k$. 非退化二次型 $Q$ 称为\textbf{表示} (represent) $c$, 若存在非零向量 $v\in V$ 使 $Q(v)=c$.

\begin{lemma}\label{lem:viii.3.3}
若非退化二次型 $Q$ 表示 $0$, 则它表示所有 $c\in k$.
\end{lemma}

\begin{proof}
注意, 对 $t\in k$,
\[
Q(tv+w)=t^{2}Q(v)+tB(v,w)+Q(w).
\]
设 $v_0$ 是满足 $Q(v_0)=0$ 的非零向量, 则因 $B$ 非退化, 存在向量 $w_0$ 使 $B(v_0,w_0)\ne 0$. 当 $t$ 取遍 $k$ 的所有值时, $Q(tv_0+w_0)=tB(v_0,w_0)+Q(w_0)$ 也取遍 $k$ 的所有值.
\end{proof}

当 $k$ 的特征 $\ne 2$ 时, 存在 $V$ 的一组基 $\{e_1,\ldots,e_n\}$ 使 $B(e_i,e_j)=0$ 对 $i\ne j$ 成立. 此时
\[
Q\Big(\sum x_ie_i\Big)=\sum_{i=1}^{n}a_ix_i^{2},\qquad n=\dim V.
\]
以下我们将二次型写作
\[
q(X_{1},\ldots,X_{n})=a_{1}X_{1}^{2}+\cdots+a_{n}X_{n}^{2},
\]
并记住: 可逆的变量代换不会改变我们的任何结论.

\begin{lemma}\label{lem:viii.3.4}
非退化二次型 $q(X_{1},\ldots,X_{n})$ 表示 $a$ 当且仅当 $r\overset{\text{def}}{=}q-aY^{2}$ 表示 $0$.
\end{lemma}

\begin{proof}
若 $q(x_1,\ldots,x_n)=a$, 则 $r(x_1,\ldots,x_n,1)=0$. 反之, 设 $r(x_1,\ldots,x_n,y)=0$. 若 $y=0$, 则 $q$ 表示 $0$, 从而表示 $k$ 的每个元素. 若 $y\ne 0$, 则 $q\big(\frac{x_1}{y},\ldots,\frac{x_n}{y}\big)=q(x_1,\ldots,x_n)/y^{2}=a$.
\end{proof}

\begin{theorem}\label{thm:viii.3.5}
设 $q$ 是系数在数域 $K$ 中的 $n$ 元非退化二次型.

(a) 若 $n\ge 3$, 则 $q$ 对除有限多个之外的所有 $v$ 在 $K_v$ 中表示 $0$.

(b) $q$ 在 $K$ 中表示 $0$ 当且仅当它对所有 $v$ 在 $K_v$ 中表示 $0$.
\end{theorem}

在开始证明之前, 我们指出一个推论.

\begin{corollary}\label{cor:viii.3.6}
设 $c\in K$. 系数在 $K$ 中的非退化二次型 $q$ 在 $K$ 中表示 $c$ 当且仅当它对所有 $v$ 在 $K_v$ 中表示 $c$.
\end{corollary}

\begin{proof}
令 $r=q-cY^{2}$. 则 $r$ 表示 $0$ 当且仅当 $q$ 表示 $c$.
\end{proof}

我们从一个纯代数的结果开始证明.

\begin{proposition}\label{prop:viii.3.7}
设 $k$ 是特征 $\ne 2$ 的域.

(a) 型 $q=X^{2}$ 不表示 $0$.

(b) 型 $q=X^{2}-aY^{2}$ 表示 $0$ 当且仅当 $a$ 是平方.

(c) 型 $q=X^{2}-aY^{2}-bZ^{2}$ 表示 $0$ 当且仅当 $b$ 是来自域 $k[\sqrt{a}]$ 的范数.

(d) 型 $q=X^{2}-bY^{2}-cZ^{2}+acT^{2}$ 在 $k$ 中表示 $0$ 当且仅当 $c$ 作为 $k[\sqrt{ab}]$ 的元素是来自 $k[\sqrt{a},\sqrt{b}]$ 的范数.
\end{proposition}

\begin{proof}
(a) 显然.

(b) 根据引理\ref{lem:viii.3.4}, $X^{2}-aY^{2}$ 表示 $0$ 当且仅当 $X^{2}$ 表示 $a$.

(c) 根据引理\ref{lem:viii.3.4}, $X^{2}-aY^{2}-bZ^{2}$ 表示 $0$ 当且仅当 $X^{2}-aY^{2}$ 表示 $b$, 即 $b$ 是来自 $k[\sqrt{a}]$ 的范数.

(d) 显然,
\[
q(x,y,z,t)=0\iff c=\frac{\Nm_{k[\sqrt{b}]/k}(x+\sqrt{b}\,y)}{\Nm_{k[\sqrt{a}]/k}(z+\sqrt{a}\,t)}.
\]
由于范数的逆仍是范数, 这表明 $q$ 表示零当且仅当 $c$ 是一个来自 $k[\sqrt{a}]$ 的范数与一个来自 $k[\sqrt{b}]$ 的范数之积. 于是 (d) 由下述引理推出.
\end{proof}

\begin{lemma}\label{lem:viii.3.8}
设 $k$ 是特征 $\ne 2$ 的域. 元素 $c\in k^{\times}$ 是一个来自 $k[\sqrt{a}]$ 的范数与一个来自 $k[\sqrt{b}]$ 的范数之积, 当且仅当它作为 $k[\sqrt{ab}]$ 的元素是来自 $L=k[\sqrt{a},\sqrt{b}]$ 的范数.
\end{lemma}

\begin{proof}
我们把 $a$、$b$ 或 $ab$ 是 $k$ 中平方的退化情形留给读者. 于是可设 $\Gal(k[\sqrt{a},\sqrt{b}]/k)=\{1,\sigma,\tau,\sigma\tau\}$, 其中 $\sigma$、$\tau$、$\sigma\tau$ 的阶都是 $2$, 并分别固定 $\sqrt{a}$、$\sqrt{b}$、$\sqrt{ab}$. 第一个条件断言:

(*) 存在 $x,y\in k[\sqrt{a},\sqrt{b}]$ 使得 $\sigma x=x$, $\tau y=y$, 且 $xy\cdot\sigma\tau(xy)=c$;

第二个条件断言:

(**) 存在 $z\in k[\sqrt{a},\sqrt{b}]$ 使得 $z\cdot\sigma\tau(z)=c$.

显然 (*) $\Rightarrow$ (**). 反之, 注意
\[
z\cdot\sigma z=\Nm_{k[\sqrt{a},\sqrt{b}]/k[\sqrt{a}]}z\in k[\sqrt{a}].
\]
再者,
\[
\Nm_{k[\sqrt{a}]/k}(z\cdot\sigma z)=z\cdot\sigma z\cdot\tau z\cdot\sigma\tau z\in k.
\]
由于 $z\cdot\sigma\tau z=c\in k$, 这蕴含 $\sigma z\cdot\tau z\in k$, 故
\[
\sigma z\cdot\tau z=\sigma(\sigma z\cdot\tau z)=z\cdot\sigma\tau z=c.
\]
因此
\[
\Nm_{k[\sqrt{a}]/k}(z\cdot\sigma z)=c^{2}.
\]
现在把 Hilbert 定理 90 ((II, \ref{cor:ii.1.23})) 应用于 $z\cdot\sigma z/c\in k[\sqrt{a}]$, 得知存在 $x\in k[\sqrt{a}]^{\times}$ 使 $\tau x/x=z\cdot\sigma z/c$. 令 $y=\sigma\tau z/x$. 则
\[
\tau y=\sigma z/\tau x=c/(z\cdot x)=z\cdot\sigma\tau z/(z\cdot x)=y
\]
(依次利用: $y$ 的定义; $x$ 的定义; $z$ 的定义; $y$ 的定义), 并且
\[
xy\cdot\sigma\tau(xy)=\sigma\tau z\cdot\sigma\tau(\sigma\tau z)=\sigma\tau z\cdot z=c
\]
(依次利用: $y$ 的定义; $(\sigma\tau)^{2}=1$; $z$ 的定义), 如所欲证.
\end{proof}

\begin{proof}[定理\ref{thm:viii.3.5}(a) 的证明]
若 $q=q_1(X_1,\ldots,X_m)+q_2(X_{m+1},\ldots,X_n)$ 且 $q_1$ 表示零, 则 $q$ 也表示零. 因此只需对三元二次型证明 (\ref{thm:viii.3.5}a). 用非零纯量乘 $q$ 之后, 可设 $q=X^{2}-aY^{2}-bZ^{2}$, 而对此种二次型, 结论由 (\ref{prop:viii.3.7}) 与定理 \ref{thm:viii.3.1} 推出.
\end{proof}

\begin{proof}[定理\ref{thm:viii.3.5}(b) 的证明]
我们对变量个数 $n$ 作归纳证明.

当 $n=1$ 时, 无需证明, 因为假设永远不会成立.

当 $n=2$ 时, 用非零纯量乘 $q$ 之后可设 $q=X^{2}-aY^{2}$, 对此种二次型, 结论由 (\ref{prop:viii.3.7}) 与定理\ref{thm:viii.1.1} 推出.

当 $n=3,4$ 时, 结论以类似方式由 (\ref{prop:viii.3.7}c,d) 与定理\ref{thm:viii.1.1} 推出.

在证明一般情形之前, 我们作一些初等观察.

(a) 非退化二次型 $q_1(X_1,\ldots,X_m)-q_2(X_{m+1},\ldots,X_n)$ 在域 $k$ 中表示 $0$ 当且仅当存在 $c\in k$ 使 $q_1$ 与 $q_2$ 都表示 $c$.

(b) 若 $q$ 在 $k^\times$ 中表示 $c$, 则 $q$ 表示傍系 $c\cdot k^{\times 2}$ 中的每个元素.

(c) 子群 $K_v^{\times 2}\subset K_v^{\times}$ 是开的. 当 $v$ 为实或复时显然. 当 $v$ 非阿基米德时, Newton 引理 (ANT, 7.32) 表明, 对任意满足 $|1-a|_{v}<|2|_{v}^{2}$ 的 $a$, $1$ 可以精化为 $X^{2}-a$ 的根.

结合 (b) 与 (c) 可见: 系数在 $K_v$ 中的二次型 $q$ 表示 $K_v^{\times}$ 的一个非空开子集中的元素.

现在设 $n\ge 5$, 并设定理\ref{thm:viii.3.5}(b) 已对 $n-1$ 获证. 令
\[
q(X_{1},\ldots,X_{n})=aX_{1}^{2}+bX_{2}^{2}-r(X_{3},\ldots,X_{n}),\qquad n-2\ge 3.
\]
由定理的 (a) 知: 除某个有限集 $S$ 中的 $v$ 之外, $r$ 在 $K_v$ 中表示 $0$. 设 $v\in S$. 由于 $q$ 在 $K_v$ 中表示 $0$, 存在同时被 $aX_1^{2}+bX_2^{2}$ 与 $r$ 表示的元素 $c_v\in K_v^{\times}$, 即存在 $x_i(v)\in K_v$ 使
\[
ax_{1}(v)^{2}+bx_{2}(v)^{2}=c_v=r(x_{3}(v),\ldots,x_{n}(v)).
\]
现在应用弱逼近定理 (ANT, 7.20), 找到对所有 $v\in S$ 都在 $K_v$ 中接近 $x_1(v)$、$x_2(v)$ 的元素 $x_1,x_2\in K$. 则
\[
c\overset{\text{def}}{=}ax_{1}^{2}+bx_{2}^{2}
\]
将对每个 $v\in S$ 接近 $c_v$; 特别地, 可设 $c/c_v\in K_v^{\times 2}$ 对所有 $v\in S$ 成立.

考虑二次型 $cY^{2}-r$. 它在 $K_v$ 中表示零, 对 $v\notin S$: 因为 $r$ 在 $K_v$ 中表示零; 对 $v\in S$: 因为 $r$ 在 $K_v$ 中表示 $c$. 由归纳, $cY^{2}-r$ 在 $K$ 中表示零. 由此得到 $q$ 在 $K$ 中表示 $0$, 因为 $aX_1^{2}+bX_2^{2}$ 与 $r$ 都在 $K$ 中表示 $c$.
\end{proof}

\begin{proposition}\label{prop:viii.3.9}
$\Q_p$ 的有限扩张 $K$ 上的四元非退化二次型 $q$ 表示 $K^{\times}$ 的每个非零元素.
\end{proposition}

\begin{proof}
若 $q$ 表示 $0$, 则它表示 $K$ 的每个元素. 我们假设相反的情形. 用 $K$ 的非零元素乘 $q$ 之后, 可设
\[
q=X^{2}-bY^{2}-cZ^{2}+acT^{2}.
\]
由于 $q$ 在 $K$ 中不表示 $0$, $b$ 与 $a$ 都不是平方.

若 $K[\sqrt{a}]\ne K[\sqrt{b}]$, 则 (由局部类域论定理 (I, \ref{thm:1.1})) $\Nm(K[\sqrt{a}]^{\times})$ 与 $\Nm(K[\sqrt{b}]^{\times})$ 是 $K^{\times}$ 中两个不同的指数 $2$ 子群, 从而 $K^{\times}=\Nm(K[\sqrt{a}]^{\times})\cdot\Nm(K[\sqrt{b}]^{\times})$. 由于范数的逆仍是范数, 这意味着可以把 $c$ 写成
\[
c=\frac{x^{2}-by^{2}}{z^{2}-at^{2}},
\]
其中 $x,y,z,t\in K$. 展开相乘后发现 $q$ 表示 $0$, 与假设矛盾. 因此 $K[\sqrt{a}]=K[\sqrt{b}]$, 且 $a=b\times(\text{平方})$ (见 (VII, \ref{prop:vii.A.1})). 该平方可被吸收进 $T^{2}$, 故可写
\[
q=X^{2}-bY^{2}-cZ^{2}+bcT^{2}.
\]
考虑四元数代数 $H(b,c)$ (见习题 IV.\ref{ex:iv.5.7}). 对
\[
\alpha=x+yi+zj+tk
\]
定义
\[
\bar{\alpha}=x-yi-zj-tk,
\]
于是
\[
\Nm(\alpha)\overset{\text{def}}{=}\alpha\bar{\alpha}=x^{2}-by^{2}-cz^{2}+bct^{2}.
\]
映射 $\alpha\mapsto\Nm(\alpha):H(b,c)^{\times}\to K^{\times}$ 是同态, 我们必须证明它是满射.

对每个 $\alpha\in H(b,c)$,
\[
P_{\alpha}(X)\overset{\text{def}}{=}(X-\alpha)(X-\bar{\alpha})=X^{2}-(\alpha+\bar{\alpha})X+\Nm(\alpha)\in K[X].
\]
因此 $P_{\alpha}(X)$ 是 $\alpha$ 在扩张 $K[\alpha]/K$ 中的特征多项式. 特别地,
\[
\Nm(\alpha)=\Nm_{K[\alpha]/K}(\alpha).
\]
现在 (见 (IV, \ref{rem:iv.4.4})) $H(b,c)$ 包含 $K$ 的每个二次扩张的同构副本, 例如 $K$ 的非分歧二次扩张与一个全分歧二次扩张. 因此 $\Nm(H(b,c)^{\times})$ 包含这两个互异二次扩张的范数群, 从而 (如上) 等于 $K^{\times}$.
\end{proof}

\begin{corollary}\label{cor:viii.3.10}
$\Q_p$ 的有限扩张上 $\ge 5$ 元的每个非退化二次型都表示 $0$.
\end{corollary}

\begin{proof}
可写 $q=r(X_1,\ldots,X_4)-aX_5^{2}+q'(X_6,\ldots)$, 其中 $r$ 是四元非退化二次型, $a\ne 0$. 则 $r$ 表示 $a$, 故 $q$ 表示 $0$.
\end{proof}

\begin{corollary}\label{cor:viii.3.11}
数域 $K$ 上 $\ge 5$ 元的非退化二次型 $q$ 表示 $0$ 当且仅当它在 $K$ 的每个实完备化中表示 $0$.
\end{corollary}

\begin{proof}
结合 (\ref{cor:viii.3.10}) 与 (\ref{thm:viii.3.5}).
\end{proof}

\begin{remark}\label{rem:viii.3.12}
命题\ref{prop:viii.3.9} 的证明对 $K=\R$ 也成立, 直到最后一步: $\R$ 的唯一的二次扩张是 $\CC$, 故
\[
\Nm(\mathbb{H}(b,c)^{\times})=\Nm(\CC^{\times})=\R_{>0}.
\]
因此它表明: $\R$ 上不表示零的四元非退化型表示所有严格正的实数.
\end{remark}
\section{基本正合列与基本类}\label{sec:VIII.4}

对 Galois 扩张 $L/K$, 记
\[
\mathrm{Br}(L/K)=H^{2}(\Gal(L/K),L^{\times});
\]
对数域的 Galois 扩张 $L/K$, 记
\[
H^{2}(L/K)=H^{2}(\Gal(L/K),\mathbf{C}_{L}).
\]
由于 $H^{1}(G,L^{\times})=0$ (见 (II, \ref{prop:ii.1.22})) 且 $H^{1}(G,\mathbf{C}_{L})=0$ (见 (VII, \ref{thm:vii.5.1})), 对每个 Galois 扩张塔 $E\supset L\supset K$, 我们得到正合列
\[
0\to\mathrm{Br}(L/K)\to\mathrm{Br}(E/K)\to\mathrm{Br}(E/L)
\]
与
\[
0\to H^{2}(L/K)\to H^{2}(E/K)\to H^{2}(E/L).
\]
对越来越大的域 $E\subset K^{\mathrm{al}}$ 取正向极限, 得到正合列
\[
0\to\mathrm{Br}(L/K)\to\mathrm{Br}(K)\to\mathrm{Br}(L)
\]
与
\[
0\to H^{2}(L/K)\to H^{2}(K^{\mathrm{al}}/K)\to H^{2}(K^{\mathrm{al}}/L).
\]
于是, 我们可以把 $\mathrm{Br}(L/K)$ 视为 $\mathrm{Br}(K)$ 中被 $L$ 分裂的元素构成的子群; 对 $H^{2}(L/K)$ 亦然.

设 $L/K$ 是 $n$ 次有限 Galois 的数域扩张, 并考虑图表:
\begin{equation*}
\begin{tikzcd}[column sep=1.2em, row sep=large]
& & & H^{2}(L/K) \\
0 \ar[r] & \mathrm{Br}(L/K) \ar[r] & \bigoplus_{v}\mathrm{Br}(L^{v}/K_{v}) \ar[ur] \ar[dr, "\varSigma"'] & \\
& & & \Q/\Z
\end{tikzcd}
\end{equation*}
上 行是
\[
0\to L^{\times}\to\mathbb{I}_{L}\to\mathbf{C}_{L}\to0
\]
的上同调列的一部分. 左端的零来自 $H^{1}(G,\mathbf{C}_L)=0$. 上 行正合, 但映射 $\bigoplus_v\mathrm{Br}(L^{v}/K_v)\to H^{2}(L/K)$ 一般不是满射——我们把它记为 $H^{2}(L/K)'$.

回顾 (III, \ref{thm:iii.2.1}): 对每个素 $v$, 有同态
\[
\mathrm{inv}_{v}:\mathrm{Br}(K_{v})\to\Q/\Z.
\]
若 $v$ 非阿基米德, 它是 $\mathrm{Br}(K_v)$ 到 $\Q/\Z$ 的同构; 若 $v$ 是实的, 它是 $\mathrm{Br}(K_v)$ 到 $\frac{1}{2}\Z/\Z$ 的同构. 此外, 若 $L_w/K_v$ 的次数为 $n_v$, 则 $\mathrm{inv}_w(\beta)=n_v\cdot\mathrm{inv}_v(\beta)$, 故 $\mathrm{inv}_v$ 定义同构
\[
\mathrm{inv}_{v}:\mathrm{Br}(L_{w}/K_{v})\to\tfrac{1}{n_{v}}\Z/\Z.
\]
图表中的东南向箭头是
\[
\varSigma:\bigoplus_{v}\mathrm{Br}(L^{v}/K_{v})\to\Q/\Z,\qquad(\beta_{v})\mapsto\sum\mathrm{inv}_{v}(\beta_{v}).
\]
$\mathrm{inv}_v$ 的像是循环群 $\frac{1}{n}\Z/\Z$ 中 $n_v$ 阶的循环子群, 故 $\varSigma$ 的像是循环子群 $\frac{1}{n_0}\Z/\Z$, 其中 $n_0=\mathrm{lcm}(n_v)$.

根据定理 (VII, \ref{thm:vii.5.1}), $H^{2}(L/K)$ 的阶整除 $n$. 根据定理 (VII, \ref{thm:vii.8.1}), 图表的下 行是复形, 故图表中的映射诱导满同态
\[
H^{2}(L/K)^{\prime}\to\tfrac{1}{n_{0}}\Z/\Z.
\]
局部次数的最小公倍数 $n_0$ 总是整除 $n$, 但不必等于 $n$ (见下面的例\ref{ex:viii.4.5}). 然而, 假设扩张 $L/K$ 具有性质 $n=n_0$. 则:

(a) 映射 $H^{2}(L/K)'\to\frac{1}{n_0}\Z/\Z$ 是同构. (它是满射, 且 $(H^{2}(L/K)':1)\le(H^{2}(L/K):1)\le n$.)

(b) $H^{2}(L/K)'=H^{2}(L/K)$, 且两者的阶都是 $n$.

(c) 下 行是正合列
\[
0\to\mathrm{Br}(L/K)\to\bigoplus_{v}\mathrm{Br}(L^{v}/K_{v})\xrightarrow{\ \varSigma\ }\tfrac{1}{n}\Z/\Z\to0
\]
(因为它同构于上 行).

\begin{lemma}\label{lem:viii.4.1}
若 $L/K$ 是循环的, 则 $n=n_0$.
\end{lemma}

\begin{proof}
设 $S\supset S_\infty$ 是包含所有在 $L$ 中分歧的素的 $K$ 的素集. 对 $v\notin S$, $(p_v,L/K)$ 是 $\Gal(L/K)$ 中 $n_v$ 阶 ($=f_v$) 的元素, 故 Artin 映射 $I^{S}\to\Gal(L/K)$ 的像的阶为 $n_0\overset{\text{def}}{=}\mathrm{lcm}(n_v)$. 根据互反律 (VII, \ref{cor:vii.4.8}), Artin 映射是满射, 这蕴含 $n_0=n$. [利用复分析可以证明更多: 对密度为 $\varphi(n)/n$ 的一组 $v$, $L^{v}/K_v$ 是 $n$ 阶循环的: 设 $\mathfrak{m}$ 为 $L/K$ 的模, 并取 $\mathfrak{a}\in I^{S}$ 使 $(\mathfrak{a},L/K)$ 生成 $\Gal(L/K)$; 则在 $C_m$ 中同余于 $\mathfrak{a}$ 的素理想集的密度为 $1/n$.]
\end{proof}

设 $\Q^{\mathrm{c}}$ 是 (I, 例 \ref{ex:A.5}(d)) 中定义的 $\Q$ 的无限循环分圆扩张 (另见 (VII, \ref{lem:vii.7.3})), 并令 $\Omega=\Q^{\mathrm{c}}\cdot K$. 对每个 $n$, $\Omega$ 包含 $K$ 的唯一的 $n$ 次循环扩张 $\Omega_n$. 前面的引理与说明表明
\[
0\to\mathrm{Br}(\Omega_{n}/K)\to\bigoplus_{v}\mathrm{Br}(\Omega_{n}^{v}/K_{v})\xrightarrow{\ \varSigma\ }\tfrac{1}{n}\Z/\Z\to0
\]
正合. 对所有 $n$ 取正向极限 (实际上是定向并), 得到正合列
\[
0\to\mathrm{Br}(\Omega/K)\to\bigoplus_{v}\mathrm{Br}(\Omega^{v}/K_{v})\to\Q/\Z\to0.
\]

\begin{theorem}\label{thm:viii.4.2}
对每个数域 $K$, 序列
\[
0\to\mathrm{Br}(K)\to\bigoplus_{v}\mathrm{Br}(K_{v})\xrightarrow{\ \varSigma\ }\Q/\Z\to0
\]
正合.
\end{theorem}

\begin{proof}
根据命题 (VII, \ref{prop:vii.7.2}), $\mathrm{Br}(\Omega/K)=\mathrm{Br}(K)$. 而且 $\mathrm{Br}(\Omega^{v}/K_v)=\mathrm{Br}(K_v)$, 因为在非阿基米德情形, 当 $n\to\infty$ 时 $[\Omega_n^{v}:K_v]\to\infty$ (见 (VII, \ref{lem:vii.7.3})).
\end{proof}

定理中的序列称为整体类域论的\textbf{基本正合列}\index{jibenzhenghelie@基本正合列(fundamental exact sequence)}.

\begin{corollary}\label{cor:viii.4.3}
对每个有限扩张 $L/K$, 序列
\[
0\to\mathrm{Br}(L/K)\to\bigoplus_{v}\mathrm{Br}(L^{v}/K_{v})\xrightarrow{\ \varSigma\ }\tfrac{1}{n_{0}}\Z/\Z\to0,\qquad n_{0}=\mathrm{lcm}(n_{v}),
\]
正合.
\end{corollary}

\begin{proof}
把 $K$ 的基本正合列映到 $L$ 的基本正合列, 对所得图表应用蛇形引理.
\end{proof}

\begin{example}\label{ex:viii.4.4}
(a) 对数域的有限循环扩张 $L/K$, 基本正合列可等同于
\[
0\to K^{\times}/\Nm(L^{\times})\to\bigoplus_{v}K_{v}^{\times}/\Nm(L^{v\times})\to\tfrac{1}{n}\Z/\Z\to0,\qquad n=[L:K].
\]

(b) 设 $D$ 是数域 $K$ 上的可除代数, 并设 $i_v=\mathrm{inv}_v(D\otimes_K K_v)$. 则: $i_v=0$ 对除有限多个之外的所有 $v$; $i_v=0$ 若 $v$ 为复; $i_v\in\frac{1}{2}\Z/\Z$ 若 $v$ 为实; 且 $\sum i_v\equiv 0 \bmod \Z$. 族 $(i_v)$ 确定 $D$ 的同构类, 而任何满足这些条件的族 $(i_v)$ 都是某个可除代数的不变量族. 显然, $D$ 的类在 $\mathrm{Br}(K)$ 中的阶是诸 $i_v$ 的最小公分母 $n$. 还可以证明 $[D:K]=n^{2}$. 例如, 给出 $\Q$ 上的一个四元数代数就等价于给出 $\Q$ 的一个由偶数个元素构成的有限素集.
\end{example}

\begin{example}\label{ex:viii.4.5}
设 $L=\Q[\sqrt{13},\sqrt{17}]$. 显然 $n=4$, 但我断言对所有 $v$ 有 $n_v=1$ 或 $2$. 由于 $13$ 与 $17$ 都同余于 $1$ 模 $4$, $2$ 在 $L$ 中不分歧. 因此, 对 $w\mid p$, $p\ne 13,17$, $L_w$ 是 $\Q_p$ 的非分歧扩张. 特别地, 其 Galois 群是循环的. 由于它是 $\Gal(L/\Q)$ 的子群, 它被 $2$ 杀死, 从而阶为 $1$ 或 $2$. 另一方面, $\left(\frac{17}{13}\right)=1$ (显然) 且 $\left(\frac{13}{17}\right)=\left(\frac{17}{13}\right)=1$. 故 $17$ 模 $13$ 是平方, 而 Hensel 引理表明它在 $\Q_{13}$ 中是平方. 同样, $13$ 在 $\Q_{17}$ 中是平方.
\end{example}

\subsection*{基本类}

由以上讨论可知, 存在同构
\[
\mathrm{inv}_{K}:H^{2}(\Omega/K)\to\Q/\Z
\]
由如下性质唯一刻画: 复合
\[
\bigoplus_{v}\mathrm{Br}(K_{v})\to H^{2}(\varOmega/K)\xrightarrow{\ \mathrm{inv}_{K}\ }\Q/\Z
\]
是 $(\beta_{v})\mapsto\sum\mathrm{inv}_{v}(\beta_{v})$.

\begin{lemma}\label{lem:viii.4.6}
对数域的每个有限扩张 $L/K$ 与每个 $\gamma\in H^{2}(\Omega/K)$, 有 $\mathrm{inv}_L(\gamma)=n\,\mathrm{inv}_K(\gamma)$, 其中 $n=[L:K]$.
\end{lemma}

\begin{proof}
利用局部次数之和等于整体次数这一事实.
\end{proof}

因此, 对每个有限 Galois 的 $L/K$, 得到同构
\[
\mathrm{inv}_{L/K}:H^{2}(L/K)\to\tfrac{1}{n}\Z/\Z.
\]
对所有 $L\subset K^{\mathrm{al}}$ 取正向极限, 得到同构
\[
\mathrm{inv}_{K}:H^{2}(K^{\mathrm{al}}/K)\to\Q/\Z.
\]

\begin{theorem}\label{thm:viii.4.7}
对数域的每个有限 Galois 扩张 $L/K$, $H^{2}(L/K)$ 是 $n=[L:K]$ 阶的循环群, 带有典范生成元 $u_{L/K}$.
\end{theorem}

\begin{proof}
取 $u_{L/K}$ 为满足 $\mathrm{inv}_{L/K}(u_{L/K})=\frac{1}{n}\bmod\Z$ 的元素.
\end{proof}

$H^{2}(L/K)$ 的生成元 $u_{L/K}$ 称为\textbf{基本类}\index{jibenlei@基本类(fundamental class)}. 可以像局部情形 (见 (III, \ref{lem:iii.2.7})) 一样证明: 对每个由有限 Galois 扩张组成的塔 $E\supset L\supset K$,
\[
\begin{aligned}
\Res(u_{E/K})&=u_{E/L},\\
\Inf(u_{L/K})&=[E:L]\,u_{E/K}.
\end{aligned}
\]
因此, 可以应用 Tate 定理 ((II, \ref{thm:ii.3.11})) 得到同构
\[
\Gal(L/K)^{\mathrm{ab}}\to\mathbf{C}_{K}/\Nm(\mathbf{C}_{L}).
\]
该同构与上一章定义的整体 Artin 映射 $\phi_{L/K}$ 互逆, 这一点由整体基本类与局部基本类相容这一事实推出.

\subsection*{范数限制定理}

\begin{theorem}\label{thm:viii.4.8}
设 $E$ 是 $K$ 的有限扩张 (不必 Galois), 并设 $M$ 是 $E$ 的使 $M/K$ 为 Abel Galois 扩张的最大子扩张. 则
\[
\Nm_{E/K}\mathbf{C}_{E}=\Nm_{M/K}\mathbf{C}_{M}.
\]
\end{theorem}

\begin{proof}
设 $L$ 是 $K$ 的包含 $E$ 的 Galois 扩张, 并令 $G=\Gal(L/K)$, $H=\Gal(L/E)$. 考虑交换图表
\[
\begin{tikzcd}[column sep=large]
H_{T}^{-2}(H,\Z) \ar[r, "\simeq"] \ar[d, "\mathrm{Cor}"'] & H_{T}^{0}(H,\mathbf{C}_{L}) \ar[d, "\mathrm{Cor}"] \\
H_{T}^{-2}(G,\Z) \ar[r, "\simeq"'] & H_{T}^{0}(G,\mathbf{C}_{L})
\end{tikzcd}
\]
其中的水平箭头是与基本类的杯积. 它可以等同于交换图表:
\[
\begin{tikzcd}[column sep=large]
H^{\mathrm{ab}} \ar[r, "\simeq"] \ar[d] & C_{E}/\Nm_{L/E}C_{L} \ar[d] \\
G^{\mathrm{ab}} \ar[r, "\simeq"'] & C_{K}/\Nm_{L/K}C_{L}.
\end{tikzcd}
\]
故 $H^{\mathrm{ab}}\to G^{\mathrm{ab}}$ 的余核同构于 $C_K/\Nm_{E/K}(C_E)$. 但该余核等于 $\Gal(M/K)$, 后者又同构于 $C_K/\Nm_{M/K}(C_M)$. 由于 $\Nm(C_M)\supset\Nm(C_E)$, 这两个群必相等.
\end{proof}
\section{高次互反律}\label{sec:VIII.5}

回顾: 对奇素数 $p$ 与不被 $p$ 整除的整数 $a$, Legendre 符号 (或二次剩余符号)
\[
\left(\frac{a}{p}\right)\overset{\text{def}}{=}
\begin{cases}
1 & \text{若 } a \text{ 模 } p \text{ 是平方},\\
-1 & \text{否则}
\end{cases}
\]
的定义. 群 $\F_p^{\times}$ 是 $p-1$ 阶循环群, 以 $-1$ 为其唯一的 $2$ 阶元素. 因此, 对 $u\in\F_p^{\times}$, $u^{\frac{p-1}{2}}$ 依 $u$ 是否为平方而等于 $1$ 或 $-1$, 故 $\left(\frac{a}{p}\right)$ 是满足
\[
\left(\frac{a}{p}\right)\equiv a^{\frac{p-1}{2}}\pmod p
\]
的 $1$ 的唯一平方根.

二次互反律说: 对奇素数 $p,q$,
\[
\left(\frac{p}{q}\right)\left(\frac{q}{p}\right)=(-1)^{\frac{p-1}{2}\cdot\frac{q-1}{2}}.
\]
二次互反律的补充说
\[
\left(\frac{-1}{p}\right)=(-1)^{\frac{p-1}{2}},\qquad
\left(\frac{2}{p}\right)=(-1)^{\frac{n^{2}-1}{8}}.
\]

对 Gauss 整数 $\alpha$ (即 $\Z[i]$ 的元素) 与奇 Gauss 素数 $\pi$ (即 $\Z[i]$ 中不整除 $2$ 的素元), Gauss 定义 $\left(\frac{\alpha}{\pi}\right)$ (四次剩余符号) 为满足
\[
\left(\frac{\alpha}{\pi}\right)\equiv\alpha^{\frac{\mathbb{N}\pi-1}{4}}\pmod{\pi}
\]
的唯一的 $1$ 的 $4$ 次根, 并对这些符号证明了一个四次互反律. 后来 Eisenstein 证明了一个三次互反律. Emil Artin 指出, 他的定理 (V, \ref{thm:v.3.5}) 蕴含所有可能的此类互反律, 因而可以视为``不含 $n$ 次单位根的域的互反律''. 本节其余部分解释这一评论.

\subsection*{幂剩余符号}

设 $K$ 是包含本原 $n$ 次单位根的数域. 对 $K$ 的元素的每个有限组 $a,b,\ldots$, 定义 $S(a,b,\ldots)$ 为满足下列条件之一的 $K$ 的素理想集: $\mathrm{ord}_{\mathfrak{p}}(n)\ne 0$, 或 $\mathrm{ord}_{\mathfrak{p}}(a)\ne 0$, 或 $\mathrm{ord}_{\mathfrak{p}}(b)\ne 0$, \ldots. 特别地, $S$ 本身仅由 $n$ 的因子组成.

回顾: $X^{n}-1$ 的判别式只被整除 $n$ 的素整除. 因此, 对任何 $\mathfrak{p}\nmid n$, $X^{n}-1$ 在 $\F_p^{\mathrm{al}}$ 中有 $n$ 个互异的根, 且映射
\[
\zeta\mapsto\zeta \bmod \mathfrak{p}:\mu_{n}(K)\to\mu_{n}(\mathcal{O}_{K}/\mathfrak{p})
\]
对任何满足 $\mathfrak{p}\nmid n$ 的素理想是双射. 对这样的素 $\mathfrak{p}$, 令 $q=\mathbb{N}\mathfrak{p}\overset{\text{def}}{=}(\mathcal{O}_{K}:\mathfrak{p})$. 则 $\F_q^{\times}$ 是 $q-1$ 阶循环群, 故 $n|q-1$, 且 $\zeta^{\frac{q-1}{n}}\in\mu_n\subset\F_q^{\times}$.

对 $a\in K^{\times}$ 与 $\mathfrak{p}\notin S(a)$, 定义 $\left(\frac{a}{\mathfrak{p}}\right)$ 为满足
\[
\left(\frac{a}{\mathfrak{p}}\right)\equiv a^{\frac{\mathbb{N}\mathfrak{p}-1}{n}}\pmod{\mathfrak{p}}
\]
的唯一的 $1$ 的 $n$ 次根.

\noindent\textbf{5.1}\quad 对所有 $a,b\in K^{\times}$ 与 $\mathfrak{p}\notin S(a,b)$,
\[
\left(\frac{ab}{\mathfrak{p}}\right)=\left(\frac{a}{\mathfrak{p}}\right)\left(\frac{b}{\mathfrak{p}}\right).
\]
由定义显然.

\noindent\textbf{5.2}\quad 对 $a\in K^{\times}$ 与 $\mathfrak{p}\notin S(a)$, 下述等价:

(a) $\left(\frac{a}{\mathfrak{p}}\right)=1$;

(b) $a$ 在 $\mathcal{O}_{K}/\mathfrak{p}$ 中成为 $n$ 次幂;

(c) $a$ 在 $K_{\mathfrak{p}}$ 中成为 $n$ 次幂.

(a) 与 (b) 的等价性由
\[
1\to\F_{q}^{\times n}\to\F_{q}^{\times}\xrightarrow{\ x\mapsto x^{\frac{q-1}{n}}\ }\mu_{n}\to1,\qquad q=\mathbb{N}\mathfrak{p}
\]
的正合性推出. 若 $X^{n}-a$ 模 $\mathfrak{p}$ 有解, 则 Hensel 引理 (ANT, 7.33) 表明它在 $K_{\mathfrak{p}}$ 中有解. 反之, 若 $a=\alpha^{n}$, $\alpha\in K_{\mathfrak{p}}$, 则 $\mathrm{ord}_{\mathfrak{p}}(\alpha)=\frac{1}{n}\mathrm{ord}_{\mathfrak{p}}(a)=0$, 故 $\alpha\in\mathcal{O}_{K_{\mathfrak{p}}}$. 映射 $\mathcal{O}_K\to\mathcal{O}_{K_{\mathfrak{p}}}/\mathfrak{p}$ 是满射, 故存在 $\alpha_0\in\mathcal{O}_K$ 映到 $\alpha$ 模 $\mathfrak{p}$ 的像.

我们用线性性把映射 $\mathfrak{p}\mapsto\left(\frac{a}{\mathfrak{p}}\right)$ 延拓到 $I^{S(a)}$: 即对 $\mathfrak{b}=\prod\mathfrak{p}_i^{r_i}\in I^{S(a)}$,
\[
\left(\frac{a}{\mathfrak{b}}\right)=\prod\left(\frac{a}{\mathfrak{p}_{i}}\right)^{r_{i}}.
\]
我们把 $\left(\frac{a}{\mathfrak{b}}\right)$ 简写为 $\left(\frac{a}{b}\right)$.

对使 $S'$ 中的素都不分歧的 Abel 扩张 $L/K$,以 $\psi_{L/K}\colon I^{S}\to\Gal(L/K)$ 表示 Artin 映射 (见第 \ref{ch:V} 章).

\noindent\phantomsection\hypertarget{item:viii.5.3}{}\textbf{5.3}\quad 对所有 $a\in K^{\times}$ 与 $\mathfrak{b}\in I^{S(a)}$,
\[
\psi_{K[a^{\frac{1}{n}}]/K}(\mathfrak{b})\big(a^{\frac{1}{n}}\big)=\left(\frac{a}{\mathfrak{b}}\right)a^{\frac{1}{n}}.
\]
由 Galois 理论, 存在 $1$ 的 $n$ 次根 $\zeta(\mathfrak{b})$ 使 $\psi(\mathfrak{b})(a^{\frac{1}{n}})=\zeta(\mathfrak{b})\cdot a^{\frac{1}{n}}$, 且映射 $\mathfrak{b}\mapsto\zeta(\mathfrak{b})$ 是同态. 因此只需对 $\mathfrak{b}=\mathfrak{p}$ 为素理想的情形证明等式. 按定义,
\[
\psi(\mathfrak{p})(x)\equiv x^{\mathbb{N}\mathfrak{p}}\pmod{\mathfrak{p}}.
\]
由
\[
\psi(\mathfrak{p})\big(a^{\frac{1}{n}}\big)=\zeta(\mathfrak{p})\cdot a^{\frac{1}{n}}
\]
得到
\[
\zeta(\mathfrak{p})\cdot a^{\frac{1}{n}}\equiv a^{\frac{\mathbb{N}\mathfrak{p}}{n}}\pmod{\mathfrak{p}},
\]
由此推出 $\zeta(\mathfrak{p})=\left(\frac{a}{\mathfrak{p}}\right)$.

\noindent\textbf{5.4}\quad 设 $a\in\mathcal{O}_{K}$, 并设 $\mathfrak{b}$ 是 $I^{S(a)}$ 中的整理想. 若 $a'\in\mathcal{O}_{K}$ 满足 $a'\equiv a \bmod \mathfrak{b}$, 则 $\mathfrak{b}\in I^{S(a')}$ 且
\[
\left(\frac{a}{\mathfrak{b}}\right)=\left(\frac{a'}{\mathfrak{b}}\right).
\]
事实上, 对每个整除 $\mathfrak{b}$ 的素理想 $\mathfrak{p}$, 有 $a'\equiv a \bmod \mathfrak{p}$, 故 $\left(\frac{a}{\mathfrak{p}}\right)=\left(\frac{a'}{\mathfrak{p}}\right)$.

Artin 互反律使我们能为 $\left(\frac{a}{\mathfrak{b}}\right)$ 作为 $\mathfrak{b}$ 的函数证明一个类似但较弱的结果.

\noindent\phantomsection\hypertarget{item:viii.5.5}{}\textbf{5.5}\quad 设 $a\in K^{\times}$. 则存在支撑在 $S(a)$ 中的模 $\mathfrak{m}$, 使得 $\left(\frac{a}{\mathfrak{b}}\right)$ 只依赖于 $\mathfrak{b}$ 在射线类群 $C_\mathfrak{m}$ 中的类.

根据命题 (VII, \ref{prop:vii.A.5}), $S(a)$ 包含所有在 $K[a^{\frac{1}{n}}]$ 中分歧的素. 因此 Artin 互反律 (V, \ref{thm:v.3.5}) 表明存在支撑在 $S(a)$ 中的模 $\mathfrak{m}$, 使得 $\psi(\mathfrak{b})$ 只依赖于 $\mathfrak{b}$ 在射线类群 $C_\mathfrak{m}$ 中的类.

\begin{NOTES}
抄自 \mo{299588}. 幂剩余符号是 $K$ 的一个酉 Hecke 特征, 无穷型平凡, 阶有限, 导子整除 $\mathfrak{m}$ 的某个幂. 这一事实可由类域论推出, 例如由本讲义 \hyperlink{item:viii.5.5}{VIII, 5.5}. (解释待补.)
\end{NOTES}

\subsection*{Hilbert 符号}

设 $K_v$ 是包含本原 $n$ 次单位根的局部域. Hilbert 符号是一个配对
\[
a,b\mapsto(a,b)_{v}:K_{v}^{\times}/K_{v}^{\times n}\times K_{v}^{\times}/K_{v}^{\times n}\to\mu_{n},
\]
其中 $\mu_{n}$ 是 $K_{v}$ 中 $1$ 的 $n$ 次根构成的群. 定义它的最自然方式也许是杯积映射
\[
H^{1}(G,\mu_{n})\times H^{1}(G,\mu_{n})\to H^{2}(G,\mu_{n}\otimes\mu_{n}),\qquad G=\Gal(K^{\mathrm{al}}/K),
\]
再接上由不变映射 $\mathrm{inv}_v$ 定义的同构
\[
H^{2}(G,\mu_{n}\otimes\mu_{n})=H^{2}(G,\mu_{n})\otimes\mu_{n}\to\mu_{n}.
\]
不过, 出于 1920 与 1930 年代的精神, 我将用中心单代数来定义它.

回顾 (IV, \S \ref{sec:IV.5}): 对任意 $a,b\in K_v^{\times}$, 定义 $A(a,b;\zeta)$ 为 $K_v$-代数, 它由元素 $i,j$ 生成, 满足关系
\[
i^{n}=a,\qquad j^{n}=b,\qquad ij=\zeta ji.
\]
它是 $K_v$ 上次数为 $n$ 的中心单代数. 当 $n=2$ 时, $A(a,b;-1)$ 就是四元数代数 $H(a,b)$. 我们定义
\[
(a,b)_{v}=\zeta^{-n\cdot\mathrm{inv}_{v}([A(a,b;\zeta)])},
\]
其中 $[A(a,b;\zeta)]$ 是 $A(a,b;\zeta)$ 在 $\mathrm{Br}(K_v)$ 中的类. 由于 $A(a,b;\zeta)$ 被一个 $n$ 次域分裂 (例如, 在 $n=2$ 的情形, $\Q[i]$ 分裂 $A(a,b;-1)$), 其不变量属于 $\frac{1}{n}\Z/\Z$, 故 $n\cdot\mathrm{inv}_v([A(a,b;\zeta)])$ 是 $\Z/n\Z$ 的元素. 显然 $A(a,b;\zeta)$ 的同构类只依赖于作为 $K_v^{\times}/K_v^{\times n}$ 中元素的 $a$ 与 $b$, 因此我们确实得到配对
\[
K_{v}^{\times}/K_{v}^{\times n}\times K_{v}^{\times}/K_{v}^{\times n}\to\mu_{n}.
\]
不过, 从这个视角并不能明显看出该配对是双线性的.

\setcounter{theorem}{5}
\begin{example}\label{ex:viii.5.6}
考虑 $K_v=\Q_p$ ($p$ 为奇素数), $n=2$ 的情形. 则 $(a,b)_p=\pm 1$, 且
\[
\begin{aligned}
(a,b)_{p}=1&\iff H(a,b)\approx M_{2}(K_{v});\\
&\iff X^{2}-aY^{2}-bZ^{2}+abT^{2} \text{ 在 } K_{v} \text{ 中表示 } 0;\\
&\iff b \text{ 是来自 } K[\sqrt{a}] \text{ 的范数};\\
&\iff X^{2}-aY^{2}-bZ^{2} \text{ 在 } K_{v} \text{ 中表示 } 0.
\end{aligned}
\]
证明这些等价时分别用到: 四元数代数的不变量为 $\frac{1}{2}$ 当且仅当它是可除代数; 习题 IV.\ref{ex:iv.5.7}; 命题\ref{prop:viii.3.7}(d); 命题\ref{prop:viii.3.7}(c). 最后一个条件表明我们对 Hilbert 符号的定义与例如 \cite{Serre1970}, III 中的定义一致.
\end{example}

\noindent\textbf{5.7}\quad 对任意 $a,b$,
\[
A(b,a;\zeta)\approx A(a,b;\zeta^{-1})\approx A(a,b;\zeta)^{\mathrm{opp}}.
\]
因此
\[
(b,a)_{v}=(a,b)_{v}^{-1}.
\]
按定义, $A(b,a;\zeta)$ 是由 $i',j'$ 生成的 $K_v$-代数, 满足关系 $i'^{n}=b$, $j'^{n}=a$, $i'j'=\zeta j'i'$. 映射 $i'\mapsto j$, $j'\mapsto i$ 是同构 $A(b,a;\zeta)\to A(a,b;\zeta^{-1})$. 映射 $i\mapsto i$, $j\mapsto j$ 是同构 $A(a,b;\zeta)^{\mathrm{opp}}\to A(a,b;\zeta^{-1})$.

\noindent\phantomsection\hypertarget{item:viii.5.8}{}\textbf{5.8}\quad 设 $a,b\in K^{\times}$. 对任意 $v\notin S(a)$,
\[
(a,b)_{v}=\left(\frac{a}{\mathfrak{p}_{v}}\right)^{\mathrm{ord}_{v}(b)}.
\]
为简单起见, 设 $A(a,b;\zeta)$ 是可除代数. 回顾 (IV, \S \ref{sec:IV.4}): 为计算局部域 $K_v$ 上中心可除代数 $D$ 的不变量, 我们

(a) 选取极大非分歧子域 $L\subset D$;

(b) 找出元素 $\beta\in D$ 使 $\alpha\mapsto\beta\alpha\beta^{-1}$ 是 $L$ 的 Frobenius 自同构 (这样的 $\beta$ 由 Noether--Skolem 定理保证存在);

(c) 令 $\mathrm{inv}_{v}([D])=\mathrm{ord}_{v}(\beta)$.

我们把它应用于 $L=K_v[i]=K_v[a^{\frac{1}{n}}]$. 注意, 由于 $v\notin S(a)$, 该扩张非分歧. 设 $\left(\frac{a}{\mathfrak{p}_v}\right)=\zeta^{r}$, 即 $(\mathfrak{p},L/K_v)(i)=\zeta^{r}i$. 由于 $jij^{-1}=\zeta^{-1}i$, 可取 $\beta=j^{-r}$. 则 $\beta^{n}=b^{-r}$, 故 $\mathrm{ord}_v(\beta)=-\frac{r}{n}\mathrm{ord}_v(b)$. 于是
\[
(a,b)_{v}\overset{\text{def}}{=}\zeta^{-n\,\mathrm{inv}_{v}(A(a,b;\zeta))}=\zeta^{r\cdot\mathrm{ord}_{v}(b)}=\left(\frac{a}{\mathfrak{p}_{v}}\right)^{\mathrm{ord}_{v}(b)}.
\]

\setcounter{theorem}{8}
\begin{remark}\label{rem:viii.5.9}
事实上,
\[
(a,b)_{v}=\frac{\phi_{v}(b)\big(a^{\frac{1}{n}}\big)}{a^{\frac{1}{n}}}
\]
对所有 $a,b,v$ 成立. 见 (III, \ref{rem:iii.4.5}).
\end{remark}

\noindent\textbf{5.10}\quad 对 $a,b\in K^{\times}$,
\[
\prod_{v}(a,b)_{v}=1.
\]
在证明互反律的过程中, 我们证明了对每个 $\beta\in\mathrm{Br}(K)$, $\sum\mathrm{inv}_v(\beta)=0$. 特别地, $\sum\mathrm{inv}_v(A(a,b;\zeta))=0$, 这就蕴含了该公式.

对 $a,b\in K^{\times}$, 定义
\[
\left(\frac{a}{b}\right)=\prod_{v\notin S(a)}\left(\frac{a}{v}\right)^{\mathrm{ord}_{v}(b)}=\left(\frac{a}{(b)^{S(a)}}\right),
\]
其中 $(b)^{S(a)}$ 是由 $b$ 生成的 $I^{S(a)}$ 中的理想. 符号 $\left(\frac{a}{b}\right)$ 对 $b$ 是乘性的, 但 $\left(\frac{aa'}{b}\right)=\left(\frac{a}{b}\right)\left(\frac{a'}{b}\right)$ 一般不成立, 除非 $S(b)\cap S(a,a')=S$.

\setcounter{theorem}{10}
\begin{theorem}[幂互反律]\label{thm:viii.5.11}
设 $a,b\in K^{\times}$ 满足 $S(a)\cap S(b)=S$ (例如 $a$ 与 $b$ 可以互素). 则
\[
\left(\frac{a}{b}\right)\left(\frac{b}{a}\right)^{-1}=\prod_{v\in S}(b,a)_{v}.
\]
此外, 若 $S(c)=S$, 则
\[
\left(\frac{c}{b}\right)=\prod_{v\in S}(b,c)_{v}.
\]
\end{theorem}

\begin{proof}
令 $S'(a)=S(a)\setminus S$, $S'(b)=S(b)\setminus S$. 我们的假设是 $S'(a)$ 与 $S'(b)$ 不相交. 则
\[
\left(\frac{a}{b}\right)=\prod_{v\in S'(b)}\left(\frac{a}{\mathfrak{p}_{v}}\right)^{\mathrm{ord}_{v}(b)}=\prod_{v\in S'(b)}(a,b)_{v}
\]
且
\[
\left(\frac{b}{a}\right)=\prod_{v\in S'(a)}\left(\frac{b}{\mathfrak{p}_{v}}\right)^{\mathrm{ord}_{v}(a)}=\prod_{v\in S'(a)}(b,a)_{v}.
\]
因此
\[
\left(\frac{a}{b}\right)\left(\frac{b}{a}\right)^{-1}=\prod_{v\in S'(a)\cup S'(b)}(a,b)_{v}.
\]
对 $v\notin S\cup S'(a)\cup S'(b)$, $(a,b)_v=1$ (例如由 \hyperlink{item:viii.5.8}{5.8}), 故乘积公式表明
\[
\prod_{v\in S'(a)\cup S'(b)}(a,b)_{v}\times\prod_{v\in S}(a,b)_{v}=1.
\]
这就完成了第一个等式的证明; 第二个是显然的.
\end{proof}

要得到完全显式的公式, 还需计算 $v\in S$ 时的 Hilbert 符号. 对无限素这很容易: 若 $v$ 为复, 则 $(a,b)_v=1$ 恒成立; 若 $v$ 为实, 则
\[
(a,b)_{v}=1\iff X^{2}-aY^{2}-bZ^{2} \text{ 表示 } 0\iff a>0 \text{ 或 } b>0.
\]
对 $K=\Q$ 与 $n=2$,
\[
(u2^{r},v2^{s})_{2}=(-1)^{\frac{u-1}{2}\cdot\frac{v-1}{2}+r\frac{v^{2}-1}{8}+s\frac{u^{2}-1}{8}},
\]
其中 $u,v$ 是 2-adic 单位, 指数模 $2$ 理解. 其初等证明见 \cite{Serre1970}, III, 1.2. 依次把这个公式应用于 $(p,q)$ ($p,q$ 为奇素数)、$(2,p)$ ($p$ 为奇素数) 与 $(-1,p)$ ($p$ 为奇素数), 便得到经典的二次互反律 (包括其补充).

对奇素数 $p$ 与 $K=\Q[\zeta]$ ($\zeta$ 为本原 $p$ 次单位根), 可以把 Hilbert 符号 $(a,b)_p$ 完全显式化. 回顾: $p$ 在 $K$ 中全分歧, 且 $(p)=(\pi)^{p-1}$, 其中 $\pi=1-\zeta$. 设 $K_\pi$ 表示 $K$ 在 $(\pi)$ 处的完备化, 并设 $U_i$ 表示 $K_\pi$ 中同余于 $1 \bmod \pi^{i}$ 的单位群. 我们有滤链
\[
\mathcal{O}_{K_{\pi}}^{\times}\supset U_{1}\supset U_{2}\supset\cdots\supset U_{p+1}\supset\cdots.
\]
若 $u\in U_{p+1}$, 则 $u$ 是 $K_\pi$ 中的 $p$ 次幂 (见 (VII, 注 \ref{rem:vii.A.6}(a))). 由此可以推出: $K_\pi^{\times}/K_\pi^{\times p}$ 由元素
\[
\pi,\ \zeta,\ 1-\pi^{2},\ \ldots,\ 1-\pi^{p}
\]
自由生成 (作为 $\F_p$-向量空间).

设 $\eta_i=1-\pi^{i}$, $i\ge 1$ (例如 $\eta_1=\zeta$).

\begin{proposition}\label{prop:viii.5.12}
Hilbert 配对
\[
a,b\mapsto(a,b)_{\pi}:K_{\pi}^{\times}\times K_{\pi}^{\times}\to\mu_{p}
\]
是满足下列条件的唯一的斜对称配对:

(a) $(\eta_{i},\eta_{j})_{\pi}=(\eta_{i},\eta_{i+j})_{\pi}\,(\eta_{i+j},\eta_{j})_{\pi}\,(\eta_{i+j},\pi)_{\pi}^{-j}$ 对所有 $i,j\ge 1$;

(b) $(\eta_{i},\pi)_{\pi}=
\begin{cases}
1 & \text{若 } 1\le i\le p-1,\\
\zeta & \text{若 } i=p;
\end{cases}$

(c) $(\cdot,\cdot)_{\pi}=1$ 在 $U_{i}\times U_{j}$ 上, 若 $i+j\ge p+1$.
\end{proposition}

提示见 \cite{CasselsFrohlich1967}, p.~354.

\begin{example}\label{ex:viii.5.13}
(三次互反律; Eisenstein). 设 $p=3$, 此时 $K=\Q[\zeta]$, $\zeta=\frac{-1+\sqrt{-3}}{2}$, $\pi=-\zeta\sqrt{-3}$. 则 $\mathcal{O}_K=\Z[\zeta]$, 且 $\mathcal{O}_K$ 的每个非零元素都可写成 $\zeta^{i}\pi^{j}a$ 的形式, 其中 $a\equiv\pm1 \bmod 3\mathcal{O}_K$. 此时互反律成为:
\[
\left(\frac{a}{b}\right)=\left(\frac{b}{a}\right)
\]
若 $a$ 与 $b$ 互素且同余于 $\pm 1 \bmod 3\mathcal{O}_K$; 而
\[
\begin{cases}
\left(\dfrac{\zeta}{a}\right)=\zeta^{-m-n},\\[2ex]
\left(\dfrac{\pi}{a}\right)=\zeta^{m},
\end{cases}
\]
若 $a=\pm\big(1+3(m+n\zeta)\big)$.

注意, 若 $a\in\Z$, 则 $a\equiv\pm1 \bmod 3\mathcal{O}_K$ 自动成立.
\end{example}

\subsection*{应用}

固定奇素数 $p$ 与本原 $p$ 次单位根 $\zeta$. 若整数 $x,y,z$ 满足 $x^{p}+y^{p}=z^{p}$, 则
\[
\prod_{i=0}^{p-1}(x+\zeta^{i}y)=z^{p}.
\]
不妨设 $x,y,z$ 无公因子. 若 $p\nmid xyz$, 则 $\Z[\zeta]$ 中的元素 $x+\zeta^{i}y$ 两两互素 (ANT, 6.9). 因此, 每一个生成的理想都是 $p$ 次幂, 对
\[
\alpha=\frac{x+\zeta y}{x+y}=1-\frac{y\pi}{x+y},\qquad \pi=1-\zeta
\]
亦然.

于是\footnote{此脚注修正了上一版本中一处令人困惑的排印错误; 见 \sx{667386}.}
$\left(\frac{\beta}{\alpha}\right)=1$ 对所有与 $\alpha$ 互素的 $\beta\in\Z[\zeta]$ 成立.

\begin{theorem}\label{thm:viii.5.14}
设 $x,y,z$ 是两两互素的正整数, 满足 $p\nmid xyz$ 且 $x^{p}+y^{p}=z^{p}$. 则对每个整除 $xyz$ 的素数 $q$, 有 $q^{p-1}\equiv 1 \bmod p^{2}$.
\end{theorem}

\begin{proof}
此时幂互反律成为
\[
\left(\frac{\beta}{\alpha}\right)\left(\frac{\alpha}{\beta}\right)^{-1}=\zeta^{\mathrm{Tr}_{\Q[\zeta]/\Q}(\eta)},
\]
其中 $\eta=\dfrac{\beta-1}{p}\cdot\dfrac{\alpha-1}{\pi}$. 我们把此式应用于 $\beta=q^{p-1}$. 不失一般性, 设 $q|y$, 于是 $\alpha\equiv 1 \bmod q$ 且 $\left(\frac{\alpha}{q}\right)=1$. 再者,
\[
\mathrm{Tr}(\eta)=\frac{q^{p-1}-1}{p}\,\mathrm{Tr}\Big(\frac{\alpha-1}{\pi}\Big),
\]
但
\[
\mathrm{Tr}\frac{\alpha-1}{\pi}=\mathrm{Tr}\Big(-\frac{y}{x+y}\Big)=-\frac{y}{x+y}(p-1)
\]
不被 $p$ 整除. 故 $\frac{q^{p-1}-1}{p}$ 被 $p$ 整除.
\end{proof}

\begin{corollary}[WIEFERICH 条件]\label{cor:viii.5.15}
若 $X^{p}+Y^{p}=Z^{p}$ 有解 $x,y,z$, 其中 $x,y,z$ 都是不被 $p$ 整除的正整数, 则 $2^{p-1}\equiv 1 \bmod p^{2}$.
\end{corollary}

\begin{proof}
若 $x^{p}+y^{p}=z^{p}$, 则 $x,y,z$ 中至少有一个是偶数.
\end{proof}

类似的论证 (用不同的 $\beta$) 证明了 Mirimanoff 条件: $3^{p-1}\not\equiv 1 \bmod p^{2}$.

满足 Wieferich 条件的 $<3\times 10^{9}$ 的素数只有 $1093$ 与 $3511$, 而它们都不满足 Mirimanoff 条件. 因此这证明了 Fermat 大定理第一情形对 $p<3\times 10^{9}$ 成立.

\begin{exercise}\label{exc:viii.5.16}
设 $p\in\Z$ 是同余于 $1$ 模 $3$ 的素数 (于是 $\F_p$ 包含 $1$ 的三次根). 证明: $2$ 模 $p$ 是立方当且仅当 $p$ 具有形式 $x^{2}+27y^{2}$, $x,y\in\Z$.
\end{exercise}

\begin{NOTES}
定理\ref{thm:viii.5.14} 由 Furtw\"angler 证明 (见 Hasse 1970, II, 22); 另见 \cite{Koch1992}, II, 6.3 与 Herbrand 1936, p.~47. 类域论还可以简化 Kummer 第二判别法在 Fermat 定理第二情形成立时的证明 (J.~Herbrand, Sur les classes des corps circulaires, \emph{J.~Math.~Pures Appl.}, IX.~S\'er.~11, 417--441).
\end{NOTES}
\section{数域上二次型的分类}\label{sec:VIII.6}

前面我们证明了: 数域上的非退化二次型在域中表示零当且仅当它在域的每个完备化中表示零. 本节对数域上的二次型给出完全的分类. 具体地, 我们将:

(a) 证明数域 $K$ 上的两个二次型等价当且仅当它们在 $K$ 的每个完备化上都等价;

(b) 给出局部域上二次型的完备不变量列表;

(c) 确定哪些局部不变量族来自一个整体二次型.

\subsection*{二次型的一般知识}

本小节中, $k$ 是特征 $\ne 2$ 的任意域. 设 $(V,Q)$ 是 $k$ 上带伴随双线性型 $B$ 的二次空间,
\[
B(v,w)=\frac{1}{2}\big(Q(v+w)-Q(v)-Q(w)\big),
\]
并设 $U_1,U_2$ 是 $V$ 的子空间. 若 $V$ 的每个元素都可以唯一地写成 $v=u_1+u_2$ 的形式, 其中 $u_1\in U_1$, $u_2\in U_2$, 则记 $V=U_1\oplus U_2$. 若再有 $B(u_1,u_2)=0$ 对所有 $u_1\in U_1$, $u_2\in U_2$ 成立, 则记 $V=U_1\perp U_2$. 对 $V$ 的每个子空间 $U$,
\[
U^{\perp}=\{v\in V\mid B(u,v)=0 \text{ 对所有 } u\in U\}.
\]
若 $Q|U$ 非退化, 则 $V=U\perp U^{\perp}$.

设 $(V,Q)$ 与 $(V',Q')$ 是 $k$ 上的二次空间. 态射 $s\colon(V,Q)\to(V',Q')$ 是 $k$-向量空间的线性映射 $s\colon V\to V'$, 使得 $Q'(s(v))=Q(v)$ 对所有 $v\in V$ 成立. 态射称为\textbf{同构} (isomorphism), 若它有同为态射的逆. 同构 $(V,Q)\to(V',Q')$ 也称为\textbf{等距} (isometry).

\begin{proposition}\label{prop:viii.6.1}
设 $(V,Q)$ 是二次空间. 若 $Q$ 表示 $a\in k^{\times}$, 则存在 $e\in V$ 使 $Q(e)=a$, 且存在子空间 $U\subset V$ 使 $V=U\perp k\cdot e$.
\end{proposition}

\begin{proof}
因为 $Q$ 表示 $a$, 存在 $e\in V$ 使 $Q(e)=a$, 取 $U$ 为 $k\cdot e$ 的正交补即可.
\end{proof}

设 $(V,Q)$ 是二次空间. 对每个满足 $Q(y)\ne 0$ 的 $y\in V$, 定义关于 $y$ (或关于直线 $k\cdot y$) 的对称 (symmetry) 为映射
\[
\tau_{y}(x)=x-\frac{2B(x,y)}{Q(y)}\,y.
\]
注意 $\tau_y$ 是态射 $(V,Q)\to(V,Q)$ 且 $\tau_y\circ\tau_y=\mathrm{id}$, 故 $\tau_y$ 是等距. 它把直线 $k\cdot y$ 中的每个向量变号, 并固定超平面 $(k\cdot y)^{\perp}$ 中的每个向量. 因此它是关于超平面 $(k\cdot y)^{\perp}$ 的反射.

\begin{proposition}\label{prop:viii.6.2}
设 $U$ 与 $W$ 是二次空间 $(V,Q)$ 的等距子空间, 并设 $Q|U$ 非退化. 则 $U^{\perp}$ 与 $W^{\perp}$ 等距.
\end{proposition}

\begin{proof}
对 $\dim U$ 作归纳证明. 先设 $U$ 与 $W$ 都是直线, 即 $U=ku$, $W=kw$. 则 $Q(u)\ne 0$, $Q(w)\ne 0$, 且不妨设 $Q(u)=Q(w)$. 由
\[
Q(u+w)+Q(u-w)=2Q(u)+2Q(w)=4Q(u)
\]
可见 $Q(u+w)$ 与 $Q(u-w)$ 中至少有一个非零, 必要时把 $w$ 换成 $-w$, 可设是后者. 于是对称 $\tau_{u-w}$ 有定义:
\[
\tau_{u-w}x=x-\frac{2B(x,u-w)}{Q(u-w)}\,(u-w).
\]
则 $\tau_{u-w}(u)=w$, 因为
\[
Q(u-w)=Q(u)+Q(w)-2B(u,w)=2Q(u)-2B(u,w)=2B(u,u-w),
\]
故 $\tau_{u-w}$ 把 $U^{\perp}$ 等距地映到 $W^{\perp}$.

于是可设 $\dim U\ge 2$, 从而 $U$ 有非平凡分解 $U=U_1\perp U_2$. 由于 $W\approx U$, 存在分解 $W=W_1\perp W_2$ 使 $U_1\approx W_1$, $U_2\approx W_2$. 注意 $Q|U_1$ 非退化, 且 $U_1^{\perp}=U_2\perp U^{\perp}$. 归纳假设蕴含 $U_2\perp U^{\perp}$ 等距于 $W_2\perp W^{\perp}$, 而等距的选取定义了分解 $U_2\perp U^{\perp}=X\perp Y$, 其中 $X\approx W_2$, $Y\approx W^{\perp}$. 但 $X\approx U_2$, 再由归纳假设得 $Y\approx U^{\perp}$. 故 $W^{\perp}\approx U^{\perp}$.
\end{proof}

对二次空间 $(V,Q)$ 选一组基 $e_i$, 便得到二次型
\[
q(X_{1},\ldots,X_{n})=\sum a_{ij}X_{i}X_{j},\qquad a_{ij}=B(e_{i},e_{j}).
\]
反之, 二次型 $q$ 定义二次空间 $(k^{n},q)$.

两个二次型 $q$ 与 $q'$ 称为\textbf{等价的} (equivalent), 记作 $q\sim q'$, 若它们定义同构的二次空间, 即其中一个可以由另一个经可逆变量代换得到. 若 $q$ 与 $q'$ 是两组不同变量的二次型, 则记 $q+q'$ 为 $q\perp q'$; 此时 $(k^{m+n},q\perp q')=(k^{m},q)\perp(k^{n},q')$.

由命题\ref{prop:viii.6.2} 得到:

设 $q=r\perp s$ 与 $q'=r'\perp s'$ 是两个二次型, 并设 $r$ 非退化. 若 $q\sim q'$ 且 $r\sim r'$, 则 $s\sim s'$.

由命题\ref{prop:viii.6.1} 得到:

$n$ 元非退化二次型 $q$ 表示 $a$ 当且仅当 $q\sim r\perp aZ^{2}$, 其中 $r$ 是 $n-1$ 元二次型.

二次空间 $(V,Q)$ 的\textbf{秩} (rank) 定义为矩阵 $(B(e_i,e_j))$ 对某组基 $e_i$ 的秩. 二次型 $q$ 的秩是相应二次空间的秩. 当 $q$ 写成对角形 $q=a_1X_1^{2}+\cdots+a_rX_r^{2}$ 时, $q$ 的秩是非零系数 $a_i$ 的个数, 即 $q$ 中实际出现的变量个数.

\subsection*{局部--整体原则}

\begin{theorem}[HASSE--MINKOWSKI]\label{thm:viii.6.3}
设 $q$ 与 $q'$ 是数域 $K$ 上的二次型. 若 $q$ 与 $q'$ 在所有素 $v$ 的 $K_v$ 上都等价, 则 $q$ 与 $q'$ 在 $K$ 上等价.
\end{theorem}

\begin{proof}
不妨设 $q$ 与 $q'$ 非退化. 对 $q$ 与 $q'$ 的公共秩 $n$ 作归纳. 若 $n=0$, 两个型都是零型, 无需证明. 否则, 存在 $q$ 表示的 $a\in K^{\times}$. 则 $q(X_1,\ldots,X_n)-aZ^{2}$ 在 $K$ 中表示 $0$, 从而对所有 $v$ 在 $K_v$ 中表示 $0$. 把定理\ref{thm:viii.3.5} 应用于 $q'-aZ^{2}$, 得知 $q'$ 在 $K$ 中表示 $a$. 因此 $q\sim q_1\perp aZ^{2}$, $q'\sim q_1'\perp aZ^{2}$, 其中 $q_1,q_1'$ 是 $n-1$ 秩二次型. 现在 (\ref{prop:viii.6.2}) 表明 $q_1$ 与 $q_1'$ 对所有 $v$ 在 $K_v$ 上等价, 故 (由归纳) 它们在 $K$ 上等价. 这蕴含 $q$ 与 $q'$ 在 $K$ 上等价.
\end{proof}

\begin{remark}\label{rem:viii.6.4}
设 $(V,Q)$ 是域 $k$ 上的二次空间, $O$ 是其等距群. 定理\ref{thm:viii.6.3} 说的是
\[
H^{1}(K,O)\to\prod_{v}H^{1}(K_{v},O)
\]
是单射.
\end{remark}

\subsection*{局部域上二次型的分类}

阿基米德情形. $\CC$ 上的任何二次型 (对任何特征 $\ne 2$ 的代数闭域亦然) 等价于唯一的二次型
\[
X_{1}^{2}+\cdots+X_{n}^{2}.
\]
因此 $\CC$ 上两个二次型等价当且仅当它们有相同的秩 $n$.

根据 Sylvester 定理, $\R$ 上的二次型 $q$ 等价于唯一的二次型
\[
X_{1}^{2}+\cdots+X_{r}^{2}-X_{r+1}^{2}-\cdots-X_{r+t}^{2}.
\]
$-1$ 的个数 $t$ 称为\textbf{负惯性指数} (index of negativity). 因此 $\R$ 上两个二次型等价当且仅当它们有相同的秩 $n$ 与相同的负惯性指数 $t$.

非阿基米德情形. 设 $K$ 是局部域. 回顾: Hilbert 符号 $(\cdot,\cdot)$ 对 $a,b\in K^{\times}$ 可定义为
\[
(a,b)=
\begin{cases}
1 & \Longleftrightarrow X^{2}-aY^{2}-bZ^{2} \text{ 表示 } 0\\
-1 & \text{否则,}
\end{cases}
\Longleftrightarrow aY^{2}+bZ^{2} \text{ 表示 } 1.
\]

\begin{lemma}\label{lem:viii.6.5}
Hilbert 符号具有如下性质:

(a) 它是双乘性的, 且 $(ac^{2},bd^{2})=(a,b)$ 对所有 $a,b,c,d\in K^{\times}$;

(b) 对任何非平方 $a\in K^{\times}$, 存在 $b\in K^{\times}$ 使 $(a,b)=-1$;

(c) $(b,a)=(a,b)^{-1}=(a,b)$;

(d) $(a,-a)=(1,a)=1$.
\end{lemma}

\begin{proof}
显然, 把 $a$ 或 $b$ 乘以平方不改变 $(a,b)$; (c) 也显然.

注意 $(a,b)=1$ 当且仅当 $b$ 是来自 $K[\sqrt{a}]$ 的范数. 由局部类域论, 若 $a$ 不是 $K$ 中的平方, 则 $\Nm(K[\sqrt{a}]^{\times})$ 是 $K^{\times}$ 中指数 $2$ 的子群, 故 $b\mapsto(a,b)$ 是同构 $K^{\times}/\Nm(K[\sqrt{a}]^{\times})\to\{\pm 1\}$. 这完成了 (a) 与 (b) 的证明. 最后, $aX^{2}-aY^{2}=a(X^{2}-Y^{2})$, 而 $X^{2}-Y^{2}$ 表示 $a^{-1}$, 因为它表示 $0$.
\end{proof}

若 $q\sim a_{1}X_{1}^{2}+\cdots+a_{n}X_{n}^{2}$, 其中 $a_{1},\ldots,a_{n}\in K^{\times}$, 则令
\[
\begin{aligned}
n(q)&=n,\\
d(q)&=a_{1}\cdots a_{n}\qquad(\text{属于 } K^{\times}/K^{\times 2}),\\
S(q)&=\prod_{1\le i\le j\le n}(a_{i},a_{j})=\prod_{1\le i\le n}(a_{i},d_{i})\qquad(\text{属于 } \{\pm 1\}),
\end{aligned}
\]
其中 $d_i=a_1\ldots a_i$. 于是 $n(q)$ 是 $q$ 的秩, $d(q)$ 是 $q$ 的判别式, 两者都只依赖于 $q$ 的等价类. 我们将证明 $S(q)$ 亦然. 它称为 $q$ 的 \textbf{Hasse 不变量}\index{hassebubianliang@Hasse 不变量(Hasse invariant)}.

\begin{remark}\label{rem:viii.6.6}
\cite{Serre1970} 定义
\[
\varepsilon(q)=\prod_{1\le i<j\le n}(a_{i},a_{j})=\prod(a_{i},d_{i-1}).
\]
注意
\[
S(q)=\varepsilon(q)\prod_{i=1}^{n}(a_{i},a_{i})=\varepsilon(q)\prod_{i=1}^{n}(-1,a_{i})(-a_{i},a_{i})=\varepsilon(q)(-1,d(q)).
\]
因此知道 $(d(q),S(q))$ 等价于知道 $(d(q),\varepsilon(q))$.
\end{remark}

\begin{proposition}\label{prop:viii.6.7}
元素 $S(q)$ 只依赖于 $q$ 的等价类.
\end{proposition}

\begin{proof}
只需证明 $\varepsilon(q)$ 只依赖于 $q$ 的等价类. 当 $q$ 的秩为 $1$ 时无需证明: 对所有 $q$, $\varepsilon(q)=1$ (空积).

其次设 $q\sim aX^{2}+bY^{2}\sim a'X^{2}+b'Y^{2}$. 由于二者等价, $aX^{2}+bY^{2}$ 与 $a'X^{2}+b'Y^{2}$ 要么都表示 $1$, 要么都不表示 $1$, 故 $(a,b)=(a',b')$.

再设 $n>2$ 且
\[
q\sim a_{1}X_{1}^{2}+\cdots+a_{i}X_{i}^{2}+a_{i+1}X_{i+1}^{2}+\cdots\sim a_{1}^{\prime}X_{1}^{2}+\cdots+a_{i}^{\prime}X_{i}^{2}+a_{i+1}^{\prime}X_{i+1}^{2}+\cdots
\]
其中 $a_{j}=a_{j}^{\prime}$ 对 $j\ne i,i+1$ 成立. 我们要证明
\[
(a_{i},d_{i-1})(a_{i+1},d_{i})=(a_{i}^{\prime},d_{i-1}^{\prime})(a_{i+1}^{\prime},d_{i}^{\prime}).
\]
但
\[
(a_{i},d_{i-1})(a_{i+1},d_{i})=(a_{i},d_{i-1})(a_{i+1},d_{i-1})(a_{i+1},a_{i})=(a_{i}a_{i+1},d_{i-1})(a_{i},a_{i+1}),
\]
而 $a_{i}a_{i+1}$ 与 $a_{i}^{\prime}a_{i+1}^{\prime}$ 相差一个平方, 故剩下只需证 $(a_{i},a_{i+1})=(a_{i}^{\prime},a_{i+1}^{\prime})$. 根据命题\ref{prop:viii.6.2}, $a_{i}X_{i}^{2}+a_{i+1}X_{i+1}^{2}\sim a_{i}^{\prime}X_{i}^{2}+a_{i+1}^{\prime}X_{i+1}^{2}$, 而前面已证明这蕴含 $(a_{i},a_{i+1})=(a_{i}^{\prime},a_{i+1}^{\prime})$.
\end{proof}

下面的初等引理完成了证明.

\begin{lemma}\label{lem:viii.6.8}
设 $B$ 与 $B'$ 是非退化二次空间 $(V,Q)$ 的正交基. 则存在一列正交基 $B_1,B_2,\ldots,B_m$ 使 $B_1=B$, $B_m=B'$, 且每个 $B_i$ 都由 $B_{i-1}$ 改动至多两个相邻元素得到.
\end{lemma}

\begin{proof}
见 \cite{OMeara1963}, 引理 58.1.
\end{proof}

\begin{proposition}\label{prop:viii.6.9}
设 $q$ 是非阿基米德局部域 $K$ 上的 $n$ 元非退化二次型, $a\in K^{\times}$. 则 $q$ 表示 $a$ 当且仅当

(a) $n=1$ 且 $a=d(q)$ (在 $K^{\times}/K^{\times 2}$ 中);

(b) $n=2$ 且 $(a,-d)(-1,d)=S(q)$ (等价地, $(a,-d)=\varepsilon(q)$);

(c) $n=3$ 且要么 $a\ne -d(q)$ (模平方), 要么 $a=-d(q)$ (模平方) 且 $(-1,-1)=S(q)$;

(d) $n\ge 4$.
\end{proposition}

\begin{proof}
(a) 显然 $dX^{2}$ 表示 $a$ 当且仅当 $a=d$ (在 $K^{\times}/K^{\times 2}$ 中).

(b) 设 $q=bX^{2}+cY^{2}$. 显然 $bX^{2}+cY^{2}$ 表示 $a$ 当且仅当 $abX^{2}+acY^{2}$ 表示 $1$, 即 $(ab,ac)=1$. 但
\[
(ab,ac)=(a,a)(a,b)(a,c)(b,c)=(a,-1)(a,d(q))(b,c)=(a,-d(q))\cdot\varepsilon(q),
\]
故条件为
\[
\varepsilon(q)=(a,-d(q)).
\]

(c) 设 $q=a_1X_1^{2}+a_2X_2^{2}+a_3X_3^{2}$. 则 $q$ 表示 $a$ 当且仅当存在 $e\in K^{\times}$ 使方程
\[
a_{1}X_{1}^{2}+a_{2}X_{2}^{2}=e=a_{3}X_{3}^{2}-aX_{4}^{2}
\]
有解. 根据 (b), 这当且仅当
\begin{equation*}
(e,-a_{1}a_{2})=(a_{1},a_{2}),\qquad (e,a_{3}a)=(a_{3},-a)
\tag{$*$}
\end{equation*}
时成立. 考虑 $\F_2$-向量空间 $V$ ($\dim V\ge 2$) 上的两个线性型 $f,g\colon V\to\F_2$. 联立线性方程 $f(x)=\varepsilon_1$, $g(x)=\varepsilon_2$ 有解, 除非它们不相容, 即: $f=0$ 且 $\varepsilon_1=-1$; 或 $g=0$ 且 $\varepsilon_2=-1$; 或 $f=g$ 且 $\varepsilon_1=-\varepsilon_2$.

把这个观察应用于线性型 $(\cdot,-a_1a_2)$, $(\cdot,a_3a)\colon K^{\times}/K^{\times 2}\to\{\pm 1\}$, 我们发现: 存在满足 (*) 的 $e$, 除非 $-a_1a_2=a_3a$ (在 $K^{\times}/K^{\times 2}$ 中) 且 $(a_1,a_2)=-(a_3,a)$. 第一个等式说 $a=-d(q)$ (模平方), 而 (当 $a=-d(q)$ 时) 第二个说 $(-1,-1)=S(q)$.

(d) 此时 $q(X_1,\ldots,X_n)-aZ^{2}$ 的秩 $\ge 5$, 故它表示 $0$ (见 \ref{cor:viii.3.10}).
\end{proof}

\begin{theorem}\label{thm:viii.6.10}
非阿基米德局部域上的两个二次型等价当且仅当它们有相同的秩、相同的判别式与相同的 Hasse 不变量.
\end{theorem}

\begin{proof}
我们在命题\ref{prop:viii.6.7} 中证明了等价的型有相同的不变量. 反之, 对两个型 $q$ 与 $q'$ 的公共秩 $n$ 作归纳. 秩为 $1$ 的两个二次型显然在判别式相同时等价, 故可设 $n>1$. 由命题\ref{prop:viii.6.9}, $q$ 与 $q'$ 表示 $K^{\times}$ 中相同的元素. 特别地, 存在同时被 $q$ 与 $q'$ 表示的 $a\in K^{\times}$. 于是
\[
q\sim q_{1}\perp aZ^{2},\qquad q'\sim q_{1}^{\prime}\perp aZ^{2},
\]
其中 $q_1,q_1'$ 是 $n-1$ 秩二次型. 现在
\[
d(q)=a\cdot d(q_{1}),\qquad S(q)=(a,d(q_{1}))\cdot S(q_{1}),
\]
对 $q'$ 与 $q_1'$ 亦然. 因此 $q_1$ 与 $q_1'$ 有相同的不变量, 归纳假设表明 $q_1\sim q_1'$.
\end{proof}

\begin{proposition}\label{prop:viii.6.11}
设 $q$ 是非阿基米德局部域 $K$ 上秩为 $n$ 的二次型.

(a) 若 $n=1$, 则 $S(q)=(-1,d)$.

(b) 若 $n=2$, 则 $d(q)=-1$ (模平方) 蕴含 $S(q)=(-1,-1)$.

除这些限制之外, 每个三元组 $n\ge 1$, $d\in K^{\times}/K^{\times 2}$, $s=\pm 1$ 都可以作为 $K$ 上某个二次型的不变量组出现.
\end{proposition}

\begin{proof}
情形 $n=1$. 此时 $q=dX^{2}$, 而 $S(q)=(d,d)=(-1,d)$.

情形 $n=2$. 对 $q=aX^{2}+bY^{2}$, $S(q)=(a,a)(b,d)$, 故
\[
d=-1\Rightarrow S(q)=(-1,a)(-1,b)=(-1,d)=(-1,-1).
\]
反之, 型 $X^{2}-Y^{2}$ 满足 $d=-1$, $S=(-1,-1)$.

现在设 $d\ne -1$ 与 $s$ 已给定. 我们找 $a\in K^{\times}$ 使 $q\overset{\text{def}}{=}aX^{2}+adY^{2}$ 满足 $S(q)=s$. 但
\[
S(q)=(a,a)(ad,d)=(a,-1)(a,d)(d,d)=(a,-d)(d,d).
\]
由于 $-d\ne 1$ (在 $K^{\times}/K^{\times 2}$ 中), 可选 $a$ 使 $(a,-d)=s\cdot(d,d)$.

情形 $n=3$. 取 $a\in K^{\times}$ 使 $a\ne -d$ (在 $K^{\times}/K^{\times 2}$ 中). 由于对 $a$ 的这个条件, 存在秩为 $2$ 的二次型 $q_1$ 使
\[
d=d(q_{1})a,\qquad s=S(q_{1})(a,d).
\]
取 $q=q_1\perp aZ^{2}$.

情形 $n\ge 4$. 存在形状为
\[
q_{1}(X_{1},X_{2},X_{3})+X_{4}^{2}+\cdots+X_{n}^{2}
\]
且具有所需不变量的二次型.
\end{proof}

\paragraph{一般说明}
对 $\R$ 或 $\CC$ 上的二次型 $q$, 用与非阿基米德情形相同的公式定义 Hasse 不变量. 对 $\CC$, 恒有 $S(q)=1$; 对 $\R$, $S(q)=(-1)^{t(t+1)/2}$, 其中 $t$ 是负惯性指数 (因为 $(-1,-1)=-1$). 注意在第二种情形 $d(q)=(-1)^{t}$ (在 $\R^{\times}/\R^{\times 2}$ 中), 且当 $r\le 3$ 时 $d(q)$ 与 $S(q)$ 确定 $t$, 而当 $r>3$ 时不能.

我们说系统 $(n,d,s,\ldots)$ ($n\in\mathbb{N}$, $d\in K^{\times}/K^{\times 2}$, $s\in\{\pm1\}$, \ldots) 是\textbf{可实现的} (realizable), 若存在二次型 $q$ 使 $n(q)=n$, $d(q)=d$, $S(q)=s$, \ldots

(a) 对非阿基米德局部域 $K$, $(n,d,s)$ 可实现, 只要 $n=1$ 时 $s=1$, 且 $n=2$、$d=-1$ 时 $s=(-1,-1)$.

(b) 对 $\R$, $(n,d,s,t)$ 可实现, 只要 $0\le t\le n$, $d=(-1)^{t}$, $s=(-1)^{t(t+1)/2}$.

\subsection*{整体域上二次型的分类}

\begin{theorem}\label{thm:viii.6.12}
设 $n\in\mathbb{N}$, 并设对数域 $K$ 的每个素 $v$ 给定了 $K_v$ 上秩为 $n$ 的非退化二次型 $q(v)$. 则存在 $K$ 上的二次型 $q_0$ 使 $(q_0)_v\sim q_v$ 对每个 $v$ 成立, 当且仅当

(a) 存在 $d_0\in K^{\times}$ 使 $d_0\equiv d(q_v) \bmod K_v^{\times 2}$ 对所有 $v$;

(b) $S(q(v))=1$ 对几乎所有 $v$, 且 $\prod_{v}S(q(v))=1$.
\end{theorem}

这些条件显然是必要的. 鉴于命题\ref{prop:viii.6.11} 及其后的说明, 可以把该定理重述如下. 设给定:

$\diamond$ 一个 $n\ge 1$ 与一个 $d_0\in K^{\times}/K^{\times 2}$;

$\diamond$ 对每个素 $v$ (有限的或实的), 一个 $s_v\in\{\pm 1\}$;

$\diamond$ 对每个实素 $v$, 一个整数 $t_v$.

则: 存在 $K$ 上秩为 $n$、判别式为 $d_0$、Hasse 不变量 $S_v(q)=s_v$ 对所有 $v$、负惯性指数 $t_v(q)=t_v$ 对所有实素 $v$ 的二次型 $q$, 当且仅当

(a) $s_v=1$ 对除有限多个之外的所有 $v$, 且 $\prod_v s_v=1$;

(b) 若 $n=1$, 则 $s_v=(-1,d)_v$; 若 $n=2$, 则要么 $d\ne -1$ (在 $K_v^{\times}/K_v^{\times 2}$ 中), 要么 $s_v=(-1,-1)_v$;

(c) 对所有实素 $v$, $0\le t_v\le n$, $d_v=(-1)^{t_v}$ (模平方), 且 $s_v=(-1)^{t_v(t_v+1)/2}$.

在 $n=1$ 的情形, $q_v=d(q_v)X^{2}$, 可取 $q_0=d_0X^{2}$, 其中 $d_0$ 是由 (a) 保证存在的 $K^{\times}$ 中的元素.

关键情形是 $n=2$, 为此需要下面的引理, 其证明需要类域论.

\begin{lemma}\label{lem:viii.6.13}
设 $T$ 是 $K$ 的由实素或有限素组成的有限集, $b\in K^{\times}$. 若 $T$ 有偶数个元素, 且 $b$ 对任何 $v\in T$ 都不在 $K_v^{\times}$ 中成为平方, 则存在 $a\in K^{\times}$ 使
\[
(a,b)_{v}=
\begin{cases}
-1 & \text{对 } v\in T,\\
1 & \text{否则.}
\end{cases}
\]
\end{lemma}

\begin{proof}
(仿照 Tate, 1976, 5.2). 设 $L$ 是 $K$ 的所有指数为 $2$ 的 Abel 扩张的复合, 并令 $G=\Gal(L/K)$. 由类域论,
\[
G\simeq C_{K}/2C_{K}\simeq\mathbb{I}/K^{\times}\cdot\mathbb{I}^{2}.
\]
列
\[
0\to\mu_{2}\to L^{\times}\xrightarrow{\ x\mapsto x^{2}\ }L^{\times 2}\to0
\]
的上同调列是正合列
\[
K^{\times}\xrightarrow{\ x\mapsto x^{2}\ }K^{\times}\cap L^{\times 2}\to\Hom_{\mathrm{conts}}(G,\mu_{2})\to0.
\]
$K^{\times}$ 的每个元素在 $L$ 中都成为平方, 故有同构
\[
K^{\times}/K^{\times 2}\to\Hom_{\mathrm{conts}}(\mathbb{I}/K^{\times}\cdot\mathbb{I}^{2},\mu_{2}).
\]
该映射把 $a\in K^{\times}$ 送到连续同态
\[
(c_{v})\mapsto\prod(a,c_{v})_{v}
\]
(因为 Hilbert 符号与局部 Artin 映射之间的关系). 于是, 找 $a$ 等价于找同态 $f\colon\mathbb{I}/\mathbb{I}^{2}\to\mu_2$ 使得

(a) $f=1$ 在主 idèle 上;

(b) $f(1,\ldots,1,i_{v}(b),1,\ldots,1)=
\begin{cases}
-1 & \text{对 } v\in T,\\
1 & \text{否则,}
\end{cases}$

其中 $i_v$ 是包含 $K\hookrightarrow K_v$. 对每个 $v$, 设 $B_v$ 是由 $i_v(b)$ 生成的 $K_v^{\times}/K_v^{\times 2}$ 的 $\F_2$-子空间, 并设 $B=\prod B_v\subset\mathbb{I}/\mathbb{I}^{2}$. 由于 $i_v(b)$ 对 $v\in T$ 不是平方, 存在 (自动连续的) 线性型 $f_1\colon B\to\mu_2$ 满足条件 (b), 而 $f_1$ 延拓为 $\mathbb{I}/\mathbb{I}^{2}$ 上满足 (a) 的连续线性型当且仅当 $f_1$ 在 $B$ 中每个主 idèle 上取值 $1$. $f_1$ 在 $b$ 的主 idèle 上的值为 $\prod_{v\in T}(-1)$, 由 $T$ 含偶数个元素的假设它等于 $1$. 设 $c\in K^{\times}$ 使其 idèle 属于 $B$. 则对每个 $v$, $i_v(c)=1$ 或 $i_v(b)$ (在 $K_v^{\times}/K_v^{\times 2}$ 中). 因此, $i_v(c)$ 对所有 $v$ 在 $K_v[\sqrt{b}]$ 中成为平方, 而这 (由 \ref{thm:viii.1.1}) 蕴含 $c$ 是 $K[\sqrt{b}]$ 中的平方. 故在 $K^{\times}/K^{\times 2}$ 中 $c=1$ 或 $b$, 从而 $f_1$ 在其 idèle 上取值 $1$.
\end{proof}

现在证明定理的 $n=2$ 情形. 设对所有 $v$ 给定二次型 $q(v)=a(v)X^{2}+a(v)d(v)Y^{2}$, 我们寻找 $q=aX^{2}+ad_{0}Y^{2}$ 使 $q\sim q(v)$ 对所有 $v$ 在 $K_v$ 上成立. 即找 $a\in K^{\times}$ 使
\[
S_{v}(q)\overset{\text{def}}{=}(a,a)_{v}(ad_{0},d_{0})_{v}=S(q(v))
\]
对所有 $v$ 成立. 现在
\[
(a,a)_{v}(ad_{0},d_{0})_{v}=(a,-1)_{v}(a,d_{0})_{v}(-1,d_{0})_{v}=(a,-d_{0})_{v}(-1,d_{0})_{v}.
\]
把引理应用于 $T$ 为满足 $S(q(v))(-1,d_{0})_{v}=-1$ 的素之集, $b=-d_0$. 由定理的条件 (b), $T$ 是有限集, 且它有偶数个元素, 因为
\[
\prod_{v}S(q(v))(-1,d_{0})_{v}=\prod_{v}S(q(v))\cdot\prod_{v}(-1,d_{0})_{v}=1\times 1=1.
\]
再者,
\[
S(q(v))\cdot(-1,d_{0})_{v}=(a(v),-d(v))_{v},
\]
故当 $v\in T$ 时 $-i_v(d_0)=-d(v)\ne 1$ (在 $K_v^{\times}/K_v^{\times 2}$ 中). 于是引理给出了所需的元素 $a$.

接着证明 $n=3$ 情形. 对型 $q=q_1+aZ^{2}$ ($a\in K^{\times}$),
\[
n(q_{1})=n(q)-1,\qquad d(q_{1})=a\cdot d(q),\qquad S_{v}(q_{1})=(a,d(q))_{v}\cdot S_{v}(q).
\]
我们找 $a$ 使不变量 $(2,a\cdot d_v,(a,d_v)_v\cdot s_v)$ 对所有 $v$ 可实现, 即 $i_v(a)d_v=-1\Rightarrow (a,d_v)_v\cdot s_v=1$. 令 $T=\{v\mid s_v\ne 1\}$——由假设它是有限集. 由弱逼近定理 (ANT, 7.20), 存在 $a\in K^{\times}$ 使 $i_v(a)d_v\ne -1$ 对所有 $v\in T$. 此时对 $v\in T$, $d(q_1)\ne -1$, 故 $(2,d_v,s_v)$ 可实现. 对 $v\notin T$, $(a,d_v)\cdot s_v=(a,d_v)$, 而 $i_v(a)d_v=-1$ 蕴含 $(a,d_v)_v=(-d_v,d_v)_v=1$. 因此, 对这样的 $a$, 存在秩为 $2$ 的二次型 $q_1$ 使 $q_1\perp aZ^{2}$ 具有所需的不变量.

$n\ge 4$ 情形用归纳法证明. 若对所有实素 $t_v<n$, 可找到形状为 $q_1(X)\perp Z^{2}$ 且具有正确局部不变量的二次型. 若没有 $t_v=0$, 可找到形状为 $q_1(X)\perp(-Z^{2})$ 且具有正确不变量的二次型. 一般情形下, 用弱逼近定理 (ANT, 7.20) 找一个在 $t_v<n$ 的实素处为正、在 $t_v=0$ 的实素处为负的元素 $a$. 于是归纳假设使我们能找到 $q_1(X)$ 使 $q_1(X)\perp aZ^{2}$ 具有正确的不变量.

\subsection*{应用}

\begin{proposition}[GAUSS]\label{prop:viii.6.14}
正整数 $n$ 是三个平方数之和当且仅当它不具有形式 $4^{a}(8b-1)$, 其中 $a,b\in\Z$.
\end{proposition}

\begin{proof}
把上述理论应用于二次型 $X_1^{2}+X_2^{2}+X_3^{2}-nZ^{2}$——见 \cite{Serre1970}, 第 IV 章.
\end{proof}

\section{密度定理}\label{sec:VIII.7}

本节始终设 $K$ 是数域.

\begin{theorem}\label{thm:viii.7.1}
对 $K$ 的任意模 $\mathfrak{m}$ 与任意非平凡 Dirichlet 特征 $\chi\colon C_m\to\CC^{\times}$, 有 $L(1,\chi)\ne 0$.
\end{theorem}

\begin{proof}
正如第 \ref{ch:VI} 章末尾所指出的, 一旦有了互反律, 这就由定理 (VI, \ref{thm:vi.4.9}) 的证明推出.
\end{proof}

\begin{theorem}\label{thm:viii.7.2}
设 $\mathfrak{m}$ 是 $K$ 的模, $H$ 是 $\mathfrak{m}$ 的同余子群: $I^{m}\supset H\supset i(K_{m,1})$. 对任意类 $\mathfrak{k}\in I^{m}/H$, $\mathfrak{k}$ 中的素理想集具有 Dirichlet 密度 $1/(I^{m}:H)$.
\end{theorem}

\begin{proof}
把定理\ref{thm:viii.7.1} 与定理 (VI, \ref{thm:vi.4.8}) 结合.
\end{proof}

\begin{corollary}\label{cor:viii.7.3}
设 $L/K$ 是以 $G$ 为 Galois 群的有限 Abel 扩张, $\sigma\in G$. 则在 $L$ 中不分歧且满足 $(\mathfrak{p},L/K)=\sigma$ 的 $K$ 的素理想 $\mathfrak{p}$ 之集具有 Dirichlet 密度 $\frac{1}{[L:K]}$.
\end{corollary}

\begin{proof}
互反律 (V, \ref{thm:v.3.5}) 说: 对某个模 $\mathfrak{m}$ 与某个 $H\supset i(K_{m,1})$, Artin 映射定义同构 $I^{m}/H\to\Gal(L/K)$. 把定理应用于 $\sigma$ 在 $I^{m}/H$ 中的原像即可.
\end{proof}

\begin{theorem}[CHEBOTAREV]\label{thm:viii.7.4}
设 $L$ 是数域 $K$ 的有限 Galois 扩张, Galois 群为 $G$, 并设 $C\subset G$ 在共轭下稳定, 即
\[
x\in C,\ \tau\in G\Rightarrow\tau x\tau^{-1}\in C.
\]
令
\[
T=\{\mathfrak{p}\mid \mathfrak{p} \text{ 在 } L \text{ 中不分歧},\ (\mathfrak{p},L/K)\subset C\}.
\]
则 $T$ 具有 Dirichlet 密度
\[
\delta(T)=\frac{C \text{ 中元素的个数}}{G \text{ 中元素的个数}}.
\]
\end{theorem}

\begin{proof}
只需对 $C$ 是单个元素 $\sigma$ 的共轭类
\[
C=\{\tau\sigma\tau^{-1}\mid\tau\in G\}
\]
的情形证明. 设 $\sigma$ 的阶为 $f$, 并令 $M=L^{\langle\sigma\rangle}$. 则 $L$ 是 $M$ 的 $f$ 次循环扩张, 故 Artin 映射给出同构
\[
C_{\mathfrak{m}}/H\to\langle\sigma\rangle
\]
对 $M$ 的某个模 $\mathfrak{m}$ 与 $H=M^{\times}\cdot\Nm_{L/M}C_{L,\mathfrak{m}}$ 成立. 我们用记号
\[
\mathfrak{P}\,|\,\mathfrak{q}\,|\,\mathfrak{p}
\]
表示 $\mathcal{O}_L$ 的素 $\mathfrak{P}$、$\mathcal{O}_M$ 的素 $\mathfrak{q}$ 与 $\mathcal{O}_K$ 的素 $\mathfrak{p}$. 令 $d=[L:K]=(G:1)$, 并令 $c$ 为 $C$ 的阶. 我们要证
\[
\delta(T)=\frac{c}{d}.
\]
证明中我们忽略 (有限多个) 不与 $\mathfrak{m}$ 互素的素理想.

令
\[
T_{M,\sigma}=\{\mathfrak{q}\subset\mathcal{O}_{M}\mid(\mathfrak{q},L/M)=\sigma,\ f(\mathfrak{q}/\mathfrak{p})=1\}.
\]
Abel 扩张的 Chebotarev 密度定理 (\ref{cor:viii.7.3}) 表明: 满足 $T_{M,\sigma}$ 定义中第一个条件的素集的密度为 $\frac{1}{f}$, 由此 (见 (VI, \ref{prop:vi.4.5})) $T_{M,\sigma}$ 本身的密度为 $\frac{1}{f}$.

令
\[
T_{L,\sigma}=\{\mathfrak{P}\subset\mathcal{O}_{L}\mid(\mathfrak{P},L/K)=\sigma\}.
\]
我们将证明:

(a) 映射 $\mathfrak{P}\mapsto\mathfrak{q}=\mathfrak{P}\cap\mathcal{O}_M$ 定义双射 $T_{L,\sigma}\to T_{M,\sigma}$;

(b) 映射 $\mathfrak{P}\mapsto\mathfrak{p}=\mathfrak{P}\cap\mathcal{O}_K\colon T_{L,\sigma}\to T$ 恰好把 $\frac{d}{cf}$ 个 $T_{L,\sigma}$ 中的素送到 $T$ 的每个素.

把这两个断言结合起来: 映射 $\mathfrak{q}\mapsto\mathfrak{p}=\mathfrak{q}\cap\mathcal{O}_K$ 定义 $\frac{d}{cf}:1$ 的映射 $T_{M,\sigma}\to T$. 对这样的 $\mathfrak{q}$, $\Nm_{M/K}\mathfrak{q}=\mathfrak{p}$, 故 $\mathbb{N}\mathfrak{q}=\mathbb{N}\mathfrak{p}$; 于是
\[
\sum_{\mathfrak{p}\in T}\frac{1}{\mathbb{N}\mathfrak{p}^{s}}=\frac{cf}{d}\sum_{\mathfrak{q}\in T_{M,\sigma}}\frac{1}{\mathbb{N}\mathfrak{q}^{s}}\sim\frac{cf}{d}\cdot\frac{1}{f}\log\frac{1}{s-1}=\frac{c}{d}\log\frac{1}{s-1},
\]
如所欲证. 剩下要证 (a) 与 (b).

设 $\mathfrak{P}\in T_{L,\sigma}$, 并令 $\mathfrak{q}=\mathfrak{P}\cap\mathcal{O}_M$, $\mathfrak{p}=\mathfrak{P}\cap\mathcal{O}_K$. 则 $L_{\mathfrak{P}}/K_{\mathfrak{p}}$ 的 Galois 群由 $\sigma$ 生成, 但 $\sigma$ 固定 $M_{\mathfrak{q}}$, 故 $M_{\mathfrak{q}}=K_{\mathfrak{p}}$. 因此 $f(\mathfrak{q}/\mathfrak{p})=1$, 这表明 $\mathfrak{q}\in T_{M,\sigma}$, 于是有映射
\[
\mathfrak{P}\mapsto\mathfrak{q}\overset{\text{def}}{=}\mathfrak{P}\cap\mathcal{O}_{M}:T_{L,\sigma}\to T_{M,\sigma}.
\]
它是单射, 因为 $f(\mathfrak{P}/\mathfrak{q})=f(\mathfrak{q}/\mathfrak{p})^{-1}f(\mathfrak{P}/\mathfrak{p})=1\times f$, 故 $\mathfrak{P}$ 是 $L$ 中位于 $\mathfrak{q}$ 之上的唯一素. 它是满射, 因为对位于 $\mathfrak{q}\in T_{M,\sigma}$ 之上的任何素 $\mathfrak{P}$,
\[
(\mathfrak{P},L/K)=(\mathfrak{P},L/K)^{f(\mathfrak{q}/\mathfrak{p})}=(\mathfrak{P},L/M)=\sigma
\]
($\mathfrak{q}$ 属于 $T_{M,\sigma}$ 的第一个条件), 故 $\mathfrak{P}\in T_{L,\sigma}$. 这证明了 (a).

固定 $\mathfrak{p}_0\in T$, 并设 $\mathfrak{P}_0\in T_{L,\sigma}$ 位于 $\mathfrak{p}_0$ 之上. 则对 $\tau\in G$,
\[
(\tau\mathfrak{P}_{0},L/K)=\tau(\mathfrak{P}_{0},L/K)\tau^{-1},
\]
故
\[
\tau(\mathfrak{P}_{0},L/K)\tau^{-1}=\sigma\iff\tau\in C_{G}(\sigma)
\]
($\sigma$ 在 $G$ 中的中心化子). 因此映射 $\tau\mapsto\tau\mathfrak{P}_0$ 是双射
\[
C_{G}(\sigma)/G(\mathfrak{P}_{0})\to\{\mathfrak{P}\in T_{L,\sigma}\mid\mathfrak{P}\cap\mathcal{O}_{K}=\mathfrak{p}_{0}\}.
\]
分解群 $G(\mathfrak{P}_0)$ 等于 $\langle\sigma\rangle$, 其阶为 $f$; 而 $C_G(\sigma)$ 的阶为 $\frac{d}{c}$, 因为存在双射
\[
\tau\mapsto\tau\sigma\tau^{-1}:G/C_{G}(\sigma)\to C.
\]
因此 $(C_G(\sigma):G(\mathfrak{P}_0))=\frac{d}{cf}$, 于是我们证明了: 对每个 $\mathfrak{p}\in T$, 恰有 $\frac{d}{cf}$ 个素 $\mathfrak{P}\in T_{L,\sigma}$ 位于 $\mathfrak{p}$ 之上. 这证明了 (b), 从而完成了定理的证明.
\end{proof}

\begin{remark}\label{rem:viii.7.5}
关于 Chebotarev 密度定理的有效形式, 见 Lagarias 与 Odlyzko (Algebraic Number Fields, Ed.~Fr\"ohlich, 1977). 设 $L$ 是 $K$ 的有限 Galois 扩张, 并令
\[
\pi_{C}(x)=|\{\mathfrak{p}\mid(\mathfrak{p},L/K)\subset C,\ \mathbb{N}\mathfrak{p}\le x\}|.
\]
则
\[
\pi_{C}(x)=\frac{c}{d}\cdot\frac{x}{\log x}+\text{具体的误差项}.
\]
\end{remark}

\section{函数域; 几何类域论}\label{sec:VIII.8}

我们还应当包括函数域 (某个 $\F_p(T)$ 的有限扩张) 的类域论. 为此, 既可以模仿数域情形的证明 (见 \cite{ArtinTate1961}), 也可以 (更好) 把证明建立在 Tsen 定理之上 (见 \cite{Milne2006}, 第 I 章附录).

设 $K$ 是域 $k$ 上一个变量的函数域, 即 $k(T)$ 的有限扩张. 几何类域论描述 $K$ 的``几何''Abel 扩张, 即``不来自 $k$''的 Abel 扩张. 关于几何类域论的文献列表, 见 \mo{73054}.

\section{数域的上同调}\label{sec:VIII.9}

我们还应当包括 Nakayama 与 Tate 的定理的证明 (例如 J.~Tate, The cohomology groups of tori in finite Galois extensions of number fields, \emph{Nagoya Math.~J.}, 27, 1966, 709--719) 以及 Poitou 与 Tate 的定理的证明 (例如 \cite{Milne2006}, 第 I 章).

\section{再谈 \texorpdfstring{$L$}{L}-级数}\label{sec:VIII.10}

设 $\chi$ 是一个 Dirichlet $L$-级数对应的特征. 则存在常数 $A(\chi)\ge 0$, $a(\chi)$, $b(\chi)\in\CC$, 使得
\[
\Phi(s,\chi)\overset{\text{def}}{=}A(\chi)^{s}\,\Gamma\Big(\frac{s}{2}\Big)^{a(\chi)}\Gamma\Big(\frac{s+1}{2}\Big)^{b(\chi)}L(s,\chi)
\]
满足函数方程
\[
\Phi(s,\chi)=W(\chi)\,\Phi(1-s,\bar{\chi}),\qquad |W(\chi)|=1.
\]
见 \cite{Narkiewicz1990}.

\subsection*{Artin $L$-级数}

设 $L$ 是 $K$ 的以 $G$ 为 Galois 群的有限 Galois 扩张. 设 $V$ 是 $\CC$ 上的有限维向量空间, 并设
\[
\rho\colon G\to\mathrm{Aut}_{\CC}(V)
\]
是 $G$ 到 $V$ 的线性自同构群的同态. 我们称 $\rho$ 为 $G$ 在 $\CC$ 上的一个 (有限维) 表示. $\rho$ 的迹是映射
\[
\sigma\mapsto\chi(\sigma)=\mathrm{Tr}(\rho(\sigma)).
\]
(回顾: $m\times m$ 矩阵 $(a_{ij})$ 的迹是 $\sum a_{ii}$, 自同态的迹是其相对任一组基的矩阵的迹.) 对 $\sigma\in G$, 令
\[
P_{\sigma}(T)=\det(1-\rho(\sigma)T\mid V)=\prod_{i=1}^{\dim V}(1-a_{i}T),\qquad a_{i}\in\CC
\]
为 $\rho(\sigma)$ 的特征多项式. 注意 $P_{\sigma}(T)$ 只依赖于 $\sigma$ 的共轭类, 故对 $K$ 的每个在 $L$ 中不分歧的素 $\mathfrak{p}$, 可定义
\[
P_{\mathfrak{p}}(T)=P_{\sigma}(T),\qquad \sigma=(\mathfrak{P},L/K),\ \mathfrak{P}|\mathfrak{p} \text{ 某个}.
\]
对这样的 $\mathfrak{p}$, 令
\[
L_{\mathfrak{p}}(s,\rho)=\frac{1}{P_{\mathfrak{p}}(\mathbb{N}\mathfrak{p}^{-s})}
\]
并令
\[
L(s,\rho)=\prod L_{\mathfrak{p}}(s,\rho).
\]
例如, 若 $L/K$ 是 Abel 的, 则表示可对角化\footnote{我们的意思是: 相对于 $V$ 的某组基,
\[
\rho(g)=\begin{pmatrix}
\chi_{1}(g) & 0 & \cdots & 0\\
0 & \chi_{2}(g) & \cdots & 0\\
\vdots & \vdots & \ddots & \vdots\\
0 & 0 & \cdots & \chi_{m}(g)
\end{pmatrix},\qquad g\in G.
\]}
\[
\rho\approx\chi_{1}\oplus\cdots\oplus\chi_{m},
\]
其中每个 $\chi_i$ 都是同态 $G\to\CC^{\times}$. 与 Artin 映射
\[
C_{m}\to G
\]
复合之后, $\chi_i$ 成为 Dirichlet 特征, 于是 Artin $L$-级数等同于 Dirichlet $L$-级数之积. 这正是 Artin 当初定义他的映射的原因.

可以证明: 若 $(V,\rho)=\Ind_{H}^{G}(V_{0},\rho_{0})$, 则
\[
L(s,\rho)=L(s,\rho_{0}).
\]
为处理更一般的 Artin $L$-级数, Artin 证明了: 有限群 $G$ 的每个特征都是循环子群的诱导特征的 $\Q$-线性组合. 故
\[
L(s,\rho)=\prod(\text{Dirichlet } L \text{-级数})^{r_{i}},\qquad r_{i}\in\Q.
\]
后来 Brauer 证明了一个更强的定理, 由此可以证明
\[
L(s,\rho)=\prod(\text{Dirichlet } L \text{-级数})^{r_{i}},\qquad r_{i}\in\Z.
\]
Artin 猜想: 只要 $\rho$ 不包含平凡表示, $L(s,\rho)$ 就延拓为整个复平面上的全纯函数. 最后一个公式蕴含: 当诸 $r_i$ 都是正整数时这是对的. 这一猜想在 Langlands 于 $V$ 的维数为 $2$ 的许多情形下证明它之前进展甚微 (见 R.~Langlands, \emph{Base Change for GL(2)}, Princeton, 1980).

\subsection*{Hecke $L$-级数}

Hecke (或 Gr\"ossen) 特征是指从 $\mathbb{I}$ 到 $\CC^{\times}$ 的单位圆的连续同态 $\psi$, 满足 $\psi(K^{\times})=1$, 且对某个有限集 $S$, $\psi$ 在集合 $\{(a_{v})\mid a_{v}=1\ (v\in S),\ a_{v}\ \text{为单位}\ (\text{对所有 } v)\}$ 上取值 $1$.

\begin{example}\label{ex:viii.10.1}
设 $D\in\Z$ 无立方因子, $\zeta$ 为本原三次单位根. 若 $p\equiv 2 \bmod 3$, 则 $p$ 在 $\Q[\zeta]$ 中保持素, 我们令 $\psi(p)=1$. 若 $p\equiv 1 \bmod 3$, 则 $p=\pi\bar{\pi}$ 分解于 $\Q[\zeta]$ 中, 我们取 $\pi$ 满足 $\pi\equiv 1 \bmod 3$. 则 $\pi=\frac{1}{2}(a+3b\sqrt{-3})$, 且 $4p=4\pi\bar{\pi}=a^{2}+27b^{2}$. 现在存在 Hecke 特征使得: 对所有奇数且 $\not\equiv 1 \bmod 3$ 的 $p$, $\psi(p)=1$; 而当 $p\equiv 1\bmod 3$ 时 $\psi(p)=\frac{\pi}{\sqrt{p}}\left(\frac{D}{\pi}\right)$.
\end{example}

对这样的特征, 定义
\[
L(s,\psi)=\prod_{v\notin S}\frac{1}{1-\psi(\pi_{v})\,\mathbb{N}\mathfrak{p}_{v}^{-s}},
\]
其中 $\pi_v$ 是在第 $v$ 个位置为素元、其余位置为 $1$ 的 idèle. Hecke $L$-级数的基本解析性质 (亚纯延拓、函数方程等) 已被充分理解 (例如 J.~Tate, \emph{Fourier analysis in number fields and Hecke's zeta function}, 博士论文, 1950; 转载于 \cite{CasselsFrohlich1967}; 或 \cite{RamakrishnanValenza1999}).

\subsection*{Weil 群与 Artin--Hecke $L$-级数}

关于这一论题, 见 \cite{Tate1979}.

\section*{Gauss 的一个定理}\label{sec:VIII.gauss}

本课程以 Gauss 首先证明的一个定理——二次互反律——开始, 以他的另一个定理结束似乎是恰当的.

考虑椭圆曲线 $E\colon X^{3}+Y^{3}+Z^{3}=0$. 设 $N_p$ 为 $E$ 上坐标在 $\F_p$ 中的点的个数. Gauss 证明了:

(a) 若 $p\not\equiv 1 \bmod 3$, 则 $N_p=p+1$;

(b) 若 $p\equiv 1 \bmod 3$, 则 $N_p=A$, 其中 $A$ 是满足 $4p=A^{2}+27B^{2}$ 且 $A\equiv 1 \bmod 3$ 的唯一整数.

见 \cite{SilvermanTate1992}, IV.2.

Gauss 定理蕴含 Weil 猜想对 $E/\F_p$ 成立, 即
\[
|N_{p}-p-1|<2\sqrt{p}.
\]
它还蕴含 Taniyama 猜想对 $E/\Q$ 成立, 因为它表明 $L$-级数 $L(s,E)$ 等于 $L(s-\frac{1}{2},\psi)$, 其中 $\psi$ 是上例中与 $D=-1$ 相关联的 Hecke 特征.

\begin{quote}
他写信告诉我, 代数数论是他所遇到过的最美的数学主题, 而他苦难中的唯一慰藉就是讲授类域论……这确实是他最钦佩的那类数学: 主要结果范围广阔、美感十足, 但证明在技术上极其艰难.

A.~Borel, 载于: \emph{Current Trends in Mathematics and Physics: A Tribute to Harish-Chandra}, 编辑 S.D.~Adhikari, Narosa Publishing House, New Delhi, Madras, 1995, p.~213.
\end{quote}


% ===== VIII 附录 A/B/C (练习、解答、历史文献来源) =====

\cleardoublepage\phantomsection
\section*{A 附录: 练习}\label{sec:VIII.A}
\addcontentsline{toc}{chapter}{A 附录: 练习}

这些是学员需要提交的练习.

\medskip\noindent\textbf{A-1}\quad
(a) 设 $K$ 是 $\Q$ 的有限 Galois 扩张. 证明 $\Gal(K/\Q)$ 由 $\mathcal{O}_K$ 中素理想的惯性群生成. (提示: 证明由这些惯性群生成的子群的固定域在 $\Z$ 的所有素理想处均不分歧. 素理想 $\mathfrak{P}$ 的惯性群是指分解群 $D(\mathfrak{P})$ 中固定剩余域的那些元素组成的子群; 它的阶等于 $\mathfrak{P}$ 的分歧指数.)

(b) 设 $p_1,\ldots,p_t$ 为互不相同的奇素数, 并令 $p_i^{*}=(-1)^{\frac{p_i-1}{2}}p_i$ (于是只有 $p_i$ 在 $\Q[\sqrt{p_i^{*}}]$ 中分歧). 设 $K=\Q[\sqrt{d}]$, $d=\prod p_i^{*}$.

设 $C_+$ 为 $K$ 的窄类群, 并设 $L$ 为 $C_+^{2}$ 的类域 (于是 $L$ 在 $\mathcal{O}_K$ 的所有素理想处不分歧, 且 $\Gal(L/K)$ 是初等 Abel $2$-群; 此外, $L$ 是使这些陈述成立的最大的域).

证明 $L$ 在 $\Q$ 上是 Galois 的, 并且 $\mathcal{O}_L$ 的每个素理想在 $\Gal(L/\Q)$ 中的惯性群的阶为 $1$ 或 $2$. 由此推出 $\Gal(L/\Q)$ 是初等 Abel $2$-群, 且 $L=\Q[\sqrt{p_1^{*}},\ldots,\sqrt{p_t^{*}}]$. 再由此推出 $(C_+:C_+^{2})=2^{t-1}$ (对照例 \ref{ex:0.7}).

\medskip\noindent\textbf{A-2}\quad
(a) 设 $d>3$ 是满足 $d\equiv 3 \bmod 8$ 的无平方因子整数. 证明存在 $\Q$ 上的三次域 $L$, 其判别式为 $-d$ 或 $-4d$.

(b) 求使 $|\Delta_{L/K}|$ 最小的 $\Q$ 上三次扩张 $L$.

(你可以引用类域论的主要结论.)

\textbf{提示:} 令 $K=\Q[\sqrt{-d}]$. 对于 $3$ 整除 $K$ 的类数的情形, 琢磨课上讲过的 $\Q[\sqrt{-31}]$ 的 Hilbert 类域这个例子. 对于相反的情形, 证明 $2\mathcal{O}_K$ 是一个素理想 $\mathfrak{p}$, 且 $3$ 整除射线类群 $C_{\mathfrak{p}}$ 的阶.

\medskip\noindent\textbf{A-3}\quad 完成练习 \ref{exc:v.3.15} (本书 p.\pageref{exc:v.3.15}). 注意, 在该练习的陈述中, $i$ 被含混地用来同时表示从 $K_{\mathfrak{m},1}$ 到 $I^{S(\mathfrak{m})}$ 与到 $I^{S\cup S(\mathfrak{m})}$ 的 (显然的) 同态.

\medskip\noindent\textbf{A-4}\quad 证明或反驳: 数域 $K$ 的 id\`ele 类群 $\mathbf{C}_K=\mathbb{I}_K/K^{\times}$ 的每个有限指数子群都是开的. (这是练习 \ref{exc:v.5.11} (本书 p.\pageref{exc:v.5.11}). 对 $K_v^{\times}$ 的类似陈述可以引用 —— 见 I \ref{para:1.7}.)

\medskip\noindent\textbf{A-5}\quad 设 $G$ 是一个群, $A$ 是一个带 $G$-作用的群 (不必交换), 即给定一个同态 $G\to\mathrm{Aut}(A)$. 与交换情形一样, 交叉同态 $\varphi\colon G\to A$ 由条件 $\varphi(\sigma\tau)=\varphi(\sigma)\cdot\sigma\varphi(\tau)$ 定义. 若存在 $a\in A$ 使得 $\varphi'(\sigma)=a^{-1}\cdot\varphi(\sigma)\cdot\sigma a$ 对所有 $\sigma\in G$ 成立, 则称两个这样的同态 $\varphi$ 与 $\varphi'$ 等价, 等价类之集记作 $H^{1}(G,A)$ (它是一个带特殊元素之集). 与 Hilbert 定理 90 的证明相类似的论证表明: 当 $L/K$ 是以 $G$ 为 Galois 群的有限 Galois 扩张时, $H^{1}(G,\mathrm{GL}_n(L))=0$.

设 $K$ 是特征 $\ne 2$ 的域. $K$ 上的\textbf{二次空间}是二元组 $(V,\Phi)$, 由 $K$ 上有限维向量空间 $V$ 与 $V$ 上非退化二次型 $\Phi$ 组成, 即 $\Phi$ 是由 $V$ 上非退化对称双线性型 $\varphi$ 伴随的映射 $x\mapsto\varphi(x,x)$. \textbf{等距} $\alpha\colon(V,\Phi)\to(V',\Phi')$ 是向量空间的同构 $\alpha\colon V\to V'$, 满足 $\Phi'(\alpha(v))=\Phi(v)$ 对所有 $v\in V$ 成立.

\textbf{练习:} 设 $L/K$ 是特征 $\ne2$ 的域的有限 Galois 扩张, 并令 $G=\Gal(L/K)$. 取定 $K$ 上的一个二次空间 $(V,\Phi)$, 并令 $O(\Phi)$ 为 $(V,\Phi)\otimes_K L$ 的自同构群 —— 它带有显然的 $G$-作用. 若 $(V',\Phi')$ 是 $K$ 上的二次空间, 且 $\alpha\colon(V,\Phi)\otimes_K L\to(V',\Phi')\otimes_K L$ 是 $L$ 上二次空间的同构, 证明 $\sigma\mapsto\alpha^{-1}\circ\sigma\alpha$ 是交叉同态 $\varphi\colon G\to O(\Phi)$, 其等价类不依赖于 $\alpha$ 的选取. 由此推出 $H^{1}(G,O(\Phi))$ 对在 $L$ 上变得与 $(V,\Phi)$ 同构的 $K$ 上二次空间的同构类之集作出分类 (即 $H^{1}(G,O(\Phi))$ 与该集之间存在一一对应). 再由此推出: 当 $K=\R$, $L=\CC$ 时 (设 $\dim V\ne0$), $H^{1}(G,O(\Phi))\ne0$. (你可以引用 $H^{1}(G,\mathrm{GL}_n(L))=0$.)

\medskip\noindent\textbf{A-6}\quad 设 $H$ 是群 $G$ 中有限指数的子群, 并选取映射 $s\colon H\backslash G\to G$ 使 $H\,s(x)=Hx$, 即为 $G$ 中 $H$ 的每个右陪集取一个代表元 —— 于是 $G=\bigcup_{x\in H\backslash G}H\,s(x)$.

(a) 对 $g\in G$, 证明 $s(x)\cdot g\cdot s(xg)^{-1}\in H$.

对 $g\in G$, 定义
\[
V(g)=\prod_{x\in H\backslash G}s(x)\cdot g\cdot s(xg)^{-1}\ \bmod\ H^{c},
\]
其中 $H^{c}$ 是 $H$ 的换位子群. (注意, 由于 $H/H^{c}$ 交换, $V(g)$ 不依赖于乘积中 $H\backslash G$ 的取序方式.)

(b) 证明 $V(g)$ 不依赖于 $s$ 的选取, 并且它是同态 $G\to H/H^{c}$.

于是 $V$ 诱导同态 $\mathrm{Ver}\colon G/G^{c}\to H/H^{c}$ —— 即 \textbf{Verlagerung} (或\textbf{转移}) 映射.

(c) 对任意 $G$-模 $M$, 证明映射 $m\mapsto N(m)=\sum_{x\in H\backslash G}s(x)\cdot m$ 诱导良定义的同态 $\mathrm{Res}\colon H_0(G,M)\to H_0(H,M)$.

(d) 证明图表
\[
\begin{tikzcd}
G/G^{c} \ar[r, "\simeq"] \ar[d, "\mathrm{Ver}"'] & I_G/I_G^{2} \ar[d, "N"] \\
H/H^{c} \ar[r, "\simeq"'] & I_G/I_H\, I_G \\
H/H^{c} \ar[r, "\simeq"'] & I_H/I_H^{2} \ar[u]
\end{tikzcd}
\]
交换 (水平的同构就是引理 \ref{lem:ii.2.6} 的同构; 向上的映射由包含映射诱导).

[假设 $\mathrm{Res}\colon H_r(G,M)\to H_r(H,M)$ 存在, 与边缘同态相容, 并且在 $r=0$ 时有 (c) 中所给的描述, 则上式表明它们在 $H_1(G,\Z)$ 上与 Verlagerung 一致. 更精确地, 上面的图表可以等同于
\[
\begin{tikzcd}
H_1(G,\Z) \ar[r, "\delta"] \ar[d, "\mathrm{Res}"'] & H_0(G, I_G) \ar[d, "\mathrm{Res}"] \\
H_1(H,\Z) \ar[r, "\delta"] & H_0(H, I_G) \\
H_1(H,\Z) \ar[r, "\delta"'] \ar[u, "\mathrm{id}"'] & H_0(H, I_H) \ar[u]
\end{tikzcd}
\]
边缘同态来自 $G$ 与 $H$ 的增广正合列. 上、下两个边缘同态是同构. 中间的边缘同态是单射, 并且当 $G$ 有限、而 $H_0(H,I_G)$ 换成 $\widehat{H}_T^{-1}(H,I_G)$ 时, 它成为同构.]

\medskip\noindent\textbf{A-7}\quad 本练习将表明上同调与逆向极限不交换 —— 事实上, 映射 $H^{1}(G,\varprojlim M_n)\to\varprojlim H^{1}(G,M_n)$ 不必是单射.

设 $G$ 是一个 profinite 群, 且
\[
\cdots\longrightarrow M_n\overset{\pi_{n-1,n}}{\longrightarrow}M_{n-1}\longrightarrow\cdots\longrightarrow M_1
\]
是一个 $G$-模的投影系统. 假设:

(a) 每个映射 $\pi_{n-1,n}\colon M_n\to M_{n-1}$ 都是满射;

(b) 存在交换图表
\[
\begin{tikzcd}[row sep=small]
\cdots \ar[r] & M_n^{G} \ar[r, "\pi_{n-1,n}"] \ar[d, "\sigma_n"', two heads] & M_{n-1}^{G} \ar[r] \ar[d, "\sigma_{n-1}"', two heads] & \cdots \\
\cdots \ar[r] & \Z \ar[r] & \Z \ar[r] & \cdots
\end{tikzcd}
\]
其中每个竖直映射都是满射.

对任意 $c_n\in\Z$, 级数 $\sum p^{n-1}c_n$ 将在 $\Z_p$ 中收敛 —— 选取序列 $c_n$, $n\ge1$, 使得 $\sum p^{n-1}c_n$ 收敛到 $\Z_p$ 中不属于 $\Z$ 的元素 $c$. 对每个 $n$, 选取元素 $x_{n,1},\ldots,x_{n,n}$ 使得:

(a) 对 $1\le i\le n-1$, $\pi_{n-1,n}(x_{n,i})=x_{n-1,i}$;

(b) $x_{n,n}\in M_n^{G}$ 且 $\sigma_n(x_{n,n})=c_n$.

定义
\[
x_n=x_{n,1}+\cdots+x_{n,n},\qquad
\varphi_n(\sigma)=\sigma x_n-x_n\quad(\sigma\in G).
\]
于是 $\varphi_n$ 是主交叉同态 $G\to M_n$. 注意 $\pi_{n-1,n}(\varphi_n(\sigma))=\varphi_{n-1}(\sigma)$ 对所有 $n$ 成立, 从而存在元素 $\varphi(\sigma)\in M\overset{\mathrm{def}}{=}\varprojlim M_n$, 它在每个 $M_n$ 中的像为 $\varphi_n(\sigma)$. 显然 $\sigma\mapsto\varphi(\sigma)$ 是交叉同态 $G\to M$.

证明 $\varphi$ 不是主交叉同态, 即证明不存在 $M$ 中的元素 $y=(y_n)$ 使得对每个 $n$ 与每个 $\sigma\in G$, 都有 $\varphi_n(\sigma)=\sigma y_n-y_n$.

提示: 令 $y_n=x_n+z_n$, 并证明 $z_1$ 被 $G$ 固定, 但 $\sigma_1(z_1)=c\notin\Z$.

你能构造出满足这些条件的系统 $(M_n,\pi_{n-1,n},\sigma_n)$ 吗?

\medskip\noindent\textbf{A-8}\quad 完成练习 \ref{ex:iv.5.7} 的 (a)--(d) 部分 (若你精力充沛, 也可以完成整个练习). 必要时可以使用第 \ref{ch:IV} 章的非上同调结果.

\medskip\noindent\textbf{A-9}\quad 设 $L_n=\Q_p[\zeta_{p^n}]$, 其中 $\zeta_{p^n}$ 是本原 $p^n$ 次单位根. 证明 $p$ 是来自 $L_n$ 的范数. 由此推出 $L_\infty=\bigcup L_n$ 是 $\Q_p^{\mathrm{ab}}$ 中被 $\mathrm{rec}_{\Q_p}(p)$ 固定的子域 (并在 $\Q_p$ 的某个取定的代数闭包内取并).

\medskip\noindent\phantomsection\hypertarget{item:viii.A-10}{}\textbf{A-10}\quad 设 $L/K$ 是 $n$ 次有限循环数域扩张. 证明: 对每个 $d\mid n$, 满足 $e(\mathfrak{P}/\mathfrak{p})=1$ 且 $f(\mathfrak{P}/\mathfrak{p})=d$ (对某个 —— 因而是所有 —— 位于 $\mathfrak{p}$ 之上的 $\mathfrak{P}$) 的 $K$ 的素 $\mathfrak{p}$ 之集有极密度 $\varphi(d)/n$. (提示: 尝试对 $d$ 作归纳.)

精力充沛者可以类似地给出 Frobenius 密度定理 (Chebotarev 密度定理的较弱、更早的形式) 的一个同样初等的证明: 设 $L/K$ 是以 $G$ 为 Galois 群的 Galois 扩张, 并设 $\sigma$ 是 $G$ 中 $n$ 阶元素. 证明
\[
\{\mathfrak{p}\mid \mathfrak{p}\ \text{不分歧且}\ (\mathfrak{p},L/K)=\sigma^{k}\ \text{对某个与}\ n\ \text{互素的}\ k\}
\]
有极密度 $c\varphi(n)/[L:K]$, 其中 $c$ 是 $\langle\sigma\rangle$ 在 $G$ 中的共轭个数. 再证明 $c\varphi(n)$ 恰是 $\sigma$ 的共轭个数, 从而这与 Chebotarev 密度定理相容. 至于提示, 见 \cite[pp.~207--208]{Marcus1977}.

\medskip\noindent\textbf{A-11}\quad 证明 $16$ 在 $\R$ 中以及在所有奇数 $p$ 的 $\Q_p$ 中是 $8$ 次幂 (但在 $\Q$ 中不是).

[提示: 证明 $\Q[\zeta_8]=\Q[i,\sqrt{2}]$ 在所有奇数 $p$ 处不分歧, 并由此推出: 对奇数 $p$, $\Q_p$ 包含 $1+i$, $\sqrt{2}$, $\sqrt{-2}$ 三者之一.]

注: 这与 Grunwald 定理相矛盾, 而该定理曾被广泛使用了十五年.

\cleardoublepage\phantomsection
\section*{B 附录: 练习解答}\label{sec:VIII.B}
\addcontentsline{toc}{chapter}{B 附录: 练习解答}

\noindent\textbf{练习 A-1 的解答.}

(a) 设 $\mathfrak{p}$ 是 $\mathcal{O}_K$ 中的素理想, 并设 $I(\mathfrak{p})$ 是它的惯性群. 则 $M=K^{I(\mathfrak{p})}$ 具有如下性质: $\mathfrak{p}\cap\mathcal{O}_M$ 在 $\mathfrak{p}\cap\Z$ 上不分歧 (ANT, 8.11). 于是被所有惯性群固定的任何域都在 $\Q$ 上不分歧, 从而等于 $\Q$ (ANT, 4.9).

(b) 域 $L$ 在 $\Q$ 上是 Galois 的, 因为它在它的 Galois 闭包的任何自同构下稳定.

显然惯性群的阶为 $1$ 或 $2$, 于是 $\Gal(L/\Q)$ 由 $2$ 阶元素生成. 这并不意味着它是交换的 (每个二面体群都由两个 $2$ 阶元素生成), 但存在扩张
\[
0\to\Gal(L/K)\to\Gal(L/\Q)\to\Gal(K/\Q)\to 0,
\]
其中 $\Gal(L/K)\simeq(\Z/2\Z)^{s}$ ($s$ 为某个数), $\Gal(K/\Q)\simeq\Z/2\Z$, 而非平凡惯性群的生成元在 $\Gal(K/\Q)$ 中的像 $\ne 1$ (否则 (a) 中的论证将表明分歧发生在 $L/K$ 而非 $K/\Q$ 中). 设 $\sigma$ 与 $\tau$ 生成两个 $2$ 阶惯性群. 则 $\sigma|_K\ne1$ 且 $\tau|_K\ne1$, 于是 $\sigma\tau|_K=1$. 从而 $\sigma\tau\in\Gal(L/K)$, 故其阶为 $1$ 或 $2$. 因此 $\sigma$ 与 $\tau$ 交换, 于是 $\Gal(L/\Q)$ 由互相交换的 $2$ 阶元素生成, 这表明它是初等 Abel $2$-群.

在 $L$ 中只有素 $p_1,\ldots,p_t$ 分歧, 且对应于每个 $p_i$ 只有一个惯性群 (在 $p_i$ 之上可能有许多 $L$ 的素理想, 但由于 $\Gal(L/\Q)$ 交换, 每个的惯性群都相同). 于是 $\Gal(L/\Q)\simeq(\Z/2\Z)^{s}$, $s\le t$. 由于 $L$ 包含 $\Q[\sqrt{p_1^{*}},\ldots]$, 而后者在 $\Q$ 上的次数为 $2^{t}$, 可见 $L=\Q[\sqrt{p_1^{*}},\ldots]$. 其余是容易的.

[关于 $L=\Q[\sqrt{p_1^{*}},\ldots]$ 的另一种证明: 由于 $\Gal(L/\Q)$ 是初等 Abel $2$-群, $L=\Q[\sqrt{m_1},\sqrt{m_2},\ldots]$, 其中 $m_i\in\Z$ (容易的 Kummer 理论); 但容易看出, $\Q[\sqrt{p_1^{*}},\ldots]$ 是这类域中其与 $K$ 的复合在 $K$ 上不分歧的最大者.]

\noindent\textbf{练习 A-2 的解答.}

由于 $d\equiv 3 \bmod 8$, 有 $-d\equiv 1 \bmod 4$, 从而 $\Delta_{K/\Q}=-d$. 我用``素''一词指``有限素''.

若 $3$ 整除 $K$ 的类数, 则存在 $K$ 的 $3$ 次不分歧扩张 $L$ (由类域论). 与 $d=31$ 的情形一样, $L$ 在 $\Q$ 上是 Galois 的, 且 $L=M\cdot K$, 其中 $M$ 是 $\Q$ 上 $3$ 次且非 Galois 的域. 由 Galois 理论, 含于 $M$ 的 Galois 闭包中的 $\Q$ 的唯一二次扩张是 $\Q[\sqrt{\Delta_{M/\Q}}]$, 于是 $\Delta_{M/\Q}=-d\times$平方 (容易的 Kummer 理论). 若素数 $p$ 不整除 $d$, 则 $p$ 在 $L$ (从而在 $M$) 中不分歧, 于是 $p$ 不整除 $\Delta_{M/K}$. 若 $p\mid d$, 则 $p\mathcal{O}_K=\mathfrak{p}_1^{2}\mathfrak{p}_2$, 这 (差理论) 表明 $p$ 整除 $\Delta_{M/K}$ 但不整除 $p^{2}$. 因此 $\Delta_{M/\Q}=-d$.

假设 $3$ 不整除类数. 由于 $-d\equiv 5 \bmod 8$, $(2)$ 在 $\mathcal{O}_K$ 中保持素\footnote{若 $-d\equiv 1 \bmod 8$, 它就会分裂.}, 记 $2\mathcal{O}_K=\mathfrak{p}$, 且 $3$ 整除射线类群 $C_{\mathfrak{p}}$ 的阶 (射线类群包含 $(\mathcal{O}_K/\mathfrak{p})^{\times}=\mathbb{F}_4^{\times}$ —— 见 V \ref{thm:v.1.7}; 这里我们利用 $d>3$ 来看出 $K$ 的单位只有 $\pm1$). 我们用与前面相同的论证, 取 $L$ 为 $K$ 的仅在 $\mathfrak{p}$ 处分歧的 $3$ 次扩张. 同样 $L=M\cdot K$, 其中 $M$ 与前面一样, 只是 $2$ 在 $M$ 中全分歧, 论证表明 $\Delta_{M/\Q}=-4d$.

从本书 p.\pageref{exc:v.3.14} 的表可见 $\Delta_{M/\Q}=-23$ 是可能的. 由 Minkowski 界, $|\Delta_{M/\Q}|\ge 13$ (有复素) 或 $|\Delta_{M/\Q}|\ge 21$ (无复素). Stickelberger 定理排除了其余情形, 只留下 $-15,-16,-19,-20,21$. 若 $M$ 的判别式 $d=-15,-19,21$, 则其 Galois 闭包 $L=M\cdot\Q[\sqrt{d}]$ 在 $\Q[\sqrt{d}]$ 上不分歧 (否则 $\Q$ 的某个素会在 $L$ 中全分歧而在 $M$ 中不然), 但 $\Q[\sqrt{d}]$ 的类数分别为 $2,1,1$. 对 $d=-16,-20$, 同样的论证表明只有位于 $2$ 之上的素可能在 $L=M\cdot\Q[\sqrt{d}]$ 中分歧, 但此处类数为 $1$ 或 $2$, 且相关的射线类数不被 $3$ 整除 ($2$ 在 $\Q[\sqrt{-1}]$ 与 $\Q[\sqrt{-5}]$ 中分歧).

\textbf{注:}

(a) 可以不用差理论, 而使用 ANT 3.39(a) 中的公式 (但该公式的证明用到了差理论).

(b) 用手计算上述类数并不太难, 但简便的办法是使用 PARI.

\noindent\textbf{练习 A-3 的解答.}

在本讲义 (及本解答) 中, $i$ 被含混地用来同时表示从 $K_{\mathfrak{m},1}$ 到 $I^{S(\mathfrak{m})}$ 与到 $I^{S\cup S(\mathfrak{m})}$ 的 (显然的) 同态. 这些映射使图表
\[
\begin{tikzcd}
 & & K_{\mathfrak{m},1} \ar[d] \ar[drr] & & \\
\Z S \ar[rr] & & I^{S(\mathfrak{m})} \ar[rr] & & I^{S\cup S(\mathfrak{m})}
\end{tikzcd}
\]
交换 —— 其中 $\Z S$ 是由 $S$ 生成的自由 Abel 群. 该三角形的核--余核正合列给出
\[
\Z S\to I^{S(\mathfrak{m})}/i(K_{\mathfrak{m},1})\to I^{S\cup S(\mathfrak{m})}/i(K_{\mathfrak{m},1})\to 0,
\]
于是
\[
I^{S(\mathfrak{m})}/\langle S\rangle\cdot i(K_{\mathfrak{m},1})\simeq I^{S\cup S(\mathfrak{m})}/i(K_{\mathfrak{m},1}).
\]
因此, $I^{S\cup S(\mathfrak{m})}$ 中包含 $i(K_{\mathfrak{m},1})$ 的子群与 $I^{S(\mathfrak{m})}$ 中包含 $\langle S\rangle\cdot i(K_{\mathfrak{m},1})$ 的子群在对应 $H\leftrightarrow H'$ 之下自然一一对应. 对这样的一对 $H$, $H'$, 设 $L$ 是 $H'$ 的类域. 我们得到同构
\[
I^{S\cup S(\mathfrak{m})}/H\to I^{S(\mathfrak{m})}/H'\to\Gal(L/K).
\]
第二个映射是 Artin 映射, 由 $\mathfrak{p}\mapsto(\mathfrak{p},L/K)$ 诱导. 当然 $L/K$ 是 Abel 的, 且 $S$ 中的素分裂, 因为它们属于 $H'$ (从而有 $(\mathfrak{p},L/K)=1$). 此外, 不在 $S(\mathfrak{m})$ 中的素不分歧. 最后, 两个映射的复合也是 Artin 映射, 于是不在 $S\cup S(\mathfrak{m})$ 中且分裂的素恰好是 $H$ 中的那些. 这表明 $L/K$ 是 $H$ 的 $S$-类域. 我们已经证明它有 Galois 群 $I^{S\cup S(\mathfrak{m})}/H$, 并且它是唯一的, 因为我们知道哪些素在其中分裂 (除有限个可能的例外).

反之, 任何满足 (a) 的 $L$ 都是某个使 $S(\mathfrak{m})\cap S=\emptyset$ 的模 $\mathfrak{m}$ 以及 $I^{S(\mathfrak{m})}$ 的包含 $\langle S\rangle\cdot i(K_{\mathfrak{m},1})$ 的某个子群 $H'$ 的类域, 由此可知, 它是 $I^{S\cup S(\mathfrak{m})}$ 中与 $H'$ 对应的子群 $H$ 的 $S$-类域.

\noindent\textbf{练习 A-4 的解答.}

唉, id\`ele 类群的每个有限指数子群都是开的这一命题并不成立, 即使对数域也如此. 如同 I \ref{cor:1.5} 中那样, 设 $V$ 是可数无穷个 $\mathbb{F}_p$ 拷贝的乘积 (或任何其他有限域的拷贝的乘积), 并设 $W\subset V$ 由那些除有限多个分量外全为零的元素组成 (于是 $W$ 是 $\mathbb{F}_p$ 拷贝的直和). 赋予乘积拓扑, $V$ 是紧的 (且不可数). 子空间 $W$ 可数, 但在 $V$ 中稠密 —— 事实上, 它满射到 $V$ 的任何有限商. 商 $V/W$ 是 $\mathbb{F}_p$ 上的向量空间, 因而有基 $B$ (这需要选择公理 —— 见 Jacobson, \emph{Lectures in Abstract Algebra}, Vol II, Chap. IX; 说 $B$ 是基, 意指 $B$ 的每个有限子集都线性无关, 且 $V/W$ 的每个元素都是 $B$ 中元素的有限线性组合). 设 $S$ 是 $B$ 的非空有限子集, 并设 $W'$ 是 $\mathrm{span}(B\setminus S)$ 在 $V$ 中的原像. 则 $V\supset W'\supset W$, 且 $W'$ 是 $V$ 中指数有限的真子群, 其闭包为 $V$. 它不可能是开的, 因为否则它将是闭的.

回忆 (本讲义 \ref{lem:v.5.9}) $\mathbf{C}_{\Q}\simeq\R_{>0}\times\prod\Z_p^{\times}$ (拓扑群同构). 每个 $\Z_p^{\times}$ 有指数 $2$ 的开子群. 取相应商映射的乘积, 我们得到连续同态 $\prod_p\Z_p^{\times}\to V$, 其中 $V$ 是以素数为指标集的 $\mathbb{F}_2$ 拷贝的乘积. 设 $W'$ 是 $V$ 中不闭 (从而不开) 的有限指数子群. 它在 $\prod\Z_p^{\times}$ 中的原像仍不闭且指数有限, 它在 $\mathbf{C}_{\Q}$ 中的原像亦然.

类似的论证表明 $\Gal(\Q^{\mathrm{al}}/\Q)$ 有不开的有限指数子群 —— 考虑其商群 $\Gal(K/\Q)$, 其中 $K=\Q[\sqrt{2},\sqrt{3},\sqrt{5},\ldots]$ (FT 7.26).

注意, 在每种情形我们都用了选择公理来构造不开的有限指数子群. 能否避免选择公理并找到显式的例子?

\noindent\textbf{练习 A-5 的解答.}

对在左 $G$-作用之下的对象之间的映射 $\alpha$, 定义 $\sigma\alpha\overset{\mathrm{def}}{=}\sigma\circ\alpha\circ\sigma^{-1}$ (这是标准的). 显然 $\sigma(\alpha\circ\beta)=\sigma\alpha\circ\sigma\beta$.

设 $V$ 是 $K$ 上的向量空间. 若线性映射 $\alpha\colon V\otimes L\to V\otimes L$ 相对于 $V$ 的某组基有矩阵 $(a_{ij})$, 则 $\sigma\alpha$ 有矩阵 $(\sigma a_{ij})$. 因此, $\alpha$ ``在 $K$ 上定义'', 即形如 $\alpha_0\otimes1$, 当且仅当 $\sigma\alpha=\alpha$ 对所有 $\sigma\in G$ 成立. 对二次型也有类似的结论.

在 $O(\Phi)$ 中, $\cdot$ 定义为 $\circ$.

现在设 $\alpha\colon(V,\Phi)\otimes L\to(V',\Phi')\otimes L$ 是二次空间的同构, 如练习中所述. 令 $\varphi(\sigma)=\alpha^{-1}\circ\sigma\alpha$. 则
\[
\varphi(\sigma)\cdot\sigma\varphi(\tau)=\alpha^{-1}\circ\sigma\alpha\circ\sigma(\alpha^{-1}\circ\tau\alpha)=\alpha^{-1}\circ\sigma\alpha\circ\sigma\alpha^{-1}\circ\sigma\tau\alpha=\varphi(\sigma\tau).
\]
因此 $\varphi$ 是交叉同态. 任何其他同构 $\alpha'\colon(V,\Phi)\otimes L\to(V',\Phi')\otimes L$ 与 $\alpha$ 相差 $(V,\Phi)\otimes L$ 的一个自同构: $\alpha'=\alpha\circ\beta$, $\beta\in O(\Phi)$. 由 $\alpha'$ 定义的交叉同态是 $\varphi'$, 满足
\[
\varphi'(\sigma)=\beta^{-1}\circ\alpha^{-1}\circ\sigma(\alpha\circ\beta)=\beta^{-1}\circ\varphi(\sigma)\circ\sigma\beta,
\]
这表明 $\varphi'$ 与 $\varphi$ 等价.

于是每个 $(V',\Phi')$ 定义了 $H^{1}(G,O(\Phi))$ 的一个元素 $[(V',\Phi')]$, 且不难验证: 若 $(V',\Phi')\simeq(V'',\Phi'')$, 则 $[(V',\Phi')]=[(V'',\Phi'')]$. 反之, 若 $[(V',\Phi')]=[(V'',\Phi'')]$, 则可以选取同构 $\alpha\colon(V,\Phi)\otimes L\to(V',\Phi')\otimes L$ 与 $\beta\colon(V,\Phi)\otimes L\to(V'',\Phi'')\otimes L$ 使它们给出同一个交叉同态 —— 因而 $\alpha^{-1}\circ\sigma\alpha=\beta^{-1}\circ\sigma\beta$ 对所有 $\sigma\in G$ 成立. 此时 $\beta\circ\alpha^{-1}\colon(V',\Phi')\to(V'',\Phi'')$ 是满足 $\sigma(\beta\circ\alpha^{-1})=\beta\circ\alpha^{-1}$ 对所有 $\sigma$ 成立的同构, 从而在 $K$ 上定义.

剩下要证明 $H^{1}(G,O(\Phi))$ 的每个元素都来自某个 $(V',\Phi')$. 设 $\varphi$ 是交叉同态 $G\to O(\Phi)$. 由于 $H^{1}(G,\mathrm{GL}_n)=0$, 有 $\varphi(\sigma)=\alpha^{-1}\circ\sigma\alpha$, 其中 $\alpha$ 是 $V\otimes L$ 的某个自同构. 设 $\Phi'$ 是 $V\otimes L$ 上满足 $\Phi'\circ\alpha=\Phi$ 的二次型. 我断言 $\Phi'$ 在 $K$ 上定义: 事实上,
\[
\sigma\Phi'=\sigma(\Phi\circ\alpha^{-1})=\Phi\circ\sigma\alpha^{-1}=\Phi\circ\varphi(\sigma)^{-1}\circ\alpha^{-1}=\Phi\circ\alpha^{-1}=\Phi'.
\]
二次型 $\Phi'$ 的选取使得 $\alpha$ 是同构 $(V,\Phi)\to(V,\Phi')$, 由此可知 $(V,\Phi')$ 由 $\varphi$ 表示.

对最后的陈述, 注意二次空间 $(\R^{2},x^{2}+y^{2})$ 不同构于二次空间 $(\R^{2},x^{2}-y^{2})$ (前者只取正值), 但在 $\CC$ 上变得与它同构.

\noindent\textbf{练习 A-6 的解答.}

(a) 按定义, $x=H\cdot s(x)$ 且 $xg=H\cdot s(xg)$. 把第一个等式两边乘以 $g$, 得 $xg=H\cdot s(x)g$. 于是 $H\cdot s(x)g=H\cdot s(xg)$, 从而 $s(x)\cdot g\cdot s(xg)^{-1}\in H$.

(b) 设 $(s^{0}(x))$ 是第二组右陪集代表元, 即 $s^{0}(x)=h(x)s(x)$. 则
\[
\prod_x s^{0}(x)\cdot g\cdot s^{0}(xg)^{-1}
=\prod_x h(x)\cdot s(x)\cdot g\cdot s(xg)^{-1}\cdot h(xg)^{-1}
\equiv\prod_x s(x)\cdot g\cdot s(xg)^{-1}\cdot\prod_x h(x)\cdot\prod_x h(xg)^{-1}\ \bmod\ H^{c}.
\]
由于当 $x$ 取遍 $H$ 在 $G$ 中的右陪集时, $xg$ 亦然, 故 $\prod_x h(x)=\prod_x h(xg)$.

对 $g,g^{0}\in G$,
\begin{multline*}
\prod_x s(x)\cdot gg^{0}\cdot s(xgg^{0})^{-1}
\equiv\prod_x s(x)\cdot g\cdot s(xg)^{-1}\cdot\prod_x s(xg)\cdot g^{0}\cdot s(xgg^{0})^{-1}\ \bmod\ H^{c}\\
=\prod_x s(x)\cdot g\cdot s(xg)^{-1}\cdot\prod_x s(x)\cdot g^{0}\cdot s(xg^{0})^{-1},
\end{multline*}
—— 我们再次利用了当 $x$ 取遍 $H$ 的右陪集时, $xg$ 亦然这一事实. 按 $\bmod\ H^{c}$ 读出, 此式成为
\[
V(gg^{0})=V(g)\cdot V(g^{0}).
\]

(c) 设 $(s^{0}(x))_x$ 是第二组右陪集代表元, 即 $s^{0}(x)=h(x)s(x)$. 对 $m\in M$,
\[
\sum_x s^{0}(x)\cdot m=\sum_x s(x)\cdot m+\sum_x(h(x)-1)\cdot s(x)\cdot m\equiv\sum_x s(x)\cdot m\ \bmod\ I_H M.
\]
因此,
\[
m\mapsto\sum_x s(x)\cdot m\colon M\to M/I_H M\qquad(*)
\]
不依赖于 $s$ 的选取.

若 $(s(x))_x$ 是右陪集代表元系, 则对任意固定的 $g\in G$, $x\mapsto s(xg^{-1})g$ 亦然, 于是
\[
\sum_x s(x)\cdot g\, m\equiv\sum_x s(xg)\, m\ \bmod\ I_H M.
\]
因此 $\sum_x s(x)\cdot(g-1)m\equiv 0 \bmod I_H M$, 这表明映射 $(*)$ 经过 $M/I_G M$ 分解.

它显然是交换群的同态.

(d) 如 \cite[p.~129]{Serre1962} 所说: ``Il n'y a plus maintenant qu'\`a calculer''.

设 $g\in G$. 它走捷径在 $I_G/I_H\, I_G$ 中的像是
\[
\sum_x s(x)(g-1)\ \bmod\ I_H\, I_G,
\]
而走远路的像是
\[
\sum_x(s(x)\cdot g\cdot s(xg)^{-1}-1)\ \bmod\ I_H\, I_G.
\]
对每个 $x$, 存在 $h_x\in H$ 使
\[
s(x)\cdot g=h_x\cdot s(xg)
\]
(见解答的第一段). 现在,
\begin{multline*}
\sum_x s(x)(g-1)=\sum_x h_x\cdot s(xg)-\sum_x s(xg)=\sum_x(h_x-1)s(xg)\\
\equiv\sum_x(h_x-1)\ \bmod\ I_H\, I_G=\sum_x(s(x)\cdot g\cdot s(xg)^{-1}-1).
\end{multline*}
或者, 注意到\footnote{更妙的做法 (Darij Grinberg): 注意 $\sum_x s(x)(g-1)=\sum_x s(x)g-\sum_x s(x)=\sum_x s(x)g-\sum_x s(xg)$; 且 $(s(x)\cdot g\cdot s(xg)^{-1}-1)-s(x)\cdot g+s(xg)=(s(x)\cdot g\cdot s(xg)^{-1}-1)(1-s(xg))$.}
\[
\sum_x s(x)(g-1)=\sum_x s(x)g-\sum_x s(x)=\sum_x s(x)g-\sum_x s(xg),
\]
以及
\[
s(x)\cdot g\cdot s(xg)^{-1}-1-s(x)\cdot g+s(xg)=(s(x)\cdot g\cdot s(xg)^{-1}-1)(1-s(xg)),
\]
后者显然属于 $I_H\cdot I_G$.

\noindent\textbf{练习 A-7 的解答.}

假设存在这样的 $y=(y_n)$. 按提示, 写 $y_n=x_n+z_n$. 由于 $\varphi_n=\sigma x_n-x_n$, 故 $\sigma z_n=z_n$ 对所有 $\sigma$ 与 $n$ 成立. 由假设, $\pi_{n-1,n}$ 把 $y_n$ 映到 $y_{n-1}$, 而经计算, 它把 $x_n$ 映到 $x_{n-1}+\pi_{n-1,n}(x_{n,n})$. 因此, 把各个 $\pi$ 逐次应用到定义 $y_n$ 的等式上, 我们得到 (对 $\pi$ 采用显然的记号)
\[
\begin{aligned}
y_{n-1}&=x_{n-1}+\pi_{n-1,n}(x_{n,n})+\pi_{n-1,n}(z_n)\\
y_{n-2}&=x_{n-2}+\pi_{n-2,n}(x_{n,n})+\pi_{n-2,n-1}(x_{n-1,n-1})+\pi_{n-2,n}(z_n)\\
&\;\;\vdots\\
y_1&=x_1+\pi_{1,n}(x_{n,n})+\cdots+\pi_{1,2}(x_{2,2})+\pi_{1,n}(z_n).
\end{aligned}
\]
把它与定义 $z_1$ 的等式比较, 我们发现\footnote{Darij Grinberg 建议在下面两个等式中删去 $x_{1,1}$ 与 $c_1$.}
\[
z_1=x_{1,1}+\pi_{1,2}(x_{2,2})+\cdots+\pi_{1,n}(x_{n,n})+\pi_{1,n}(z_n).
\]
对它应用 $\sigma_1$, 得
\[
\sigma_1(z_1)=c_1+pc_2+\cdots+p^{n-1}c_n+\sigma_1(\pi_{1,n}(z_n)).
\]
现在, $\sigma_1(\pi_{1,n}(z_n))=p^{n-1}\sigma_n(z_n)$, 于是令 $n\to\infty$, 我们发现 $\sigma_1(z_1)=c\notin\Z$, 矛盾.

我们可以如下构造满足条件的系统 $(M_n,\pi_{n-1,n},\sigma_n)$. 设 $G$ 是一个有商 $G_n\to G_{n-1}$ 的 profinite 群, 其中 $G_n$ 的阶为 $p^{n}$ 且 $G_n\to G_{n-1}$ 为满射 (例如 $G=\Z_p$, $G_n=\Z/p^n\Z$), 并令 $M_n=\Z[G_n]$. 映射 $\pi_{n-1,n}$ 由 $G_n\to G_{n-1}$ 诱导. 由于 $M_n^{G}=\Z\cdot(\sum_{\sigma\in G_n}\sigma)$, 我们可以把 $\sigma_n$ 定义为把 $\sum_{\sigma\in G_n}\sigma$ 映到 $1$ 的映射.

\noindent\textbf{练习 A-8 的解答.}

(a) 设 $A$ 是 $F$ 上的中心单代数. 若 $A$ 是矩阵代数, 则 $A\simeq H(1,1)$ (见下面的 (c)). 否则 (Wedderburn 定理) 它是可除代数, 我们现在假设这一点. 设 $i\in A$, $i\notin F$. 则 $F[i]$ 是 $F$ 的二次域扩张. 配方之后, 可设 $i^{2}=a\in F$. 根据 Noether--Skolem 定理, 存在 $j\in A$ 使 $jij^{-1}=\bar{i}$. 现在 $j^{2}$ 与 $F[i]$ 交换, 故 $j^{2}\in F[i]$. 由于 $j^{2}$ 又与 $j$ 交换, 它属于 $F$, 记 $j^{2}=b\in F$. 剩下要验证 $1,i,j,ij$ 线性无关. 设
\[
c+di+ej+fij=0,\qquad c,d,e,f\in F.
\]
用 $i$ 共轭, 得
\[
c+di-ej-fij=0,
\]
于是 $c+di=0$, 从而 $c=0=d$. 再由 $0=ej+fij=(e+fi)j$ 得 $e=0=f$. 这就证明了 $A\simeq H(a,b)$.

(b) 设 $A=H(a,b)$. 若对某组不全为零的 $w,x,y,z\in F$ 有 $w^{2}-ax^{2}-by^{2}+abz^{2}=0$, 则 $\alpha\overset{\mathrm{def}}{=}w+xi+yj+zk$ 是非零零因子, 从而 $A$ 不是可除代数. 反之, 假设对所有非零 $\alpha\in A$ 都有 $\alpha\bar{\alpha}\ne0$. 由于 $\alpha\bar{\alpha}\in F\setminus\{0\}$, 它有逆, 从而 $\alpha$ 可逆, 其逆为 $\bar{\alpha}\cdot(\alpha\bar{\alpha})^{-1}$.

(c) 设 $\alpha=\begin{pmatrix}0&1\\1&0\end{pmatrix}$ 与 $\beta=\begin{pmatrix}1&0\\0&-1\end{pmatrix}$. 直接验证可知 $I,\alpha,\beta,\alpha\beta=\begin{pmatrix}0&-1\\1&0\end{pmatrix}$ 构成 $M_2(F)$ 的一组基, $\alpha^{2}=I$, $\beta^{2}=I$, 且 $\alpha\beta=-\beta\alpha$. 因此 $M_2(F)\simeq H(1,1)$. 或者, 这也可以由 (b) 推出 (利用 Wedderburn 定理).

(d) 映射 $i\mapsto xi$, $j\mapsto yj$ 定义了 $F$-代数同构 $H(a,b)\to H(ax^{2},by^{2})$. (DG 认为映射应为 $i\mapsto i/x$, $j\mapsto j/y$.)

(e) 与 $L$ 作张量积不改变 $H(a,b)$ 的一组基的乘法表.

(f) 设 $A=H(a,b)$. 设 $I$ 是 $A$ 中的真双侧理想. 由于 $F\cap I=\{0\}$, 陪集 $1+I$, $i+I$, $j+I$, $ij+I$ 具有与原元素相同的乘法表, 于是 $A/I$ 是四元数代数 $\simeq H(a,b)$. 特别地, 它在 $F$ 上的次数为 $4$, 这蕴含 $I$ 为零. 因此 $A$ 是单代数. 容易看出 $F$ 是 $A$ 的中心.

(g) 二次型
\[
w^{2}-ax^{2}-(1-a)y^{2}+a(1-a)z^{2}
\]
有非平凡零点 $(t,t,t,0)$, 其中 $t$ 为 $F$ 中任意非零元. 于是可应用 (b).

(h) 二次型
\[
w^{2}-x^{2}-by^{2}+bz^{2}
\]
有非平凡零点 $(t,t,s,s)$, 其中 $s,t\in F$ 不全为零. 二次型
\[
w^{2}-ax^{2}+ay^{2}-a^{2}z^{2}
\]
有非平凡零点 $(at,s,s,t)$, 其中 $s,t\in F$ 不全为零. 于是可应用 (b). (当然, 不假设 Wedderburn 定理而直接证明 (g) 与 (h) 是更有意思的.)

(i) 映射 $1\mapsto1$, $i\mapsto i$, $j\mapsto j$, $k\mapsto-k$ 是同构 $H(a,b)\to H(a,b)^{\mathrm{opp}}$.

(j) 若 $b$ 是 $F$ 中的平方数, 则由 (d) 与 (h) 得 $H(a,b)\simeq M_2(F)$, 且 $a\in\mathrm{Nm}(F[\sqrt{b}])$. 于是假设它不是平方数. 设 $(w,x,y,z)$ 是 $w^{2}-ax^{2}-by^{2}+abz^{2}$ 的非平凡零点. 若 $x^{2}-bz^{2}=0$, 则 $w^{2}-by^{2}=0$, 从而 $w,x,y,z$ 全为零, 矛盾. 因此 $x^{2}-bz^{2}\ne0$, 且
\[
a=\frac{w^{2}-by^{2}}{x^{2}-bz^{2}},
\]
它是来自 $F[\sqrt{b}]$ 的范数 (作为两个范数之商). 反之, 若 $a=w^{2}-by^{2}$, 则 $w^{2}-a\cdot1^{2}-by^{2}+ab\cdot0^{2}=0$, 于是 $(w,1,y,0)$ 是该二次型的非平凡零点.

(k) 用引理 \ref{lem:iv.3.15} 证明 $H(a,b)\otimes_k H(a,c)\approx H(a,bc)$.

\noindent\textbf{练习 A-9 的解答.}

$p^{n}$ 次分圆多项式为
\[
\Phi(X)=\frac{X^{p^{n}}-1}{X^{p^{n-1}}-1}=X^{p^{n-1}(p-1)}+X^{p^{n-1}(p-2)}+\cdots+X^{p^{n-1}}+1.
\]
它的根是本原 $p^{n}$ 次单位根.

令 $F(X)=\Phi(X+1)$. 则 $F(0)=\Phi(1)=p$, 且
\[
F(X)=\frac{(X+1)^{p^{n}}-1}{(X+1)^{p^{n-1}}-1}\equiv\frac{X^{p^{n}}}{X^{p^{n-1}}}=X^{p^{n-1}(p-1)} \bmod p.
\]
因此 $F(X)$ 是 Eisenstein 多项式, 从而在 $\Q_p$ 上不可约. 它的分裂域是 $L_n$. 显然 $L_n=\Q_p[\zeta_{p^n}-1]$, 且 $F(X)$ 是 $\zeta_{p^n}-1$ 的极小多项式. 因此
\[
\mathrm{Nm}_{L_n/\Q_p}(\zeta_{p^n}-1)=(-1)^{\deg F}F(0)=p
\]
(这里利用了 $\deg F=p^{n-1}(p-1)$ 为偶数 —— 除平凡情形 $n=1$, $p=2$ 外均如此).

第二个陈述是显然的.

\noindent\textbf{练习 A-10 的解答.}

我们必须证明: 对每个 $d\mid n$, 在 $L$ 中不分歧且 $(\mathfrak{p},L/K)$ 的阶为 $d$ 的 $K$ 的素 $\mathfrak{p}$ 之集有极密度 $\varphi(d)/n$. 我们对 $d$ 作归纳.

若 $d=1$, 该集由在 $L$ 中完全分裂的素组成, 而我们知道这样的素之集的密度为 $1/n=\varphi(1)/n$.

设 $H$ 是 $G$ 中阶整除 $d$ 的元素组成的子群 —— 它是 $d$ 阶循环子群. 元素 $\sigma\in G=\Gal(L/K)$ 的阶整除 $d$ 当且仅当 $\sigma|_{L^{H}}=\mathrm{id}$. 因此, $(\mathfrak{p},L/K)$ 的阶整除 $d$ 当且仅当 $(\mathfrak{p},L/K)|_{L^{H}}=1$. 但 $(\mathfrak{p},L/K)|_{L^{H}}=(\mathfrak{p},L^{H}/K)$, 于是 $(\mathfrak{p},L/K)$ 的阶整除 $d$ 当且仅当 $\mathfrak{p}$ 在 $L^{H}$ 中分裂 —— 这样的素之集的密度为 $d/n$. 由归纳假设, 其中 $(\mathfrak{p},L/K)$ 的阶为 $d'\mid d$, $d'\ne d$ 的那些素具有密度 $\varphi(d')/n$. 于是, 使 $(\mathfrak{p},L/K)$ 的阶恰好为 $d$ 的素 $\mathfrak{p}$ 之集的密度为
\[
\frac{d}{n}-\sum_{\substack{d'\mid d\\ d'\ne d}}\frac{\varphi(d')}{n}
=\frac{d-\sum\varphi(d')}{n}.
\]
但 $\sum_{d'\mid d}\varphi(d')=d$ (若 $\zeta$ 是本原 $d$ 次单位根, 则对 $d$ 的每个因子 $d'$, 恰有一个幂 $\zeta^{m}$ 是本原 $d'$ 次单位根), 这就完成了证明.

\noindent\textbf{练习 A-11 的解答.}

注意 $\zeta_8=\frac{1+i}{\sqrt{2}}$, 于是 $\Q[\zeta_8]\subset\Q[i,\sqrt{2}]$. 由于二者在 $\Q$ 上的次数相同, 必然相等.

$X^{2}-1$, $X^{2}+2$, $X^{2}-2$ 的判别式都只被 $2$ 整除.

设 $p$ 为奇素数. $-1,2,-2$ 中至少有一个是 $\mathbb{F}_p$ 中的平方数 —— 例如, 若 $-1$ 与 $2$ 都不是平方数, 则它们的乘积 $-2$ 必是 (平方数构成 $\mathbb{F}_p^{\times}$ 中指数为 $2$ 的子群). 因此这些多项式中至少有一个在 $\mathbb{F}_p[X]$ 中分裂, 且由 Hensel 引理也在 $\Q_p[X]$ 中分裂. 这表明 $-1,2,-2$ 中至少有一个是 $\Q_p$ 中的平方数. 由于 $16=(1+i)^{8}=(\sqrt{2})^{8}=(\sqrt{-2})^{8}$, 可知 $16$ 是 $\Q_p$ 中的 $8$ 次幂.

它也是 $\R$ 中的 $8$ 次幂.

\cleardoublepage\phantomsection
\section*{C 附录: 类域论历史的文献来源}\label{sec:VIII.C}
\addcontentsline{toc}{chapter}{C 附录: 类域论历史的文献来源}

\medskip\noindent Artin, Emil; Nesbitt, Cecil J.; Thrall, Robert M.: \emph{Rings with Minimum Condition}. University of Michigan Publications in Mathematics, no.~1. University of Michigan Press, Ann Arbor, Mich., 1944. x+123~pp. MR0010543.\par\medskip

\noindent Artin, Emil; Tate, John: \emph{Class field theory}. AMS Chelsea Publishing, Providence, RI, 2009. viii+194~pp.\par\medskip

\noindent Brauer, Richard: Emil Artin. \emph{Bull. Amer. Math. Soc.} \textbf{73} (1967), 27--43.\par\medskip

\noindent Conrad, Keith: History of class field theory (见作者网页上的在线讲义).\par\medskip

\noindent Fenster, Della D.; Schwermer, Joachim: A delicate collaboration: Adrian Albert and Helmut Hasse and the principal theorem in division algebras in the early 1930's. \emph{Arch. Hist. Exact Sci.} \textbf{59} (2005), no.~4, 349--379.\par\medskip

\noindent Fenster, D.~D.: 2007. Artin in America (1937--1958): A time of transition. In: Reich, K., Kreuzer, A. (Eds.), \emph{Emil Artin (1898--1962): Beitr\"age zu Leben, Werk und Pers\"onlichkeit}. Dr.~Erwin Rauner Verlag, Augsburg, pp.~99--118.\par\medskip

\noindent Frei, G\"unther: The reciprocity law from Euler to Eisenstein. \emph{The intersection of history and mathematics}, 67--90, Sci.~Networks Hist.~Stud., \textbf{15}, Birkh\"auser, Basel, 1994. MR1308080.\par\medskip

\noindent Frei, G\"unther: On the history of the Artin reciprocity law in abelian extensions of algebraic number fields: how Artin was led to his reciprocity law. \emph{The legacy of Niels Henrik Abel}, 267--294, Springer, Berlin, 2004.\par\medskip

\noindent Hasse, Helmut: History of class field theory. 1967 \emph{Algebraic Number Theory} (Proc.~Instructional Conf., Brighton, 1965) pp.~266--279 Thompson, Washington, D.C. MR0218330.\par\medskip

\noindent Herbrand, J.: \emph{Le D\'eveloppment Moderne de la Th\'eorie des Corps Alg\'ebriques: Corps de Classes et Lois de R\'eciprocit\'e}. Gauthier-Villars, Paris, 1936.\par\medskip

\noindent Iyanaga, Sh\=okichi: \emph{Collected papers}. With a survey of Iyanaga's work on number theory by Ichir\=o Satake. Iwanami Shoten, Tokyo, 1994. xx+363~pp. MR1272246.\par\medskip

\noindent Iyanaga, Shokichi: Travaux de Claude Chevalley sur la th\'eorie du corps de classes: introduction. \emph{Jpn. J. Math.} \textbf{1} (2006), no.~1, 25--85. MR2261061.\par\medskip

\noindent Jacobson, Nathan: Abraham Adrian Albert (1905--1972). \emph{Bull. Amer. Math. Soc.} \textbf{80} (1974), 1075--1100. MR0342341.\par\medskip

\noindent Kisilevsky, H.: Olga Taussky-Todd's work in class field theory. \emph{Olga Taussky-Todd: in memoriam}. \emph{Pacific J. Math.} 1997, Special Issue, 219--224.\par\medskip

\noindent Lemmermeyer, Franz: \emph{Reciprocity laws. From Euler to Eisenstein}. Springer Monographs in Mathematics. Springer-Verlag, Berlin, 2000. xx+487~pp.\par\medskip

\noindent Lewis, David W.: Quaternion algebras and the algebraic legacy of Hamilton's quaternions. \emph{Irish Math. Soc. Bull.} No.~57 (2006), 41--64.\par\medskip

\noindent Lorenz, Falko; Roquette, Peter: The theorem of Grunwald--Wang in the setting of valuation theory. \emph{Valuation theory and its applications}, Vol.~II (Saskatoon, SK, 1999), 175--212, Fields Inst.~Commun., \textbf{33}, Amer.~Math.~Soc., Providence, RI, 2003.\par\medskip

\noindent Mac Lane, Saunders: Origins of the cohomology of groups. \emph{Enseign. Math.} (2) \textbf{24} (1978), no.~1--2, 1--29. MR0497280.\par\medskip

\noindent Milne, J.~S.: The Work of John Tate. In Helge Holden and Ragni Piene, editors, \emph{The Abel Prize 2008--2012}, pages 259--340. Springer, Heidelberg, 2014.\par\medskip

\noindent Miyake, Katsuya: The establishment of the Takagi--Artin class field theory. \emph{The intersection of history and mathematics}, 109--128, Sci.~Networks Hist.~Stud., \textbf{15}, Birkh\"auser, Basel, 1994. MR1308082.\par\medskip

\noindent \emph{Class field theory---its centenary and prospect}. Papers from the 7th International Research Institute of the Mathematical Society of Japan (MSJ) held in Tokyo, June 3--12, 1998. Edited by Katsuya Miyake. Advanced Studies in Pure Mathematics, \textbf{30}. Mathematical Society of Japan, Tokyo, 2001. xii+631~pp.\par\medskip

\noindent Miyake, Katsuya: Teiji Takagi, founder of the Japanese School of Modern Mathematics. \emph{Jpn. J. Math.} \textbf{2} (2007), no.~1, 151--164.\par\medskip

\noindent Parshall, Karen Hunger: Joseph H.~M. Wedderburn and the structure theory of algebras. \emph{Arch. Hist. Exact Sci.} \textbf{32} (1985), no.~3--4, 223--349.\par\medskip

\noindent Petri, Birgit; Schappacher, Norbert: From Abel to Kronecker: episodes from 19th century algebra. \emph{The legacy of Niels Henrik Abel}, 227--266, Springer, Berlin, 2004.\par\medskip

\noindent Roquette, Peter: Class field theory in characteristic $p$, its origin and development. \emph{Class field theory---its centenary and prospect} (Tokyo, 1998), 549--631, Adv.~Stud.~Pure Math., \textbf{30}, Math.~Soc.~Japan, Tokyo, 2001.\par\medskip

\noindent Roquette, Peter: \emph{The Brauer--Hasse--Noether theorem in historical perspective}. Schriften der Mathematisch-Naturwissenschaftlichen Klasse der Heidelberger Akademie der Wissenschaften [Publications of the Mathematics and Natural Sciences Section of Heidelberg Academy of Sciences], \textbf{15}. Springer-Verlag, Berlin, 2005. vi+92~pp. ISBN: 3-540-23005-X.\par\medskip

\noindent Schappacher, Norbert: On the history of Hilbert's twelfth problem: a comedy of errors. \emph{Mat\'eriaux pour l'histoire des math\'ematiques au XX\textsuperscript{e} si\`ecle} (Nice, 1996), 243--273, S\'emin.~Congr., \textbf{3}, Soc.~Math.~France, Paris, 1998.\par\medskip

\noindent Snyder, Noah: Artin's $L$-functions: A Historical Approach. Undergrad thesis (B.~Gross). Unpublished.\par\medskip

\noindent Stevenhagen, P.; Lenstra, H.~W.: Chebotar\'ev and his density theorem. \emph{Math. Intelligencer} \textbf{18} (1996), no.~2, 26--37. MR1395088.\par\medskip

\noindent Stevenhagen, Peter: Hilbert's 12th problem, complex multiplication and Shimura reciprocity. \emph{Class field theory---its centenary and prospect} (Tokyo, 1998), 161--176, Adv.~Stud.~Pure Math., \textbf{30}, Math.~Soc.~Japan, Tokyo, 2001.\par\medskip

\noindent Tate, J.: Problem 9: The general reciprocity law. \emph{Mathematical developments arising from Hilbert problems} (Proc.~Sympos.~Pure Math., Northern Illinois Univ., De Kalb, Ill., 1974), pp.~311--322. Proc.~Sympos.~Pure Math., Vol.~XXVIII, Amer.~Math.~Soc., Providence, R.I., 1976. MR0429839.\par\medskip

\noindent Tate, John: Number theory in the 20th century: Part 1. \emph{Bull. Amer. Math. Soc.} (N.S.) \textbf{54} (2017), no.~4, 547--550. MR3683622.
 % VIII 补遗

% 索引
\backmatter
% 正文引用已全部改为 \cite 交叉引用 (authoryear 样式, 印出 "作者 年份" 与原书正文
% 体例一致, 并带超链接; backref 会在书目条目末尾显示 "引用于 p. X").
% \nocite{*} 使书目中未被正文引用的条目也打印, 以与原书完整书目一致.
\nocite{*}
\printbibliography[heading=bibintoc]
\printindex

\end{document}
